\documentclass[10pt,twoside,openright]{book}
\input{glyphtounicode}
\pdfgentounicode=1
\usepackage[utf8]{inputenc}
\DeclareUnicodeCharacter{2081}{\textsubscript{1}}
\DeclareUnicodeCharacter{2082}{\textsubscript{2}}
\DeclareUnicodeCharacter{03B9}{\ensuremath{\iota}}
\usepackage[T1]{fontenc}
\usepackage{emptypage}
\usepackage[
  paperwidth=6.14in,
  paperheight=9.21in,
  textwidth=4.55in,
  top=0.82in,
  bottom=0.92in,
  inner=0.85in,
  headheight=14pt,
  headsep=0.22in,
  footskip=0.36in
]{geometry}

\usepackage{amsmath,amsthm,mathtools}
\providecommand{\llbracket}{\mathopen{[\![}}
\providecommand{\rrbracket}{\mathclose{]\!]}}
\IfFileExists{newpxtext.sty}{%
  \usepackage{newpxtext}%
  \usepackage{newpxmath}%
}{%
  \usepackage{amssymb}%
  \usepackage[sc]{mathpazo}%
  \usepackage{tgpagella}%
}
\IfFileExists{zi4.sty}{%
  \usepackage[scaled=0.90]{zi4}%
}{%
  \usepackage[scale=0.84]{tgcursor}%
}
\linespread{1.07}
\usepackage{microtype}
\usepackage{booktabs,tabularx,xltabular,array}
\newcolumntype{Y}{>{\raggedright\arraybackslash\hspace{0pt}}X}
\newcolumntype{P}[1]{>{\raggedright\arraybackslash\hspace{0pt}}p{#1}}
\usepackage{enumitem}
\usepackage{xcolor}
\definecolor{CRMBlue}{HTML}{1B365D}
\usepackage{etoolbox}
\usepackage{tikz}
\usepackage{tikz-cd}
\usetikzlibrary{arrows.meta,positioning,calc,fit,backgrounds,shapes.geometric}
\usepackage[font=footnotesize,labelfont={bf,color=CRMBlue},labelsep=period]{caption}
\usepackage{titlesec}
\usepackage{fancyhdr}
\usepackage{mdframed}
\usepackage[noautomatic]{imakeidx}
\makeindex[columns=2,title=Subject Index,intoc]
\usepackage[colorlinks=true,linkcolor=CRMBlue,citecolor=CRMBlue,urlcolor=CRMBlue]{hyperref}
\usepackage{bookmark}
\pdfstringdefDisableCommands{\def\delta{delta}\def\infty{infinity}\def\mathbb#1{#1}\def\mathcal#1{#1}\def\mathbf#1{#1}\def\mathsf#1{#1}\def\mathrm#1{#1}\def\text#1{#1}\def\textsubscript#1{_#1}\def\Eff{Eff}\def\Prism{Prism}}

\setlength{\parskip}{0.42em}
\setlength{\parindent}{1.35em}
\emergencystretch=2.5em
\hbadness=10000
\vbadness=10000

% Part, Chapter, and Section Typography
\titleformat{\part}[display]
  {\thispagestyle{empty}\centering\normalfont\bfseries\color{CRMBlue}}
  {\MakeUppercase{\textls[140]{\small Part \thepart}}}
  {0.9em}
  {\LARGE\color{black}}

\titleformat{\chapter}[display]
  {\normalfont\raggedright}
  {\footnotesize\bfseries\color{CRMBlue}\MakeUppercase{\textls[120]{\chaptername\ \thechapter}}\\[-0.2em]\color{CRMBlue}\rule{\textwidth}{0.8pt}}
  {0.65em}
  {\LARGE\bfseries\color{black}}
\titlespacing*{\chapter}{0pt}{-10pt}{18pt}

\titleformat{name=\chapter,numberless}[display]
  {\normalfont\raggedright}
  {\color{CRMBlue}\rule{\textwidth}{0.8pt}}
  {0.55em}
  {\LARGE\bfseries\color{black}}
\titlespacing*{name=\chapter,numberless}{0pt}{-10pt}{18pt}

\titleformat{\section}
  {\normalfont\large\bfseries\raggedright\color{black}}
  {\color{CRMBlue}\thesection\hspace{0.65em}}{0pt}{}
\titleformat{\subsection}
  {\normalfont\normalsize\bfseries\raggedright\color{black}}
  {\color{CRMBlue}\thesubsection\hspace{0.65em}}{0pt}{}

% Running Headers and Footers
\pagestyle{fancy}
\fancyhf{}
\fancyhead[LE]{\footnotesize\scshape\textls[50]{\color{CRMBlue}Categorical Reverse Mathematics}}
\fancyhead[RO]{\footnotesize\scshape\textls[40]{\color{black}\nouppercase{\leftmark}}}
\fancyfoot[LE,RO]{\footnotesize\thepage}
\renewcommand{\headrulewidth}{0.4pt}
\makeatletter
\renewcommand{\headrule}{{\color{CRMBlue!45}\hrule\@height\headrulewidth\@width\headwidth\vskip-\headrulewidth}}
\fancypagestyle{plain}{%
  \fancyhf{}%
  \fancyfoot[C]{\footnotesize\thepage}%
  \renewcommand{\headrulewidth}{0pt}%
  \renewcommand{\headrule}{}%
}
\renewcommand{\chaptermark}[1]{\markboth{\if@mainmatter Chapter\ \thechapter:\ #1\else #1\fi}{}}
\makeatother

% Chapter & Section Epigraph Macros
\newcommand{\chapterepigraph}[2]{%
  \vspace{-0.4em}\begin{flushright}\begin{minipage}{0.84\textwidth}%
    \raggedleft\small\itshape #1\par\vspace{0.3em}%
    \normalfont\footnotesize\upshape\color{CRMBlue}---\,#2%
  \end{minipage}\end{flushright}\vspace{1.05em}%
}
\newcommand{\sectionepigraph}[2]{%
  \nopagebreak\vspace{-0.15em}\begin{flushright}\begin{minipage}{0.82\textwidth}%
    \raggedleft\footnotesize\itshape #1\par\vspace{0.2em}%
    \normalfont\scriptsize\upshape\color{CRMBlue}---\,#2%
  \end{minipage}\end{flushright}\vspace{0.55em}%
}
% Left-Ruled Theorems/Definitions & Diamond Remarks
\newtheoremstyle{crmplain}{0pt}{0pt}{\itshape}{}{\bfseries\color{CRMBlue}}{.}{0.5em}{}
\newtheoremstyle{crmdef}{0pt}{0pt}{\normalfont}{}{\bfseries\color{CRMBlue}}{.}{0.5em}{}
\newtheoremstyle{crmremark}{0.55\baselineskip}{0.55\baselineskip}{\normalfont}{}{\bfseries\color{CRMBlue}}{.}{0.5em}{}

\mdfdefinestyle{crmleftbar}{%
  leftline=true,rightline=false,topline=false,bottomline=false,%
  linecolor=CRMBlue,linewidth=1.5pt,%
  innerleftmargin=8.5pt,innerrightmargin=0pt,%
  innertopmargin=1.5pt,innerbottommargin=1.5pt,%
  skipabove=0.55\baselineskip,skipbelow=0.55\baselineskip,%
  nobreak=false%
}

\theoremstyle{crmplain}
\newtheorem{theorem}{Theorem}[section]
\newtheorem{proposition}[theorem]{Proposition}
\newtheorem{lemma}[theorem]{Lemma}
\newtheorem{conjecture}[theorem]{Conjecture}
\newtheorem{principle}[theorem]{Principle}
\newtheorem{problem}{Problem}[chapter]

\theoremstyle{crmdef}
\newtheorem{definition}[theorem]{Definition}

\surroundwithmdframed[style=crmleftbar]{theorem}
\surroundwithmdframed[style=crmleftbar]{proposition}
\surroundwithmdframed[style=crmleftbar]{lemma}
\surroundwithmdframed[style=crmleftbar]{conjecture}
\surroundwithmdframed[style=crmleftbar]{principle}
\surroundwithmdframed[style=crmleftbar]{problem}
\surroundwithmdframed[style=crmleftbar]{definition}

\theoremstyle{crmremark}
\newtheorem{remarkx}[theorem]{Remark}
\newenvironment{remark}[1][]{%
  \ifstrempty{#1}{\begin{remarkx}}{\begin{remarkx}[#1]}%
  \pushQED{\qed}\renewcommand{\qedsymbol}{\textcolor{CRMBlue}{$\diamond$}}%
}{%
  \popQED\end{remarkx}%
}

\renewcommand{\qedsymbol}{\textcolor{CRMBlue}{\rule{0.52em}{0.52em}}}
\newcommand{\Eff}{\mathcal{E}\mathrm{ff}}
\newcommand{\Prism}{\mathbf{\Delta}}
\newcommand{\CohesiveQuad}{\Pi \dashv \mathrm{Disc} \dashv \Gamma \dashv \mathrm{Codisc}}
\providecommand{\sslash}{\mathbin{/\mkern-5mu/}}
\newcommand{\lean}[1]{\ignorespaces}
\newcommand{\discussion}[1]{\ignorespaces}
\newcommand{\leanok}{\ignorespaces}
\newcommand{\mathlibok}{\ignorespaces}
\newcommand{\notready}{\ignorespaces}
\ExplSyntaxOn
\NewDocumentCommand{\uses}{m}
 {\clist_map_inline:nn{#1}{\vphantom{\ref{##1}}}\ignorespaces}
\NewDocumentCommand{\proves}{m}
 {\clist_map_inline:nn{#1}{\vphantom{\ref{##1}}}\ignorespaces}
\ExplSyntaxOff

\begin{document}

\frontmatter
\setcounter{secnumdepth}{-1}

% Page i (recto, odd): Half-Title Page
\thispagestyle{empty}
\vspace*{0.28\textheight}
\begin{center}
  {\LARGE\bfseries Categorical Reverse Mathematics with Applications\par}
\end{center}
\cleardoublepage

% Page iii (recto, odd; page ii is blank via emptypage): Title Page
\thispagestyle{empty}
{\centering
  \vspace*{0.14\textheight}
  {\Huge\bfseries Categorical Reverse Mathematics with Applications\par}
  \vspace{1.1em}
  {\color{CRMBlue}\noindent\rule{\textwidth}{0.8pt}}\par
  \vspace{1.1em}
  {\LARGE\itshape\color{CRMBlue} The Structure of Necessity\par}
  \vfill
  {\large James F. Walters\par}
  \vspace{0.35em}
  {\normalsize Port Townsend, Washington\par}
  \vspace{0.6em}
  {\large November 2026\par}
  \vspace*{0.05\textheight}\par}
\clearpage

% Page iv (verso, even): Title-Page Verso Colophon
\thispagestyle{empty}
\vspace*{\fill}
\begingroup
\raggedright
\scriptsize
\setlength{\parskip}{0.10em}
\textbf{Categorical Reverse Mathematics with Applications: The Structure of Necessity}\\
\emph{A Treatise on the Structural Residuums of Foundational Axioms}

\vspace{0.18em}
\noindent By James F.\ Walters\\
Port Townsend, Washington\\
Typeset in \LaTeXe\ via \texttt{pdfLaTeX} in Royal Octavo ($6.14\,\mathrm{in} \times 9.21\,\mathrm{in}$) using the Palatino typeface family, AMS-\LaTeX\ mathematics, and \texttt{TikZ}/\texttt{tikz-cd} vector diagrams.\\
Accompanying Formal Verification: \texttt{HomotopyCRM} (Lean~4 / \texttt{mathlib4}).

\vspace{0.18em}
\noindent\textbf{Cataloging Metadata (Authorial Open-Access Edition)}\\
Author: Walters, James F.\\
Title: Categorical Reverse Mathematics with Applications: The Structure of Necessity / James F.\ Walters.\\
Description: First edition. $|$ Port Townsend, Washington, November 2026. $|$ Includes bibliographical references, notation glossary, and two-column Subject Index (\texttt{imakeidx}/\texttt{makeindex}).\\
Subjects: Topos theory; Reverse mathematics; Category theory; Homotopy type theory; Constructive mathematics; Sheaf theory; Mathematical logic.

\vspace{0.18em}
\noindent\textbf{Mathematics Subject Classification (MSC 2020):}\\
\emph{Primary}: 18B25 (Topoi), 03B30 (Foundations of classical theories, reverse mathematics), 03G30 (Categorical logic, topoi), 18N60 ($\infty$-categories).\\
\emph{Secondary}: 18N70 ($(\infty,\infty)$-categories), 18F10 (Grothendieck topologies and sheaf topoi), 55U35 (Abstract homotopy theory), 03D80 (Applications of computability), 03E25 (Axiom of choice and related propositions), 03F35 (Second-order systems and subsystems), 68V15 (Formal proof in Lean~4).

\vspace{0.12em}
\noindent\textbf{Copyright and Open Access}\\
\copyright\ 2026 James F.\ Walters (Port Townsend, Washington).\\
Released under the \emph{Creative Commons Attribution 4.0 International License} (CC~BY~4.0): \url{https://creativecommons.org/licenses/by/4.0/}.\\
Archival DOI (Zenodo): \href{https://doi.org/10.5281/zenodo.24238087}{\texttt{10.5281/zenodo.24238087}} $|$ \href{https://github.com/tryggth/categorical-reverse-mathematics}{GitHub Source \& Blueprint}.

\vspace{0.12em}
\noindent\textbf{Machine-Checked Formalization Ledger}\\
Verified in Lean~4 (\texttt{leanprover/lean4:v4.28.0}) over \texttt{mathlib4} across 19,961 lines of kernel-checked code in eleven production modules in \texttt{HomotopyCRM} (\texttt{InfinityResiduum}, \texttt{DeltaRingPerfection}, \texttt{BorelTailResiduum}, \texttt{OrdinalRankCompleteness}, \texttt{StepProductPhantom}, \texttt{DirectedFraisseSAP}, \texttt{NominalSimplicity}, \texttt{CubicalAmalgamation}, \texttt{NominalICaTT}, \texttt{SetsResiduum}, and \texttt{DoctrinalFibration}; see Chapter~\ref{ch:master-atlas}, \S\ref{subsec:lean4-verification-ledger} for module-by-module counts). All eleven modules pass the 4-Layer Definition Check Gate with zero \texttt{sorry}, zero custom \texttt{axiom}, and zero \texttt{opaque} declarations, audited via 216 \texttt{\#guard\_msgs in \#print axioms} checks using solely Lean's core logical primitives (\texttt{propext}, \texttt{Classical.choice}, \texttt{Quot.sound}).

\vspace{0.12em}
\noindent\textbf{Typographic Colophon}\\
Set in 10pt Royal Octavo ($6.14\,\mathrm{in} \times 9.21\,\mathrm{in}$) Palatino (\texttt{newpxtext}/\texttt{newpxmath} with Inconsolata when installed; \TeX~Gyre Pagella with Pazo mathematics and scaled \TeX~Gyre Cursor otherwise), \texttt{microtype} optical adjustments, AMS-\LaTeX\ mathematics, \texttt{TikZ}/\texttt{tikz-cd} diagrams, and Unicode PDF text extraction via \texttt{glyphtounicode} (\texttt{\textbackslash pdfgentounicode=1}).\\
First Edition: November 2026 $|$ Port Townsend, Washington $|$ Open-Access Digital Release.\par
\endgroup
\cleardoublepage
\tableofcontents
\chapter{Preface: Form, Symmetry, and the Geometry of Missing Axioms}
\markboth{Preface}{Form, Symmetry, and the Geometry of Missing Axioms}
\chapterepigraph{``Mathematics is the art of giving the same name to different things.''}{Henri Poincar\'e, \emph{Science and Method} (1908)}

\noindent\emph{An Invitation for the General Mathematical Reader}
\medskip

Every undergraduate student of mathematics encounters a curious double standard between \emph{existence by fiat} and \emph{existence by construction}.

Consider two familiar surjective maps:

\begin{enumerate}
  \item \textbf{Reduction modulo $2$:} $q : \mathbb{Z} \twoheadrightarrow \mathbb{Z}/2\mathbb{Z}$, sending each integer to its parity. Viewed merely as a function between bare sets, $q$ admits infinitely many right inverses $s : \mathbb{Z}/2\mathbb{Z} \to \mathbb{Z}$; for instance, set $s(0) = 0$ and $s(1) = 1$, so $q(s(x)) = x$. Viewed as a homomorphism of abelian groups, however, \emph{no right inverse exists}: any homomorphism $s$ would have to satisfy $2 \cdot s(1) = s(1 + 1) = s(0) = 0$ inside the torsion-free group $\mathbb{Z}$, forcing $s(1) = 0$ and contradicting $q(s(1)) = 1$.
  \item \textbf{Wrapping the line around the circle:} $p : \mathbb{R} \twoheadrightarrow S^1$, given by $p(t) = (\cos 2\pi t, \sin 2\pi t)$. As a map of bare sets, $p$ again splits---choose for each point on the circle its unique angle $t \in [0, 1)$. Yet every student of introductory topology learns that \emph{no continuous right inverse $s : S^1 \to \mathbb{R}$ can exist}: walking once around $S^1$ forces $s$ to jump discontinuously by an integer, and that integer---the \textbf{winding number} $1 \in \mathbb{Z} \cong \pi_1(S^1)$---is the founding invariant of algebraic topology.
\end{enumerate}

Notice what happened in both examples. When we stripped away structure and retreated to the universe of bare sets ($\mathbf{Set}$) governed by the Axiom of Choice, both maps split effortlessly, and both obstructions vanished. When we insisted on preserving structure---group addition in the first case, continuity in the second---the maps failed to split. Crucially, however, mathematics did not simply halt with a negative verdict (\emph{``no section exists''}). Instead, the failure of the section precipitated a \textbf{new mathematical object} measuring the exact obstruction: the short exact sequence $0 \to 2\mathbb{Z} \to \mathbb{Z} \to \mathbb{Z}/2\mathbb{Z} \to 0$ in algebra, and the winding-number group $\mathbb{Z}$ in topology.

Alain Connes famously emphasized a parallel phenomenon throughout geometry~\cite[Ch.~II]{Connes94}: when one forms the naive set-theoretic quotient of the torus $\mathbb{T}^2$ by a foliation of lines of irrational slope $\theta$, every leaf is dense, so the classical quotient topology collapses the entire space of leaves to an indiscrete, featureless lump indistinguishable from a single point. Nothing is wrong with the foliation; what failed was a \textbf{premature collapse} that threw away \emph{how} points were identified while keeping only \emph{whether} they were identified. Likewise, Saunders Mac Lane taught generations of mathematicians to ask why a finite-dimensional vector space $V$ is canonically isomorphic to its double dual $V^{**}$ yet isomorphic to its single dual $V^*$ only after an arbitrary choice of basis: the ``missing'' data between $V$ and $V^*$ is not a logical defect, but the geometry of inner products and frames.

\textbf{This book is an exploration of what happens across modern mathematics when we stop treating foundational axioms as binary truth switches and start measuring the geometric objects that appear when those switches are turned off.}

\section*{1. From Classical Reverse Mathematics to Categorical Reverse Mathematics}

In the 1970s, Harvey Friedman and Stephen Simpson inverted the traditional direction of mathematical logic~\cite{Friedman75,Simpson09}. Rather than assuming a strong list of set-theoretic axioms and deducing theorems, \textbf{Reverse Mathematics} begins with a theorem $T$ of ordinary undergraduate mathematics---the Bolzano--Weierstrass theorem, the Heine--Borel covering theorem, K\"onig's infinity lemma, the existence of maximal ideals---and asks: \emph{working over a weak computational base theory, which set-existence axiom is logically equivalent to $T$?} Remarkably, hundreds of classical theorems sort themselves into a small handful of equivalence classes.

Yet classical Reverse Mathematics remains \textbf{propositional}: it asks whether an implication $A \implies B$ is \emph{true} or \emph{false} in a model of arithmetic or set theory. Category theory invites a deeper, structural question. Suppose two mathematical constructions $\Phi$ and $\Psi$ are connected by a canonical comparison map
\[
  \kappa : \Phi \longrightarrow \Psi
\]
which becomes an isomorphism (or equivalence) in classical Zermelo--Fraenkel set theory with the Axiom of Choice ($\mathsf{ZFC}$). For example:
\begin{itemize}
  \item $\Phi$ might be the set of global continuous sections of a bundle $E \twoheadrightarrow B$, and $\Psi$ the one-point set witnessing that every fiber $E_b$ is non-empty;
  \item $\Phi$ might be the polynomials or finite trees built from the bottom up in finitely many steps, and $\Psi$ the formal series or infinite trees unfolded from the top down; or
  \item $\Phi$ might be an infinite-dimensional directed system, and $\Psi$ the limit of all its finite-dimensional shadows.
\end{itemize}

When we move out of $\mathbf{Set}_{\mathsf{ZFC}}$ into mathematical universes whose objects carry built-in \textbf{continuity}, \textbf{symmetry}, \textbf{definability}, or \textbf{irreversible direction}, $\kappa : \Phi \to \Psi$ ceases to be an isomorphism. \textbf{Categorical Reverse Mathematics (CRM)} asks two paired questions:

\begin{enumerate}
  \item \textbf{What is the \emph{structural residuum} of $\kappa$?} That is, what exact algebraic, cohomological, or combinatorial invariant $\mathrm{Res}(\kappa)$ lives in the gap between $\Phi$ and $\Psi$?
  \item \textbf{Which foundational principle trivializes $\mathrm{Res}(\kappa)$?} That is, which exact rung of Choice, arithmetic infinity, $p$-torsion freeness, symmetry-breaking, or arrow-inversion holds in our ambient doctrine\footnote{Introduced by F.~W.~Lawvere~\cite{Lawvere69}, a \emph{categorical doctrine} specifies a family of structured categories or mathematical universes equipped with chosen categorical operations---such as finite limits, exponentials, logical quantifiers, Frobenius lifts, or higher cells---together with the functors preserving that structure. Whereas an ordinary \emph{theory} (such as the axioms for groups or Peano arithmetic) classifies models \emph{inside} a fixed ambient logic, a \emph{doctrine} specifies the structural rules of the ambient framework itself (see Section~2 below and Definition~\ref{def:comparison-triple}).} \emph{if and only if} $\mathrm{Res}(\kappa)$ vanishes?
\end{enumerate}

The word \textbf{residuum} carries two complementary meanings here. In algebra (following Emmy Noether's ideal residuals), the residual of $P$ inside $Q$ is the exact missing factor required to upgrade $P$ to $Q$. In chemistry, a residue is the crystalline solid that precipitates out of solution when a solvent boils off. In classical mathematics, discontinuous Choice ($\mathsf{AC}$) and algebraic group-completion act as universal solvents: they dissolve twisted bundles and one-way processes into flat, reversible ones. Boil those solvents away, and the hidden geometry crystallizes as the residuum.

\section*{2. A Pocket Dictionary Before Entering the Text}

Although Parts~I--IV of this monograph engage contemporary tools---categorical doctrines, Grothendieck topoi, sheaf cohomology, $\delta$-rings, Fra\"iss\'e limits, and $(\infty, \infty)$-categories---every one of those concepts grows directly out of undergraduate algebra, analysis, and topology:

\begin{itemize}
  \item \textbf{A Categorical Doctrine (``The Structural Rules of the Ambient Universe''):}\index{Categorical doctrine} In undergraduate algebra, a \emph{theory} (such as the axioms for groups, rings, or vector spaces) describes objects living \emph{inside} a fixed background like $\mathbf{Set}$. Following F.~W.~Lawvere~\cite{Lawvere69}, a \textbf{doctrine} steps up one level to specify the structural capabilities of the \emph{ambient category itself}: which products, quotients, function spaces, logical quantifiers, Frobenius lifts, or higher-dimensional arrows are available, and which functors must preserve them. Categories with finite limits, Heyting categories, elementary topoi, Joyal $\delta$-rings, and weak $(\infty, \infty)$-categories are each distinct doctrines; passing between a broad base doctrine $\mathcal{D}_0$ and a stronger subdoctrine $\mathcal{D}_1 \subseteq \mathcal{D}_0$ is how CRM turns foundational axioms into geometry.
  \item \textbf{A Topos (``Sets with Built-In Geometry or Symmetry''):} Imagine doing ordinary set theory---forming subsets, products, function sets, quotients, and logical quantifiers $\forall, \exists$---except that \emph{every set in the room is required to carry a continuous action of a symmetry group $\mathcal{G}$}, and \emph{every function must commute with $\mathcal{G}$}. Or imagine that every ``set'' is actually a continuously varying family of sets parameterized by open subsets of a topological space $X$. Such a self-contained mathematical universe is a \textbf{Grothendieck topos}. Inside a topos, constructive reasoning works seamlessly, while principles like the Law of Excluded Middle or the Axiom of Choice become \emph{geometric properties of the space or group} rather than universal dogmas.
  \item \textbf{A Torsor (``A Group That Forgot Its Identity Element''):} In multivariable calculus, a Euclidean plane without a marked origin $(0, 0)$ is an \emph{affine space} over $\mathbb{R}^2$: you cannot add two points $P + Q$, because there is no origin, but given any two points $P, Q$ there is a unique displacement vector $Q - P \in \mathbb{R}^2$ carrying $P$ to $Q$. Replace the vector space $\mathbb{R}^2$ by any group $G$, and an inhabited space with a free, transitive $G$-action is a \textbf{$G$-torsor}. Choosing a single point in a $G$-torsor immediately identifies it with $G$; when symmetry or continuity forbids choosing a point globally, the torsor represents a non-zero cohomology class in $H^1$.
  \item \textbf{Initial Algebras ($\mu F$) vs.\ Terminal Coalgebras ($\nu F$):} Every student of algebra knows the difference between the polynomial ring $\mathbb{R}[x]$ (expressions $\sum_{k=0}^N a_k x^k$ built in \emph{finitely} many steps from $0$ and $x$) and the formal power series ring $\mathbb{R}[[x]]$ (infinite sequences $(a_0, a_1, a_2, \dots)$ whose coefficients can be inspected indefinitely one term at a time). Category theory unifies these as the \textbf{initial algebra} $\mu F$ (inductive, bottom-up, finite depth) and \textbf{terminal coalgebra} $\nu F$ (coinductive, top-down, potentially infinite depth) of a step functor $F$.
  \item \textbf{An $(\infty, \infty)$-Category (``One-Way Arrows All the Way Up''):} A directed graph or partially ordered set has objects ($0$-cells) and one-way arrows $f : x \to y$ ($1$-cells). A $2$-category adds $2$-cells $\alpha : f \Rightarrow g$ deforming one arrow into another (like homotopies between paths, or natural transformations between functors). Continuing through all dimensions $n = 0, 1, 2, 3, \dots$ while allowing composition to be associative \emph{up to coherent higher cells} yields a weak \textbf{$(\infty, \infty)$-category}. When every cell above dimension $1$ is invertible, one gets an $(\infty, 1)$-category; when even $1$-cells are invertible, one gets a topological space (an \emph{anima} or $\infty$-groupoid), where arrows are simply reversible paths. Keeping higher cells \textbf{non-invertible} preserves the arrow of time---rewrites, cobordisms, and causal flows that cannot be run backward.
  \item \textbf{Noetherianity and Postnikov Towers:} In 1921, Emmy Noether revolutionized ring theory through the Ascending Chain Condition: every nested chain of ideals $I_1 \subseteq I_2 \subseteq \cdots$ stabilizes at a finite step. Lifted to $(\infty, \infty)$-categories, \textbf{Noetherianity} demands that any tower of ever-higher parallel cells between two arrows eventually becomes contractible at some finite dimension---ruling out wild spaces like the $2$-sphere $S^2$, whose higher homotopy groups $\pi_m(S^2)$ remain non-zero infinitely often. Meanwhile, a \textbf{Postnikov tower} approximates a higher category by its $n$-dimensional truncations $\pi_{\le 1}(\mathcal{C}), \pi_{\le 2}(\mathcal{C}), \dots$---much like approximating a real number by its finite decimal truncations $3,\; 3.1,\; 3.14,\; \dots$ or an integer by its residues modulo $p^n$---and asks whether every compatible sequence of finite shadows comes from an actual object of $\mathcal{C}$.
\end{itemize}

\section*{3. A Guided Tour of the Book}

With that dictionary in hand, the nine chapters of Volume~I of this monograph unfold along a single governing law, which we call \textbf{Base--Fiber Conservation}:

\begin{quote}
\emph{Mathematical complexity is conserved. Whenever a strong foundational axiom (such as discontinuous Choice or algebraic group-completion) flattens an obstruction to zero in the foreground, that complexity has merely been absorbed into unconstructible choices or formal inverses in the background; conversely, switching that axiom off makes the background constructive and forces the conserved complexity back into view as an explicit geometric residuum.}
\end{quote}

\subsection*{Chapters 1 and 2: The Choice--Geometry Atlas and Base--Fiber Conservation}

Chapters~1 and~2 set up the foundational laboratory. In $\mathbf{Set}_{\mathsf{ZFC}}$, the statement \emph{``every surjective map $E \twoheadrightarrow B$ has a right inverse''} is equivalent to the Axiom of Choice. Inside a general topos $\mathcal{E}$, however, ``Choice'' splits into a hierarchy of genuinely distinct principles:
\begin{itemize}
  \item \textbf{External Choice ($\mathsf{EAC}$ / $\mathrm{AC}$):} Every epimorphism $E \twoheadrightarrow B$ (and in particular every globally surjective bundle $E \twoheadrightarrow 1_{\mathcal{E}}$) admits a global section.
  \item \textbf{Internal Choice ($\mathsf{IAC}$):} Every dependent product functor $\Pi_f$ (equivalently every exponential functor $(-)^Z$) preserves epimorphisms---so bundles split locally and in slices even when global sections fail.
  \item \textbf{Countable Dependent Choice ($\mathsf{DC}_\omega$ / $\mathsf{IDC}_\omega$):} Given an infinite sequence of step-by-step choices $X_0 \twoheadleftarrow X_1 \twoheadleftarrow X_2 \twoheadleftarrow \cdots$, the inverse limit projection $\varprojlim_n X_n \twoheadrightarrow X_0$ remains an epimorphism.
\end{itemize}

Chapter~1 also widens the domain of the residuum operator beyond Choice alone to arbitrary \emph{doctrinal comparison triples}, placing arithmetic infinity in $\mathbf{FinSet}$ alongside logic, $p$-torsion, and higher invertibility. Chapter~2 builds the \textbf{Choice--Geometry Residuum Atlas}, showing how the failure of $\mathsf{EAC}$ on group-equivariant bundles manifests precisely as non-trivial torsors in first sheaf cohomology $H^1(\mathcal{E}; G)$, and separating internal two-valuedness from external projectivity of the terminal object---while calibrating pure well-founded $\mathsf{ZF}$ choice separations without urelements via symmetric Boolean-valued submodels $\mathrm{HS}_{\mathcal{F}} \subset V^{\mathbb{B}}$ and $\mathcal{F}$-continuous $\mathcal{G}$-equivariant Boolean sheaves $\mathbf{Sh}_{\mathcal{G}, \mathcal{F}}(\mathbb{B})$.

\subsection*{Chapters 3 and 4: Fixed Points, Fermat Quotients on Torsion Rings, and Infinite Towers}

Chapters~3 and~4 investigate what happens when iterative constructions run infinitely many times.

On the algebraic side, we examine \textbf{$\delta$-rings}\index{delta-ring@$\delta$-ring}---an unexpectedly elementary refinement of Fermat's Little Theorem ($x^p \equiv x \pmod p$) that sits at the heart of modern $p$-adic geometry. On any $\delta$-ring $R$, the associated Frobenius lift satisfies $\phi(x) = x^p + p\,\delta(x)$. Over the integers $\mathbb{Z}$, which is initial among commutative rings, the unique ring endomorphism is the identity map $\phi = \mathrm{id}_{\mathbb{Z}}$ (since $x \mapsto x^p$ is not additive for $p > 1$), forcing a unique, non-zero Fermat quotient operator
\[
  \begin{aligned}
    \delta(x) &\;=\; \frac{\phi(x) - x^p}{p} \;=\; \frac{x - x^p}{p} \;\neq\; 0 \\
    &\qquad \bigl(|x| \ge 2;\; \text{e.g., } \delta(2) = -1 \text{ at } p = 2\bigr).
  \end{aligned}
\]
What happens if the ring has \textbf{$p$-torsion}---say, $R_2 = \mathbb{Z}[\eta]/(2\eta, \eta^2)$, which decomposes as $\mathbb{Z} \oplus \mathbb{F}_2$ as an additive abelian group while remaining strictly distinct from the reduced product ring $\mathbb{Z} \times \mathbb{F}_2$ because $\eta \neq 0$ and $\eta^2 = 0$? Naive division by $p = 2$ collapses, because multiplying by $2$ erases the $\eta$-coordinate. Replacing division with Andr\'e Joyal's~\cite{Joyal85} universal binomial identity $\delta(x + y) = \delta(x) + \delta(y) + \frac{x^p + y^p - (x+y)^p}{p}$ (for $p = 2$, simply $\delta(x+y) = \delta(x) + \delta(y) - xy$) rescues the theory without ever dividing by $2$---and uncovers \emph{two distinct $\delta$-ring structures} ($\delta_0(\eta) = 0$ and $\delta_1(\eta) = \eta$) lifting the exact same \textbf{augmentation Frobenius endomorphism} $\phi(a + b\eta) = a$ (where $\phi(\eta) = 0$). Passing to the ``perfection'' colimit where $\phi$ becomes invertible further reveals a clean \textbf{shift chirality}: the bottom-up initial algebra map shifts forward by $+1$, whereas the induced Frobenius coalgebra map is forced by naturality to shift backward by $-1$, while our torsion-free dual-integer witness $\mathbb{Z}[\varepsilon]/(\varepsilon^2)$ demonstrates how forcing perfection obliterates the Nygaard filtration\index{Nygaard filtration}.

On the topological and descriptive-set-theoretic side, Chapter~4 examines infinite towers of surjections $X_0 \twoheadleftarrow X_1 \twoheadleftarrow X_2 \twoheadleftarrow \cdots$ of Polish (complete separable metric) spaces~\cite[Ch.~I--II]{Kechris95}. With discontinuous Choice, every such tower has a point, so the classical derived limit $\varprojlim^1$ vanishes. When we forbid discontinuous Choice---embedding Polish spaces into Clausen and Scholze's \emph{light condensed sets} $\mathrm{Sh}(\mathrm{Prof}_{\aleph_0})$, sheaves over Cantor-like profinite spaces---a non-zero \textbf{definable $\varprojlim^1_{\mathrm{def}}$ residuum} appears. Remarkably, this same inductive-versus-coinductive gap ($\mu F < \nu F$) models one of the deepest miracles of $4$-dimensional manifold topology: finite handlebody trees ($\mu F$) untangle smoothly to standard $\mathbb{R}^4$, whereas Michael Freedman's infinite Casson towers ($\nu F \setminus \mu F$) collapse \emph{topologically} to standard handles while retaining turbulent end-invariants that support \textbf{exotic smooth structures on $\mathbb{R}^4$}. Indeed, Alain Connes's irrational-slope torus foliation $\mathbb{T}^2/\mathcal{F}_\theta$ from our opening pages reappears in Chapter~\ref{ch:base-fiber} (\S\ref{sec:exotic-r4-e0}) in its profinite hyperfinite twin $\mathcal{Q}_{E_0} = 2^{\mathbb{N}}/E_0 = \prod_{\mathbb{N}} \mathbb{F}_2 / \bigoplus_{\mathbb{N}} \mathbb{F}_2$\index{Connes quotient Q_E0@Connes quotient $\mathcal{Q}_{E_0} = 2^{\mathbb{N}}/E_0$} as the exact Gompf--Panagiotopoulos~\cite{GompfPanagiotopoulos26} obstruction to classifying smooth structures on $\mathbb{R}^4$, whose étale groupoid convolution algebra is Glimm's CAR algebra $\mathrm{UHF}(2^\infty)$~\cite{Glimm60} ($K_0 \cong \mathbb{Z}[\tfrac{1}{2}]$)---completing, just like Rieffel's irrational rotation algebra $A_\theta$, to the hyperfinite $\mathrm{II}_1$ factor $\mathcal{R}$~\cite[Ch.~V]{Connes94}.

\subsection*{Chapter 5: The Four-Tier $(\infty, \infty)$-Ladder, Pure Phantoms, and Sphere Homotopy from Forgetting Direction}

Chapter~5 applies the CRM lens to weak $(\infty, \infty)$-categories, where it produces three of the book's central discoveries.

First, in an infinite-dimensional directed category, when should a $1$-arrow $f : x \to y$ count as an ``equivalence''? Asking for an inverse $g : y \to x$ up to $2$-arrows $g \circ f \Rightarrow \mathrm{id}_x$ and $f \circ g \Rightarrow \mathrm{id}_y$ merely pushes the question up one dimension: those $2$-arrows need $3$-arrow inverses, which need $4$-arrow inverses, and so on. Work by Ozornova--Rovelli--Walde and Gepner--Heine~\cite{ORW26,GepnerHeine26} revealed that inside the unconstrained right-truncation projective limit $\mathrm{Cat}^R_{(\infty,\infty)} := \varprojlim_\kappa \mathrm{Cat}_{(\infty, n)}$, Left Bousfield localization exhibits sharp \textbf{Bousfield contravariance}: enlarging the saturated class of localizing equivalences
\[
  W_{\mathrm{ind}} \;\subsetneq\; W_{\mathrm{coind}} \;\subsetneq\; W_\omega \;\subsetneq\; W_{\mathrm{Post}}
\]
(from inductive equivalences $W_{\mathrm{ind}}$ generated bottom-up from identities, to coinductive bi-invertible equivalences $W_{\mathrm{coind}}$, to Street--Verity $\omega$-equivalences $W_\omega$, up to all finite-dimensional Postnikov equivalences $W_{\mathrm{Post}}$) shrinks the reflective subcategory of local objects in the \emph{opposite} direction, yielding a \textbf{strict 4-tier localization hierarchy}\index{Four-tier localization ladder}:
\[
  \mathrm{Cat}^L_{(\infty,\infty)} \;\subsetneq\; \mathrm{Cat}^\omega_{(\infty,\infty)} \;\subsetneq\; \mathrm{Cat}^{\mathrm{coind}}_{(\infty,\infty)} \;\subsetneq\; \mathrm{Cat}^R_{(\infty,\infty)},
\]
separating (from largest to smallest) the ambient right-truncation limit $\mathrm{Cat}^R_{(\infty,\infty)}$ (built over inductive equivalences $W_{\mathrm{ind}}$, and strictly larger than the Postnikov-complete tier), the coinductively saturated categories $\mathrm{Cat}^{\mathrm{coind}}_{(\infty,\infty)}$, the Street--Verity $\omega$-equivalence tier $\mathrm{Cat}^\omega_{(\infty,\infty)}$, and the Ad\'amek initial-algebra / Postnikov-complete tier $\mathrm{Cat}^L_{(\infty,\infty)}$ (which match their finite-dimensional truncations via $\mathcal{C} \xrightarrow{\sim} RL(\mathcal{C})$).

Asking our standard CRM questions immediately uncovers three hidden phenomena:

\begin{enumerate}
  \item \textbf{The Dependent Choice Cost of Algebraic Witness Streams (Theorem~\ref{thm:coinductive-stream-choice}):} Syntactic type theories for $\omega$-categories (\texttt{ICaTT}) package coinductive invertibility into an explicit infinite stream of algebraic witness terms. Over $\mathbf{Set}_{\mathsf{ZFC}}$, having a relational post-fixed point and having a witness stream are identical. Varying the base topos proves they are separated by \textbf{Countable Dependent Choice ($\mathsf{IDC}_\omega$)}: inside a permutation topos where $\mathsf{DC}_\omega$ fails, we construct an explicit arrow $u : 0 \to 1$ that is relationally invertible at every finite stage yet possesses \emph{zero} symmetric witness streams.
  \item \textbf{Pure Combinatorial Phantoms Without Homotopy Spheres (Theorems~\ref{thm:ordinal-rank-completeness}, \ref{thm:classifier-scc-trapping}--\ref{thm:step-product-counterexample}):} Published examples separating the upper tiers rely on topological configuration spaces ($\mathbb{R}\mathbb{P}^2$, yielding $E^{(\omega)} \in \mathrm{Cat}^{\mathrm{coind}}_{(\infty,\infty)} \setminus \mathrm{Cat}^\omega_{(\infty,\infty)}$)~\cite{HenryLoubaton25,ORW26} or infinite loop hubs ($W \in \mathrm{Cat}^\omega_{(\infty,\infty)} \setminus \mathrm{Cat}^L_{(\infty,\infty)}$)~\cite{ORW26}. Stripping away all topological spheres, we restrict to categories where every mapping object between two points has \emph{finite} dimension ($\mathrm{Cat}_{\le \omega}$). We prove that finite mapping dimensions automatically prevent ``phantom arrows'' (the Postnikov comparison map is always fully faithful), and if $1$-arrows never form directed cycles, \emph{no phantom objects can exist either}. Yet the moment directed cycles are allowed---even in a completely rigid, sphere-free, Noetherian category built from finite posets with an absorbing minimum operation $\min(x, y)$---the restricted-tail step-product $\mathcal{Q} \subset \prod_{k \ge 1} \mathcal{A}_k$\index{Step-product Q@Step-product $\mathcal{Q}$} on the walking arrow $[1] = \{0 < 1\}$ (where each factor $\mathcal{A}_k = \mathcal{M}_{k+1}$ has contractible hom-objects $\mathrm{Hom}(0_k, 0_k) \simeq \mathrm{Hom}(0_k, 1_k) \simeq \mathrm{Hom}(1_k, 0_k) \simeq *$ while $\mathrm{Hom}(1_k, 1_k) = \mathcal{M}_k \not\simeq *$, producing $\mathrm{Hom}_{\mathcal{Q}}(e^{(m)}, e^{(r)}) \simeq \prod_{k=1}^{\min(m, r)} \mathcal{M}_k$ of rank $\min(m, r)$) glues finite-prefix indicator sequences $e^{(0)}, e^{(1)}, e^{(2)}, \dots$ via their $n$-truncations $y_n := \ell_n(e^{(n-1)}) \in \pi_{\le n}(\mathcal{Q})$ across dimensions into a pure Postnikov \textbf{phantom object} $X_\bullet \in RL(\mathcal{Q}) \setminus \mathcal{Q}$ absent from $\mathcal{Q}$ itself, much like $3,\; 3.1,\; 3.14,\dots$ converging to $\pi \notin \mathbb{Q}$.
  \item \textbf{How Forgetting Direction Births the Homotopy Groups of Spheres (Chapter~\ref{ch:higher-crm}, \S\ref{subsec:crm-discovery-engine-inf-inf}):} Consider the category of finite sets and bijections, $\mathbf{FinSet}^{\simeq} = \bigsqcup_{m \ge 0} B\Sigma_m$, under disjoint union $\amalg$. Only the empty set ($m = 0$) has an inverse under $\amalg$: you cannot add a $3$-element set to another set to get $\emptyset$. When we stabilize this structure \emph{directedly} (keeping non-invertibility intact), every single level of the resulting categorical spectrum is finitely truncated and Noetherian---completely devoid of higher homotopy groups above its finite ceiling (cf.\ Gei\ss's Oriented Freudenthal Suspension Theorem~\cite{Geiss26}). What happens when we switch on the collapse functor $\pi_{\le 0}$ that forces all directed arrows to become invertible paths? Inverting non-empty finite sets forces $(\mathbb{N}, +)$ to group-complete into $(\mathbb{Z}, +)$ and triggers Quillen's plus-construction on the infinite symmetric group $B\Sigma_\infty$. By the celebrated \textbf{Barratt--Priddy--Quillen theorem}~\cite{BarrattPriddy72,Quillen71}, the result is the \textbf{classical sphere spectrum $\mathbb{S}_{\mathrm{top}}$}---meaning that the entire wild zoo of stable homotopy groups of spheres $\pi_*^s(\mathbb{S}_{\mathrm{top}})$ emerges purely as the \textbf{collapse residuum} of inverting directed morphisms!
\end{enumerate}

\subsection*{Chapter 6: Fra\"iss\'e Symmetry, the Odd-Parity Torsor, and the $n = 3$ Triangle Wall}

Chapter~6 brings all of these ideas down to concrete combinatorics and permutation groups.

Take the group $\mathrm{Sym}(\mathbb{N})$ of all bijections of $\mathbb{N} = \{0, 1, 2, \dots\}$, and consider sets $X$ equipped with a $\mathrm{Sym}(\mathbb{N})$-action such that each element $x \in X$ depends on only finitely many numbers $F \subset \mathbb{N}$ (any permutation fixing $F$ fixes $x$). This is \textbf{Schanuel's nominal topos $\mathbf{Nom}$}~\cite[Ch.~III, \S 9]{MacLaneMoerdijk92}---the mathematical home of bound variables ($\alpha$-equivalence in calculus and computer science, where $\int_0^1 f(x)\,dx$ and $\int_0^1 f(y)\,dy$ name the same integral).

Now ask an elementary puzzle: \emph{Can you write down a rule, invariant under all permutations of $\mathbb{N}$, that selects a single finite subset $S \subset \mathbb{N}$ of \textbf{odd} size?}
Certainly not: if $S$ is non-empty and finite, pick some $a \in S$ and let $b = \max(S) + 1 \notin S$; the transposition $(a \; b)$ swaps $a$ out for $b$, moving $S$ to a different set $S' \neq S$. Thus the set $\mathcal{T}_{\mathrm{odd}}$ of all odd-sized finite subsets of $\mathbb{N}$ has \textbf{no global invariant element}. Yet notice:
\begin{itemize}
  \item Under symmetric difference $A \triangle B := (A \setminus B) \cup (B \setminus A)$, \emph{all} finite subsets form an $\mathbb{F}_2$-vector space, and the \emph{even}-sized subsets $\mathrm{EvenFinset}$ form a subgroup (the kernel of $|A| \bmod 2$).
  \item Whenever we allow ourselves to name \emph{even a single number} $a \in \mathbb{N}$, $\mathcal{T}_{\mathrm{odd}}$ immediately acquires the canonical element $\{a\}$; and if someone else chooses $\{b\}$, their difference $\{a\} \triangle \{b\} = \{a, b\}$ has size $2$, landing squarely inside $\mathrm{EvenFinset}$!
\end{itemize}

Consequently, $\mathcal{T}_{\mathrm{odd}}$ is a non-trivial $\mathrm{EvenFinset}$-torsor representing a non-zero cohomology class $[\mathcal{T}_{\mathrm{odd}}] \neq 0 \in H^1(\mathbf{Nom}; \mathrm{EvenFinset})$---even though a classic theorem of Schreier and Ulam implies that $\mathrm{Sym}(\mathbb{N})$ has \emph{zero} cohomology with coefficients in any constant discrete group ($H^1(\mathbf{Nom}; \mathbb{Z}/2\mathbb{Z}) = 0$, \textbf{Pro-Discrete Blindness})! Chapter~6 situates $\mathbf{Nom}$ at the base of Takayuki Kihara's \textbf{strict 5-rung Choice Spectrum}~\cite{Kihara26} for Fra\"iss\'e permutation topoi, contrasts finite discrete symmetry groups $G$ (where $\{1\}$ is open) against infinite oligomorphic groups $\mathrm{Aut}(\mathcal{M})$ (where stabilizers shrink endlessly without bottoming out), and proves a sharp \textbf{dimensional amalgamation wall} for directed cell complexes: gluing two acyclic directed complexes along a flow-convex overlap always stays acyclic ($n = 2$), whereas gluing around the boundary of the facet-supported directed triangle $F_3 = \partial\Delta^2$ or a directed $n$-cube fails for \emph{every} $n \ge 3$ by closing a global directed cycle.

\subsection*{Chapter 7: Residuums and ``$\mathbf{Sets}$''}

Up through Chapter~6, the ambient doctrine of ``$\mathbf{Sets}$'' itself stands quietly offstage behind folkloric appeals to $\mathsf{ZFC}$--$\mathsf{NBG}$ conservativity. Chapter~7 turns the residuum operator directly onto ``$\mathbf{Sets}$'':
\begin{itemize}
  \item \textbf{The Baseline Adjunction and Fixed-Set Class Kaleidoscope:} Above Zermelo models $\mathbf{Mod}(\mathsf{Z})$, parameter-definable classes furnish a minimal coreflective left adjoint $\mathrm{Def} \dashv \pi_{\mathrm{Set}}$ into Replacement-free class models $\mathbf{Mod}(\mathsf{GB}_0)$, while Kameryn Williams's~\cite{Williams18} theorem shows that even over a completely frozen countable model $M \models \mathsf{ZFC}$, the vertical fiber $\pi_{\mathrm{Set}}^{-1}(M)$ of $\mathsf{GBC}$-realizations embeds every countable partial order---with Global Choice ($\mathsf{NBG}$) calibrated precisely by splitting the large weak equivalence $\mathbf{Card} \hookrightarrow \mathbf{Set}$.
  \item \textbf{The $\mathsf{ETCS} \simeq \mathsf{Mac}$ Plot Twist, Glazer's Leak, and Borel Determinacy:} Lawvere's Elementary Theory of the Category of Sets ($\mathsf{ETCS}$~\cite{Lawvere64})---the Boolean elementary-topos baseline of Chapters~1--6---is bi-interpretable with Bounded Zermelo / Mac~Lane set theory ($\mathsf{Mac}$~\cite{Cole73,Mitchell72,Osius74,Mathias01}), quietly lacking both Replacement and Unbounded $\mathcal{L}_\in$-Separation. Once Replacement drops out from under Felgner's~\cite{Felgner71} insulating gasket, Global Choice ceases to be $\mathcal{L}_\in$-conservative over $\mathsf{ZC}$ (Elliot Glazer~\cite{Glazer23}), while inside $V_{\omega+\omega} \models \mathsf{ZC}$ the countable colimit $\bigcup_{n \in \mathbb{N}} \mathcal{P}^n(\mathbb{N})$ escapes the universe and blocks Harvey Friedman's~\cite{Friedman71} rank-$(\omega + 3)$ Borel Determinacy.
  \item \textbf{Double-Clamped $\mathbf{Set}_{\mathsf{ZFC}}$ Along Axes~VI and~VII:} Inside Joyal--Moerdijk~\cite{JoyalMoerdijk95} small-map class ambients $(\mathcal{C}, \mathcal{S})$ (paired with univalent $(\infty, 1)$-topos universes $(\mathcal{X}, \mathcal{U})$~\cite{Lurie09,Gylterud18} and a Volume~II higher-AST horizon $(\mathcal{C}_\infty, \mathcal{S}_\infty)$), $\mathbf{Set}_{\mathsf{ZFC}}$ emerges as a double-clamped locus where upper existence/logic compressions (Infinity, $\mathrm{LEM}$, Power Set, Choice, and \textbf{Axis~VI} Replacement / quotient-image descent $\mathrm{im}(f) \in \mathcal{S}$, plus external well-pointedness) meet lower structural truncations along \textbf{Axis~VII}: Foundation as the well-founded coalgebra restriction $\mathbf{Coalg}_{\mathrm{wf}}(\mathcal{P}_{\mathcal{S}}) \hookrightarrow \mathbf{Coalg}(\mathcal{P}_{\mathcal{S}})$ inducing strict $\kappa_{\mathrm{Reg}} : \nu^{\mathrm{wf}}\mathcal{P}_{\mathcal{S}} \simeq \mu\mathcal{P}_{\mathcal{S}} \hookrightarrow \nu\mathcal{P}_{\mathcal{S}}$~\cite{Osius74,Aczel88}, Osius pure-set fullness, Maurice Boffa's~\cite{Boffa69} Quine-atom Rieger--Bernays~\cite{Bernays54,Rieger57} slice and parity torsors, and H\aa kon Gylterud's~\cite{Gylterud18} hereditary fiberwise $(-1)$-truncation $V_0 \hookrightarrow M = \mu\mathcal{P}_{\mathrm{Anima}}$ enforcing Extensionality.
\end{itemize}

\subsection*{Part IV (Chapters 8 and 9): The Programmatic Panorama, Master Atlas, and Four-Layer Honesty Gate}

When a mathematical program deliberately works \emph{without} familiar crutches like the Axiom of Choice or arrow invertibility, hidden assumptions become doubly dangerous: a human author's muscle memory silently invokes $\mathsf{DC}_\omega$, divides by $2$ in a ring with $2$-torsion, or states a theorem $P(X) \implies Q(X)$ whose hypotheses are subtly vacuous on the intended examples.

Part~IV separates external benchmark syntheses from our proved internal spine: Chapter~\ref{ch:panorama} consolidates the programmatic panorama of classical, constructive, homological, synthetic-homotopical, and differential-cohesive benchmark lattices across literature traditions, while Chapter~\ref{ch:master-atlas} synthesizes the five core conjectures and seventeen open problems of the program alongside the architecture of our accompanying \textbf{Lean~4} formalization library (\texttt{HomotopyCRM}). Rather than treating formal verification merely as a stamp of approval after the fact, every formalized module is subjected to a \textbf{4-Layer Definition Check Gate}\index{Four-Layer Definition Check Gate}:
\begin{enumerate}
  \item \textbf{Inhabitation:} Prove a genuinely non-trivial example exists (ruling out vacuous theorems about the empty set).
  \item \textbf{Bidirectional Separation:} Exhibit both an object satisfying each core property and a concrete counterexample violating it (ruling out collapsed or tautological definitions).
  \item \textbf{Textbook Projection \& Derived Sanity:} Prove compatibility with standard \texttt{mathlib4} structures alongside non-trivial boundary theorems not baked into the axioms.
  \item \textbf{Kernel Axiom Audit:} Verify directly in Lean's trusted kernel that zero \texttt{sorry}s and zero unstated axioms were used.
\end{enumerate}

Time and again during the writing of this book, that formal microscope caught genuine mathematical phenomena that informal intuition overlooked---forcing the two-clause definition of directed convexity (\texttt{cell\_mem} + \texttt{inter\_mem}) needed to block $2$-step Steiner cycles, uncovering the exact pair of $\delta$-structures on $\mathbb{Z}[\eta]/(2\eta, \eta^2)$, and pinpointing why Marcus Nicolas's pinched $2$-cube remains flow-acyclic even where partial Steiner duality breaks down.

\section*{How to Read This Book}

A reader coming from an undergraduate background in abstract algebra, point-set/algebraic topology, and introductory logic need not master higher topos theory before turning to Chapter~1. Whenever a high-altitude categorical machine appears in the text, look first for its \textbf{anchor witness}---the parity torsor $\mathcal{T}_{\mathrm{odd}}$ on finite sets, the $2$-torsion ring $\mathbb{Z}[\eta]/(2\eta, \eta^2)$, the step-product $\mathcal{Q}$ on $[1] = \{0 < 1\}$, the facet-supported directed triangle $F_3 = \partial\Delta^2$, and the Cantor smooth-end quotient $\mathcal{Q}_{E_0} = 2^{\mathbb{N}}/E_0$ with its noncommutative CAR algebra $\mathrm{UHF}(2^\infty)$. Like Connes's irrational rotation algebra or Mac Lane's basis-free vector spaces, those small, completely explicit witnesses contain the entire subject in miniature.

\bigskip
\noindent\emph{James F.\ Walters}\\
\emph{Port Townsend, Washington, November 2026}

\chapter{Notation and Symbol Glossary}
\markboth{Notation and Symbol Glossary}{Notation and Symbol Glossary}
\label{ch:notation-glossary}

This glossary provides a unified index of mathematical symbols, categorical doctrines, topos-theoretic invariants, formal axioms, and anchor witnesses employed throughout Parts~I--IV of this monograph.

\section*{Categories, Topoi, Doctrines, and Ambient Universes}
\noindent{\scriptsize\setlength{\tabcolsep}{3pt}\renewcommand{\arraystretch}{1.13}%
\begin{xltabular}{\textwidth}{@{}>{\bfseries}P{0.24\textwidth}YP{0.21\textwidth}@{}}
\toprule
Symbol & Mathematical Meaning & Defining Context \\
\midrule
\endfirsthead
\toprule
Symbol & Mathematical Meaning & Defining Context \\
\midrule
\endhead
\bottomrule
\endfoot
$\mathbf{Set}$ & Topos of classical or constructive sets & Ch.~1, \S1.1 \\
$\mathbf{FinSet}$ & Elementary topos of hereditarily finite sets ($V_\omega$, $\neg\mathsf{Inf}$) & Ch.~1, \S\ref{sec:finset-infinity} \\
$\mathbf{Doct}$ & $2$-category (or $(\infty, 2)$-category) of categorical doctrines and structure-preserving functors & Preface; Ch.~1, \S1.2 \\
$p : \mathcal{E} \to \mathcal{B}$ & Cloven doctrinal fibration over base doctrine/category $\mathcal{B}$ with Beck--Chevalley comparison squares & Ch.~1, \S1.2 \\
$\mathcal{D}_0 \supseteq \mathcal{D}_1$ & Ambient base doctrine and stronger subdoctrine in $\mathbf{Doct}$ & Preface; Ch.~1, \S1.2 \\
$\mathbf{FrobLiftRing}_p$ & Base doctrine of commutative rings equipped with a Frobenius lift $\phi(x) \equiv x^p \pmod{pR}$ & Ch.~1, \S1.2; Ch.~\ref{ch:delta-rings}, \S\ref{sec:r2-vs-dualint} \\
$\mathbf{Topos}$ & $2$-category of elementary topoi and geometric morphisms & Ch.~1, \S1.2 \\
$\mathbf{GTop}$ & $2$-category of Grothendieck $1$-topoi & Ch.~1, \S1.2, \S1.4 \\
$\mathcal{E}, \mathcal{E}'$ & Variable elementary or Grothendieck topos & Ch.~1, \S1.1 \\
$\mathcal{S}$ & Base topos in relative topos theory $\mathcal{E}/\mathcal{S}$ & Ch.~1, \S\ref{sec:base-fiber} \\
$\mathbf{Sh}(\mathbb{C}, J)$ & Grothendieck sheaf topos on site $(\mathbb{C}, J)$ & Ch.~1, \S1.3--\S1.4 \\
$\mathbf{Cont}(\mathcal{G})$ & Atomic permutation topos of continuous $\mathcal{G}$-sets & Ch.~1, \S1.3; Ch.~6, \S\ref{subsec:kihara-choice-nominal} \\
$\mathbf{Sh}_{\mathcal{G}, \mathcal{F}}(\mathbb{B})$,\, $\mathrm{HS}_{\mathcal{F}}$ & $\mathcal{F}$-continuous $\mathcal{G}$-equivariant Boolean sheaves and hereditarily symmetric submodel $\mathrm{HS}_{\mathcal{F}} \subset V^{\mathbb{B}}$ & Ch.~2, \S\ref{sec:symmetric-bv-zf} \\
$\mathbf{Nom}$ & Topos of nominal sets ($\mathbf{Cont}(\mathrm{Sym}(\mathbb{N}))$ with finite supports) & Preface; Ch.~6, \S\ref{subsec:kihara-choice-nominal} \\
$\Eff$ & Martin Hyland's Effective Topos $\mathbf{RT}(\mathcal{K}_1)$ & Ch.~6, \S\ref{sec:kleene-eff} \\
$\mathbf{RT}(\mathcal{A})$ & Realizability topos over partial combinatory algebra $\mathcal{A}$ & Ch.~6, \S\ref{sec:kleene-eff} \\
$(\infty, 1)\text{-}\mathbf{Topos}$ & The $(\infty, 2)$-category of $(\infty, 1)$-topoi & Ch.~5, \S\ref{sec:higher-topoi-intro} \\
$\mathrm{Cat}_{(\infty, n)}$ & The $(\infty, 1)$-category of $(\infty, n)$-categories & Ch.~5, \S\ref{subsec:directed-inf-inf} \\
$\mathrm{Cat}^L_{(\infty, \infty)}$ & Postnikov-complete / Ad\'amek initial-algebra $(\infty, \infty)$-categories ($\mathcal{C} \simeq RL(\mathcal{C})$) & Preface; Ch.~5, \S\ref{subsec:directed-inf-inf} \\
$\mathrm{Cat}^\omega_{(\infty, \infty)}$ & Street--Verity $\omega$-equivalence tier of the 4-tier ladder (3rd tier below $\mathrm{Cat}^R$) & Ch.~5, \S\ref{subsec:directed-inf-inf} \\
$\mathrm{Cat}_{\le \omega}$ & Locally finite-dimensional $(\infty, \infty)$-categories (finite hom-ranks; distinct from $\mathrm{Cat}^\omega$) & Ch.~5, Prop.~\ref{prop:cat-le-omega-lfd} \\
$\mathrm{Cat}^{\mathrm{coind}}_{(\infty, \infty)}$ & Coinductively saturated $(\infty, \infty)$-categories (Ozornova--Rovelli--Walde~\cite{ORW26}) & Ch.~5, \S\ref{subsec:directed-inf-inf} \\
$\mathrm{Cat}^R_{(\infty, \infty)}$ & Right-truncation projective limit $\varprojlim_\kappa \mathrm{Cat}_{(\infty, n)}$ & Ch.~5, \S\ref{subsec:directed-inf-inf} \\
$\mathrm{CatSp}$ & Directed spectrum category $\mathrm{Sp}(\mathrm{Cat}^R_*)$ avoiding Giraud collapse & Ch.~5, \S\ref{subsec:crm-discovery-engine-inf-inf} \\
$\mathbf{Poset}_{\mathrm{og,fin}}^{\mathrm{closed}}$ & Finite closed oriented graded posets with closed embeddings & Ch.~6, \S\ref{subsec:cubical-amalgamation} \\
$\mathbf{DirComp}_{\mathrm{acyc}}^{\mathrm{convex}}$ & Finite flow-acyclic convex directed complexes & Ch.~6, \S\ref{subsec:cubical-amalgamation} \\
$\mathbf{Clan}(\mathcal{S}, \mathbf{Nom})$ & Grothendieck--Maltsiniotis weak $\omega$-category / clan models & Ch.~6, \S\ref{subsec:coinductive-streams-choice} \\
$(\mathcal{C}, \mathcal{S})$,\, $(\mathcal{C}_\infty, \mathcal{S}_\infty)$ & $1$-categorical Joyal--Moerdijk small-map AST ambient and Volume~II class-anima horizon (modeled via $(\mathcal{X}, \mathcal{U}_0 \in \mathcal{U}_1)$) & Ch.~\ref{ch:sets-residuum}, \S\ref{sec:ast-plumbing-compression}, \S\ref{sec:scale-invariance-synthesis} \\
$\mathbf{Cat}(\mathcal{S})$,\, $\mathbf{Cat}(\mathcal{C})$ & $2$-categories of internal small categories vs.\ large class categories & Ch.~\ref{ch:sets-residuum}, \S\ref{sec:conservativity-kaleidoscope} \\
$V_{\mathsf{ZF}} = \mu\mathcal{P}_{\mathcal{S}}$, $V_{\mathsf{AFA}} = \nu\mathcal{P}_{\mathcal{S}}$, $M = \mu\mathcal{P}_{\mathrm{Anima}}$ & Well-founded, Aczel coalgebraic anti-foundation, and Gylterud multiset hierarchies & Ch.~\ref{ch:sets-residuum}, \S\ref{sec:anti-existence-residua} \\
\end{xltabular}\par}

\vspace{1.0em}
\section*{CRM Operators, Residua, and Homological Invariants}
\noindent{\scriptsize\setlength{\tabcolsep}{3pt}\renewcommand{\arraystretch}{1.13}%
\begin{xltabular}{\textwidth}{@{}>{\bfseries}P{0.24\textwidth}YP{0.21\textwidth}@{}}
\toprule
Symbol & Mathematical Meaning & Defining Context \\
\midrule
\endfirsthead
\toprule
Symbol & Mathematical Meaning & Defining Context \\
\midrule
\endhead
\bottomrule
\endfoot
$\mathrm{Dom}(\mathrm{Res})$ & Domain of doctrinal comparison triples $\mathfrak{T} = (\mathcal{D}_0 \supseteq \mathcal{D}_1, \kappa : \Phi \Rightarrow \Psi)$ induced by comparison squares & Ch.~1, \S1.2 \\
$\mathbf{Mod}(T)$ & Subclass of universes satisfying property or axiom $T$ & Ch.~1, \S1.2 \\
$\mathrm{Def}(\alpha),\, \mathrm{Res}(\kappa),\, \mathcal{R}(\mathcal{E})$ & Universal Doctrinal Square Defect of comparison square $\alpha$ and induced categorical residuum operator & Preface; Ch.~1, \S1.2 \\
$\mathfrak{P} = (\mathbb{D}, \mathcal{P}),\, \mathrm{Res}_{\mathcal{P}}^{\min}$ & Doctrinal Philosophy (framed equipment paired with probe suite $\mathcal{P}$) and minimal Pareto Residuum Frontier & Ch.~1, \S\ref{sec:admissibility} \\
$\widehat{\mathbb{N}}_{\mathrm{lasso}},\, L_{p,q}$ & Finite $S$-algebra lassos ($s^{p+q}(x_0) = s^p(x_0)$) and profinite limit $\varprojlim L_{p,q} \twoheadrightarrow \widehat{\mathbb{Z}}$ & Ch.~1, \S\ref{sec:finset-infinity} \\
$\mathrm{Res}_{\mathrm{Ded}}(X)$ & Dedekind/pigeonhole non-surjective monomorphism gap $\mathrm{Mono}(X,X)\setminus\mathrm{Iso}(X,X)$ & Ch.~1, \S\ref{sec:finset-infinity} \\
$\mathcal{Q}_{E_0} = 2^{\mathbb{N}}/E_0$ & Profinite hyperfinite tail quotient $\prod_{\mathbb{N}}\mathbb{F}_2 / \bigoplus_{\mathbb{N}}\mathbb{F}_2$ (Connes twin of $\mathbb{T}^2/\mathcal{F}_\theta$) and CAR convolution algebra $\mathrm{UHF}(2^\infty)$ & Preface; Ch.~1, \S\ref{sec:finset-infinity}; Ch.~4, \S\ref{sec:exotic-r4-e0} \\
$\Omega_{\mathcal{E}},\, 1_{\mathcal{E}}$ & Subobject classifier and terminal object of category or topos $\mathcal{E}$ & Ch.~1, \S1.1 \\
$j : \Omega \to \Omega$ & Lawvere--Tierney modal topology & Ch.~1, \S\ref{sec:admissibility} \\
$\CohesiveQuad$ & Lawvere cohesive adjoint quadruple ($\Pi \dashv \mathrm{Disc} \dashv \Gamma \dashv \mathrm{Codisc}$) & Ch.~1, \S1.5; Ch.~\ref{ch:panorama}, \S\ref{sec:cohesive-lattice} \\
$H^1(\mathcal{E}/Y, G)$ & \v{C}ech cohomology / torsor classification set with coefficients in $G$ & Ch.~2, \S2.4 \\
$H^2_{\mathrm{gerbe}}(\mathcal{X}/Y, \mathfrak{L})$ & Giraud non-abelian $1$-gerbe obstruction set over $Y$ (for band $\mathfrak{L}$ or $L(G)$) & Ch.~5, \S\ref{sec:postnikov-residuum} \\
$H^{n+2}(\mathcal{X}/Y, A)$ & Postnikov $k$-invariant obstruction classified by $K(A, n+2)$ ($n \ge 1$) & Ch.~5, \S\ref{sec:postnikov-residuum} \\
$\tau_{\le n}, \pi_{\le n}$ & Postnikov $n$-truncation operator and homotopy group sheaves & Ch.~5, \S\ref{sec:postnikov-residuum} \\
$\pi_{\le 0}$ & Directed $0$-truncation / Left Bousfield groupoidification functor & Preface; Ch.~5, \S\ref{subsec:directed-inf-inf} \\
$\Sigma_{\mathrm{red}} \dashv \Omega_{\mathrm{dir}}$ & Pointed directed suspension--loop adjunction on $\mathrm{Cat}^R_*$ & Ch.~5, \S\ref{subsec:crm-discovery-engine-inf-inf} \\
$\mathcal{P}_{\mathcal{S}}$ & Small-powerset functor (contravariant on $(\mathcal{C}, \mathcal{S})$; covariant via $\exists_g$ under Axis~VI Replacement) & Ch.~\ref{ch:sets-residuum}, \S\ref{sec:ast-plumbing-compression} \\
$\mathrm{Def} \dashv \pi_{\mathrm{Set}}$ & Minimal definable-class left adjoint to set projection on $\mathbf{Mod}(\mathsf{GB}_0)$ & Ch.~\ref{ch:sets-residuum}, \S\ref{sec:conservativity-kaleidoscope} \\
\end{xltabular}\par}

\vspace{1.0em}
\section*{Choice Principles, Logical Systems, and Formal Axioms}
\noindent{\scriptsize\setlength{\tabcolsep}{3pt}\renewcommand{\arraystretch}{1.13}%
\begin{xltabular}{\textwidth}{@{}>{\bfseries}P{0.24\textwidth}YP{0.21\textwidth}@{}}
\toprule
Symbol & Description & Defining Context \\
\midrule
\endfirsthead
\toprule
Symbol & Description & Defining Context \\
\midrule
\endhead
\bottomrule
\endfoot
$\mathrm{AC},\, \mathrm{EAC}$ & Classical External Axiom of Choice (all epimorphisms split globally) & Ch.~1, \S1.1; Ch.~2, \S2.1 \\
$\mathrm{IAC}$ & Internal Axiom of Choice (dependent product $\Pi_f$ / $(-)^Z$ preserves epimorphisms) & Ch.~2, \S2.1 \\
$\mathrm{DC}, \mathrm{IDC}$ & Dependent Choice and Internal Dependent Choice ($\varprojlim_n X_n \twoheadrightarrow X_0$ epic) & Ch.~1, \S1.1; Ch.~2, \S2.3 \\
$\mathrm{COSHEP}, \mathrm{ICOSHEP}$ & Enough Projectives (external and internal Presentation Axiom) & Ch.~2, \S2.3; Ch.~4, \S4.3 \\
$\mathrm{WISC}$ & Weakly Initial Set of Covers (Axiom of Multiple Choice) & Ch.~2, \S2.3 \\
$\mathrm{SVC}, \mathrm{SVC}_{\mathrm{ret}}, \mathrm{ISVC}_{\mathrm{ret}}$ & Blass's Small Violations of Choice (subquotient and external/internal retract forms) & Ch.~2, \S2.3; Ch.~4, \S4.3 \\
$\mathrm{BPI}, \mathrm{IBPI}$ & Boolean Prime Ideal Theorem (external and internal) & Ch.~2, \S2.3 \\
$\mathrm{FAC}_k$ & Finite Choice on families of cardinality $k$ & Ch.~6, \S\ref{sec:permutation-topoi} \\
$\mathrm{AC}_n, \mathrm{AC}_\infty$ & Higher truncated choice principles in $(\infty, 1)$-topoi & Ch.~5, \S\ref{sec:higher-choice} \\
$\mathrm{LEM}, \mathrm{WLEM}$ & Law of Excluded Middle and Weak LEM ($\neg P \lor \neg\neg P$) & Ch.~1, \S1.1 \\
$\mathrm{LLPO}$ & Lesser Limited Principle of Omniscience & Ch.~1, \S1.1 \\
$\mathrm{NNO}$ & Natural Numbers Object (initial algebra $\mu(\mathrm{Id}+1)$) & Ch.~1, \S\ref{sec:finset-infinity}; Ch.~2, \S2.1 \\
$\mathrm{RCA}_0$--$\Pi^1_1\text{-}\mathrm{CA}_0$ & The classical Friedman--Simpson ``Big Five'' subsystems & Ch.~1, \S1.1 \\
$\mathsf{Mac}$ ($\mathsf{ZBQC} \simeq \mathsf{ETCS}$), $\mathsf{ZC}$ & Bounded Zermelo / Mac~Lane set theory ($\Delta_0$-Separation) and full Zermelo set theory with Choice & Ch.~\ref{ch:sets-residuum}, \S\ref{sec:cracks-in-the-floor} \\
$\mathsf{ZFC}^-$,\, $\mathsf{ZFC}_{\neg\mathrm{Reg}}$ & Power-Set-free $\mathsf{ZFC}$ (with Collection~\cite{GitmanHamkinsJohnstone16}) and Foundation-free $\mathsf{ZFC}$~\cite{Felgner76} & Ch.~\ref{ch:sets-residuum}, \S\ref{sec:cracks-in-the-floor}--\S\ref{sec:ast-plumbing-compression} \\
$\mathsf{GB}_0$, $\mathsf{GBC}_0$, $\mathsf{GB}$, $\mathsf{GBC}$, $\mathsf{NBG}$ & Replacement-free ($\mathsf{GB}_0, \mathsf{GBC}_0$), Bernays--G\"odel ($\mathsf{GB}, \mathsf{GBC}$), and Global-Choice ($\mathsf{NBG}$) class theories & Ch.~\ref{ch:sets-residuum}, \S\ref{sec:conservativity-kaleidoscope} \\
$\mathsf{LoS}$ & Von Neumann's Limitation of Size ($C \notin \mathcal{S} \iff C \twoheadrightarrow V$; equivalent over $\mathsf{GB}_0$ to $\mathrm{Rep} \wedge \mathsf{GC}$) & Ch.~\ref{ch:sets-residuum}, \S\ref{sec:cracks-in-the-floor} \\
$\mathsf{ETR}_{\mathrm{Ord}}$,\, $\mathsf{KM}$ & Elementary Transfinite Recursion along $\mathrm{Ord}$~\cite{Gitman20} and impredicative Kelley--Morse class theory & Ch.~\ref{ch:sets-residuum}, \S\ref{sec:conservativity-kaleidoscope} \\
$\mathsf{AFA}$,\, $\mathsf{BAFA}$ & Aczel's and Boffa's Anti-Foundation Axioms ($\nu\mathcal{P}_{\mathcal{S}}$ bisimulation vs.\ distinct Quine-atom extensions) & Ch.~\ref{ch:sets-residuum}, \S\ref{sec:anti-existence-residua} \\
\end{xltabular}\par}

\vspace{1.0em}
\section*{The Five Foundational Anchor Witnesses}
\noindent{\scriptsize\setlength{\tabcolsep}{3pt}\renewcommand{\arraystretch}{1.14}%
\begin{tabularx}{\textwidth}{@{}>{\bfseries}P{0.23\textwidth}YP{0.29\textwidth}@{}}
\toprule
Symbol & Mathematical Role and Obstruction Property & Verified Lean Module \\
\midrule
$\mathcal{T}_{\mathrm{odd}}$ & Odd-finset parity torsor in $\mathbf{Nom}$; $[\mathcal{T}_{\mathrm{odd}}] = \partial(1) \neq 0 \in H^1(\mathbf{Nom}; \mathrm{EvenFinset})$, proving non-projectivity of $1_{\mathbf{Nom}}$ & \texttt{HomotopyCRM/\allowbreak NominalSimplicity.lean} \\
$R_2 = \mathbb{Z}[\eta]/(2\eta, \eta^2)$ & $2$-torsion $\delta$-ring ($\mathbb{Z} \oplus \mathbb{F}_2$ additively; $R_2 \not\cong \mathbb{Z}\times\mathbb{F}_2$ as rings) with two $\delta$-lifts ($\delta_0(\eta)=0, \delta_1(\eta)=\eta$) over augmentation Frobenius $\phi(a+b\eta)=a$; companion \texttt{DualInt} ($\mathbb{Z}[\varepsilon]/(\varepsilon^2)$) separates Nygaard filtration from perfection & \texttt{HomotopyCRM/\allowbreak DeltaRingPerfection.lean} \\
$\mathcal{Q} \subset \prod_{k \ge 1} \mathcal{A}_k$ & Gaunt Noetherian rank-$\omega$ step-product on $[1]$ ($\mathrm{Hom}(1_k,1_k) = \mathcal{M}_k \not\simeq *$ vs.\ contractible off-diagonals, rank $\min(m,r)$); generates phantom $X_\bullet \in RL(\mathcal{Q}) \setminus \mathcal{Q}$ & \texttt{HomotopyCRM/\allowbreak StepProductPhantom.lean} \\
$F_3 = \partial\Delta^2,\, F_n$ & Facet cycle complex; directed flow-acyclic obstruction to cubical $n$-amalgamation ($n \ge 3$) despite binary ($n = 2$) SAP & \texttt{HomotopyCRM/\allowbreak CubicalAmalgamation.lean} \\
$\mathcal{Q}_{E_0} = 2^{\mathbb{N}}/E_0,\, \mathrm{UHF}(2^\infty)$ & Cantor smooth-end quotient $\prod_{\mathbb{N}}\mathbb{F}_2 / \bigoplus_{\mathbb{N}}\mathbb{F}_2$; $\pi_0^{\mathrm{Pol}} \simeq 1$ vs.\ $\mathsf{ZFC}$ Hamel section, CAR convolution algebra $K_0 \cong \mathbb{Z}[\tfrac{1}{2}]$, and $E_0^{\mathrm{even}}$ parity torsor & \texttt{HomotopyCRM/\allowbreak BorelTailResiduum.lean} \\
\bottomrule
\end{tabularx}\par}

\mainmatter
\setcounter{secnumdepth}{2}

\part{Foundations of Necessity and the Residuum Calculus}

\chapter{Foundations of Categorical Reverse Mathematics}
\chaptermark{Foundations of Categorical Reverse Mathematics}
\label{ch:foundations}
\chapterepigraph{``When the theorem is proved from the right axioms, the axioms can be proved from the theorem.''}{Harvey Friedman, ICM Vancouver, 1974; \emph{Proc.}, 1975}

\section{From Syntactic Axioms to Geometric Universes}

For half a century, the foundational program of Reverse Mathematics has sought to determine the exact logical strength of mathematical theorems. In the classical tradition inaugurated by Harvey Friedman and Stephen Simpson~\cite{Friedman75,Simpson09}, one fixes a single weak base theory---typically $\mathrm{RCA}_0$, a subsystem of second-order arithmetic equipped with recursive comprehension, $\Sigma^0_1$ induction, and the classical Law of the Excluded Middle ($\mathrm{LEM}$)---and calibrates a theorem $T$ by identifying the weakest set-existence axiom $A$ such that $\mathrm{RCA}_0 \vdash A \leftrightarrow T$. Hundreds of theorems across classical analysis, algebra, and combinatorics cluster into a linearly ordered hierarchy of five canonical subsystems, known as the ``Big Five''\index{Big Five subsystems}: $\mathrm{RCA}_0$, $\mathrm{WKL}_0$, $\mathrm{ACA}_0$, $\mathrm{ATR}_0$, and $\Pi^1_1\text{-}\mathrm{CA}_0$.

Constructive Reverse Mathematics\index{Constructive Reverse Mathematics}, developed by Hajime Ishihara and his collaborators~\cite{Ishihara06}, refines this calibration by lowering the base theory from $\mathrm{RCA}_0$ to Errett Bishop's informal constructive mathematics ($\mathrm{BISH}$)~\cite{Bishop67}. Working without $\mathrm{LEM}$ or the Axiom of Choice ($\mathrm{AC}$), the constructive program reveals that many classical equivalences factor into a constructive analytical core combined with semi-classical logical principles---such as the Weak Law of Excluded Middle ($\mathrm{WLEM}$) or the Lesser Limited Principle of Omniscience ($\mathrm{LLPO}$)---or weak choice axioms like Countable Choice ($\mathrm{CC}$) and Dependent Choice ($\mathrm{DC}$).

Yet both classical and constructive reverse mathematics share a fundamental limitation: they treat the ambient mathematical universe as a fixed syntactic background and project all mathematical phenomena onto a one-dimensional ladder of formal axioms. Modern mathematics, however, does not live in a single formal language of second-order arithmetic. Since Alexander Grothendieck's work in algebraic geometry and F.~William Lawvere and Myles Tierney's axiomatization of elementary topoi~\cite{MacLaneMoerdijk92,Johnstone02}, and subsequently Jacob Lurie's synthesis of higher topos theory~\cite{Lurie09}, mathematicians have understood that mathematical structures are interpreted across a vast landscape of variable categorical universes and doctrines---sheaf topoi over topological spaces and locales, realizability topoi in computability theory, smooth topoi in synthetic differential geometry, commutative algebra with Frobenius lifts, and higher $(\infty, 1)$-topoi and $(\infty, \infty)$-categories in homotopy theory.

\textbf{Categorical Reverse Mathematics}\index{Categorical Reverse Mathematics} ($\mathrm{CRM}$) shifts the calibration of mathematical theorems from syntactic axioms in a fixed formal system to the structural, algebraic, topological, and cohomological geometry of abstract mathematical doctrines. Rather than asking solely which set-existence axiom must be added to arithmetic to prove a theorem, $\mathrm{CRM}$ treats the ambient universe or algebraic doctrine $X \in \mathcal{D}_0$ as a variable space and asks: given two mathematical constructions $\Phi$ and $\Psi$ connected by a comparison transformation $\kappa : \Phi \Rightarrow \Psi$ that becomes an equivalence inside a stronger subdoctrine $\mathcal{D}_1 \subseteq \mathcal{D}_0$, what exact structural invariant $\mathrm{Res}(\kappa)(X)$ measures the gap separating $\mathcal{D}_0$ from $\mathcal{D}_1$?

\section{The Fundamental Equivalence: Doctrinal Fibrations, Comparison Squares, and the Universal Defect}\label{sec:fundamental-equivalence}

To establish Categorical Reverse Mathematics on a genuinely universal foundation---capable of treating Choice, constructive logic, the Axiom of Infinity, $p$-torsion algebra, and directed higher invertibility on equal footing without hard-coding specialized syntactic templates---we formulate the residuum operator via concrete cloven \emph{doctrinal fibrations} in the tradition of Lawvere~\cite{Lawvere69,Maietti05,Johnstone02} and Grothendieck's calculus of base change.

\begin{definition}[Doctrinal Fibration, Comparison Square, and $\mathrm{Dom}(\mathrm{Res})$]\label{def:comparison-triple}%
\lean{HomotopyCRM.DoctrinalFibration, HomotopyCRM.MateDoubleSquare, HomotopyCRM.mateDoubleSquare_mate_involutive}%
\leanok%
A \textbf{categorical doctrine}\index{Categorical doctrine}\index{Doctrine|see{Categorical doctrine}} is a category, $2$-category, or $(\infty, 2)$-category $\mathcal{D}$ whose objects $X \in \mathcal{D}$ are mathematical universes or structured algebraic/higher-categorical objects equipped with a specified signature of categorical operations, limits, colimits, quantifiers, modalities, or coherent higher cells, and whose morphisms preserve that structure up to canonical isomorphism. Let $\mathbf{Doct}$\index{Doct@$\mathbf{Doct}$ ($2$-category of doctrines)} denote the meta-$2$-category (or $(\infty, 2)$-category) of categorical doctrines and structure-preserving functors.
\begin{enumerate}[label=\textup{(\roman*)},leftmargin=2.0em]
  \item A \textbf{cloven doctrinal fibration}\index{Doctrinal fibration} is a Grothendieck fibration $p : \mathcal{E} \to \mathcal{B}$ (\texttt{DoctrinalFibration E B}) whose fibers $\mathcal{E}_B$ over $B \in \mathcal{B}$ carry doctrinal structure preserved by Cartesian reindexing $u^* : \mathcal{E}_{B'} \to \mathcal{E}_B$ along $u : B \to B'$ (specializing via $\mathrm{id}_{\mathcal{D}} : \mathcal{D} \to \mathcal{D}$ to any bare doctrine $\mathcal{D} \in \mathbf{Doct}$).
  \item A \textbf{comparison square} $\alpha$\index{Comparison square} (\texttt{MateDoubleSquare}) between adjoint doctrinal functors $(L \dashv R : \mathcal{D} \rightleftarrows \mathcal{E})$ and $(L' \dashv R' : \mathcal{D}' \rightleftarrows \mathcal{E}')$ paired with functors $u : \mathcal{D} \to \mathcal{D}'$ and $F : \mathcal{E} \to \mathcal{E}'$ is a $2$-cell $\alpha : L' \circ u \Longrightarrow F \circ L$, together with its \textbf{Beck--Chevalley mate}\index{Beck--Chevalley mate}
  \[
    \mu_\alpha \;:\; u \circ R \Longrightarrow R' \circ F
  \]
  whose unit--counit transposition $\alpha \mapsto \mu_\alpha$ is strictly involutive (\texttt{mateDoubleSquare\_\allowbreak mate\_\allowbreak involutive}). More generally, a comparison square between fibrations $p : \mathcal{E} \to \mathcal{B}$ and $p' : \mathcal{E}' \to \mathcal{B}'$ consists of functors $(F, u)$ and a comparison $2$-cell $\alpha : p' \circ F \Longrightarrow u \circ p$ (or between fiberwise adjoints along $u$).
  \item Evaluating $\Phi := u \circ R$ and $\Psi := R' \circ F$ across $\mathcal{D}_0 \in \mathbf{Doct}$ against the Beck--Chevalley equivalence locus $\mathcal{D}_1 := \mathrm{Loc}_{\mathrm{BC}}(\alpha) \subseteq \mathcal{D}_0$ induces a \textbf{doctrinal comparison triple}\index{Doctrinal comparison triple}
  \[
    \mathfrak{T} \;=\; \bigl(\mathcal{D}_0 \supseteq \mathcal{D}_1, \quad \kappa := \mu_\alpha : \Phi \Longrightarrow \Psi\bigr),
  \]
  where $\kappa|_{\mathcal{D}_1} : \Phi|_{\mathcal{D}_1} \xrightarrow{\;\sim\;} \Psi|_{\mathcal{D}_1}$ is an equivalence. The \textbf{domain of the categorical residuum}, denoted $\mathrm{Dom}(\mathrm{Res})$\index{Dom(Res)@$\mathrm{Dom}(\mathrm{Res})$}, is the collection of all comparison squares and induced triples across mathematics.
\end{enumerate}
\end{definition}

\begin{definition}[Universal Doctrinal Square Defect and Categorical Residuum]\label{def:residuum}%
\uses{def:comparison-triple}%
Let $\alpha$ be a comparison square with Beck--Chevalley mate $\mu_\alpha : \Phi \Longrightarrow \Psi$ and associated triple $\mathfrak{T} = (\mathcal{D}_0 \supseteq \mathcal{D}_1, \; \kappa : \Phi \Longrightarrow \Psi) \in \mathrm{Dom}(\mathrm{Res})$. The \textbf{Universal Doctrinal Square Defect}\index{Universal Doctrinal Square Defect}\index{Categorical residuum} of $\alpha$, denoted $\mathrm{Def}(\alpha)$ (written $\mathrm{Res}(\kappa) := \mathrm{Def}(\alpha_\kappa)$ on $\mathfrak{T}$), is the single structural obstruction assignment
\[
  \mathrm{Def}(\alpha) : \mathcal{D}_0 \longrightarrow \mathbf{Obstr}
\]
measuring, for each $X \in \mathcal{D}_0$, the obstruction to $\mu_{\alpha, X} : \Phi(X) \to \Psi(X)$ being an equivalence with $(-1)$-connected section space ($\|\Gamma_X(\alpha)\|_{-1} \simeq 1$) inside a pointed obstruction category $\mathbf{Obstr}$ with distinguished trivial object $0 \in \mathbf{Obstr}$:
\[
  \mathrm{Def}(\alpha)(X) \simeq 0 \iff \kappa_X : \Phi(X) \xrightarrow{\;\sim\;} \Psi(X) \text{ with } \|\Gamma_X(\alpha)\|_{-1} \simeq 1.
\]
This single defect validates the \textbf{Fundamental Equivalence of Categorical Reverse Mathematics}:
\[
  X \in \mathcal{D}_0 \quad \text{and} \quad \mathrm{Def}(\alpha)(X) \simeq 0 \iff X \in \mathcal{D}_1,
\]
written compactly as $\mathcal{D}_0 + \mathrm{Res}(\kappa) \iff \mathcal{D}_1$.

As an immediate corollary in the topos-logical specialization $\mathcal{D}_0 = \mathbf{Mod}(A) := \{\mathcal{E} \in \mathbf{Topos} \mid \mathcal{E} \models A\}$ and $\mathcal{D}_1 = \mathbf{Mod}(B) \subsetneq \mathbf{Mod}(A)$ for internal principles $B \implies A$, setting $\mathcal{R} = (A \Rightarrow_{\mathbf{Topos}} B) := \mathrm{Def}(\alpha)$ yields the factorization $\mathbf{Mod}(A) \cap \mathbf{Mod}(\mathcal{R}) = \mathbf{Mod}(B)$, or $A + \mathcal{R}(\mathcal{E}) \iff B$.
\end{definition}

In the topos specialization, the categorical residuum $(A \Rightarrow_{\mathbf{Topos}} B)$ acts order-theoretically as a relative pseudocomplement in the lattice of universe classes (Figure~\ref{fig:comparison-triple-residuum}(c)): a strong theorem $B$ factors orthogonally into its internal deductive core $A$ together with the ambient geometry $\mathcal{R}(\mathcal{E})$ required of the surrounding world. Crucially, Categorical Reverse Mathematics must be distinguished from Olivia Caramello's program of ``toposes as bridges'' via Morita equivalence~\cite{Caramello18}. Caramello's methodology compares two \emph{different syntactic geometric theories} $T_1$ and $T_2$ sharing the \emph{exact same classifying topos} $\mathcal{E}[T_1] \simeq \mathcal{E}[T_2]$, transferring properties across alternative sites $(\mathbb{C}_{T_1}, J_1)$ and $(\mathbb{C}_{T_2}, J_2)$ of a single invariant space. By contrast, $\mathrm{CRM}$ evaluates $\mathrm{Def}(\alpha)$ across doctrines $\mathcal{D}_0 \supsetneq \mathcal{D}_1$ of \emph{genuinely unequal foundational strength}.

\begin{proposition}[Derived Tri-Modal Factorization of Square Defects]\label{prop:tri-modal-factorization}%
\lean{HomotopyCRM.intToZMod2TypeMode_sectionSpace_nontrivial, HomotopyCRM.not_intToZMod2TypeMode_contractible, HomotopyCRM.isEmpty_intToZMod2AbMode_sectionSpace, HomotopyCRM.profiniteLassoTowerMode_hasPhantom, HomotopyCRM.stepProductTowerMode_hasPhantom}%
\leanok%
\uses{def:comparison-triple,def:residuum}%
Under comprehensive Horizontal / $\mathcal{J}$-Completion / Vertical factorization of comparison squares in cloven fibrations, the single Universal Doctrinal Square Defect $\mathrm{Def}(\alpha)$ projects canonically onto three obstruction modes (Figure~\ref{fig:comparison-triple-residuum}(b)):
\begin{enumerate}[label=\textbf{Mode~(\roman*).},leftmargin=4.2em]
  \item \textbf{Horizontal Base-Change / Section Obstruction ($\mathrm{Res}_{\mathrm{sec}}$):} Measuring $(-1)$-connected section inhabitation $\|\Pi_f(p)\|_{-1} \simeq 1$ (\texttt{SectionObstructionMode}) along dependent product functors $f^* \dashv \Pi_f$. In particular, $(-1)$-connected inhabitation is strictly weaker than $(-2)$-contractibility: for $q : \mathbb{Z} \twoheadrightarrow \mathbb{Z}/2\mathbb{Z}$, the set-level section space is non-trivial (\texttt{intToZMod2TypeMode\_\allowbreak sectionSpace\_\allowbreak nontrivial}) and therefore inhabited yet not contractible (\texttt{not\_\allowbreak intToZMod2TypeMode\_\allowbreak contractible}), whereas its additive group lift is empty (\texttt{isEmpty\_\allowbreak intToZMod2AbMode\_\allowbreak sectionSpace}).
  \item \textbf{$\mathcal{J}$-Completion / Tower Phantom ($\mathrm{Res}_{\mathrm{lim}}$):} Measuring the completion gap $\widehat{\mathcal{E}}_{\mathcal{J}} \setminus \mathcal{E}$ (\texttt{TowerPhantomMode}) and derived $\varprojlim^1$ when a pro-system or Postnikov truncation tower fails to stabilize in $\mathcal{E}$, witnessed by $-1 \in \widehat{\mathbb{Z}} \setminus \mathrm{range}(\mathbb{N})$ for $\widehat{\mathbb{N}}_{\mathrm{lasso}}$ (\texttt{profiniteLassoTowerMode\_\allowbreak hasPhantom}) and by $X_\bullet \in RL(\mathcal{Q}) \setminus \mathcal{Q}$ (\texttt{stepProductTowerMode\_\allowbreak hasPhantom}).
  \item \textbf{Vertical Forgetful Homotopy-Fiber Lift ($\mathrm{Res}_{\mathrm{lift}}$):} Evaluating a forgetful doctrinal functor $U : \mathcal{E}' \to \mathcal{E}$ directly at a \emph{base} object $B \in \mathcal{E}$ via its homotopy fiber
  \[
    \mathrm{hofib}_B(U) \;:=\; \sum_{E' \in \mathcal{E}'} \mathrm{Iso}_{\mathcal{E}}\bigl(U(E'), B\bigr)
  \]
  (\texttt{ForgetfulLiftMode}), whose quotient $\pi_0\,\mathrm{hofib}_B(U)$ classifies distinct enriched lifts modulo $\mathrm{Aut}_{\mathcal{E}}(B)$ conjugation.
\end{enumerate}
\end{proposition}

\begin{proof}
Every comparison square $\alpha : p' \circ F \Rightarrow u \circ p$ in the $2$-category $\mathbf{Fib}$ of cloven doctrinal fibrations factors up to canonical equivalence into: (1) a Cartesian base-change square along $u$, whose defect is the obstruction to inverting the right adjoint Beck--Chevalley mate $\Pi_u$ via global/slice section inhabitation $\|\Pi_f(p)\|_{-1} \simeq 1$ (Mode~(i)); (2) a $\mathcal{J}$-dense completion embedding $\mathcal{E} \hookrightarrow \widehat{\mathcal{E}}_{\mathcal{J}}$ into pro-objects or Postnikov limits, whose defect is the strict complement $\widehat{\mathcal{E}}_{\mathcal{J}} \setminus \mathcal{E}$ and associated $\varprojlim^1$ phantom locus (Mode~(ii)); and (3) a vertical functor $U : \mathcal{E}' \to \mathcal{E}$ over $\mathrm{id}_{\mathcal{B}}$, whose defect over $B \in \mathcal{E}$ is precisely the homotopy fiber $\mathrm{hofib}_B(U)$ and its $\mathrm{Aut}_{\mathcal{E}}(B)$-orbit quotient (Mode~(iii)). The concrete witnesses for Modes~(i)--(iii) follow from direct evaluation: parity sections $\{0,1\}^{\mathbb{Z}/2\mathbb{Z}}$ split $q : \mathbb{Z} \twoheadrightarrow \mathbb{Z}/2\mathbb{Z}$ in $\mathbf{Set}$ with cardinality $>1$ while $\mathrm{Hom}_{\mathbf{Ab}}(\mathbb{Z}/2\mathbb{Z}, \mathbb{Z}) = \emptyset$; $-1 \in \widehat{\mathbb{Z}}$ and the inductive prefix chain $(e^{(n-1)})_{n \ge 1}$ precipitate phantoms outside $\mathbb{N}$ and $\mathcal{Q}$; and the vertical forgetful fiber over $\mathbf{FrobLiftRing}_2$ yields the three-stage lift stratification proved in Chapter~\ref{ch:delta-rings} (Proposition~\ref{prop:delta-three-stage-tower}).
\end{proof}

\begin{remark}[Framed Equipments, Axiom Encoding, and Morita Invariance]\label{rem:framed-equipments}%
\uses{def:comparison-triple,def:residuum,prop:tri-modal-factorization}%
Organizing cloven doctrinal fibrations globally yields a \emph{framed bicategory}---or more generally a framed $(\infty, 2)$-equipment $\mathbb{D}$ in the sense of Shulman~\cite{Shulman08}---governed by the higher calculus of mates for two-variable fibrations established by Haugseng, Hebestreit, Linskens, and Nuiten~\cite{HHLN23}. Leading our exposition with concrete comparison squares rather than axiomatic equipments keeps the foundation directly computable, while separating two logically distinct steps: first, the \emph{modeling choice} that encodes an informal principle $\Phi$ into an explicit comparison square $\alpha_\Phi$ (for instance, choosing which covers test Choice or which tower truncates infinity); and second, the \emph{Morita-invariant defect} $\mathrm{Def}(\alpha_\Phi)$, which depends solely on the equivalence class of $\alpha_\Phi$ in the framed equipment $\mathbb{D}$.
\end{remark}

\begin{figure}[htbp]
\centering
\begin{tikzpicture}[
  >=Stealth,
  doc/.style={draw=CRMBlue!80,rounded corners=2.5pt,fill=CRMBlue!5,inner sep=3pt,font=\scriptsize,align=center},
  mod/.style={draw=CRMBlue,rounded corners=2.5pt,fill=CRMBlue!9,inner sep=3.2pt,font=\scriptsize,align=center},
  edge/.style={->,thick,draw=black!75},
  cbedge/.style={->,thick,draw=CRMBlue}
]
  % Panel (a): Top-Left — Master Comparison Square alpha & Mate
  \node[font=\scriptsize\bfseries,color=CRMBlue,anchor=west] at (-0.15,4.95) {(a) Master Square $\alpha$ \& Mate $\mu_\alpha$};
  \node[doc] (E)  at (0.45,4.25) {$\mathcal{E}$};
  \node[doc] (Ep) at (3.05,4.25) {$\mathcal{E}'$};
  \node[doc] (B)  at (0.45,2.95) {$\mathcal{B}$};
  \node[doc] (Bp) at (3.05,2.95) {$\mathcal{B}'$};
  \draw[edge] (E) -- node[above,font=\tiny]{$F$} (Ep);
  \draw[edge] (E) -- node[left,font=\tiny]{$p$} (B);
  \draw[edge] (Ep) -- node[right,font=\tiny]{$p'$} (Bp);
  \draw[edge] (B) -- node[below,font=\tiny]{$u$} (Bp);
  \draw[-{Implies},double distance=1.6pt,thick,draw=CRMBlue] (1.35,3.85) -- node[above right,font=\tiny,inner sep=1pt]{$\alpha$} (2.15,3.35);
  \node[font=\tiny,color=CRMBlue,align=center] at (1.75,2.30) {Mate $\mu_\alpha : u \circ R \Rightarrow R' \circ F$ \\ Master defect: $\mathrm{Def}(\alpha)$};

  % Panel (c): Top-Right — Pullback & Topos Specialization
  \node[font=\scriptsize\bfseries,color=CRMBlue,anchor=west] at (5.15,4.95) {(c) Universal Defect Pullback ($\mathbf{Topos}$)};
  \node[mod] (ModB) at (6.45,4.25) {$\mathcal{D}_1 = \mathbf{Mod}(B)$};
  \node[doc] (Zero) at (9.45,4.25) {$0_{\mathbf{Obstr}}$};
  \node[doc] (ModA) at (6.45,2.95) {$\mathcal{D}_0 = \mathbf{Mod}(A)$};
  \node[doc] (Obstr) at (9.45,2.95) {$\mathbf{Obstr}$};
  \draw[cbedge] (ModB) -- node[above,font=\tiny]{$!$} (Zero);
  \draw[{Hooks[right]}->,thick,draw=CRMBlue] (ModB) -- node[left,font=\tiny]{$\iota$} (ModA);
  \draw[{Hooks[right]}->,edge] (Zero) -- (Obstr);
  \draw[cbedge] (ModA) -- node[below,font=\tiny]{$\mathrm{Def}(\alpha)$} (Obstr);
  \draw[CRMBlue] (6.80,3.82) -- (7.02,3.82) -- (7.02,4.02);
  \node[font=\tiny,color=CRMBlue,align=center] at (7.95,2.30) {$\mathbf{Mod}(A) \cap \mathrm{Def}(\alpha)^{-1}(0) = \mathbf{Mod}(B)$ \\ ($A + \mathcal{R}(\mathcal{E}) \iff B$)};

  % Divider line
  \draw[draw=CRMBlue!35,dashed] (-0.15,1.82) -- (10.35,1.82);

  % Panel (b): Bottom — Tri-Modal Orientations
  \node[font=\scriptsize\bfseries,color=CRMBlue,anchor=west] at (-0.15,1.52) {(b) Derived Tri-Modal Factorization in $\mathbf{Fib}$};
  \node[doc,minimum width=3.15cm,inner sep=3.5pt] (M1) at (1.55,0.62) {\textbf{Mode~(i): Horizontal}\\$\|\Pi_f(p)\|_{-1} \simeq 1$\\Section inhabitation $\mathrm{Res}_{\mathrm{sec}}$};
  \node[doc,minimum width=3.15cm,inner sep=3.5pt] (M2) at (5.10,0.62) {\textbf{Mode~(ii): $\mathcal{J}$-Completion}\\$\widehat{\mathcal{E}}_{\mathcal{J}} \setminus \mathcal{E},\; \varprojlim^1$\\Tower phantom $\mathrm{Res}_{\mathrm{lim}}$};
  \node[doc,minimum width=3.15cm,inner sep=3.5pt] (M3) at (8.65,0.62) {\textbf{Mode~(iii): Vertical}\\$\mathrm{hofib}_B(U)$ over $B \in \mathcal{E}$\\Forgetful lift $\mathrm{Res}_{\mathrm{lift}}$};
\end{tikzpicture}
\caption{Comparison square $\alpha$, Beck--Chevalley mate $\mu_\alpha$, derived tri-modal orientations of the Universal Doctrinal Square Defect $\mathrm{Def}(\alpha)$, and pullback extraction in $\mathbf{Topos}$.}\label{fig:comparison-triple-residuum}
\end{figure}

Placing general comparison squares and their derived defects on equal footing organizes Categorical Reverse Mathematics along \textbf{Five Foundational Coordinate Axes}:

\begin{enumerate}[label=\textbf{Axis \Roman*.},leftmargin=2.6em]
  \item \textbf{Choice / Splitting ($\mathrm{EAC}, \mathrm{IAC}, \mathrm{IDC}_\omega$):} Mode~(i) horizontal base change over $\mathcal{D}_0 = \mathbf{Topos}$ ($\mathbf{GTop}$) vs.\ $\mathcal{D}_1 = \mathbf{Topos}_{\mathrm{AC}}$, testing $(-1)$-connectedness of section spaces $\|\Pi_f(p)\|_{-1} \simeq 1$ on covers $E \twoheadrightarrow B$ and $\mathrm{im}(\mathrm{InvStr}) \hookrightarrow E_{\mathrm{rel}}$. Residua include non-abelian torsors $H^1(\mathcal{E}/Y; G)$ ($\mathcal{T}_{\mathrm{odd}}$ in $\mathbf{Nom}$), definable $\varprojlim^1_{\mathrm{def}}$, and the \texttt{ICaTT} witness-stream gap (Chapter~\ref{ch:realizability-sites}, Theorem~\ref{thm:coinductive-stream-choice}).
  \item \textbf{Constructive Logic ($\mathrm{LEM}, \mathrm{WLEM}$):} Modes~(i) and~(ii) across $\mathcal{D}_0 = \mathbf{HeytingCat} \supseteq \mathbf{BooleanCat} = \mathcal{D}_1$, comparing classical truth values to $\Omega$ via $1+1 \hookrightarrow \Omega$ and $\mathrm{Sub}_{\mathrm{dec}}(X) \hookrightarrow \mathrm{Sub}(X)$. Residua isolate boundary sheaves $\partial X$, $\neg\neg$-dense nuclei, and Dedekind--MacNeille monicity failure.
  \item \textbf{Arithmetic Infinity ($\mathrm{NNO}$, $\mu F$ vs.\ $\nu F$):} Splitting cleanly into Mode~(ii) pro-initial completion $\widehat{\mathbb{N}}_{\mathrm{lasso}} \notin \mathbf{FinSet}$ (trivialized by adjoining an $\mathrm{NNO}$ in $\mathbf{Pretopos}_{\mathrm{NNO}}$) and finitary/well-founded stabilization $\mathrm{Res}_{\mathrm{Ded}}(X),\, \prod_{\mathbb{N}} A / \bigoplus_{\mathbb{N}} A$ (trivialized by $\neg\mathsf{Inf}$ in $\mathbf{FinSet}$), together with step-product phantoms $X_\bullet \in RL(\mathcal{Q}) \setminus \mathcal{Q}$ (Theorem~\ref{thm:step-product-counterexample}).
  \item \textbf{Characteristic / $p$-Torsion:} Mode~(iii) vertical forgetful fiber along $U : \delta\text{-}\mathbf{Ring}_p \to \mathbf{FrobLiftRing}_p$ over base Frobenius-lift rings $(R, \phi)$ via Joyal's~\cite{Joyal85} identity $\phi(x) = x^p + p\,\delta(x)$, governed over $p$-torsion bases by a three-stage lift tower in $\mathrm{hofib}_{(R,\phi)}(U)$ (Proposition~\ref{prop:delta-three-stage-tower}), paired with Mode~(ii) colimit completion $A \to A_{\mathrm{perf}}$ vs.\ the Nygaard filtration separating relative perfection from general prisms.
  \item \textbf{Directional Invertibility ($\pi_{\le 0}$):} Mode~(ii) completion $\mathcal{C} \hookrightarrow RL(\mathcal{C})$ across $\mathrm{Cat}^R_{(\infty,\infty)}$ alongside Mode~(i)/(iii) comparison $\mathrm{id} \Rightarrow \pi_{\le 0}$ comparing directed stabilization $\Omega_{\mathrm{dir}}$ to groupoidification. Residua recover stable homotopy stems $\pi_*^s(\mathbb{S}_{\mathrm{top}})$ from $\mathbf{B}^\infty(\mathbf{FinSet}^{\simeq})$, step-product phantoms ($\mathcal{Q}$), and directed facet cycles $F_n$ ($n \ge 3$, Theorem~\ref{thm:directed-cubical-amalgamation}).
\end{enumerate}

\begin{remark}[Grounding the Five Axes in Five Anchor Witnesses]\label{rem:anchor-witnesses}
To prevent the residuum calculus from remaining abstract formalism, every coordinate axis in Categorical Reverse Mathematics is calibrated against five concrete, self-contained \emph{anchor witnesses} previewed in the Preface:
\begin{itemize}[leftmargin=1.5em]
  \item \textbf{The Parity Torsor $\mathcal{T}_{\mathrm{odd}}$ (Axes I \& II):} Inside Schanuel's nominal topos $\mathbf{Nom} = \mathbf{Cont}(\mathrm{Sym}(\mathbb{N}))$, the set of odd-sized finite subsets of $\mathbb{N}$ has no global invariant element ($\Gamma(\mathcal{T}_{\mathrm{odd}}) = \emptyset$), yet splits locally as soon as any single atom $a \in \mathbb{N}$ is named. It represents a non-zero cohomology class $[\mathcal{T}_{\mathrm{odd}}] = \partial(1) \neq 0 \in H^1(\mathbf{Nom}; \mathrm{EvenFinset})$, witnessing the exact obstruction separating $\mathrm{IAC}$ from $\mathrm{AC}$ in a Boolean, two-valued universe despite pro-discrete blindness ($H^1(\mathbf{Nom}; \mathbb{Z}/2\mathbb{Z}) = 0$).
  \item \textbf{The $2$-Torsion Ring $R_2 = \mathbb{Z}[\eta]/(2\eta, \eta^2)$ (Axis IV):} In the doctrine of $\delta$-rings at $p = 2$, $R_2$ is isomorphic as an additive abelian group to $\mathbb{Z} \oplus \mathbb{F}_2$ (though \emph{not} as commutative rings to $\mathbb{Z} \times \mathbb{F}_2$, because $\eta \neq 0$ while $\eta^2 = 0$), and $2\eta = 0$ blocks naive division by $2$. Joyal's binomial identity $\delta(x + y) = \delta(x) + \delta(y) - xy$ reveals two distinct $\delta$-structures $\delta_0(\eta) = 0$ and $\delta_1(\eta) = \eta$ lifting the \textbf{augmentation Frobenius endomorphism} $\phi(a + b\eta) = a$ (where $\phi(\eta) = 0$), proving that $\mathrm{hofib}_{(R,\phi)}(U)$ over $p$-torsion rings is a non-trivial $\mathrm{Der}_\phi$-torsor, while the torsion-free dual-integer ring \texttt{DualInt} ($\mathbb{Z}[\varepsilon]/(\varepsilon^2)$) witnesses the collapse of the Nygaard filtration under relative colimit perfection.
  \item \textbf{The Step-Product $\mathcal{Q}$ on $[1] = \{0 < 1\}$ (Axes III \& V):} In weak $(\infty, \infty)$-categories, build the finite gaunt $k$-categories $\mathcal{A}_k = \mathcal{M}_{k+1}$ with two objects $\{0_k, 1_k\}$ such that $\mathrm{Hom}(0_k, 0_k) \simeq \mathrm{Hom}(0_k, 1_k) \simeq \mathrm{Hom}(1_k, 0_k) \simeq *$ whereas the diagonal hom-object $\mathrm{Hom}(1_k, 1_k) = \mathcal{M}_k \not\simeq *$ carries non-trivial $k$-dimensional structure. Inside $\prod_{k \ge 1} \mathcal{A}_k$, the restricted-tail step-product $\mathcal{Q}$ spanned by finite-prefix indicator sequences $e^{(m)} = (1_1, \dots, 1_m, 0_{m+1}, \dots)$ satisfies $\mathrm{Hom}_{\mathcal{Q}}(e^{(m)}, e^{(r)}) \simeq \prod_{k=1}^{\min(m, r)} \mathcal{M}_k$ of rank $\min(m, r)$, yet its $n$-truncations $y_n := \ell_n(e^{(n-1)}) \in \pi_{\le n}(\mathcal{Q})$ assemble into a pure Postnikov phantom $X_\bullet \in RL(\mathcal{Q}) \setminus \mathcal{Q}$, separating $\mathrm{Cat}^\omega_{(\infty,\infty)}$ from $\mathrm{Cat}^L_{(\infty,\infty)}$ without topological spheres (Theorem~\ref{thm:step-product-counterexample}).
  \item \textbf{The Facet-Supported Directed Triangle $F_3 = \partial\Delta^2$ (Axis V \& Site Geometry):} In directed combinatorial geometry, gluing two acyclic directed complexes along a flow-convex overlap preserves acyclicity ($n = 2$), whereas gluing around the boundary of $F_3 = \partial\Delta^2$ or a directed $n$-cube fails for all $n \ge 3$ by closing an irreversible directed cycle (Theorem~\ref{thm:directed-cubical-amalgamation}).
  \item \textbf{The Cantor Smooth-End Quotient $\mathcal{Q}_{E_0} = 2^{\mathbb{N}}/E_0$ and $\mathrm{UHF}(2^\infty)$ (Axes I, III, \& V):} In light condensed sets $\mathbf{Sh}(\mathrm{Prof}_{\aleph_0})$, the Vitali/Hamel tail quotient $\mathcal{Q}_{E_0} = \prod_{\mathbb{N}} \mathbb{F}_2 / \bigoplus_{\mathbb{N}} \mathbb{F}_2$ sits suspended between Freedman's topological collapse ($*$), discontinuous $\mathsf{ZFC}$ splitting ($\mathbb{R}$), and trivial Polish reflection ($\pi_0^{\mathrm{Pol}}(\mathcal{Q}_{E_0}) \simeq 1$), precipitating a non-splitting condensed class $[q] \in \mathrm{Ext}^1_{\mathrm{CondAb}}(\mathcal{Q}_{E_0}, \bigoplus_{\mathbb{N}}\mathbb{F}_2)$ alongside the Connes noncommutative $C^*$-modulus $C^*(\mathcal{G}_{E_0}) \cong \mathrm{UHF}(2^\infty)$ ($K_0 \cong \mathbb{Z}[\tfrac{1}{2}]$) and its Pauli-$Z$ $E_0^{\mathrm{even}}$ parity grading (Theorem~\ref{thm:double-solvent-r4} and Section~\ref{subsec:uhf-cstar-e0}).\qedhere
\end{itemize}
\end{remark}

\section{The Admissibility Problem, Doctrinal Philosophies, and Orthogonal Invariants}\label{sec:admissibility}
\sectionepigraph{``There are more things in heaven and earth, Horatio, than are dreamt of in your philosophy.''}{William Shakespeare, \emph{Hamlet}, Act~I, Scene~5 (Second Quarto)}

At the heart of Categorical Reverse Mathematics lies a subtle methodological hazard that we call the \emph{Triviality Trap}\index{Triviality Trap}. If we place no restrictions on what constitutes a valid condition $\mathcal{R}(\mathcal{E})$, one could simply define $\mathcal{R}(\mathcal{E})$ to be the internal material implication $A \to B$ (or the external implication $(\mathcal{E} \models A) \implies (\mathcal{E} \models B)$), or set $\mathrm{Def}(\alpha)(X) := \kappa_X$ literally without structural reduction. Syntactically, $A \wedge (A \to B) \leftrightarrow B$ is a tautology, but it provides zero structural insight into why $B$ fails in general models of $A$.

To prevent tautological restatements, $\mathrm{CRM}$ enforces an \emph{Admissibility Criterion}\index{Admissibility Criterion}: the residuum must be formulated purely in terms of \emph{orthogonal structural invariants} of the ambient doctrine $X \in \mathcal{D}_0$ rather than the specific sentences or instances appearing in $\mathcal{D}_1$. In a general elementary or Grothendieck topos $\mathcal{E}$, admissible invariants fall into four foundational families (supplemented on Axes~III--V by initial-algebra/tail quotients, $\delta$-torsor and Nygaard invariants, and Postnikov/amalgamation obstructions):

\begin{itemize}[leftmargin=1.5em]

  \item \textbf{Subobject Classifier Geometry:} Order-theoretic and algebraic properties of the truth-value object $\Omega$, such as Booleanness ($\Omega \simeq 1+1$), Stone/De Morgan structure ($\neg\phi \vee \neg\neg\phi = \top$), G\"odel--Dummett linearity, or internal compactness of $\Omega$.

  \item \textbf{Generator and Exactness Conditions:} Projectivity and generating properties of canonical objects, such as whether the terminal object $1_{\mathcal{E}}$ is \emph{externally projective} ($\Gamma(X) \neq \emptyset$ for every epimorphism $X \twoheadrightarrow 1_{\mathcal{E}}$), whether $\mathcal{E}$ is \emph{two-valued} ($\mathrm{Sub}_{\mathcal{E}}(1_{\mathcal{E}}) \simeq \{0, 1\}$) or \emph{well-pointed} (two-valued with $1_{\mathcal{E}}$ an externally projective generator), whether subobjects of $1_{\mathcal{E}}$ generate $\mathcal{E}$ (localicity), or whether $\mathcal{E}$ possesses enough projectives ($\mathrm{COSHEP}$). Crucially, two-valuedness does \emph{not} imply external projectivity of $1_{\mathcal{E}}$: in any atomic permutation topos $\mathbf{Cont}(\mathcal{G})$ possessing a proper open subgroup $U \subsetneq \mathcal{G}$ (such as Schanuel's nominal topos $\mathbf{Nom}$), $\mathrm{Sub}(1_{\mathcal{E}}) = \{0, 1\}$, yet the canonical coset epimorphism $\mathcal{G}/U \twoheadrightarrow 1_{\mathcal{E}}$ has empty global sections ($\Gamma(\mathcal{G}/U) = (\mathcal{G}/U)^{\mathcal{G}} = \emptyset$).

  \item \textbf{Site and Locale Topology:} For a Grothendieck sheaf topos $\mathcal{E} \simeq \mathbf{Sh}(\mathbb{C}, J)$, topological invariants of the underlying site or locale, including localicity, coherence (compactness of coverings), zero-dimensionality (refinement by clopen partitions), and the countable chain condition ($\omega$-presentability).

  \item \textbf{Cohomological Vanishing Obstructions:} Homological and sheaf-cohomological invariants, such as the vanishing of $\mathrm{Ext}^1_{\mathcal{E}}$ or first sheaf cohomology $H^1(\mathcal{E}, G) = 0$ for canonical classes of internal structure groups $G$, diagnosing when locally split structures glue into global sections.

\end{itemize}

Rather than leaving admissibility informal, we internalize the interaction between universes and invariant probes as a precise mathematical structure inspired by Shulman's~\cite{Shulman20} $2$-Chu-Dialectica pairing.

\begin{definition}[Doctrinal Philosophy, Pareto Frontier, and Epistemic Regimes]\label{def:doctrinal-philosophy}%
\uses{def:comparison-triple,def:residuum,rem:framed-equipments}%
A \textbf{(doctrinal) philosophy}\index{Doctrinal philosophy} is a triple $\mathfrak{P} = (\mathbb{D}, \mathcal{P}, \langle -, - \rangle)$ (\texttt{DoctrinalPhilosophy D P}) consisting of a locally small, fibrant framed doctrinal equipment $\mathbb{D}$~\cite{Shulman08,HHLN23}, a stratified suite $\mathcal{P}$ of structural probes, and a pairing $\langle X, P \rangle \in \mathcal{V}$ ($X \models P$). The induced Galois connection
\[
  \begin{aligned}
    \mathrm{Th}_{\mathcal{P}}(\mathcal{U}) &:=\; \{P \in \mathcal{P} \mid \forall X \in \mathcal{U},\; X \models P\}, \\
    \mathrm{Ext}_{\mathcal{P}}(Q) &:=\; \{X \in \mathbb{D} \mid \forall P \in Q,\; X \models P\}
  \end{aligned}
\]
generates a Moore closure $\mathcal{M}(\mathcal{P}) := \mathrm{Ext}_{\mathcal{P}}(\mathrm{Th}_{\mathcal{P}}(-))$ on subsets of $\mathbb{D}$ (\texttt{mooreClosureOperator}). Given a hypothesis class $K \subseteq \mathbb{D}$ and target class $T \subseteq \mathbb{D}$ with $K \not\subseteq T$ (such as $K = \mathbf{Mod}(\mathrm{IAC})$ and $T = \mathbf{Mod}(\mathrm{AC})$), the \textbf{Pareto Residuum Frontier}\index{Pareto residuum frontier} $\mathrm{Res}_{\mathcal{P}}^{\min}(K \Rightarrow T)$ consists of all inclusion-maximal (deductively weakest) $\mathcal{M}(\mathcal{P})$-closed extents $E$ satisfying $K \cap E \subseteq T$. A philosophy $\mathfrak{P}$ natively calibrates five epistemic regimes:
\begin{enumerate}[label=\textup{(\arabic*)},leftmargin=2.0em]
  \item \textbf{Dogma (Triviality Trap):} $T$ (or $K \Rightarrow T$) is syntactically encoded as a tautological probe in $\mathcal{P}$ without weaker witness separation, collapsing $\mathrm{Res}_{\mathcal{P}}^{\min}(K \Rightarrow T) = \{\mathrm{Cl}_{\mathcal{P}}(T)\}$.
  \item \textbf{Blindness:}\index{Doctrinal blindness} $\mathrm{Res}_{\mathcal{P}}^{\min}(K \Rightarrow T) = \emptyset$ despite non-trivial defects (for instance, pro-discrete finite cohomology $H^1(-;\mathbb{Z}/2\mathbb{Z})$ is blind to $\mathcal{T}_{\mathrm{odd}}$ in $\mathbf{Nom}$, and Polish $\pi_0^{\mathrm{Pol}}$ is blind to $\mathcal{Q}_{E_0}$ and $\mathrm{UHF}(2^\infty)$).
  \item \textbf{Schism:} $|\mathrm{Res}_{\mathcal{P}}^{\min}(K \Rightarrow T)| \ge 2$ forms a genuine antichain of incomparable geometric trade-offs.
  \item \textbf{Heresy:} A concrete witness $W \in E_1 \setminus E_2$ separating two candidate residua $E_1, E_2$.
  \item \textbf{Conversions and Bi-interpretation:} A mixed-variance Chu morphism~\cite{Shulman20} $(F, G) : \mathfrak{P}_1 \to \mathfrak{P}_2$ with adjoint pairing $\langle F(X), P_2 \rangle \simeq \langle X, G(P_2) \rangle$; $(F, G)$ is a \textbf{bi-interpretation} when its induced closure maps identify $\mathcal{M}(\mathcal{P}_1) \simeq \mathcal{M}(\mathcal{P}_2)$.
\end{enumerate}
\end{definition}

Could negative stability rules alone (such as closure under binary products, open subtopoi, or Lawvere--Tierney localizations) filter out Dogma without specifying an explicit probe suite $\mathcal{P}$? The following proposition settles that question negatively while establishing an exact finite Pareto Schism inside Choice Geometry.

\begin{proposition}[Inquisitor's Impossibility and Finite Pareto Schism]\label{prop:anti-casuistry-schism}%
\lean{HomotopyCRM.DoctrinalPhilosophy, HomotopyCRM.inquisitor_impossibility_casuistry, HomotopyCRM.monoid_proj1_iff_rightZero, HomotopyCRM.m1_projSub_and_not_localic, HomotopyCRM.vPoset_localic_not_proj1, HomotopyCRM.diamond_proj1_not_projSub, HomotopyCRM.freydCover_dogma_banishment, HomotopyCRM.pareto_incomparable_projSub_localic, HomotopyCRM.M1SliceLTT.m1_slice_lt_count}%
\leanok%
\uses{def:doctrinal-philosophy}%
Let $\mathfrak{P} = (\mathbb{D}, \mathcal{P}, \langle -, - \rangle)$ be a doctrinal philosophy on a product-complete meet-semilattice of doctrines.
\begin{enumerate}[label=\textup{(\roman*)},leftmargin=2.0em]
  \item \textbf{Inquisitor's Impossibility (Anti-Casuistry Lemma):} For any monoid $\Sigma$ of product-preserving doctrinal endomorphisms (\texttt{ProductEndo}) and any product-closed target class $T \subseteq \mathbb{D}$, its modal $S4$-interior $\Box_\Sigma(T) := \{X \in \mathbb{D} \mid \forall \sigma \in \Sigma,\; \sigma(X) \in T\}$ is itself product-closed, $\Sigma$-invariant, contained in $T$, and uniquely maximal among all $\Sigma$-closed subclasses of $T$ (\texttt{inquisitor\_\allowbreak impossibility\_\allowbreak casuistry}). Hence negative stability laws alone cannot banish Dogma without a kataphatic probe suite $\mathcal{P}$.
  \item \textbf{Open-Subtopos Stability and Supports Split ($\mathrm{SS}$):} On presheaves $\mathbf{Set}^M$ over a monoid $M$, external projectivity of $1$ ($\mathrm{Proj}_1$) holds constructively if and only if $M$ has a right zero $e \in M$ (\texttt{monoid\_\allowbreak proj1\_\allowbreak iff\_\allowbreak rightZero}, cf.~\cite{HemelaerRogers21}). Requiring open-subtopos stability banishes raw $\mathrm{Proj}_1$-dogma via the Freyd cover $(\mathbf{B}(\mathbb{Z}/2\mathbb{Z}))_\bot$: $\mathrm{Proj}_1$ holds on $(\mathbf{B}(\mathbb{Z}/2\mathbb{Z}))_\bot$ (using $\bot$ as a right zero) yet fails on its open subtopos $\mathbf{Set}^{\mathbb{Z}/2\mathbb{Z}}$ (\texttt{freydCover\_\allowbreak dogma\_\allowbreak banishment}), tightening $\mathrm{Proj}_1$ to hereditary subterminal projectivity $\mathrm{Proj}_{\mathrm{sub}} = \mathrm{SS}$ (Freyd--{\v{S}}{\v{c}}edrov's~\cite{FreydScedrov90} \emph{Supports Split}), separated constructively from $\mathrm{Proj}_1$ on the $4$-element diamond poset $D_4$ (\texttt{diamond\_\allowbreak proj1\_\allowbreak not\_\allowbreak projSub}).
  \item \textbf{Constructive Finite Pareto Schism and Slice Count:} Evaluated across the four-topos test suite $\{\mathbf{Set}, \mathbf{Set}^{\mathbb{Z}/2\mathbb{Z}}, \mathbf{Set}^{\mathcal{M}_1}, \mathbf{Set}^{V^{\mathrm{op}}}\}$ with $K = \mathrm{Ext}(\{\mathrm{IAC}\})$ and target $T = \mathbf{Mod}(\mathrm{AC}) = K \cap \mathrm{Ext}(\{\mathrm{Proj}_{\mathrm{sub}}\})$ (equivalently $K \cap E \subseteq \mathrm{Ext}(\{\mathrm{Proj}_{\mathrm{sub}}\})$), the Pareto frontier $\mathrm{Res}_{\mathcal{P}}^{\min}(K \Rightarrow T)$ splits constructively into the incomparable pair of extents $\mathrm{Ext}(\{\mathrm{Proj}_{\mathrm{sub}}\})$ and $\mathrm{Ext}(\{\mathrm{Localic}\})$ (\texttt{pareto\_\allowbreak incomparable\_\allowbreak projSub\_\allowbreak localic}):
  \begin{itemize}[leftmargin=1.4em]
    \item the $2$-element absorbing base monoid $\mathcal{M}_1 = (\{e_1, 1_1\}, \wedge)$ of Chapter~\ref{ch:higher-crm} (\texttt{AbsorbingMonoidCarrier 1}) satisfies $\mathrm{Proj}_{\mathrm{sub}}$ via $u \mapsto e_1 \cdot x_0$ while failing localicity since $e_1 \neq 1_1$ (\texttt{m1\_\allowbreak projSub\_\allowbreak and\_\allowbreak not\_\allowbreak localic}); and
    \item the $3$-element V-poset $V = \{a > c < b\}$ is localic by Yoneda generation yet fails $\mathrm{Proj}_1$ (hence fails $\mathrm{Proj}_{\mathrm{sub}}$) on $\mathbf{y}(a) \amalg \mathbf{y}(b)$ (\texttt{vPoset\_\allowbreak localic\_\allowbreak not\_\allowbreak proj1}).
  \end{itemize}
  Moreover, over any finite idempotent $\mathcal{M}_1$-set $(X, \rho)$ with $k$ fixed points and $m$ non-fixed points, the slice topos $\mathbf{Set}^{\mathcal{M}_1}/(X, \rho)$ admits exactly $3^k 2^m$ Lawvere--Tierney topologies, splitting cleanly into $2^{k+m}$ localic slices and $(3^k - 2^k) 2^m$ non-localic (and non-Boolean) slices (\texttt{M1SliceLTT.\allowbreak m1\_\allowbreak slice\_\allowbreak lt\_\allowbreak count}).
\end{enumerate}
\end{proposition}

\begin{proof}
For~(i), $\mathbf{1} \in \Sigma$ gives $\Box_\Sigma(T) \subseteq T$; closure of $\Sigma$ under composition gives $\sigma(\tau(X)) \in T \implies \tau(X) \in \Box_\Sigma(T)$ ($\Sigma$-invariance); and $\sigma(\top) = \top$, $\sigma(X \sqcap Y) = \sigma(X) \sqcap \sigma(Y)$ transport product closure from $T$ to $\Box_\Sigma(T)$, while every $\Sigma$-invariant subclass $C \subseteq T$ lies inside $\Box_\Sigma(T)$. For~(ii), the right-zero characterization ($\mathbf{Set}^M \models \mathrm{Proj}_1 \iff \exists e \in M,\, \forall m,\, me = e$) is immediate by testing the free cyclic $M$-set $M$ at $1 \in M$ and conversely mapping any inhabited $X$ to $e \cdot x_0$. Adjoining a zero $\bot$ to $\mathbb{Z}/2\mathbb{Z}$ makes $\bot$ a right zero while its open complement $\mathbf{Set}^{\mathbb{Z}/2\mathbb{Z}}$ has no fixed points on the free $\mathbb{Z}/2\mathbb{Z}$-torsor; on the diamond $D_4 = \{\bot < l, r < \top\}$, the top element $\top$ forces $\mathrm{Proj}_1$, whereas over the open subterminal $U = \{\bot, l, r\}$ the presheaf $\mathbf{y}(l) \amalg \mathbf{y}(r)$ splits locally without gluing at $\bot$. For~(iii), on $\mathcal{M}_1$ every surjection onto an inhabited subterminal $1$ picks $e_1 \cdot x_0$ constructively while non-triviality of $\mathcal{M}_1$ blocks subterminal generation; on $V$, Yoneda principal sieves $\downarrow p$ generate $\mathbf{Set}^{V^{\mathrm{op}}}$ choice-freely while compatibility on $a$ and $b$ diverges at $c$; thus neither $\mathrm{Proj}_{\mathrm{sub}}$ nor $\mathrm{Localic}$ lies inside the other, yet $\mathrm{IAC}$ holds on finite presheaves iff all morphisms are invertible, leaving only $\mathbf{Set}$ in $\mathrm{Ext}(\{\mathrm{IAC}, \mathrm{Localic}\}) = \mathrm{Ext}(\{\mathrm{IAC}, \mathrm{Proj}_{\mathrm{sub}}\})$. Finally, decomposing fiberwise modal closures $j_x : \Omega_{\mathcal{M}_1} \to \Omega_{\mathcal{M}_1}$ on $\Omega_{\mathcal{M}_1} = \{\emptyset < \{e_1\} < 1\}$ under $e_1$-equivariance gives $3$ maps on each fixed point ($\mathrm{id}, j_{\mathrm{loc}}, \top$, with two localic choices killing $\mathrm{id}$) and $2$ extensions on each non-fixed point once $j_{\rho(x)}(\emptyset)$ is fixed, yielding $3^k 2^m = 2^{k+m} + (3^k - 2^k)2^m$.
\end{proof}

\section{Paradigmatic Case Study: The Choice-Geometry Residuum Atlas}

The archetypal testing ground for Axis~I of Categorical Reverse Mathematics is the \emph{Choice-Geometry Residuum Atlas} (developed in full in Chapter~\ref{ch:choice-atlas}): given a weak form of choice $A$, what is the weakest structural condition $\mathcal{R}(\mathcal{E})$ on a topos $\mathcal{E}$ such that $A + \mathcal{R}(\mathcal{E}) \iff \mathrm{AC}$?

In any elementary topos $\mathcal{E}$, the \emph{External Axiom of Choice} ($\mathrm{AC}$) states that every epimorphism $e: X \twoheadrightarrow Y$ admits a global section $s: Y \to X$, whereas the \emph{Internal Axiom of Choice} ($\mathrm{IAC}$) asserts that dependent product functors $\Pi_f : \mathcal{E}/X \to \mathcal{E}/Y$ (equivalently exponentials $(-)^Z$) preserve epimorphisms~\cite{Johnstone02}, and \emph{Local Choice} splits epimorphisms locally along covering sieves in every Grothendieck topos $\mathbf{Sh}(\mathbb{C}, J)$. Evaluating $(A \Rightarrow_{\mathbf{Topos}} \mathrm{AC})$ isolates five benchmark transitions:

\begin{enumerate}[label=\textbf{\arabic*.},leftmargin=1.8em]
  \item \textbf{Trivial Base ($A = \top \Rightarrow \mathrm{AC}$):} By Diaconescu's theorem~\cite{Diaconescu75}, $\mathrm{IAC}$ forces Booleanness ($\Omega \simeq 1+1$), yet Freyd's two-valued Boolean topos $\mathcal{F} = \mathbf{Sh}_{\neg\neg}(\mathbb{A})$ with $\mathrm{NNO}$ still refutes $\mathrm{AC}$~\cite{Freyd72}. Over $\mathbf{Set}_{\mathrm{AC}}$, the Barr--Diaconescu--Freyd theorem~\cite{Barr74,Diaconescu75,Freyd72} (\cite[\S D4.5]{Johnstone02}) identifies $\mathcal{E} \models \mathrm{AC} \iff \mathcal{E} \simeq \mathbf{Sh}(B)$, isolating \emph{Booleanness plus localicity} as $(\top \Rightarrow_{\mathbf{GTop}} \mathrm{AC})$.
  \item \textbf{Internal to External Choice ($\mathrm{IAC} \Rightarrow \mathrm{AC}$):} Bridging internal splittings to global sections $\Gamma(X) = \mathrm{Hom}_{\mathcal{E}}(1_{\mathcal{E}}, X)$ isolates \emph{external projectivity of $1_{\mathcal{E}}$}: $\mathrm{IAC} + (1_{\mathcal{E}} \text{ externally projective}) \iff \mathrm{AC}$ (Theorem~\ref{thm:iac-to-ac}). In two-valued atomic permutation topoi $\mathbf{Cont}(\mathcal{G})$, any proper open subgroup $U \subsetneq \mathcal{G}$ forces $\Gamma(\mathcal{G}/U) = \emptyset$, separating two-valuedness from external projectivity of $1_{\mathcal{E}}$ (Chapter~\ref{ch:realizability-sites}).
  \item \textbf{Enough Projectives ($\mathrm{COSHEP} \Rightarrow \mathrm{AC}$):} Combining projective covers $\mathrm{COSHEP}$ with Blass's~\cite{Blass79} Small Violations of Choice faces both non-projectivity of $1$ in $\mathbf{Set}^{\mathbb{Z}/2\mathbb{Z}}$ and non-projectivity of unsplit quotients; Theorem~\ref{thm:coshep-svc} establishes $\mathrm{COSHEP} + (\mathrm{SVC}_{\mathrm{ret}} \wedge \text{Booleanness} \wedge 1_{\mathcal{E}} \text{ externally projective}) \iff \mathrm{AC}$.
  \item \textbf{Boolean Prime Ideal Theorem ($\mathrm{IBPI} \Rightarrow \mathrm{IAC}$):} Where Deligne's theorem~\cite[Exp.~VI, Appendice]{Deligne72} equates external $\mathrm{BPI}$ to coherent spatiality and $\mathrm{IBPI}$ yields ultrafilters on internal Boolean algebras~\cite{Johnstone82,Johnstone02}, Conjecture~\ref{conj:bpi-ac} posits Booleanness plus \emph{\'etale clopen-partition refinement} of coherent epimorphic covers over $\mathbf{Set}_{\mathrm{AC}}$.
  \item \textbf{Dependent Choice ($\mathrm{IDC} \Rightarrow \mathrm{IAC}$) and Higher Choice ($\mathrm{AC}_n \Rightarrow \mathrm{AC}_\infty$):} While $\mathbf{Sh}(B) \models \mathrm{AC}$ over $\mathbf{Set}_{\mathrm{AC}}$ regardless of $\mathrm{CCC}$~\cite{Barr74,Johnstone02} (which enters over $\mathcal{S} \models \mathrm{DC}$ in Theorem~\ref{thm:bdf-deformed}), Conjecture~\ref{conj:dc-ac} isolates \emph{\'etale $\omega$-presentability of localic groupoid isotropy}~\cite{Johnstone02,Kihara26} over $\mathbf{Set}_{\mathrm{AC}}$; in $(\infty, 1)$-topoi $\mathcal{X}$~\cite{Lurie09}, bridging $\mathrm{AC}_0$ to $\mathrm{AC}_\infty$ isolates $0$-truncated generation (Chapter~\ref{ch:higher-crm}, Theorem~\ref{thm:global-inf-choice}).
\end{enumerate}

\section{Beyond Choice: The Five-Axis Map and a Panorama of Structural Lattices}

Although the Axiom of Choice (Axis~I) provides the archetypal historical arena for $\mathrm{CRM}$, the residuum operator $\mathrm{Res}(\kappa)$ on $\mathrm{Dom}(\mathrm{Res})$ treats all five foundational axes symmetrically, as summarized in Table~\ref{tab:five-axis-master}.

\begin{table}[htbp]
\centering
\scriptsize
\setlength{\tabcolsep}{2.4pt}
\renewcommand{\arraystretch}{1.10}
\begin{tabularx}{\textwidth}{@{}P{0.14\textwidth}P{0.18\textwidth}YP{0.22\textwidth}P{0.20\textwidth}@{}}
\toprule
\textbf{Foundational Axis} & \textbf{Doctrine $\mathcal{D}_0$} & \textbf{Mode \& Comparison ($\alpha$)} & \textbf{Universal Defect $\mathrm{Def}(\alpha)$} & \textbf{Trivializing Property} \\
\midrule
\textbf{Axis I: Choice / Splitting} & $\mathbf{Topos} / \mathbf{GTop}$ & Mode~(i): $\|\Pi_f(p)\|_{-1} \simeq 1$ on $E \twoheadrightarrow B$; $\mathrm{im}(\mathrm{InvStr}) \hookrightarrow E_{\mathrm{rel}}$ & Non-abelian torsors $H^1(\mathcal{E}/Y; G)$, $\varprojlim^1_{\mathrm{def}}$, \texttt{ICaTT} gap & Booleanness + ($\mathrm{Localic}$ or $\mathrm{Proj}_{\mathrm{sub}}$ / $\mathrm{AC}$; $\mathcal{T}_{\mathrm{odd}}$ in $\mathbf{Nom}$) \\
\textbf{Axis II: Constructive Logic} & $\mathbf{HeytingCat}$ & Modes~(i)/(ii): $1+1 \hookrightarrow \Omega$; $\mathrm{Sub}_{\mathrm{dec}}(X) \hookrightarrow \mathrm{Sub}(X)$ & Boundary sheaves $\partial X$, $\neg\neg$-dense nuclei & Booleanness ($\mathrm{LEM}$), De Morgan ($\mathrm{WLEM}$) \\
\textbf{Axis III: Arithmetic Infinity} & $\mathbf{LexCat}$ ($\mathbf{FinSet}$) & Mode~(ii): pro-initial $\mu S$ vs.\ $\widehat{\mathbb{N}}_{\mathrm{lasso}}$; vs.\ $\mathrm{Res}_{\mathrm{Ded}}$, $\bigoplus_{\mathbb{N}} A \hookrightarrow \prod_{\mathbb{N}} A$ & Profinite lasso $\widehat{\mathbb{N}}_{\mathrm{lasso}} \twoheadrightarrow \widehat{\mathbb{Z}}$; Dedekind \& tail $\prod/\bigoplus$ gaps & $\mathrm{NNO}$ kills $\widehat{\mathbb{N}}_{\mathrm{lasso}}$; $\neg\mathsf{Inf}$ kills $\mathrm{Res}_{\mathrm{Ded}}$, $\prod/\bigoplus$ \\
\textbf{Axis IV: Characteristic / Torsion} & $\mathbf{FrobLift}\allowbreak\mathbf{Ring}_p$ & Mode~(iii): $U : \delta\text{-}\mathbf{Ring}_p \to \mathcal{D}_0$; Mode~(ii): $A \to A_{\mathrm{perf}}$ & $3$-stage tower $\mathrm{hofib}_{(R,\phi)}(U)$ ($\delta_0, \delta_1$), Nygaard $\mathcal{N}^{\ge i}A$ & $p$-torsion freeness / perfection ($R_2$ with $\phi(\eta)=0$; \texttt{DualInt}) \\
\textbf{Axis V: Directional Invertibility} & $\mathrm{Cat}^R_{(\infty,\infty)}$ & Mode~(ii) $\mathcal{C} \hookrightarrow RL(\mathcal{C})$ \& Mode~(i)/(iii) $\mathrm{id} \Rightarrow \pi_{\le 0}$ & Stable stems $\pi_*^s(\mathbb{S}_{\mathrm{top}})$, Postnikov phantoms & Arrow invertibility ($\infty$-groupoid collapse; $\mathcal{Q}$ and $F_3$) \\
\bottomrule
\end{tabularx}
\caption{The Five Foundational Coordinate Axes of Categorical Reverse Mathematics.}
\label{tab:five-axis-master}
\end{table}

Beyond Axis~I, the residuum calculus organizes classical and constructive bifurcations across seven major non-$\mathrm{AC}$ structural lattices---proved internally along Axes~IV--V in Chapters~\ref{ch:delta-rings} and~\ref{ch:higher-crm} and synthesized against external literature benchmarks in the Programmatic Panorama of Chapter~\ref{ch:panorama}:
\begin{itemize}[leftmargin=1.5em,noitemsep,topsep=2pt]
  \item \textbf{Internal Continuum \& Pointfree Spatiality (Axes~II--III, \S\ref{sec:panorama-continuum-spatial}):} Separating $\mathbb{R}_C \rightarrowtail \mathbb{R}_D \to \mathbb{R}_M$, $\varepsilon\text{-}\mathrm{IVT} \Rightarrow \mathrm{IVT}$, and Johnstone's localic Tychonoff / Mulvey--Pelletier Hahn--Banach theorems from spatial point-set counterparts via countable cover refinement, local connectedness, and internal spatiality~\cite{Bishop67,Ishihara06,Johnstone81,Johnstone82,MulveyPelletier91,Deligne72,Johnstone02}.
  \item \textbf{Intermediate Logics \& Homological Exactness (Axes~I--III, \S\ref{sec:panorama-omega-homological}):} Calibrating $\mathrm{WLEM}$ and G\"odel--Dummett linearity via $\Omega$-geometry~\cite{Johnstone82,Johnstone02}, and separating Bourn protomodularity~\cite{Bourn91}, Barr-exactness~\cite[\S\S A1.3--A1.4]{Johnstone02}, semi-abelian Snake Lemma naturality~\cite{JMT02}, and $\mathbf{AB5}$ vs.\ $\mathbf{AB4}^*$ ($\varprojlim^1$).
  \item \textbf{Predicative Pretopoi, $W$-Types, \& Smooth Cohesion (Axes~II--IV, Ch.~\ref{ch:delta-rings} \& \S\S\ref{sec:predicativity-wtypes}, \ref{sec:cohesive-lattice}):} Separating Maietti--Joyal Arithmetic Universes~\cite{Maietti05}, two-sorted Lawvere hyperdoctrines~\cite{Lawvere69} for the classical ``Big Five''~\cite{Friedman75,Simpson09} (below $\mathsf{CZF}/\mathsf{KP}$~\cite{Aczel78} and $\mathsf{Mac}$~\cite{Cole73}; Core Conjecture~5), Kock--Lawvere microlinearity and Lawvere--Schreiber cohesion ($\CohesiveQuad$)~\cite{Kock06,MoerdijkReyes91,Diaconescu75,Lawvere07,Schreiber13}, and Joyal $\delta$-ring $2$-torsion lift torsors ($R_2$ vs.\ \texttt{DualInt})~\cite{Joyal85}.
  \item \textbf{Directed $(\infty, \infty)$-Localizations \& $\mathrm{CatSp}$ (Axis~V, Ch.~\ref{ch:higher-crm} \& \S\ref{subsec:excision-catsp}):} Navigating Ozornova--Rovelli--Walde and Gepner--Heine's four-tier ladder $\mathrm{Cat}^L_{(\infty,\infty)} \subsetneq \mathrm{Cat}^\omega_{(\infty,\infty)} \subsetneq \mathrm{Cat}^{\mathrm{coind}}_{(\infty,\infty)} \subsetneq \mathrm{Cat}^R_{(\infty,\infty)}$~\cite{ORW26,GepnerHeine26}, step-product phantoms ($\mathcal{Q}$), and $\pi_{\le 0}$ group-completion birthing $\pi_*^s(\mathbb{S}_{\mathrm{top}})$~\cite{Heine26,Geiss26,BarrattPriddy72,Quillen71}.
\end{itemize}

\section{\texorpdfstring{Pure Axiom-of-Infinity Residua in $\mathbf{FinSet}$}{Pure Axiom-of-Infinity Residua in FinSet}}\label{sec:finset-infinity}%
\lean{HomotopyCRM.isEmpty_hom_of_fintype_to_succ_lasso, HomotopyCRM.profiniteLassoToZHat_surjective, HomotopyCRM.zHatToProfiniteLasso_neg_one_not_mem_range_nat, HomotopyCRM.dedekindResiduum_eq_empty_of_fintype, HomotopyCRM.nat_succ_mem_dedekindResiduum, HomotopyCRM.nominalDedekindResiduum_eq_empty, HomotopyCRM.not_hasFiniteSupportMap_nat_succ, HomotopyCRM.subsingleton_tailCokernel_of_fintype, HomotopyCRM.nontrivial_tailCokernel_nat, HomotopyCRM.chapter1_pure_infinity_residua_summary}%

Before relativizing Choice in Chapter~\ref{ch:choice-atlas}, the comparison-square framework of Definitions~\ref{def:comparison-triple}--\ref{def:residuum} reveals a phenomenon invisible inside choice-centric topos theory: \emph{pure structural residuums can crystallize even when $\mathrm{LEM}$ and External Choice ($\mathrm{AC}$) hold $100\%$}. Indeed, Proposition~\ref{prop:tri-modal-factorization} explains why two opposite Mode~(ii) completion polarities diverge along Axis~III: the pro-initial lasso phantom $\widehat{\mathbb{N}}_{\mathrm{lasso}}$ vanishes precisely upon adjoining a Natural Numbers Object ($\mathrm{NNO}$), whereas the Dedekind gap $\mathrm{Res}_{\mathrm{Ded}}(X)$ and biproduct-tail cokernel $\prod_{\mathbb{N}} A / \bigoplus_{\mathbb{N}} A$ vanish precisely under $\neg\mathsf{Inf}$ in $\mathbf{FinSet}$.

Consider the elementary topos $\mathcal{E} = \mathbf{FinSet} \in \mathbf{Topos}$ of finite sets and maps (internally modeling hereditarily finite sets $V_\omega$, i.e., $\mathsf{ZFC} - \mathsf{Inf} + \neg\mathsf{Inf}$). Inside $\mathbf{FinSet}$, the subobject classifier is strictly Boolean ($\Omega \simeq 1+1$), $1$ is externally projective, and every epimorphism $e : X \twoheadrightarrow Y$ admits a global section: \emph{every Choice residuum vanishes identically}. Yet because $\mathbf{FinSet}$ lacks a Natural Numbers Object ($\mathrm{NNO}$), comparing finitary constructions against infinitary doctrines along Axis~III precipitates three distinct, sharp Axiom-of-Infinity residuums---each machine-checked in \texttt{HomotopyCRM/\allowbreak InfinityResiduum.lean}\index{InfinityResiduum@\texttt{InfinityResiduum} Lean~4 module} (949 lines, 20 \texttt{\#print axioms} checks under the 4-Layer Definition Check Gate):

\begin{enumerate}[label=\textbf{\arabic*.},leftmargin=1.8em]

  \item \textbf{The Profinite Lasso Initial-Algebra Phantom ($\widehat{\mathbb{N}}_{\mathrm{lasso}} \twoheadrightarrow \widehat{\mathbb{Z}}$).}\index{Profinite lasso@Profinite lasso $\widehat{\mathbb{N}}_{\mathrm{lasso}}$} Consider the polynomial successor endofunctor $S(X) = 1 \amalg X$. Inside any doctrine possessing an $\mathrm{NNO}$ ($\mathbf{Pretopos}_{\mathrm{NNO}}$), the initial algebra $\mu S = \mathbb{N}$ exists. Inside $\mathbf{FinSet}$, every $S$-algebra $(X, x_0 : 1 \to X, s : X \to X)$ generated by $x_0$ must eventually collide, forming a finite \emph{lasso} $L_{p,q}$ ($p \ge 0$ stem length, $q \ge 1$ cycle period, formalized as \texttt{lassoDynSys}) satisfying $s^{p+q}(x_0) = s^p(x_0)$ and $|L_{p,q}| = p + q$; no finite carrier $A$ admits a pointed dynamical homomorphism into $L_{|A|, 1}$ (\texttt{isEmpty\_\allowbreak hom\_\allowbreak of\_\allowbreak fintype\_\allowbreak to\_\allowbreak succ\_\allowbreak lasso}). Passing to the profinite completion $\mathrm{Pro}(\mathbf{FinSet})$ along the coherent surjective transition system $L_{p', q'} \twoheadrightarrow L_{p, q}$ for $p \le p'$ and $q \mid q'$ (\texttt{lassoTransition}) precipitates the initial-algebra failure residuum as the profinite lasso inverse limit
  \[
    \mathrm{Res}_{\mathrm{init}}(S; \mathbf{FinSet}) \;=\; \widehat{\mathbb{N}}_{\mathrm{lasso}} \;:=\; \varprojlim_{(p,q)} L_{p,q} \;\notin\; \mathbf{FinSet},
  \]
  whose periodic quotient ($p = 0$) surjects onto the ring of \textbf{profinite integers} $\widehat{\mathbb{Z}} = \varprojlim_{q \ge 1} \mathbb{Z}/q\mathbb{Z}$\index{Profinite integers@Profinite integers $\widehat{\mathbb{Z}}$} (\texttt{profiniteLassoToZHat\_\allowbreak surjective}), with non-natural element $-1 \in \widehat{\mathbb{Z}}$ separated from $\mathbb{N}$ in \texttt{zHatToProfiniteLasso\_\allowbreak neg\_\allowbreak one\_\allowbreak not\_\allowbreak mem\_\allowbreak range\_\allowbreak nat}. Thus $\widehat{\mathbb{N}}_{\mathrm{lasso}}$ crystallizes externally in $\mathrm{Pro}(\mathbf{FinSet})$ as the exact obstruction between $\mathbf{FinSet}$ and an $\mathrm{NNO}$-bearing doctrine.

  \item \textbf{The Dedekind / Pigeonhole Residuum and the Nominal Crossing Point.} For any object $X$ in a Cartesian category, compare the subobject of automorphisms to monomorphic endomorphisms via $\kappa_{\mathrm{Ded}, X} : \mathrm{Iso}(X, X) \hookrightarrow \mathrm{Mono}(X, X)$, defining the Dedekind residuum\index{Dedekind residuum@Dedekind residuum $\mathrm{Res}_{\mathrm{Ded}}$}
  \[
    \mathrm{Res}_{\mathrm{Ded}}(X) \;:=\; \mathrm{Mono}(X, X) \setminus \mathrm{Iso}(X, X).
  \]
  By the Pigeonhole Principle, $\mathrm{Res}_{\mathrm{Ded}}(X) = \emptyset$ unconditionally for every $X \in \mathbf{FinSet}$ (\texttt{dedekindResiduum\_\allowbreak eq\_\allowbreak empty\_\allowbreak of\_\allowbreak fintype}), whereas $\mathrm{Res}_{\mathrm{Ded}}(\mathbb{N}) \neq \emptyset$ via $n \mapsto n + 1$ (\texttt{nat\_\allowbreak succ\_\allowbreak mem\_\allowbreak dedekindResiduum}) whenever an $\mathrm{NNO}$ is present. Remarkably, inside Schanuel's nominal topos $\mathbf{Nom} = \mathbf{Cont}(\mathrm{Sym}(\mathbb{N}))$ (Chapter~\ref{ch:realizability-sites}), the constant sheaf $\Delta(\mathbb{N})$ exists---so the Axiom of Infinity and an $\mathrm{NNO}$ hold---yet on the Kuratowski-infinite object of atoms $\mathbb{A}$, finite support forces every injective equivariant map $f : \mathbb{A} \to \mathbb{A}$ to be surjective (any omitted atom outside $\mathrm{supp}(f)$ can be transposed into the range, blocking \texttt{Nat.succ} via \texttt{not\_\allowbreak hasFiniteSupportMap\_\allowbreak nat\_\allowbreak succ}), yielding
  \[
    \mathrm{Mono}_{\mathbf{Nom}}(\mathbb{A}, \mathbb{A}) \;=\; \mathrm{Iso}_{\mathbf{Nom}}(\mathbb{A}, \mathbb{A}), \qquad \mathrm{Res}_{\mathrm{Ded}}(\mathbb{A}) \;=\; \emptyset
  \]
  (\texttt{nominalDedekindResiduum\_\allowbreak eq\_\allowbreak empty}). Thus $\mathbb{A} \in \mathbf{Nom}$ exhibits a pure crossing point between Axis~III and Axis~I, cleanly isolating $\mathsf{Inf} \wedge \neg\mathsf{AC}_\omega$ (an internally Kuratowski-infinite object that remains Dedekind-finite).

  \item \textbf{The Biproduct-Tail Cokernel ($\prod_{\mathbb{N}} A / \bigoplus_{\mathbb{N}} A$).} For any additive or abelian coefficient object $A \neq 0$ and index object $I$, consider the canonical inclusion of finitely supported sums into arbitrary products:
  \[
    \kappa_{I, A} : \bigoplus_{i \in I} A \;\hookrightarrow\; \prod_{i \in I} A.
  \]
  For every finite index set $I \in \mathbf{FinSet}$ ($\neg\mathsf{Inf}$), $\kappa_{I, A}$ is an isomorphism and \texttt{TailCokernel I A} is subsingleton (\texttt{subsingleton\_\allowbreak tailCokernel\_\allowbreak of\_\allowbreak fintype}). Adjoining $I = \mathbb{N}$ via an $\mathrm{NNO}$ breaks the isomorphism and precipitates the non-trivial asymptotic \textbf{tail cokernel} (\texttt{nontrivial\_\allowbreak tailCokernel\_\allowbreak nat} and \texttt{chapter1\_\allowbreak pure\_\allowbreak infinity\_\allowbreak residua\_\allowbreak summary}):
  \[
    \mathrm{Res}_{\mathrm{tail}}(\mathbb{N}, A) \;:=\; \prod_{n \in \mathbb{N}} A \,\Big/\, \bigoplus_{n \in \mathbb{N}} A \;\neq\; 0.
  \]
  This single tail quotient\index{Tail quotient residuum} foreshadows two of the monograph's central constructions across Axes~I, III, and~V: in Chapter~\ref{ch:base-fiber} (\S\ref{sec:exotic-r4-e0}, specializing to $A = \mathbb{F}_2$), it gives the profinite hyperfinite Connes quotient $\mathcal{Q}_{E_0} = 2^{\mathbb{N}}/E_0 = \prod_{\mathbb{N}} \mathbb{F}_2 / \bigoplus_{\mathbb{N}} \mathbb{F}_2$ obstructing Borel classification of radial exotic $\mathbb{R}^4$'s~\cite{GompfPanagiotopoulos26}; and in Chapter~\ref{ch:higher-crm} (Theorem~\ref{thm:step-product-counterexample}), replacing additive direct sums by finite-prefix sequences $e^{(m)} = (1_1, \dots, 1_m, 0_{m+1}, \dots)$ inside $\mathcal{Q} \subset \prod_{k \ge 1} \mathcal{A}_k$ generates the pure Postnikov phantom $X_\bullet \in RL(\mathcal{Q}) \setminus \mathcal{Q}$.

\end{enumerate}

\section{Relativization and the Base-Fiber Conservation Principle}\label{sec:base-fiber}

Unlike classical reverse mathematics, which is anchored to a single immutable metatheory, Categorical Reverse Mathematics is doubly relative across comparison squares of cloven doctrinal fibrations and adjunctions $F \dashv U : \mathcal{D}_{\mathcal{E}} \rightleftarrows \mathcal{D}_{\mathcal{S}}$: along a geometric morphism $\gamma : \mathcal{E} \to \mathcal{S}$, the Beck--Chevalley mate square comparing internal slice splitting against direct-image pushforward $\gamma_* : \mathcal{E} \to \mathcal{S}$ couples deformations of the ambient universe $\mathcal{E}$ itself (Axis~1) with variations of the metamathematical base universe $\mathcal{S}$ (Axis~2).

Along Axis~1, weakening the Lawvere--Tierney axioms of $\mathcal{E}$ tests the robustness of categorical residuums. Dropping $\Omega$ moves us from elementary topoi to predicative pretopoi, where Diaconescu's theorem ($\mathrm{IAC} \implies \mathrm{LEM}$) still holds~\cite{Diaconescu75,Maietti05} (\cite[\S D4.5]{Johnstone02}), demonstrating that the choice-to-Booleanness bridge does not depend on impredicative comprehension. Dropping exponentials $Y^X$ moves us to coherent and regular categories, forcing local Cartesian closure into the residuum of $\mathrm{IAC}$, while dropping the Natural Numbers Object $\mathrm{NNO}$ isolates the finitary residuums of Section~\ref{sec:finset-infinity} in $\mathbf{FinSet}$.

Along Axis~2, varying the base topos $\mathcal{S}$ exposes the distinction between \emph{absolute} (Beck--Chevalley mate-invariant under pullback along $\gamma^*$) and \emph{relative} (direct-image $\gamma_*$-sensitive) residuums. Purely internal equivalences, such as $\mathrm{IAC} \implies \Omega_{\mathcal{E}} \simeq 1+1$, are preserved unconditionally by $\gamma^*$. By contrast, external equivalences---such as the Barr--Diaconescu--Freyd theorem ($\mathcal{E} \models \mathrm{AC} \iff \mathcal{E} \simeq \mathbf{Sh}_{\mathcal{S}}(B)$)~\cite{Barr74,Diaconescu75,Freyd72} or Deligne's spatiality theorem for coherent topoi~\cite{Deligne72}---transpose via the Beck--Chevalley mate $\gamma^* \dashv \gamma_*$ to depend on $\mathrm{AC}$ or $\mathrm{BPI}$ holding in the base topos $\mathcal{S}$.

When $\mathrm{AC}$ or $\mathrm{BPI}$ is weakened in the base $\mathcal{S}$, Beck--Chevalley mate transposition does not allow total strength to disappear; instead, the Universal Doctrinal Square Defect obeys the \emph{Base-Fiber Conservation Principle}:

\begin{principle}[Conservation of Logical Strength Across Geometric Morphisms]\label{prin:conservation}%
\uses{def:comparison-triple,def:residuum}%

For any external mathematical theorem $B$ evaluated in a relative topos $\gamma: \mathcal{E} \to \mathcal{S}$ over a weaker principle $A$, the total logical and structural strength required to validate $B$ is conserved across the base $\mathcal{S}$ and the fiber $\mathcal{E}$. Any choice, compactness, or classical logic axiom removed from the metamathematical base $\mathcal{S}$ reappears as a strictly stronger topological, atomic, or cohomological rigidity condition inside the relative residuum $\mathcal{R}(\mathcal{E}/\mathcal{S})$.

\end{principle}

For example, if $\mathrm{BPI}$ is dropped in the base $\mathcal{S}$, coherent $\mathcal{S}$-topoi lose spatiality, and the residuum required to bridge pointfree localic theorems to spatial theorems must upgrade from \emph{finitary coherence} (every cover has a finite subcover) to \emph{supercompactness} (every cover refines to a single element), bypassing the need for base ultrafilters~\cite{Deligne72} (\cite[Ch.~II, \S\S 2--4]{Johnstone82}, \cite[\S D3.3]{Johnstone02}). Through this conservation principle (developed quantitatively in Chapter~\ref{ch:base-fiber}), Categorical Reverse Mathematics unifies the internal logic of topoi, the geometry of sites, and the metamathematics of set theory into a single structural framework.

\chapter{The Choice-Geometry Residuum Atlas and Cohomological Obstruction Theory}
\chaptermark{Choice-Geometry Atlas and Cohomological Obstructions}
\label{ch:choice-atlas}
\chapterepigraph{``The Axiom of Choice is necessary to select a set from an infinite number of pairs of socks, but not an infinite number of pairs of shoes.''}{Bertrand Russell, \emph{Introduction to Mathematical Philosophy} (1919)}

\section{The Anatomy of Choice in Topos Theory}

Interpreting Axis~I of the doctrinal comparison framework (Definitions~\ref{def:comparison-triple}--\ref{def:residuum}) across $\mathcal{D}_0 = \mathbf{Topos} \supseteq \mathbf{Topos}_{\mathrm{AC}} = \mathcal{D}_1$ compares global epimorphism sections against local covers. In classical set theory over $\mathrm{ZF}$, the Axiom of Choice ($\mathrm{AC}$) and its weaker relatives---such as Dependent Choice ($\mathrm{DC}$), Countable Choice ($\mathrm{CC}$), the Presentation Axiom ($\mathrm{COSHEP}$), $\mathrm{WISC}$, and the Boolean Prime Ideal Theorem ($\mathrm{BPI}$)---are formulated as existential assertions about subsets and functions within a single cumulative hierarchy. When we pass to an arbitrary elementary or Grothendieck topos $\mathcal{E}$, these principles bifurcate into distinct internal and external structural properties governing how dependent products, exponentials, and global section functors behave.

To construct the \emph{Choice-Geometry Residuum Atlas}, we establish the precise categorical definitions of the choice hierarchy inside an ambient universe $\mathcal{E}$:

\begin{itemize}[leftmargin=1.5em]

  \item \textbf{External Axiom of Choice ($\mathrm{AC}$):} Every epimorphism $e: X \twoheadrightarrow Y$ in $\mathcal{E}$ admits a global section $s: Y \to X$ satisfying $e \circ s = \mathrm{id}_Y$, meaning every object $Y \in \mathcal{E}$ is externally projective.

  \item \textbf{Internal Axiom of Choice ($\mathrm{IAC}$):} For every morphism $f: X \to Y$, the dependent product functor $\Pi_f: \mathcal{E}/X \to \mathcal{E}/Y$ (right adjoint to pullback $f^*$) preserves epimorphisms~\cite{Johnstone02}---equivalently, for every epimorphism $e: X \twoheadrightarrow Y$ and object $Z \in \mathcal{E}$, the exponential map $e^Z: X^Z \to Y^Z$ is an epimorphism (every object is internally projective).

  \item \textbf{Enough Projectives ($\mathrm{COSHEP}$):} Every object $X \in \mathcal{E}$ can be covered by an epimorphism $P \twoheadrightarrow X$ from an externally projective object $P$, corresponding to the categorical Presentation Axiom.

  \item \textbf{Weakly Initial Set of Covers ($\mathrm{WISC}$):} For every object $X \in \mathcal{E}$ there exists a small (set-indexed) family of epimorphisms $\{q_i : Y_i \twoheadrightarrow X\}_{i \in I}$ through at least one of which every epimorphism $E \twoheadrightarrow X$ factors (van den Berg--Moerdijk--Streicher, weakening $\mathrm{COSHEP}$ by removing projectivity of $Y_i$).

  \item \textbf{Small Violations of Choice ($\mathrm{SVC}$ and $\mathrm{SVC}_{\mathrm{ret}}$):} Andreas Blass's~\cite{Blass79} $\mathrm{SVC}$ posits a single parameterizing object $S \in \mathcal{E}$ such that every object $X \in \mathcal{E}$ is a subquotient of $C \times S$ for a projective choice object $C$; its retract refinement $\mathrm{SVC}_{\mathrm{ret}}$ requires a parameter $S$ with projective cover $q : P \twoheadrightarrow S$ such that every $X \in \mathcal{E}$ is a retract of a subobject $U \rightarrowtail C \times P$ with $C \times P$ projective (equivalently, the presentation quotient over $C \times P$ splits).

  \item \textbf{Sequential Choice ($\mathrm{DC}$, $\mathrm{IDC}$, and $\mathrm{CC}$):} In a topos with a Natural Numbers Object ($\mathrm{NNO}$), external $\mathrm{DC}$ asserts that every countable inverse tower of epimorphisms $\cdots \twoheadrightarrow X_2 \twoheadrightarrow X_1 \twoheadrightarrow X_0$ admits a global limit projection splitting, while Internal Dependent Choice ($\mathrm{IDC}$) asserts that the inverse limit projection $\varprojlim_{n \in \omega} X_n \twoheadrightarrow X_0$ remains an epimorphism in every slice $\mathcal{E}/E$.

  \item \textbf{Boolean Prime Ideal Theorem ($\mathrm{BPI}$ and $\mathrm{IBPI}$):} Externally, every coherent topos has enough points (Deligne's theorem~\cite{Deligne72}); internally ($\mathrm{IBPI}$), every non-trivial internal Boolean algebra $B \in \mathbf{BA}(\mathcal{E})$ admits an ultrafilter (equivalently, internal Stone/coherent locales are compact~\cite{Johnstone82,Johnstone02}).

  \item \textbf{Universal Local Choice:} In any Grothendieck sheaf topos $\mathcal{E} \simeq \mathbf{Sh}(\mathbb{C}, J)$, every epimorphism $e: X \twoheadrightarrow Y$ splits locally along a covering sieve of the site $(\mathbb{C}, J)$.

\end{itemize}

A foundational structural asymmetry separates external $\mathrm{AC}$ ($\mathrm{EAC}$) from internal choice ($\mathrm{IAC}$). While external $\mathrm{AC}$ implies $\mathrm{IAC}$ in every elementary topos, $\mathrm{IAC}$ is a purely internal local-Cartesian property: $\mathcal{E} \models \mathrm{IAC}$ if and only if for every morphism $f: X \to Y$, the \emph{dependent product functor} $\Pi_f: \mathcal{E}/X \to \mathcal{E}/Y$ (the right adjoint to pullback $f^*$) preserves epimorphisms~\cite{Johnstone02} (Elephant \S D4.5), or equivalently every exponential $(-)^Z : \mathcal{E} \to \mathcal{E}$ preserves epimorphisms. (Notice that the pullback functor $f^*$ itself preserves epimorphisms unconditionally in \emph{every} regular category and elementary topos because colimits are universal.) Because $\Pi_f$ commutes with slice restriction, $\mathrm{IAC}$ is inherited by all slice topoi $\mathcal{E}/E$. External $\mathrm{AC}$, by contrast, tests whether internal splittings survive the global sections functor $\Gamma = \mathrm{Hom}_{\mathcal{E}}(1_{\mathcal{E}}, -): \mathcal{E} \to \mathbf{Set}$---which depends on external projectivity of $1_{\mathcal{E}}$, a property strictly stronger than two-valuedness ($\mathrm{Sub}_{\mathcal{E}}(1_{\mathcal{E}}) \simeq \{0, 1\}$), as atomic permutation topoi demonstrate.

\section{Diaconescu's Theorem, Freyd's Countermodel, and the Fundamental Residuums}

Before calibrating intermediate choice principles, Categorical Reverse Mathematics must determine the absolute baseline residuums: $(\top \Rightarrow_{\mathbf{Topos}} \mathrm{AC})$ and $(\mathrm{IAC} \Rightarrow_{\mathbf{Topos}} \mathrm{AC})$. The first bridge between choice and internal logic is given by Diaconescu's theorem\index{Diaconescu's theorem}~\cite{Diaconescu75}.

\begin{theorem}[Diaconescu's Theorem and Predicative Stability]\label{thm:diaconescu}%
\uses{def:comparison-triple,def:residuum}%

In any elementary topos $\mathcal{E}$~\cite{Diaconescu75}---and more generally in any locally Cartesian closed pretopos~\cite[\S D4.5, Thm.~D4.5.1--Cor.~D4.5.2]{Johnstone02} (cf.~\cite{Maietti05})---the Internal Axiom of Choice implies that the internal logic is Boolean ($\mathrm{IAC} \implies \Omega \simeq 1+1$, or internal $\mathrm{LEM}$).

\end{theorem}

\begin{proof}[Proof Sketch]
Given any subobject $U \rightarrowtail X$, one constructs the pushout $X +_U X$ of two copies of $X$ glued along $U$, equipped with the canonical codiagonal epimorphism $e: X + X \twoheadrightarrow X +_U X$. Pulling back or exponentiating $e$ via $\mathrm{IAC}$ produces an internal section whose preimages decompose $X$ into $U \amalg \neg U$, proving that every subobject lattice $\mathrm{Sub}_{\mathcal{E}}(X)$ is complemented. Because this construction requires only finite colimits, pullbacks, and local Cartesian closure rather than impredicative power objects $\Omega^X$, the implication $\mathrm{IAC} \implies \mathrm{LEM}$ is predicative and base-invariant across any base topos $\mathcal{S}$.
\end{proof}

Diaconescu's theorem proves that Booleanness ($\Omega \simeq 1+1$) is an obligatory factor of the residuum $(A \Rightarrow_{\mathbf{Topos}} \mathrm{AC})$ whenever the weak principle $A$ is consistent with non-Boolean intuitionistic logic. However, Booleanness and $\mathrm{IAC}$ do not suffice to split epimorphisms externally.

\begin{theorem}[The Internal-to-External Choice Residuum: $\mathrm{IAC} \Rightarrow \mathrm{AC}$]\label{thm:iac-to-ac}%
\uses{thm:diaconescu}%

In an elementary topos $\mathcal{E}$, the categorical residuum bridging the Internal Axiom of Choice to the External Axiom of Choice is the external projectivity of the terminal object $1_{\mathcal{E}}$\index{External projectivity of $1_{\mathcal{E}}$}:

\[
  \mathrm{IAC} + \bigl(1_{\mathcal{E}} \text{ is externally projective}\bigr) \iff \mathrm{AC}
\]

In particular, whenever $1_{\mathcal{E}}$ is externally projective (and a fortiori if $\mathcal{E}$ is a well-pointed topos), $\mathrm{IAC} \iff \mathrm{AC}$ unconditionally.

\end{theorem}

\begin{proof}
If $\mathcal{E} \models \mathrm{AC}$, every object (including $1_{\mathcal{E}}$) is externally projective, and $\mathrm{AC} \implies \mathrm{IAC}$. Conversely, suppose $\mathcal{E} \models \mathrm{IAC}$ and $1_{\mathcal{E}}$ is externally projective. For any epimorphism $e: X \twoheadrightarrow Y$, internal choice ($\mathrm{IAC}$) guarantees that the exponential morphism $e^Y: X^Y \twoheadrightarrow Y^Y$ is an epimorphism. Pulling $e^Y$ back along the name of the identity morphism $\ulcorner \mathrm{id}_Y \urcorner: 1_{\mathcal{E}} \to Y^Y$ yields the internal object of sections $\mathrm{Sect}(e) \subseteq X^Y$ together with a canonical pullback epimorphism $\mathrm{Sect}(e) \twoheadrightarrow 1_{\mathcal{E}}$. Because $1_{\mathcal{E}}$ is externally projective, the epimorphism $\mathrm{Sect}(e) \twoheadrightarrow 1_{\mathcal{E}}$ splits via a global point $1_{\mathcal{E}} \to \mathrm{Sect}(e)$, which Cartesian adjunction transposes into a global section $s: Y \to X$ satisfying $e \circ s = \mathrm{id}_Y$.
\end{proof}

\noindent\textbf{Remark on Subterminal Generation (Localicity).} In any elementary topos satisfying $\mathrm{AC}$, subterminal generation (localicity\index{Localicity (subterminal generation)}: subobjects of $1_{\mathcal{E}}$ generate $\mathcal{E}$) holds automatically as a \emph{consequence} of $\mathrm{AC}$ (since supports of slice objects split). Rather than being needed to bridge $\mathrm{IAC}$ to $\mathrm{AC}$, localicity serves as the exact geometric residuum bridging internal Booleanness ($\mathrm{LEM}$) to $\mathrm{AC}$ in Grothendieck topoi over $\mathbf{Set}_{\mathrm{AC}}$ (Theorem~\ref{thm:bdf}).

When we specialize to Grothendieck topoi over $\mathbf{Set}$ (assuming $\mathrm{AC}$ in the base $\mathbf{Set}$), one might ask whether Booleanness together with two-valuedness\index{Two-valuedness} ($\mathrm{Sub}(1_{\mathcal{E}}) \simeq \{0, 1\}$) or external projectivity of $1_{\mathcal{E}}$ already forces localicity. Crucially, as detailed in Chapter~\ref{ch:realizability-sites} (\S\ref{subsec:kihara-choice-nominal}), two-valuedness must never be conflated with external projectivity of the terminal object $1_{\mathcal{E}}$: in any connected group-action or atomic permutation topos $\mathbf{Cont}(\mathcal{G})$~\cite[Ch.~III, \S 9]{MacLaneMoerdijk92} (such as $\mathbf{Cont}(\mathrm{Aut}(\mathcal{M}))$ for an infinite Fra\"iss\'e structure $\mathcal{M}$ or Schanuel's nominal topos $\mathbf{Nom}$), $1_{\mathcal{E}}$ is strictly two-valued, yet $1_{\mathcal{E}}$ is \emph{never} externally projective whenever a proper open subgroup $U \subsetneq \mathcal{G}$ exists, since the coset epimorphism $\mathcal{G}/U \twoheadrightarrow 1_{\mathcal{E}}$ has zero fixed points ($\Gamma(\mathcal{G}/U) = (\mathcal{G}/U)^{\mathcal{G}} = \emptyset$). Conversely, in sheaves on a non-trivial complete Boolean algebra $\mathbf{Sh}(B)$ over $\mathbf{Set}_{\mathrm{AC}}$, $1$ is externally projective while failing two-valuedness ($\mathrm{Sub}(1) \simeq B$)~\cite{Barr74} (\cite[\S D4.5]{Johnstone02}; \cite[Chs.~VI, IX]{MacLaneMoerdijk92}). Peter Freyd's classic counterexample\index{Freyd's countermodel}~\cite{Freyd72} demonstrates that even in a two-valued Boolean Grothendieck topos with an $\mathrm{NNO}$, localicity need not hold.

\begin{theorem}[Freyd's Countermodel and the Barr--Diaconescu--Freyd Residuum]\label{thm:bdf}%
\uses{thm:diaconescu,thm:iac-to-ac}%

There exists a two-valued, Boolean Grothendieck topos $\mathcal{F} = \mathbf{Sh}_{\neg\neg}(\mathbb{A})$ equipped with a Natural Numbers Object ($\mathrm{NNO}$) in which $\mathrm{AC}$ fails~\cite{Freyd72}. Over $\mathbf{Set}_{\mathrm{AC}}$, a Grothendieck topos $\mathcal{E}$ satisfies $\mathrm{AC}$ if and only if $\mathcal{E} \simeq \mathbf{Sh}(B)$ is equivalent to sheaves on a complete Boolean algebra $B$~\cite{Barr74,Diaconescu75,Freyd72} (\cite[\S D4.5, Thms.~D4.5.3--D4.5.7]{Johnstone02}). Consequently, the baseline Grothendieck residuum is:

\[
  \top + \bigl(\text{Booleanness } (\Omega \simeq 1+1) \wedge \text{Localicity}\bigr) \iff \mathrm{AC}
\]

\end{theorem}

In Freyd's topos $\mathcal{F} = \mathbf{Sh}_{\neg\neg}(\mathbb{A})$, the nontrivial automorphism group of the site prevents subterminal objects from generating the category, creating a countable sequence of epimorphisms that cannot be split by global sections~\cite{Freyd72} (\cite[\S D4.5]{Johnstone02}). Thus, in the lattice of Grothendieck topoi, \emph{localicity} is completely orthogonal to Booleanness and two-valuedness. Moreover, by the \emph{Base-Fiber Conservation Principle} (Chapter~\ref{ch:foundations}, Principle~\ref{prin:conservation}), if $\mathrm{AC}$ is dropped in the base topos $\mathcal{S}$, localicity and Booleanness over $\mathcal{S}$ no longer suffice for $\mathcal{E} \models \mathrm{AC}$; the residuum $\mathcal{R}(\mathcal{E}/\mathcal{S})$ must strengthen to require that the underlying locale $B$ is an internally projective atomic Boolean algebra over a choice slice of $\mathcal{S}$.

\section{\texorpdfstring{Residuums of Intermediate Choice Principles ($\mathrm{COSHEP}$, $\mathrm{WISC}$, $\mathrm{IBPI}$, and $\mathrm{IDC}$)}{Residuums of Intermediate Choice Principles (COSHEP, WISC, IBPI, and IDC)}}

We now compute and conjecture the categorical residuums $(A \Rightarrow_{\mathbf{Topos}} \mathrm{AC})$ when $A$ is a non-trivial weak choice axiom. Each weak choice axiom absorbs a specific dimension of projective or topological complexity, leaving a complementary geometric invariant as the residuum $\mathcal{R}(\mathcal{E})$.

\subsection{\texorpdfstring{Projective Covers, Small Violations, and $\mathrm{WISC}$ ($\mathrm{COSHEP} \Rightarrow \mathrm{AC}$)}{Projective Covers, Small Violations, and WISC (COSHEP => AC)}}

In constructive and categorical logic, the Presentation Axiom $\mathrm{COSHEP}$\index{Axiom of Choice!COSHEP@$\mathrm{COSHEP}$ (Presentation Axiom)} (``the category of sets has enough projectives'') is strictly weaker than $\mathrm{AC}$ because projective covers $P \twoheadrightarrow X$ do not force $X$ itself to be projective. Between bare topos exactness and $\mathrm{COSHEP}$ sits an even weaker structural principle, \textbf{$\mathrm{WISC}$ (Weakly Initial Set of Covers / Axiom of Multiple Choice)}\index{Axiom of Choice!WISC@$\mathrm{WISC}$ (Weakly Initial Set of Covers)}, asserting that every object $X \in \mathcal{E}$ possesses a small family of epimorphisms $\{q_i : Y_i \twoheadrightarrow X\}_{i \in I}$ through at least one of which every epimorphism $E \twoheadrightarrow X$ factors. Whereas $\mathrm{WISC}$ suffices for predicative coinductive and $W$-type semantics without needing projective covers, $\mathrm{COSHEP}$ strengthens $\mathrm{WISC}$ by asking for a single projective cover $P \twoheadrightarrow X$ (as holds in Martin Hyland's Effective Topos $\Eff$~\cite{Hyland82}, \cite[\S\S 3.1--3.2]{vanOosten08}, Chapter~\ref{ch:realizability-sites}, without $\mathrm{LEM}$ or $\mathrm{AC}$).

Dual to $\mathrm{COSHEP}$ is Andreas Blass's principle of Small Violations of Choice ($\mathrm{SVC}$)\index{Axiom of Choice!SVC@$\mathrm{SVC}$ (Small Violations of Choice)}~\cite{Blass79}, which bounds non-projectivity across the universe by a single parameter object $S$. Two sharp obstructions govern any attempt to bridge $\mathrm{COSHEP}$ to external $\mathrm{AC}$: first, the discrete group topos $\mathbf{Set}^{\mathbb{Z}/2\mathbb{Z}}$ is Boolean ($\Omega \simeq 1+1$), satisfies $\mathrm{COSHEP}$ (free regular covers $|X| \times \mathbb{Z}/2\mathbb{Z} \twoheadrightarrow X$), and satisfies $\mathrm{SVC}$ ($S = \mathbb{Z}/2\mathbb{Z}$), yet fails $\mathrm{AC}$ because $1$ is not externally projective ($\Gamma(\mathbb{Z}/2\mathbb{Z}) = \emptyset$); second, whereas complemented subobjects and retracts of projectives are always projective in a Boolean topos, arbitrary \emph{quotients} of projectives are not projective in general (the \emph{quotient fallacy}). Incorporating both external projectivity of $1_{\mathcal{E}}$ and retract presentations yields the exact rectified residuum theorem:

\begin{theorem}[The Rectified $\mathrm{COSHEP}$--$\mathrm{SVC}$ Residuum]\label{thm:coshep-svc}%
\uses{thm:diaconescu,thm:iac-to-ac}%

In an elementary topos $\mathcal{E}$, external $\mathrm{AC}$ holds if and only if $\mathcal{E}$ satisfies $\mathrm{COSHEP}$, Booleanness ($\Omega \simeq 1+1$), external projectivity of the terminal object $1_{\mathcal{E}}$, and Blass's retract-bounded Small Violations of Choice ($\mathrm{SVC}_{\mathrm{ret}}$):

\[
  \mathrm{COSHEP} + \bigl(\mathrm{SVC}_{\mathrm{ret}} \wedge \text{Booleanness} \wedge (1_{\mathcal{E}} \text{ is externally projective})\bigr) \iff \mathrm{AC}
\]

where $\mathrm{SVC}_{\mathrm{ret}}$ guarantees a parameter $S$ with projective cover $q : P \twoheadrightarrow S$ such that every object $X \in \mathcal{E}$ is a retract of a subobject $U \rightarrowtail C \times P$ for a choice object $C$ with $C \times P$ externally projective.

\end{theorem}

\begin{proof}
Under $\mathrm{AC}$, every object (including $1_{\mathcal{E}}$) is externally projective (so $\mathrm{COSHEP}$ holds), Theorem~\ref{thm:diaconescu} yields Booleanness, and $S = P = 1_{\mathcal{E}}$ with $C = X$ witnesses $\mathrm{SVC}_{\mathrm{ret}}$ ($X \cong X \times 1_{\mathcal{E}}$).

Conversely, assume $\mathrm{COSHEP}$, $\mathrm{SVC}_{\mathrm{ret}}$, Booleanness ($\Omega \simeq 1+1$), and external projectivity of $1_{\mathcal{E}}$. We first prove two exact closure lemmas for externally projective objects:
\begin{enumerate}[label=(\alph*),leftmargin=1.8em]
  \item \emph{Boolean Subobject Lemma:} Every subobject $m : U \rightarrowtail Q$ of an externally projective object $Q$ is externally projective. Indeed, Booleanness furnishes a complement $\neg U \rightarrowtail Q$ such that $Q \cong U \amalg \neg U$. Given any epimorphism $e : E \twoheadrightarrow U$, the coproduct morphism $e \amalg \mathrm{id}_{\neg U} : E \amalg \neg U \twoheadrightarrow U \amalg \neg U \cong Q$ is an epimorphism, hence admits a section $\sigma : Q \to E \amalg \neg U$ because $Q$ is projective. Consider the composite $\sigma \circ \mathrm{in}_1 : U \to E \amalg \neg U$, and decompose $U = U_E \amalg U_{\neg U}$ by pulling back along the coproduct summands $\mathrm{in}_1 : E \rightarrowtail E \amalg \neg U$ and $\mathrm{in}_2 : \neg U \rightarrowtail E \amalg \neg U$. On $U_{\neg U}$, $(e \amalg \mathrm{id}_{\neg U}) \circ \sigma \circ \mathrm{in}_1|_{U_{\neg U}}$ simultaneously factors through $\mathrm{in}_2 : \neg U \rightarrowtail U \amalg \neg U$ and equals $\mathrm{in}_1|_{U_{\neg U}} : U_{\neg U} \to U \rightarrowtail U \amalg \neg U$. By disjointness of coproducts in any elementary topos ($U \times_{U \amalg \neg U} \neg U \cong \emptyset$), $U_{\neg U} \cong \emptyset$, so $\sigma \circ \mathrm{in}_1$ factors uniquely through $E$ via a morphism $s : U \to E$ satisfying $e \circ s = \mathrm{id}_U$.
  \item \emph{Retract Preservation Lemma:} Every retract $r : U \twoheadrightarrow X$ ($r \circ i = \mathrm{id}_X$) of an externally projective object $U$ is externally projective. For any epimorphism $e : E \twoheadrightarrow X$, pulling back $e$ along $r$ yields an epimorphism $e' : E \times_X U \twoheadrightarrow U$ (universality of epimorphisms in a topos), which splits via $s' : U \to E \times_X U$; setting $s := \mathrm{pr}_1 \circ s' \circ i : X \to E$ satisfies $e \circ s = r \circ e' \circ s' \circ i = r \circ i = \mathrm{id}_X$.
\end{enumerate}
Now let $X \in \mathcal{E}$ be any object. By $\mathrm{COSHEP}$ and $\mathrm{SVC}_{\mathrm{ret}}$, $X$ is a retract of a subobject $U \rightarrowtail C \times P$ where $Q := C \times P$ is externally projective. Lemma~(a) shows that $U$ is externally projective, and Lemma~(b) transfers external projectivity from $U$ to its retract $X$. Since $X$ was arbitrary, every object of $\mathcal{E}$ is externally projective, establishing $\mathrm{AC}$. Note that without external projectivity of $1_{\mathcal{E}}$, the theorem fails in $\mathbf{Set}^{\mathbb{Z}/2\mathbb{Z}}$; and without the retract condition on presentations over $C \times P$, mere subquotients of projectives fail to be projective unless arbitrary quotients of projectives already split.
\end{proof}

\subsection{\texorpdfstring{From Logical Compactness to Zero-Dimensionality ($\mathrm{IBPI} \Rightarrow \mathrm{IAC}$ and $\mathrm{BPI} \Rightarrow \mathrm{AC}$)}{From Logical Compactness to Zero-Dimensionality (IBPI => IAC and BPI => AC)}}

In classical set theory, Halpern and L\"auchli (with atoms) and Halpern and L\'evy~\cite{HalpernLevy71} (in $\mathrm{ZF}$; see~\cite[Ch.~II, \S\S 3--4]{Johnstone82}, \cite[\S D3.3]{Johnstone02}) proved that the Boolean Prime Ideal Theorem\index{Boolean Prime Ideal Theorem ($\mathrm{BPI}$)} ($\mathrm{BPI}$, or the Ultrafilter Principle $\mathrm{UF}$) is strictly weaker than $\mathrm{AC}$. Categorically, Pierre Deligne's theorem~\cite[Exp.~VI, Appendice]{Deligne72} characterizes $\mathrm{BPI}$ as the exact external condition ensuring that every coherent topos has enough points, whereas the internal Boolean Prime Ideal Theorem ($\mathrm{IBPI}$) asserts internally that every non-trivial internal Boolean algebra $B \in \mathbf{BA}(\mathcal{E})$ admits an ultrafilter (internal compactness of Stone locales~\cite[Ch.~II, \S\S 3--4]{Johnstone82}, \cite[\S D3.3, Thm.~D3.3.13]{Johnstone02}).

Scoping $\mathrm{IBPI} \Rightarrow \mathrm{IAC}$ requires two structural rectifications. First, over a general base $\mathcal{S} \models \mathrm{BPI}$, any coherent Boolean localic topos $\mathbf{Sh}_{\mathcal{S}}(L)$ collapses to a finite power $\mathcal{S}^m$. Indeed, any atomless summand $c > 0$ carries an internal ultrafilter $u \subseteq \mathcal{O}(L){\downarrow}c$ whose complements $\{c \wedge \neg x : x \in u\}$ form a directed join to $c$ without a finite subcover, violating coherent compactness of $1_L$; hence $L$ is finite atomic and $\mathbf{Sh}_{\mathcal{S}}(L) \simeq \mathcal{S}^m$ satisfies $\mathrm{IAC}$ if and only if $\mathcal{S} \models \mathrm{IAC}$, failing over Halpern--L\'evy bases $\mathcal{S} \models \mathrm{BPI}+\neg\mathrm{AC}$~\cite{HalpernLevy71} (constructed in~\S\ref{sec:symmetric-bv-zf}). Second, over $\mathbf{Set}_{\mathrm{AC}}$ in a Boolean topos, finite \emph{subobject} covers $\{U_i \rightarrowtail Y\}_{i=1}^n$ trivially disjointify via $V_i := U_i \wedge \neg\bigvee_{j<i} U_j$~\cite{Johnstone02}. The genuine pasting obstruction lives instead on \emph{non-monic} coherent epimorphic covers $\{u_i : U_i \to Y\}$ where groupoids twist local sections, as in Halpern's Boolean coherent atomic permutation topos $\mathbf{Cont}(\mathrm{Aut}(\mathcal{T}))$~\cite{Caramello18,Kihara26}, which satisfies $\mathrm{IBPI} + (\Omega \simeq 1+1) + \neg\mathrm{IAC}$. Replacing subobject refinement with \emph{\'etale clopen-partition refinement}---for which $\mathrm{IAC}$-necessity holds unconditionally---we state Conjecture~\ref{conj:bpi-ac} over $\mathbf{Set}_{\mathrm{AC}}$ (or $\mathcal{S} \models \mathrm{IAC}$):

\begin{conjecture}[The Zero-Dimensionality Residuum Conjecture: $\mathrm{IBPI} \Rightarrow \mathrm{IAC}$]\label{conj:bpi-ac}

In a coherent Grothendieck topos $\mathcal{E}$ over $\mathbf{Set}_{\mathrm{AC}}$ (or over a Boolean base $\mathcal{S} \models \mathrm{IAC}$), the categorical residuum bridging the Internal Boolean Prime Ideal Theorem ($\mathrm{IBPI}$) to Internal Choice ($\mathrm{IAC}$) is the conjunction of Booleanness ($\Omega_{\mathcal{E}} \simeq 1+1$) and \emph{\'etale internal zero-dimensionality}:

\[
  \begin{aligned}
    \mathrm{IBPI} + \bigl(&\text{Booleanness } (\Omega_{\mathcal{E}} \simeq 1+1) \\
    &\wedge \text{\'Etale Internal Zero-Dimensionality}\bigr) \iff \mathrm{IAC}
  \end{aligned}
\]

where \emph{\'etale internal zero-dimensionality} asks for an inhabited cover $E \twoheadrightarrow 1_{\mathcal{E}}$ such that in $\mathcal{E}/E$, every finite coherent epimorphic family $\{u_i : U_i \to Y\}_{i=1}^n$ ($\coprod_{i=1}^n U_i \twoheadrightarrow Y$) refines to a finite complemented subobject partition $Y \cong \coprod_{j=1}^m V_j$ admitting sections $V_j \to U_{i(j)}$. Necessity ($\mathrm{IAC} \implies$ \'etale internal zero-dimensionality) holds unconditionally at $E = 1_{\mathcal{E}}$ by pulling back coproduct summands along a global section $Y \to \coprod_{i=1}^n U_i$; sufficiency remains open (equivalently vanishing of $H^1(\mathcal{E}/Y, K_X)$ at $A = \mathrm{IBPI}$ in Conjecture~\ref{conj:cohom-vanishing}, yielding external $\mathrm{AC}$ when $1_{\mathcal{E}}$ is externally projective).

\end{conjecture}

\subsection{\texorpdfstring{From Internal Dependent Choice to Internal Choice ($\mathrm{IDC} \Rightarrow \mathrm{IAC}$)}{From Internal Dependent Choice to Internal Choice (IDC => IAC)}}

In a topos $\mathcal{E}$ equipped with an $\mathrm{NNO}$, Countable Choice ($\mathrm{CC}$) and Dependent Choice ($\mathrm{DC}$, or its internal local-Cartesian counterpart $\mathrm{IDC}$) govern sequential induction: $\mathrm{IDC}$ asserts that every internal inverse $\omega$-tower of epimorphisms $\cdots \twoheadrightarrow X_2 \twoheadrightarrow X_1 \twoheadrightarrow X_0$ admits an internal limit projection $\varprojlim_{n \in \omega} X_n \twoheadrightarrow X_0$ that remains epic in every slice $\mathcal{E}/E$.

Comparing $\mathrm{IDC}$ to $\mathrm{IAC}$ requires three structural distinctions. First, over $\mathbf{Set}_{\mathrm{AC}}$, \emph{every} Boolean localic topos $\mathbf{Sh}(B)$ already satisfies $\mathrm{AC}$ by Theorem~\ref{thm:bdf}~\cite{Barr74} (\cite[\S D4.5]{Johnstone02}) regardless of the countable chain condition ($\mathrm{CCC}$), which enters externally only over $\mathcal{S} \models \mathrm{DC}$ (Theorem~\ref{thm:bdf-deformed}). Second, asking every $e : X \twoheadrightarrow Y$ to be a \emph{retract} in $\mathcal{E}/Y$ of a decidable-finite-fiber $\omega$-tower $\varprojlim_{n \in \omega} X_n \twoheadrightarrow Y$ fails in both directions. In $\mathbf{Set}_{\mathrm{ZFC}} \models \mathrm{IAC}$, finite-fiber $\omega$-towers over $1$ satisfy $|\varprojlim_{n \in \omega} X_n| \le 2^{\aleph_0}$, blocking retracts whenever $|X| > 2^{\aleph_0}$; conversely, under $\mathrm{IDC}$ even $2$-to-$1$ covers fail to split over uncountable codomains, both on $2^{\mathbb{N}}/E_0^{\mathrm{even}} \twoheadrightarrow \mathcal{Q}_{E_0}$ in $\mathrm{ZF}+\mathrm{DC}+\mathrm{BP}$~\cite{Serafin23} and on ordered-pair covers $V^{(2)} \twoheadrightarrow \binom{V}{2}$ in Kihara's Rung~1--2 atomic topoi $\mathbf{Sh}(\mathcal{K}_\lambda^{\mathrm{op}}, J_{\mathrm{at}})$ at $\lambda^{<\lambda}=\lambda>\aleph_0$~\cite{Kihara26}. Third, at the epimorphism level in any Boolean topos $\mathcal{E}$ with an $\mathrm{NNO}$, least-number minimization~\cite[\S D5.1]{Johnstone02} and prefix-section chains prove unconditionally that $\mathrm{IDC}$ plus \emph{\'etale forward refinement} ($E^* e \circ r = \pi_0$ on $E \twoheadrightarrow 1_{\mathcal{E}}$) along pullbacks of subcountable-codomain epimorphisms $q_n : A_n \twoheadrightarrow B_n$ ($B_n \rightarrowtail \mathbb{N}_{\mathcal{E}/E}$) is equivalent to $\mathrm{IAC}$. This refinement condition is strictly non-tautological without $\mathrm{IDC}$: in profinite $\mathbf{Cont}(\mathbb{Z}_2)$, every epimorphism refines along $\coprod_{m \in \mathbb{N}}(\mathbb{Z}_2/2^m\mathbb{Z}_2) \twoheadrightarrow \Delta\mathbb{N}$ via $(h + ((g-h)\bmod 2^{a_i})) \cdot x_i$, whereas $\mathrm{IDC}$ fails. Seeking an \emph{intrinsic} site invariant over $\mathbf{Set}_{\mathrm{AC}}$, we conjecture:

\begin{conjecture}[The Internal $\omega$-Presentability Residuum Conjecture: $\mathrm{IDC} \Rightarrow \mathrm{IAC}$]\label{conj:dc-ac}

In a Boolean Grothendieck topos $\mathcal{E}$ with an $\mathrm{NNO}$ over $\mathbf{Set}_{\mathrm{AC}}$ (or $\mathcal{S} \models \mathrm{IAC}$), the intrinsic residuum bridging Internal Dependent Choice ($\mathrm{IDC}$) to Internal Choice ($\mathrm{IAC}$) is \emph{\'etale $\omega$-presentability of localic groupoid isotropy}:

\[
  \mathrm{IDC} + \bigl(\text{\'Etale } \omega\text{-presentability of localic groupoid isotropy}\bigr) \iff \mathrm{IAC}
\]

meaning $\mathcal{E}$ admits a Joyal--Tierney open localic groupoid presentation $\mathbb{G} = (G_1 \rightrightarrows G_0)$~\cite[\S C5.3]{Johnstone02} whose identity section $e(G_0) \subseteq G_1$ is an $\omega$-cofinal intersection of open subgroupoid neighborhoods over $G_0$. Necessity holds unconditionally on Boolean \'etendues ($e(G_0)$ clopen), and sufficiency holds on normal-basis pro-discrete $\mathbf{Cont}(\mathcal{G})$ ($\bigcap_{n \in \omega} g_n U_n g_n^{-1} = \bigcap_{n \in \omega} U_n = \{1_{\mathcal{G}}\}$ open, forcing $\mathcal{G}$ discrete); general sufficiency remains open (giving external $\mathrm{AC}$ when $1_{\mathcal{E}}$ is externally projective, killing uncountable-base $K_X, \underline{\mathrm{Aut}}([c])$ in Conjecture~\ref{conj:cohom-vanishing}).

\end{conjecture}

\section{Cohomological Obstruction Theory for Categorical Reversals}

A central methodological insight of Categorical Reverse Mathematics is that the failure of a stronger theorem $B$ in a universe $\mathcal{E}$ satisfying a weaker principle $A$ ($\mathcal{E} \models A \not\models B$) is not merely a logical counterexample---it is a \emph{topological obstruction} that can be measured by sheaf cohomology.

Consider any Grothendieck topos $\mathcal{E} \simeq \mathbf{Sh}(\mathbb{C}, J)$. By Universal Local Choice, every epimorphism $e: X \twoheadrightarrow Y$ in $\mathcal{E}$ splits locally: there exists a covering family $\{u_i: U_i \to Y\}_{i \in I} \in J(Y)$ and local sections $s_i: U_i \to X$ such that $e \circ s_i = u_i$ for each $i \in I$. The External Axiom of Choice ($\mathrm{AC}$) asks whether these local sections $s_i$ can be modified to agree on all pairwise fiber products $U_{ij} = U_i \times_Y U_j$, thereby descending to a single global section $s: Y \to X$.

Suppose $e: X \twoheadrightarrow Y$ is a principal homogeneous space ($G$-torsor) over $Y$ under an internal structure group $G \in \mathbf{Grp}(\mathcal{E}/Y)$ (or, in the abelian case, an extension $0 \to A \to X \xrightarrow{e} Y \to 0$ classified by $\mathrm{Ext}^1_{\mathcal{E}}(Y, A)$). On each overlap $U_{ij} = U_i \times_Y U_j$, the two local sections $\mathrm{pr}_1^*(s_i)$ and $\mathrm{pr}_2^*(s_j)$ differ by a unique group element $g_{ij} \in G(U_{ij})$ satisfying:

\[
  \mathrm{pr}_2^*(s_j) = \mathrm{pr}_1^*(s_i) \cdot g_{ij} \quad \text{on } U_i \times_Y U_j
\]

On triple overlaps $U_{ijk} = U_i \times_Y U_j \times_Y U_k$, associativity forces the cocycle condition $g_{ik} = g_{ij} \cdot g_{jk}$, defining a \v{C}ech 1-cocycle $[g_{ij}] \in \check{H}^1(\{U_i \to Y\}, G)$ and hence a sheaf cohomology class $\omega(e) \in H^1(\mathcal{E}/Y, G)$. Following Jean Giraud~\cite[Chs.~III--IV]{Giraud71}, any short exact sequence of internal groups $1 \to K \to G \to Q \to 1$ in $\mathcal{E}/Y$ induces a $6$-term non-abelian cohomology exact sequence
\[
  1 \longrightarrow H^0(K) \longrightarrow H^0(G) \longrightarrow H^0(Q) \xrightarrow{\;\partial\;} H^1(K) \longrightarrow H^1(G) \longrightarrow H^1(Q),
\]
where the torsor $e: X \twoheadrightarrow Y$ admits a global section $s: Y \to X$ if and only if its obstruction class $\omega(e)$ vanishes in $H^1(\mathcal{E}/Y, G)$.

\medskip
\noindent\textbf{Anchor Witness Example ($\mathcal{T}_{\mathrm{odd}}$ in Schanuel's Nominal Topos $\mathbf{Nom}$).} Inside $\mathbf{Nom} = \mathbf{Cont}(\mathrm{Sym}(\mathbb{N}))$, let $\mathrm{Finset}(\mathbb{A})$ denote the $\mathbb{F}_2$-group object (under symmetric difference $\triangle$) of finite subsets of atoms $\mathbb{A} = \mathbb{N}$, sitting in the parity short exact sequence
\[
  1 \longrightarrow \mathrm{EvenFinset} \longrightarrow \mathrm{Finset}(\mathbb{A}) \xrightarrow{\;|\cdot|\bmod 2\;} \mathbb{Z}/2\mathbb{Z} \longrightarrow 1.
\]
Taking global sections $H^0(\mathbf{Nom}; -) = \Gamma$, every singleton $\{a\} \in \mathrm{Finset}(\mathbb{A})$ maps to $1 \in \mathbb{Z}/2\mathbb{Z}$ (so $|\cdot|\bmod 2$ is an epimorphism in $\mathbf{Nom}$), while the connecting boundary map sends $1 \in H^0(\mathbf{Nom}; \mathbb{Z}/2\mathbb{Z}) \cong \mathbb{Z}/2\mathbb{Z}$ precisely to the odd-finset parity torsor:
\[
  \partial(1) \;=\; [\mathcal{T}_{\mathrm{odd}}] \;\neq\; 0 \;\in\; H^1(\mathbf{Nom}; \mathrm{EvenFinset}).
\]
No permutation-invariant global element of $\mathcal{T}_{\mathrm{odd}}$ exists ($\Gamma(\mathcal{T}_{\mathrm{odd}}) = \emptyset$, since any transposition $(a\; b)$ moving $a \in S$ to $b \notin S$ moves $S$), whereas Schreier--Ulam rigidity kills all constant-coefficient cohomology ($H^1(\mathbf{Nom}; \mathbb{Z}/2\mathbb{Z}) = 0$, \emph{pro-discrete blindness}). Thus $\mathcal{T}_{\mathrm{odd}}$ proves concretely why the class $\mathcal{G}_{\top}$ below must include non-constant internal group objects (Chapter~\ref{ch:realizability-sites}, \S\ref{subsec:kihara-choice-nominal}).

Generalizing Andreas Blass's theorem that first cohomology detects failures of $\mathrm{AC}$ over discrete sets in $\mathrm{ZF}$~\cite{Blass83}, when a Grothendieck topos $\mathcal{E}$ already satisfies a weak choice principle $A$, the classes $\omega(e)$ for certain structure-group families already vanish, leaving a restricted slice assignment $\mathcal{G}_A$ to govern $(A \Rightarrow_{\mathbf{Topos}} \mathrm{AC})$:

\begin{conjecture}[The Cohomological Vanishing Conjecture for Choice Residuums]\label{conj:cohom-vanishing}

For every standard weak choice axiom $A$ (such as $\top$, $\mathrm{BPI}$, $\mathrm{DC}$, or $\mathrm{SVC}$), there exists a canonical assignment $Y \mapsto \mathcal{G}_A(Y) \subseteq \mathbf{Grp}(\mathcal{E}/Y)$ of internal structure groups such that $(A \Rightarrow_{\mathbf{Topos}} \mathrm{AC})$ in a Grothendieck topos $\mathcal{E}$ (with an $\mathrm{NNO}$) is equivalent to internal Booleanness ($\Omega_{\mathcal{E}} \simeq 1+1$), external projectivity of $1_{\mathcal{E}}$, and the cohomological vanishing condition:

\[
  H^1(\mathcal{E}/Y, G) = 0 \quad \text{for all } Y \in \mathcal{E} \text{ and all } G \in \mathcal{G}_A(Y).
\]

\end{conjecture}

Remarkably, across all four benchmark axioms $A \in \{\top, \mathrm{IBPI}, \mathrm{IDC}, \mathrm{SVC}\}$ (while leaving open intermediate axioms such as $\mathrm{WISC}$, $\mathrm{OP}$, and Kihara's $\mathrm{AC}_\lambda$ spectrum), canonical assignments $\mathcal{G}_A(Y)$ admit unconditional reductions:

\begin{itemize}[leftmargin=1.5em]

  \item \textbf{Universal Local Choice ($A = \top$, Unconditional Baseline):} Taking $\mathcal{G}_\top(Y) = \mathbf{Grp}(\mathcal{E}/Y)$, Conjecture~\ref{conj:cohom-vanishing} holds unconditionally (and in fact requires only free-abelian augmentation kernels and finite symmetric groups). Necessity is immediate: $\mathrm{AC}$ implies $\Omega_{\mathcal{E}} \simeq 1+1$ and external projectivity of $1_{\mathcal{E}}$ (Theorems~\ref{thm:diaconescu} and~\ref{thm:iac-to-ac}), and splits every $G$-torsor $P \twoheadrightarrow Y$~\cite{Giraud71}. Conversely, let $e : X \twoheadrightarrow Y$ be epic in $\mathcal{E}$, i.e.\ $X \twoheadrightarrow 1_{\mathcal{E}'}$ in the Boolean slice topos $\mathcal{E}' := \mathcal{E}/Y$ (where equality is decidable). Inside $\mathbb{Z}^X$, the subgroup $\mathbb{Z}^{(X)}$ of functions bounded by a Kuratowski-finite $S \in K(X)$ has decidable, cardinal-finite supports $\mathrm{supp}(f) = \{x \in X \mid f(x) \neq 0\} \in K(X)$~\cite[\S\S D5.2,~D5.4]{Johnstone02}, inducing an augmentation epimorphism $\varepsilon : \mathbb{Z}^{(X)} \twoheadrightarrow \mathbb{Z}_{\mathcal{E}'}$ ($\delta_x \mapsto 1$) with abelian kernel $K_X := \ker(\varepsilon) \in \mathbf{Ab}(\mathcal{E}')$. Since $H^1(\mathcal{E}/Y, K_X) = 0$, the $K_X$-torsor $\varepsilon^{-1}(1) \twoheadrightarrow 1_{\mathcal{E}'}$ admits a global section $\sigma$, whose support $F := \mathrm{supp}(\sigma) \rightarrowtail X$ is inhabited ($\varepsilon(\sigma) = 1$) of unique finite cardinal $c : 1_{\mathcal{E}'} \to \mathbb{N}_{\ge 1}$. Vanishing of $H^1(\mathcal{E}/Y, \underline{\mathrm{Aut}}_{\mathcal{E}'}([c]))$ then splits the finite-frame torsor $\underline{\mathrm{Iso}}_{\mathcal{E}'}([c], F) \twoheadrightarrow 1_{\mathcal{E}'}$ via $\varphi$, and $s := \varphi(0) : Y \to F \rightarrowtail X$ splits $e$ (forcing external projectivity of $1_{\mathcal{E}}$ at $Y = 1_{\mathcal{E}}$, whereas for $A = \mathrm{IAC}$ with $\mathcal{G}_{\mathrm{IAC}}(Y) = \emptyset$, external projectivity of $1_{\mathcal{E}}$ remains independent by $\mathbf{Set}^{\mathbb{Z}/2\mathbb{Z}}$).

  \item \textbf{Internal Boolean Prime Ideal Theorem ($A = \mathrm{IBPI}$, Unconditional Reduction):} Mere compactness of $\mathrm{Sub}_{\mathcal{E}}(1_{\mathcal{E}})$ does not suffice (failing in $\mathbf{Nom}$ and $\mathcal{F}$), whereas $\mathrm{IBPI}$ (every non-trivial internal Boolean algebra $B \in \mathbf{BA}(\mathcal{E})$ admits an ultrafilter~\cite{Deligne72,Johnstone82}) together with $\Omega_{\mathcal{E}} \simeq 1+1$ and external projectivity of $1_{\mathcal{E}}$ kills all finitary and cofiltered-profinite torsor obstructions, reducing $\mathcal{G}_{\mathrm{IBPI}}(Y)$ unconditionally to the discrete torsion-free abelian augmentation kernels $K_X = \ker(\mathbb{Z}^{(X)} \twoheadrightarrow \mathbb{Z}_{\mathcal{E}/Y}) \in \mathbf{Ab}(\mathcal{E}/Y)$. Indeed, for any epimorphism $p : F \twoheadrightarrow Y$ with Kuratowski-finite fibers in $\mathcal{E}/Y$ (such as $F = \mathrm{supp}(\sigma)$ from $H^1(\mathcal{E}/Y, K_X) = 0$ above), let $A \rightarrowtail 2^{2^F}$ be the internal Boolean subalgebra generated by $\{\mathrm{ev}_x \mid x \in F\}$ and let $I \rightarrowtail A$ be the ideal generated by $\mathrm{ev}_{x_1} \wedge \mathrm{ev}_{x_2}$ ($x_1 \neq x_2$ in $F \times_Y F$) and $\neg \bigvee_{x \in F_y} \mathrm{ev}_x$ ($y \in Y$). By Kuratowski induction under $\Omega_{\mathcal{E}} \simeq 1+1$~\cite[\S D5.4]{Johnstone02}, every $S \in K(Y)$ admits an internal partial section $\sigma_S : S \to F$ whose indicator $h \in 2^F$ yields an evaluation homomorphism $\mathrm{Ev}_h : 2^{2^F} \to 2$ killing all generators over $S$ with $\mathrm{Ev}_h(1_A) = 1$; hence $\mathcal{E} \models \neg(0_{A/I} = 1_{A/I})$. Applying $\mathrm{IBPI}$ to $B := A/I$ renders $\mathrm{Hom}_{\mathbf{BA}}(B, 2) \twoheadrightarrow 1_{\mathcal{E}}$ epic, and sending each ultrafilter $u$ to the unique $x \in F_y$ with $u([\mathrm{ev}_x]) = 1$ yields an epimorphism $\mathrm{Sect}_Y(F) \twoheadrightarrow 1_{\mathcal{E}}$ that splits by external projectivity of $1_{\mathcal{E}}$ (and identically along cofiltered profinite towers via transition ideals, while Halpern--L\"auchli models show $K_X$ itself cannot be omitted).

  \item \textbf{Internal Dependent Choice ($A = \mathrm{IDC}$, Unconditional Base and Tower Reduction):} Under $\mathrm{IDC}$, $\Omega_{\mathcal{E}} \simeq 1+1$, and external projectivity of $1_{\mathcal{E}}$, $H^1(\mathcal{E}/Y, G) = 0$ vanishes for all $G \in \mathbf{Grp}(\mathcal{E}/Y)$ over every internally subcountable codomain $Y$---that is, any quotient $U \twoheadrightarrow Y$ of $U \rightarrowtail \mathbb{N}_{\mathcal{E}}$. Indeed, least-number minimization~\cite[\S D5.1]{Johnstone02} embeds $Y \rightarrowtail \mathbb{N}_{\mathcal{E}}$ as a complemented summand, and $\mathrm{IDC}$ splits the prefix-section $\omega$-tower $\mathrm{Sect}_{[n]}(\widetilde{X}|_{[n]})$, yielding $\mathcal{G}_{\mathrm{IDC}}(Y) = \emptyset$. Moreover, sequential $\varprojlim^1$ obstructions vanish along any inverse $\omega$-tower $G \cong \varprojlim_{n \in \omega} G_n$ ($G_0 = 1$) whose pullback-stable kernels $K_n = \ker(G_{n+1} \twoheadrightarrow G_n)$ have vanishing $H^1$, because $\mathrm{Sect}_Y(P_{n+1}) \twoheadrightarrow \mathrm{Sect}_Y(P_n)$ is then epic. By the $A = \top$ decomposition above, $\mathcal{G}_{\mathrm{IDC}}(Y)$ reduces unconditionally to augmentation kernels $K_X$ and finite symmetric groups $\underline{\mathrm{Aut}}_{\mathcal{E}/Y}([c])$ over internally uncountable codomains $Y$. Conversely, fiber groups cannot be restricted to uncountably presented groups: in $\mathrm{ZF} + \mathrm{DC}$ models where all subsets of $2^{\mathbb{N}}$ have the Property of Baire, the Pauli-$Z$ parity cover $2^{\mathbb{N}}/E_0^{\mathrm{even}} \twoheadrightarrow \mathcal{Q}_{E_0}$ (Remark~\ref{rem:pauli-z-parity}) is a non-split torsor over $Y = \mathcal{Q}_{E_0}$ under the finite group $\Delta(\mathbb{Z}/2\mathbb{Z})$~\cite{Serafin23}, since any flip-alternating hat predictor $h : 2^{\mathbb{N}} \to \mathbb{Z}/2\mathbb{Z}$ swaps comeager preimages under an unconstrained bit flip on a basic cylinder, contradicting Baire category under $\mathrm{DC}$.

  \item \textbf{Small Violations of Choice ($A = \mathrm{SVC}_{\mathrm{ret}}$ and Bare $\mathrm{SVC}$, Unconditional Single-Slice Reductions):} Under retract-bounded $\mathrm{SVC}_{\mathrm{ret}}(S)$ over $\mathbf{Set}_{\mathrm{AC}}$---where every $Y \in \mathcal{E}$ is a retract of $U \rightarrowtail \Delta I \times S$ (Theorem~\ref{thm:coshep-svc})---setting $\mathcal{G}_{\mathrm{SVC}_{\mathrm{ret}}}(Y) = \emptyset$ for $Y \neq S$ and restricting $\mathcal{G}_{\mathrm{SVC}_{\mathrm{ret}}}(S)$ to $K_X$ and $\underline{\mathrm{Aut}}_{\mathcal{E}/S}([c])$ suffices unconditionally. Indeed, $H^1(\mathcal{E}/S, G) = 0$ makes $S$ externally projective by the $A = \top$ reduction, whence $\Delta I \times S \cong \coprod_{i \in I} S$, its complemented subobjects, and their retracts are externally projective (Lemmas~\ref{thm:coshep-svc}(a)--(b)). For Blass's bare subquotient $\mathrm{SVC}(S)$~\cite{Blass79}, where every $X \in \mathcal{E}$ is a quotient $q_X : U_X \twoheadrightarrow X$ of $U_X \rightarrowtail \Delta\kappa \times S$ for an ordinal $\kappa$, unsplit quotients block $Y = S$ alone (the quotient fallacy), yet shifting the single test slice to the inhabited-subobject classifier $P^+(S) \rightarrowtail 2^S$ ($A \in P^+(S) \iff \exists s \in A$, with $\mathcal{G}_{\mathrm{SVC}}(Y) = \emptyset$ for $Y \neq P^+(S)$) suffices unconditionally over any base $\mathbf{Set}$. Vanishing of $H^1(\mathcal{E}/P^+(S), G)$ splits membership onto $P^+(S)$ via $\gamma : P^+(S) \to S$, internally well-ordering $S$ by Zermelo's theorem under $\Omega_{\mathcal{E}} \simeq 1+1$~\cite[\S D4.5]{Johnstone02}. Transfinite induction on $V_\alpha = U_\alpha \wedge \neg \bigvee_{\beta < \alpha} U_\beta$ in the complete Boolean algebra $\mathrm{Sub}_{\mathcal{E}}(Z)$ internally well-orders $\Delta\kappa = \coprod_{\alpha < \kappa} 1_{\mathcal{E}}$ across slices, so $<_{\mathrm{lex}}$ well-orders $U_X$, $x \mapsto \min_{<_{\mathrm{lex}}} q_X^{-1}(x)$ embeds $X \rightarrowtail U_X$, and $y \mapsto \min_{<_X} e^{-1}(y)$ splits every epimorphism $e : X \twoheadrightarrow Y$. Meanwhile, $\mathbf{Set}^{\mathbb{Z}/2\mathbb{Z}}$ confirms both the necessity of external projectivity of $1_{\mathcal{E}}$ at $Y = S$ and $H^1(\mathcal{E}/P^+(S), \Delta(\mathbb{Z}/2\mathbb{Z})) \neq 0$ on $\{0,1\} \in P^+(S)$.

\end{itemize}

\section{\texorpdfstring{Symmetric Boolean-Valued Submodels, Equivariant Boolean Sheaves, and Pure-$\mathsf{ZF}$ Choice Separations}{Symmetric Boolean-Valued Submodels, Equivariant Boolean Sheaves, and Pure-ZF Choice Separations}}\label{sec:symmetric-bv-zf}

A foundational metamathematical subtlety governs how the intermediate rungs of Sections~2.1--2.4 are separated by concrete models. Over a classical ground universe $V \models \mathsf{ZFC}$, for \emph{every} complete Boolean algebra $\mathbb{B}$, the full Boolean-valued universe $V^{\mathbb{B}}$ (and its localic sheaf topos $\mathbf{Sh}(\mathbb{B})$) validates the full Axiom of Choice ($\llbracket \mathsf{ZFC} \rrbracket^{\mathbb{B}} = 1_{\mathbb{B}}$) by the Barr--Diaconescu--Freyd theorem (Theorem~\ref{thm:bdf}): raw $V^{\mathbb{B}}$ \emph{never} separates weak choice principles over $V \models \mathsf{ZFC}$. Conversely, Fraenkel--Mostowski--Specker atomic permutation topoi $\mathbf{Cont}(\mathcal{G})$ (Chapter~\ref{ch:realizability-sites}, \S\ref{sec:permutation-topoi}) separate choice principles transparently via open-subgroup stabilizers, but model $\mathsf{ZFA} + \neg\mathrm{AC}$ \emph{with urelements} (violating Osius pure-set fullness on the atom object $\mathbb{A}$; Chapter~\ref{ch:sets-residuum}, Theorem~\ref{thm:scale-invariance}(i)).

To calibrate pure well-founded set theory ($\mathsf{ZF}$ without atoms) alongside its Grothendieck-topos incarnation, Cohen--Scott--Solovay symmetric forcing marries complete Boolean algebras $\mathbb{B}$ with continuous group symmetry $(\mathcal{G}, \mathcal{F})$ inside $V^{\mathbb{B}}$~\cite[Ch.~5]{Jech73}.

\begin{definition}[Symmetric Forcing Datum and Hereditarily Symmetric Submodel $\mathrm{HS}_{\mathcal{F}} \subset V^{\mathbb{B}}$]\label{def:symmetric-bv-submodel}%
\uses{thm:bdf}%
Let $V \models \mathsf{ZFC}$ be the ground universe. A \emph{symmetric datum} $(\mathbb{B}, \mathcal{G}, \mathcal{F})$ consists of a complete Boolean algebra $\mathbb{B}$, a group $\mathcal{G}$ acting on $\mathbb{B}$ by complete Boolean algebra automorphisms ($\mathcal{G} \to \mathrm{Aut}(\mathbb{B})$, faithful $\mathcal{G} \le \mathrm{Aut}(\mathbb{B})$ in pure-set forcing), and a \emph{normal filter} $\mathcal{F}$ of subgroups of $\mathcal{G}$---a non-empty upward-closed family ($\mathcal{G} \in \mathcal{F}$) closed under finite intersections ($H_1, H_2 \in \mathcal{F} \implies H_1 \cap H_2 \in \mathcal{F}$) and conjugation ($\pi \in \mathcal{G},\, H \in \mathcal{F} \implies \pi H \pi^{-1} \in \mathcal{F}$). Extending each $\pi \in \mathcal{G}$ recursively across $V^{\mathbb{B}} = \bigcup_{\alpha \in \mathrm{Ord}} V^{\mathbb{B}}_\alpha$ by $\mathrm{dom}(\pi\tau) := \{\pi\sigma \mid \sigma \in \mathrm{dom}(\tau)\}$ and $(\pi\tau)(\pi\sigma) := \pi(\tau(\sigma))$, a $\mathbb{B}$-name $\tau \in V^{\mathbb{B}}$ is \emph{$\mathcal{F}$-symmetric} if its stabilizer $\mathrm{sym}_{\mathcal{G}}(\tau) := \{\pi \in \mathcal{G} \mid \pi\tau = \tau\}$ lies in $\mathcal{F}$, and \emph{hereditarily $\mathcal{F}$-symmetric} ($\tau \in \mathrm{HS}_{\mathcal{F}}$) if $\mathrm{sym}_{\mathcal{G}}(\tau) \in \mathcal{F}$ and $\mathrm{dom}(\tau) \subseteq \mathrm{HS}_{\mathcal{F}}$~\cite[\S 5.2]{Jech73}.
\end{definition}

Every canonical ground name $\check{x} := \{(\check{y}, 1_{\mathbb{B}}) \mid y \in x\}$ ($x \in V$) satisfies $\pi\check{x} = \check{x}$ for all $\pi \in \mathcal{G}$, so $\mathrm{sym}_{\mathcal{G}}(\check{x}) = \mathcal{G} \in \mathcal{F}$ and $\check{V} \subseteq \mathrm{HS}_{\mathcal{F}}$. Crucially, Boolean truth in $\mathrm{HS}_{\mathcal{F}}$ must be \emph{relativized} to $\mathrm{HS}_{\mathcal{F}}$ rather than inherited unrestrictedly from $V^{\mathbb{B}}$: for $\mathcal{L}_\in$-formulas with parameters in $\mathrm{HS}_{\mathcal{F}}$, we evaluate Boolean connectives ($\neg, \wedge, \vee, \to$) in $\mathbb{B}$ and define unbounded quantifiers over $\mathrm{HS}_{\mathcal{F}}$ (well-defined in $\mathbb{B}$ because $\mathbb{B} \in V$ is a set, so $\{\llbracket \psi(\tau, \vec{\sigma}) \rrbracket^{\mathrm{HS}} \mid \tau \in \mathrm{HS}_{\mathcal{F}}\} \subseteq \mathbb{B}$ is set-sized by Separation in $V$):
\[
  \begin{aligned}
    \llbracket \exists x\,\psi(x, \vec{\sigma}) \rrbracket^{\mathrm{HS}} &\;:=\; \bigvee_{\tau \in \mathrm{HS}_{\mathcal{F}}} \llbracket \psi(\tau, \vec{\sigma}) \rrbracket^{\mathrm{HS}}, \\
    \llbracket \forall x\,\psi(x, \vec{\sigma}) \rrbracket^{\mathrm{HS}} &\;:=\; \bigwedge_{\tau \in \mathrm{HS}_{\mathcal{F}}} \llbracket \psi(\tau, \vec{\sigma}) \rrbracket^{\mathrm{HS}}.
  \end{aligned}
\]

\begin{theorem}[Relativized Symmetric Boolean Truth, Equivariant Sheaf Realization, and Pure-$\mathsf{ZF}$ Choice Calibrations]\label{thm:symmetric-bv-zf}\index{Symmetric Boolean-valued submodels ($\mathrm{HS}_{\mathcal{F}}$)}\index{Equivariant Boolean sheaves@$\mathcal{F}$-continuous $\mathcal{G}$-equivariant Boolean sheaves $\mathbf{Sh}_{\mathcal{G},\mathcal{F}}(\mathbb{B})$}%
\uses{def:symmetric-bv-submodel,thm:bdf,thm:coshep-svc}%
Let $(\mathbb{B}, \mathcal{G}, \mathcal{F})$ be a symmetric datum over $V \models \mathsf{ZFC}$. Then:
\begin{enumerate}[label=\textup{(\roman*)},leftmargin=2.0em]
  \item \textbf{$\Delta_0$-Absoluteness, Symmetry Lemma, and Full $\mathsf{ZF}$ Soundness:} For all $\sigma, \tau \in \mathrm{HS}_{\mathcal{F}}$ and every bounded ($\Delta_0$) $\mathcal{L}_\in$-formula $\varphi_0(\vec{\sigma})$, relativized truth agrees on the nose with ambient Boolean truth in $V^{\mathbb{B}}$,
  \[
    \begin{aligned}
      \llbracket \sigma \in \tau \rrbracket^{\mathrm{HS}} &= \llbracket \sigma \in \tau \rrbracket^{\mathbb{B}}, \qquad
      \llbracket \sigma = \tau \rrbracket^{\mathrm{HS}} = \llbracket \sigma = \tau \rrbracket^{\mathbb{B}}, \\
      \llbracket \varphi_0(\vec{\sigma}) \rrbracket^{\mathrm{HS}} &= \llbracket \varphi_0(\vec{\sigma}) \rrbracket^{\mathbb{B}},
    \end{aligned}
  \]
  whereas unbounded quantifiers do \emph{not} in general agree with $\llbracket - \rrbracket^{\mathbb{B}}$ (permitting $\llbracket \mathsf{AC} \rrbracket^{\mathrm{HS}} < 1_{\mathbb{B}} = \llbracket \mathsf{AC} \rrbracket^{\mathbb{B}}$). For every $\mathcal{L}_\in$-formula $\varphi(x_1, \dots, x_n)$, all $\sigma_1, \dots, \sigma_n \in \mathrm{HS}_{\mathcal{F}}$, and all $\pi \in \mathcal{G}$, relativized truth is $\mathcal{G}$-equivariant:
  \[
    \pi\bigl(\llbracket \varphi(\sigma_1, \dots, \sigma_n) \rrbracket^{\mathrm{HS}}\bigr) \;=\; \llbracket \varphi(\pi\sigma_1, \dots, \pi\sigma_n) \rrbracket^{\mathrm{HS}},
  \]
  and $\mathrm{HS}_{\mathcal{F}}$ validates all axioms of $\mathsf{ZF}$ ($\llbracket \chi \rrbracket^{\mathrm{HS}} = 1_{\mathbb{B}}$ for every $\chi \in \mathsf{ZF}$;~\cite[Thm.~5.14]{Jech73}).
  \item \textbf{$\mathcal{F}$-Continuous Equivariant Boolean Sheaf Realization $\mathbf{Sh}_{\mathcal{G}, \mathcal{F}}(\mathbb{B})$:} Let $\mathbf{Sh}_{\mathcal{G}, \mathcal{F}}(\mathbb{B})$ denote the Grothendieck topos of \emph{$\mathcal{F}$-continuous $\mathcal{G}$-equivariant Boolean sheaves} on $\mathbb{B}$: sheaves $\mathcal{X} \in \mathbf{Sh}(\mathbb{B})$ equipped with a semilinear $\mathcal{G}$-action $\pi_b : \mathcal{X}(b) \xrightarrow{\sim} \mathcal{X}(\pi b)$ such that every section $s \in \mathcal{X}(b)$ is covered ($b = \bigvee_{i \in I} b_i$) by restrictions $s|_{b_i}$ whose joint stabilizer $\mathrm{Stab}_{\mathcal{G}}(b_i, s|_{b_i}) := \{\pi \in \mathcal{G} \mid \pi b_i = b_i \text{ and } \pi_{b_i}(s|_{b_i}) = s|_{b_i}\}$ lies in $\mathcal{F}$~\cite[\S D4.5]{Johnstone02}. Then $\mathbf{Sh}_{\mathcal{G}, \mathcal{F}}(\mathbb{B})$ is a Boolean Grothendieck topos unifying $\mathbf{Sh}(\mathbb{B})$ (at $\mathcal{G} = \{1\}$) and $\mathbf{Cont}(\mathcal{G})$ (at $\mathbb{B} = 2$ with open-subgroup filter $\mathcal{F} = \mathcal{U}(\mathcal{G})$), whose internal well-founded pure-set hierarchy $\mu\mathcal{P}$ represents $(\mathrm{HS}_{\mathcal{F}}, \llbracket - \rrbracket^{\mathrm{HS}})$. Whenever $\mathcal{G}$ acts on an atomless complete Boolean algebra $\mathbb{B}$ with symmetric generic reals separating cosets, $\mathbf{Sh}_{\mathcal{G}, \mathcal{F}}(\mathbb{B})$ discharges local Osius pure-set fullness without urelements.
  \item \textbf{Pure-$\mathsf{ZF}$ Choice-Separation Benchmarks:} Varying $(\mathbb{B}, \mathcal{G}, \mathcal{F})$ separates the rungs of Chapter~\ref{ch:choice-atlas} inside pure well-founded $\mathsf{ZF}$:
  \begin{enumerate}[label=\textup{(\alph*)},leftmargin=1.6em]
    \item \emph{$\neg\mathrm{AC}_\omega$ (and $\neg\mathrm{DC}$):} Cohen's first symmetric model on $\mathbb{B} = \mathrm{RO}(\mathrm{Add}(\omega, \omega))$ with $\mathcal{G} = \mathrm{Sym}(\omega)$ and finite-support filter $\mathcal{F}_{\mathrm{fin}}$ refutes $\mathrm{AC}_\omega$, while the Feferman--L\'evy symmetric collapse of $\aleph_n$ ($\mathrm{Col}(\omega, <\!\aleph_\omega)$) makes $\mathbb{R}$ a countable union of countable sets, refuting \emph{both} $\mathrm{AC}_\omega$ and $\mathrm{DC}$~\cite[Chs.~5, 10]{Jech73}.
    \item \emph{$\mathrm{DC} + \neg\mathrm{AC}$:} Jech's countably supported symmetric model on countably closed forcing $\mathbb{B} = \mathrm{RO}(\mathrm{Add}(\omega_1, \omega_1))$ with $\mathcal{G} = \mathrm{Sym}(\omega_1)$ and the $\sigma$-complete normal filter $\mathcal{F}_{\le \aleph_0}$ of countable supports validates $\mathsf{ZF} + \mathrm{DC}$ (and $\mathrm{AC}_\omega$) while refuting $\mathrm{AC}$~\cite[Ch.~8]{Jech73}.
    \item \emph{$\mathrm{BPI} + \neg\mathrm{AC}$ and $\mathrm{SVC}$:} The Halpern--L\'evy model on $\mathrm{RO}(\mathrm{Add}(\omega, \omega))$ under $\mathcal{G} \cong \mathrm{Aut}(\mathbb{Q}, <)$ with finite supports validates $\mathsf{ZF} + \mathrm{BPI} + \neg\mathrm{AC}$~\cite{HalpernLevy71} (\cite[Ch.~7]{Jech73}), and Blass's theorem~\cite[Thm.~4.3]{Blass79} characterizes $\mathsf{ZF} + \mathrm{SVC}$ precisely as symmetric extensions $\mathrm{HS}_{\mathcal{F}}$ of $\mathsf{ZFC}$ ground models.
  \end{enumerate}
\end{enumerate}
\end{theorem}

\begin{proof}
For~(i), domain-transitivity ($\tau \in \mathrm{HS}_{\mathcal{F}} \implies \mathrm{dom}(\tau) \subseteq \mathrm{HS}_{\mathcal{F}}$) ensures that the mutual well-founded rank recursions defining $\llbracket \sigma \in \tau \rrbracket^{\mathbb{B}} = \bigvee_{u \in \mathrm{dom}(\tau)} (\tau(u) \wedge \llbracket \sigma = u \rrbracket^{\mathbb{B}})$ and $\llbracket \sigma = \tau \rrbracket^{\mathbb{B}} = \bigwedge_{u \in \mathrm{dom}(\sigma)} (\sigma(u) \to \llbracket u \in \tau \rrbracket^{\mathbb{B}}) \wedge \bigwedge_{v \in \mathrm{dom}(\tau)} (\tau(v) \to \llbracket v \in \sigma \rrbracket^{\mathbb{B}})$ range exclusively over names already in $\mathrm{HS}_{\mathcal{F}}$, giving $\llbracket \sigma \in \tau \rrbracket^{\mathrm{HS}} = \llbracket \sigma \in \tau \rrbracket^{\mathbb{B}}$ and $\llbracket \sigma = \tau \rrbracket^{\mathrm{HS}} = \llbracket \sigma = \tau \rrbracket^{\mathbb{B}}$ by rank induction. Structural induction on $\psi$ verifies equality substitution $\llbracket \rho = u \wedge \psi(\rho, \vec{\sigma}) \rrbracket^{\mathrm{HS}} \le \llbracket \psi(u, \vec{\sigma}) \rrbracket^{\mathrm{HS}}$ across Boolean connectives and $\bigvee_{\xi \in \mathrm{HS}_{\mathcal{F}}}$, reducing bounded quantification $\llbracket (\exists x \in \tau)\,\psi(x, \vec{\sigma}) \rrbracket^{\mathrm{HS}}$ to $\bigvee_{u \in \mathrm{dom}(\tau)} (\tau(u) \wedge \llbracket \psi(u, \vec{\sigma}) \rrbracket^{\mathrm{HS}})$ and establishing $\Delta_0$-absoluteness. Next, if $\tau \in \mathrm{HS}_{\mathcal{F}}$ and $\pi \in \mathcal{G}$, rank induction gives $\mathrm{dom}(\pi\tau) = \pi[\mathrm{dom}(\tau)] \subseteq \mathrm{HS}_{\mathcal{F}}$, while $\gamma(\pi\tau) = \pi\tau \iff (\pi^{-1}\gamma\pi)\tau = \tau$ yields $\mathrm{sym}_{\mathcal{G}}(\pi\tau) = \pi\,\mathrm{sym}_{\mathcal{G}}(\tau)\,\pi^{-1} \in \mathcal{F}$ by normality of $\mathcal{F}$; thus $\tau \mapsto \pi\tau$ is a rank-preserving bijection $\mathrm{HS}_{\mathcal{F}} \xrightarrow{\sim} \mathrm{HS}_{\mathcal{F}}$. Since $\pi \in \mathrm{Aut}(\mathbb{B})$ preserves arbitrary joins and meets, induction on formula complexity gives
\[
  \begin{aligned}
    \pi\bigl(\llbracket \exists x\,\psi(x, \vec{\sigma}) \rrbracket^{\mathrm{HS}}\bigr)
    &\;=\; \bigvee_{\tau \in \mathrm{HS}_{\mathcal{F}}} \pi\bigl(\llbracket \psi(\tau, \vec{\sigma}) \rrbracket^{\mathrm{HS}}\bigr) \\
    &\;=\; \bigvee_{\tau \in \mathrm{HS}_{\mathcal{F}}} \llbracket \psi(\pi\tau, \pi\vec{\sigma}) \rrbracket^{\mathrm{HS}}
    \;=\; \llbracket \exists x\,\psi(x, \pi\vec{\sigma}) \rrbracket^{\mathrm{HS}},
  \end{aligned}
\]
proving the Symmetry Lemma. Extensionality and Foundation hold in $\mathrm{HS}_{\mathcal{F}}$ because $\mathrm{dom}(\sigma), \mathrm{dom}(\tau) \subseteq \mathrm{HS}_{\mathcal{F}}$; Empty Set and Infinity are witnessed by $\check{\emptyset}, \check{\omega} \in \mathrm{HS}_{\mathcal{F}}$ ($\mathrm{sym}_{\mathcal{G}} = \mathcal{G}$ and $\Delta_0$-absoluteness); and $\{\sigma, \tau\}_{\mathbb{B}}$ and $\bigcup_{\mathbb{B}} \sigma$ have stabilizers containing $\mathrm{sym}_{\mathcal{G}}(\sigma) \cap \mathrm{sym}_{\mathcal{G}}(\tau) \in \mathcal{F}$ and $\mathrm{sym}_{\mathcal{G}}(\sigma) \in \mathcal{F}$. For unbounded Separation on $\dot{a}, p_1, \dots, p_k \in \mathrm{HS}_{\mathcal{F}}$, set $\mathrm{dom}(\tau) := \mathrm{dom}(\dot{a}) \subseteq \mathrm{HS}_{\mathcal{F}}$ and $\tau(u) := \dot{a}(u) \wedge \llbracket \varphi(u, \vec{p}) \rrbracket^{\mathrm{HS}}$; for every $\pi \in H := \mathrm{sym}_{\mathcal{G}}(\dot{a}) \cap \bigcap_{i=1}^k \mathrm{sym}_{\mathcal{G}}(p_i) \in \mathcal{F}$, the Symmetry Lemma yields $(\pi\tau)(\pi u) = \pi(\dot{a}(u)) \wedge \llbracket \varphi(\pi u, \pi\vec{p}) \rrbracket^{\mathrm{HS}} = \tau(\pi u)$, so $H \le \mathrm{sym}_{\mathcal{G}}(\tau)$ and $\tau \in \mathrm{HS}_{\mathcal{F}}$. For Power Set on $\dot{a} \in \mathrm{HS}_{\mathcal{F}}$, every $\rho \in \mathrm{HS}_{\mathcal{F}}$ with $\llbracket \rho \subseteq \dot{a} \rrbracket^{\mathrm{HS}}$ restricts to $\tilde{\rho} \in \mathbb{B}^{\mathrm{dom}(\dot{a})} \cap \mathrm{HS}_{\mathcal{F}}$ via $\tilde{\rho}(u) := \dot{a}(u) \wedge \llbracket u \in \rho \rrbracket^{\mathrm{HS}}$, and $P := \{(\sigma, 1_{\mathbb{B}}) \mid \sigma \in \mathbb{B}^{\mathrm{dom}(\dot{a})} \cap \mathrm{HS}_{\mathcal{F}}\}$ is set-sized with $\mathrm{sym}_{\mathcal{G}}(\dot{a}) \le \mathrm{sym}_{\mathcal{G}}(P) \in \mathcal{F}$. For Collection/Replacement on $\dot{a}, \vec{p} \in \mathrm{HS}_{\mathcal{F}}$, $\mathrm{dom}(\dot{a}) \times \mathbb{B} \in V$ is a small set; letting $\alpha(u, b)$ be the least rank of a witness $\tau \in \mathrm{HS}_{\mathcal{F}}$ achieving $\llbracket \varphi(u, \tau, \vec{p}) \rrbracket^{\mathrm{HS}} = b$ when one exists, Replacement in $V \models \mathsf{ZFC}$ bounds $\beta := \sup_{(u,b) \in \mathrm{dom}(\dot{a}) \times \mathbb{B}} (\alpha(u, b) + 1) \in \mathrm{Ord}$, and the rank-truncated level $B_\beta := \{(\tau, 1_{\mathbb{B}}) \mid \tau \in \mathrm{HS}_{\mathcal{F}} \cap V_\beta\}$ satisfies $\mathrm{sym}_{\mathcal{G}}(B_\beta) = \mathcal{G} \in \mathcal{F}$ because every $\pi \in \mathcal{G}$ preserves rank and $\mathrm{HS}_{\mathcal{F}}$.

For~(ii), $\mathbf{Sh}_{\mathcal{G}, \mathcal{F}}(\mathbb{B})$ is equivalent to sheaves for the join topology on the site of coset-opens $(b, H)$ ($b \in \mathbb{B}^+$, $H \in \mathcal{F}$, $H \le \mathrm{Stab}_{\mathcal{G}}(b)$), inheriting $\Omega \simeq 1+1$ from Boolean $\mathbb{B}$~\cite[\S D4.5]{Johnstone02}. Quotienting $\mathrm{HS}_{\mathcal{F}}$ by $\llbracket \sigma = \tau \rrbracket^{\mathrm{HS}}$ with Boolean $\mathbb{B}$-extensionality gluing and $\mathcal{G}$-action $[\pi\sigma] = \pi[\sigma]$ embeds $\mathrm{HS}_{\mathcal{F}}$ as the internal cumulative hierarchy $\mu\mathcal{P}$ of $\mathbf{Sh}_{\mathcal{G}, \mathcal{F}}(\mathbb{B})$ (matching Kripke--Joyal forcing over $(b, H)$ with $\llbracket - \rrbracket^{\mathrm{HS}}$); whereas $\mathbb{B} = 2$ leaves $\mathrm{HS}_{\mathcal{F}} = \check{V}$ unable to represent non-trivial coset sheaves $\mathcal{G}/H \in \mathbf{Cont}(\mathcal{G})$, atomless Cohen/measure algebras $\mathbb{B}$ furnish distinct symmetric generic subsets $\{\dot{c}_i\} \subset \mathrm{HS}_{\mathcal{F}}$ whose stabilizer orbits embed $\mathcal{G}/H$ into $\mu\mathcal{P}$. Finally, (iii)(a)--(c) instantiate~(i) on the classical symmetric data of~\cite[Chs.~5, 7, 8, 10]{Jech73}, \cite{HalpernLevy71}, and~\cite{Blass79} (in particular, $\sigma$-completeness of $\mathcal{F}_{\le \aleph_0}$ together with countable closure of $\mathrm{Add}(\omega_1, \omega_1)$ closes $\mathrm{HS}_{\mathcal{F}_{\le \aleph_0}}$ under $\omega$-sequences $\bigcap_{n<\omega} \mathrm{sym}_{\mathcal{G}}(\tau_n) \in \mathcal{F}_{\le \aleph_0}$, validating $\mathrm{DC}$ while leaving the unordered $\omega_1$-sequence without a global well-ordering).
\end{proof}

Having mapped the $1$-categorical and symmetric-forcing residuums of Axis~I, dropping $0$-truncation of mapping spaces fractures Choice into the bigraded homotopy hierarchy $\mathrm{AC}_{n, m}$ over $(\infty, 1)$-topoi $\mathcal{X}$, developed in full in Chapter~\ref{ch:higher-crm} (\S\S\ref{sec:higher-choice}--\ref{sec:postnikov-residuum}).

\part{\texorpdfstring{Algebraic Frobenius Lifts, Base--Fiber Conservation, and Infinite Towers}{Algebraic Frobenius Lifts, Base-Fiber Conservation, and Infinite Towers}}

\chapter{\texorpdfstring{Joyal $\delta$-Rings, $2$-Torsion Lift Torsors ($R_2$ vs.\ \texttt{DualInt}), and Relative Slice Perfection}{Joyal delta-Rings, 2-Torsion Lift Torsors (R2 vs. DualInt), and Relative Slice Perfection}}
\chaptermark{Joyal $\delta$-Rings and Relative Slice Perfection}
\label{ch:delta-rings}
\chapterepigraph{``The notion of $\delta$-ring is an algebraic version of the notion of $p$-adic Witt vectors\dots\ free of any $p$-torsion pathology.''}{Andr\'e Joyal, \emph{$\delta$-anneaux et vecteurs de Witt} (1985)}

\section{\texorpdfstring{Joyal's $\delta$-Ring Axioms and Polynomial Cofactors $C_p(x, y)$}{Joyal's delta-Ring Axioms and Polynomial Cofactors Cp(x, y)}}\label{sec:delta-rings}

Axis~IV of $\mathrm{Dom}(\mathrm{Res})$ (Chapter~\ref{ch:foundations}, Definition~\ref{def:comparison-triple} and Table~\ref{tab:five-axis-master}) demonstrates how Categorical Reverse Mathematics operates natively inside commutative algebra without requiring any ambient topos of sets. Comparing Andr\'e Joyal's~\cite{Joyal85} universal $\delta$-ring functor to naive Fermat quotient division $\frac{\phi(x) - x^p}{p}$ exposes a sharp $p$-torsion bifurcation between lift torsors ($R_2$), non-perfect torsion-free obstructions (\texttt{DualInt}), and relative colimit perfection~\cite{BhattScholze22,Goldthorpe23}.

Fix a prime $p$. Following Joyal~\cite{Joyal85} and Bhatt--Scholze~\cite{BhattScholze22} (\S2), a \emph{$p$-typical $\delta$-ring}\index{delta-ring@$\delta$-ring!Joyal axioms} is a commutative ring $A$ equipped with a set map $\delta : A \to A$ satisfying $\delta(0) = 0$, $\delta(1) = 0$, and the integral polynomial identities
\[
  \begin{aligned}
    \delta(x + y) &\;=\; \delta(x) + \delta(y) + C_p(x, y), \\
    \delta(xy) &\;=\; x^p\delta(y) + y^p\delta(x) + p\,\delta(x)\delta(y),
  \end{aligned}
\]
where $C_p(x, y) := \frac{x^p + y^p - (x+y)^p}{p} = -\sum_{k=1}^{p-1} \frac{1}{p}\binom{p}{k} x^k y^{p-k} \in \mathbb{Z}[x, y]$ is the universal binomial cofactor polynomial. Because $C_p(x, y)$ has integral coefficients, these identities make no use of division by $p$, yet unconditionally force the associated map $\phi(x) := x^p + p\,\delta(x)$ to be a ring endomorphism $A \to A$ (the \emph{Frobenius lift}) satisfying $\phi(x) \equiv x^p \pmod{pA}$ and $\delta \circ \phi = \phi \circ \delta$.

At $p = 2$, $C_2(x, y) = -xy$, simplifying the additive axiom to $\delta(x + y) = \delta(x) + \delta(y) - xy$ (\texttt{DeltaRing2} in \texttt{HomotopyCRM/\allowbreak DeltaRing\allowbreak Perfection.lean}). Over $\mathbb{Z}$---which is initial in $\mathbf{CRing}$---the unique ring endomorphism is $\phi = \mathrm{id}_{\mathbb{Z}}$ (since $x \mapsto x^2$ fails additivity), forcing the unique $2$-typical $\delta$-structure $\delta_{\mathbb{Z}}(a) = \frac{a - a^2}{2}$ (\texttt{instDeltaRing2Int}, \texttt{instUniqueDeltaRing2Int}), which is already perfect ($\phi = \mathrm{id}_{\mathbb{Z}}$, \texttt{isPerfectDeltaRing\_int}).

\section{\texorpdfstring{The $2$-Torsion Lift Torsor $R_2$ vs.\ Torsion-Free \texttt{DualInt}}{The 2-Torsion Lift Torsor R2 vs. Torsion-Free DualInt}}\label{sec:r2-vs-dualint}

When a commutative ring possesses $p$-torsion, the comparison $\kappa : \delta(x) \mapsto \frac{\phi(x) - x^p}{p}$ breaks down, precipitating two complementary obstructions:

\begin{proposition}[Frobenius-Twisted Derivation Pseudo-Torsor and Non-Isomorphic $2$-Torsion $\delta$-Lifts on $R_2$ vs.\ \texttt{DualInt}]\label{prop:r2-delta-torsor}%
\uses{def:comparison-triple,def:residuum}%
Let $p$ be a prime and $(A, \phi)$ a commutative ring equipped with a Frobenius lift $\phi(x) \equiv x^p \pmod{pA}$.
\begin{enumerate}[label=\textup{(\roman*)},leftmargin=2.0em]
  \item \textbf{$\mathrm{Der}_\phi(A, \mathrm{Ann}_A(p))$-Pseudo-Torsor of Lifts and $R_2$ Separation:}\index{R2@$R_2 = \mathbb{Z}[\eta]/(2\eta, \eta^2)$} Whenever the fiber of $\delta\text{-}\mathbf{Ring} \to \mathbf{Ring}^\phi$ over $(A, \phi)$ is inhabited, it is a torsor under the $\mathbb{F}_p$-module $\mathrm{Der}_\phi(A, \mathrm{Ann}_A(p))$ of $\phi$-derivations into $p$-torsion. On the $2$-torsion witness $R_2 := \mathbb{Z}[\eta]/(2\eta, \eta^2)$---which decomposes as $\mathbb{Z} \oplus \mathbb{F}_2$ as an additive abelian group while remaining non-isomorphic as a ring to $\mathbb{Z} \times \mathbb{F}_2$ ($\eta \neq 0$, $\eta^2 = 0$)---fixing the \emph{augmentation Frobenius endomorphism} $\phi(a + b\eta) := a$ ($\phi(\eta) = 0$; never $\phi = \mathrm{id}_{R_2}$) yields \emph{precisely two} $\delta$-ring structures,
  \[
    \begin{aligned}
      \delta_0(a + b\eta) &\;=\; \frac{a - a^2}{2} + ab\,\eta &&(\delta_0(\eta) = 0), \\
      \delta_1(a + b\eta) &\;=\; \frac{a - a^2}{2} + (a+1)b\,\eta &&(\delta_1(\eta) = \eta),
    \end{aligned}
  \]
  and $(R_2, \delta_0) \not\cong (R_2, \delta_1)$ in $\delta\text{-}\mathbf{Ring}$.
  \item \textbf{Torsion-Free Non-Perfection in \texttt{DualInt} $= \mathbb{Z}[\varepsilon]/(\varepsilon^2)$:}\index{DualInt@$\texttt{DualInt} = \mathbb{Z}[\varepsilon]/(\varepsilon^2)$} On the torsion-free dual integers $\mathtt{DualInt} := \mathbb{Z}[\varepsilon]/(\varepsilon^2)$ (carrier $\mathbb{Z} \times \mathbb{Z}$), non-zero-divisibility of $2$ forces the augmentation Frobenius $\phi(a + b\varepsilon) := a$ to admit the \emph{unique} $2$-typical $\delta$-structure $\delta(a + b\varepsilon) = \frac{a - a^2}{2} - ab\,\varepsilon$, yet $\mathtt{DualInt}$ is strictly non-perfect ($\phi(\varepsilon) = 0$ prevents both injectivity and surjectivity of $\phi$).
\end{enumerate}
\end{proposition}

\begin{proof}
For~(i), given two $p$-typical $\delta$-structures $\delta_0, \delta_1$ lifting the same $\phi$ on $A$, set $D := \delta_1 - \delta_0 : A \to A$. Subtracting $x^p + p\,\delta_i(x) = \phi(x)$ gives $p\,D(x) = 0$, so $D$ lands in $\mathrm{Ann}_A(p)$. Subtracting Joyal's additive identities cancels $C_p(x, y)$ to give $D(x + y) = D(x) + D(y)$, while subtracting the multiplicative identities and substituting $p\,\delta_1(x) = \phi(x) - x^p$ yields
\[
  \begin{aligned}
    D(xy) &\;=\; x^p D(y) + y^p D(x) + p\,\delta_1(x) D(y) + p\,\delta_0(y) D(x) \\
    &\;=\; \phi(x) D(y) + y^p D(x) + (\phi(y) - y^p - p\,D(y)) D(x) \\
    &\;=\; \phi(x) D(y) + \phi(y) D(x),
  \end{aligned}
\]
since $p\,D(y) = 0$; conversely, adding any $D \in \mathrm{Der}_\phi(A, \mathrm{Ann}_A(p))$ with $D(1) = 0$ preserves Joyal's axioms and $\phi$. Specializing to $R_2 = \mathbb{Z}[\eta]/(2\eta, \eta^2)$ with $\phi(a + b\eta) = a$ ($a \in \mathbb{Z}$, $b \in \{0, 1\}$), $\delta(a) = \frac{a - a^2}{2}$ is forced on $\mathbb{Z}$ while $2\delta(\eta) = \phi(\eta) - \eta^2 = 0$ leaves $c := \delta(\eta) \in \mathrm{Ann}_{R_2}(2) = \{0, \eta\} \cong \mathbb{F}_2$ unconstrained, and $\delta(a + b\eta) = \delta(a) + \delta(b\eta) - ab\eta = \frac{a - a^2}{2} + (ab + c\cdot b)\eta$ (using $2ab\eta = 0$ in $R_2$, so $-ab\eta = +ab\eta$) checks both $\delta_0$ ($c = 0$) and $\delta_1$ ($c = \eta$). Any ring automorphism $\alpha : R_2 \xrightarrow{\sim} R_2$ fixes $1$ (hence $\mathbb{Z}$) and maps $\eta$ injectively into $\mathrm{Ann}_{R_2}(2) \setminus \{0\} = \{\eta\}$, forcing $\alpha = \mathrm{id}_{R_2}$; since $\delta_0(\eta) = 0 \neq \eta = \delta_1(\eta)$, no $\delta$-ring isomorphism can exist between $(R_2, \delta_0)$ and $(R_2, \delta_1)$.

For~(ii), in $\mathtt{DualInt} = \mathbb{Z}[\varepsilon]/(\varepsilon^2)$, multiplication by $2$ is injective on $\mathbb{Z} \times \mathbb{Z}$, so $\phi(a + b\varepsilon) = a$ uniquely forces $2\delta(a + b\varepsilon) = a - (a + b\varepsilon)^2 = (a - a^2) - 2ab\varepsilon$ and hence $\delta(a, b) = \bigl(\frac{a - a^2}{2}, -ab\bigr)$ (\texttt{DualInt.dualDelta}, \texttt{DualInt.instDeltaRing2} in \texttt{HomotopyCRM/\allowbreak DeltaRing\allowbreak Perfection.lean}, 2,113 lines, 26 kernel checks), while $\phi(0, 1) = (0, 0)$ witnesses both non-injectivity and non-surjectivity (\texttt{DualInt.not\_isPerfectDeltaRing}, \texttt{Checks.dualInt\_frobeniusLift\_eps}).
\end{proof}

\begin{proposition}[Three-Stage Mode~(iii) Obstruction Tower Over $\mathbf{FrobLiftRing}_2$]\label{prop:delta-three-stage-tower}%
\lean{HomotopyCRM.r2_aut_unique, HomotopyCRM.deltaRing2_sub_frobDerivation2, HomotopyCRM.r2Delta0_not_iso_r2Delta1, HomotopyCRM.delta_mul_p_eq_frob_sub, HomotopyCRM.frob_annihilates_p_torsion, HomotopyCRM.empty_deltaLift_R3, HomotopyCRM.isEmpty_deltaRing2_zmod4, HomotopyCRM.empty_deltaLift_R220, HomotopyCRM.empty_deltaLift_R2Idemp, HomotopyCRM.r2_lift_card_eq_two, HomotopyCRM.r22_lift_card_eq_sixteen, HomotopyCRM.r22_iso_orbits_eq_six}%
\leanok%
\uses{def:comparison-triple,def:residuum,prop:tri-modal-factorization,prop:r2-delta-torsor}%
Along the vertical forgetful functor $U : \delta\text{-}\mathbf{Ring}_2 \to \mathbf{FrobLiftRing}_2$ (Mode~(iii) of Proposition~\ref{prop:tri-modal-factorization}), the homotopy fiber $\mathrm{hofib}_{(R,\phi)}(U)$ over a base ring $R$ with candidate Frobenius lift $\phi : R \to R$ stratifies into a three-stage obstruction tower:
\begin{enumerate}[label=\textup{(\arabic*)},leftmargin=2.0em]
  \item \textbf{Stage~1 (Primary $p$-Torsion Annihilation):} Every $p$-derivation $\delta$ lifting $\phi(x) = x^p + p\,\delta(x)$ satisfies $p\,\delta(p) = \phi(p) - p^p$ (\texttt{delta\_\allowbreak mul\_\allowbreak p\_\allowbreak eq\_\allowbreak frob\_\allowbreak sub}). Whenever $\delta(0) = 0$, Joyal's product law holds, and $\phi(p) = p$, any $p$-torsion element $a \in R[p]$ must satisfy $p\,\phi(a) = 0$ (\texttt{frob\_\allowbreak annihilates\_\allowbreak p\_\allowbreak torsion}). Consequently, on $R_2$ the Fermat-congruent map $\phi_3(a + b\eta) := a + 2(b \bmod 2)$ has empty lift fiber ($2\,\phi_3(\eta) = 4 \neq 0$, \texttt{empty\_\allowbreak deltaLift\_\allowbreak R3}), and $\mathbb{Z}/4\mathbb{Z}$ admits zero $2$-typical $\delta$-structures (\texttt{isEmpty\_\allowbreak deltaRing2\_\allowbreak zmod4}).
  \item \textbf{Stage~2 (Secondary $2$-Cocycle Obstruction):} Passing Stage~1 ($2\,\phi(\eta) = 0$) does not guarantee inhabitation: setting $\phi = \mathrm{id}$ yields empty lift fibers both on $R_2 = \mathbb{Z}[\eta]/(2\eta, \eta^2)$ (\texttt{empty\_\allowbreak deltaLift\_\allowbreak R220}) and on $R_{2,\mathrm{idemp}} := \mathbb{Z}[\eta]/(2\eta, \eta^2 - \eta) \cong \mathbb{Z} \times \mathbb{F}_2$ (where $\mathrm{id}$ even satisfies $x - x^2 \in 2R$ pointwise, yet $\delta(2\eta) = 0$ fails; \texttt{empty\_\allowbreak deltaLift\_\allowbreak R2Idemp}).
  \item \textbf{Stage~3 (Tertiary $\mathrm{Der}_\phi$-Torsor and $\mathrm{GL}_2(\mathbb{F}_2)$ Monodromy):} Over $(R, \phi)$ admitting a base lift, the fiber $\mathrm{hofib}_{(R,\phi)}(U)$ is a $\mathrm{Der}_\phi(R, R[2])$-torsor (\texttt{deltaRing2\_\allowbreak sub\_\allowbreak frobDerivation2}): it is a singleton on $\mathbb{Z}$, a $2$-element set of non-isomorphic lifts on $(R_2, \phi_{\mathrm{aug}})$ ($\mathrm{Aut}(R_2) = \{\mathrm{id}\}$, \texttt{r2\_\allowbreak aut\_\allowbreak unique}, \texttt{r2\_\allowbreak lift\_\allowbreak card\_\allowbreak eq\_\allowbreak two}, \texttt{r2Delta0\_\allowbreak not\_\allowbreak iso\_\allowbreak r2Delta1}), and on $R_{2,2} := \mathbb{Z}[\eta_1, \eta_2]/(2\eta_i, \eta_i\eta_j)$ with augmentation Frobenius $\phi_{\mathrm{aug}}(a + b_1\eta_1 + b_2\eta_2) := a$, $|\mathrm{hofib}_{(R_{2,2},\phi_{\mathrm{aug}})}(U)| = 16$ (\texttt{r22\_\allowbreak lift\_\allowbreak card\_\allowbreak eq\_\allowbreak sixteen}), carving under $\mathrm{Aut}(R_{2,2}, \phi_{\mathrm{aug}}) \cong \mathrm{GL}_2(\mathbb{F}_2)$ conjugation into precisely $6$ isomorphism orbits in $\pi_0\,\mathrm{hofib}_{(R_{2,2},\phi_{\mathrm{aug}})}(U)$ (\texttt{r22\_\allowbreak iso\_\allowbreak orbits\_\allowbreak eq\_\allowbreak six}).
\end{enumerate}
\end{proposition}

\begin{proof}
For~(1), $\phi(p) = p^p + p\,\delta(p)$ gives $p\,\delta(p) = \phi(p) - p^p$; if $pa = 0$ ($p \ge 1$), then $p a^p = 0$, and multiplying $0 = \delta(pa) = p^p\delta(a) + a^p\delta(p) + p\,\delta(p)\delta(a)$ by $p$ yields $0 = (p^p + p\,\delta(p))(p\,\delta(a)) = \phi(p)(p\,\delta(a)) = p^2\delta(a)$ (since $\phi(p) = p$), whence $p\,\phi(a) = p a^p + p^2\delta(a) = 0$. On $R_2$, $2\,\phi_3(\eta) = 4 \neq 0$ blocks any lift of $\phi_3$, while on $\mathbb{Z}/4\mathbb{Z}$ we have $\delta(2) = \delta(1+1) = 2\delta(1) - 1 = -1$ and $0 = \delta(2+2) = 2\delta(2) - 4 = -2 \neq 0$.

For~(2), on $(R_2, \mathrm{id})$, writing $\delta(\eta) = a + b\eta$ gives $\eta^2 + 2\delta(\eta) = 2a \neq \eta$; on $(R_{2,\mathrm{idemp}}, \mathrm{id})$ where $\eta^2 = \eta$ and $2\eta = 0$, adding $\eta + \eta = 0$ gives $0 = \delta(\eta + \eta) = 2\delta(\eta) - \eta^2 = 2a - \eta \neq 0$ (mod~$2$), contradicting existence of $\delta$.

For~(3), subtracting any two lifts lies in $\mathrm{Der}_\phi(R, R[2])$ by Proposition~\ref{prop:r2-delta-torsor}(i). Over $R_2$, $\mathrm{Aut}(R_2) = \{\mathrm{id}_{R_2}\}$ preserves $\delta_0$ and $\delta_1$ individually, giving $2$ orbits. Over $R_{2,2}$, any automorphism fixes $\mathbb{Z}$ pointwise and maps $\mathbb{F}_2\eta_1 \oplus \mathbb{F}_2\eta_2$ linearly, identifying $\mathrm{Aut}(R_{2,2}) \cong \mathrm{GL}_2(\mathbb{F}_2)$, while $\phi_{\mathrm{aug}}$-derivations into $R_{2,2}[2] = \mathbb{F}_2\eta_1 \oplus \mathbb{F}_2\eta_2$ correspond bijectively to $2 \times 2$ matrices $M \in \mathrm{Mat}_{2\times 2}(\mathbb{F}_2)$ via $M_{ij} = \delta(\eta_j)_i$ ($16$ lifts). Conjugation by $g \in \mathrm{GL}_2(\mathbb{F}_2)$ transforms $M \mapsto g M g^{-1}$, whose rational canonical forms over $\mathbb{F}_2$ fall into exactly $6$ conjugacy classes: $0$, $I_2$, nilpotent $E_{12}$, rank-$1$ idempotent $E_{11}$, order-$2$ transvection $\bigl(\begin{smallmatrix}1 & 1\\ 0 & 1\end{smallmatrix}\bigr)$, and order-$3$ elliptic $\bigl(\begin{smallmatrix}0 & 1\\ 1 & 1\end{smallmatrix}\bigr)$.
\end{proof}

\section{\texorpdfstring{Goldthorpe Theorems~A and~B, Shift Chirality, and Relative Slice Perfection}{Goldthorpe Theorems A and B, Shift Chirality, and Relative Slice Perfection}}\label{sec:slice-perfection}

Passing from non-perfect $\delta$-rings to their universal perfections exhibits an exact colimit initial-algebra duality with an intrinsic $+1$/$-1$ shift chirality:

\begin{theorem}[Undercategory Fixed-Point Initiality, Relative Perfection $A_{\mathrm{perf}}$, and $+1$/$-1$ Shift Chirality]\label{thm:slice-perfection-chirality}\index{Shift chirality ($+1$/$-1$)}%
\uses{def:comparison-triple,prop:r2-delta-torsor}%
Let $\mathcal{C}$ be a category, $F : \mathcal{C} \to \mathcal{C}$ an endofunctor, and $c = (V, \mathrm{str} : V \to F(V))$ an $F$-coalgebra.
\begin{enumerate}[label=\textup{(\roman*)},leftmargin=2.0em]
  \item \textbf{Undercategory Fixed-Point Initiality (Goldthorpe~\cite[Prop.~2.3.13]{Goldthorpe23}, Theorem~A):} On the coslice $\mathcal{C}_{V/}$ (\texttt{Under c.V}), define the slice endofunctor $F_c(X, \eta : V \to X) := (F(X), F(\eta) \circ \mathrm{str})$. Any initial $F_c$-algebra $(A_{\mathrm{init}}, \iota_0, \nu_\infty : F(A_{\mathrm{init}}) \to A_{\mathrm{init}})$ has an invertible structure morphism $\nu_\infty$ by Lambek's Lemma, and $(A_{\mathrm{init}}, \iota_0, \nu_\infty^{-1})$ is initial among fixed-point $F$-coalgebras under $c$.
  \item \textbf{Sequential Colimit Relative Perfection and $+1$/$-1$ Shift Chirality (Theorem~B):} Specializing to $F = \mathbf{1}_{\mathcal{C}}$ and a natural endomorphism $\varphi : \mathbf{1}_{\mathcal{C}} \Rightarrow \mathbf{1}_{\mathcal{C}}$ (such as $\varphi = \phi$ on $\delta\text{-}\mathbf{Ring}_p$) evaluated on $c = (A, \varphi_A)$, let $A_{\mathrm{perf}} := \varinjlim \bigl(A \xrightarrow{\varphi_A} A \xrightarrow{\varphi_A} \cdots\bigr)$ be a sequential colimit with cocone legs $\iota_n : A \to A_{\mathrm{perf}}$ ($\iota_{n+1} \circ \varphi_A = \iota_n$). Then $A_{\mathrm{perf}}$ carries a forward shift automorphism $s : A_{\mathrm{perf}} \xrightarrow{\sim} A_{\mathrm{perf}}$ satisfying $s \circ \iota_n = \iota_{n+1}$ such that:
  \begin{enumerate}[label=\textup{(\alph*)},leftmargin=1.6em]
    \item $(A_{\mathrm{perf}}, \iota_0, s)$ ($+1$ shift) is the initial $(\mathbf{1}_{\mathcal{C}})_c$-algebra in $\mathcal{C}_{A/}$; and
    \item naturality forces the induced endomorphism on $A_{\mathrm{perf}}$ to equal the \emph{inverse} shift, $\varphi_{A_{\mathrm{perf}}} = s^{-1}$ ($-1$ shift), making $\varphi_{A_{\mathrm{perf}}}$ an isomorphism. On $\mathtt{DualInt}$, idempotence $\phi(\phi(x)) = \phi(x)$ collapses $\mathtt{DualInt}_{\mathrm{perf}} \xrightarrow{\sim} \mathbb{Z}$ in one step.
  \end{enumerate}
\end{enumerate}
\end{theorem}

\begin{proof}
For~(i), applying Lambek's Lemma (\texttt{lambek\_isIso}) to $F_c$ on $\mathcal{C}_{V/}$ shows $\nu_\infty : F_c(A_{\mathrm{init}}, \iota_0) \xrightarrow{\sim} (A_{\mathrm{init}}, \iota_0)$ is an isomorphism in $\mathcal{C}_{V/}$ (\texttt{sliceLambek\_right\_isIso}); inverting the structure map bijects $F_c$-algebra homomorphisms out of $(A_{\mathrm{init}}, \iota_0, \nu_\infty)$ with coalgebra morphisms out of $(A_{\mathrm{init}}, \iota_0, \nu_\infty^{-1})$ into fixed-point $F$-coalgebras under $c$ (\texttt{theoremA\_homEquiv}, \texttt{theoremA\_isInitial\_fixedPointCoalgebraUnder}).

For~(ii), the shifted cocone legs $n \mapsto \iota_{n+1}$ satisfy $\iota_{n+2} \circ \varphi_A = \iota_{n+1}$ and induce $s : A_{\mathrm{perf}} \to A_{\mathrm{perf}}$ with $s \circ \iota_n = \iota_{n+1}$, while the backward cocone legs $n \mapsto \iota_n \circ \varphi_A$ ($= \iota_{n-1}$ for $n \ge 1$) induce $t : A_{\mathrm{perf}} \to A_{\mathrm{perf}}$ satisfying $t \circ s \circ \iota_n = t \circ \iota_{n+1} = \iota_{n+1} \circ \varphi_A = \iota_n$ and $s \circ t \circ \iota_n = s \circ \iota_n \circ \varphi_A = \iota_{n+1} \circ \varphi_A = \iota_n$, so $t = s^{-1}$ (\texttt{endomorphismColimitShiftIso}). Given any slice algebra $(B, f : A \to B, \psi : B \to B)$ satisfying $\psi \circ f = f \circ \varphi_A$, induction on $n$ shows $\psi^n \circ f = \psi^{n+1} \circ f \circ \varphi_A$ (\texttt{comp\_endomorphismPow\_eq\_iterate}), inducing a unique colimit morphism $h : A_{\mathrm{perf}} \to B$ with $h \circ \iota_n = \psi^n \circ f$ and $h \circ s \circ \iota_n = h \circ \iota_{n+1} = \psi^{n+1} \circ f = \psi \circ h \circ \iota_n$, proving~(a) (\texttt{theoremB\_perfectionSliceAlg\_isInitial}). Finally, naturality of $\varphi$ along $\iota_{n+1} : A \to A_{\mathrm{perf}}$ yields $\varphi_{A_{\mathrm{perf}}} \circ \iota_{n+1} = \iota_{n+1} \circ \varphi_A = \iota_n = s^{-1} \circ \iota_{n+1}$ for all $n \ge 0$; since $\{\iota_{n+1}\}_{n \ge 0}$ is jointly epimorphic ($\iota_0 = \iota_1 \circ \varphi_A$), $\varphi_{A_{\mathrm{perf}}} = s^{-1}$ (\texttt{theoremB\_natEndo\_app\_eq\_shift\_inv}, \texttt{theoremB\_isIso\_natEndo\_app}). On $\mathtt{DualInt}$, $\phi(a, b) = (a, 0)$ satisfies $\phi \circ \phi = \phi$ (\texttt{Checks.dualInt\_frobeniusLift\_idem}), identifying the colimit with $(\mathbb{Z}, \mathtt{DualInt.fstDeltaHom})$ (\texttt{DualInt.perfectionDeltaCocone\_isColimit}, \texttt{DualInt.perfectionSliceAlgDualInt\_isInitial}).
\end{proof}

\begin{remark}[Mathematical Red-Team Guardrail~2: Relative Perfection vs.\ General Prisms]\label{rem:guardrail-2}%
\lean{HomotopyCRM.lambek_isIso, HomotopyCRM.sliceLambek_right_isIso, HomotopyCRM.theoremA_homEquiv, HomotopyCRM.theoremA_isInitial_fixedPointCoalgebraUnder, HomotopyCRM.endomorphismColimitShiftIso, HomotopyCRM.theoremB_natEndo_app_eq_shift_inv, HomotopyCRM.theoremB_isIso_natEndo_app, HomotopyCRM.theoremB_perfectionSliceAlg_isInitial, HomotopyCRM.instDeltaRing2Int, HomotopyCRM.isPerfectDeltaRing_int, HomotopyCRM.instUniqueDeltaRing2Int, HomotopyCRM.DualInt.instDeltaRing2, HomotopyCRM.DualInt.not_isPerfectDeltaRing, HomotopyCRM.DualInt.perfectionDeltaCocone_isColimit, HomotopyCRM.DualInt.perfectionSliceAlgDualInt_isInitial}%
\leanok%
\uses{def:comparison-triple,prop:r2-delta-torsor,thm:slice-perfection-chirality}%
The sequential colimit construction $A_{\mathrm{perf}} = \varinjlim_\varphi A$ of Theorem~\ref{thm:slice-perfection-chirality}(ii) models \textbf{strictly} the \emph{relative slice perfection}\index{Relative slice perfection@Relative slice perfection $A_{\mathrm{perf}}$} (the left-adjoint perfect reflector $A \to A_{\mathrm{perf}}$ on which $\varphi$ becomes invertible). A Lambek initial-algebra model \textbf{must never} be imposed on general Bhatt--Scholze prisms $(A, I)$ or on the prismatic site $(R)_{\Prism}$~\cite[\S\S 2--3, \S 12]{BhattScholze22}. In a general prism, non-invertibility of $\phi$ is essential to define the \textbf{Nygaard filtration}\index{Nygaard filtration} $\mathcal{N}^{\ge i} A := \phi^{-1}(I^i)$ (\texttt{nygaardLevel} and \texttt{nygaard\_antitone} in \texttt{DeltaRingPerfection.lean}); forcing $\phi$ to become an isomorphism collapses $\mathcal{N}^{\ge i} A$ to $I^i$. Indeed, Layer~3 of \texttt{DeltaRingPerfection.lean} (\texttt{Checks.dualInt\_nygaard\_separation}) proves that for $I = (\varepsilon)$ in \texttt{DualInt}, $\varepsilon \in \mathcal{N}^{\ge 1}\mathrm{DualInt}$ whereas $1 \notin \mathcal{N}^{\ge 1}\mathrm{DualInt}$---a non-trivial filtration completely obliterated upon passing to $\mathrm{DualInt}_{\mathrm{perf}} \cong \mathbb{Z}$ (where $\phi = \mathrm{id}_{\mathbb{Z}}$).
\end{remark}

\chapter{Base-Topos Variation, Geometric Morphisms, and the Base-Fiber Conservation Principle}
\chaptermark{Base-Topos Variation and Base-Fiber Conservation}
\label{ch:base-fiber}
\chapterepigraph{``We describe a basic construction which extends the familiar duality between ordinary spaces and commutative algebras to a duality between quotient spaces and noncommutative algebras.''}{Alain Connes, \emph{Noncommutative Geometry---Year 2000} (2000)}

\section{The Two-Axis Relativization of Categorical Reverse Mathematics}\label{sec:two-axis-relativization}

Having calibrated the $1$-categorical and symmetric-forcing Choice atlas in Chapter~\ref{ch:choice-atlas} and characteristic-$p$ Joyal $\delta$-rings with relative slice perfection in Chapter~\ref{ch:delta-rings}, we now examine the multi-axis dynamism of Categorical Reverse Mathematics. Classical Friedman--Simpson reverse mathematics over $\mathrm{RCA}_0$ and Ishihara's constructive reverse mathematics over $\mathrm{BISH}$ both fix a single static metamathematical background. In $\mathrm{CRM}$, by contrast, every deductive reversal $A + \mathcal{R} \iff B$ is governed by a doctrinal comparison triple
\[
  \mathfrak{T} = \bigl(\mathcal{D}_0 \supseteq \mathcal{D}_1,\; \kappa : \Phi \Longrightarrow \Psi\bigr) \in \mathrm{Dom}(\mathrm{Res})
\]
in the sense of Chapter~\ref{ch:foundations} (Definitions~\ref{def:comparison-triple}--\ref{def:residuum}, Sections~\ref{sec:finset-infinity}--\ref{sec:base-fiber}, and Table~\ref{tab:five-axis-master}), where $\mathrm{Res}(\kappa)$ measures the universal obstruction to the canonical comparison $\kappa$ restricting to an equivalence outside the stronger subdoctrine $\mathcal{D}_1$. Because $\mathrm{Dom}(\mathrm{Res})$ ranges over structured categorical doctrines equipped with comparisons across ambient and base universes, we can deform any residuum along two orthogonal axes:

\begin{itemize}[leftmargin=1.5em]

  \item \textbf{Axis~1 (Structural Deformation of the Ambient Universe $\mathcal{E}$):} We weaken or strengthen the Lawvere--Tierney categorical axioms of $\mathcal{E}$ itself---dropping impredicativity ($\Omega$), local Cartesian closure ($Y^X$ and $\Pi_f$), arithmetic infinity ($\mathrm{NNO}$), or $0$-truncation (into $(\infty,1)$-topoi)---to observe how $\mathcal{R}(\mathcal{E})$ fractures across doctrinal boundaries.

  \item \textbf{Axis~2 (Metamathematical Variation of the Base Universe $\mathcal{S}$):} Given an adjunction $F \dashv U : \mathcal{D}_{\mathcal{E}} \rightleftarrows \mathcal{D}_{\mathcal{S}}$---prototypically a bounded geometric morphism $\gamma : \mathcal{E} \to \mathcal{S}$ ($\gamma^* \dashv \gamma_*$, classically $\mathcal{S} = \mathbf{Set}$ with $\mathrm{AC}$ and $\mathrm{LEM}$), or a ring base change $R \to S$---we vary the logical and choice axioms assumed in $\mathcal{S}$ and measure the compensatory shift in the relative fiber residuum $\mathcal{R}(\mathcal{E}/\mathcal{S}) = \mathrm{Res}(\kappa_{\mathcal{E}/\mathcal{S}})$.

\end{itemize}

Together, these two axes transform Reverse Mathematics from a static taxonomy of formal subsystems into a quantitative theory of \emph{logical-geometric covariance} across mathematical universes.

\section{\texorpdfstring{Axis 1 Deformations: Weakening the Lawvere--Tierney Axioms of $\mathcal{E}$}{Axis 1 Deformations: Weakening the Lawvere--Tierney Axioms of E}}\label{sec:axis1-deformations}

An elementary topos $\mathcal{E}$ with a Natural Numbers Object ($\mathrm{NNO}$) is governed by four structural pillars: finite limits and colimits, exponential function objects $Y^X$ (Cartesian closure in every slice $\mathcal{E}/X$), an impredicative subobject classifier $\Omega$, and an infinitary $\mathrm{NNO}$ together with $0$-truncation of hom-sets. Removing each pillar systematically exposes which components of $\mathcal{R}(\mathcal{E})$ originate from impredicativity, higher-order function spaces, arithmetic infinity, or strict set-truncation.

\subsection{\texorpdfstring{Dropping the Subobject Classifier $\Omega$ (Elementary Topos $\to$ Pretopos and Arithmetic Universe)}{Dropping the Subobject Classifier Omega (Elementary Topos -> Pretopos and Arithmetic Universe)}}

In an elementary topos $\mathcal{E}$, the subobject classifier $1 \xrightarrow{\top} \Omega$ and exponential objects $Y^X$ combine to form impredicative power objects $\mathcal{P}(X) = \Omega^X$. Dropping $\Omega$ places us in the predicative doctrine of locally Cartesian closed pretopoi ($\Pi$-pretopoi) and \emph{Maietti--Joyal Arithmetic Universes}~\cite{Maietti05}.

Without $\Omega$, any residuum $\mathcal{R}(\mathcal{E})$ previously formulated as an internal order-algebraic property of $\Omega$---such as Booleanness ($\Omega \simeq 1+1$) or De Morgan's Law ($\Omega$ is a Stone lattice)---must be \emph{externalized} into a uniform schematic condition across the fibration of subobject lattices $\mathrm{Sub}_{\mathcal{E}}(X)$ for all $X \in \mathcal{E}$. Note carefully which functor encodes internal choice: in any locally Cartesian closed pretopos, pullback $f^* : \mathcal{E}/Y \to \mathcal{E}/X$ already possesses a right adjoint $\Pi_f$ and therefore preserves colimits and effective epimorphisms unconditionally; internal choice instead constrains the right adjoint $\Pi_f$~\cite{Diaconescu75,Johnstone02}.

\begin{proposition}[Predicative Stability of Diaconescu's Residuum Factor]\label{prop:predicative-diaconescu}%
\uses{thm:diaconescu}%

In any locally Cartesian closed pretopos $\mathcal{E}$ lacking a subobject classifier $\Omega$, the Internal Axiom of Choice ($\mathrm{IAC}$, asserting that the dependent product functors $\Pi_f : \mathcal{E}/X \to \mathcal{E}/Y$, right adjoint to pullback $f^*$, preserve effective epimorphisms) implies that every subobject lattice $\mathrm{Sub}_{\mathcal{E}}(X)$ is a complemented Boolean algebra.

\end{proposition}

\textbf{Significance for $\mathrm{CRM}$.} Proposition~\ref{prop:predicative-diaconescu} proves that the obligatory Booleanness factor in $(A \Rightarrow \mathrm{AC})$ is \emph{predicative}. The collapse of intuitionistic logic under choice does not depend on impredicative comprehension $\Omega^X$, as Diaconescu's pushout construction uses only finite colimits ($X +_U X$), pullbacks, and dependent products ($\Pi$-types). Conversely, as established in Chapter~\ref{ch:panorama} (\S\ref{sec:predicativity-wtypes}), impredicative comprehension $\Omega$ is isolated as the exact Axis~1 residuum bridging predicative tree recursion ($\mathrm{ATR}_0$ in $\Pi W$-pretopoi) to $\Pi^1_1$-comprehension ($\Pi^1_1\text{-}\mathrm{CA}_0$ in elementary topoi)~\cite{Maietti05}.

\subsection{\texorpdfstring{Dropping Exponentials $Y^X$ (Topos $\to$ Coherent and Regular Categories)}{Dropping Exponentials YX (Topos -> Coherent and Regular Categories)}}

Dropping local Cartesian closure ($Y^X$ and dependent products $\Pi_f$) descends from pretopoi to \emph{coherent} and \emph{regular categories}~\cite{Johnstone02}. Two structural deformations occur in this regime:

\begin{itemize}[leftmargin=1.5em]

  \item \textbf{Loss of Internal Choice ($\mathrm{IAC}$):} Because $\mathrm{IAC}$ is defined by preservation of epimorphisms under $\Pi_f$ (or exponentiation $(-)^Z : X^Z \to Y^Z$), $\mathrm{IAC}$ cannot even be formulated internally without function objects. Consequently, bridging local splitting of covers to internal choice in a coherent category $\mathcal{E}$ forces \emph{local Cartesian closure} directly into the residuum $\mathcal{R}(\mathcal{E})$.

  \item \textbf{Bifurcation of Epimorphisms:} In an elementary topos, every epimorphism is effective (the coequalizer of its kernel pair). In a general regular or coherent category $\mathcal{A}$ that is not balanced, cancellation epimorphisms are strictly weaker than \emph{regular (effective) epimorphisms}. External choice therefore splits into $\mathrm{AC}_{\mathrm{reg}}$ (every regular epimorphism splits) versus full $\mathrm{AC}_{\mathrm{epi}}$, with $(\mathrm{AC}_{\mathrm{reg}} \Rightarrow_{\mathbf{RegCat}} \mathrm{AC}_{\mathrm{epi}})$ isolating \emph{balancedness} (or Barr-exactness/pretopos structure).

\end{itemize}

\subsection{\texorpdfstring{Dropping the Natural Numbers Object $\mathrm{NNO}$ ($\mathbf{FinSet}$ vs.\ Base--Fiber Arithmetic)}{Dropping the Natural Numbers Object NNO (FinSet vs. Base-Fiber Arithmetic)}}

Removing the Natural Numbers Object ($\mathrm{NNO}$) moves $\mathrm{CRM}$ from infinitary universes to \emph{finitary topoi}, such as $\mathbf{FinSet}$ or finite group-action topoi $\mathbf{FinSet}^G$ for a finite group $G$. As calibrated in Chapter~\ref{ch:foundations} (Section~\ref{sec:finset-infinity}), $\mathbf{FinSet} \in \mathbf{ElemTop}$ satisfies both $\mathrm{LEM}$ and full external $\mathrm{AC}$ everywhere, so \emph{every choice residuum vanishes identically}, yet the failure of $\mathrm{NNO}$ ($\neg\mathsf{Inf}$) isolates three pure Axis~III infinity residua inside $\mathrm{Pro}(\mathbf{FinSet})$: the profinite lasso initial-algebra phantom $\widehat{\mathbb{N}}_{\mathrm{lasso}} = \varprojlim_{(p,q)} L_{p,q} \notin \mathbf{FinSet}$ (whose periodic $p=0$ quotient is $\widehat{\mathbb{Z}}$), the Dedekind--Pigeonhole gap $\mathrm{Mono}(X,X)\setminus\mathrm{Iso}(X,X)$, and the biproduct-tail comparison $\bigoplus_{i \in I} A \to \prod_{i \in I} A$ (isomorphic for all $I \in \mathbf{FinSet}$, but precipitating $\prod_{\mathbb{N}} A / \bigoplus_{\mathbb{N}} A \neq 0$ once $\mathrm{NNO}$ enters).

Adding non-trivial symmetries inside a finitary base produces orthogonal mixtures: in $\mathbf{FinSet}^G$ for $G \neq \{1\}$, the topos remains Boolean ($\Omega \simeq 1+1$) and Kuratowski-finite, yet external $\mathrm{AC}$ fails because $1$ is not projective (nontrivial $G$-sets without fixed points admit no global sections). Conversely, in Schanuel's nominal topos $\mathbf{Nom}$, an $\mathrm{NNO}$ exists ($\Delta(\mathbb{N})$) so Axis~III infinity holds, yet on the Kuratowski-infinite atom object $\mathbb{A}$ finite support forces $\mathrm{Mono}_{\mathbf{Nom}}(\mathbb{A},\mathbb{A}) = \mathrm{Iso}_{\mathbf{Nom}}(\mathbb{A},\mathbb{A})$, cleanly crossing Axis~I ($\neg\mathrm{AC}_\omega$) with Axis~III ($\mathrm{Inf}$).

\subsection{\texorpdfstring{Dropping $0$-Truncation ($1$-Topos $\to$ $(\infty, 1)$-Topos $\mathcal{X}$)}{Dropping 0-Truncation (1-Topos -> (infinity, 1)-Topos X)}}\label{sec:axis1-truncation}

Finally, dropping $0$-truncation of mapping spaces elevates $\mathrm{CRM}$ from $1$-topoi to an $(\infty, 1)$-topos $\mathcal{X}$ equipped with homotopy pullbacks and an univalent object classifier $\mathcal{U}$~\cite{Lurie09,GepnerKock17}. As shown in Chapter~\ref{ch:higher-crm} (\S\S\ref{sec:higher-choice}--\ref{sec:postnikov-residuum}; Definition~\ref{def:bigraded-choice}, Theorem~\ref{thm:global-inf-choice}, Theorem~\ref{thm:postnikov-stepwise-proved}, and Conjecture~\ref{conj:postnikov-stepwise}), removing $0$-truncation fractures any single $1$-categorical choice residuum into a descending deductive chain of truncated choice principles
\[
  \mathrm{AC}_\infty \implies \cdots \implies \mathrm{AC}_n \implies \cdots \implies \mathrm{AC}_1 \implies \mathrm{AC}_0,
\]
with model classes satisfying $\mathbf{Mod}(\mathrm{AC}_\infty) \subseteq \bigcap_{k < \infty} \mathbf{Mod}(\mathrm{AC}_k)$ and $\mathbf{Mod}(\mathrm{AC}_{n+1}) \subsetneq \mathbf{Mod}(\mathrm{AC}_n)$ for every $n \ge 0$ (Chapter~\ref{ch:higher-crm}, \S\ref{sec:higher-choice}). The residuum bridging set-level choice $\mathrm{AC}_0$ to $\mathrm{AC}_\infty$ is precisely the geometric condition that $0$-truncated objects generate $\mathcal{X}$ under effective epimorphisms, while each adjacent step $(\mathrm{AC}_n \Rightarrow \mathrm{AC}_{n+1})$ is governed by the vanishing of Giraud non-abelian gerbe cohomology $H^2_{\mathrm{gerbe}}$~\cite[Ch.~IV]{Giraud71} ($n=0$) and degree-$(n+2)$ Eilenberg--MacLane obstructions $H^{n+2}$~\cite[\S 7.2.2]{Lurie09} ($n \ge 1$).

\section{Axis 2 Relativization and Deforming Classical Residuums Without Base Choice or Ultrafilters}

We now turn to Axis~2: varying the metamathematical axioms of the base topos $\mathcal{S}$ over which a Grothendieck or relative topos $\mathcal{E}$ is constructed~\cite{Johnstone02,MacLaneMoerdijk92}.

\subsection{Bounded Geometric Morphisms and Absolute vs.\ Relative Residuums}

Let $\mathcal{S}$ be an elementary base topos, and let $\gamma: \mathcal{E} \to \mathcal{S}$ be a bounded geometric morphism\index{Geometric morphism!bounded} consisting of a direct image functor $\gamma_*: \mathcal{E} \to \mathcal{S}$ and a left exact inverse image functor $\gamma^*: \mathcal{S} \to \mathcal{E}$ with $\gamma^* \dashv \gamma_*$, so that $\mathcal{E} \simeq \mathbf{Sh}_{\mathcal{S}}(\mathbb{C}, J)$ for an internal site $(\mathbb{C}, J)$ in $\mathcal{S}$. Because $\gamma^*$ preserves the terminal object ($\gamma^*(1_{\mathcal{S}}) \simeq 1_{\mathcal{E}}$), the relative global sections functor $\Gamma_{\mathcal{E}/\mathcal{S}}: \mathcal{E} \to \mathcal{S}$ is canonically isomorphic to the direct image functor $\gamma_* = \mathrm{Hom}_{\mathcal{E}}(1_{\mathcal{E}}, -)$.

\begin{definition}[Absolute vs.\ Relative Categorical Residuums]\label{def:abs-rel-residuum}

Let $A$ and $B$ be structural or internal principles evaluated across relative topoi $\gamma: \mathcal{E} \to \mathcal{S}$.

\begin{itemize}[leftmargin=1.5em]

  \item A categorical residuum $\mathcal{R} = (A \Rightarrow B)$ is \emph{absolute (base-invariant)}\index{Categorical residuum!absolute vs.\ relative} if the equivalence $\mathcal{E} \models A \wedge \mathcal{R}(\mathcal{E}) \iff \mathcal{E} \models B$ holds internally for every topos $\mathcal{E}$ over $\mathcal{S}$ regardless of which logical or choice axioms hold in the base topos $\mathcal{S}$.

  \item A categorical residuum $\mathcal{R}(\mathcal{E}/\mathcal{S}) = (A \Rightarrow_{\mathbf{Topos}/\mathcal{S}} B)$ is \emph{relative (base-sensitive)} if its characterization depends on the metamathematical axioms (such as $\mathrm{AC}$, $\mathrm{DC}$, $\mathrm{BPI}$, or $\mathrm{LEM}$) assumed in the base topos $\mathcal{S}$.

\end{itemize}

\end{definition}

\textbf{Why Internal Residuums are Absolute and External Residuums are Base-Sensitive.} Purely internal implications formulated inside the Mitchell--Bénabou language of $\mathcal{E}$ using finite limits, colimits, and exponentials are \emph{absolute}. Thus Diaconescu's theorem ($\mathcal{E} \models \mathrm{IAC} \implies \Omega_{\mathcal{E}} \simeq 1+1$) holds over \emph{any} base topos $\mathcal{S}$, even if $\mathcal{S}$ is intuitionistic and non-Boolean; indeed, $\mathbf{Sh}_{\neg\neg}(\mathcal{S}) \hookrightarrow \mathcal{S}$ is always Boolean. Likewise, internalizing Theorem~\ref{thm:coshep-svc} by calling $Q \in \mathcal{E}$ \emph{internally projective} when $(-)^Q : \mathcal{E} \to \mathcal{E}$ preserves epimorphisms---yielding internal enough projectives ($\mathrm{ICOSHEP}$\index{Axiom of Choice!ICOSHEP@$\mathrm{ICOSHEP}$ (Internal Presentation Axiom)}) and internally retract-bounded Small Violations of Choice ($\mathrm{ISVC}_{\mathrm{ret}}$\index{Axiom of Choice!ISVC_ret@$\mathrm{ISVC}_{\mathrm{ret}}$ (Internal retract-bounded $\mathrm{SVC}$)}, positing an internally projective cover $P \twoheadrightarrow S$ such that every $X \in \mathcal{E}$ is a retract $r : U \twoheadrightarrow X$ of $U \rightarrowtail C \times P$ for an internally projective $C$)---gives the base-invariant equivalence $\mathrm{ICOSHEP} + (\mathrm{ISVC}_{\mathrm{ret}} \wedge \Omega_{\mathcal{E}} \simeq 1+1) \iff \mathrm{IAC}$ over \emph{any} base $\mathcal{S}$. Indeed, $(-)^{C \times P} \cong ((-)^P)^C$ preserves epimorphisms; if $U \rightarrowtail Q$ in a Boolean topos ($Q \cong U \amalg \neg U$) and $e : A \twoheadrightarrow B$ is epic, setting $e_1 := e \amalg \mathrm{id}_1 : A \amalg 1 \twoheadrightarrow B \amalg 1$ makes $e_1^Q \cong e_1^U \times e_1^{\neg U}$ epic and $\mathrm{pr}_1 : (B \amalg 1)^U \times (B \amalg 1)^{\neg U} \twoheadrightarrow (B \amalg 1)^U$ split epic via $1 \xrightarrow{\mathrm{in}_2} B \amalg 1$, whence $e_1^U$ is epic and its pullback along $\mathrm{in}_1^U : B^U \rightarrowtail (B \amalg 1)^U$ recovers $e^U : A^U \twoheadrightarrow B^U$ as epic by extensivity; and if $r \circ i = \mathrm{id}_X$, then $B^i \circ B^r = \mathrm{id}_{B^X}$ makes $B^i \circ e^U = e^X \circ A^i$ epic, forcing $e^X$ epic (whereas bare subquotient $\mathrm{SVC}$ still fails internally by the quotient fallacy~\cite{Blass79}). By contrast, any external statement $B$ asserting global sections, external projectivity, or spatial points---such as external $\mathrm{COSHEP} \Rightarrow \mathrm{AC}$ (Theorem~\ref{thm:coshep-svc}, which additionally requires external projectivity of $1_{\mathcal{E}}$, separated in $\mathbf{Set}^{\mathbb{Z}/2\mathbb{Z}}$ where $\mathrm{IAC}$ holds yet $\Gamma(\mathbb{Z}/2\mathbb{Z}) = \emptyset$)---invokes the direct image functor $\Gamma_{\mathcal{E}/\mathcal{S}} = \gamma_* : \mathcal{E} \to \mathcal{S}$; because the right adjoint $\gamma_*$ fails to preserve epimorphisms in general, external residuums couple $\mathcal{E}$ directly to the splitting and compactness axioms holding in $\mathcal{S}$.

\subsection{\texorpdfstring{Case 1: Deforming the Barr--Diaconescu--Freyd Residuum Without Base $\mathrm{AC}$}{Case 1: Deforming the Barr--Diaconescu--Freyd Residuum Without Base AC}}

Over a classical base $\mathcal{S} = \mathbf{Set}$ satisfying $\mathrm{AC}$, Theorem~\ref{thm:bdf} (the Barr--Diaconescu--Freyd theorem~\cite{Barr74,Diaconescu75,Freyd72}) states that a Grothendieck topos $\mathcal{E}$ satisfies external $\mathrm{AC}$ if and only if $\mathcal{E}$ is Boolean and localic over $\mathbf{Set}$ ($\mathcal{E} \simeq \mathbf{Sh}_{\mathbf{Set}}(B)$ for a complete Boolean algebra $B$):
\[
  \top + \bigl(\text{Booleanness } \wedge \text{ Localicity over } \mathbf{Set}_{\mathrm{AC}}\bigr) \iff \mathrm{AC}.
\]
When we drop $\mathrm{AC}$ in $\mathcal{S}$ ($\mathcal{S} \not\models \mathrm{AC}$), two fundamental constraints govern $\mathcal{E} \models \mathrm{AC}$:
\begin{enumerate}[label=\textup{(\arabic*)},leftmargin=2.0em]
  \item \textbf{Connected-Base Reflection Obstruction:} Suppose $\gamma : \mathcal{E} \to \mathcal{S}$ is \emph{connected}\index{Geometric morphism!connected} ($\gamma^*$ is fully faithful, so $\gamma_* \gamma^* \simeq \mathrm{id}_{\mathcal{S}}$) and $\mathcal{E} \models \mathrm{AC}$. Given any epimorphism $e : X \twoheadrightarrow Y$ in $\mathcal{S}$, the pullback $\gamma^* e$ is an epimorphism in $\mathcal{E}$ (since the left adjoint $\gamma^*$ preserves all colimits and hence epimorphisms) and therefore admits a section $\sigma : \gamma^* Y \to \gamma^* X$ in $\mathcal{E}$. Fullness of $\gamma^*$ reflects $\sigma$ back to a section $s : Y \to X$ of $e$ in $\mathcal{S}$, forcing $\mathcal{S} \models \mathrm{AC}$. Consequently, over a base $\mathcal{S} \not\models \mathrm{AC}$, no connected $\mathcal{S}$-topos---not even $\mathrm{id} : \mathcal{S} \to \mathcal{S}$ on a Boolean Cohen--Halpern--Läuchli model---can satisfy $\mathrm{AC}$; $\gamma$ must factor through a non-trivial \emph{choice slice} $\mathcal{S}/U \models \mathrm{AC}$ (or an $\omega$-choice slice when $\mathcal{S} \models \mathrm{DC}$).
  \item \textbf{Atomic Permutation Collapse (Non-Projectivity of $1$):} Inside any connected two-valued atomic permutation topos $\mathbf{Cont}_{\mathcal{S}}(\mathcal{G})$ ($\mathrm{Sub}(1) = \{0,1\}$) for a topological group $\mathcal{G}$, the terminal object $1$ is \emph{not} externally projective whenever a proper open subgroup $U \subsetneq \mathcal{G}$ exists, because the quotient $\mathcal{G}/U \twoheadrightarrow 1$ is an epimorphism whose global sections vanish ($\Gamma(\mathcal{G}/U) = (\mathcal{G}/U)^{\mathcal{G}} = \emptyset$). Thus external projectivity of $1$, localicity, $\mathrm{COSHEP}$, $\mathrm{IAC}$, and $\mathrm{AC}$ in $\mathbf{Cont}_{\mathcal{S}}(\mathcal{G})$ all collapse simultaneously to the trivial filter of open subgroups $\mathcal{U}(\mathcal{G}) = \{\mathcal{G}\}$ ($\mathbf{Cont}_{\mathcal{S}}(\mathcal{G}) \simeq \mathcal{S}$)~\cite{Johnstone02,Kihara26}.
\end{enumerate}

To recover $\mathcal{E} \models \mathrm{AC}$ over $\mathcal{S} \not\models \mathrm{AC}$, $\gamma$ must therefore be non-connected over a choice slice and the fiber residuum $\mathcal{R}(\mathcal{E}/\mathcal{S})$ must absorb the missing base strength:

\begin{theorem}[Deformation of the Barr--Diaconescu--Freyd Residuum Over an Arbitrary Base $\mathcal{S}$]\label{thm:bdf-deformed}%
\uses{thm:bdf,prin:conservation,def:abs-rel-residuum}%

Let $\gamma: \mathcal{E} \to \mathcal{S}$ be a Grothendieck topos over a base topos $\mathcal{S}$ in which $\mathrm{AC}$ is not assumed.

\begin{itemize}[leftmargin=1.5em]

  \item \textbf{Arbitrary Base ($\mathcal{S} \models \top$):} $\mathcal{E}$ satisfies external $\mathrm{AC}$ if and only if $\gamma$ factors through a non-trivial choice slice $\mathcal{S}/U \models \mathrm{AC}$ over which $\mathcal{E}$ is Boolean and localic, and its underlying complete Boolean algebra $B$ is an \emph{internally projective atomic Boolean locale} over a projective generating family in $\mathcal{S}/U$ (so every cover refines to disjoint projective atoms whose global sections split in $\mathcal{S}/U$). In particular, if $\gamma$ is connected (for example $\mathbf{Cont}_{\mathcal{S}}(\mathcal{G})$ over $\mathcal{S} \not\models \mathrm{AC}$), $\mathcal{E} \models \mathrm{AC}$ is impossible, and in any two-valued atomic permutation topos $\mathbf{Cont}_{\mathcal{S}}(\mathcal{G})$ any proper open subgroup $U \subsetneq \mathcal{G}$ forces $\Gamma(\mathcal{G}/U) = \emptyset$ and breaks external projectivity of $1$, localicity, $\mathrm{COSHEP}$, $\mathrm{IAC}$, and $\mathrm{AC}$.

  \item \textbf{Sequential Base ($\mathcal{S} \models \mathrm{DC}$):} Over $\mathcal{S} \models \mathrm{ZF} + \mathrm{DC}$, a Boolean localic $\mathcal{S}$-topos $\mathcal{E} \simeq \mathbf{Sh}_{\mathcal{S}}(B)$ satisfies external $\mathrm{AC}$ if and only if $B$ is an internally projective countably-generated Boolean locale \emph{over an $\omega$-choice slice $\mathcal{S}/U$}, reducing both fiber coverings and slice epimorphisms to countable towers split by base $\mathrm{DC}$.

\end{itemize}

\end{theorem}

\subsection{\texorpdfstring{Case 2: Deforming the Deligne--$\mathrm{BPI}$ Residuum Without Base Ultrafilters}{Case 2: Deforming the Deligne--BPI Residuum Without Base Ultrafilters}}

Over $\mathcal{S} = \mathbf{Set}$, Deligne's theorem~\cite[Exp.~VI, Appendice]{Deligne72} (\cite[\S D3.3, Thm.~D3.3.13]{Johnstone02}) states that every coherent Grothendieck topos $\mathcal{E}$ has \emph{enough points} if and only if the Boolean Prime Ideal Theorem ($\mathrm{BPI}$ / Ultrafilter Principle) holds in $\mathcal{S}$. As shown in Chapter~\ref{ch:panorama} (\S\ref{sec:panorama-continuum-spatial}), over $\mathcal{S} \models \mathrm{BPI}$ the residuum elevating pointfree localic theorems (such as Johnstone's localic Tychonoff theorem~\cite{Johnstone81} or Mulvey--Pelletier's localic Hahn--Banach theorem~\cite{MulveyPelletier91}) to classical spatial theorems is \emph{finitary coherence} of the underlying locale.

When $\mathcal{S} \not\models \mathrm{BPI}$, finitary compactness no longer yields stalk points $p: \mathcal{S} \to \mathcal{E}$. To guarantee enough points constructively without base ultrafilters, the fiber residuum $\mathcal{R}(\mathcal{E}/\mathcal{S})$ must strengthen from finitary subcovers to \emph{single-element refinements}:

\begin{theorem}[The Coherence-to-Supercompactness Shift Without Base $\mathrm{BPI}$]\label{thm:deligne-deformed}%
\uses{prin:conservation,def:abs-rel-residuum}%

Let $\gamma: \mathcal{E} \to \mathcal{S}$ be a localic Grothendieck topos over a base topos $\mathcal{S}$ in which $\mathrm{BPI}$ fails. To bridge pointfree localic compactness theorems ($A$) to spatial point-set theorems ($B$), the topological residuum $\mathcal{R}(\mathcal{E}/\mathcal{S})$ must strengthen from \emph{finitary coherence} (every covering refines to a finite subcovering) to \emph{supercompactness}\index{Supercompactness} (every covering of a generating open refines to a single element, yielding completely prime filters and stalk points constructively without invoking $\mathrm{BPI}$ in $\mathcal{S}$~\cite[Ch.~II, \S\S 2--4]{Johnstone82}, \cite[\S D3.3]{Johnstone02}).

\end{theorem}

\section{The Base-Fiber Conservation Principle and the Quantitative Leray Sequence}

The deformations of Theorems~\ref{thm:bdf-deformed} and~\ref{thm:deligne-deformed} realize Principle~\ref{prin:conservation} from Chapter~\ref{ch:foundations} (Section~\ref{sec:base-fiber}) as an exact order-theoretic and cohomological conservation law\index{Base--Fiber Conservation Principle}. Crucially, because the inverse image functor $\gamma^* : \mathcal{S} \to \mathcal{E}$ is left adjoint, it preserves all colimits and epimorphisms unconditionally; non-preservation of epimorphisms and cohomological obstruction strength live instead in the direct image (pushforward) functor $\gamma_* : \mathcal{E} \to \mathcal{S}$ and its right derived functors $R^q\gamma_*$ ($q \ge 1$).

\begin{principle}[Order-Theoretic Base-Fiber Conservation Principle]\label{thm:conservation-law}

Let $\gamma: \mathcal{E} \to \mathcal{S}$ be a relative Grothendieck topos over a base topos $\mathcal{S}$, and let $A \implies B$ be an external deductive reversal evaluated in $\mathcal{E}$. The assignment mapping each metamathematical axiom system $\mathcal{M}(\mathcal{S})$ assumed in the base topos $\mathcal{S}$ to the relative fiber residuum $\mathcal{R}(\mathcal{E}/\mathcal{S}) = (A \Rightarrow_{\mathbf{Topos}/\mathcal{S}} B)$ defines an order-reversing correspondence: weakening the base axiom system from $\mathcal{M}_1(\mathcal{S})$ to $\mathcal{M}_2(\mathcal{S})$ ($\mathbf{Mod}(\mathcal{M}_1) \subsetneq \mathbf{Mod}(\mathcal{M}_2)$) induces a compensatory strengthening of the fiber residuum, transferring the missing metamathematical strength from failures of the pushforward functor $\gamma_*$ and its derived functors $R^q\gamma_*$ to force a strictly stronger topological, atomic, or cohomological rigidity condition:
\[
  \mathbf{Mod}(\mathcal{R}_2(\mathcal{E}/\mathcal{S})) \subsetneq \mathbf{Mod}(\mathcal{R}_1(\mathcal{E}/\mathcal{S})).
\]

\end{principle}

Passing from order theory to sheaf cohomology makes this conservation law completely quantitative: the total cohomological obstruction in $\mathcal{E}$ splits cleanly into base and fiber components coupled by a transgression differential.

\begin{theorem}[Quantitative Base--Fiber Conservation via the Leray Sequence]\label{thm:quantitative-conservation}%
\uses{prin:conservation,thm:conservation-law}%

Let $\gamma : \mathcal{E} \to \mathcal{S}$ be a bounded geometric morphism and $A \in \mathbf{Ab}(\mathcal{E})$ an internal abelian sheaf. The Grothendieck--Leray spectral sequence $E_2^{p,q} = H^p(\mathcal{S}; R^q\gamma_* A) \Longrightarrow H^{p+q}(\mathcal{E}; A)$ induces the exact five-term conservation sequence (where $\Gamma_{\mathcal{S}}(R^1\gamma_* A) = H^0(\mathcal{S}; R^1\gamma_* A)$):
\[
  \begin{aligned}
    0 \longrightarrow H^1_{\mathrm{base}}(\mathcal{S}; \gamma_* A) &\longrightarrow H^1(\mathcal{E}; A) \longrightarrow \Gamma_{\mathcal{S}}(R^1\gamma_* A) \\
    &\xrightarrow{\;d_2\;} H^2_{\mathrm{base}}(\mathcal{S}; \gamma_* A) \longrightarrow H^2(\mathcal{E}; A).
  \end{aligned}
\]
Consequently, for $A$-torsor splitting obstructions $\mathrm{Res}_{\mathrm{total}}(\mathcal{E}; A) := H^1(\mathcal{E}; A)$, base obstructions $\mathrm{Res}_{\mathrm{base}}(\mathcal{S}; \gamma_* A) := H^1_{\mathrm{base}}(\mathcal{S}; \gamma_* A)$, and fiber obstructions $\mathrm{Res}_{\mathrm{fiber}}(\mathcal{E}/\mathcal{S}; A) := \Gamma_{\mathcal{S}}(R^1\gamma_* A)$, exactness yields the partition
\[
  \begin{aligned}
    0 &\longrightarrow \mathrm{Res}_{\mathrm{base}}(\mathcal{S}; \gamma_* A) \longrightarrow \mathrm{Res}_{\mathrm{total}}(\mathcal{E}; A) \\
    &\longrightarrow \ker\bigl(d_2 : \mathrm{Res}_{\mathrm{fiber}}(\mathcal{E}/\mathcal{S}; A) \to H^2_{\mathrm{base}}(\mathcal{S}; \gamma_* A)\bigr) \longrightarrow 0.
  \end{aligned}
\]
Weakening the base topos $\mathcal{S}$ enlarges $H^1_{\mathrm{base}}(\mathcal{S}; \gamma_* A)$ and $H^2_{\mathrm{base}}(\mathcal{S}; \gamma_* A)$; exactness yields
\[
  \begin{aligned}
    H^1(\mathcal{E}; A) = 0 \;\iff\;& H^1_{\mathrm{base}}(\mathcal{S}; \gamma_* A) = 0 \\
    &\wedge\; \ker\bigl(d_2 : \Gamma_{\mathcal{S}}(R^1\gamma_* A) \to H^2_{\mathrm{base}}(\mathcal{S}; \gamma_* A)\bigr) = 0,
  \end{aligned}
\]
so $H^1(\mathcal{E}; A) = 0$ holds either when $\Gamma_{\mathcal{S}}(R^1\gamma_* A)$ vanishes (such as forcing $B$ in $\mathbf{Sh}_{\mathcal{S}}(B)$ to be projective atomic or supercompact) or when the transgression differential $d_2$ injects onto $H^2_{\mathrm{base}}(\mathcal{S}; \gamma_* A)$.

\end{theorem}

\begin{remark}[Transgressive Cancellation Along the $\mathrm{SL}_2(\mathbb{F}_5) \twoheadrightarrow A_5$ Schur Cover]\label{rem:leray-transgressive-cancellation}%
\lean{HomotopyCRM.commutator_sl2_f5_eq_top, HomotopyCRM.sl2_f5_leray_transgression_injective}%
\leanok%
\uses{thm:quantitative-conservation}%
Because the kernel of $d_2$ governs total $H^1(\mathcal{E}; A)$ vanishing in Theorem~\ref{thm:quantitative-conservation}, $H^1(\mathcal{E}; A) = 0$ does \emph{not} force $\Gamma_{\mathcal{S}}(R^1\gamma_* A) = 0$ in general: fiber obstructions can cancel transgressively into base~$H^2$. Consider the central group extension
\[
  1 \longrightarrow \{\pm I\} \longrightarrow \mathrm{SL}_2(\mathbb{F}_5) \xrightarrow{\;\gamma\;} A_5 \longrightarrow 1
\]
over classifying topoi $\gamma : \mathbf{B}\mathrm{SL}_2(\mathbb{F}_5) \to \mathbf{B}A_5$ with trivial coefficients $A = \mathbb{Z}/2\mathbb{Z}$. Since $\mathrm{SL}_2(\mathbb{F}_5)$ is perfect ($[\mathrm{SL}_2(\mathbb{F}_5), \mathrm{SL}_2(\mathbb{F}_5)] = \mathrm{SL}_2(\mathbb{F}_5)$, \texttt{commutator\_\allowbreak sl2\_\allowbreak f5\_\allowbreak eq\_\allowbreak top}), total first cohomology vanishes identically: $H^1(\mathrm{SL}_2(\mathbb{F}_5); \mathbb{Z}/2\mathbb{Z}) = \mathrm{Hom}(\mathrm{SL}_2(\mathbb{F}_5), \mathbb{Z}/2\mathbb{Z}) = 0$. Nevertheless, the invariant fiber torsor group $\Gamma(R^1\gamma_* \mathbb{Z}/2\mathbb{Z}) \cong \mathrm{Hom}(\{\pm I\}, \mathbb{Z}/2\mathbb{Z}) \cong \mathbb{Z}/2\mathbb{Z}$ is non-zero, and the five-term exact sequence closes precisely because the Schur transgression differential $d_2 : \mathrm{Hom}(\{\pm I\}, \mathbb{Z}/2\mathbb{Z}) \hookrightarrow H^2(A_5; \mathbb{Z}/2\mathbb{Z})$ is injective (\texttt{sl2\_\allowbreak f5\_\allowbreak leray\_\allowbreak transgression\_\allowbreak injective}).
\end{remark}

\subsection{Iterated Geometric Morphisms and Tower Conservation}

Because bounded geometric morphisms compose, Base-Fiber Conservation extends directly to \emph{towers of relative topoi}:
\[
  \mathcal{E}_n \xrightarrow{\gamma_n} \mathcal{E}_{n-1} \xrightarrow{\gamma_{n-1}} \cdots \xrightarrow{\gamma_2} \mathcal{E}_1 \xrightarrow{\gamma_1} \mathcal{S}.
\]
Along any composite geometric morphism $\gamma = \gamma_1 \circ \cdots \circ \gamma_n : \mathcal{E}_n \to \mathcal{S}$, splitting an epimorphism or constructing a spatial point in $\mathcal{E}_n$ is governed step-by-step by the iterated Leray sequences of $\gamma_k : \mathcal{E}_k \to \mathcal{E}_{k-1}$: a theorem $B$ holds in $\mathcal{E}_n$ either by assuming strong choice in $\mathcal{S}$ together with mild localic Booleanness in each fiber, or over a purely constructive base $\mathcal{S}$ by enforcing supercompactness, atomicity, and cohomological vanishing $H^1(\mathcal{E}_k/\mathcal{E}_{k-1}; G_k) = 0$ at every level.

\section{\texorpdfstring{The $E_0$ Tail Residuum, Noncommutative Smooth Moduli, and the Unclassifiability of Exotic $\mathbb{R}^4$'s}{The E0 Tail Residuum, Noncommutative Smooth Moduli, and the Unclassifiability of Exotic R4's}}\label{sec:exotic-r4-e0}

The Base-Fiber Conservation Principle is not confined to classical topos extensions of $\mathrm{ZF}$: in topological and geometric analysis, forbidding discontinuous Choice ($\neg\mathrm{AC}$) over Polish spaces~\cite{Kechris95} immediately reveals hyperfinite residua governing the classification of smooth $4$-manifolds. Recall from the opening pages of our Preface that Alain Connes's irrational-slope torus foliation $\mathbb{T}^2 / \mathcal{F}_\theta$ collapses in $\mathbf{Set}$ while remaining a sharp noncommutative quotient stack; in the profinite setting, its exact hyperfinite twin is the Cantor tail quotient $\mathcal{Q}_{E_0} = 2^{\mathbb{N}} \sslash E_0$\index{Connes quotient Q_E0@Connes quotient $\mathcal{Q}_{E_0} = 2^{\mathbb{N}}/E_0$}.

\subsection{\texorpdfstring{Compactness vs.\ Infinite Ends ($=_{\mathbb{N}}$ vs.\ $E_0$)}{Compactness vs. Infinite Ends (=N vs. E0)}}

Comparing closed smooth manifolds to open $4$-manifolds isolates a dramatic phase transition along Axis~III ($\mathrm{Inf}$) and Axis~I ($\neg\mathrm{AC}$):
\begin{itemize}[leftmargin=1.5em]
  \item \textbf{Compact Smooth Manifolds ($\neg\mathrm{Inf}$ on charts $\implies \, =_{\mathbb{N}}$):} Bergfalk and Smythe~\cite[Theorem~A]{BergfalkSmythe25} prove that diffeomorphism of \emph{compact} smooth $n$-manifolds ($n \ge 0$) is Borel equivalent to equality $=_{\mathbb{N}}$ on $\mathbb{N}$\index{Borel reducibility!compact manifolds@Compact manifolds ($=_{\mathbb{N}}$)}---represented internally by the discrete $0$-sheaf $\Delta(\mathbb{N})$ with vanishing residuum.
  \item \textbf{Radial Exotic $\mathbb{R}^4$'s vs.\ Full $\mathcal{R}^4$ ($\mathrm{Inf} + \neg\mathrm{AC} \implies E_0$):} In stark contrast, on Gompf and Panagiotopoulos's~\cite{GompfPanagiotopoulos26} explicit radial family $\{M_x\}_{x \in 2^{\mathbb{N}}}$ parameterized along a compact exhaustion $K_0 \Subset K_1 \Subset \cdots \subset \mathbb{R}^4$, oriented diffeomorphism satisfies $M_x \cong_{\mathrm{Diff}} M_y \iff x \,E_0\, y$. On the full moduli space $\mathcal{R}^4$ of all smoothings of $\mathbb{R}^4$, eventual tail equality $E_0$\index{Exotic smooth R4@Exotic smooth $\mathbb{R}^4$!E0@$E_0$ tail equivalence} Borel-reduces into oriented diffeomorphism ($E_0 \le_B \mathrm{Diff}(\mathcal{R}^4)$), embedding $2^{\mathbb{N}}/E_0$ as a Borel sub-moduli problem (climbing further above $=^+$ on general open $4$-manifolds, while remaining bounded above by continuous $S_\infty = \mathrm{Sym}(\mathbb{N})$ orbit equivalence). By the Harrington--Kechris--Louveau $E_0$-dichotomy~\cite[Theorem~1.1]{HKL90}, no Borel Polish or real-valued invariant can classify smooth structures on $\mathbb{R}^4$.
\end{itemize}

\subsection{\texorpdfstring{The Double-Solvent Theorem for $4$-Dimensional Smoothings}{The Double-Solvent Theorem for 4-Dimensional Smoothings}}

Identifying $2^{\mathbb{N}} \cong \prod_{n \in \mathbb{N}} \mathbb{F}_2$ recovers the exact Axis~III biproduct-tail quotient of Chapter~\ref{ch:foundations} (Section~\ref{sec:finset-infinity}) at $A = \mathbb{F}_2$:
\[
  0 \longrightarrow \bigoplus_{n \in \mathbb{N}} \mathbb{F}_2 \longrightarrow \prod_{n \in \mathbb{N}} \mathbb{F}_2 \xrightarrow{\;q\;} \mathcal{Q}_{E_0} \longrightarrow 0.
\]
Smooth ends of $\mathbb{R}^4$ sit suspended inside Clausen and Scholze's light condensed topos $\mathbf{Sh}(\mathrm{Prof}_{\aleph_0})$\index{Light condensed sets@Light condensed sets $\mathbf{Sh}(\mathrm{Prof}_{\aleph_0})$} (sheaves on countably presented profinite sets) between two opposite foundational collapses:

\begin{theorem}[The Double-Solvent Theorem for $4$-Dimensional Smoothings]\label{thm:double-solvent-r4}\index{Double-Solvent Theorem}%
\lean{HomotopyCRM.exists_zfc_linear_section_e0Quotient, HomotopyCRM.nonempty_linearEquiv_cantor_prod_e0Quotient, HomotopyCRM.dense_e0Coset, HomotopyCRM.continuous_e0Invariant_const, HomotopyCRM.not_exists_continuous_t1_e0Separator, HomotopyCRM.chapter4_borel_tail_residuum_summary}%
\leanok%
\uses{def:comparison-triple,def:residuum,prin:conservation}%

In the light condensed topos $\mathbf{Sh}(\mathrm{Prof}_{\aleph_0})$, the $E_0$ smooth-end modulus $\mathcal{Q}_{E_0} = 2^{\mathbb{N}} \sslash E_0 = \bigl[\prod_{n \in \mathbb{N}} \mathbb{F}_2 \sslash \bigoplus_{n \in \mathbb{N}} \mathbb{F}_2\bigr]_{\mathrm{Cond}}$ mediates between topological and set-theoretic collapse:
\[
  \mathbf{Top}(\mathbb{R}^4) \simeq * \;\longleftarrow\; \Bigl[\prod_{n \in \mathbb{N}} \mathbb{F}_2 \;\sslash\; \bigoplus_{n \in \mathbb{N}} \mathbb{F}_2\Bigr]_{\mathrm{Cond}} \;\xrightarrow{\;\mathrm{AC}\;}\; \mathbf{Set}_{\mathrm{ZFC}}(\pi_0) \cong \mathbb{R},
\]
according to three complementary regimes:
\begin{enumerate}[label=\textup{(\roman*)},leftmargin=2.0em]
  \item \textbf{Left Solvent (Freedman $\nu F_{\mathrm{CH}}$ Collapse):} Passing from finite smooth handlebody trees ($\mu F$) to topological Casson handles absorbs infinitely many stages simultaneously, collapsing $\mathcal{Q}_{E_0}$ to $*$ in $\mathbf{Top}(\mathbb{R}^4)$.
  \item \textbf{Right Solvent ($\mathrm{ZFC}$ Vitali Collapse):} Invoking full external $\mathrm{AC}$ (Zorn's Lemma) splits $q$ via a discontinuous $\mathbb{F}_2$-linear Hamel section $s_{\mathrm{AC}}$, producing an artificial cardinal classifier $|\mathcal{Q}_{E_0}| = |\mathbb{R}|$ that destroys smooth-invariant definability.
  \item \textbf{Light Condensed / Baire Topos Residuum:} Over $\mathbf{Sh}(\mathrm{Prof}_{\aleph_0})$, compact interior isotopies absorb $\bigoplus_{n \in \mathbb{N}} \mathbb{F}_2$. Density of $\bigoplus_{n \in \mathbb{N}} \mathbb{F}_2$ forces separated Polish reflection collapse $\pi_0^{\mathrm{Pol}}(\mathcal{Q}_{E_0}) \simeq 1$ (Baire/sheaf blindness of Polish real invariants), while Base--Fiber Conservation precipitates the non-splitting condensed torsor class
  \[
    [q] \neq 0 \;\in\; \mathrm{Ext}^1_{\mathrm{CondAb}}\Bigl(\mathcal{Q}_{E_0},\, \bigoplus_{n \in \mathbb{N}} \mathbb{F}_2\Bigr),
  \]
  paired with Panagiotopoulos's non-vanishing definable derived limit $\varprojlim^1_{\mathrm{def}}$ along compact exhaustions $K_0 \Subset K_1 \Subset \cdots \subset \mathbb{R}^4$.
\end{enumerate}

\end{theorem}

The algebraic splitting, Cantor-space density collapse, and continuous $T_1$ blindness of Theorem~\ref{thm:double-solvent-r4} are machine-checked in Lean~4 (\texttt{HomotopyCRM/\allowbreak BorelTailResiduum.lean}\index{BorelTailResiduum@\texttt{BorelTailResiduum} Lean~4 module}, 682 lines, 19 \texttt{\#print axioms} checks under the 4-Layer Definition Check Gate): the module formalizes the short exact sequence $0 \to (\mathbb{N} \to_0 \mathrm{ZMod}\;2) \to (\mathbb{N} \to \mathrm{ZMod}\;2) \xrightarrow{q} \mathcal{Q}_{E_0} \to 0$ (\texttt{exact\_\allowbreak finFlipIncl\_\allowbreak e0QuotientMk}), certifies the discontinuous $\mathsf{ZFC}$ linear right-inverse section via \texttt{Submodule.exists\_\allowbreak isCompl} (\texttt{exists\_\allowbreak zfc\_\allowbreak linear\_\allowbreak section\_\allowbreak e0Quotient} and \texttt{nonempty\_\allowbreak linearEquiv\_\allowbreak cantor\_\allowbreak prod\_\allowbreak e0Quotient}), proves coordinatewise coset density in the product Cantor topology without Hamel bases (\texttt{dense\_\allowbreak e0Coset}), and deduces that every continuous $E_0$-invariant map into a $T_1$ space is constant (\texttt{continuous\_\allowbreak e0Invariant\_\allowbreak const} and \texttt{not\_\allowbreak exists\_\allowbreak continuous\_\allowbreak t1\_\allowbreak e0Separator}).

\subsection{\texorpdfstring{Noncommutative $C^*$-Moduli ($\mathrm{UHF}(2^\infty)$), Glimm's Dichotomy, and the $E_0^{\mathrm{even}}$ Parity Grading}{Noncommutative C*-Moduli (UHF(2\^{}infty)), Glimm's Dichotomy, and the E0-even Parity Grading}}\label{subsec:uhf-cstar-e0}

The Polish blindness $\pi_0^{\mathrm{Pol}}(\mathcal{Q}_{E_0}) \simeq 1$ admits an exact noncommutative geometric reformulation~\cite[Ch.~II, \S\S 7--8]{Connes94}. Viewing Cantor space $2^{\mathbb{N}} \cong \prod_{n \in \mathbb{N}} \mathbb{F}_2$ under the action of compact interior isotopies (finite coordinate flips $\bigoplus_{n \in \mathbb{N}} \mathbb{F}_2$), orbit density collapses the commutative fixed-point subalgebra of continuous invariants to scalars:
\[
  C(2^{\mathbb{N}})^{\bigoplus_{n \in \mathbb{N}} \mathbb{F}_2} \;=\; \mathbb{C}\cdot 1,
\]
reflecting $\pi_0^{\mathrm{Pol}}(\mathcal{Q}_{E_0}) \simeq 1$. In Connes's noncommutative geometry, rather than discarding the equivalences, one forms the étale groupoid\index{Connes groupoid G_E0@Connes groupoid $\mathcal{G}_{E_0}$} $\mathcal{G}_{E_0} := 2^{\mathbb{N}} \rtimes \bigoplus_{n \in \mathbb{N}} \mathbb{F}_2$ and evaluates its groupoid convolution $C^*$-algebra~\cite[Ch.~II]{Connes94} in two elementary steps:
\begin{enumerate}[label=\textbf{\arabic*.},leftmargin=1.8em]
  \item At each collar stage $n \in \mathbb{N}$, diagonal functions $C(\mathbb{F}_2)$ together with the Pauli-$X$ flip unitary $U_n = \bigl(\begin{smallmatrix} 0 & 1 \\ 1 & 0 \end{smallmatrix}\bigr)$ satisfy the covariant crossed-product relations generating $C(\mathbb{F}_2) \rtimes \mathbb{F}_2 \cong M_2(\mathbb{C})$.
  \item Forming the infinite inductive tensor product across all collar stages $n \in \mathbb{N}$ identifies the groupoid $C^*$-algebra of $\mathcal{G}_{E_0}$ with Glimm's \textbf{uniformly hyperfinite CAR algebra}\index{UHF(2^infty) CAR algebra@$\mathrm{UHF}(2^\infty)$ (CAR algebra)}\index{CAR algebra|see{$\mathrm{UHF}(2^\infty)$}} $\mathrm{UHF}(2^\infty)$~\cite{Glimm60}:
  \[
    \begin{aligned}
      C^*(\mathcal{G}_{E_0}) \;=\; C(2^{\mathbb{N}}) \rtimes \bigoplus_{n \in \mathbb{N}} \mathbb{F}_2
      &\;\cong\; \bigotimes_{n=0}^{\infty} M_2(\mathbb{C}) \\
      &\;=\; \varinjlim_N M_{2^N}(\mathbb{C}) \;=\; \mathrm{UHF}(2^\infty).
    \end{aligned}
  \]
\end{enumerate}

Unlike the collapsed commutative invariant ring $\mathbb{C}\cdot 1$, the noncommutative modulus $\mathrm{UHF}(2^\infty)$ is a simple $C^*$-algebra with trivial center $\mathbb{C}\cdot 1$ that retains complete ordered $K$-theoretic memory of dyadic smooth-handle subdivisions~\cite[Ch.~II]{Connes94}:
\[
  \begin{aligned}
    \bigl(K_0(\mathrm{UHF}(2^\infty)),\; K_0(\mathrm{UHF}(2^\infty))^+,\; [1]\bigr) &\;\cong\; \bigl(\mathbb{Z}[\tfrac{1}{2}],\; \mathbb{Z}[\tfrac{1}{2}]_{\ge 0},\; 1\bigr), \\
    K_1(\mathrm{UHF}(2^\infty)) &\;=\; 0.
  \end{aligned}
\]
Moreover, under GNS completion with respect to its unique tracial state, both Rieffel's irrational rotation algebra $A_\theta$ (representing Connes's torus foliation $\mathbb{T}^2/\mathcal{F}_\theta$ from the Preface) and $\mathrm{UHF}(2^\infty)$ complete to the Murray--von~Neumann \textbf{hyperfinite $\mathrm{II}_1$ factor} $\mathcal{R}$\index{Hyperfinite II1 factor@Hyperfinite $\mathrm{II}_1$ factor $\mathcal{R}$}~\cite[Ch.~V]{Connes94}---turning the Preface analogy between $\mathbb{T}^2/\mathcal{F}_\theta$ and $\mathcal{Q}_{E_0}$ into an exact von~Neumann isomorphism on tracial closures.

This operator-algebraic translation illuminates the historical origin of $E_0$: Glimm~\cite[Theorem~1, \S 9]{Glimm61}\index{Glimm Type-I dichotomy} proved that a separable $C^*$-algebra $A$ is non-Type-I (equivalently, its unitary dual $\widehat{A}$ fails to be a standard Borel space) if and only if $\mathrm{UHF}(2^\infty)$ embeds as a subquotient; Effros and Harrington--Kechris--Louveau~\cite[Theorem~1.1]{HKL90}\index{Glimm--Effros--HKL E0 dichotomy@Glimm--Effros--HKL $E_0$ dichotomy} translated Glimm's theorem directly into the $E_0$-dichotomy of descriptive set theory. Thus $E_0 \le_B \mathrm{Diff}(\mathcal{R}^4)$ expresses in precise geometric form that $4$-dimensional smoothings cross Glimm's boundary from commutative Type-I moduli into genuinely noncommutative geometry.

\begin{remark}[Pauli-$Z$ Parity Grading ($E_0^{\mathrm{even}}$) and Anchor Witness~I ($\mathcal{T}_{\mathrm{odd}}$)]\label{rem:pauli-z-parity}%
\lean{HomotopyCRM.dense_evenE0Coset, HomotopyCRM.not_exists_continuous_flipAlternating, HomotopyCRM.exists_zfc_linear_flipAlternating, HomotopyCRM.evenToE0Quotient_surjective, HomotopyCRM.exists_zfc_linear_section_evenToE0Quotient}%
\leanok%
\uses{thm:double-solvent-r4}%
Restricting finite coordinate flips to the kernel $\mathrm{Even}_0(\mathbb{N}) := \ker\bigl(\Sigma : \bigoplus_{\mathbb{N}} \mathbb{F}_2 \to \mathbb{F}_2\bigr)$ of even-cardinality supports defines the subrelation $E_0^{\mathrm{even}} \subset E_0$ and the $\mathbb{F}_2$-torsor exact sequence\index{Parity torsor!E0 even@$2^{\mathbb{N}}/E_0^{\mathrm{even}} \twoheadrightarrow \mathcal{Q}_{E_0}$}\index{E0 even@$E_0^{\mathrm{even}}$|see{Parity torsor}}
\[
  0 \longrightarrow \mathbb{F}_2 \longrightarrow 2^{\mathbb{N}}/E_0^{\mathrm{even}} \xrightarrow{\;\pi\;} \mathcal{Q}_{E_0} \longrightarrow 0.
\]
In light-condensed or Baire universes ($\mathsf{ZF}+\mathsf{DC}_\omega+\mathsf{BP}$), neither $q : 2^{\mathbb{N}} \twoheadrightarrow \mathcal{Q}_{E_0}$ nor the double cover $\pi$ splits; by contrast, the Boolean Prime Ideal Theorem ($\mathsf{BPI}$) extends the finite-support parity sum $\Sigma$ to a discontinuous homomorphism splitting $\pi$ while remaining strictly independent of full $E_0$-transversals (Geschke--Lubarsky--Rahn; Serafin~\cite{Serafin23}). In \texttt{HomotopyCRM/\allowbreak BorelTailResiduum.lean}, we machine-check both sides of this parity dichotomy: every $E_0^{\mathrm{even}}$-coset is already dense in Cantor space (\texttt{dense\_\allowbreak evenE0Coset}), so \emph{no continuous} single-flip-alternating functional $\chi : 2^{\mathbb{N}} \to \mathbb{F}_2$ can exist (\texttt{not\_\allowbreak exists\_\allowbreak continuous\_\allowbreak flipAlternating}), whereas Zorn's Lemma ($\mathsf{ZFC}$ / \texttt{Classical.choice}) produces discontinuous $\mathbb{F}_2$-linear sections splitting both $\pi$ and parity evaluation (\texttt{exists\_\allowbreak zfc\_\allowbreak linear\_\allowbreak flipAlternating}, \texttt{exists\_\allowbreak zfc\_\allowbreak linear\_\allowbreak section\_\allowbreak evenToE0Quotient}, \texttt{chapter4\_\allowbreak borel\_\allowbreak tail\_\allowbreak residuum\_\allowbreak summary}). On the $C^*$-algebra $\mathrm{UHF}(2^\infty)$, the diagonal Pauli-$Z$ conjugation automorphism
\[
  \alpha_{-1} \;:=\; \bigotimes_{n=0}^{\infty} \mathrm{Ad}\,\begin{pmatrix} 1 & 0 \\ 0 & -1 \end{pmatrix}
\]
acts on finite-flip unitaries $U_S = \prod_{n \in S} U_n$ by $\alpha_{-1}(U_S) = (-1)^{|S|} U_S$: its $+1$ fixed-point subalgebra is precisely $C^*(\mathcal{G}_{E_0^{\mathrm{even}}})$, while its $-1$ eigenspace forms the odd-flip imprimitivity bimodule---yielding the exact operator-algebraic counterpart of Anchor Witness~I, the nominal odd-finset parity torsor $\mathcal{T}_{\mathrm{odd}} \in H^1(\mathbf{Nom}; \mathrm{EvenFinset})$ of Chapter~\ref{ch:realizability-sites}.
\end{remark}

\subsection{The Five-Rung CRM--Borel Complexity Dictionary}

Embedding Borel equivalence relations as condensed and nominal quotient stacks calibrates Borel reducibility directly against the Categorical Reverse Mathematics residuum hierarchy, summarized in Table~\ref{tab:crm-borel-dictionary} and illustrated around the $\mathcal{Q}_{E_0}$ hub in Figure~\ref{fig:exotic-r4-hub}. Above smooth Polish invariants ($\le_B\, =_{\mathbb{R}}$, such as isometry of complete hyperbolic $2$-manifolds with finitely generated $\pi_1$, Bergfalk--Smythe~\cite[Theorem~B]{BergfalkSmythe25}) and the hyperfinite $E_0$ threshold ($E_0 \le_B$ already for complete hyperbolic $3$-manifolds with finitely generated $\pi_1$,~\cite[Theorem~C]{BergfalkSmythe25}), Rung~4 spans equivalence relations \emph{classifiable by countable structures} ($\le_B E_{S_\infty}$, Borel-reducible to the universal logic action of $S_\infty = \mathrm{Sym}(\mathbb{N})$ on countable $\mathcal{L}$-structures~\cite[\S2.7]{BeckerKechris96}): isometry of arbitrary complete hyperbolic $n$-manifolds ($n \ge 2$) is essentially countable universal ($\sim_B E_\infty$~\cite[Theorem~D]{BergfalkSmythe25}, generalizing Riemann surfaces~\cite{HjorthKechris00}), while homeomorphism and diffeomorphism of connected noncompact $2$-manifolds achieve full $S_\infty$-completeness~\cite[Theorem~E]{BergfalkSmythe25}, internalized via nominal and closed-subgroup Fra\"iss\'e quotient stacks over $\mathbf{Nom} \simeq \mathbf{Cont}(\mathrm{Sym}(\mathbb{N}))$ along Kihara's 5-rung $\mathbf{Cont}(\mathrm{Aut}(\mathcal{M}))$ spectrum~\cite{Kihara26}. Crossing into Hjorth turbulence ($\not\le_B E_{S_\infty}$)~\cite[Ch.~3, Theorem~3.21]{Hjorth00} at Rung~5---exemplified by biholomorphism of complex $n$-manifolds for $n \ge 2$~\cite{HjorthKechris00}---forces every Baire-measurable or Borel invariant into a continuous $S_\infty$-action to stabilize generically on orbits (\emph{pro-oligomorphic blindness}~\cite{Hjorth00,Kihara26}). Remarkably, replacing additive direct sums by finite-prefix sequences inside $\mathcal{Q} \subset \prod_{k \ge 1} \mathcal{A}_k$ reproduces the hyperfinite tail residuum at the summit of directed higher category theory: the step-product Postnikov phantom $X_\bullet \in RL(\mathcal{Q}) \setminus \mathcal{Q}$ of Chapter~\ref{ch:higher-crm} (Theorem~\ref{thm:step-product-counterexample}) is the directed $(\infty,\infty)$-categorical twin of $\mathcal{Q}_{E_0}$, while Horizon~VII in Chapter~\ref{ch:master-atlas} isolates the open problems linking $\mathrm{Ext}^1_{\mathrm{CondAb}}$, $\varprojlim^1_{\mathrm{def}}$, and pro-oligomorphic blindness.

\begin{figure}[htbp]
\centering
\begin{tikzpicture}[
  >=Stealth,
  hub/.style={draw=CRMBlue,line width=1.1pt,rounded corners=4pt,fill=CRMBlue!10,inner sep=5pt,font=\scriptsize,align=center},
  wing/.style={draw=CRMBlue!75,rounded corners=3pt,fill=CRMBlue!4,inner sep=4pt,font=\scriptsize,align=center},
  arw/.style={->,thick,draw=CRMBlue}
]
  \node[hub] (C) at (0,0) {$\mathcal{Q}_{E_0} = \bigl[\prod_{\mathbb{N}}\mathbb{F}_2 \sslash \bigoplus_{\mathbb{N}}\mathbb{F}_2\bigr]_{\mathrm{Cond}}$\\[-1pt]\tiny $\in \mathbf{Sh}(\mathrm{Prof}_{\aleph_0})$;\; $[q]\neq 0$,\; $\varprojlim^1_{\mathrm{def}}\neq 0$};
  \node[wing] (W) at (-4.1,0) {\textbf{Freedman Collapse}\\[1pt]$\mathbf{Top}(\mathbb{R}^4)\simeq *$\\[-1pt]\tiny$\nu F_{\mathrm{CH}}$ handles};
  \node[wing] (E) at (4.1,0) {\textbf{$\mathsf{ZFC}$ Vitali}\\[1pt]$\mathbf{Set}_{\mathrm{ZFC}}(\pi_0)\cong\mathbb{R}$\\[-1pt]\tiny Hamel ($\mathsf{AC}$)};
  \node[wing] (S) at (0,-2.05) {\textbf{Polish Reflection (Baire Blindness)}\\[1pt]$\pi_0^{\mathrm{Pol}}(\mathcal{Q}_{E_0})\simeq 1,\qquad C(2^{\mathbb{N}})^{\bigoplus\mathbb{F}_2}=\mathbb{C}\cdot 1$};
  \node[wing] (N) at (0,2.15) {\textbf{Connes Noncommutative $C^*$-Modulus}\\[1pt]$C^*(\mathcal{G}_{E_0}) \cong \mathrm{UHF}(2^\infty) \xrightarrow{\;K_0\;} \mathbb{Z}[\tfrac{1}{2}]$\\[-1pt]\tiny Pauli-$Z$ parity $\alpha_{-1} \leftrightarrow E_0^{\mathrm{even}}$\quad ($\mathrm{II}_1$ factor $\mathcal{R}$)};

  \draw[arw] (C) -- node[above,font=\tiny]{left} (W);
  \draw[arw] (C) -- node[above,font=\tiny]{$\mathsf{AC}$} (E);
  \draw[arw] (C) -- node[right,font=\tiny]{$\pi_0^{\mathrm{Pol}}$} (S);
  \draw[arw] (C) -- node[right,font=\tiny]{$\mathcal{G}_{E_0}$} (N);
\end{tikzpicture}
\caption{Exotic-$\mathbb{R}^4$ Double-Solvent and Noncommutative $\mathrm{UHF}(2^\infty)$ Moduli Hub.}\label{fig:exotic-r4-hub}
\end{figure}

\begin{table}[htbp]
\centering
\scriptsize
\setlength{\tabcolsep}{3pt}
\renewcommand{\arraystretch}{1.14}
\begin{tabularx}{\textwidth}{@{}P{0.19\textwidth}P{0.23\textwidth}Y@{}}
\toprule
\textbf{Borel Complexity Rung} & \textbf{Geometric / Manifold Witness} & \textbf{CRM Quotient Stack \& Residuum Witness} \\
\midrule
1.\ Countable ($=_{\mathbb{N}}$) & Compact smooth $n$-manifolds~\cite[Thm.~A]{BergfalkSmythe25} & Discrete $0$-sheaf $\Delta(\mathbb{N})$; zero residuum ($\neg\mathrm{Inf}$ on finite atlases). \\
2.\ Smooth ($\le_B\, =_{\mathbb{R}}$) & Finite-type hyperbolic $2$-manifolds~\cite[Thm.~B]{BergfalkSmythe25} and Kleinian groups & Separated Polish $0$-sheaf $\mathcal{M} \xrightarrow{\sim} \pi_0^{\mathrm{Pol}}(\mathcal{M})$; $H^1$ torsor vanishes. \\
3.\ Hyperfinite ($E_0$) & Radial exotic $\mathbb{R}^4$'s~\cite{GompfPanagiotopoulos26} & Profinite Connes stack $\bigl[\prod_{\mathbb{N}} \mathbb{F}_2 \sslash \bigoplus_{\mathbb{N}} \mathbb{F}_2\bigr]_{\mathrm{Cond}}$; $\pi_0^{\mathrm{Pol}} \simeq 1$, non-trivial $\mathrm{Ext}^1_{\mathrm{CondAb}}$ \& $\varprojlim^1_{\mathrm{def}}$ tail class; $C^*(\mathcal{G}_{E_0}) \cong \mathrm{UHF}(2^\infty)$ ($K_0 \cong \mathbb{Z}[\tfrac{1}{2}]$). \\
4.\ Countable Structures ($E_\infty \le_B E \le_B E_{S_\infty}$) & Complete hyperbolic $n$-manifolds ($n \ge 2$, $E_\infty$~\cite[Thm.~D]{BergfalkSmythe25}); open $2$-manifolds ($S_\infty$-complete~\cite[Thm.~E]{BergfalkSmythe25}) & Borel reduction to $S_\infty = \mathrm{Sym}(\mathbb{N})$ logic actions~\cite{BeckerKechris96}; nominal / Fra\"iss\'e quotient stacks over $\mathbf{Nom} \simeq \mathbf{Cont}(\mathrm{Sym}(\mathbb{N}))$ along Kihara's 5-rung spectrum~\cite{Kihara26} ($H^1(\mathbf{Nom}; G)$ torsors). \\
5.\ Turbulent ($\not\le_B E_{S_\infty}$) & Complex $n$-manifolds ($n \ge 2$) up to biholomorphism~\cite{HjorthKechris00,Hjorth00} & \emph{Pro-oligomorphic blindness}\index{Pro-oligomorphic blindness}: every definable functor into $\mathbf{Nom}$ or any Fra\"iss\'e groupoid collapses generically~\cite{Hjorth00,Kihara26}. \\
\bottomrule
\end{tabularx}
\caption{The Five-Rung CRM--Borel Complexity Dictionary calibrating manifold classification complexity against condensed and nominal residua.}
\label{tab:crm-borel-dictionary}
\end{table}

\part{\texorpdfstring{Higher Categories, Site Combinatorics, and Set-Theoretic Foundations}{Higher Categories, Site Combinatorics, and Set-Theoretic Foundations}}

\chapter{\texorpdfstring{Higher Categorical Reverse Mathematics in $(\infty, 1)$-Topoi, Univalent Foundations, and Postnikov Obstruction Towers}{Higher Categorical Reverse Mathematics in (infinity, 1)-Topoi, Univalent Foundations, and Postnikov Obstruction Towers}}
\chaptermark{Higher CRM in $(\infty, 1)$-Topoi and Postnikov Towers}
\label{ch:higher-crm}
\chapterepigraph{``Good general theory does not search for the maximum generality, but for the right generality.''}{Saunders Mac~Lane, \emph{Categories for the Working Mathematician} (2nd ed., 1998)}

\section{\texorpdfstring{From 1-Topoi to $(\infty, 1)$-Topoi and Univalent Universes}{From 1-Topoi to (infinity, 1)-Topoi and Univalent Universes}}\label{sec:higher-topoi-intro}

In the first four chapters of this book, our ambient mathematical universes were primarily elementary and Grothendieck 1-topoi $\mathcal{E}$. In a 1-topos, hom-sets $\mathrm{Hom}_{\mathcal{E}}(X, Y)$ are discrete $0$-truncated sets, equality of morphisms is a strict proposition, and the subobject classifier $\Omega$ classifies $(-1)$-truncated monomorphisms. As established along Axis 1 of our relativization framework (Chapter~\ref{ch:base-fiber}, Sections~\ref{sec:two-axis-relativization} and~\ref{sec:axis1-truncation}), removing the requirement of $0$-truncation elevates Categorical Reverse Mathematics ($\mathrm{CRM}$) from 1-topoi to an $(\infty, 1)$-topos $\mathcal{X}$ (in the sense of Joyal, Lurie~\cite{Lurie09}, Rezk, and To\"en--Vezzosi).

An $(\infty, 1)$-topos $\mathcal{X}$\index{infinity-topos@$(\infty, 1)$-topos} is a locally presentable (or elementary) $(\infty, 1)$-category equipped with homotopy limits and colimits, universal colimits (pullback-stable effective epimorphisms), descent (van Kampen colimits), and a hierarchy of \emph{object classifiers} $\mathcal{U}$~\cite[\S 6.1.6]{Lurie09}. Its internal language is \emph{Homotopy Type Theory (HoTT)} and Voevodsky's \emph{Univalent Foundations}~\cite{HoTT13,GepnerKock17}, in which objects $X \in \mathcal{X}$ are interpreted not as static sets, but as homotopy types ($\infty$-groupoids) stratified by their \emph{truncation level}:

\begin{itemize}[leftmargin=1.5em]

  \item \textbf{$(-2)$-Truncated Objects:} Contractible types $X \simeq 1$, possessing a unique point up to coherent higher homotopy.

  \item \textbf{$(-1)$-Truncated Objects:} Mere propositions, whose equivalence classes form the internal subobject classifier $\Omega = \mathcal{U}_{\le -1}$~\cite[\S 6.1.6]{Lurie09}.

  \item \textbf{$0$-Truncated Objects:} Internal discrete sets, whose identity types are $(-1)$-truncated propositions and which form the underlying 1-topos $\tau_{\le 0}\mathcal{X}$~\cite[\S 6.4.5]{Lurie09}.

  \item \textbf{$n$-Truncated Objects ($n \ge 1$):} Homotopy $n$-types ($n$-groupoids) whose homotopy groups $\pi_k(X, x)$ vanish for all $k > n$.

\end{itemize}

For each integer $n \ge -2$, the inclusion of $n$-truncated objects $\tau_{\le n}\mathcal{X} \hookrightarrow \mathcal{X}$ admits a left adjoint truncation functor $\tau_{\le n}: \mathcal{X} \to \tau_{\le n}\mathcal{X}$~\cite[Prop.~5.5.6.18]{Lurie09}. Similarly, a morphism $f: X \to Y$ in $\mathcal{X}$ is $n$-truncated if its homotopy fibers $\mathrm{fib}_y(f) = X \times_Y \{y\}$ are $n$-truncated objects in the slice $(\infty, 1)$-topos $\mathcal{X}/Y$.

Passing from $\tau_{\le 0}\mathcal{X}$ to the full $(\infty, 1)$-topos $\mathcal{X}$ transforms the Fundamental Equivalence of Categorical Reverse Mathematics, $A + \mathcal{R}(\mathcal{X}) \iff B$. In a 1-topos, an epimorphism $e: X \twoheadrightarrow Y$ is simply a map with non-empty $(-1)$-truncated support, and splitting $e$ is a single 1-categorical problem. In an $(\infty, 1)$-topos $\mathcal{X}$, an \emph{effective epimorphism} $e: X \twoheadrightarrow Y$ is a morphism whose $(-1)$-truncation $\tau_{\le -1}(e): \tau_{\le -1}(X \to Y) \xrightarrow{\sim} Y$ is an equivalence in $\mathcal{X}/Y$~\cite[\S 6.2.3, Prop.~7.2.1.14]{Lurie09}. Because the homotopy fibers of $e$ can possess infinitely many non-vanishing higher homotopy groups $\pi_k(\mathrm{fib}_y(e))$, any 1-categorical reversal fractures into an infinite tower of higher homotopical and cohomological residuums.

\section{\texorpdfstring{The Higher Choice Hierarchy ($\mathrm{AC}_{n, m}$) and the Global Infinity-Choice Residuum}{The Higher Choice Hierarchy (ACn, m) and the Global Infinity-Choice Residuum}}\label{sec:higher-choice}

To map the Choice-Geometry Residuum Atlas inside an $(\infty, 1)$-topos $\mathcal{X}$, we must first observe how truncation levels interact with the domain and codomain of an effective epimorphism $e: X \twoheadrightarrow Y$.

\begin{definition}[Bigraded Higher Choice]\label{def:bigraded-choice}

Let $\mathcal{X}$ be an $(\infty, 1)$-topos, and let $-1 \le n, m \le \infty$. We say that $\mathcal{X}$ satisfies the Axiom of $(n, m)$-Choice\index{Axiom of Choice!bigraded@Bigraded choice $\mathrm{AC}_{n,m}$}, denoted $\mathrm{AC}_{n, m}$, if every $n$-truncated effective epimorphism $e: X \twoheadrightarrow Y$ whose codomain $Y$ is an $m$-truncated object admits a section $s: Y \to X$ ($e \circ s \simeq \mathrm{id}_Y$).

\end{definition}

Why must the higher choice spine $\mathrm{AC}_n := \mathrm{AC}_{n, 0}$ restrict the codomain $Y$ to be $0$-truncated ($m = 0$)? As Vladimir Voevodsky's \textbf{Univalence Collapse Theorem}~\cite[Ex.~3.8.5]{HoTT13} (cf.~\cite[\S\S 5.5.6, 6.5]{Lurie09}) demonstrates, demanding sections over even $1$-truncated or untruncated codomains ($m \ge 1$) collides directly with univalence and destroys higher groupoids:

\begin{proposition}[Degeneracy of Un-Truncated Choice]\label{prop:untruncated-degeneracy}

If an $(\infty, 1)$-topos $\mathcal{X}$ satisfies $\mathrm{AC}_{0, 1}$ (every $0$-truncated effective epimorphism onto a $1$-truncated codomain splits), then every connected $1$-groupoid in $\mathcal{X}$ is contractible, collapsing all $1$-truncated homotopy types into $0$-truncated sets.

\end{proposition}

\begin{proof}
Let $G$ be any non-trivial $0$-truncated group in $\mathcal{X}$ (for example, $G = \mathbb{Z}/2\mathbb{Z}$), and let $BG = K(G, 1)$ be its delooping $1$-groupoid. Because $BG$ is $0$-connected ($\tau_{\le 0}(BG) \simeq 1$), the basepoint inclusion $e: 1 \twoheadrightarrow BG$ is a $0$-truncated effective epimorphism (its homotopy fiber is the discrete set $G$) onto a $1$-truncated codomain $BG$. If $\mathrm{AC}_{0, 1}$ held in $\mathcal{X}$, $e$ would admit a section $s: BG \to 1$ with $e \circ s \simeq \mathrm{id}_{BG}$, forcing $BG \simeq 1$ and $G \simeq \Omega(BG) \simeq 1$.
\end{proof}

Proposition~\ref{prop:untruncated-degeneracy} proves that demanding sections over un-truncated codomains ($m \ge 1$) is an anti-homotopical truncation axiom rather than a choice principle. Consequently, genuine Higher Categorical Reverse Mathematics calibrates the spine where the codomain $Y$ is a $0$-truncated set ($m = 0$) while the fiber truncation level $n$ climbs from $0$ to $\infty$:

\[
  \mathrm{AC}_\infty \implies \cdots \implies \mathrm{AC}_{n+1} \implies \mathrm{AC}_n \implies \cdots \implies \mathrm{AC}_1 \implies \mathrm{AC}_0,
\]

yielding the model hierarchy $\mathbf{Mod}(\mathrm{AC}_\infty) \subseteq \allowbreak \bigcap_{k < \infty} \mathbf{Mod}(\mathrm{AC}_k) \subsetneq \allowbreak \cdots \subsetneq \allowbreak \mathbf{Mod}(\mathrm{AC}_{n+1}) \subsetneq \allowbreak \mathbf{Mod}(\mathrm{AC}_n) \subsetneq \allowbreak \cdots \subsetneq \allowbreak \mathbf{Mod}(\mathrm{AC}_0)$, where $\mathrm{AC}_n := \mathrm{AC}_{n, 0}$ and $\mathrm{AC}_\infty := \mathrm{AC}_{\infty, 0}$. Indeed, over $\mathbf{Set}_{\mathrm{ZFC}}$, for each integer $n \ge 0$ let $K_{n+2} := K(\mathbb{Z}/2\mathbb{Z}, n+2) \in \mathcal{S}$ be the unstable Eilenberg--MacLane classifying space and consider the slice $(\infty,1)$-topos $\mathcal{X}^{(n)} := \mathcal{S}_{/K_{n+2}} \simeq \mathbf{Fun}(K_{n+2}, \mathcal{S})$~\cite[Proposition~6.3.5.1]{Lurie09}. Because $K_{n+2}$ is $(n+1)$-connected ($\tau_{\le n+1}K_{n+2} \simeq *$) while the core $(\tau_{\le n}\mathcal{S})^{\simeq}$ is $(n+1)$-truncated and mapping spaces in $\tau_{\le n}\mathcal{S}$ are $n$-truncated, the constant-functor embedding $\Delta : \tau_{\le n}\mathcal{S} \xrightarrow{\sim} \tau_{\le n}(\mathcal{X}^{(n)})$ is an equivalence~\cite[Proposition~5.5.6.18]{Lurie09} preserving and reflecting effective epimorphisms. Any $n$-truncated effective epimorphism $e : X \twoheadrightarrow Y$ in $\mathcal{X}^{(n)}$ with $Y \in \tau_{\le 0}(\mathcal{X}^{(n)})$ has $X \in \tau_{\le n}(\mathcal{X}^{(n)})$~\cite[Lemma~5.5.6.14]{Lurie09} and therefore corresponds under $\Delta$ to an effective epimorphism onto a discrete set in $\mathcal{S}$, which splits over $\mathbf{Set}_{\mathrm{ZFC}}$; hence $\mathcal{X}^{(n)} \models \mathrm{AC}_n$. Conversely, in $\mathcal{S}_{/K_{n+2}}$, the basepoint map $x_0 : * \twoheadrightarrow K_{n+2}$ ($n+2 \ge 2$) is $(-1)$-connected---hence an effective epimorphism onto $1_{\mathcal{X}^{(n)}} \in \tau_{\le 0}(\mathcal{X}^{(n)})$~\cite[Proposition~7.2.1.14]{Lurie09}---with $(n+1)$-truncated homotopy fiber $\Omega K_{n+2} \simeq K(\mathbb{Z}/2\mathbb{Z}, n+1)$, yet admits no section in $\mathcal{S}_{/K_{n+2}}$ since $\pi_{n+2}(K_{n+2}) \cong \mathbb{Z}/2\mathbb{Z} \neq 0$; thus $\mathcal{X}^{(n)} \in \mathbf{Mod}(\mathrm{AC}_n) \setminus \mathbf{Mod}(\mathrm{AC}_{n+1})$. At the transfinite limit, $\mathbf{Mod}(\mathrm{AC}_\infty) = \bigcap_{n < \infty} \mathbf{Mod}(\mathrm{AC}_n)$ on hypercomplete $(\infty,1)$-topoi~\cite[\S 6.5.2]{Lurie09} over $\mathbf{Set}_{\mathrm{ZFC}}$ by Theorem~\ref{thm:limit-residuum}, whereas non-hypercomplete separation is governed by Problem~\ref{prob:open-9}.

\begin{itemize}[leftmargin=1.5em]

  \item \textbf{Set-Level Choice ($\mathrm{AC}_0$):} Every $0$-truncated effective epimorphism $X \twoheadrightarrow Y$ onto a $0$-truncated set $Y$ splits; this is equivalent to the classical External Axiom of Choice ($\mathrm{AC}$) in the underlying 1-topos $\tau_{\le 0}\mathcal{X}$.

  \item \textbf{Stack/Groupoid Choice ($\mathrm{AC}_1$):} Every $1$-truncated effective epimorphism $X \twoheadrightarrow Y$ (a stack of 1-groupoids with non-empty connected components) over a $0$-truncated set $Y$ admits a global section.

  \item \textbf{Infinity-Choice ($\mathrm{AC}_\infty$):} Every effective epimorphism $X \twoheadrightarrow Y$ from an arbitrary $\infty$-groupoid $X$ onto a $0$-truncated set $Y$ admits a section (equivalently, every $0$-truncated object $Y \in \tau_{\le 0}\mathcal{X}$ is projective in the full $(\infty, 1)$-topos $\mathcal{X}$).

\end{itemize}

When we ask for the global categorical residuum $(\mathrm{AC}_0 \Rightarrow_{(\infty,1)\text{-}\mathbf{Topos}} \mathrm{AC}_\infty)$ that elevates set-level choice in $\tau_{\le 0}\mathcal{X}$ to full infinity-choice across $\mathcal{X}$, we obtain the higher geometric counterpart of localicity:

\begin{theorem}[Global Infinity-Choice Residuum]\label{thm:global-inf-choice}%
\uses{def:bigraded-choice,prop:untruncated-degeneracy,thm:bdf}%

In any $(\infty, 1)$-topos $\mathcal{X}$, the categorical residuum bridging $0$-truncated choice ($\mathrm{AC}_0$) to full infinity-choice ($\mathrm{AC}_\infty$)\index{Axiom of Choice!infinity-choice@Infinity-choice $\mathrm{AC}_\infty$} is the geometric condition that $0$-truncated objects (sets) cover $\mathcal{X}$:

\[
  \mathrm{AC}_0 + \bigl(\forall X \in \mathcal{X},\, \exists S \in \tau_{\le 0}\mathcal{X} \text{ with } S \twoheadrightarrow X\bigr) \iff \mathrm{AC}_\infty
\]

\end{theorem}

\begin{proof}
Suppose $\mathcal{X} \models \mathrm{AC}_\infty$. Then $\mathrm{AC}_0$ holds as a special case. Moreover, for any object $X \in \mathcal{X}$, the $0$-truncation unit $\eta_X: X \twoheadrightarrow \tau_{\le 0}X$ is an effective epimorphism onto the $0$-truncated object $\tau_{\le 0}X \in \tau_{\le 0}\mathcal{X}$. By $\mathrm{AC}_\infty$, $\eta_X$ admits a section $s: \tau_{\le 0}X \to X$ with $\eta_X \circ s \simeq \mathrm{id}_{\tau_{\le 0}X}$. Because $\tau_{\le 0}(s) = \mathrm{id}_{\tau_{\le 0}X}$ is an effective epimorphism, $s: \tau_{\le 0}X \twoheadrightarrow X$ is an effective epimorphism from a $0$-truncated object covering $X$ (Lurie~\cite{Lurie09}, \emph{Higher Topos Theory}, Prop.~7.2.1.14). Conversely, suppose $\mathcal{X} \models \mathrm{AC}_0$ and every object $X \in \mathcal{X}$ is covered by a $0$-truncated set $q: S \twoheadrightarrow X$ with $S \in \tau_{\le 0}\mathcal{X}$. Let $e: X \twoheadrightarrow Y$ be an effective epimorphism onto a $0$-truncated object $Y \in \tau_{\le 0}\mathcal{X}$. Composing $e$ with a $0$-truncated cover $q: S \twoheadrightarrow X$ yields an effective epimorphism $e \circ q: S \twoheadrightarrow Y$ between $0$-truncated objects, which is a $0$-truncated effective epimorphism in $\tau_{\le 0}\mathcal{X}$. By $\mathrm{AC}_0$, $e \circ q$ admits a section $\sigma: Y \to S$, so $s = q \circ \sigma: Y \to X$ is a section of $e$ ($e \circ s \simeq \mathrm{id}_Y$), proving $\mathcal{X} \models \mathrm{AC}_\infty$.
\end{proof}

\section{\texorpdfstring{Postnikov Towers, $k$-Invariants, and the Stepwise Residuum ($\mathrm{AC}_n \Rightarrow \mathrm{AC}_{n+1}$)}{Postnikov Towers, k-Invariants, and the Stepwise Residuum (ACn => ACn+1)}}\label{sec:postnikov-residuum}

Theorem~\ref{thm:global-inf-choice} characterizes the global transition from $\mathrm{AC}_0$ to $\mathrm{AC}_\infty$. Between adjacent rungs $n$ and $n+1$ ($n \ge 0$), translating the stepwise residuum $\mathcal{R}_n(\mathcal{X}) = (\mathrm{AC}_n \Rightarrow_{(\infty,1)\text{-}\mathbf{Topos}} \mathrm{AC}_{n+1})$ into sheaf cohomology without smuggling $\mathrm{COSHEP}$ (as a $0$-truncated \emph{projective} cover would) or contradicting $1$-localic topoi lacking $\mathrm{AC}_0$ (such as $\mathbf{Sh}_\infty(S^{n+2})$, where $H^{n+2}(S^{n+2}; \mathbb{Z}) \cong \mathbb{Z} \neq 0$ yet vanishes on the $0$-truncated hemisphere cover $U_+ \amalg U_- \twoheadrightarrow 1$) singles out \emph{$0$-truncated effacement}\index{0-truncated effacement@$0$-truncated effacement}: every $(n+1)$-gerbe class over $Y \in \tau_{\le 0}\mathcal{X}$ trivializes after pullback along some $0$-truncated effective epimorphism $q : S \twoheadrightarrow Y$ ($S \in \tau_{\le 0}\mathcal{X}$). Under $\mathrm{AC}_n$ ($n \ge 0 \implies \mathrm{AC}_0$), every such $q$ splits globally in $\tau_{\le 0}\mathcal{X}$, collapsing $0$-truncated effacement onto global slice vanishing over $\mathcal{X}/Y$; without $\mathrm{AC}_0$, effacement holds unconditionally on every $1$-localic $(\infty, 1)$-topos $\mathbf{Sh}_\infty(L)$~\cite[\S6.4.5]{Lurie09} while failing strictly at degree $n+2$ on $\mathcal{X}^{(n)} = \mathcal{S}_{/K(\mathbb{Z}/2\mathbb{Z}, n+2)}$.

Let $e: X \twoheadrightarrow Y$ be an $(n+1)$-truncated effective epimorphism onto $Y \in \tau_{\le 0}\mathcal{X}$. In the slice $(\infty, 1)$-topos $\mathcal{X}/Y$, relative Postnikov truncation\index{Postnikov tower} factors $e$ into a finite tower over $1_{\mathcal{X}/Y}$:
\[
  X \simeq \tau_{\le n+1}(X/Y) \xrightarrow{p_{n+1}} \tau_{\le n}(X/Y) \xrightarrow{p_n} \cdots \xrightarrow{p_1} \tau_{\le 0}(X/Y) \xrightarrow{e_0} 1_{\mathcal{X}/Y},
\]
and pulling $p_{n+1}$ back along a partial section $s_n : 1_{\mathcal{X}/Y} \to \tau_{\le n}(X/Y)$ yields an $n$-connected, $(n+1)$-truncated $(n+1)$-gerbe $E_{n+1} := 1_{\mathcal{X}/Y} \times_{\tau_{\le n}(X/Y)} \tau_{\le n+1}(X/Y) \xrightarrow{\pi_{n+1}} 1_{\mathcal{X}/Y}$~\cite[Def.~7.2.2.20]{Lurie09}. Separating the abelian rungs $n \ge 1$ from the non-abelian base rung $n = 0$ establishes the proved unconditional core prior to Conjecture~\ref{conj:postnikov-stepwise}:

\begin{theorem}[Stepwise Generator Bridge, Unconditional Abelian $k$-Invariant Effacement ($n \ge 1$), and Nilpotent/\'Etale $1$-Gerbe Rungs]\label{thm:postnikov-stepwise-proved}%
\uses{def:bigraded-choice,thm:global-inf-choice}%
Let $\mathcal{X}$ be an $(\infty, 1)$-topos and $n \ge 0$ an integer.
\begin{enumerate}[label=\textup{(\roman*)},leftmargin=2.0em]
  \item \textbf{Generator-Level Bridge ($n \ge 0$):} $\mathrm{AC}_n + \bigl(\forall X \in \tau_{\le n+1}\mathcal{X},\; \exists W \in \tau_{\le n}\mathcal{X}\text{ with } W \twoheadrightarrow X\bigr) \iff \mathrm{AC}_{n+1}$.
  \item \textbf{Abelian Postnikov Rungs ($n \ge 1$):} $\mathrm{AC}_n + \mathcal{R}_n(\mathcal{X}) \iff \mathrm{AC}_{n+1}$ holds unconditionally with $\mathcal{R}_n(\mathcal{X})$ given by $0$-truncated effacement of $H^{n+2}(\mathcal{X}/Y; A)$ for all $Y \in \tau_{\le 0}\mathcal{X}$ and all $A \in \mathbf{Ab}(\tau_{\le 0}(\mathcal{X}/Y))$~\cite[Constr.~7.2.2.21, Thm.~7.2.2.26]{Lurie09}.
  \item \textbf{Non-Abelian Base Rung ($n = 0$):} Over all Giraud bands (\emph{liens}) $\mathfrak{L} \in \mathbf{Lien}^\kappa(\mathcal{X}/Y)$~\cite[Ch.~IV]{Giraud71}, $0$-truncated effacement is unconditionally equivalent to $(\mathrm{AC}_0 \Rightarrow \mathrm{AC}_1)$. Restricting strictly to global group bands $L(G)$ for $G \in \mathbf{Grp}(\tau_{\le 0}(\mathcal{X}/Y))$ via $H^2_{\mathrm{gerbe}}(\mathcal{X}/Y, L(G))$, necessity holds unconditionally under $\mathrm{AC}_1$, and sufficiency holds unconditionally for all locally bounded-nilpotent $1$-gerbes ($\mathrm{Inn}^c(G_x) = \{1\}$, including all locally abelian $1$-gerbes at $c = 1$) as well as on all $1$-localic and \'etale $\mathbf{Sh}_\infty(B)_{/W}$ topoi over $\mathbf{Set}_{\mathrm{AC}}$~\cite[Prop.~6.4.5.7]{Lurie09}.
\end{enumerate}
\end{theorem}

\begin{proof}
For~(i), under $\mathrm{AC}_{n+1}$ the $(n+1)$-truncated unit $\eta_X : X \twoheadrightarrow \tau_{\le 0}X$ ($X \in \tau_{\le n+1}\mathcal{X}$~\cite[Lemma~5.5.6.15]{Lurie09}) splits via a $0$-truncated cover $s : \tau_{\le 0}X \twoheadrightarrow X$~\cite[Prop.~7.2.1.14]{Lurie09}; conversely, any $(n+1)$-truncated $e : X \twoheadrightarrow Y$ ($Y \in \tau_{\le 0}\mathcal{X}$) satisfies $X \in \tau_{\le n+1}\mathcal{X}$~\cite[Lemma~5.5.6.14]{Lurie09}, so covering $X$ by $q : W \twoheadrightarrow X$ ($W \in \tau_{\le n}\mathcal{X}$) renders $e \circ q : W \twoheadrightarrow Y$ $n$-truncated and splits $e$ via $q \circ \sigma$.

For~(ii), let $n \ge 1$. Given $\alpha = [c] \in H^{n+2}(\mathcal{X}/Y; A) =\allowbreak \pi_0\,\mathrm{Map}_{\mathcal{X}/Y}(1_{\mathcal{X}/Y},\allowbreak K(A, n+2))$, pulling the basepoint $0_Y : 1_{\mathcal{X}/Y} \to K(A, n+2)$ back along $c$ yields an $(n+1)$-truncated effective epimorphism $\pi_c : P_c \twoheadrightarrow Y$ with fiber $K(A, n+1)$; splitting $\pi_c$ under $\mathrm{AC}_{n+1}$ forces $\alpha = 0$ (while on a $1$-localic topos, covering $P_c$ by $r : S \twoheadrightarrow P_c$ effaces $q^*\alpha = 0$ along $q = \pi_c \circ r$). Conversely, because $n \ge 1$, the section-lifting gerbe $E_{n+1}$ is \emph{$1$-connected} ($\tau_{\le 1}(E_{n+1}) \simeq 1_{\mathcal{X}/Y}$). Consequently, the local loop-group functor into the $1$-truncated moduli $\mathbf{Ab}(\mathcal{U}_{\le 0})^\simeq$ descends uniquely across $E_{n+1} \twoheadrightarrow 1_{\mathcal{X}/Y}$ to a global abelian group object $\Pi_{n+1} \in \mathbf{Ab}(\tau_{\le 0}(\mathcal{X}/Y))$ classified by $\kappa_{n+1}(e, s_n) \in\allowbreak \pi_0\,\mathrm{Map}_{\mathcal{X}/Y}(1_{\mathcal{X}/Y},\allowbreak K(\Pi_{n+1}, n+2)) =\allowbreak H^{n+2}(\mathcal{X}/Y;\allowbreak \Pi_{n+1})$~\cite[Constr.~7.2.2.21, Thm.~7.2.2.26]{Lurie09}. Effacing $\kappa_{n+1}$ along $q : S \twoheadrightarrow Y$ and splitting $q$ by $\mathrm{AC}_0$ kills $\kappa_{n+1}(e, s_n)$, furnishing $s_{n+1} : 1_{\mathcal{X}/Y} \to E_{n+1} \to X$.

For~(iii), at $n = 0$ the pullback $E_1 := 1_{\mathcal{X}/Y} \times_{\tau_{\le 0}(X/Y)} X$ is only $0$-connected: conjugation by $\mathrm{Inn}(G) \cong G/Z(G)$ prevents $x \mapsto \Omega_x(E_1)$ from descending directly to $\mathbf{Grp}(\tau_{\le 0}(\mathcal{X}/Y))$, descending instead to a Giraud band $\mathfrak{L}(E_1) : 1_{\mathcal{X}/Y} \to \mathbf{Lien}^\kappa := \tau_{\le 1}(\mathcal{M})$ over $\mathcal{M} = \mathcal{U}^\kappa_{\le 1, \ge 1}$~\cite[Ch.~IV]{Giraud71} ($\pi_1(\mathbf{Lien}^\simeq, L(G)) \simeq \mathrm{Out}(G)$, $\pi_2(\mathcal{M}, BG) \simeq Z(G)$;~\cite[Thm.~6.1.6.8, Prop.~7.2.2.11, Rem.~7.2.2.25]{Lurie09}), which gives the unconditional equivalence over all \emph{liens}. Restricted to global $G \in \mathbf{Grp}(\tau_{\le 0}(\mathcal{X}/Y))$, $\mathrm{AC}_1$ splits $\mathcal{G} \simeq BH$ and $\mathrm{AC}_0$ splits the $0$-truncated $\mathrm{Inn}(G)$-torsor $\underline{\mathrm{Isom}}_{\mathbf{Grp}}^{\mathrm{band}}(H, G) \twoheadrightarrow Y$, proving necessity. Pulling the pointed classifier $L : \mathcal{M}_* \simeq \mathbf{Grp}^\kappa(\mathcal{U}_{\le 0})^\simeq \twoheadrightarrow \mathbf{Lien}^\kappa$ back along $\mathfrak{L}(E_1)$ forms the \emph{representation $1$-gerbe} $R(E_1) := 1_{\mathcal{X}/Y} \times_{\mathbf{Lien}^\kappa} \mathcal{M}_* \twoheadrightarrow 1_{\mathcal{X}/Y}$ with local fiber $B\mathrm{Inn}(G_x) = K(G_x/Z(G_x), 1)$ (from $K(G_x, 1) \to \mathrm{fib}(L) \to K(Z(G_x), 2)$). Iterating $R^{(k)}(E_1)$ along upper central quotients $\mathrm{Inn}^k(G_x) = G_x/Z_k(G_x)$ gives $R^{(c)}(E_1) \simeq 1_{\mathcal{X}/Y}$ whenever $\mathrm{Inn}^c(G_x) = \{1\}$, descending step-by-step to $E_1$. Finally, via $\pi^* \dashv \pi_* : \mathcal{X} \rightleftarrows \mathbf{Sh}_\infty(B)$~\cite[Prop.~6.4.5.7]{Lurie09}, any \'etale $\mathcal{X} \simeq \mathbf{Sh}_\infty(B)_{/W}$ over $\mathbf{Set}_{\mathrm{AC}}$ satisfying $\mathrm{AC}_0$ has $W \to \tau_{\le 0}W$ $1$-connected, forcing $\tau_{\le 1}W \simeq \tau_{\le 0}W \in \mathbf{Sh}(B) \models \mathrm{AC}_\infty$ and splitting $E_1$ unconditionally.
\end{proof}

\begin{conjecture}[Postnikov Stepwise Residuum]\label{conj:postnikov-stepwise}%
\uses{thm:postnikov-stepwise-proved}%

For every integer $n \ge 0$ in an $(\infty, 1)$-topos $\mathcal{X}$, the stepwise categorical residuum $\mathcal{R}_n(\mathcal{X}) = (\mathrm{AC}_n \Rightarrow_{(\infty,1)\text{-}\mathbf{Topos}} \mathrm{AC}_{n+1})$ satisfying

\[
  \mathrm{AC}_n + \mathcal{R}_n(\mathcal{X}) \iff \mathrm{AC}_{n+1}
\]

is equivalent to \emph{$0$-truncated effacement} over every $0$-truncated object $Y \in \tau_{\le 0}\mathcal{X}$ (collapsing under $\mathrm{AC}_0$ to global slice vanishing): the effacement of Eilenberg--MacLane cohomology $H^{n+2}(\mathcal{X}/Y; A)$ for every abelian group object $A \in \mathbf{Ab}(\tau_{\le 0}(\mathcal{X}/Y))$ when $n \ge 1$ (proved unconditionally above), and the effacement of Giraud's non-abelian gerbe cohomology $H^2_{\mathrm{gerbe}}(\mathcal{X}/Y, L(G))$ restricted to global $0$-truncated group objects $G \in \mathbf{Grp}(\tau_{\le 0}(\mathcal{X}/Y))$ when $n = 0$ (proved unconditionally for locally bounded-nilpotent bands and on all $1$-localic and \'etale $\mathbf{Sh}_\infty(B)_{/W}$ topoi over $\mathbf{Set}_{\mathrm{AC}}$, and conjecturing that $R(E_1) \twoheadrightarrow 1_{\mathcal{X}/Y}$ splits for non-nilpotent bands across arbitrary non-\'etale $(\infty, 1)$-topoi).

\end{conjecture}

Accordingly, the rungs of the Postnikov choice ladder stratify into three geometric regimes:

\begin{itemize}[leftmargin=1.5em]

  \item \textbf{Rung 0 ($\top \Rightarrow \mathrm{AC}_0$, Proved):} Splitting $0$-truncated surjections in $\tau_{\le 0}\mathcal{X}$, governed over $\mathbf{Set}_{\mathrm{AC}}$ by Booleanness and localicity (Theorem~\ref{thm:bdf}) and cohomologically by $H^1(\tau_{\le 0}\mathcal{X}/Y, G) = 0$ alongside external projectivity of $1$ (Chapter~\ref{ch:choice-atlas}, \S2.4).

  \item \textbf{Rung 1 ($\mathrm{AC}_0 \Rightarrow \mathrm{AC}_1$, Proved for \emph{Liens} / Open for Group Bands):} Splitting $1$-groupoid stacks over sets $Y$, governed unconditionally by $0$-truncated effacement of Giraud $\mathfrak{L}$-gerbes~\cite[Ch.~IV]{Giraud71} and conjecturally by global group bands $L(G)$ ($G \in \mathbf{Grp}(\tau_{\le 0}(\mathcal{X}/Y))$, Conjecture~\ref{conj:postnikov-stepwise}).

  \item \textbf{Higher Rungs ($\mathrm{AC}_n \Rightarrow \mathrm{AC}_{n+1}$ for $n \ge 1$, Proved):} Splitting $(n+1)$-truncated epimorphisms over sets $Y$, governed unconditionally by $0$-truncated effacement (and under $\mathrm{AC}_0$, global vanishing $H^{n+2}(\mathcal{X}/Y; A) = 0$) for internal abelian groups $A \in \mathbf{Ab}(\tau_{\le 0}(\mathcal{X}/Y))$.

\end{itemize}

\section{\texorpdfstring{Hypercompleteness, Whitehead's Principle, and the $\infty$-Connected Residuum}{Hypercompleteness, Whitehead's Principle, and the infinity-Connected Residuum}}

A subtle phenomenon emerges at the transfinite limit $n \to \infty$ of the Postnikov tower. Suppose an $(\infty, 1)$-topos $\mathcal{X}$ satisfies $\mathrm{AC}_n$ for \emph{every} finite integer $n \ge 0$ ($\forall n < \infty,\, \mathcal{X} \models \mathrm{AC}_n$). Does $\mathcal{X}$ automatically satisfy full Infinity-Choice ($\mathrm{AC}_\infty$)?

In general, the answer is \emph{no}, and the categorical residuum separating bounded truncated choice ($\forall n < \infty,\, \mathrm{AC}_n$) from $\mathrm{AC}_\infty$ exposes two fundamental invariants of higher topos theory: \emph{Postnikov convergence (hypercompleteness)} and $\varprojlim^1$ phantom vanishing.

A morphism $f: X \to Z$ in an $(\infty, 1)$-topos $\mathcal{X}$ is called $\infty$-connected if its $n$-truncation $\tau_{\le n}(f): \tau_{\le n}X \xrightarrow{\sim} \tau_{\le n}Z$ is an equivalence for every finite $n \ge 0$ (equivalently, all homotopy group sheaves $\pi_k(\mathrm{fib}(f))$ vanish). In the canonical $(\infty, 1)$-topos of spaces $\mathbf{Spaces} = \infty\text{-}\mathbf{Grpd}$, J. H. C. Whitehead's theorem states that every $\infty$-connected map is an equivalence~\cite[\S 6.5.2]{Lurie09}. In an arbitrary Grothendieck $(\infty, 1)$-topos $\mathcal{X}$ (such as sheaves of spaces over a topological space of infinite covering dimension), however, non-trivial $\infty$-connected maps can exist~\cite[\S 6.5.4, \S 7.2.1]{Lurie09}.

\begin{definition}[Hypercompleteness and Postnikov Convergence]\label{def:hypercomplete}

Let $\mathcal{X}$ be an $(\infty, 1)$-topos.

\begin{itemize}[leftmargin=1.5em]

  \item $\mathcal{X}$ is \emph{hypercomplete} (satisfies the internal Whitehead Principle) if every $\infty$-connected morphism $f: X \to Z$ in $\mathcal{X}$ is an equivalence.

  \item Postnikov towers \emph{converge} in $\mathcal{X}$ if for every object $X \in \mathcal{X}$, the canonical map to the homotopy inverse limit of its Postnikov tower, $X \xrightarrow{\sim} \varprojlim_{n \in \omega} \tau_{\le n}X$, is an equivalence.

\end{itemize}

\end{definition}

Even when Postnikov towers converge in $\mathcal{X}$ ($X \simeq \varprojlim_n \tau_{\le n}X$), having a section $s_n: Y \to \tau_{\le n}X$ at each finite stage $n$ does not guarantee a compatible thread $(s_n)_{n \in \omega}$ in the homotopy inverse limit $\varprojlim_n \tau_{\le n}X$ unless sequential Dependent Choice ($\mathrm{DC}$) or Milnor $\varprojlim^1$ vanishing holds along the tower (and conversely, under $\mathrm{AC}_\infty$, $0$-truncated generation from Theorem~\ref{thm:global-inf-choice} ensures that Postnikov convergence over $0$-truncated targets coincides with hypercompleteness). Combining these ingredients yields the limit residuum:

\begin{theorem}[Limit Residuum for Infinity-Choice]\label{thm:limit-residuum}%
\uses{thm:global-inf-choice,def:hypercomplete,conj:postnikov-stepwise}%

In an $(\infty, 1)$-topos $\mathcal{X}$ equipped with a Natural Numbers Object ($\mathrm{NNO}$), the categorical residuum bridging bounded truncated choice ($\forall n < \infty,\, \mathrm{AC}_n$) to full Infinity-Choice ($\mathrm{AC}_\infty$) is the conjunction of \emph{Postnikov convergence} (hypercompleteness under $0$-truncated generation) and sequential $\mathrm{DC}$ (killing $\varprojlim^1$ phantom obstructions along Postnikov towers):

\[
  \bigl(\forall n < \infty,\, \mathrm{AC}_n\bigr) + \bigl(\text{Postnikov Convergence } \wedge \text{ Sequential } \mathrm{DC}\bigr) \iff \mathrm{AC}_\infty
\]

\end{theorem}

\begin{proof}
$(\Longrightarrow)$ Assume $\mathcal{X} \models \mathrm{AC}_\infty$. Restricting $\mathrm{AC}_\infty$ to $n$-truncated effective epimorphisms over $0$-truncated targets immediately gives $\mathcal{X} \models \mathrm{AC}_n$ for every finite $n < \infty$, and restricting to $n = 0$ shows $\tau_{\le 0}\mathcal{X} \models \mathrm{AC}_0$ (external choice on $0$-truncated sets). Over the internal $\mathrm{NNO}$ $\mathbb{N}$ of $\mathcal{X}$, splitting $\coprod_{n \in \mathbb{N}} A_{n+1} \twoheadrightarrow \coprod_{n \in \mathbb{N}} A_n$ in $\tau_{\le 0}\mathcal{X}$ (together with external $\mathrm{AC}_0$ and $\mathrm{DC}$) extracts compatible threads across countable surjective towers, establishing sequential $\mathrm{DC}$. Moreover, by Theorem~\ref{thm:global-inf-choice}, $\mathrm{AC}_\infty$ guarantees that every object $W \in \mathcal{X}$ admits a $0$-truncated cover $S \twoheadrightarrow W$ with $S \in \tau_{\le 0}\mathcal{X}$. For any $\infty$-connected morphism $f : X \to Z$ in $\mathcal{X}$ (in particular the canonical Postnikov limit map $\theta_X : X \to \varprojlim_{n \in \omega} \tau_{\le n} X$), every pullback $z^* f : z^* X \to S$ along a map $z : S \to Z$ from any $S \in \tau_{\le 0}\mathcal{X}$ remains $\infty$-connected, hence $(-1)$-connected (an effective epimorphism onto $S$), and splits by $\mathrm{AC}_\infty$, proving $\pi_0$-surjectivity of $\mathrm{Map}_{\mathcal{X}}(S, f)$. Iterating this same pullback-splitting argument along the path-space pullbacks of the $\infty$-connected diagonals $\Delta_{z^* f}$ shows by induction on $k \ge 0$ that $\pi_k(\mathrm{Map}_{\mathcal{X}/S}(S, z^* X)) = 0$, so $\mathrm{Map}_{\mathcal{X}}(S, X) \xrightarrow{\sim} \mathrm{Map}_{\mathcal{X}}(S, Z)$ is a weak equivalence of $\infty$-groupoids for every $S \in \tau_{\le 0}\mathcal{X}$. Under $0$-truncated generation from Theorem~\ref{thm:global-inf-choice}, this implies Postnikov convergence $X \xrightarrow{\sim} \varprojlim_{n \in \omega} \tau_{\le n} X$ (hypercompleteness).

$(\Longleftarrow)$ Conversely, assume $\mathcal{X} \models \mathrm{AC}_n$ for all $n < \infty$, Postnikov towers converge in $\mathcal{X}$, and sequential $\mathrm{DC}$ holds. Let $e : X \twoheadrightarrow Y$ be an arbitrary effective epimorphism onto a $0$-truncated object $Y \in \tau_{\le 0}\mathcal{X}$. In the slice $(\infty, 1)$-topos $\mathcal{X}/Y$, factor $e$ through its relative Postnikov tower
\[
  \cdots \xrightarrow{p_{n+2}} \tau_{\le n+1}(X/Y) \xrightarrow{p_{n+1}} \tau_{\le n}(X/Y) \xrightarrow{p_n} \cdots \xrightarrow{p_0} \tau_{\le -1}(X/Y) \simeq Y,
\]
and set $M_n := \mathrm{Map}_{\mathcal{X}/Y}(Y, \tau_{\le n}(X/Y))$ for each $n \ge 0$ (presented as a tower of Kan fibrations of mapping spaces). By Postnikov convergence in $\mathcal{X}$ (pulled back to the slice over $Y \in \tau_{\le 0}\mathcal{X}$), $X/Y \xrightarrow{\sim} \varprojlim_{n \in \omega^{\mathrm{op}}} \tau_{\le n}(X/Y)$, giving $\mathrm{Map}_{\mathcal{X}/Y}(Y, X) \xrightarrow{\sim} \varprojlim_{n \in \omega^{\mathrm{op}}} M_n$. First, for any stage-$n$ section $s_n : Y \to \tau_{\le n}(X/Y)$ in $M_n$, pulling back $\tau_{\le n+1}(X/Y) \twoheadrightarrow \tau_{\le n}(X/Y)$ along $s_n$ produces an $(n+1)$-truncated effective epimorphism $\pi_{n+1} : E_{n+1} := Y \times_{\tau_{\le n}(X/Y)} \tau_{\le n+1}(X/Y) \twoheadrightarrow Y$ in $\mathcal{X}/Y$, which splits by $\mathrm{AC}_{n+1}$; hence $(p_{n+1})_* : \pi_0(M_{n+1}) \twoheadrightarrow \pi_0(M_n)$ is surjective for every $n \ge 0$. Applying $\mathrm{AC}_0$ ($\pi_0(M_0) \neq \emptyset$) and sequential $\mathrm{DC}$ on the surjective tower $(\pi_0(M_n))_{n \ge 0}$ (lifting compatible stages via $\mathrm{DC}$ across the fibration tower) yields a compatible basepoint sequence $(s_n)_{n \ge 0}$ with $([s_n])_{n \ge 0} \in \varprojlim_n \pi_0(M_n)$. Second, for every $n \ge 1$, the homotopy fiber $\mathrm{fib}_{s_n}(M_{n+1} \to M_n) \simeq \mathrm{Map}_{\mathcal{X}/Y}(Y, E_{n+1}) \simeq \mathrm{Map}_{\mathcal{X}/Y}(Y, K(\Pi_{n+1}, n+1))$ has lower homotopy groups $\pi_k\bigl(\mathrm{Map}_{\mathcal{X}/Y}(Y, K(\Pi_{n+1}, n+1))\bigr) \cong H^{n+1-k}(\mathcal{X}/Y, \Pi_{n+1}) = 0$ for all $0 \le k \le n$ because $\mathrm{AC}_{n-k}$ splits every $(n-k)$-gerbe banded by $\Pi_{n+1}$. Thus every transition fiber for $n \ge 1$ is $n$-connected, so the long exact homotopy sequence yields $\pi_1(M_{n+1}, s_{n+1}) \xrightarrow{\cong} \pi_1(M_n, s_n)$ for all $n \ge 1$, ensuring the Mittag--Leffler condition and killing $\varprojlim^1_m \pi_1(M_m, s_m) = *$. By the Milnor exact sequence for the homotopy limit tower $(M_m)$,
\[
  * \longrightarrow \varprojlim^1_m \pi_1(M_m, s_m) \longrightarrow \pi_0\Bigl(\varprojlim_m M_m\Bigr) \longrightarrow \varprojlim_m \pi_0(M_m) \longrightarrow *,
\]
the compatible $\pi_0$-thread lifts to a global section $s \in \pi_0\,\mathrm{Map}_{\mathcal{X}/Y}(Y, X)$, proving $\mathcal{X} \models \mathrm{AC}_\infty$.
\end{proof}

\section{Methodological Bridge: Categorical Reverse Mathematics as a Discovery Engine in \texorpdfstring{$(\infty,\infty)$}{(infinity,infinity)}-Category Theory}\label{subsec:crm-discovery-engine-inf-inf}

Passing along Axis~1 from $(\infty, 1)$-categories to directed $(\infty, \infty)$-categories while varying the ambient choice or topos base along Axis~2 turns Categorical Reverse Mathematics into an active \emph{discovery engine}: rather than only classifying pre-existing classical theorems, extracting the exact residuum $\mathcal{R}$ validating $A + \mathcal{R} \iff B$ systematically isolates hidden geometric hypotheses, prevents false equivalences, and engineers sharp separating countermodels across the $(\infty, \infty)$-landscape. Four concrete $(\infty, \infty)$-categorical arenas (Table~\ref{tab:crm-inf-inf-residuums}) demonstrate how this residuum calculus drove the breakthroughs verified across Chapters~\ref{ch:higher-crm}--\ref{ch:realizability-sites} and the \texttt{HomotopyCRM} Lean~4 suite.

\begin{table}[htbp]
\centering
\scriptsize
\setlength{\tabcolsep}{2.8pt}
\renewcommand{\arraystretch}{1.12}
\begin{tabularx}{\textwidth}{@{}P{0.19\textwidth}P{0.18\textwidth}P{0.19\textwidth}Y@{}}
\toprule
\textbf{$(\infty, \infty)$-Discovery Arena} & \textbf{Weak Principle $A$} & \textbf{Strong Theorem $B$} & \textbf{Exact CRM Residuum $\mathcal{R}$ \& Countermodel} \\
\midrule
(1) Dependent Choice \& \texttt{ICaTT} Stream Adequacy & Relational coinduction $E_{\mathrm{rel}} = \nu \Phi$ & Witness-stream image $E_{\mathrm{rel}} = \mathrm{im}(\mathrm{InvStr})$ & Internal Dependent Choice $\mathrm{IDC}_\omega$ (Theorem~\ref{thm:coinductive-stream-choice}; fails at Rung~0 on $\mathrm{Free}_{\mathbf{CaTT}}(E[\mathbb{A}]) \in \mathbf{Clan}(\mathcal{S}_{\mathbf{CaTT}}, \mathbf{Nom})$) \\
(2) Pure Postnikov Phantoms \& Conjecture~C1 & Finite hom-ranks $\mathcal{C} \in \mathrm{Cat}_{\le \omega}$ ($\alpha = 0$ in $\mathrm{Cat}^\omega_{(\infty,\infty)}$) & Oriented Postnikov convergence ($\mathcal{C} \in \mathrm{Cat}^L_{(\infty,\infty)}$) & Absence of absorbing-chain Postnikov phantoms / $1$-acyclicity of $\preceq_1$ (Proposition~\ref{prop:cat-le-omega-lfd}, Theorems~\ref{thm:ordinal-rank-completeness} and~\ref{thm:classifier-scc-trapping}--\ref{thm:step-product-counterexample}; refuted by gaunt Noetherian rank-$\omega$ step-product $\mathcal{Q} \subseteq \prod_{k \ge 1} \mathcal{A}_k$) \\
(3) Directed Delooping Tameness vs.\ $\pi_{\le 0}$ Collapse & Pointed directed delooping $\mathbf{B}^\infty(\mathbf{FinSet}^{\simeq}) \in \mathrm{CatSp}^L$, $\mathbf{B}^\infty(M_4^\#) \in \mathrm{CatSp}^R$ & Classical sphere spectrum $\mathbb{S}_{\mathrm{top}}$ / Wall stable diffeomorphism collapse & Left Bousfield groupoidification $\pi_{\le 0}$ triggering group-completion $\mathbb{N} \to \mathbb{Z}$ (Barratt--Priddy--Quillen, avoiding Giraud collapse; square-zero monoids $M(X)$ transport strictness across $\mathrm{CatSp}^L \subsetneq \mathrm{CatSp}^\omega \subsetneq \mathrm{CatSp}^{\mathrm{coind}} \subsetneq \mathrm{CatSp}^R$) \\
(4) Oriented Convexity \& Cubical $n \ge 3$ Wall & Binary ($n = 2$) Strong Amalgamation (\texttt{sap\_\allowbreak flowConvex}) & All-$n$ flow-acyclic cubical $n$-amalgamation ($n \ge 3$) & Two-clause \texttt{IsFlowConvex} (\texttt{cell\_\allowbreak mem} + \texttt{inter\_\allowbreak mem}) \& 1-sided cobase-change sufficiency (keeping Nicolas's pinched $2$-cube $P, P^{\mathrm{op}}, P^{\mathrm{co}}$ strictly \texttt{IsFlowAcyclic}, vs.\ $n \ge 3$ directed $n$-cycle $F_n$, Theorem~\ref{thm:directed-cubical-amalgamation}) \\
\bottomrule
\end{tabularx}
\caption{Methodological Residuum Bridges in $(\infty, \infty)$-Category Theory.}\label{tab:crm-inf-inf-residuums}
\end{table}

Arenas~(1), (2), and~(4) of Table~\ref{tab:crm-inf-inf-residuums} receive full formal proofs in their home chapters: Section~\ref{subsec:ordinal-rank-icatt-perfection} below establishes ordinal-rank $\alpha = 0$ stabilization (Theorem~\ref{thm:ordinal-rank-completeness}), $[1]$-classifier SCC trapping and $1$-acyclic completeness (Theorems~\ref{thm:classifier-scc-trapping}--\ref{thm:ff-acyclic-completeness}), and the step-product counterexample $\mathcal{Q}$ disproving Conjecture~C1 (Theorem~\ref{thm:step-product-counterexample}), while Chapter~\ref{ch:realizability-sites} proves the cubical $n = 2$ vs.\ $n \ge 3$ flow-convexity dichotomy (Theorem~\ref{thm:directed-cubical-amalgamation}, keeping Marcus Nicolas's pinched $2$-cube $P = \square^2 \cup_{\{10, 01\}} \square^1$, $P^{\mathrm{op}}$, and $P^{\mathrm{co}}$ strictly \texttt{IsFlowAcyclic} despite Masuda's v2 Steiner partial-duality erratum~\cite[\S 2.4, Footnote~5]{Masuda26}) and the $\mathrm{IDC}_\omega$ calibration of internal $\mathbf{ICaTT}$ coinductive witness streams on $\mathrm{Free}_{\mathbf{CaTT}}(E[\mathbb{A}]) \in \mathbf{Clan}(\mathcal{S}_{\mathbf{CaTT}}, \mathbf{Nom})$ (Theorem~\ref{thm:coinductive-stream-choice}). Arena~(3) synthesizes the homotopical collapse residuum directly on categorical spectra:

\subsection*{Directed Delooping Tameness versus \texorpdfstring{$\pi_{\le 0}$}{pi\_<=0} Group-Completion and Smooth \texorpdfstring{$4$}{4}-Manifold Stabilization}

Enforcing Guardrail~1---Giraud strict initiality ($\emptyset \times X \simeq \emptyset$)~\cite[Thm.~6.1.0.6, \S 6.1.1]{Lurie09} prevents any non-trivial ``stable $(\infty,1)$-topos,'' and $\Sigma^\infty_+(*) \simeq \mathbb{S}$ in $\mathrm{Pr}^L$~\cite[\S\S 1.4.4, 4.8.2]{Lurie17} (never $\Sigma^\infty_+(U) \simeq \mathbb{S}$ for a general object $U$)---forces directed stabilization to pass to pointed $(\infty,\infty)$-categories $\mathrm{Cat}^R_* := (\mathrm{Cat}^R_{(\infty,\infty)})_{*/}$, forming $\mathrm{CatSp}^R := \mathrm{Sp}(\mathrm{Cat}^R_*)$~\cite{Masuda26,Heine26}. Because directed looping $\Omega_{\mathrm{dir}}(B\mathcal{C}) \simeq \mathcal{C}$ commutes with $\pi_{\le n}$ for all $n \ge 1$ ($\Omega_{\mathrm{dir}}(\pi_{\le n}(X, x)) \xrightarrow{\sim} \pi_{\le n-1}(\Omega_{\mathrm{dir}}(X, x))$) while avoiding $n = 0$ group-completion ($\Omega(\pi_{\le 0} B\mathbb{N}) \simeq \mathbb{Z} \not\simeq \mathbb{N} = \pi_{\le 0}(\Omega_{\mathrm{dir}} B\mathbb{N})$, cf.\ Chapter~\ref{ch:panorama}, \S\ref{subsec:excision-catsp} and~\cite{Heine26,Geiss26}), CRM identifies Left Bousfield groupoidification $\pi_{\le 0}$ itself as the exact residuum generating classical stable topological complexity (Chapter~\ref{ch:panorama}, \S\ref{subsec:excision-catsp}): the categorical delooping $\mathbf{B}^\infty(\mathbf{FinSet}^{\simeq}) \in \mathrm{CatSp}^L$ remains $(k+1, k)$-truncated and Noetherian of rank $k + 2$ at every level $k$ (since only $\emptyset$ is $\amalg$-invertible), until levelwise $\pi_{\le 0}$ forces group-completion $\mathbb{N} \to \mathbb{Z}$ and Quillen's plus-construction $B\Sigma_\infty \to B\Sigma_\infty^+$ (Barratt--Priddy--Quillen~\cite{BarrattPriddy72,Quillen71}) to birth the classical topological sphere spectrum $\mathbb{S}_{\mathrm{top}}$ and all $\pi_*^s(\mathbb{S}_{\mathrm{top}})$. By the exact same mechanism, $\mathbf{B}^\infty(M_4^\#) \in \mathrm{CatSp}^R$ under connected sum $\#$ preserves smooth $4$-manifold non-invertibility (avoiding Wall's stable diffeomorphism collapse triggered by $\pi_{\le 0}$), while square-zero absorbing monoids $M(X) := *_0 \sqcup X \sqcup *_2$ (where $(B^k M(X))^{\simeq} \simeq *$ for all $k \ge 1$, killing Rezk gluing and $\varprojlim^1$ twists) transport strictness across the lifted four-tier chain $\mathrm{CatSp}^L \subsetneq \mathrm{CatSp}^\omega \subsetneq \mathrm{CatSp}^{\mathrm{coind}} \subsetneq \mathrm{CatSp}^R$~\cite{ORW26,GepnerHeine26}.

\section{\texorpdfstring{Inductive versus Coinductive Equivalences in Directed $(\infty, \infty)$-Categories}{Inductive versus Coinductive Equivalences in Directed (infinity, infinity)-Categories}}\label{subsec:directed-inf-inf}

When we further drop $1$-invertibility along Axis~1 (Chapter~\ref{ch:base-fiber}, Section~\ref{sec:axis1-deformations}) to pass from $(\infty, 1)$-categories to $(\infty, \infty)$-categories, the Whitehead hypercompleteness gap acquires a directed counterpart---together with a foundational distinction between inductive initial-algebra fixed points and coinductive reflective localizations.

On the inductive side, Zach Goldthorpe~\cite{Goldthorpe23}, via an $\infty$-categorical generalization of Ad\'amek's free endofunctor algebra\index{Adamek initial algebra@Ad\'amek initial algebra} construction for weak Gepner--Haugseng enrichment $\mathrm{Enr}: \mathcal{V} \mapsto \mathcal{V}\mathrm{Cat}$~\cite{GepnerHaugseng15}, proved that over presentably symmetric monoidal $(\infty, 1)$-categories $\mathrm{SymMon}^{\mathrm{Pres}}_\infty$, $(\infty, \infty)$-categories equipped with \emph{inductively defined equivalences} (generated bottom-up from identity morphisms, modeled by the right-truncation limit $\mathrm{Cat}^R_{(\infty, \infty)} := \varprojlim_\kappa \mathrm{Cat}_{(\infty, n)}$, \cite[Theorem~3.3.5]{Goldthorpe23}, with its Postnikov-complete / left-generated subcategory $\mathrm{Cat}^L_{(\infty, \infty)} := \varprojlim_\pi \mathrm{Cat}_{(\infty, n)}$ embedding reflectively inside $\mathrm{Cat}^R_{(\infty, \infty)}$) form the initial algebra of $\mathrm{Enr}$, while over $\mathrm{SymMon}_\infty$, Ad\'amek iteration across the large ordinal $\Lambda$ stabilizes on Noetherian $(\infty, \infty)$-categories $\mathrm{Cat}^{\mathrm{Noeth}}_{(\infty, \infty)}$ lacking infinite non-trivial parallel morphism towers ($\mathrm{Cell}_\infty(\mathcal{C}) \simeq \emptyset$). By contrast, an earlier dual claim of Goldthorpe~\cite{Goldthorpe23} asserting that $\mathrm{Cat}_\omega$ yields a terminal coalgebra of $\mathrm{Enr}$ on $\mathrm{SymMon}_\infty$ was withdrawn in revision per Simon Henry's erratum\index{Henry erratum (Goldthorpe)}: the groupoid-reflection left adjoint $\pi_{\le 0}(\mathcal{C}) = |\mathcal{C}| \not\simeq \mathcal{C}^{\simeq}$ fails to agree with the maximal Kan groupoid core (whereas Theorem~3.2.6 of~\cite{Goldthorpe23} (\href{https://arxiv.org/abs/2307.00442v5}{\texttt{arXiv:2307.00442v5}}) is the surviving weakly-contractible limit-preservation theorem supporting Theorem~3.3.5, and Simon Henry's counterexample excludes naive left-localization limits $\varprojlim_\pi \mathrm{Cat}_{(\infty, n)}$ from coinciding with $\mathrm{Cat}^R_{(\infty, \infty)}$).

Instead, inside the ambient right-truncation tower limit $\mathrm{Cat}^R_{(\infty, \infty)} := \varprojlim_\kappa \mathrm{Cat}_{(\infty, n)}$ taken along the dimensional core/truncation right adjoints $\kappa_{\le n}$, the verified landscape of $(\infty, \infty)$-categorical homotopy theories organizes into a strict four-tier chain of reflective localizations\index{Four-tier localization ladder} established by Ozornova--Rovelli--Walde~\cite[Thm.~4.25, \S 5]{ORW26} and Gepner--Heine~\cite{GepnerHeine26} (illustrated alongside the step-product phantom in Figure~\ref{fig:inf-inf-ladder-phantom}):

\[
  \mathrm{Cat}^L_{(\infty, \infty)} \subsetneq \mathrm{Cat}^\omega_{(\infty, \infty)} \subsetneq \mathrm{Cat}^{\mathrm{coind}}_{(\infty, \infty)} \subsetneq \mathrm{Cat}^R_{(\infty, \infty)}
\]

\begin{itemize}[leftmargin=1.5em]

  \item \textbf{Inductive / Postnikov-Complete Tier ($\mathrm{Cat}^L_{(\infty, \infty)}$):} Left-generated $(\infty, \infty)$-categories ($\mathcal{C} \xrightarrow{\sim} RL(\mathcal{C}) \simeq \varprojlim_{k \in \omega^{\mathrm{op}}} i_k \pi_{\le k}(\mathcal{C})$) equipped with inductively generated equivalences built bottom-up from finite-dimensional truncations and identity morphisms.

  \item \textbf{Street--Verity $\omega$-Tier ($\mathrm{Cat}^\omega_{(\infty, \infty)}$):} Complicial / $\omega$-categorical models whose $\omega$-equivalences sit strictly between finite-dimensional Postnikov convergence and full coinductive saturation.

  \item \textbf{Coinductive Tier ($\mathrm{Cat}^{\mathrm{coind}}_{(\infty, \infty)}$):} $(\infty, \infty)$-categories localized with respect to \emph{coinductive equivalences} (morphisms admitting two-sided inverses up to coinductive equivalence at every level).

  \item \textbf{Right-Truncation Limit ($\mathrm{Cat}^R_{(\infty, \infty)} := \varprojlim_\kappa \mathrm{Cat}_{(\infty, n)}$):}\index{CatR@$\mathrm{Cat}^R_{(\infty,\infty)}$} The ambient projective limit along right-adjoint core truncations $\kappa_{\le n}$, on which the composite reflective localization endofunctor $RL : \mathrm{Cat}^R_{(\infty, \infty)} \to \mathrm{Cat}^R_{(\infty, \infty)}$ commutes with categorical suspension $\Sigma$ and set-indexed coproducts~\cite{ORW26,GepnerHeine26}.

\end{itemize}

Within this four-tier reflective hierarchy, two verified placement results calibrate the exact scope of Noetherian tower termination relative to inductive and coinductive hypercompleteness:

\begin{itemize}[leftmargin=1.5em]

  \item \textbf{Noetherian Coinductive Hypercompleteness ($\mathrm{Cat}^{\mathrm{Noeth}}_{(\infty, \infty)} \subseteq \mathrm{Cat}^{\mathrm{coind}}_{(\infty, \infty)}$):} Vanishing of infinite non-trivial parallel cell towers ($\mathrm{Cell}_\infty(\mathcal{C}) \simeq \emptyset$) implies coinductive hypercompleteness, embedding $\mathrm{Cat}^{\mathrm{Noeth}}_{(\infty, \infty)}$ inside $\mathrm{Cat}^{\mathrm{coind}}_{(\infty, \infty)}$~\cite{Goldthorpe23,ORW26,GepnerHeine26}.

  \item \textbf{Strict Separation by $S^2$ ($S^2 \in \mathrm{Cat}^L_{(\infty, \infty)} \setminus \mathrm{Cat}^{\mathrm{Noeth}}_{(\infty, \infty)}$):} Viewing the $2$-sphere $S^2$ as an $\infty$-groupoid inside $\mathrm{Cat}^L_{(\infty, \infty)}$, its infinitely many non-trivial higher homotopy groups $\pi_n(S^2)$ witness $\mathrm{Cell}_\infty(S^2) \not\simeq \emptyset$, strictly separating inductive $(\infty, \infty)$-categories from Noetherian ones~\cite{ORW26,GepnerHeine26}.

\end{itemize}

In an arbitrary $(\infty, \infty)$-category $\mathcal{C}$, every inductive equivalence is coinductively invertible, but the converse fails outside hypercomplete regimes: for example, the higher cobordism $(\infty, \infty)$-category $\mathrm{Cob}_\infty$ has duals for morphisms in all dimensions---forcing every morphism to be a coinductive equivalence and collapsing $\mathrm{Cob}_\infty$ to an $\infty$-groupoid under coinductive saturation in $\mathrm{Cat}^{\mathrm{coind}}_{(\infty, \infty)}$---while remaining a non-trivial directed $(\infty, \infty)$-category under inductive equivalences in $\mathrm{Cat}^R_{(\infty, \infty)}$~\cite{ORW26,GepnerHeine26} (and separating $\mathbf{B}^\infty \mathrm{Cob}_\infty \in \mathrm{CatSp}^R \setminus \mathrm{CatSp}^{\mathrm{coind}}$~\cite{Masuda26,Heine26}). Consequently, bridging principles valid for coinductive equivalences ($A$, in $\mathrm{Cat}^{\mathrm{coind}}_{(\infty, \infty)}$) to strict inductive / Postnikov-convergent invertibility ($B$, in $\mathrm{Cat}^L_{(\infty, \infty)}$) inside $\mathrm{Cat}^R_{(\infty, \infty)}$ isolates \emph{coinductive hypercompleteness} across the localization tower (with $\mathrm{Cell}_\infty(\mathcal{C}) \simeq \emptyset$ on $\mathrm{Cat}^{\mathrm{Noeth}}_{(\infty, \infty)} \subseteq \mathrm{Cat}^{\mathrm{coind}}_{(\infty, \infty)}$ furnishing a sufficient Noetherian criterion, though not necessary in light of $S^2 \in \mathrm{Cat}^L_{(\infty, \infty)} \setminus \mathrm{Cat}^{\mathrm{Noeth}}_{(\infty, \infty)}$) as the exact directed Whitehead residuum.

\section{\texorpdfstring{Ordinal-Rank Completeness, Negative Resolution of Conjecture~C1, \texttt{ICaTT} Clan Semantics, and Slice-Algebra Perfection}{Ordinal-Rank Completeness, Negative Resolution of Conjecture C1, ICaTT Clan Semantics, and Slice-Algebra Perfection}}\label{subsec:ordinal-rank-icatt-perfection}

Finer calibrations sharpen the geometric mechanisms separating Tiers~1, 2, and~3 of the reflective localization tower, resolve \textbf{Conjecture~C1} negatively via an explicit gaunt Noetherian rank-$\omega$ step-product phantom, and anchor the exact boundary of Ad\'amek initial-algebra fixed points.

\begin{proposition}[Inductive Lali Split-Idempotent Ceiling and Cartesian-vs.-Gray Obstruction]\label{prop:lali-ceiling}%
\lean{HomotopyCRM.prop11A_lali_stabilization_above_ceiling, HomotopyCRM.prop11A_ceiling_sharp_at_max_add_two}%
\leanok%
\uses{def:comparison-triple}%
Let $l : \mathcal{A} \rightleftarrows \mathcal{D} : r$ be a Gray-enriched (oriented~\cite[Definition~3.5.1]{GepnerHeine26}) left-adjoint-left-inverse (\emph{Lali})\index{Lali adjunction} adjunction with idempotent endofunctor $E := l \circ r : \mathcal{D} \to \mathcal{D}$, over objects $x, y \in \mathcal{D}$ possessing finite local hom cell-stabilization heights $K_x = \mathrm{height}(\mathrm{Hom}_{\mathcal{D}}(x, y)) < \omega$ and $K_y = \mathrm{height}(\mathrm{Hom}_{\mathcal{D}}(E x, E y)) < \omega$ (\texttt{HasHomCellHeight}).
\begin{enumerate}[label=(\arabic*),leftmargin=2em]
  \item \textbf{Finite-Rank Split-Idempotent Stabilization:} For every pair of parallel $(m-1)$-cell boundaries $(\Phi, \Psi)$ over $\mathrm{Hom}_{\mathcal{D}}(x, y)$ (\texttt{HomBoundary}) in ambient dimension $m > \max(K_x, K_y) + 2$, the Lali idempotent $E$ induces equivalences on ambient $m$-cell animae
  \[
    \mathrm{Hom}_{\mathcal{D}}(\Phi, \Psi) \xrightarrow{\;\sim\;} \mathrm{Hom}_{\mathcal{D}}(E(\Phi), E(\Psi)) \simeq *,
  \]
  and this ceiling is sharp at $m = \max(K_x, K_y) + 2$ on the oriented join $\{a\} \star \Sigma^{K+1}(S^0)$ (owing to the $(+1) + (+1) = +2$ shift from global mapping suspension plus lax counit modification coherence). Both the stabilization bound and its exact $+2$ boundary are machine-checked in Lean~4 (\texttt{HomotopyCRM/\allowbreak OrdinalRank\allowbreak Completeness.lean}, 1,774 lines, 4-Layer Gate verified) via \texttt{prop11A\_\allowbreak lali\_\allowbreak stabilization\_\allowbreak above\_\allowbreak ceiling} and \texttt{prop11A\_\allowbreak ceiling\_\allowbreak sharp\_\allowbreak at\_\allowbreak max\_\allowbreak add\_\allowbreak two} on \texttt{JoinSuspState.joinCone K}.
  \item \textbf{Cartesian-vs.-Gray Naturality Obstruction ($n \ge 2$):} On finite-stage walking $n$-equivalences $E^{(n)}$ ($n \ge 2$), Cartesian functor categories $\mathrm{Fun}^{\times}(E^{(n)}, E^{(n)})$ cannot contract $E^{(n)}$ ($\mathrm{id}_{E^{(2)}} \not\simeq \mathrm{const}_*$ in $\mathrm{Fun}^{\times}(E^{(2)}, E^{(2)})$, blocking $E^{(\omega)}$ from Cartesian contractibility) because Cartesian cylinder cells satisfy $\dim(\mathbf{D}^{n-1} \times \mathbf{D}^1) = \max(n-1, 1) = n - 1 < n$, whereas lax/oplax Gray functor categories $\mathrm{Fun}^{\mathrm{lax}}(E^{(n)}, E^{(n)})$ supply the additive prism cell $\mathbf{D}^{n-1} \boxtimes \mathbf{D}^1$ of dimension $(n - 1) + 1 = n$~\cite{GepnerHeine26,Heine26,Geiss26}.
\end{enumerate}
\end{proposition}

Why do both Henry--Loubaton's~\cite{HenryLoubaton25} wide pushout $E^{(\omega)}$ (separating $\mathrm{Cat}^\omega_{(\infty, \infty)} \subsetneq \mathrm{Cat}^{\mathrm{coind}}_{(\infty, \infty)}$ in~\cite[Corollary~5.56]{ORW26}) and Ozornova--Rovelli--Walde's full subcategory $W \hookrightarrow V(\infty) := * \amalg_{\coprod_{n \in \mathbb{N}} \{1\}} \coprod_{n \in \mathbb{N}} E^{(n+1)}$ (separating $\mathrm{Cat}^L_{(\infty, \infty)} \subsetneq \mathrm{Cat}^\omega_{(\infty, \infty)}$ in~\cite[\S 5]{ORW26}) evade Proposition~\ref{prop:lali-ceiling} ($\sup K_x = \infty$) and exit $\mathrm{Cat}^{\mathrm{Noeth}}_{(\infty, \infty)}$? Two independent mechanisms force $\mathrm{Cell}_\infty \not\simeq \emptyset$:
\begin{itemize}[leftmargin=1.5em]
  \item \textbf{$\mathbb{E}_3$ Configuration-Space Spheres in Weak Cellular Pushouts:} Attaching weak cells at codimension $k \ge 3$ over loop boundaries activates Eckmann--Hilton / syllepsis braiding on $\Omega^3_x$, whose $2$-strand component has underlying anima $\mathrm{Conf}_2(\mathbb{R}^3)_{/\Sigma_2} \simeq \mathbb{R}\mathbb{P}^2$ with $\pi_m(\mathbb{R}\mathbb{P}^2) \cong \pi_m(S^2) \neq 0$ for infinitely many $m \ge 2$~\cite{HenryLoubaton25,ORW26}.
  \item \textbf{Hub Excursions in $W \subseteq V(\infty)$:} Even on skeletal cell heights, each wedge summand $E^{(n+1)}$ identifies $0 \mapsto n \in \mathrm{Ob}(W)$ and $1 \mapsto *$, contributing $\eta_n : n \to *$ and candidate inverses $v_n : * \to n$. Already at object $1 \in \mathrm{Ob}(W)$, the hub excursion $f_1 := v_1 \circ \eta_1 : 1 \to * \to 1$ and iterated loops $f_{1, n} := v_1 \circ (\eta_n \circ v_n) \circ \eta_1 \in W(1, 1)$ embed unbounded $(n+1)$-cell towers into the single hom-object $W(1, 1)$, forcing $\mathrm{rank}(W(1, 1)) \ge \omega$ and $\mathrm{rank}(W) \ge \omega + 1$~\cite[\S 5]{ORW26}.
\end{itemize}

Define the recursive ordinal hom-rank\index{Ordinal hom-rank@Ordinal hom-rank ($\mathrm{Cat}_{\le\omega}$)} by
\[
  \mathrm{rank}(\mathcal{D}) \;:=\; \sup_{x, y \in \mathrm{Ob}(\mathcal{D})} \bigl(\mathrm{rank}(\mathcal{D}(x, y)) + 1\bigr),
\]
and let $\mathrm{Cat}_{\le \omega}$ denote the class of $(\infty, \infty)$-categories whose hom-objects have finite ordinal rank ($\mathrm{rank}(\mathcal{C}(x, y)) < \omega$ for all $x, y$, blocking both $\mathbb{R}\mathbb{P}^2$ homotopy towers and hub excursions). Calibrating Ozornova--Rovelli--Walde's~\cite[Questions~5.60 and~5.62 (v1 Question~5.58)]{ORW26} transfinite reflector $L_\omega : \mathrm{Cat}^R_{(\infty, \infty)} \to \mathrm{Cat}^\omega_{(\infty, \infty)}$ built via $\mathcal{C}_{\alpha+1} :=\allowbreak \mathcal{C}_\alpha[\mathrm{IsOmegaInv}\allowbreak(\mathcal{C}_\alpha)^{-1}]$ yields complete stabilization at $\alpha = 0$ on $\mathrm{Cat}_{\le \omega}$---machine-checked in \texttt{HomotopyCRM/\allowbreak Ordinal\allowbreak Rank\allowbreak Completeness.lean} (1,774 lines, 4-Layer Gate verified) via universe-safe globular state coalgebras \texttt{Globular\allowbreak Inf\allowbreak System} and transfinite ordinal invertibility \texttt{Is\allowbreak Ordinal\allowbreak Inv}, together with \texttt{Omega\allowbreak Tower\allowbreak State} separating ranks $0 < 1 < 2 < \omega < \omega + 1$:

\begin{theorem}[Ordinal-Rank Completeness and Stabilization of ORW Question~5.58 at $\alpha = 0$]\label{thm:ordinal-rank-completeness}%
\lean{HomotopyCRM.isOmegaInv_iff_isOrdinalInv_omega0, HomotopyCRM.theorem9A_ordinal_rank_completeness, HomotopyCRM.ranked_coindEquiv_iff_isEquiv, HomotopyCRM.theorem9A_omega_completeness, HomotopyCRM.orw_question558_stabilizes_at_zero, HomotopyCRM.apexArrow_isOmegaInv, HomotopyCRM.apexArrow_not_isEquiv, HomotopyCRM.apexArrow_not_isCoindEquiv}%
\leanok%
\uses{prop:lali-ceiling}%
For every ordinal $\theta$ and every ranked $(\infty, \infty)$-category $\mathcal{D}$ satisfying $\mathrm{rank}(\mathcal{D}) \le \theta$, transfinite $\theta$-inductive invertibility coincides with genuine equivalence, $\mathrm{IsOrdinalInv}_\theta(\mathcal{D}) = \mathrm{IsEquiv}(\mathcal{D})$ (\texttt{theorem9A\_\allowbreak ordinal\_\allowbreak rank\_\allowbreak completeness}, sharp at rank $\omega + 1$ via the additive non-surjective inclusion $\mathrm{Vect} \to \mathrm{Vect}^2$, $V \mapsto (V, 0)$, and \texttt{OmegaTowerState.apex}). Specializing to $\theta = \omega$, every $\mathcal{C} \in \mathrm{Cat}_{\le \omega}$ satisfies $\mathrm{IsOmegaInv}(\mathcal{C}) = \mathrm{IsEquiv}(\mathcal{C})$ (\texttt{theorem9A\_\allowbreak omega\_\allowbreak completeness}), so v2 Questions~5.60 and~5.62 (v1 Question~5.58) stabilize immediately at $\alpha = 0$ (\texttt{orw\_\allowbreak question558\_\allowbreak stabilizes\_\allowbreak at\_\allowbreak zero}, giving $\mathrm{Cat}_{\le \omega} \subseteq \mathrm{Cat}^\omega_{(\infty, \infty)}$) and narrows \textbf{Conjecture~C1}\index{Conjecture C1 (Phantom-Freeness)} strictly to \emph{Postnikov Phantom-Freeness}:
\begin{gather*}
  \text{\textbf{Conjecture~C1 (Phantom-Freeness):}} \qquad \mathrm{Cat}_{\le \omega} \;\stackrel{?}{\subseteq}\; \mathrm{Cat}^L_{(\infty, \infty)} \\
  \text{[Resolved negatively in Theorem~\ref{thm:step-product-counterexample}].}
\end{gather*}
\end{theorem}

\begin{proposition}[Local Finite-Dimensionality of $\mathrm{Cat}_{\le \omega}$ and Failure of $\omega$-Filtered Colimit Closure of $\mathrm{Cat}^L_{(\infty, \infty)}$]\label{prop:cat-le-omega-lfd}
An $(\infty, \infty)$-category $\mathcal{C} \in \mathrm{Cat}^R_{(\infty, \infty)}$ lies in $\mathrm{Cat}_{\le \omega}$ ($\mathrm{rk}_S(\mathcal{C}) \le \omega$) if and only if $\mathcal{C}$ is \emph{locally finite-dimensional}: for every pair of objects $x, y \in \mathrm{Ob}(\mathcal{C})$, there exists a finite integer $d(x, y) < \omega$ such that $\mathrm{Hom}_{\mathcal{C}}(x, y) \in \mathrm{Cat}_{(\infty, d(x, y))}$. Moreover, while Ozornova--Rovelli--Walde's~\cite[\S 5]{ORW26} category $W = \varinjlim_d W(d)$ shows that $\mathrm{Cat}^L_{(\infty, \infty)}$ fails $\omega$-filtered colimit closure inside $\mathrm{Cat}^R_{(\infty, \infty)}$, non-fully-faithful bonding maps $W(d) \to W(d+1)$ inject unbounded hub excursions $f_{1, n} \in W(1, 1)$, forcing $\mathrm{rk}_S(W) \ge \omega + 1$; consequently $W \notin \mathrm{Cat}_{\le \omega}$ does not test Conjecture~C1.
\end{proposition}

\begin{proof}
By definition of recursive hom-rank, $\mathrm{rk}_S(\mathcal{C}) \le \omega$ holds if and only if $\mathrm{rk}_S(\mathcal{C}(x, y)) < \omega$ for all $x, y \in \mathrm{Ob}(\mathcal{C})$, which holds precisely when iterated hom-towers over $\mathcal{C}(x, y)$ have contractible cell animae above some finite depth $d(x, y) < \omega$, i.e., $\mathrm{Hom}_{\mathcal{C}}(x, y) \in \mathrm{Cat}_{(\infty, d(x, y))}$. In \S 5 of~\cite{ORW26}, each finite stage $W(d) \subseteq V(d) \in \mathrm{Cat}_{(\infty, d+1)} \subseteq \mathrm{Cat}^L_{(\infty, \infty)}$ lies in $\mathrm{Cat}^L_{(\infty, \infty)}$, yet $W = \varinjlim_d W(d) \in \mathrm{Cat}^\omega_{(\infty, \infty)} \setminus \mathrm{Cat}^L_{(\infty, \infty)}$. Because the iterated hub loops $f_{1, n} = v_1 \circ (\eta_n \circ v_n) \circ \eta_1$ populate $W(1, 1)$ at unbounded heights $n + 1$, $\mathrm{rk}_S(W(1, 1)) \ge \omega$ and $\mathrm{rk}_S(W) \ge \omega + 1$, excluding $W$ from $\mathrm{Cat}_{\le \omega}$.
\end{proof}

\begin{theorem}[{$[1]$}-Classifier SCC Trapping and Grading Rigidity]\label{thm:classifier-scc-trapping}%
\lean{HomotopyCRM.theorem12B_classifier_scc_trapping, HomotopyCRM.theorem12B_tower_trapped_in_single_scc}%
\leanok%
\uses{thm:ordinal-rank-completeness}%
Let $\mathcal{C} \in \mathrm{Cat}^R_{(\infty, \infty)}$, and define $1$-cell reachability $a \preceq_1 b$ on $\mathrm{Ob}(\mathcal{C})$ by $\mathrm{Hom}_{\mathcal{C}}(a, b) \neq \emptyset$.
\begin{enumerate}[label=(\arabic*),leftmargin=2em]
  \item \textbf{Reachability SCC Trapping:} For every Postnikov-compatible tower $X_\bullet = (x_n)_{n \ge 1} \in RL(\mathcal{C})$ (satisfying $\ell_n(x_{n+1}) \simeq \ell_n(x_n)$ for all $n \ge 1$), all stage representatives $x_n$ are trapped inside a single strongly connected component (SCC) of $(\mathrm{Ob}(\mathcal{C}), \preceq_1)$.
  \item \textbf{Grading Rigidity:} On any infinite $1$-SCC $\mathcal{S}$ lacking strict absorbing/zero-like morphisms under composition, local finite hom-dimensions $d(a, b) < \omega$ cannot grow unboundedly compatibly with $1$-cell composition whiskering under $\pi_{\le M(a,b)}$ or coreflective $\kappa_{\le M(a,b)}$; hence any counterexample in $\mathrm{Cat}_{\le \omega} \setminus \mathrm{Cat}^L_{(\infty, \infty)}$ must decouple an infinite chain of objects via an absorbing hierarchy rather than an infinite non-absorbing $1$-SCC or homotopy spheres.
\end{enumerate}
\end{theorem}

\begin{proof}
For each $x \in \mathrm{Ob}(\mathcal{C})$, the principal upper set $U_x := \{z \in \mathrm{Ob}(\mathcal{C}) \mid x \preceq_1 z\}$ induces a well-defined functor $\chi_x : \mathcal{C} \to [1] = \{0 < 1\}$ sending $U_x \mapsto 1$ and $\mathcal{C} \setminus U_x \mapsto 0$, because $[1] \in \mathrm{Cat}_{(0,1)}$ has contractible non-empty homs and $\mathrm{Hom}_{[1]}(1, 0) = \emptyset$. For every $n \ge 1$, the universal property of $\ell_n : \mathcal{C} \to \pi_{\le n}(\mathcal{C})$ factors $\chi_x \simeq \overline{\chi}_{x,n} \circ \ell_n$. Since $[1]$ is gaunt ($0 \not\simeq 1$),
\[
  \ell_n(x) \simeq \ell_n(y) \;\Longrightarrow\; \overline{\chi}_{x,n}(\ell_n y) = 1 \;\wedge\; \overline{\chi}_{y,n}(\ell_n x) = 1 \;\Longrightarrow\; x \preceq_1 y \;\wedge\; y \preceq_1 x.
\]
Setting $n = 1$ along $\ell_1(x_m) \simeq \ell_1(x_1)$ places every $x_m$ in the $1$-SCC of $x_1$, proving (1). For (2), whiskering a morphism in $\mathrm{Hom}_{\mathcal{C}}(a, b)$ across invertible or non-absorbing paths in $\mathcal{S}$ transfers cell dimensions across objects unless higher composites collapse through an absorbing element.
\end{proof}

\begin{theorem}[Automatic Full Faithfulness, Vanishing Constant-Tower $\varprojlim^1$, and $1$-Acyclic Completeness on $\mathrm{Cat}_{\le \omega}$]\label{thm:ff-acyclic-completeness}%
\lean{HomotopyCRM.theorem12C_antisymmetric_postnikov_completeness}%
\leanok%
\uses{thm:classifier-scc-trapping}%
For every locally finite-dimensional $(\infty, \infty)$-category $\mathcal{C} \in \mathrm{Cat}_{\le \omega}$:
\begin{enumerate}[label=(\arabic*),leftmargin=2em]
  \item For all $x, y \in \mathrm{Ob}(\mathcal{C})$ and every $n \ge d(x, y) + 1$, $\mathrm{Hom}_{\pi_{\le n}(\mathcal{C})}(\ell_n x, \ell_n y) \simeq \pi_{\le n-1}(\mathrm{Hom}_{\mathcal{C}}(x, y)) \simeq \mathrm{Hom}_{\mathcal{C}}(x, y)$, so the Postnikov unit $\eta_{\mathcal{C}} : \mathcal{C} \hookrightarrow RL(\mathcal{C}) = \varprojlim_n i_n \pi_{\le n}(\mathcal{C})$ is always fully faithful.
  \item Every Postnikov-compatible tower $X_\bullet \in RL(\mathcal{C})$ with an eventually constant representative $X_n \simeq \ell_n(x)$ for $n \gg 0$ has stabilizing equivalence bonding maps on $\mathrm{Isom}_{\pi_{\le n}\mathcal{C}}(\ell_n x, X_n)$ for $n \ge d(x, x) + 1$; hence all constant-tower $\varprojlim^1$ twists vanish ($X_\bullet \simeq \eta_{\mathcal{C}}(x)$).
  \item If $\preceq_1$ is $1$-acyclic (antisymmetric up to equivalence in $\mathcal{C}$: $x \preceq_1 y \land y \preceq_1 x \implies x \simeq y$), then $\mathcal{C} \in \mathrm{Cat}^L_{(\infty, \infty)}$.
\end{enumerate}
\end{theorem}

\begin{proof}
Since $\pi_{\le n-1} : \mathrm{Cat}^R_{(\infty, \infty)} \to \mathrm{Cat}_{(\infty, n-1)}$ preserves finite products~\cite[Lemma~2.13 and Theorem~4.25]{ORW26} and Rezk completion $\pi_{\le n}(\mathcal{C}) \simeq L_{\mathrm{Rezk}}(\pi_{\le n-1}^{\mathrm{hom}}(\mathcal{C}))$ is Dwyer--Kan fully faithful on hom-objects~\cite[Theorem~5.6.2]{GepnerHaugseng15}, any $n \ge d(x, y) + 1$ gives $\pi_{\le n-1}(\mathrm{Hom}_{\mathcal{C}}(x, y)) \simeq \mathrm{Hom}_{\mathcal{C}}(x, y)$, proving (1). Passing to maximal groupoid cores stabilizes $\mathrm{Isom}_{\pi_{\le n}\mathcal{C}}(\ell_n x, X_n) \simeq \mathrm{Isom}_{\mathcal{C}}(x, x)$ for $n \ge d(x, x) + 1$, giving an eventually constant Mittag--Leffler tower with $\varprojlim^1 = *$ and proving (2). Under $1$-acyclicity, Theorem~\ref{thm:classifier-scc-trapping}(1) forces $x_n \simeq x_1$ for all $n \ge 1$, so (2) yields $X_\bullet \simeq \eta_{\mathcal{C}}(x_1)$ and $\mathcal{C} \in \mathrm{Cat}^L_{(\infty, \infty)}$, proving (3).
\end{proof}

\begin{theorem}[Gaunt Noetherian Rank-$\omega$ Step-Product Counterexample $\mathcal{Q}$ Disproving Conjecture~C1]\label{thm:step-product-counterexample}%
\lean{HomotopyCRM.isHereditarilyGaunt_monoid, HomotopyCRM.absorbingMonoid_collapse_le, HomotopyCRM.stepProductQ_omega_completeness, HomotopyCRM.isHereditarilyGaunt_qFull, HomotopyCRM.homQ_card_eq, HomotopyCRM.qFull_not_nTruncated, HomotopyCRM.theorem12D_stepProductQ_phantom_not_realized, HomotopyCRM.conjecture_c1_disproof_summary}%
\leanok%
\uses{thm:ordinal-rank-completeness,thm:classifier-scc-trapping,thm:ff-acyclic-completeness}%
There exists a gaunt, locally finite-dimensional, Noetherian $(\infty, \infty)$-category $\mathcal{Q}$\index{Step-product Q@Step-product $\mathcal{Q}$} of global rank $\mathrm{rk}_S(\mathcal{Q}) = \omega$ satisfying $\alpha = 0$ stabilization in $\mathrm{Cat}^\omega_{(\infty, \infty)}$ yet admitting a pure Postnikov phantom\index{Phantom object@Phantom object $X_\bullet \in RL(\mathcal{Q}) \setminus \mathcal{Q}$} $X_\bullet \in RL(\mathcal{Q}) \setminus \mathcal{Q}$, witnessing
\[
  \mathcal{Q} \;\in\; \bigl(\mathrm{Cat}_{\le \omega} \cap \mathrm{Cat}^{\mathrm{Noeth}}_{(\infty, \infty)} \cap \mathrm{Cat}^\omega_{(\infty, \infty)}\bigr) \setminus \mathrm{Cat}^L_{(\infty, \infty)}
\]
and resolving Conjecture~C1 negatively.
\end{theorem}

\begin{proof}
Construct inductively a sequence of finite gaunt commutative absorbing monoids $(\mathcal{M}_k, m_k, 1_k, e_k) \in \mathrm{Cat}_{(0, k)}$: start at $k = 1$ with $\mathcal{M}_1 := [1] = \{e_1 < 1_1\}$, product $m_1 = \wedge$, unit $1_1$, and absorbing element $e_1$. Given $\mathcal{M}_k$, let $\mathcal{M}_{k+1} \in \mathrm{Cat}_{(0, k+1)}$ have object set $\{e_{k+1}, 1_{k+1}\}$ with $\mathrm{Hom}(1_{k+1}, 1_{k+1}) := \mathcal{M}_k$ and $\mathrm{Hom}(e_{k+1}, e_{k+1}) = \mathrm{Hom}(e_{k+1}, 1_{k+1}) = \mathrm{Hom}(1_{k+1}, e_{k+1}) := *$ (unique $1$-cells $f_k : e_{k+1} \to 1_{k+1}$ and $g_k : 1_{k+1} \to e_{k+1}$), setting $g_k \circ f_k = \mathrm{id}_{e_{k+1}}$ and $f_k \circ g_k = e_k$; strict absorption of $e_k$ under $m_k$ ensures associativity and induces a commutative absorbing monoid structure $(m_{k+1}, 1_{k+1}, e_{k+1})$ on $\mathcal{M}_{k+1}$. Set $\mathcal{A}_k := \mathcal{M}_{k+1} \in \mathrm{Cat}_{(0, k+1)}$ with $0_k := e_{k+1}$, so that in every coordinate factor $\mathcal{A}_k$,
\[
  \begin{aligned}
    \mathrm{Hom}_{\mathcal{A}_k}(0_k, 0_k) &\simeq \mathrm{Hom}_{\mathcal{A}_k}(0_k, 1_k) \simeq \mathrm{Hom}_{\mathcal{A}_k}(1_k, 0_k) \simeq *, \\
    \mathrm{Hom}_{\mathcal{A}_k}(1_k, 1_k) &= \mathcal{M}_k \not\simeq *.
  \end{aligned}
\]
Each $\mathcal{A}_k$ is finite gaunt ($0_k \not\simeq 1_k$), and induction on $k$ gives $\pi_{\le n}(\mathcal{A}_k) \simeq *$ for all $0 \le n \le k$: for $1 \le n \le k$, $\pi_{\le n-1}(\mathcal{M}_k) \simeq *$ collapses $\pi_{\le n-1}^{\mathrm{hom}}(\mathcal{A}_k)$ to the walking isomorphism $J$, whose Rezk completion $\pi_{\le n}(\mathcal{A}_k)$ is $*$, while $\pi_{\le 0}(\mathcal{A}_k) \simeq \pi_{\le 0}(\pi_{\le 1}(\mathcal{A}_k)) \simeq *$.

Define $\mathcal{Q} \subseteq \prod_{k \ge 1} \mathcal{A}_k$ as the full subcategory spanned by the finite-prefix indicator sequences $e^{(m)} := (1_1, \dots, 1_m, 0_{m+1}, 0_{m+2}, \dots)$ ($m \in \mathbb{N}$). Combining tail contractibility $\prod_{k > \max(m, r)} \mathrm{Hom}_{\mathcal{A}_k}(0_k, 0_k) \simeq *$ with off-diagonal absorbing contraction on the indices $\min(m, r) < k \le \max(m, r)$ (where $\{e^{(m)}_k, e^{(r)}_k\} = \{0_k, 1_k\}$ forces $\mathrm{Hom}_{\mathcal{A}_k}(e^{(m)}_k, e^{(r)}_k) \simeq *$, formalized as \texttt{qHomAt\_\allowbreak eq\_\allowbreak pt\_\allowbreak of\_\allowbreak le}) leaves only the diagonal coordinates $1 \le k \le \min(m, r)$ where $e^{(m)}_k = e^{(r)}_k = 1_k$:
\[
  \begin{aligned}
    \mathrm{Hom}_{\mathcal{Q}}(e^{(m)}, e^{(r)}) &\;\simeq\; \prod_{k=1}^{\max(m, r)} \mathrm{Hom}_{\mathcal{A}_k}(e^{(m)}_k, e^{(r)}_k) \\
    &\;\simeq\; \prod_{k=1}^{\min(m, r)} \mathcal{M}_k \;\in\; \mathrm{Cat}_{(0, \min(m, r))},
  \end{aligned}
\]
which has exact finite hom-rank $\min(m, r)$. Hence $\mathcal{Q}$ is gaunt, locally finite-dimensional of global rank $\mathrm{rk}_S(\mathcal{Q}) = \omega$, Noetherian, and lies in $\mathrm{Cat}_{\le \omega} \cap \mathrm{Cat}^{\mathrm{Noeth}}_{(\infty, \infty)} \cap \mathrm{Cat}^\omega_{(\infty, \infty)}$ by Theorem~\ref{thm:ordinal-rank-completeness}. Yet for each $n \ge 1$, the slice inclusion $\sigma_n = (c_{1_1}, \dots, c_{1_{n-1}}, \mathrm{id}_{\mathcal{A}_n}, c_{0_{n+1}}, \dots) : \mathcal{A}_n \hookrightarrow \mathcal{Q}$ maps $0_n \mapsto e^{(n-1)}$ and $1_n \mapsto e^{(n)}$, pushing $\pi_{\le n}(\mathcal{A}_n) \simeq *$ to an equivalence $\gamma_n : \ell_n(e^{(n)}) \xrightarrow{\sim} \ell_n(e^{(n-1)})$ in $\pi_{\le n}(\mathcal{Q})$. Thus $y_n := \ell_n(e^{(n-1)})$ with transition $\pi_{\le n}(y_{n+1}) \xrightarrow{\sim} \ell_n(e^{(n)}) \xrightarrow{\gamma_n} y_n$ defines a compatible tower $X_\bullet \in RL(\mathcal{Q})$. For every $e^{(m)} \in \mathrm{Ob}(\mathcal{Q})$, the $(m+1)$-st projection $\mathrm{pr}_{m+1} : \mathcal{Q} \to \mathcal{A}_{m+1} \in \mathrm{Cat}_{(0, m+2)}$ factors universally through $\overline{\mathrm{pr}}_{m+1} : \pi_{\le m+2}(\mathcal{Q}) \to \mathcal{A}_{m+1}$, separating $\overline{\mathrm{pr}}_{m+1}((X_\bullet)_{m+2}) = 1_{m+1} \not\simeq 0_{m+1} = \overline{\mathrm{pr}}_{m+1}(\ell_{m+2}(e^{(m)}))$ in gaunt $\mathcal{A}_{m+1}$. Hence $X_\bullet \not\simeq \eta_{\mathcal{Q}}(e^{(m)})$ for all $m \in \mathbb{N}$, proving $\mathcal{Q} \notin \mathrm{Cat}^L_{(\infty, \infty)}$. Every component of this construction and separation is machine-checked in Lean~4 (\texttt{HomotopyCRM/\allowbreak StepProduct\allowbreak Phantom.lean}, 3,505 lines, 22 \texttt{\#guard\_msgs in \#print axioms} checks in \texttt{KernelAxiomGate}, including \texttt{stepProductQ\_\allowbreak rank\_\allowbreak hom}, \texttt{homQ\_\allowbreak card\_\allowbreak eq}, and \texttt{theorem12D\_\allowbreak stepProductQ\_\allowbreak phantom\_\allowbreak not\_\allowbreak realized}).
\end{proof}

\begin{remark}[The Step-Product Phantom as the $(\infty,\infty)$-Categorical Twin of $\mathcal{Q}_{E_0}$]\label{rem:step-product-e0-twin}
Comparing $\mathcal{Q} \subset \prod_{k \ge 1} \mathcal{A}_k$ with the smooth-$4$-manifold modulus $\mathcal{Q}_{E_0} = \prod_{\mathbb{N}} \mathbb{F}_2 / \bigoplus_{\mathbb{N}} \mathbb{F}_2 \simeq 2^{\mathbb{N}} \sslash E_0$ of Chapter~\ref{ch:base-fiber} (Section~\ref{sec:exotic-r4-e0}, Theorem~\ref{thm:double-solvent-r4}) reveals an exact structural parallel across opposite ends of Categorical Reverse Mathematics. In both settings, individual objects carry only finite prefixes ($e^{(m)} \in \mathrm{Ob}(\mathcal{Q})$ vs.\ finite-support elements $\bigoplus_{\mathbb{N}} \mathbb{F}_2$), while passing to asymptotic invariants---finite Postnikov truncations $\pi_{\le n}$ for $\mathcal{Q}$, or compact interior isotopies and finite-stage Polish quotients $\pi_0^{\mathrm{Pol}}$ for exotic $\mathbb{R}^4$'s---collapses finite prefixes into a single coherent phantom object ($X_\bullet \in RL(\mathcal{Q}) \setminus \mathcal{Q}$ vs.\ the non-splitting condensed class $[q] \in \mathrm{Ext}^1_{\mathrm{CondAb}}(\mathcal{Q}_{E_0}, \bigoplus_{\mathbb{N}} \mathbb{F}_2)$). Whereas $\mathbf{Set}_{\mathrm{ZFC}}$ obliterates the geometric structure of $\mathcal{Q}_{E_0}$ via a Vitali Hamel section, $\mathrm{Cat}^R_{(\infty, \infty)}$ retains $X_\bullet$ inside $RL(\mathcal{Q}) \setminus \mathcal{Q}$ as a pure Postnikov residuum (Figure~\ref{fig:inf-inf-ladder-phantom}).
\end{remark}

\begin{figure}[htbp]
\centering
\begin{tikzpicture}[
  >=Stealth,
  tier/.style={draw=CRMBlue,rounded corners=3pt,fill=CRMBlue!7,inner sep=3.5pt,font=\scriptsize\bfseries},
  box/.style={draw=black!50,rounded corners=2.5pt,fill=black!2,inner sep=3pt,font=\scriptsize,align=center},
  wit/.style={font=\scriptsize,color=CRMBlue,align=center},
  phan/.style={draw=CRMBlue,line width=0.9pt,rounded corners=3pt,fill=CRMBlue!9,inner sep=3.5pt,font=\scriptsize,align=center}
]
  % Upper Band: 4-Tier Localization Ladder
  \node[font=\scriptsize\bfseries,color=CRMBlue] at (0,4.85) {Four-Tier $(\infty,\infty)$-Reflective Localization Ladder (Contravariant Bousfield)};
  \node[tier] (T1) at (-4.05,3.75) {$\mathrm{Cat}^L_{(\infty,\infty)}$};
  \node[tier] (T2) at (-1.35,3.75) {$\mathrm{Cat}^\omega_{(\infty,\infty)}$};
  \node[tier] (T3) at (1.35,3.75) {$\mathrm{Cat}^{\mathrm{coind}}_{(\infty,\infty)}$};
  \node[tier] (T4) at (4.05,3.75) {$\mathrm{Cat}^R_{(\infty,\infty)}$};

  \draw[->,thick,draw=black!75,shorten <=2pt,shorten >=2pt] (T1) -- node[midway,fill=white,inner sep=1.5pt,font=\scriptsize]{$\subsetneq$} (T2);
  \draw[->,thick,draw=black!75,shorten <=2pt,shorten >=2pt] (T2) -- node[midway,fill=white,inner sep=1.5pt,font=\scriptsize]{$\subsetneq$} (T3);
  \draw[->,thick,draw=black!75,shorten <=2pt,shorten >=2pt] (T3) -- node[midway,fill=white,inner sep=1.5pt,font=\scriptsize]{$\subsetneq$} (T4);

  \draw[->,dashed,thick,draw=CRMBlue,bend right=25] (T2.north) to node[above,font=\tiny]{$L_L$} (T1.north);
  \draw[->,dashed,thick,draw=CRMBlue,bend right=25] (T3.north) to node[above,font=\tiny]{$L_\omega$} (T2.north);
  \draw[->,dashed,thick,draw=CRMBlue,bend right=25] (T4.north) to node[above,font=\tiny]{$L_{\mathrm{coind}}$} (T3.north);

  \node[wit,draw=CRMBlue!35,rounded corners=2pt,fill=CRMBlue!4,inner sep=2pt] at (-2.7,3.02) {$\mathcal{Q},\; W$};
  \node[wit,draw=CRMBlue!35,rounded corners=2pt,fill=CRMBlue!4,inner sep=2pt] at (0.0,3.02) {$E^{(\omega)}$};
  \node[wit,draw=CRMBlue!35,rounded corners=2pt,fill=CRMBlue!4,inner sep=2pt] at (2.7,3.02) {$\mathrm{Cob}_\infty$};
  \node[font=\tiny,color=black!70,align=center] at (-4.1,2.98) {$S^2 \in \mathrm{Cat}^L \setminus \mathrm{Cat}^{\mathrm{Noeth}}$};

  \draw[CRMBlue!35] (-5.0,2.52) -- (5.0,2.52);

  % Lower Band: Step-Product Q and Phantom X_bullet
  \node[font=\scriptsize\bfseries,color=CRMBlue] at (0,2.22) {Step-Product $\mathcal{Q} \subset \prod_{k\ge 1}\mathcal{A}_k$ and Pure Postnikov Phantom $X_\bullet \in RL(\mathcal{Q})\setminus\mathcal{Q}$};

  \node[box] (e0) at (-4.1,1.3) {$e^{(0)}$};
  \node[box] (e1) at (-2.2,1.3) {$e^{(1)}$};
  \node[box] (e2) at (-0.3,1.3) {$e^{(2)}$};
  \node[font=\scriptsize] (edots) at (0.85,1.3) {$\cdots$};
  \draw[<->,black!65] (e0) -- (e1);
  \draw[<->,black!65] (e1) -- (e2);
  \draw[<->,black!65] (e2) -- (edots);

  \node[box] (y1) at (-4.1,0.12) {$y_1 = \ell_1(e^{(0)})$};
  \node[box] (y2) at (-1.85,0.12) {$y_2 = \ell_2(e^{(1)})$};
  \node[box] (y3) at (0.4,0.12) {$y_3 = \ell_3(e^{(2)})$};
  \node[phan] (X) at (3.65,0.52) {$X_\bullet = (y_n)_{n\ge 1}$\\[-1pt]\tiny $\in RL(\mathcal{Q})\setminus\mathcal{Q}$};

  \draw[->,thick,CRMBlue] (e0) -- node[left,font=\tiny]{$\ell_1$} (y1);
  \draw[->,thick,CRMBlue] (e1) -- node[left,font=\tiny,pos=0.45]{$\ell_2$} (y2);
  \draw[->,thick,CRMBlue] (e2) -- node[left,font=\tiny,pos=0.45]{$\ell_3$} (y3);
  \draw[->,black!75] (y2) -- node[below,font=\tiny]{$\gamma_1$} (y1);
  \draw[->,black!75] (y3) -- node[below,font=\tiny]{$\gamma_2$} (y2);
  \draw[->,dashed,CRMBlue] (1.45,0.12) -- node[below right,font=\tiny]{$\varprojlim$} (X.south west);

  \node[font=\tiny,align=left,color=black!80] at (3.55,1.58) {$\mathrm{rk}\,\mathrm{Hom}_{\mathcal{Q}}(e^{(m)},e^{(r)})=\min(m,r)$\\$\overline{\mathrm{pr}}_{m+1}((X_\bullet)_{m+2})=1_{m+1}\neq 0_{m+1}$};
\end{tikzpicture}
\caption{Four-tier $(\infty,\infty)$-localization ladder and step-product phantom $X_\bullet \in RL(\mathcal{Q})\setminus\mathcal{Q}$.}\label{fig:inf-inf-ladder-phantom}
\end{figure}

On the type-theoretic and algebraic fixed-point fronts, two complementary calibrations anchor Tier~3 ($\mathrm{Cat}^{\mathrm{coind}}_{(\infty, \infty)}$) and relative undercategory perfection:
\begin{itemize}[leftmargin=1.5em]
  \item \textbf{Grothendieck--Maltsiniotis Weak-$\omega$ Clan Calibration of Infinitely Coherent Type Theory (\texttt{ICaTT}):} Benjamin, Champin, and Markakis~\cite{BCM26} formulate Finster--Mimram's weak $\omega$-categorical type theory $\mathbf{CaTT}$~\cite{FinsterMimram17} and its coinductive-invertibility extension $\mathbf{ICaTT}$\index{ICaTT@$\mathbf{ICaTT}$ coinductive type theory} via syntactic clans $\mathcal{S}_{\mathbf{CaTT}} \hookrightarrow \mathcal{S}_{\mathbf{ICaTT}}$ whose models in $\mathbf{Clan}(\mathcal{S}_{\mathbf{CaTT}}, \mathbf{Set})$ and over a Grothendieck $1$-topos $\mathbf{Clan}(\mathcal{S}_{\mathbf{CaTT}}, \mathcal{E})$ are \emph{Grothendieck--Maltsiniotis weak $\omega$-category / clan models}~\cite{BCM26,ORW26} (distinct from strict $\omega$-categories). Extending the 3-variable walking-arrow context $(d_0^- : \star) \rhd (d_0^+ : \star) \rhd (d_1 : d_0^- \to d_0^+)$ by a single coinductive witness stream $e_1 : \mathrm{Inv}(d_1)$ under suspension-substitution yields the finite walking-equivalence context $\Gamma_E \equiv \mathcal{E}_1$, whose clan semantics evaluates strictly to Ozornova--Rovelli--Walde's walking coinductive equivalence computad ($\mathcal{J}^*[\![\mathcal{E}_1]\!] \cong \mathcal{E}_1^{\mathrm{OR}}$, \cite[Theorem~6.5]{BCM26}). Bounded-$\eta$ plus $\beta$-rewriting proves $\mathbf{ICaTT}$ syntactically conservative over $\mathbf{CaTT}$ (\cite[Proposition~4.6, Theorem~6.2]{BCM26}), while localness with respect to $\Sigma^{n-1}\mathcal{E}_1^{\mathrm{OR}} \to \mathbf{D}^{n-1}$ cuts out \emph{exact Tier~3} ($\mathrm{Cat}^{\mathrm{coind}}_{(\infty, \infty)}$, failing to reach Tier~2 since $E^{(\omega)}$ is $\mathcal{E}_1$-local~\cite[Corollary~5.56]{ORW26}); internally over Grothendieck $1$-topoi $\mathcal{E}$, Theorem~\ref{thm:coinductive-stream-choice} isolates $\mathrm{IDC}_\omega$ as the exact choice residuum for witness-stream adequacy on Grothendieck--Maltsiniotis weak $\omega$-category / clan models $\mathbf{Clan}(\mathcal{S}_{\mathbf{CaTT}}, \mathcal{E})$.

  \item \textbf{Relative Undercategory Slice-Algebra Perfection (\texttt{HomotopyCRM/\allowbreak DeltaRing\allowbreak Perfection.lean}):} Complementing Tier~3 coinductive localization, Goldthorpe's relative undercategory initial-algebra theorem~\cite[Proposition~2.3.13]{Goldthorpe23} models strictly the \emph{inductive} perfect reflector $A \to A_{\mathrm{perf}} = \varinjlim_\phi A$ on Joyal $\delta$-rings~\cite{Joyal85} with $+1$/$-1$ shift chirality ($\phi_{A_{\mathrm{perf}}} = s^{-1}$), separated from $2$-torsion lift torsors ($R_2$) and non-perfect Nygaard filtrations (\texttt{DualInt},~\cite{BhattScholze22}) in Chapter~\ref{ch:delta-rings} (Proposition~\ref{prop:r2-delta-torsor}, Theorem~\ref{thm:slice-perfection-chirality}, and Remark~\ref{rem:guardrail-2}).
\end{itemize}

\chapter{Site-Level Combinatorics, Realizability Topoi, and Computational Reverse Mathematics}
\chaptermark{Site Combinatorics and Realizability Topoi}
\label{ch:realizability-sites}
\chapterepigraph{``Jump on calculations with both feet; group the operations, classify them according to their difficulty and not according to their form; such, according to me, is the task of future geometers.''}{\'Evariste Galois, \emph{Pr\'eface aux deux m\'emoires d'analyse pure} (1831)}

\section{From Sheaf Topoi to Realizability Universes}\label{sec:kleene-eff}

Having calibrated the summit of higher categorical Reverse Mathematics in $(\infty, 1)$-topoi and weak $(\infty, \infty)$-categories in Chapter~\ref{ch:higher-crm}, we now investigate the combinatorial and computational engines that generate them. Base $1$-topoi furnish the foundational universes inside which weak $\omega$-categories, relational coinductive streams, and discrete or profinite symmetry invariants are internalized. Chapter~\ref{ch:realizability-sites} unites three distinct combinatorial pillars:
\begin{enumerate}[label=\textup{(\Roman*)},leftmargin=2.0em]
  \item \textbf{Computability and Realizability Topoi (\S\ref{sec:kleene-eff}--\S\ref{sec:ltt-topologies}):} Calibrating constructive reverse mathematics ($\mathrm{WKL}_0, \mathrm{FAN}$) over Partial Combinatory Algebras via Lawvere--Tierney modal operators in $\mathbf{LTT}(\Eff)$.
  \item \textbf{Site Combinatorics and Fra\"iss\'e Permutation Topoi (\S\ref{sec:site-combinatorics}--\S\ref{subsec:kihara-choice-nominal}):} Discovering the Fra\"iss\'e atomic-topology escape from Left Bousfield groupoidification, classifying Kihara's five-rung atomic choice spectrum, and computing nominal parity Galois cohomology.
  \item \textbf{Directed Combinatorial Geometry and Higher Coinductive Streams (\S\ref{subsec:cubical-amalgamation}--\S\ref{subsec:coinductive-streams-choice}):} Machine-checking binary Strong Amalgamation ($\mathrm{SAP}$) along flow-convex embeddings, proving the all-$n \ge 3$ cubical amalgamation phase transition obstructed by facet cycle diagrams $F_n$, and calibrating internal coinductive invertibility streams ($\mathbf{ICaTT}$) against internal Dependent Choice ($\mathrm{IDC}_\omega$).
\end{enumerate}

In the preceding chapters, our primary models of variable mathematical universes were Grothendieck sheaf topoi $\mathcal{E} \simeq \mathbf{Sh}(\mathbb{C}, J)$ over topological spaces, locales, and algebraic sites, as well as their higher $(\infty, 1)$-topos generalizations. Yet as noted in Chapter~\ref{ch:foundations}, Categorical Reverse Mathematics ($\mathrm{CRM}$) is designed to transfer functorially across \emph{all} structured elementary topoi, including a second major domain of topos theory: \emph{realizability topoi} arising from computability theory.

Whereas a Grothendieck sheaf topos measures \emph{spatial and topological variation} (whether locally defined sections glue coherently across open coverings), a realizability topos measures \emph{algorithmic and computational witnessability} (whether existential and implicational statements are witnessed by uniform partial recursive programs).

To construct a realizability universe, one begins with a \emph{Partial Combinatory Algebra (PCA)}\index{Partial Combinatory Algebra (PCA)} $\mathcal{A} = (A, \cdot)$---a set $A$ equipped with a partial binary application operation $(a, b) \mapsto a \cdot b$ possessing combinators $\mathsf{k}, \mathsf{s} \in A$ that satisfy functional completeness (the $\lambda$-abstraction theorem):

\[
  \mathsf{k} \cdot a \cdot b = a, \qquad \mathsf{s} \cdot a \cdot b \downarrow \quad \text{and} \quad \mathsf{s} \cdot a \cdot b \cdot c \simeq (a \cdot c) \cdot (b \cdot c),
\]
where $\downarrow$ indicates definedness (termination of the partial application) and $\simeq$ denotes Kleene strong equality (meaning that if either side is defined, then both sides are defined and equal~\cite[Ch.~1, \S\S 1.1--1.4]{vanOosten08}).

Three canonical PCAs anchor computational reverse mathematics:

\begin{itemize}[leftmargin=1.5em]

  \item \textbf{Kleene's First Algebra ($\mathcal{K}_1$):} The underlying set is $\mathbb{N}$, and $e \cdot n \simeq \varphi_e(n)$ is the result of applying the $e$-th Turing machine (partial recursive function) to input $n \in \mathbb{N}$.

  \item \textbf{Oracle / Relative Kleene Algebras ($\mathcal{K}_1^O$):} The underlying set is $\mathbb{N}$, with application $e \cdot_O n \simeq \varphi_e^O(n)$ given by Turing machines equipped with an external oracle $O \subseteq \mathbb{N}$ (such as the halting jump $\emptyset'$ or hyperarithmetical oracles).

  \item \textbf{Kleene's Second Algebra ($\mathcal{K}_2$):} The underlying set is Baire space $\mathbb{N}^\mathbb{N}$, where $\alpha \cdot \beta$ is continuous tracking application (governed by Brouwerian continuity and fan functionals).

\end{itemize}

Via the Hyland--Johnstone--Pitts~\cite[pp.~205--232]{HJP80} \emph{tripos-to-topos construction}, any PCA $\mathcal{A}$ over a base topos $\mathbf{Set}$ generates an elementary \emph{realizability topos} $\mathbf{RT}(\mathcal{A})$ whose objects $(X, =_X)$ are sets $X$ equipped with $\mathcal{P}(\mathcal{A})$-valued partial equivalence relations $=_X: X \times X \to \mathcal{P}(\mathcal{A})$ assigning to each pair $(x, x')$ the set of realizers witnessing $x = x'$~\cite[Ch.~2, \S\S 2.1--2.2]{vanOosten08}.

When $\mathcal{A} = \mathcal{K}_1$, $\mathbf{RT}(\mathcal{K}_1)$ is Martin Hyland's Effective Topos ($\Eff$)\index{Effective Topos@Effective Topos $\mathcal{E}\mathrm{ff}$}~\cite[pp.~165--216]{Hyland82}, \cite[Chs.~2--3]{vanOosten08}. From the perspective of Categorical Reverse Mathematics, $\Eff$ exhibits a striking combination of structural invariants that cannot occur in any spatial Grothendieck topos:

\begin{itemize}[leftmargin=1.5em]

  \item \textbf{Two-Valuedness Without Booleanness:} The terminal object $1 \in \Eff$ is externally projective and indecomposable ($\Eff$ is two-valued), yet $\Eff$ is strongly non-Boolean ($\Omega \not\simeq 1+1$) because a closed sentence $\phi$ is true in $\Eff$ if and only if it possesses a computable realizer (with general propositions valued in $\Omega = \mathcal{P}(\mathbb{N})/{\sim}$)~\cite{Hyland82}, \cite[\S 3.1]{vanOosten08}.

  \item \textbf{Non-Grothendieck Exactness:} $\Eff$ is an elementary topos with a Natural Numbers Object ($\mathrm{NNO}$), but it has no points (no geometric morphisms $\mathbf{Set} \to \Eff$) and is not cocomplete over $\mathbf{Set}$~\cite{Hyland82}, \cite[\S\S 3.1--3.2]{vanOosten08}.

  \item \textbf{Internal Church's Thesis ($\mathrm{CT}$) and Markov's Principle ($\mathrm{MP}$):} Internally in $\Eff$, \emph{every} function $f: \mathbb{N} \to \mathbb{N}$ is computable ($\Eff \models \mathrm{CT}$), and double-negation elimination holds for decidable existential predicates over $\mathbb{N}$ ($\Eff \models \mathrm{MP}: \forall n\,(\phi(n) \vee \neg\phi(n)) \wedge \neg\neg\exists n\,\phi(n) \to \exists n\,\phi(n)$~\cite{Hyland82}, \cite[\S 3.2]{vanOosten08}).

  \item \textbf{Countable and Dependent Choice Without Full Choice:} Because a realizer for $\forall n \in \mathbb{N}\,\exists x \in X\, R(n, x)$ is already a computable function extracting a witness for each $n$, $\Eff$ validates Countable Choice ($\mathrm{CC}_{\mathbb{N}}$) and Dependent Choice ($\mathrm{DC}_{\mathbb{N}}$) over $\mathbb{N}$---and more generally over any projective object (partitioned assembly, such as $\mathbb{N}$)--- while refuting full $\mathrm{AC}$ (since $\mathrm{AC} \implies \mathrm{LEM}$ by Diaconescu's theorem, Chapter~\ref{ch:choice-atlas}, Theorem~\ref{thm:diaconescu}, \cite{Diaconescu75};~\cite{Hyland82}, \cite[\S 3.2]{vanOosten08}).

\end{itemize}

\section{The Computability Residuum Atlas: Kleene Trees, Weak König's Lemma, and Lawvere--Tierney Oracles}\label{sec:ltt-topologies}

Because classical Friedman--Simpson reverse mathematics is built over recursive comprehension ($\mathrm{RCA}_0$), calibrating the classical ``Big Five'' inside realizability topoi $\mathbf{RT}(\mathcal{A})$ reveals the exact categorical bridge between \emph{computability theory} and \emph{topos geometry}.

\subsection{\texorpdfstring{The Failure of Weak König's Lemma ($\mathrm{WKL}$) and the Kleene Tree Obstruction}{The Failure of Weak König's Lemma (WKL) and the Kleene Tree Obstruction}}

In classical reverse mathematics over $\mathrm{RCA}_0$, the first major reversal tier above computable mathematics is Weak König's Lemma\index{Weak Konig's Lemma@Weak K\"onig's Lemma ($\mathrm{WKL}_0$)} ($\mathrm{WKL}_0$): every infinite binary tree $T \subseteq 2^{<\mathbb{N}}$ possesses an infinite path $\alpha \in 2^{\mathbb{N}}$. In recursion theory, $\mathrm{WKL}_0$ fails in the minimal $\omega$-model $\mathrm{REC}$ of computable sets because of \emph{Kleene's recursive tree}\index{Kleene tree@Kleene tree $T_{\mathrm{Kleene}}$}~\cite{Kleene52} $T_{\mathrm{Kleene}} \subseteq 2^{<\mathbb{N}}$---an infinite, computable binary tree separating two disjoint recursively inseparable enumerable sets, which possesses no computable infinite branch~\cite[\S IV.1, \S VIII.2]{Simpson09}.

Inside the Effective Topos $\Eff = \mathbf{RT}(\mathcal{K}_1)$~\cite{Hyland82}, internal $\mathrm{WKL}$ fails for the exact same reason: $T_{\mathrm{Kleene}}$ is an internal infinite binary tree in $\Eff$, but because $\Eff \models \mathrm{CT}$, any internal infinite branch $\alpha: \mathbb{N} \to 2$ would have to be a computable function in $\mathbf{Set}$, which does not exist. Similarly, the contrapositive of $\mathrm{WKL}$---Brouwer's Fan Theorem ($\mathrm{FAN}$), asserting that every bar on $2^{\mathbb{N}}$ is uniform (equivalent to the internal compactness of the Cantor space object $2^{\mathbb{N}}$)---\emph{also} fails in $\Eff$~\cite[\S 3.2]{vanOosten08}, whereas $\mathrm{FAN}$ holds in Kleene's function-realizability topos $\mathbf{RT}(\mathcal{K}_2)$~\cite[\S 4.3]{vanOosten08} and in the sheaf topos $\mathbf{Sh}(\mathbb{N}^\mathbb{N})$~\cite[Theorem~3.2]{FourmanHyland79} (cf.~\cite{LietzStreicher12}).

What is the categorical residuum $\mathcal{R}(\mathcal{E})$ that bridges constructive recursive arithmetic ($A = \mathrm{HA}^\omega + \mathrm{DC}_{\mathbb{N}}$) to $\mathrm{WKL}$ or $\mathrm{FAN}$ inside a realizability universe $\mathcal{E} = \mathbf{RT}(\mathcal{A})$?

\begin{theorem}[The Realizability Residuums for $\mathrm{FAN}$ and $\mathrm{WKL}$]\label{thm:realizability-fan-wkl}%
\uses{def:comparison-triple,def:residuum}%

Let $\mathcal{E} = \mathbf{RT}(\mathcal{A})$ be the realizability topos over a Partial Combinatory Algebra $\mathcal{A}$.

\begin{itemize}[leftmargin=1.5em]

  \item \textbf{The Compactness / Fan Residuum ($\mathrm{RCA}_0 \Rightarrow \mathrm{FAN}$):} Elevating basic recursive arithmetic ($A$) to internal Cantor compactness / the Fan Theorem ($B = \mathrm{FAN}$, sufficient for uniform continuity and the Extreme Value Theorem on $[0,1]$) isolates the \emph{Baire-continuity / fan-functional structure} of the PCA $\mathcal{A}$ (passing from discrete Turing application $\mathcal{K}_1$ to continuous Kleene application $\mathcal{K}_2$, or restricting to the localic reflection of $2^{\mathbb{N}}$~\cite{Hyland82}, \cite[\S 3.2, \S 4.3, \S 4.5]{vanOosten08}, \cite[Thm.~3.2]{FourmanHyland79}).

  \item \textbf{The Lifschitz / Low-Oracle Residuum ($\mathrm{RCA}_0 \Rightarrow \mathrm{WKL}_0$):} Elevating recursive arithmetic ($A$) to actual branch selection in infinite binary trees ($B = \mathrm{WKL}$) without collapsing to full classical logic isolates \emph{Lifschitz bounded-non-determinism realizability}~\cite[\S 4.4]{vanOosten08} (cf.~\cite{LeeVanOosten13}, where realizers encode finite non-empty candidate sets $e \subseteq \mathbb{N}$ closed under König tree pruning, or where the PCA $\mathcal{A}$ is closed under Jockusch--Soare diagonally non-computable / $\mathrm{PA}$-degree oracles).

\end{itemize}

\end{theorem}

\subsection{Lawvere--Tierney Topologies and Weihrauch Degrees}\label{subsec:ltt-weihrauch}

In modern computable analysis, the finer structure of reverse mathematics is calibrated by \emph{Weihrauch reducibility}\index{Weihrauch reducibility} ($f \le_W g$): a computational problem (multi-valued map) $f: \subseteq X \rightrightarrows Y$ is Weihrauch reducible to $g: \subseteq Z \rightrightarrows W$ if there exist computable pre- and post-processors $\Phi, \Psi$ such that $f(x)$ is solved by a single oracle call to $g(\Phi(x))$.

In Categorical Reverse Mathematics, Sori Lee and Jaap van Oosten~\cite{LeeVanOosten13,vanOosten08} characterized the lattice of basic local operators ($j: \Omega \to \Omega$) on the Effective Topos, while Peter Lietz and Thomas Streicher~\cite{LietzStreicher12} exhibited realizability separations refuting semi-classical principles such as $\mathrm{BD}$-$\mathbb{N}$; together they demonstrated that the hierarchy of computational principles inside a realizability topos $\mathbf{RT}(\mathcal{A})$ is governed by the lattice of \emph{Lawvere--Tierney topologies (local operators)}\index{Lawvere--Tierney topology} $j: \Omega \to \Omega$. Every Lawvere--Tierney topology $j$ on $\mathbf{RT}(\mathcal{A})$ defines a full reflective \emph{sheaf subtopos} $\mathbf{Sh}_j(\mathbf{RT}(\mathcal{A})) \hookrightarrow \mathbf{RT}(\mathcal{A})$ whose internal logic forces all $j$-dense propositions to become true. More recently, Takayuki Kihara~\cite{Kihara23} proved a lattice isomorphism $\mathbf{LTT}(\Eff) \cong \mathcal{W}_{\mathrm{gen,ext}}$ with extended generalized Weihrauch degrees (where $\cap$-preserving and $(\cup, \cap)$-preserving sublattices classify generalized Weihrauch and sub-Turing / realizability subtoposes, with no minimal topology strictly above $\mathrm{id}_\Omega$), while Samuele Maschio and Davide Trotta~\cite{MaschioTrotta26} showed that the extended Weihrauch doctrine $\mathsf{eW}$ forms a tripos whose topos $\mathbf{EW}[\mathcal{A}, \mathcal{A}']$ recovers relative realizability topoi as sheaf subtoposes $\mathbf{Sh}_j$.

\begin{theorem}[The Lawvere--Tierney Subtopos Residuum Principle]\label{thm:ltt-residuum}%
\uses{thm:realizability-fan-wkl,thm:omega-dictionary}%

Let $\mathcal{E}$ be an elementary topos (such as $\Eff = \mathbf{RT}(\mathcal{K}_1)$), and let $A \implies B$ be a deductive reversal between internal mathematical principles where $\mathcal{E} \models A$, $\mathcal{E} \not\models B$, and $B$ is a \emph{$j$-hereditary} principle (closed upward along passing to smaller sheaf subtopoi $\mathbf{Sh}_{j'}(\mathcal{E}) \hookrightarrow \mathbf{Sh}_j(\mathcal{E})$, such as a geometric sequent or modal forcing condition). Within the lattice of subtopoi $\mathbf{SubTopos}(\mathcal{E}) \simeq \mathbf{LTT}(\mathcal{E})^{\mathrm{op}}$, the categorical residuum $(A \Rightarrow_{\mathbf{SubTopos}(\mathcal{E})} B)$ is the \emph{minimal Lawvere--Tierney modal operator} $j_{(A \Rightarrow B)}: \Omega \to \Omega$ generated by forcing $B$, such that for any subtopos $\mathcal{E}_j = \mathbf{Sh}_j(\mathcal{E}) \hookrightarrow \mathcal{E}$:

\[
  \mathcal{E}_j \models B \iff j \ge j_{(A \Rightarrow B)} \quad \text{in } \mathbf{LTT}(\mathcal{E}).
\]
(By contrast, anti-classical principles that are \emph{not} $j$-hereditary, such as Church's Thesis $\mathrm{CT}$, hold at $j = \mathrm{id}_\Omega$ in $\Eff$ yet fail in the strictly smaller subtopoi $\mathbf{Sh}_{j_{\mathrm{Lif}}}(\Eff)$ and $\mathbf{Sh}_{\neg\neg}(\Eff) \simeq \mathbf{Set}$.)

\end{theorem}

In accordance with the Admissibility Criterion of Chapter~\ref{ch:foundations} (\S\ref{sec:admissibility}), while $j_{(A \Rightarrow B)}: \Omega \to \Omega$ is abstractly the relative pseudocomplement in $\mathbf{LTT}(\Eff)$, genuinely solving the $\mathrm{CRM}$ problem requires characterizing $j_{(A \Rightarrow B)}$ as an explicit computational or oracle modality on the underlying PCA (such as the Lifschitz operator $j_{\mathrm{Lif}}$ or a Turing-jump operator $j_d$). Specifically, ascending the Lawvere--Tierney lattice inside $\Eff$ calibrates the exact computational cost of classical reversals:

\begin{itemize}[leftmargin=1.5em]

  \item \textbf{Identity Topology ($j = \mathrm{id}_\Omega$):} Validates $\mathrm{CT}$, $\mathrm{MP}$, $\mathrm{DC}_{\mathbb{N}}$, and the Cauchy real line $\mathbb{R}_C \simeq \mathbb{R}_D$ (because $\mathrm{CC}_{\mathbb{N}}$ holds in $\Eff$), while refuting $\mathrm{WKL}$, $\mathrm{LLPO}$, and $\mathrm{LEM}$~\cite{Hyland82}, \cite[\S 3.2]{vanOosten08}.

  \item \textbf{Lifschitz / $\mathrm{LLPO}$ Topology ($j_{\mathrm{Lif}}$):} Forces the Lesser Limited Principle of Omniscience ($\mathrm{LLPO}$) and Weak König's Lemma ($\mathrm{WKL}$) for bounded-multiplicity trees, refuting $\mathrm{CT}$ while remaining non-Boolean.

  \item \textbf{Turing-Jump Topologies ($j_O$ for Turing degrees $d = \mathrm{deg}(O)$):} Each Turing degree $d$ induces a Lawvere--Tierney topology $j_d: \Omega \to \Omega$ whose subtopos $\mathbf{Sh}_{j_d}(\Eff)$ is equivalent to the relative realizability topos $\mathbf{RT}(\mathcal{K}_1^O)$; iterating the Turing jump $d \mapsto d'$ along $\mathbb{N}$ climbs from $\mathrm{RCA}_0$ to arithmetical comprehension $\mathrm{ACA}_0$.

  \item \textbf{Double-Negation Topology ($j_{\neg\neg} = \neg\neg: \Omega \to \Omega$):} The maximum dense Lawvere--Tierney topology on $\Eff$, whose sheaf subtopos $\mathbf{Sh}_{\neg\neg}(\Eff) \simeq \mathbf{Set}$ collapses the entire computational hierarchy back into classical Boolean set theory.

\end{itemize}

\section{The Site-Level Combinatorial Dictionary for Grothendieck Topoi}\label{sec:site-combinatorics}

When we return from realizability topoi to Grothendieck sheaf topoi $\mathcal{E} \simeq \mathbf{Sh}(\mathbb{C}, J)$ over a small category $\mathbb{C}$ equipped with a Grothendieck topology $J$, a natural question arises: how do abstract categorical residuums $\mathcal{R}(\mathcal{E})$---such as Stone/De Morgan structure on $\Omega$, Booleanness, or internal choice---translate into concrete combinatorial properties on the morphisms and covering sieves of the site $(\mathbb{C}, J)$?

Drawing on Peter Johnstone's \emph{Sketches of an Elephant}~\cite{Johnstone02} and Olivia Caramello's theory of site-level invariants~\cite{Caramello18}, Categorical Reverse Mathematics establishes a direct \emph{Site-Combinatorial Residuum Dictionary}:

\begin{theorem}[Combinatorial Characterization of Logical and Choice Residuums on Sites]\label{thm:site-dictionary}%
\uses{thm:omega-dictionary,thm:bdf}%

Let $\mathcal{E} \simeq \mathbf{Sh}(\mathbb{C}, J)$ be the Grothendieck topos of sheaves on a site $(\mathbb{C}, J)$. Then the primary structural residuums $\mathcal{R}(\mathcal{E})$ pull back to the following combinatorial properties on $(\mathbb{C}, J)$:

\begin{enumerate}[label=(\arabic*),leftmargin=2em]

  \item \textbf{The De Morgan ($\mathrm{WLEM}$) Site Residuum (The $J$-Amalgamation / Ore Condition):} The subobject classifier $\Omega_{\mathcal{E}}$ is an internal De Morgan (Stone) lattice ($\mathcal{E} \models \neg\phi \vee \neg\neg\phi$) if and only if for every object $U \in \mathbb{C}$ and every pair of $J$-closed sieves $R_1, R_2$ on $U$ whose intersection $R_1 \cap R_2$ contains only $J$-empty morphisms, there exists a $J$-covering sieve $S \in J(U)$ such that for every $f: V \to U$ in $S$, either $f^*(R_1)$ is $J$-empty or $f^*(R_2)$ is $J$-empty. In a contravariant presheaf topos $\mathbf{Set}^{\mathbb{C}^{\mathrm{op}}}$ ($J$ trivial), this holds if and only if $\mathbb{C}$ satisfies the \emph{Right Ore condition on cospans}: every cospan $V_1 \xrightarrow{f_1} U \xleftarrow{f_2} V_2$ in $\mathbb{C}$ can be completed to a commutative square $V_1 \xleftarrow{g_1} W \xrightarrow{g_2} V_2$ with $f_1 g_1 = f_2 g_2$ (note that completing spans $V_1 \leftarrow U \rightarrow V_2$ is the dual condition for covariant presheaves $\mathbf{Set}^{\mathbb{C}}$; see Johnstone, \emph{Sketches of an Elephant}, \S D4.6, D4.6.3~\cite{Johnstone02}).

  \item \textbf{The Gödel--Dummett Linearity Site Residuum (Directedness / Tree Condition):} $\mathcal{E} \models (\phi \to \psi) \vee (\psi \to \phi)$ if and only if for every pair of sieves $R_1, R_2$ on $U$, there is a $J$-cover $S \in J(U)$ on which either $f^*(R_1) \subseteq f^*(R_2)$ or $f^*(R_2) \subseteq f^*(R_1)$. For a poset site $\mathcal{P}$ with the trivial topology, this holds if and only if $\mathcal{P}$ is a \emph{forest of trees} (every principal down-set/cone is linearly ordered).

  \item \textbf{The Booleanness ($\mathrm{LEM}$) Site Residuum (Dense-Covering Collapse / Groupoid Condition):} $\mathcal{E}$ is Boolean ($\Omega_{\mathcal{E}} \simeq 1+1$) if and only if every $J$-dense sieve on any object $U \in \mathbb{C}$ is already a $J$-covering sieve in $J(U)$ (meaning $J = J_{\neg\neg}$ coincides with the double-negation topology on $\mathbb{C}$). For a presheaf topos $\mathbf{Set}^{\mathbb{C}^{\mathrm{op}}}$ with the trivial topology, $\mathbf{Set}^{\mathbb{C}^{\mathrm{op}}}$ is Boolean if and only if the category $\mathbb{C}$ is a \emph{groupoid} (every morphism $f: V \to U$ in $\mathbb{C}$ is an isomorphism).

  \item \textbf{The Global Choice ($\mathrm{AC}$) Site Residuum (Split-Cover Refinement):} Over $\mathbf{Set}_{\mathrm{AC}}$, $\mathbf{Sh}(\mathbb{C}, J) \models \mathrm{AC}$ if and only if $(\mathbb{C}, J)$ is Morita-equivalent to the canonical topology on a complete Boolean algebra $B$, so that every $J$-covering sieve refines to a family of mutually disjoint subterminal projections that split globally.

\end{enumerate}

\end{theorem}

\section{\texorpdfstring{Group-Action and Permutation Topoi ($\mathbf{Set}^G$, Schanuel's Topos, and Fraïssé Automorphism Groups)}{Group-Action and Permutation Topoi (SetG, Schanuel's Topos, and Fraïssé Automorphism Groups)}}\label{sec:permutation-topoi}

Theorem~\ref{thm:site-dictionary}(3) identifies an important family of Boolean universes: whenever $\mathcal{G}$ is a group (viewed as a one-object groupoid), the presheaf topos $\mathbf{Set}^{\mathcal{G}}$ of $\mathcal{G}$-sets is automatically a \emph{Boolean topos} ($\Omega \simeq 1+1$). More generally, if $\mathcal{G}$ is a \emph{topological group} equipped with a neighborhood base of open subgroups $\mathcal{U}(\mathcal{G})$, the category $\mathbf{Cont}(\mathcal{G})$ of discrete sets equipped with a \emph{continuous} $\mathcal{G}$-action (where every element $x \in X$ has an open stabilizer subgroup $\mathrm{Stab}_{\mathcal{G}}(x) \in \mathcal{U}(\mathcal{G})$) forms an \emph{atomic, two-valued, Boolean Grothendieck topos}~\cite[Ch.~III, \S 9]{MacLaneMoerdijk92}.

In foundational logic, continuous group-action topoi $\mathbf{Cont}(\mathcal{G})$ are the exact categorical incarnation of \emph{Fraenkel--Mostowski--Specker permutation models} of set theory with atoms ($\mathrm{ZFA} + \neg\mathrm{AC}$)~\cite[Ch.~4]{Jech73} and \emph{Fraïssé homogeneous model theory}~\cite{Fraisse54}, \cite[\S 7.1]{Hodges93}, exhibiting a sharp dichotomy between discrete groups possessing a minimal open subgroup $\{1\}$ and non-discrete oligomorphic Polish groups possessing no minimal open subgroup (Kihara~\cite[Theorems~6.2 and~6.19]{Kihara26}):

\begin{itemize}[leftmargin=1.5em]

  \item \textbf{Discrete Group Topoi ($\mathbf{Set}^G$):} For any non-trivial discrete group $G \neq \{1\}$, the trivial subgroup $\{1\} \subsetneq G$ \emph{is} open, so the principal torsor $G = G/\{1\}$ is a minimal-stabilizer atom and an internally projective generator ($G \times A \twoheadrightarrow A_{\mathrm{triv}} \cong A$). Consequently, over $\mathbf{Set}_{\mathrm{AC}}$, $\mathbf{Set}^G$ is Boolean and satisfies both the Internal Axiom of Choice ($\mathrm{IAC}$) and $\mathrm{COSHEP}$ (Kihara~\cite[Theorem~6.2]{Kihara26}). However, $\mathbf{Set}^G$ still fails External Choice ($\mathrm{EAC}$ / $\mathrm{AC}$) and localicity because the terminal object $1$ is \emph{not} externally projective: the unique map $G \twoheadrightarrow 1$ from the free transitive $G$-set $G$ onto the trivial $G$-set $1$ is an epimorphism with empty global sections $\Gamma(G) = G^G = \emptyset$ (Kihara~\cite[Theorem~6.19]{Kihara26}).

  \item \textbf{Schanuel's Topos and Non-Discrete Fraïssé Permutation Topoi ($\mathbf{Cont}(\mathrm{Aut}(\mathcal{M}))$):} Let $\mathcal{M}$ be a countable ultrahomogeneous relational structure (a \emph{Fraïssé limit}~\cite[\S 7.1]{Hodges93}, such as an infinite countable set $I$ with equality, dense linear order $(\mathbb{Q}, <)$, or the Rado graph), and let $\mathcal{G} = \mathrm{Aut}(\mathcal{M})$ be its Polish topological automorphism group equipped with the pointwise convergence topology (where open subgroups are pointwise stabilizers $\mathcal{G}_{(F)}$ of finite substructures $F \subseteq_{\mathrm{fin}} \mathcal{M}$). Then $\mathcal{E} = \mathbf{Cont}(\mathrm{Aut}(\mathcal{M}))$ (known as \emph{Schanuel's nominal topos}\index{Nominal topos@Nominal topos $\mathbf{Nom}$} $\mathbf{Nom}$ when $\mathcal{M}$ is a pure infinite set~\cite[Ch.~III, \S 9]{MacLaneMoerdijk92}) is a two-valued ($\mathrm{Sub}(1) = \{0, 1\}$), Boolean, atomic Grothendieck topos. For \emph{every} topological group $\mathcal{G}$ possessing a proper open subgroup $U \subsetneq \mathcal{G}$, the homogeneous coset surjection $\mathcal{G}/U \twoheadrightarrow 1$ is an epimorphism in $\mathbf{Cont}(\mathcal{G})$ with empty global sections $\Gamma(\mathcal{G}/U) = (\mathcal{G}/U)^{\mathcal{G}} = \emptyset$ (\texttt{isEmpty\_\allowbreak globalSections\_\allowbreak quotient\_\allowbreak of\_\allowbreak lt\_\allowbreak top} and \texttt{Checks.\allowbreak layer2\_\allowbreak sep\_\allowbreak guardrail5\_\allowbreak coset\_\allowbreak globalSections} in \texttt{HomotopyCRM/\allowbreak NominalSimplicity.lean}), so $1$ is \emph{not} externally projective and localicity fails unless $\mathcal{U}(\mathcal{G}) = \{\mathcal{G}\}$ ($\mathbf{Cont}(\mathcal{G}) \simeq \mathbf{Set}$). Moreover, in non-discrete oligomorphic Fraïssé permutation topoi $\mathbf{Cont}(\mathrm{Aut}(\mathcal{M}))$ without minimal open subgroups (where $\{1\}$ is not open, \texttt{Checks.not\_isOpenSubgroup\_bot}, since $\mathcal{G}_{(F \cup \{a\})} \subsetneq \mathcal{G}_{(F)}$), no projective atom exists, so external projectivity of $1$ ($\mathrm{EAC}$), localicity, $\mathrm{COSHEP}$, $\mathrm{IAC}$, and $\mathrm{AC}_{\mathrm{Set}}$ (for unstable $\mathcal{M}$) or $\mathrm{AC}$ all fail, with external projectivity of $1$, localicity, $\mathrm{COSHEP}$, $\mathrm{IAC}$, and $\mathrm{AC}$ collapsing together to $\mathcal{U}(\mathcal{G}) = \{\mathcal{G}\}$ ($\mathbf{Cont}(\mathcal{G}) \simeq \mathbf{Set}$).

\end{itemize}

Why do external projectivity of $1$ and Choice fail in $\mathbf{Cont}(\mathrm{Aut}(\mathcal{M}))$ even though the topos is Boolean and two-valued? Because whenever a proper open subgroup $U \subsetneq \mathcal{G}$ exists (such as $\mathcal{G}_{(F)}$ for any non-empty finite support $F \subseteq_{\mathrm{fin}} \mathcal{M}$), the inhabited homogeneous orbit $\mathcal{G}/U \twoheadrightarrow 1$ admits no $\mathcal{G}$-invariant global section ($\Gamma(\mathcal{G}/U) = (\mathcal{G}/U)^{\mathcal{G}} = \emptyset$), and non-trivial finite-support orbits $\mathcal{G}/\mathcal{G}_{(F)}$ possess strictly smaller open stabilizers $\mathcal{G}_{(F')} \subsetneq \mathcal{G}_{(F)}$ with no canonical invariant selector---preventing $1$ from being externally projective whenever $U \subsetneq \mathcal{G}$ exists, and preventing any inhabited object of $\mathbf{Cont}(\mathrm{Aut}(\mathcal{M}))$ from being internally projective in the absence of a minimal open subgroup.

Calibrating the hierarchy of weak choice principles inside permutation topoi $\mathbf{Cont}(\mathcal{G})$ yields the \emph{Symmetry-Breaking / Support Residuum Theorem}:

\begin{theorem}[The Permutation-Topos Residuum Atlas in $\mathbf{Cont}(\mathcal{G})$]\label{thm:permutation-topos}%
\lean{HomotopyCRM.pointwiseStabilizer_lt_top_of_nonempty, HomotopyCRM.isEmpty_globalSections_quotient_of_lt_top, HomotopyCRM.open_normal_eq_top, HomotopyCRM.open_finiteIndex_eq_top, HomotopyCRM.instModuleZMod2EvenFinset, HomotopyCRM.instAddTorsorEvenFinsetTOdd, HomotopyCRM.isEmpty_t_odd_sections_empty, HomotopyCRM.nonempty_t_odd_sections_singleton, HomotopyCRM.atomSection_vsub_card_and_ne_zero}%
\leanok%
\uses{thm:site-dictionary,thm:coshep-svc}%

Let $\mathcal{G} = \mathrm{Aut}(\mathcal{M})$ be the topological automorphism group of a Fraïssé structure $\mathcal{M}$ over $\mathbf{Set}_{\mathrm{AC}}$, and let $\mathcal{E} = \mathbf{Cont}(\mathcal{G})$ be the atomic Boolean Grothendieck topos of continuous $\mathcal{G}$-sets:

\begin{itemize}[leftmargin=1.5em]

  \item \textbf{Small Violations of Choice ($\mathcal{E} \models \mathrm{SVC}$):} Every Fraïssé permutation topos $\mathbf{Cont}(\mathcal{G})$ unconditionally satisfies Blass's Small Violations of Choice ($\mathrm{SVC}$), because every continuous $\mathcal{G}$-set is a subquotient of $C \times S$ where $C$ is a trivial (choice) $\mathcal{G}$-set and $S = \coprod_{F \subseteq_{\mathrm{fin}} \mathcal{M}} \mathcal{G}/\mathcal{G}_{(F)}$ is the universal object of finite supports.

  \item \textbf{From $\mathrm{SVC}$ to Enough Projectives and External Choice ($\mathrm{SVC} \Rightarrow \mathrm{COSHEP}$ and $\mathrm{AC}$):} By Theorem~\ref{thm:coshep-svc} and Kihara~\cite[Theorems~6.2 and~6.19]{Kihara26}, in any atomic, two-valued permutation topos $\mathbf{Cont}(\mathcal{G})$ (where $\mathrm{Sub}(1) = \{0, 1\}$), external projectivity of $1$ ($\mathrm{EAC}$ / $\mathrm{AC}$) and localicity hold if and only if $\mathcal{U}(\mathcal{G}) = \{\mathcal{G}\}$ ($\mathbf{Cont}(\mathcal{G}) \simeq \mathbf{Set}$), whereas $\mathrm{COSHEP}$ and $\mathrm{IAC}$ hold if and only if $\mathcal{U}(\mathcal{G})$ possesses a minimal open subgroup (achieved by $\{1\} \in \mathcal{U}(G)$ in discrete $\mathbf{Set}^G$, but failing across all non-discrete Fraïssé groups $\mathcal{G} = \mathrm{Aut}(\mathcal{M})$ without minimal open stabilizers, where $\mathrm{EAC}$, localicity, $\mathrm{COSHEP}$, $\mathrm{IAC}$, and $\mathrm{AC}$ all collapse to $\mathcal{U}(\mathcal{G}) = \{\mathcal{G}\}$).

  \item \textbf{Finite $k$-Choice ($\mathrm{FAC}_k$) and Ramsey / Amalgamation Symmetries:} For each integer $k \ge 2$, $\mathbf{Cont}(\mathcal{G})$ satisfies Finite $k$-Choice ($\mathrm{FAC}_k$) on constant finite fibers if and only if no open stabilizer subgroup $U \in \mathcal{U}(\mathcal{G})$ admits a continuous transitive action on a $k$-element set---isolating the absence of finite index-$k$ symmetry obstructions $H^1_{\mathrm{cont}}(U, S_k) = 0$ in the Fraïssé automorphism group $\mathrm{Aut}(\mathcal{M})$.

\end{itemize}

\end{theorem}

\section{\texorpdfstring{The Fraïssé Atomic-Topology Escape from Presheaf Groupoidification}{The Fraïssé Atomic-Topology Escape from Presheaf Groupoidification}}\label{subsec:fraisse-atomic-escape}

Combining the Site-Combinatorial Dictionary (Theorem~\ref{thm:site-dictionary}) with the Fraïssé permutation topoi of Theorem~\ref{thm:permutation-topos} resolves a foundational tension between directed combinatorial embeddings and classical Boolean logic: how can a category $\mathcal{K}$ of finite structures and strict, non-invertible embeddings $A \hookrightarrow B$ generate a Boolean classifying topos without forcing its embeddings to become isomorphisms?

\begin{theorem}[The Atomic-Topology Escape from Left Bousfield Collapse for Fraïssé Classes]\label{thm:fraisse-atomic-escape}%
\lean{HomotopyCRM.sap_closed_oriented_poset, HomotopyCRM.steiner_cycle_obstruction, HomotopyCRM.IsFlowReachReflecting.toIsClosedEmbedding, HomotopyCRM.IsFlowConvex.toIsFlowReachReflecting, HomotopyCRM.sap_flowReachReflecting, HomotopyCRM.sap_flowConvex}%
\leanok%
\uses{thm:permutation-topos}%
Let $\mathcal{K}$ be a countable Fraïssé class of finite relational structures and strict embeddings satisfying the Joint Embedding Property ($\mathrm{JEP}$) and the Amalgamation Property ($\mathrm{AP}$), with countable ultrahomogeneous Fraïssé limit $\mathcal{M}$, and write $\mathbb{C} = \mathcal{K}^{\mathrm{op}}$ so that covariant diagrams $\mathbf{Set}^{\mathcal{K}}$ coincide with contravariant presheaves $\mathbf{Set}^{\mathbb{C}^{\mathrm{op}}}$.
\begin{enumerate}[label=(\arabic*),leftmargin=2em]
  \item \textbf{The Presheaf Groupoidification Obstruction:} By Theorem~\ref{thm:site-dictionary}(3), the bare presheaf topos $\mathbf{Set}^{\mathcal{K}} = \mathbf{Set}^{\mathbb{C}^{\mathrm{op}}}$ (with the trivial topology) is Boolean ($\Omega \simeq 1+1$) if and only if $\mathcal{K}$ is a groupoid; forcing classical logic onto $\mathbf{Set}^{\mathcal{K}}$ without a non-trivial Grothendieck topology---or applying groupoidification $\pi_{\le 0}(\mathcal{K})$---inverts all embeddings $A \hookrightarrow B$ into isomorphisms.
  \item \textbf{Amalgamation as the Right Ore Residuum for $\mathrm{WLEM}$:} Under $\mathbb{C} = \mathcal{K}^{\mathrm{op}}$, a cospan $B_1 \to A \leftarrow B_2$ in $\mathbb{C}$ is a span of embeddings $B_1 \hookleftarrow A \hookrightarrow B_2$ in $\mathcal{K}$, so the Fraïssé Amalgamation Property ($\mathrm{AP}$) on $\mathcal{K}$ is equivalent to the Right Ore condition\index{Right Ore condition} on cospans in $\mathbb{C}$, validating De Morgan's Law ($\mathrm{WLEM}$) in $\mathbf{Set}^{\mathcal{K}}$ by Theorem~\ref{thm:site-dictionary}(1).
  \item \textbf{Atomic Sheafification Preserves Directed Embeddings:} Because $\mathbb{C} = \mathcal{K}^{\mathrm{op}}$ satisfies the Right Ore condition, every non-empty sieve on $\mathbb{C}$ is $\neg\neg$-dense, so $J_{\neg\neg}$ coincides with the atomic Grothendieck topology $J_{\mathrm{at}}$\index{Atomic Grothendieck topology@Atomic Grothendieck topology $J_{\mathrm{at}}$} (generated by single morphisms). Passing to the sheaf subtopos $\mathbf{Sh}(\mathcal{K}^{\mathrm{op}}, J_{\mathrm{at}}) \simeq \mathbf{Cont}(\mathrm{Aut}(\mathcal{M}))$~\cite{Caramello18} achieves full Booleanness ($\mathrm{LEM}$) and two-valuedness (via $\mathrm{JEP}$) without inverting the embeddings of $\mathcal{K}$, trading morphism invertibility for the non-localic stabilizer obstructions of Theorem~\ref{thm:permutation-topos}.
\end{enumerate}
\end{theorem}

\section{\texorpdfstring{Kihara's Five-Rung Atomic Choice Spectrum and Nominal Parity Cohomology}{Kihara's Five-Rung Atomic Choice Spectrum and Nominal Parity Cohomology}}\label{subsec:kihara-choice-nominal}

Calibrating internal and external choice principles across Boolean atomic Grothendieck topoi $\mathcal{E} = \mathbf{Sh}(\mathcal{K}_\lambda^{\mathrm{op}}, J_{\mathrm{at}})$ associated to subcanonical amalgamation ($\mathrm{AP}$) classes $\mathcal{K}_\lambda$ parameterized by an infinite regular cardinal $\lambda$ satisfying $\lambda^{<\lambda} = \lambda$ (with $\lambda = \aleph_0$ recovering countable Fra\"iss\'e classes of finite structures), Takayuki Kihara~\cite[Sections~3, 5--8, and Table~1]{Kihara26} establishes a \emph{strict five-rung Choice Spectrum}\index{Kihara five-rung Choice Spectrum} (charted with its Lean-verified obstructions in Figure~\ref{fig:kihara-choice-spectrum}):
\[
  \mathrm{LEM} \;\subsetneq\; \mathrm{DC}_\omega + \mathrm{AC}_{\le \lambda} \;\subsetneq\; \mathrm{AC}_{\mathrm{Set}} \;\subsetneq\; \mathrm{IAC} + \mathrm{COSHEP} \;\subsetneq\; \mathrm{EAC}.
\]
\begin{itemize}[leftmargin=1.5em]
  \item \textbf{Rung~0 ($\mathrm{LEM}$ holds; $\mathrm{DC}_\omega + \mathrm{AC}_{\le \lambda}$ fails):} Schanuel's nominal topos $\mathbf{Nom} = \mathbf{Sh}(\mathbf{Set}_{\mathrm{fin}}^{\mathrm{op}}, J_{\mathrm{at}}) \simeq \mathbf{Cont}(\mathrm{Sym}(\mathbb{N}))$ validates $\mathrm{LEM}$, $\mathrm{SVC}$, and constant-fiber $\mathrm{FAC}_k$, yet refutes both $\mathrm{AC}_\omega$ and $\mathrm{DC}_\omega$ ($\mathrm{IDC}_\omega$) because $\mathbf{Set}_{\mathrm{fin}}$ has unbounded finite cardinalities (\cite[Theorem~7.9]{Kihara26}); similarly, $\mathbf{Sh}(\mathrm{WO}_\lambda^{\mathrm{op}}, J_{\mathrm{at}})$ validates $\mathrm{LEM} + \mathrm{AC}_\lambda$ (\cite[Proposition~7.8]{Kihara26}) while refuting $\mathrm{DC}_\omega$ via the shift $s(n) = n + 1 : \omega \hookrightarrow \omega$ (\cite[Proposition~7.5]{Kihara26}) and refuting $\mathrm{AC}_{\lambda^+}$ (\cite[Proposition~7.7]{Kihara26}).
  \item \textbf{Rung~1 ($\mathrm{DC}_\omega + \mathrm{AC}_{\le \lambda}$ holds; $\mathrm{AC}_{\mathrm{Set}}$ fails):} Unstable classes admitting $\omega$-chain and $\le \lambda$-fan cocones at $\lambda^{<\lambda} = \lambda$---linear orders $\mathrm{LO}_\lambda$, groups $\mathrm{Grp}_\lambda$, Boolean algebras $\mathrm{BA}_\lambda$, and graphs $\mathrm{Graph}_\lambda$---validate $\mathrm{DC}_\omega$ and $\mathrm{AC}_{\le \lambda}$ (\cite[Theorems~3.3--3.4 and Table~1]{Kihara26}), yet refute $\mathrm{AC}_{\mathrm{Set}} \equiv \forall \kappa, \mathrm{AC}_\kappa$: under $\lambda$-$\mathrm{TV}$ ($\lambda$-two-valuedness, meaning every non-trivial slice over a $\mathcal{K}_\lambda$-object splits into at most $\lambda$ pieces, ensuring $\mathcal{K}_\lambda$ generates without hyper-split obstructions), $\mathrm{AC}_{\mathrm{Set}} \iff \mathrm{AC}_{\lambda^+} \iff$ \emph{Galois $\lambda$-stability} $s(\mathcal{K}; \lambda) = \lambda$, where $s(\mathcal{K}; \lambda)$ denotes the cardinality of the space of Galois $1$-types over a base structure of size $\lambda$ (\cite[Theorems~5.6 and~5.10; see also Theorem~6.34 for Polish groups]{Kihara26}).
  \item \textbf{Rung~2 ($\mathrm{AC}_{\mathrm{Set}} + \mathrm{DC}_\omega$ holds; $\mathrm{IAC} + \mathrm{COSHEP}$ fails):} Unbounded Galois-stable ($\omega$-stable) classes---$\mathrm{Set}_\lambda$, $\mathrm{Vec}_{F, \lambda}$, $\mathrm{AbGrp}_\lambda$, and $\mathrm{Mod}(\mathrm{ACF}_0)_\lambda$ ($\lambda \ge \aleph_0$)---validate $\mathrm{DC}_\omega + \mathrm{AC}_{\mathrm{Set}}$ (\cite[Theorem~5.6, Corollary~7.1]{Kihara26}), yet lack maximal structures, refuting $\mathrm{COSHEP} \iff \mathrm{AtPAx} \iff \mathrm{IAC} \iff$ \emph{Maximal Extension Property}, where $\mathrm{AtPAx}$ is the Atomic Presentation / Projectivity Axiom asserting that every object is covered by a projective atom (\cite[Proposition~6.6 and Theorems~6.2, 6.7, and~6.8]{Kihara26}).
  \item \textbf{Rung~3 ($\mathrm{IAC} + \mathrm{COSHEP}$ holds; $\mathrm{EAC}$ fails):} Bounded $\mathrm{AP}$ classes with non-rigid maximal structures---$\mathrm{Set}_{\le n}$ ($n \ge 2$, yielding $\mathbf{Set}^{S_n}$) and finite discrete $\mathbf{Set}^G$ ($G \neq \{1\}$)---satisfy Maximal Extension ($\mathrm{IAC} + \mathrm{COSHEP}$) while failing rigidity ($\mathrm{EAC}$ fails by~\cite[Theorem~6.19]{Kihara26}).
  \item \textbf{Rung~4 ($\mathrm{EAC}$ and localicity hold):} Maximal Extension plus rigidity of maximal structures ($\mathrm{Aut}(M_i) = \{1\}$, e.g., $\mathrm{Set}_{\le 1}$), collapsing connected $\mathcal{E}$ to $\mathbf{Set}$ ($\mathcal{U}(\mathcal{G}) = \{\mathcal{G}\}$, \cite[Theorem~6.19]{Kihara26}).
\end{itemize}

\begin{figure}[htbp]
\centering
\begin{tikzpicture}[
  >=Stealth,
  rung/.style={draw=CRMBlue,rounded corners=2.5pt,fill=CRMBlue!6,minimum width=3.9cm,inner sep=3.5pt,font=\scriptsize,align=center},
  obs/.style={draw=black!40,rounded corners=2.5pt,fill=black!2,text width=6.1cm,inner sep=3.5pt,font=\scriptsize,align=left}
]
  \node[rung] (R4) at (-3.55,5.0) {\textbf{Rung 4:} $\mathrm{EAC}$ + Localicity};
  \node[rung] (R3) at (-3.55,3.75) {\textbf{Rung 3:} $\mathrm{IAC} + \mathrm{COSHEP}$};
  \node[rung] (R2) at (-3.55,2.5) {\textbf{Rung 2:} $\mathrm{AC}_{\mathrm{Set}} + \mathrm{DC}_\omega$};
  \node[rung] (R1) at (-3.55,1.25) {\textbf{Rung 1:} $\mathrm{DC}_\omega + \mathrm{AC}_{\le\lambda}$};
  \node[rung,fill=CRMBlue!12,line width=0.9pt] (R0) at (-3.55,0.0) {\textbf{Rung 0:} $\mathrm{LEM} + \mathrm{SVC},\;\neg\mathrm{DC}_\omega$};

  \draw[->,thick,CRMBlue] (R0) -- (R1);
  \draw[->,thick,CRMBlue] (R1) -- (R2);
  \draw[->,thick,CRMBlue] (R2) -- (R3);
  \draw[->,thick,CRMBlue] (R3) -- (R4);

  \node[obs,right=0.45cm of R4] (O4) {$\mathbf{Set}$ ($\mathrm{Set}_{\le 1}$): rigid maximal objects, $\mathrm{Aut}(M)=\{1\}$.};
  \node[obs,right=0.45cm of R3] (O3) {$\mathbf{Set}^{\mathbb{Z}/2\mathbb{Z}},\;\mathbf{Set}^{S_n}$: $1_{\mathcal{E}}$ fails external projectivity ($\Gamma(\mathbb{Z}/2\mathbb{Z})=\emptyset$).};
  \node[obs,right=0.45cm of R2] (O2) {$\mathrm{Set}_\lambda,\;\mathrm{Vec}_{F,\lambda},\;\mathrm{ACF}_0$: Galois $\omega$-stable; no maximal structures.};
  \node[obs,right=0.45cm of R1] (O1) {$\mathrm{LO}_\lambda,\;\mathrm{Graph}_\lambda$: $\mathrm{IDC}_\omega$ holds; \texttt{ICaTT} streams adequate ($\mathrm{im}(\mathrm{InvStr})=E_{\mathrm{rel}}$).};
  \node[obs,draw=CRMBlue!80,right=0.45cm of R0] (O0) {$\mathbf{Nom}=\mathbf{Cont}(\mathrm{Sym}(\mathbb{N}))$: $[\mathcal{T}_{\mathrm{odd}}]\neq 0\in H^1$; $\mathrm{InvStr}_1(\mathrm{Free}_{\mathbf{CaTT}}(E[\mathbb{A}]))=\emptyset$; $\mathrm{Mono}(\mathbb{A},\mathbb{A})=\mathrm{Iso}(\mathbb{A},\mathbb{A})$.};

  \draw[-,dotted,CRMBlue] (R4) -- (O4);
  \draw[-,dotted,CRMBlue] (R3) -- (O3);
  \draw[-,dotted,CRMBlue] (R2) -- (O2);
  \draw[-,dotted,CRMBlue] (R1) -- (O1);
  \draw[-,dotted,CRMBlue] (R0) -- (O0);
\end{tikzpicture}
\caption{Kihara's five-rung atomic choice spectrum calibrated against $\mathbf{Nom}$ and discrete $\mathbf{Set}^{\mathbb{Z}/2\mathbb{Z}}$.}\label{fig:kihara-choice-spectrum}
\end{figure}

At Rung~0, the sheaf-cohomological separation between constant and non-constant coefficients in Schanuel's nominal topos $\mathbf{Nom} \simeq \mathbf{Cont}(\mathrm{Sym}(\mathbb{N}))$ has been machine-checked in Lean~4 (\texttt{HomotopyCRM/\allowbreak Nominal\allowbreak Simplicity.lean}, 1,275 lines, 4-Layer Gate verified): because every open normal or open finite-index subgroup of $\mathrm{Sym}(\mathbb{N})$ is improper (\texttt{open\_\allowbreak normal\_\allowbreak eq\_\allowbreak top} via explicit block-swap conjugators, \texttt{open\_\allowbreak finiteIndex\_\allowbreak eq\_\allowbreak top}), constant-coefficient Galois cohomology vanishes identically, $H^1(\mathbf{Nom}; \Delta(\mathbb{Z}/2\mathbb{Z})) = \mathrm{Hom}_{\mathrm{cts}}(\mathrm{Sym}(\mathbb{N}), \mathbb{Z}/2\mathbb{Z}) = 0$. By contrast, equipping finite subsets of $\mathbb{N}$ with symmetric difference $\triangle$ forms the $\mathbb{F}_2$-module sheaf $\mathrm{Finset}_\Delta(\mathbb{N}) = \mathbb{F}_2[\mathbb{A}]$ (\texttt{natCast\_\allowbreak card\_\allowbreak symmDiff\_\allowbreak zmod2}) and yields the short exact sequence of abelian group sheaves in $\mathbf{Nom}$:
\[
  0 \longrightarrow \mathrm{EvenFinset} \longrightarrow \mathrm{Finset}_\Delta(\mathbb{N}) \xrightarrow{\;|\cdot|\bmod 2\;} \Delta(\mathbb{Z}/2\mathbb{Z}) \longrightarrow 0,
\]
locally split over $1$-atom covers $\mathbf{y}(\{a\}) \twoheadrightarrow 1$ by \texttt{atomSection a} $= \{a\}$ (\texttt{nonempty\_\allowbreak t\_\allowbreak odd\_\allowbreak sections\_\allowbreak singleton}) with odd-cardinality torsor fiber $\mathcal{T}_{\mathrm{odd}} := (|\cdot|\bmod 2)^{-1}(1)$\index{Odd-finset parity torsor@Odd-finset parity torsor $\mathcal{T}_{\mathrm{odd}}$} (\texttt{T\_\allowbreak odd}, \texttt{Checks.\allowbreak tOdd\allowbreak AddTorsor\allowbreak Check}). Since no non-empty finite subset of $\mathbb{N}$ is $\mathrm{Sym}(\mathbb{N})$-invariant, $H^0(\mathbf{Nom}; \mathrm{EvenFinset}) = H^0(\mathbf{Nom}; \mathrm{Finset}_\Delta(\mathbb{N})) = \{\emptyset\} = 0$ while $H^0(\mathbf{Nom}; \Delta(\mathbb{Z}/2\mathbb{Z})) = \mathbb{Z}/2\mathbb{Z}$, so the connecting homomorphism $\delta : \mathbb{Z}/2\mathbb{Z} \hookrightarrow H^1(\mathbf{Nom}; \mathrm{EvenFinset})$ injects $1 \mapsto [\mathcal{T}_{\mathrm{odd}}] = \delta(1) \neq 0$ (\texttt{isEmpty\_\allowbreak t\_\allowbreak odd\_\allowbreak sections\_\allowbreak empty}, witnessed by the Čech $1$-cocycle $\{a\} \triangle \{b\} = \{a, b\} \in \mathrm{EvenFinset}$ for $a \neq b$ in \texttt{atomSection\_\allowbreak vsub\_\allowbreak card\_\allowbreak and\_\allowbreak ne\_\allowbreak zero} and \texttt{Checks.\allowbreak derived\_\allowbreak cech\_\allowbreak cocycle\_\allowbreak atom\allowbreak Section}).

\section{\texorpdfstring{Cubical $n$-Amalgamation and $2 \times 2$ Orthogonality}{Cubical n-Amalgamation and 2x2 Orthogonality}}\label{subsec:cubical-amalgamation}

While Kihara's Theorem~5.6~\cite{Kihara26} equates $\mathrm{AC}_{\mathrm{Set}}$ with 1-dimensional Galois fan stability, higher $n$-amalgamation ($n \ge 3$) measures higher-dimensional descent across punctured $n$-cubes $\mathcal{P}^-(n) := \{S \subseteq [n] : S \subsetneq [n]\}$:

\begin{proposition}[Internal Matching Criterion for $n$-Amalgamation]\label{prop:matching-amalgamation}%
\lean{HomotopyCRM.PuncturedCube.card_eq, HomotopyCRM.cubical_ndap_closed_oriented_poset, HomotopyCRM.hasCubicalNDAP_iff_isIntersective}%
\leanok%
\uses{thm:fraisse-atomic-escape}%
In $\mathcal{E} = \mathbf{Sh}(\mathcal{K}_\lambda^{\mathrm{op}}, J_{\mathrm{at}})$, a punctured $n$-cube $D : \mathcal{P}^-(n) \to \mathcal{K}_\lambda$ ($n \ge 2$) admits an $n$-amalgamation cocone in $\mathcal{K}_\lambda$ if and only if the canonical projection from the internal matching limit $\mathrm{Match}_n(D) :=\allowbreak \varprojlim_{S \in \mathcal{P}^-(n)^{\mathrm{op}}} \mathbf{y}_{D(S)} \to \mathbf{y}_{D(\emptyset)}$ is an epimorphism in $\mathcal{E}$. Moreover, $D$ admits Disjoint $n$-Amalgamation ($n$-$\mathrm{DAP}$) if and only if $\mathrm{Match}_n^{\mathrm{disj}}(D) \allowbreak\twoheadrightarrow \mathbf{y}_{D(\emptyset)}$ is epimorphic, where $\mathrm{Match}_n^{\mathrm{disj}}(D) \hookrightarrow \mathrm{Match}_n(D)$ is the disjoint matching subobject carving out tuples $(j_S)_{S \subsetneq [n]}$ whose images satisfy exact face intersection $j_S(D(S)) \cap \allowbreak j_T(D(T)) = \allowbreak j_{S \cap T}(D(S \cap T))$ (matching \texttt{D.\allowbreak Is\allowbreak Intersective} on $\mathbf{Poset}_{\mathrm{og,fin}}^{\mathrm{closed}}$).
\end{proposition}

Consequently, Galois stability ($\mathrm{AC}_{\mathrm{Set}}$) and higher $n$-amalgamation ($n \ge 3$) form an orthogonal $2 \times 2$ classification:

\medskip
\noindent{\scriptsize\setlength{\tabcolsep}{3.2pt}\renewcommand{\arraystretch}{1.15}%
\begin{tabularx}{\textwidth}{@{}P{0.24\textwidth}YY@{}}
\toprule
\textbf{Stability vs.\ $n$-AP} & \textbf{All-$n$ Amalgamation Holds ($\forall n \ge 2$)} & \textbf{Higher $m$-Amalgamation Fails ($m \ge 3$, $(< m)$-AP Holds)} \\
\midrule
Galois Stable ($\mathcal{E} \models \mathrm{AC}_{\mathrm{Set}}$) & $\mathrm{ACF}_0$ ($\omega$-stable; linear disjointness of $\overline{K(\bar{a}_S)}$ resolves all $n$-cubes) & Goodrick--Kim--Kolesnikov groupoid covers $T_m$~\cite{GKK14} ($\omega$-stable so $\mathrm{AC}_{\mathrm{Set}}$ holds, yet abelian $(m-2)$-cocycle twists on $\partial\Delta^{m-1} \simeq S^{m-2}$ obstruct $\mathrm{Match}_m$) \\
Galois Unstable ($\mathcal{E} \not\models\allowbreak \mathrm{AC}_{\mathrm{Set}}$) & Random / Rado graph $\mathrm{RG}$ (has the Independence Property $\mathrm{IP}$, Kihara~\cite[Example~8.3(3)]{Kihara26}; direct vertex/edge union satisfies $n$-$\mathrm{DAP}$ for all $n \ge 2$) & Henson $K_m$-free homogeneous graph $\mathcal{H}_m$ (has $\mathrm{IP}$; satisfies $(< m)$-$\mathrm{DAP}$, yet fails $m$-amalgamation on $K_{m-1}$-facets) \\
\bottomrule
\end{tabularx}\par}

In directed combinatorial geometry (\texttt{HomotopyCRM/\allowbreak Directed\allowbreak Fraisse\allowbreak SAP.lean}, 2,125 lines, 4-Layer Gate verified), binary ($n = 2$) Strong Amalgamation via \texttt{Set\allowbreak Pushout\allowbreak Square} holds both on finite oriented graded posets along closed cellular embeddings (\texttt{Oriented\allowbreak Graded\allowbreak Poset}, \texttt{Is\allowbreak Closed\allowbreak Embedding}) and on flow-acyclic complexes (\texttt{IsFlow\allowbreak Acyclic}, using unbiased equal-dimension steps $\dim(u) = \dim(v)$ in \texttt{FlowStep.\allowbreak dim\_\allowbreak eq} that keep Marcus Nicolas's pinched $2$-cube $P = \square^2 \cup_{\{10, 01\}} \square^1$, $P^{\mathrm{op}}$, and $P^{\mathrm{co}}$ all strictly flow-acyclic, even though $P^{\mathrm{op}}$ breaks strong Steiner partial duality via mixed-dimension chain coupling~\cite{Masuda26}) along two-clause flow-convex embeddings (\texttt{IsFlow\allowbreak Convex}\index{Flow-convex embedding@Flow-convex embedding (\texttt{IsFlowConvex})}: \texttt{cell\_\allowbreak mem} + \texttt{inter\_\allowbreak mem}, strictly separated across three embedding tiers in \texttt{Checks.\allowbreak layer2\_\allowbreak sep\_\allowbreak closed\_\allowbreak not\_\allowbreak flow\allowbreak Reach\allowbreak Reflecting} and \texttt{Checks.\allowbreak layer2\_\allowbreak sep\_\allowbreak flow\allowbreak Reach\allowbreak Reflecting\_\allowbreak not\_\allowbreak flow\allowbreak Convex} and closed under 1-sided cobase-change in \texttt{Amalgamated\allowbreak Poset\allowbreak Square.\allowbreak isFlow\allowbreak Reach\allowbreak Reflecting\_\allowbreak j₁}, \texttt{Amalgamated\allowbreak Poset\allowbreak Square.\allowbreak isFlow\allowbreak Convex\_\allowbreak j₁}, and \texttt{sap\_\allowbreak flow\allowbreak Convex} to block \texttt{steiner\_\allowbreak cycle\_\allowbreak obstruction}). Ascending from $n = 2$ to $n \ge 3$---machine-checked in \texttt{HomotopyCRM/\allowbreak Cubical\allowbreak Amalgamation.lean} (1,636 lines, 4-Layer Gate verified via \texttt{Punctured\allowbreak Cube n}, the quotient-free \texttt{Nat.find} minimal-support carrier \texttt{Cubical\allowbreak Amalgam D}, \texttt{cubical\_\allowbreak ndap\_\allowbreak closed\_\allowbreak oriented\_\allowbreak poset}, \texttt{hasCubical\allowbreak NDAP\_\allowbreak iff\_\allowbreak isIntersective}, \texttt{facet\allowbreak Cycle\allowbreak Diagram}, \texttt{no\_\allowbreak flow\allowbreak Acyclic\_\allowbreak cubical\_\allowbreak amalgam\_\allowbreak of\_\allowbreak three\_\allowbreak le}, and \texttt{theorem10B\_\allowbreak cubical\_\allowbreak amalgamation\_\allowbreak dichotomy})---uncovers a sharp cubical phase transition:

\begin{theorem}[Cubical Disjoint $n$-Amalgamation Dichotomy]\label{thm:directed-cubical-amalgamation}%
\lean{HomotopyCRM.facetCycleDiagram_isIntersective, HomotopyCRM.facetCycleFace_isFlowAcyclic, HomotopyCRM.facetCycleMap_isFlowConvex, HomotopyCRM.no_flowAcyclic_cubical_amalgam_of_three_le, HomotopyCRM.theorem10B_cubical_amalgamation_dichotomy}%
\leanok%
\uses{prop:matching-amalgamation,thm:fraisse-atomic-escape}%
\leavevmode
\begin{enumerate}[label=(\arabic*),leftmargin=2em]
  \item \textbf{All-$n$ Disjoint $n$-Amalgamation on $\mathbf{Poset}_{\mathrm{og,fin}}^{\mathrm{closed}}$:} Every punctured $n$-cube $D : \mathcal{P}^-(n) \to \mathbf{Poset}_{\mathrm{og,fin}}^{\mathrm{closed}}$ ($n \ge 2$) of finite oriented graded posets and closed cellular embeddings satisfying exact face intersections admits a finite oriented graded poset cocone $(j_S : D(S) \hookrightarrow Q)_{S \subsetneq [n]}$ of closed cellular embeddings satisfying global $n$-$\mathrm{DAP}$ ($j_S(D(S)) \cap j_T(D(T)) = j_{S \cap T}(D(S \cap T))$) and downward closure preservation $\mathrm{cl}_Q(j_S(x)) \simeq_o \mathrm{cl}_{D(S)}(x)$.
  \item \textbf{All-$n \ge 3$ Failure on $\mathbf{DirComp}_{\mathrm{acyc}}^{\mathrm{convex}}$:} For every $n \ge 3$, the facet-supported directed $n$-cycle $F_n$\index{Facet-supported directed n-cycle@Facet-supported directed $n$-cycle $F_n$} : $\mathcal{P}^-(n) \to \mathbf{DirComp}_{\mathrm{acyc}}^{\mathrm{convex}}$---assigning vertex $v_i$ of the directed $n$-cycle $C_n$ support $A_i = [n] \setminus \{i, i+1\}$ ($|A_i| = n - 2$) and directed edge $e_i : v_i \to v_{i+1}$ support $M_i = A_i \cup A_{i+1} = [n] \setminus \{i+1\}$ over $\mathbb{Z}/n\mathbb{Z}$---consists strictly of flow-acyclic proper faces ($\emptyset$, $*$, or \texttt{WalkingArrow}) and flow-convex embeddings, yet admits \emph{no} flow-acyclic oriented graded poset cocone $(j_S : F_n(S) \hookrightarrow Q)_{S \subsetneq [n]}$.
\end{enumerate}
\end{theorem}

The geometric origin of this sharp $n = 2$ versus $n \ge 3$ transition is contrasted side-by-side in Figure~\ref{fig:cubical-amalgamation-dichotomy}.

\begin{figure}[htbp]
\centering
\begin{tikzpicture}[
  >=Stealth,
  vtx/.style={circle,draw=CRMBlue,fill=white,inner sep=1.6pt,font=\tiny\bfseries}
]
  % Left panel: n=2 Acyclic Flow-Convex SAP
  \begin{scope}[xshift=-2.85cm]
    \node[font=\scriptsize\bfseries,color=CRMBlue] at (0,3.05) {$n = 2$: Flow-Convex SAP Holds};
    \draw[draw=black!70,fill=black!4,rounded corners=4pt] (-2.2,0.55) rectangle (0.55,2.55);
    \draw[draw=CRMBlue,fill=CRMBlue!6,rounded corners=4pt] (-0.55,0.55) rectangle (2.2,2.55);
    \fill[CRMBlue!16] (-0.55,0.6) rectangle (0.55,2.5);
    \draw[CRMBlue,thick] (-0.55,0.55) rectangle (0.55,2.55);

    \node[font=\scriptsize\bfseries] at (-1.4,2.28) {$B_1$};
    \node[font=\scriptsize\bfseries,color=CRMBlue] at (0,2.28) {$A$};
    \node[font=\scriptsize\bfseries] at (1.4,2.28) {$B_2$};

    \node[vtx] (a1) at (-0.18,1.78) {};
    \node[vtx] (a2) at (0.22,0.95) {};
    \node[vtx] (b1) at (-1.52,1.35) {};
    \node[vtx] (b2) at (1.52,1.35) {};
    \draw[->,thick,CRMBlue] (a1) -- (a2);
    \draw[->,thick,black!75] (a2) -- (b2);
    \draw[->,dashed,red!65!black] (a2) to[bend left=18] (b1);
    \draw[->,dashed,red!65!black] (b1) to[bend left=18] node[above left,font=\tiny,inner sep=1pt]{blocked} (a1);

    \node[font=\tiny,align=center,color=black!80] at (0,0.12) {\texttt{IsFlowConvex}: \texttt{cell\_mem} + \texttt{inter\_mem}\\keeps $B_1 \cup_A B_2$ strictly \texttt{IsFlowAcyclic}};
  \end{scope}

  % Divider
  \draw[CRMBlue!30] (0,0.0) -- (0,3.2);

  % Right panel: n=3 Facet Triangle F_3
  \begin{scope}[xshift=2.85cm]
    \node[font=\scriptsize\bfseries,color=CRMBlue] at (0,3.05) {$n = 3$: Facet Triangle $F_3 = \partial\Delta^2$ Cycle};
    \node[vtx,label={above:\tiny $v_0\,(A_0\!=\!\{2\})$}] (v0) at (0,2.35) {};
    \node[vtx,label={below left:\tiny $v_2\,(A_2\!=\!\{1\})$}] (v2) at (-1.35,0.72) {};
    \node[vtx,label={below right:\tiny $v_1\,(A_1\!=\!\{0\})$}] (v1) at (1.35,0.72) {};

    \draw[->,very thick,CRMBlue] (v0) -- node[above right,font=\tiny,inner sep=1.5pt]{$e_0\,(M_0\!=\!\{0,1\})$} (v1);
    \draw[->,very thick,CRMBlue] (v1) -- node[below,font=\tiny]{$e_1\,(M_1\!=\!\{1,2\})$} (v2);
    \draw[->,very thick,CRMBlue] (v2) -- node[above left,font=\tiny,inner sep=1.5pt]{$e_2\,(M_2\!=\!\{2,0\})$} (v0);

    \draw[->,thick,red!70!black] (0.35,1.25) arc (0:290:0.35);
    \node[font=\tiny\bfseries,color=red!70!black] at (0,1.25) {$C_3$};
    \node[font=\tiny,align=center,color=black!80] at (0,0.05) {Proper facets acyclic (\texttt{WalkingArrow}),\\cocone closes global $3$-cycle ($n\ge 3$)};
  \end{scope}
\end{tikzpicture}
\caption{Side-by-side geometric contrast of $n=2$ flow-convex amalgamation versus $n=3$ facet-cycle obstruction $F_3 = \partial\Delta^2$.}\label{fig:cubical-amalgamation-dichotomy}
\end{figure}

\section{Choice Calibration of Internal Coinductive Invertibility Streams}\label{subsec:coinductive-streams-choice}

Connecting Kihara's Rung~0 $\to$ Rung~1 choice threshold ($\mathrm{IDC}_\omega \equiv \mathrm{DC}_\omega$, Figure~\ref{fig:kihara-choice-spectrum}) to Chapter~\ref{ch:higher-crm}'s $\mathbf{ICaTT}$ coinductive invertibility streams~\cite{BCM26,FinsterMimram17,ORW26}, let $\mathcal{C} \in \mathbf{Clan}(\mathcal{S}_{\mathbf{CaTT}}, \mathcal{E})$ be an internal Grothendieck--Maltsiniotis weak $\omega$-category / clan model inside a Grothendieck 1-topos $\mathcal{E}$. Let $\Phi : \mathrm{Sub}_{\mathcal{E}}(\bigsqcup_{k \ge 1} C_k) \to \mathrm{Sub}_{\mathcal{E}}(\bigsqcup_{k \ge 1} C_k)$ be the monotone one-step bi-invertibility operator ($\Phi(R)_k$ selecting $k$-cells $f : x \to y$ admitting a $k$-cell $g : y \to x$ together with $(k+1)$-cells $\eta_L : g \circ_{k-1} f \Rightarrow \mathrm{id}_x$ and $\eta_R : f \circ_{k-1} g \Rightarrow \mathrm{id}_y$ lying in $R_{k+1}$), and compare the relational greatest fixed point $E_{\mathrm{rel}} := \nu\Phi$ against the witness-stream image $\mathrm{im}(\mathrm{InvStr}) := \mathrm{im}\bigl(\pi_0 : \varprojlim_{d \in \omega^{\mathrm{op}}} \mathrm{InvStage}^{(d)}_k(\mathcal{C}) \to C_k\bigr)$.

Inside Schanuel's nominal topos $\mathbf{Nom} = \mathbf{Cont}(\mathrm{Sym}(\mathbb{A}))$ ($\mathbb{A} = \mathbb{N}$), define the internal \emph{fresh-atom computad} $E[\mathbb{A}]$ across dimensions by taking two invariant $0$-cells $0, 1$ ($\mathrm{supp} = \emptyset$), one invariant root $1$-cell $u : 0 \to 1$ ($\mathrm{supp}(u) = \emptyset$) alongside atom-indexed candidate $1$-inverses $v(a_1) : 1 \to 0$ ($a_1 \in \mathbb{A}$, $\mathrm{supp}(v(a_1)) = \{a_1\}$), and at each dimension $d + 1 \ge 2$, for every injective $d$-tuple $(a_1, \dots, a_d) \in \mathbb{A}^{\#(d)}$, every stage-$d$ cancellation boundary tag $\epsilon$ on $(a_1, \dots, a_d)$, and every fresh atom $a_{d+1} \in \mathbb{A} \setminus \{a_1, \dots, a_d\}$, a pair of generating $(d+1)$-cells $\eta_\epsilon(a_1, \dots, a_{d+1})$ and $\bar{\eta}_\epsilon(a_1, \dots, a_{d+1})$ with coordinatewise $\mathrm{Sym}(\mathbb{A})$-action $\pi \cdot \eta_\epsilon(a_1, \dots, a_{d+1}) := \eta_\epsilon(\pi a_1, \dots, \pi a_{d+1})$ (so $\mathrm{supp}(g) = \{a_1, \dots, a_k\}$ of cardinality $k$ for every $k$-generator $g \neq u$).

\begin{lemma}[Polygraphic Finite-Support Equality on {$\mathrm{Free}_{\mathbf{CaTT}}(E[\mathbb{A}])$}]\label{lem:polygraphic-support}%
\lean{HomotopyCRM.cattTerm_exact_support, HomotopyCRM.cattTerm_freshCell_exact_support}%
\leanok%
\uses{thm:permutation-topos}%
Let $\mathrm{Free}_{\mathbf{CaTT}}(E[\mathbb{A}]) \in \mathbf{Clan}(\mathcal{S}_{\mathbf{CaTT}}, \mathbf{Nom})$ be the free Grothendieck--Maltsiniotis weak $\omega$-category generated by $E[\mathbb{A}]$ under Finster--Mimram's finitary polynomial $\mathbf{CaTT}$ constructors~\cite{FinsterMimram17,BCM26}. Every cell $t$ of $\mathrm{Free}_{\mathbf{CaTT}}(E[\mathbb{A}])$ is represented by a unique $\alpha$-equivalence syntax tree over a finite set of generator leaves $\mathrm{Gen}(t) \subset E[\mathbb{A}]$, and its least nominal support satisfies the exact two-sided equality
\[
  \mathrm{supp}(t) \;=\; \bigcup_{g \in \mathrm{Gen}(t)} \mathrm{supp}(g).
\]
\end{lemma}

\begin{proof}
Each $\mathbf{CaTT}$ coherence constructor $\mathrm{coh}_{\Gamma, A}[t_1, \dots, t_m]$ applies a fixed finite globular context $\Gamma$ (whose bound context variables carry trivial $\mathrm{Sym}(\mathbb{A})$-action up to $\alpha$-equivalence) to a finite substitution tuple $(t_1, \dots, t_m)$, with equivariant action $\pi \cdot \mathrm{coh}_{\Gamma, A}[t_1, \dots, t_m] = \mathrm{coh}_{\Gamma, A}[\pi \cdot t_1, \dots, \pi \cdot t_m]$~\cite{FinsterMimram17,BCM26}. Set $U(t) := \bigcup_{g \in \mathrm{Gen}(t)} \mathrm{supp}(g)$, which is finite by structural induction on $t$. For the forward inclusion $\mathrm{supp}(t) \subseteq U(t)$, any permutation $\pi \in \mathrm{Fix}(U(t))$ fixes $\mathrm{supp}(g)$ pointwise for every $g \in \mathrm{Gen}(t)$, hence satisfies $\pi \cdot g = g$ on all generator leaves and $\pi \cdot t = t$ by structural induction on $t$ (establishing simultaneously that $\mathrm{Free}_{\mathbf{CaTT}}(E[\mathbb{A}])$ is an internal object of $\mathbf{Nom}$ with least support $\mathrm{supp}(t) \subseteq U(t)$). Conversely, to prove $U(t) \subseteq \mathrm{supp}(t)$, let $a \in \mathrm{supp}(g_0)$ for some leaf $g_0 \in \mathrm{Gen}(t)$, and let $S \subset_{\mathrm{fin}} \mathbb{A}$ be any finite subset with $a \notin S$. Because $\mathbb{A}$ is infinite while $S \cup U(t)$ is finite, choose a fresh atom $b \in \mathbb{A} \setminus (S \cup U(t))$ and form the transposition $\tau := (a \; b) \in \mathrm{Fix}(S)$. Because $\mathbf{CaTT}$ imposes zero definitional term equations between polygraphic terms beyond syntactic typing and $\alpha$-renaming of bound variables in $\Gamma$~\cite{FinsterMimram17,BCM26}, syntax trees in $\mathrm{Free}_{\mathbf{CaTT}}(E[\mathbb{A}])$ have rigid generator sets: $\mathrm{Gen}(\tau \cdot t) = \{\tau \cdot g \mid g \in \mathrm{Gen}(t)\}$. On $g_0$, equivariance of least supports gives $b = \tau(a) \in \tau[\mathrm{supp}(g_0)] = \mathrm{supp}(\tau \cdot g_0)$, whereas $b \notin U(t)$ ensures $b \notin \mathrm{supp}(g)$ for every $g \in \mathrm{Gen}(t)$. Hence $\tau \cdot g_0 \notin \mathrm{Gen}(t)$, forcing $\mathrm{Gen}(\tau \cdot t) \neq \mathrm{Gen}(t)$ and therefore $\tau \cdot t \neq t$. Since $S \not\ni a$ was arbitrary, $a \in \mathrm{supp}(t)$, proving $U(t) \subseteq \mathrm{supp}(t)$.
\end{proof}

\begin{theorem}[Choice Calibration of Internal Coinductive Invertibility Streams]\label{thm:coinductive-stream-choice}\index{Coinductive invertibility streams}%
\lean{HomotopyCRM.isStepSound_coindStepRel, HomotopyCRM.surjective_toRel_of_stepSound, HomotopyCRM.coindStepRel_iff_nonempty_pureInvStream, HomotopyCRM.globularInvTree_isCoindEquiv, HomotopyCRM.isCoindEquiv_iff_nonempty_globularInvTree, HomotopyCRM.exists_freshChain_supported_of_card_le, HomotopyCRM.not_hasFiniteSupport_freshStream, HomotopyCRM.isEmpty_nominal_freshRootStream, HomotopyCRM.theorem9B_choice_calibration_of_coinductive_streams}%
\leanok%
\uses{lem:polygraphic-support,thm:ordinal-rank-completeness,thm:permutation-topos}%
\leavevmode
\begin{enumerate}[label=(\arabic*),leftmargin=2em]
  \item \textbf{$\mathrm{IDC}_\omega$ Adequacy at Rung~$\ge 1$:} If $\mathcal{E} \models \mathrm{IDC}_\omega$ ($\mathrm{DC}_\omega$), then $E_{\mathrm{rel}} = \mathrm{im}(\mathrm{InvStr})$ for every internal Grothendieck--Maltsiniotis weak $\omega$-category / clan model $\mathcal{C} \in \mathbf{Clan}(\mathcal{S}_{\mathbf{CaTT}}, \mathcal{E})$.
  \item \textbf{Rung~0 Failure on $\mathrm{Free}_{\mathbf{CaTT}}(E[\mathbb{A}])$ in $\mathbf{Nom}$:} In Schanuel's Boolean atomic topos $\mathbf{Nom} = \mathbf{Cont}(\mathrm{Sym}(\mathbb{A}))$ ($\mathbb{A} = \mathbb{N}$, validating $\mathrm{LEM} + \neg\mathrm{DC}_\omega$), evaluating $\Phi$ and $\mathrm{InvStr}_1$ on $\mathrm{Free}_{\mathbf{CaTT}}(E[\mathbb{A}]) \in \mathbf{Clan}(\mathcal{S}_{\mathbf{CaTT}}, \mathbf{Nom})$ yields $u \in \Gamma(E_{\mathrm{rel}, 1})$ by Knaster--Tarski coinduction, whereas $\mathrm{InvStr}_1(\mathrm{Free}_{\mathbf{CaTT}}(E[\mathbb{A}]))_u = \emptyset$ ($\mathrm{im}(\mathrm{InvStr})_1 \subsetneq E_{\mathrm{rel}, 1}$).
\end{enumerate}
\end{theorem}

\begin{proof}
For~(1), in any Grothendieck topos $\mathcal{E}$ without choice, unrolling streams $\mathrm{InvStr}_k(\mathcal{C}) = \varprojlim_{d \in \omega^{\mathrm{op}}} \mathrm{InvStage}^{(d)}_k(\mathcal{C})$ one step shows $\mathrm{im}(\mathrm{InvStr}) \subseteq \Phi(\mathrm{im}(\mathrm{InvStr}))$, whence $\mathrm{im}(\mathrm{InvStr}) \subseteq \nu\Phi = E_{\mathrm{rel}}$ by Tarski's fixed-point theorem on $\mathrm{Sub}_{\mathcal{E}}(\bigsqcup_{k \ge 1} C_k)$. Conversely, assume $\mathcal{E} \models \mathrm{IDC}_\omega$, and define the \emph{frontier-filtered stage subobject} $\widetilde{\mathrm{Stage}}^{(d)}(\mathcal{C}) \rightarrowtail \mathrm{InvStage}^{(d)}(\mathcal{C})$ ($d \ge 0$, with $\widetilde{\mathrm{Stage}}^{(0)}(\mathcal{C}) := E_{\mathrm{rel}}$) consisting of depth-$d$ witness trees whose $2^d$ frontier cancellation cells at depth $d$ all belong to $E_{\mathrm{rel}} = \Phi(E_{\mathrm{rel}})$. Because $E_{\mathrm{rel}} \subseteq \Phi(E_{\mathrm{rel}})$ holds internally and finite $2^d$-fold conjunctions of epimorphisms split locally in $\mathcal{E}$ without infinite choice, each transition map $\widetilde{\mathrm{Stage}}^{(d+1)}(\mathcal{C}) \twoheadrightarrow \widetilde{\mathrm{Stage}}^{(d)}(\mathcal{C})$ is an epimorphism in $\mathcal{E}$. Applying $\mathrm{IDC}_\omega$ renders the limit projection $\varprojlim_{d \in \omega^{\mathrm{op}}} \widetilde{\mathrm{Stage}}^{(d)}(\mathcal{C}) \twoheadrightarrow E_{\mathrm{rel}}$ epic; composing with $\varprojlim_d \widetilde{\mathrm{Stage}}^{(d)}(\mathcal{C}) \rightarrowtail \mathrm{InvStr}(\mathcal{C})$ gives $E_{\mathrm{rel}} \subseteq \mathrm{im}(\mathrm{InvStr})$.

For~(2), set $\mathcal{F} := \mathrm{Free}_{\mathbf{CaTT}}(E[\mathbb{A}]) \in \mathbf{Clan}(\mathcal{S}_{\mathbf{CaTT}}, \mathbf{Nom})$ (internalized by Lemma~\ref{lem:polygraphic-support}). Define the graded $\mathrm{Sym}(\mathbb{A})$-invariant subsheaf $R \subseteq \bigsqcup_{k \ge 1} \mathcal{F}_k$ spanned by $u$ and $\{v(a_1)\}_{a_1 \in \mathbb{A}}$ at $k = 1$, and by $\{\eta_\epsilon(a_1, \dots, a_k), \bar{\eta}_\epsilon(a_1, \dots, a_k)\}_{(a_1, \dots, a_k) \in \mathbb{A}^{\#(k)}}$ at each $k \ge 2$. For any generator $c \in R_d$ with finite support $S = \mathrm{supp}(c)$ ($|S| \le d$), picking any fresh atom $a_{d+1} \in \mathbb{A} \setminus S$ refines $\mathrm{Fix}(S)$ to the open subgroup $\mathrm{Fix}(S \cup \{a_{d+1}\}) \le_{\mathrm{open}} \mathrm{Sym}(\mathbb{A})$, over which the partner inverse and depth-$(d+1)$ cancellation witnesses $\eta_{\epsilon'}(a_1, \dots, a_{d+1}) \in R_{d+1}$ are fixed. By Kripke--Joyal semantics on the atomic site of $\mathbf{Nom}$, $\mathbf{Nom} \models R \subseteq \Phi(R)$; hence $R \subseteq \nu\Phi = E_{\mathrm{rel}}$ and $u \in \Gamma(R_1) \subseteq \Gamma(E_{\mathrm{rel}, 1})$.

Suppose now for contradiction that $\mathrm{InvStr}_1(\mathcal{F})_u \in \mathbf{Nom}$ contains an element $\xi = (\xi_d)_{d \ge 1}$ supported by some finite subset $F \subset_{\mathrm{fin}} \mathbb{A}$ of cardinality $n := |F| < \omega$. Equivariance of the stage-$(n+1)$ projection $\mathrm{pr}_{n+1}$ gives $\mathrm{supp}(\xi_{n+1}) \subseteq F$. By Lemma~\ref{lem:polygraphic-support}, every generator leaf $g \in E[\mathbb{A}]$ occurring in $\xi_{n+1}$ satisfies $\mathrm{supp}(g) \subseteq F$, hence $|\mathrm{supp}(g)| \le n$. Since every generator of $E[\mathbb{A}]$ in dimension $k \ge n + 1$ has $|\mathrm{supp}| = k \ge n + 1 > n$, $\xi_{n+1}$ lies inside the sub-clan $\mathrm{Free}_{\mathbf{CaTT}}(E_{\le n}[F])$ generated exclusively in dimensions $\le n$. Apply the canonical strictification clan morphism $q_n : \mathrm{Free}_{\mathbf{CaTT}}(E_{\le n}[F]) \to \mathcal{S}_n := \mathrm{Str}_n(E_{\le n}[F])$ into the free strict $n$-category on the $n$-computad $E_{\le n}[F]$ (sending all unbiased $\mathbf{CaTT}$ coherence operations to strict composites and all cells above dimension $n$ to strict identities~\cite[\S 3]{FinsterMimram17}, \cite{BCM26}). On any free strict $k$-category $\mathcal{S}_k = \mathrm{Str}_k(E_{\le k}[F])$ ($1 \le k \le n$), sending identity $k$-cells to $0$, generating $k$-cells to $1$, and compositions $x \circ_i y$ ($0 \le i < k$) to $\mathrm{wt}_k(x) + \mathrm{wt}_k(y)$ defines a strict homomorphism $\mathrm{wt}_k : (\mathcal{S}_k)_k \to (\mathbb{N}, +)$ (invariant under associativity, unitality, functoriality of identities, and middle-four interchange) with $\mathrm{wt}_k(z) = 0 \iff z = \mathrm{id}$. Because $(n+1)$-cells in $\mathcal{S}_n$ are identities, any top $n$-cell witness $w_n$ in $q_n(\xi_{n+1})$ satisfies $w_n' \circ_{n-1} w_n = \mathrm{id}$, forcing $\mathrm{wt}_n(w_n') + \mathrm{wt}_n(w_n) = 0$ and $w_n = \mathrm{id}$; descending by downward induction from $k = n$ to $k = 1$ collapses all witness cells in dimensions $2, \dots, n$ to identities and forces $w_1 \circ_0 u = \mathrm{id}_0$ in $\mathcal{S}_1$, contradicting $\mathrm{wt}_1(w_1) + \mathrm{wt}_1(u) \ge 1 \neq 0$ (or $\mathrm{Hom}_{\mathcal{S}_0}(1, 0) = \emptyset$ when $n = 0$). Equivalently, every stage-$(n+1)$ witness $\xi_{n+1}$ satisfies $\mathrm{supp}(\xi_{n+1}) \not\subseteq F$: picking $a_m \in \mathrm{supp}(\xi_{n+1}) \setminus F$ and $b \in \mathbb{A} \setminus (F \cup \mathrm{supp}(\xi_{n+1}))$ yields $(a_m \; b) \in \mathrm{Fix}(F)$ with $(a_m \; b) \cdot \xi \neq \xi$, proving $\mathrm{InvStr}_1(\mathcal{F})_u = \emptyset$.
\end{proof}

Theorem~\ref{thm:coinductive-stream-choice} and Lemma~\ref{lem:polygraphic-support} are machine-checked in Lean~4 (\texttt{HomotopyCRM/\allowbreak Nominal\allowbreak ICaTT.lean}\index{NominalICaTT@\texttt{NominalICaTT} Lean~4 module}, 1,524 lines, 21 \texttt{\#print axioms} checks under the 4-Layer Definition Check Gate): classical Dependent Choice (\texttt{Classical.choice}) proves surjectivity of \texttt{InvStream.\allowbreak toRel} onto any reflexive step-sound relation (\texttt{surjective\_\allowbreak toRel\_\allowbreak of\_\allowbreak step\allowbreak Sound}) and globular coinductive tree adequacy (\texttt{globular\allowbreak InvTree\_\allowbreak isCoind\allowbreak Equiv}), while over the nominal atom permutation action on $\mathbb{A} = \mathbb{N}$, finitary \texttt{CaTTTerm} syntax trees obey two-sided minimal support equality (\texttt{cattTerm\_\allowbreak exact\_\allowbreak support}, \texttt{cattTerm\_\allowbreak freshCell\_\allowbreak exact\_\allowbreak support}), every finite depth-$n$ truncation \texttt{Fresh\allowbreak Chain n} splits locally over any $n$-atom support (\texttt{exists\_\allowbreak fresh\allowbreak Chain\_\allowbreak supported\_\allowbreak of\_\allowbreak card\_\allowbreak le} and \texttt{hasFinite\allowbreak Support\_\allowbreak fresh\allowbreak Chain}), and no infinite fresh-atom stream \texttt{Fresh\allowbreak Stream} or root witness stream \texttt{Fresh\allowbreak Root\allowbreak Stream} possesses finite support (\texttt{not\_\allowbreak hasFinite\allowbreak Support\_\allowbreak fresh\allowbreak Stream}, \texttt{isEmpty\_\allowbreak nominal\_\allowbreak fresh\allowbreak Root\allowbreak Stream}, and \texttt{theorem9B\_\allowbreak choice\_\allowbreak calibration\_\allowbreak of\_\allowbreak coinductive\_\allowbreak streams}).

\IfFileExists{ch07_sets_residuum.tex}{\chapter{\texorpdfstring{Residuums and ``$\mathbf{Sets}$''}{Residuums and "Sets"}}
\chaptermark{Residuums and ``Sets''}
\label{ch:sets-residuum}
\chapterepigraph{``The inequalities in the world are like the combinations in a kaleidoscope; at every turn a fresh picture strikes the eye; and yet, in reality, you see only the same bits of glass as you saw before.''}{Arthur Schopenhauer, \emph{Counsels and Maxims}, \S 47}
\chapterepigraph{``Sometimes I have had the task to write reviews of articles, where set theoretic notions and theorems are used without any explanation of what kind of set theory is meant. It is a disagreeable job to write reviews in such cases. One does not know what the author really means.''}{Thoralf Skolem, \emph{A Critical Remark on Foundational Research} (1955)~\cite{Skolem55}}

Throughout Chapters~\ref{ch:foundations}--\ref{ch:realizability-sites}, Categorical Reverse Mathematics ($\mathrm{CRM}$) has calibrated mathematical necessity \emph{inside} variable universes of small maps~\cite{Lawvere69,Johnstone02}. Across all Five Foundational Coordinate Axes established in Chapter~\ref{ch:foundations} (Definitions~\ref{def:comparison-triple}--\ref{def:residuum} and Table~\ref{tab:five-axis-master})---Choice/Splitting (\textbf{Axis~I}), Constructive Logic (\textbf{Axis~II}), Arithmetic Infinity (\textbf{Axis~III}), Characteristic/$p$-Torsion (\textbf{Axis~IV}), and Directional Invertibility (\textbf{Axis~V})---doctrinal comparison triples $\mathfrak{T} = (\mathcal{D}_0 \supseteq \mathcal{D}_1,\; \kappa : \Phi \Rightarrow \Psi) \in \mathrm{Dom}(\mathrm{Res})$ were measured through concrete internal invariants: non-abelian torsor cohomology $H^1(\mathcal{E}/Y; G)$, derived limit phantoms $\varprojlim^1$, Heyting boundary lattices, $\delta$-ring lift torsors, and Postnikov obstructions. Throughout those calibrations, the background doctrine of ``$\mathbf{Sets}$'' remained quietly offstage---even in the face of Thoralf Skolem's warning~\cite{Skolem55} against unspecified background set theories---shielded by the folkloric consensus that Zermelo--Fraenkel set theory ($\mathsf{ZFC}$) and von Neumann--Bernays--G\"odel class theory ($\mathsf{NBG}$) are functionally interchangeable ways of naming the same mathematical bedrock.

Turning the residuum operator $\mathrm{Res}(\kappa)$ onto ``$\mathbf{Sets}$'' itself---via Andr\'e Joyal and Ieke Moerdijk's \textbf{Algebraic Set Theory}\index{Algebraic Set Theory} of class categories equipped with small-map ideals $(\mathcal{C}, \mathcal{S})$~\cite{JoyalMoerdijk95}---shatters that consensus along two orthogonal directions. First, holding an underlying universe of small sets completely rigid still leaves an entire vertical kaleidoscope of $2$-categorical class residua overhead. Second, and far more subversively, weakening the horizontal set axioms below full $\mathsf{ZFC}$ breaks class-to-set conservativity outright---exposing the structural plot twist that Lawvere's Elementary Theory of the Category of Sets ($\mathsf{ETCS}$)\index{ETCS@$\mathsf{ETCS}$ (Elementary Theory of the Category of Sets)}~\cite{Lawvere64}, the Boolean elementary-topos baseline inhabited across Chapters~\ref{ch:foundations}--\ref{ch:realizability-sites}, is bi-interpretable with Bounded Zermelo / Mac~Lane set theory ($\mathsf{Mac}$)~\cite{Cole73,Mitchell72,Osius74} rather than $\mathsf{ZFC}$ and has quietly operated without Replacement (and even without unbounded $\mathcal{L}_\in$-Separation) from page one.

\section{\texorpdfstring{The Conservativity Illusion and the Fixed-Set Kaleidoscope}{The Conservativity Illusion and the Fixed-Set Kaleidoscope}}\label{sec:conservativity-kaleidoscope}
\sectionepigraph{``Our present assumption of `one universe' is an adequate stopgap, not a forecast of the future.''}{Saunders Mac~Lane, \emph{Categories for the Working Mathematician}~\cite[Ch.~I, \S 6]{MacLane98}}

Despite Saunders Mac~Lane's caution against treating a single universe as final~\cite[Ch.~I, \S 6]{MacLane98}, foundational folklore routinely dismisses the gap between set theory and class theory on the strength of Ulrich Felgner's 1971 class-forcing theorem~\cite{Felgner71}: $\mathsf{NBG}$\index{NBG@$\mathsf{NBG}$ (von Neumann--Bernays--G\"odel class theory)}---formulated as Bernays--G\"odel class theory $\mathsf{GBC}$ augmented with Global Choice ($\mathsf{GC}$)---proves the exact same first-order $\mathcal{L}_\in$-sentences about \emph{sets} as $\mathsf{ZFC}$. Proof-theoretically, $\mathsf{NBG}$ stands to $\mathsf{ZFC}$ just as arithmetical comprehension $\mathrm{ACA}_0$ stands to Peano Arithmetic $\mathsf{PA}$ in classical Reverse Mathematics~\cite{Simpson09}.

The catch is that \textbf{first-order conservativity is strictly weaker than bi-interpretability} or categorical Morita equivalence~\cite{Caramello18}. To track class-to-set projections across both $\mathsf{ZFC}$ and weaker Replacement-free bases without circularity, we distinguish Replacement-free class algebra from full Bernays--G\"odel class theory at the outset.

\begin{definition}[Class-Set Model Spaces and the Baseline Adjunction $\mathrm{Def} \dashv \pi_{\mathrm{Set}}$]\label{def:set-projection-fibration}%
\lean{HomotopyCRM.SetsResiduum.defSetProjAdjunction, HomotopyCRM.SetsResiduum.finToPowElemHom, HomotopyCRM.SetsResiduum.finToPowClassHom, HomotopyCRM.SetsResiduum.naive_pullback_breaks_subclass_separation}%
\leanwitness%
\uses{def:comparison-triple,def:residuum}%
Let $\mathsf{GB}_0$\index{GB0@$\mathsf{GB}_0$ (Replacement-free class theory)} denote \emph{Replacement-free Bernays--G\"odel class theory}: the two-sorted $(\mathcal{L}_\in, \mathcal{C})$-theory over Zermelo set theory $\mathsf{Z}$ (Empty Set, Pairing, Union, Power Set, Separation, Infinity, Foundation, and Extensionality) whose class sort $\mathcal{C}$ is closed under G\"odel's eight finite predicative class operations \textbf{G1--G8} together with Subclass Separation ($C \subseteq A \in \mathcal{S} \implies C \in \mathcal{S}$, inducing full $\mathcal{L}_\in$-Separation on sets), \emph{omitting} Bernays--G\"odel Axiom~D (Class Replacement). Write:
\[
  \begin{aligned}
    \mathsf{GBC}_0 &:= \mathsf{GB}_0 + \mathrm{AC}, &\qquad
    \mathsf{GB} &:= \mathsf{GB}_0 + \mathrm{Rep}, \\
    \mathsf{GBC} &:= \mathsf{GB} + \mathrm{AC}, &\qquad
    \mathsf{NBG} &:= \mathsf{GBC} + \mathsf{GC}.
  \end{aligned}
\]
Following Definitions~\ref{def:comparison-triple}--\ref{def:residuum} of Chapter~\ref{ch:foundations}, let $\mathbf{Mod}(\mathsf{GB}_0)$ denote the category whose objects are two-sorted models $(M, \mathcal{C}) \models \mathsf{GB}_0$ and whose $1$-morphisms $(f, F) : (M, \mathcal{C}) \to (N, \mathcal{D})$ consist of an $\mathcal{L}_\in$-elementary embedding $f : M \hookrightarrow N$ paired with a $\Sigma^1_0$-elementary class map $F : \mathcal{C} \to \mathcal{D}$ preserving singleton parameters ($\{p\} \mapsto \{f(p)\}$), G\"odel operations \textbf{G1--G8}, and membership ($m \in C \iff f(m) \in F(C)$, restricting over $f = \mathrm{id}_M$ to fixed-set class inclusions $\mathcal{C} \subseteq \mathcal{C}'$), and let $\mathbf{Mod}(\mathsf{Z})$ be the category of models of $\mathsf{Z}$ and $\mathcal{L}_\in$-elementary embeddings.
\begin{enumerate}[label=\textup{(\roman*)},leftmargin=1.8em]
  \item The forgetful functor $\pi_{\mathrm{Set}} : \mathbf{Mod}(\mathsf{GB}_0) \twoheadrightarrow \mathbf{Mod}(\mathsf{Z})$ sending $(M, \mathcal{C}) \mapsto M$ is surjective on objects, with vertical fiber $\pi_{\mathrm{Set}}^{-1}(M)$ over $\mathrm{id}_M$ consisting of the poset of all $\mathsf{GB}_0$-realizations $(M, \mathcal{C})$ sharing the exact same underlying set universe $M$. (Note that naive class pullback $f^*\mathcal{D} := \{f^{-1}[D] \mid D \in \mathcal{D}\}$ along a proper $\mathcal{L}_\in$-elementary embedding $f : M \hookrightarrow V_\kappa$ does \emph{not} make $\pi_{\mathrm{Set}}$ a model-space Grothendieck fibration, as any small set $r \in \mathcal{P}(\omega)^{V_\kappa} \setminus M$ over countable $M \prec V_\kappa$ satisfies $f^{-1}[r] = r \notin M$, violating Subclass Separation while failing to commute with \textbf{G4} domain projection on $\{(n, r) \mid n \in r\}$; syntactically, however, deductive theory pairs $(T_{\mathrm{Set}}, T_{\mathrm{Class}}) \mapsto T_{\mathrm{Set}}$ do form a base-change opfibration via $T_{\mathrm{Class}} \mapsto T_{\mathrm{Class}} + \sigma$.)
  \item Sending each set model $M \models \mathsf{Z}$ to its first-order parameter-definable class expansion $\mathrm{Def}(M)$ defines a full coreflective baseline embedding $\mathrm{Def} : \mathbf{Mod}(\mathsf{Z}) \hookrightarrow \mathbf{Mod}(\mathsf{GB}_0)$ ($\pi_{\mathrm{Set}} \circ \mathrm{Def} = \mathrm{Id}_{\mathbf{Mod}(\mathsf{Z})}$) left adjoint to $\pi_{\mathrm{Set}}$:
  \[
    \mathrm{Def} \;\dashv\; \pi_{\mathrm{Set}} : \mathbf{Mod}(\mathsf{GB}_0) \twoheadrightarrow \mathbf{Mod}(\mathsf{Z}).
  \]
  Indeed, for any $(N, \mathcal{D}) \models \mathsf{GB}_0$ and $\mathcal{L}_\in$-elementary $f : M \hookrightarrow N$, $\mathcal{L}_\in$-elementarity and G\"odel's Class Existence Theorem send each definable class $\varphi^M(-, p) \in \mathrm{Def}(M)$ to a unique well-defined class $\varphi^N(-, f(p)) \in \mathrm{Def}(N) \subseteq \mathcal{D}$. Over $\mathbf{Mod}(\mathsf{ZFC}) \hookrightarrow \mathbf{Mod}(\mathsf{ZC})$, $\mathrm{Def} \dashv \pi_{\mathrm{Set}}$ restricts to $\mathbf{Mod}(\mathsf{GBC}) \twoheadrightarrow \mathbf{Mod}(\mathsf{ZFC})$.
\end{enumerate}
\end{definition}

Suppose we clamp an underlying countable model $M \models \mathsf{ZFC}$ completely rigid, freezing every set-level $\mathsf{ZFC}$ residuum to zero so that not a single element, ordinal, real number, or small morphism can vary. Even above this frozen set floor, the vertical fiber $\pi_{\mathrm{Set}}^{-1}(M)$ opens into an intricate second-order poset governing which proper classes $\mathcal{C}$ exist above $M$.

\begin{theorem}[{External Benchmark: Fixed-Set Class Kaleidoscope; Williams~2018~\cite[Ch.~3, \S 3.2]{Williams18}}]\label{thm:williams-kaleidoscope}%
\uses{def:set-projection-fibration}%
Fix any countable model $M \models \mathsf{ZFC}$. The partial order of $\mathsf{GBC}$-realizations $(M, \mathcal{C})$ having the exact same underlying set universe $M$, ordered by class inclusion $\mathcal{C} \subseteq \mathcal{C}'$, embeds every countable partial order.
\end{theorem}

Standing inside a fixed model $M \models \mathsf{ZFC}$ while varying $\mathcal{C}$ realizes Schopenhauer's kaleidoscope literally: the set-level bits of glass never move, yet turning through admissible class skies $(M, \mathcal{C})$ reconfigures which large categorical constructions succeed or obstruct.

Recall how Choice appears categorically at the small scale: complementing Chapter~\ref{ch:choice-atlas}'s epimorphism-splitting Choice-Geometry Atlas along Axis~I, Michael Makkai~\cite{Makkai96} observed that inside the $2$-category $\mathbf{Cat}(\mathcal{S})$ of small categories over $\mathbf{Set}$, $\mathrm{AC}$ is equivalent to requiring every \emph{weak categorical equivalence} (every functor that is fully faithful and essentially surjective on objects) to admit a strong quasi-inverse functor without passing to anafunctors. Crucially, the very first step up our vertical class fiber ($\mathsf{GBC} \subsetneq \mathsf{NBG}$) is this exact equivalence-splitting residuum shifted one universe upward from internal small categories $\mathbf{Cat}(\mathcal{S})$ to large categories $\mathbf{Cat}(\mathcal{C})$.

\begin{proposition}[Large Equivalence-Splitting Calibrates Global Choice]\label{prop:large-equiv-splitting}%
\lean{HomotopyCRM.SetsResiduum.card_split_iff_global_choice, HomotopyCRM.SetsResiduum.type_equiv_splitting_global_choice}%
\leanwitness%
\uses{def:set-projection-fibration}%
Work over $\mathsf{GBC}$\index{Global Choice ($\mathsf{GC}$)} in a class-set model $(M, \mathcal{C})$, and let $I : \mathbf{Card} \hookrightarrow \mathbf{Set}$ denote the full inclusion in $\mathbf{Cat}(\mathcal{C})$ of the skeletal subcategory of small von Neumann cardinals into small sets. Then:
\begin{enumerate}[label=\textup{(\roman*)},leftmargin=1.8em]
  \item $I : \mathbf{Card} \hookrightarrow \mathbf{Set}$ is a weak categorical equivalence (fully faithful and essentially surjective on objects).
  \item The existence of a quasi-inverse class functor $F : \mathbf{Set} \to \mathbf{Card}$ equipped with a class natural isomorphism $\eta : \mathrm{Id}_{\mathbf{Set}} \xrightarrow{\sim} I \circ F$ is strictly equivalent over $\mathsf{GBC}$ to Global Choice ($\mathsf{NBG} := \mathsf{GBC} + \mathsf{GC}$).
  \item Whenever $M \models \mathsf{ZFC} + \forall p\,(V \neq \mathrm{HOD}(\{p\}))$, as obtained by homogeneity of tail Easton class-forcing extensions~\cite{Easton70}, the minimal baseline fiber $(M, \mathrm{Def}(M)) \models \mathsf{GBC}$ violates $\mathsf{GC}$ and therefore leaves $I : \mathbf{Card} \hookrightarrow \mathbf{Set}$ unsplit. (Bare $V \neq \mathrm{HOD}$ does not suffice: e.g.\ adding a single Cohen real $c$ over $L$ gives $L[c] \models V \neq \mathrm{HOD}$, yet the parameter-definable well-ordering $<_{L[c]} \in \mathrm{Def}(L[c])$ still splits $I$.)
\end{enumerate}
\end{proposition}

\begin{proof}
Under set-level $\mathrm{AC}$ and Replacement in $\mathsf{GBC}$, every small set $X \in \mathbf{Set}$ admits a bijection with a unique von Neumann initial ordinal $|X| \in \mathbf{Card}$. Because $\mathbf{Card}$ is a full subcategory of $\mathbf{Set}$, $I$ is fully faithful and essentially surjective, proving~(i).

For~(ii), suppose first that $I$ splits via a class functor $F : \mathbf{Set} \to \mathbf{Card}$ and a class natural isomorphism $\eta : \mathrm{Id}_{\mathbf{Set}} \xrightarrow{\sim} I \circ F$. Since cardinals are skeletal, $F(X) = |X|$ for every set $X$. Whenever $X \neq \emptyset$, its cardinal $|X|$ is a nonzero von Neumann ordinal and hence satisfies $\emptyset \in |X|$. Because $\eta_X : X \xrightarrow{\sim} |X|$ is a bijection for each set $X$, the predicative class assignment
\[
  G(X) \;:=\; \eta_X^{-1}(\emptyset) \;\in\; X
\]
selects a uniform element of every nonempty small set $X \in \mathcal{S}$. Over $\mathsf{GBC}$, Foundation decomposes $V = \bigcup_{\alpha \in \mathrm{Ord}} V_\alpha$, so setting $\tilde{G}(C) := G(C \cap V_{\alpha_0})$ for $\alpha_0 := \min\{\alpha \in \mathrm{Ord} \mid C \cap V_\alpha \neq \emptyset\}$ immediately extends $G$ to a class-form Global Choice selector on all nonempty classes $C \in \mathcal{C}$, establishing $\mathsf{GC}$. Conversely, assume $\mathsf{GC}$. Applying global choice to the class family of nonempty bijection sets $\{\mathrm{Bij}(X, |X|)\}_{X \in \mathbf{Set}}$ simultaneously selects one canonical bijection $\eta_X : X \xrightarrow{\sim} |X|$ for every small set $X$. Defining $F(X) := |X|$ on objects and $F(f) := \eta_Y \circ f \circ \eta_X^{-1} : |X| \to |Y|$ on morphisms $f : X \to Y$ assembles $F : \mathbf{Set} \to \mathbf{Card}$ into a genuine class functor with natural isomorphism $\eta : \mathrm{Id}_{\mathbf{Set}} \xrightarrow{\sim} I \circ F$.

Finally, for~(iii), any Global Choice selector $G \in \mathrm{Def}(M)$ definable over $M$ with a set parameter $p \in M$ induces by transfinite recursion over each $V_\alpha$ a well-ordering of $V$ definable from $p$ and ordinals, forcing $M \models V = \mathrm{HOD}(\{p\})$. Notice that bare $V \neq \mathrm{HOD}$ allows parameter-definable well-orderings (for instance, $<_{L[c]} \in \mathrm{Def}(L[c])$ for a single Cohen real $c$ over $L$, where $L[c] \models V \neq \mathrm{HOD}$ yet $L[c] = \mathrm{HOD}(\{c\})$ still splits $I : \mathbf{Card} \hookrightarrow \mathbf{Set}$). By homogeneity of tail Easton class-forcing extensions~\cite{Easton70} (where every set parameter $p$ lies in an initial segment and is fixed by automorphisms of the tail forcing), models satisfying $M \models \mathsf{ZFC} + \forall p\,(V \neq \mathrm{HOD}(\{p\}))$ exist; in any such $M$, no parameter-definable global choice selector lies in $\mathrm{Def}(M)$, so $I : \mathbf{Card} \hookrightarrow \mathbf{Set}$ cannot split in $(M, \mathrm{Def}(M))$.
\end{proof}

Climbing above $\mathsf{NBG}$ along the vertical fiber calibrates the implicit class commitments invoked across modern $2$-category theory and homotopy theory:
\begin{itemize}[leftmargin=1.5em]
  \item \textbf{Class Forcing and Transfinite Ad\'amek Chains ($\mathsf{NBG} \subsetneq \mathsf{NBG} + \mathsf{ETR}_\omega \subsetneq \mathsf{NBG} + \mathsf{ETR}_{\mathrm{Ord}} \subsetneq \mathsf{NBG} + \mathsf{ETR}$):} Gitman, Hamkins, Holy, Schlicht, and Williams~\cite[Thm.~1.2]{Gitman20} (and Williams~\cite{Williams18}; note that Gitman et al.\ write ``$\mathsf{GBC}$'' for $\mathsf{NBG} = \mathsf{GBC} + \mathsf{GC}$) proved that the Class Forcing Theorem (definability of $\Vdash_{\mathbb{P}}$ and the Truth Lemma for proper-class forcing posets) is equivalent over $\mathsf{NBG} = \mathsf{GBC} + \mathsf{GC}$ to Elementary Transfinite Recursion of length $\mathrm{Ord}$ ($\mathsf{ETR}_{\mathrm{Ord}}$), failing in $\mathsf{NBG}$---and likewise encompassing $\mathrm{Ord}$-length Ad\'amek chains~\cite{Goldthorpe23} when encoded via first-order class step relations. By contrast, \emph{full $\mathsf{ETR}$} along arbitrary class well-orderings is strictly stronger, proving $\mathrm{Con}(\mathsf{NBG} + \mathsf{ETR}_{\mathrm{Ord}})$~\cite{Gitman20}, whereas constructing full Tarskian truth predicates relative to arbitrary class parameters holds already at $\mathsf{ETR}_\omega$, which already proves $\mathrm{Con}(\mathsf{ZFC})$ and triggers a syntactic $\mathcal{L}_\in$-reflection leak over $\mathsf{ZFC}$ while remaining $\mathcal{L}_\in$-conservative over its first-order trace $T^* := \mathrm{Th}_{\mathcal{L}_\in}(\mathsf{NBG} + \mathsf{ETR}_{\mathrm{Ord}})$ (strictly below full $\mathsf{ETR}$ and impredicative $\mathsf{KM}$).
  \item \textbf{Impredicative Class Comprehension vs.\ Grothendieck Universes ($\mathsf{NBG} + \mathsf{ETR}_{\mathrm{Ord}} \subsetneq \mathsf{NBG} + \mathsf{ETR} \subsetneq \mathsf{NBG} + \Pi^1_1\text{-}\mathsf{CA} \subsetneq \mathsf{KM}$):}\index{Kelley--Morse class theory ($\mathsf{KM}$)} By Freyd's theorem~\cite[Ch.~V, \S 2, Prop.~3]{MacLane98}, no small category admitting products indexed by its own morphism size can avoid collapsing to a preorder; hence $\mathbf{Set}$ cannot admit arbitrary proper-class-indexed limits, colimits, or unrestricted pointwise large Kan extensions even in Kelley--Morse ($\mathsf{KM}$), which instead require passing to a larger Grothendieck universe $V_\lambda \ni V_\kappa$ ($\kappa < \lambda$ inaccessible) as the ambient ``$\mathbf{Set}$''. Rather, $\mathsf{KM}$ calibrates full impredicative ($\Pi^1_\infty$) class comprehension over $V$---supporting complete class-lattice meets/joins and impredicative second-order categorical invariants over $V$---driving the syntactic reflection leak far past $T^*$ to $\mathrm{Th}_{\mathcal{L}_\in}(\mathsf{KM})$ (proving $\mathrm{Con}(\mathsf{NBG} + \mathsf{ETR})$ and hence $\mathrm{Con}(\mathsf{NBG} + \mathsf{ETR}_{\mathrm{Ord}})$ as modeled by two-sorted $(V_\kappa, V_{\kappa+1})$ for inaccessible $\kappa$), mirroring the classical climb from $\mathrm{ACA}_0$ through $\mathrm{ATR}_0$ and $\Pi^1_1\text{-}\mathrm{CA}_0$ to full second-order arithmetic $\mathsf{Z}_2$~\cite{Simpson09}.
\end{itemize}

\section{\texorpdfstring{Cracks in the Floor: Why ``$\mathbf{Sets}$'' Cannot Hold the Base}{Cracks in the Floor: Why "Sets" Cannot Hold the Base}}\label{sec:cracks-in-the-floor}

At this juncture, a pragmatic category theorist might still hope to quarantine the entire phenomenon upstairs: let specialists pay vertical class tolls in $\mathsf{NBG}$, $\mathsf{ETR}_{\mathrm{Ord}}$, or $\mathsf{KM}$ whenever they split large equivalences, iterate $\mathrm{Ord}$-chains, or evaluate impredicative class comprehensions (moving to higher Grothendieck universes whenever large limits or Kan extensions are formed), while ordinary small-map mathematics proceeds undisturbed inside a canonical ground floor of ``$\mathbf{Sets}$.''

Yet our \textbf{Base--Fiber Conservation Principle} (Chapter~\ref{ch:foundations}, Principle~\ref{prin:conservation}, and Chapter~\ref{ch:base-fiber}, Principle~\ref{thm:conservation-law}) prompts a pointed structural question about Felgner's 1971 conservativity theorem~\cite{Felgner71}: \emph{what prevented proper-class Global Choice from leaking down into set-level theorems in the first place?} Inspecting Felgner's forcing proof reveals two load-bearing set axioms acting as an insulating gasket underneath $\pi_{\mathrm{Set}}$: \textbf{Replacement}, which keeps the class forcing of set-sized partial well-orderings of $V$ $\kappa$-closed for every cardinal $\kappa$ so that no new small sets are added, and \textbf{Foundation}, which anchors cumulative rank induction across $V$.

Syntactic $\mathcal{L}_\in$-conservativity across $\pi_{\mathrm{Set}}$ is consequently not an intrinsic harmlessness of proper classes; it is a \emph{local flatness property} over $\mathbf{Mod}(\mathsf{ZFC})$ that shatters on wider small-set bases where Replacement fails:
\[
\begin{tikzcd}[row sep=2.2em, column sep=2.8em, font=\footnotesize]
\mathbf{Mod}(\mathsf{NBG})
  \arrow[r, hook, "\tilde{\iota}"]
  \arrow[d, "\substack{\pi_{\mathrm{Set}}|_{\mathsf{ZFC}}\\ \text{($\mathcal{L}_\in$-conserv.: Felgner 1971)}}"']
&
\mathbf{Mod}(\mathsf{GBC}_0 + \mathsf{GC})
  \arrow[d, "\substack{\pi_{\mathrm{Set}}|_{\mathsf{ZC}}\\ \text{(non-conserv.: Glazer 2023)}}"]
\\
\mathbf{Mod}(\mathsf{ZFC})
  \arrow[r, hook, "\iota"]
&
\mathbf{Mod}(\mathsf{ZC})
\end{tikzcd}
\]
Here horizontal inclusions run contravariantly to syntactic theory strength ($\mathsf{ZC} \subsetneq \mathsf{ZFC}$ induces $\mathbf{Mod}(\mathsf{ZFC}) \hookrightarrow \mathbf{Mod}(\mathsf{ZC})$). Moreover, unlike $\pi_{\mathrm{Set}} : \mathbf{Mod}(\mathsf{GBC}) \twoheadrightarrow \mathbf{Mod}(\mathsf{ZFC})$ (surjective via $\mathrm{Def}(M)$), neither vertical $\mathsf{GC}$-projection above is surjective on all models: $\pi_{\mathrm{Set}}|_{\mathsf{ZC}}$ omits $\mathsf{ZC}$-models of $\theta_{\mathrm{Glz}}$ by Theorem~\ref{thm:glazer-nonconservativity}, and $\pi_{\mathrm{Set}}|_{\mathsf{ZFC}}$ fails surjectivity on uncountable rather classless models of Easton's theory (cf.\ Shelah~\cite{Shelah78}). As soon as $\mathrm{CRM}$ relaxes the set base below $\mathsf{ZFC}$, class principles stop tracking the set base along three independent fault lines.

\begin{remark}[Set-Theoretic Notation Alignment: $\mathsf{Mac}$ vs.\ $\mathsf{ZC}$ and $\mathsf{ZFC}^-$ vs.\ $\mathsf{ZFC}_{\neg\mathrm{Reg}}$]\label{rem:zfc-minus-notation}%
\uses{def:set-projection-fibration}%
Two standard terminology collisions must be fixed cleanly before dissecting weaker set bases:
\begin{enumerate}[label=\textup{(\alph*)},leftmargin=1.8em]
  \item \textbf{Bounded Zermelo ($\mathsf{Mac}$) versus Full Zermelo ($\mathsf{ZC}$):}\index{Mac Lane set theory ($\mathsf{Mac}$)} Following Adrian Mathias~\cite{Mathias01}, $\mathsf{ZC}$ denotes \emph{full Zermelo set theory with Choice}, including the \emph{unbounded} first-order $\mathcal{L}_\in$-Separation schema, whereas \emph{Bounded Zermelo / Mac~Lane set theory} ($\mathsf{Mac}$, also denoted $\mathsf{ZBQC}$) restricts Separation to bounded ($\Delta_0$) formulas. Crucially, the Cole--Mitchell--Osius equivalence~\cite{Cole73,Mitchell72,Osius74} (\cite[Ch.~VI, \S 10]{MacLaneMoerdijk92}) proves that Lawvere's $\mathsf{ETCS}$~\cite{Lawvere64} is bi-interpretable with $\mathsf{Mac}$ ($\Delta_0$-Separation), \emph{not} full $\mathsf{ZC}$, and Mathias~\cite[pp.~107--234]{Mathias01} proved that $\mathsf{ZC}$ has strictly greater consistency strength than $\mathsf{Mac} \simeq \mathsf{ETCS}$.
  \item \textbf{Disambiguating $\mathsf{ZFC}^-$ versus $\mathsf{ZFC}_{\neg\mathrm{Reg}}$:} Set-theoretic literature overloads $\mathsf{ZFC}^-$ for both Power-Set-free $\mathsf{ZFC}$ and Foundation-free $\mathsf{ZFC}$~\cite{Felgner76}. Throughout this chapter and Chapter~\ref{ch:master-atlas}, we reserve $\mathsf{ZFC}^-$ strictly for \textbf{Power-Set-free} $\mathsf{ZFC}$ (formulated with Collection in place of functional Replacement~\cite{GitmanHamkinsJohnstone16}), and write $\mathsf{ZFC}_{\neg\mathrm{Reg}}$ and $\mathsf{NBG}_{\neg\mathrm{Reg}}$ when \textbf{Foundation (Regularity)} is dropped.
\end{enumerate}
\end{remark}

\begin{theorem}[{External Benchmark: Syntactic $\mathcal{L}_\in$-Non-Conservativity of Global Choice Over $\mathsf{ZC}$; Glazer~2023~\cite[Main Theorem]{Glazer23}}]\label{thm:glazer-nonconservativity}%
\uses{def:set-projection-fibration,rem:zfc-minus-notation}%
Global Choice is \emph{not} syntactically $\mathcal{L}_\in$-conservative over Zermelo set theory with Choice ($\mathsf{ZC}$, with or without Foundation, over $\mathsf{GBC}_0 + \mathsf{GC}$, i.e.\ $\mathsf{GB}_0 + \mathsf{AC} + \mathsf{GC}$): there exists an explicit first-order $\mathcal{L}_\in$-sentence $\theta_{\mathrm{Glz}} := (\mathsf{CH} \wedge \psi_{\mathrm{fnd}})$ consistent with $\mathsf{ZC}$ (realized in a symmetric extension model satisfying the Generalized Continuum Hypothesis $\mathsf{GCH}$) such that $\mathsf{GBC}_0 + \mathsf{GC} \vdash \neg\theta_{\mathrm{Glz}}$. Explicitly, $\psi_{\mathrm{fnd}}$ asserts that parameterless $\mathcal{L}_\in$-definable classes encode:
\begin{enumerate}[label=\textup{(\alph*)},leftmargin=1.8em]
  \item an $\omega$-indexed sequence $\langle X_n \rangle_{n < \omega}$ of nonempty countable sets (admitting no small uniform enumeration sequence in the absence of Replacement);
  \item a class $A$ separated via a binary relation $R$ by subsets of $\bigcup_{n < \omega} X_n$ (so $\{x \in \bigcup_n X_n \mid x \mathrel{R} a_0\} = \{x \in \bigcup_n X_n \mid x \mathrel{R} a_1\} \implies a_0 = a_1$ for $a_0, a_1 \in A$); and
  \item a class surjection $P_2 : A \twoheadrightarrow W_2$ onto the canonical small prewellordering $W_2 \in \mathcal{P}^5(\omega)$ of $\mathcal{P}^2(\omega)$ of second uncountable order type $\omega_2$ (``$\aleph_2$'').
\end{enumerate}
In $\mathsf{GBC}_0 + \mathsf{GC}$, Global Choice simultaneously selects surjections $f_n : \omega \twoheadrightarrow X_n$ for all $n < \omega$, allowing unbounded Subclass Separation (projecting out the unbounded existential quantifier $a \in A$ via \textbf{G4} and intersecting with $\mathcal{P}(\omega \times \omega) \times W_2 \in \mathcal{S}$) to carve out a small-set partial surjection $\mathcal{P}(\omega \times \omega) \rightharpoonup W_2$ sending $\{(m, n) \mid f_n(m) \mathrel{R} a\} \mapsto P_2(a)$ ($a \in A$), refuting $\mathsf{CH}$. (Separately, adjoining unbounded $\mathsf{GB}_0$ Subclass Separation already breaks $\mathcal{L}_\in$-conservativity over $\Delta_0$-bounded $\mathsf{Mac} \simeq \mathsf{ETCS}$~\cite{Mathias01}, whereas Theorem~\ref{thm:glazer-nonconservativity} proves $\mathsf{GC}$ itself breaks $\mathcal{L}_\in$-conservativity over full $\mathsf{ZC}$.)
\end{theorem}

Notice two structural subtleties highlighted by Theorem~\ref{thm:glazer-nonconservativity}. First, syntactic $\mathcal{L}_\in$-non-conservativity is strictly stronger than mere \emph{model non-expandability}: for example, compositional truth $\mathsf{CT}^-[\mathsf{ZFC}]$~\cite{Fujimoto12} is syntactically $\mathcal{L}_\in$-conservative over $\mathsf{ZFC}$ yet fails to expand countable non-recursively-saturated models $M \models \mathsf{ZFC}$ (and likewise $\mathsf{NBG} + \mathsf{ETR}_{\mathrm{Ord}}$ exhibits countable model non-expandability over its own first-order trace $T^*$, while already leaking $\mathrm{Con}(\mathsf{ZFC})$ over $\mathsf{ZFC}$~\cite[Thm.~1.2]{Gitman20}), whereas over $\mathsf{ZC}$, $\mathsf{GBC}_0 + \mathsf{GC}$ refutes the explicit first-order $\mathcal{L}_\in$-sentence $\theta_{\mathrm{Glz}}$ consistent with $\mathsf{ZC}$~\cite{Glazer23}. Second, inside any model $M$ of ordinal height $\mathrm{Ord}^M = \omega + \omega$, uncountable von Neumann initial ordinals do not exist at all (so Proposition~\ref{prop:large-equiv-splitting}(i)'s inclusion $I : \mathbf{Card} \hookrightarrow \mathbf{Set}$ also loses essential surjectivity once Replacement is dropped!), requiring uncountable cardinalities such as $\aleph_n$ to be represented via canonical small prewellorderings inside $\mathcal{P}^{n+3}(\omega)$.

Alongside Glazer's Replacement-free leak, dropping either Foundation or Choice exposes two further horizontal cracks:
\begin{itemize}[leftmargin=1.5em]
  \item \textbf{Dropping Foundation ($\mathsf{ZFC}_{\neg\mathrm{Reg}}$; Felgner~1976~\cite{Felgner76}):} Over Foundation-free set theory $\mathsf{ZFC}_{\neg\mathrm{Reg}} + \mathrm{Collection}$, Ulrich Felgner proved that $\mathsf{NBG}_{\neg\mathrm{Reg}}$ augmented with class Equipollence is non-conservative, forcing a strictly stronger set-level $\mathrm{Ord}$-indexed Dependent Choice schema $\mathsf{DC}_{\mathrm{Ord}}$.
  \item \textbf{Dropping Choice ($\mathsf{GB}_0$ vs.\ Limitation of Size $\mathsf{LoS}$):}\index{Limitation of Size ($\mathsf{LoS}$)} In von Neumann's original 1925 axiomatization of class theory, Replacement and Global Choice were fused into a single mega-axiom, \textbf{Limitation of Size} ($\mathsf{LoS}$): \emph{a class $C \in \mathcal{C}$ fails to belong to the smallness ideal $\mathcal{S}$ if and only if $C$ surjects onto the universe $V$}.
\end{itemize}

\begin{proposition}[Unbundling Limitation of Size Over $\mathsf{GB}_0$]\label{prop:los-smuggles-gc}%
\uses{def:set-projection-fibration}%
Over Replacement-free, Choice-free class theory $\mathsf{GB}_0$, von Neumann's Limitation of Size ($\mathsf{LoS}$) is strictly equivalent to the conjunction of Class Replacement ($\mathrm{Rep}$) and Global Choice ($\mathsf{GC}$):
\[
  \mathsf{GB}_0 \;\vdash\; \mathsf{LoS} \;\Longleftrightarrow\; (\mathrm{Rep} \wedge \mathsf{GC}).
\]
In particular, $\mathsf{GB}_0 + \mathsf{LoS}$ coincides with full $\mathsf{NBG}$ ($\mathrm{Ord} \xrightarrow{\sim} V$), smuggling Global Choice across Choice-free bases and satisfying $\mathsf{GB} + \mathsf{GC} \vdash \mathsf{LoS}$ whenever Replacement holds; conversely, over $\mathsf{GB}_0$ without Replacement, $\mathsf{GB}_0 + \mathsf{GC} \not\vdash \mathsf{LoS}$.
\end{proposition}

\begin{proof}
Work first in $\mathsf{GB}_0 + \mathsf{LoS}$. By predicative class comprehension, the class $\mathrm{Ord}$ of von Neumann ordinals is well-ordered by $\in$, and by the Burali-Forti paradox $\mathrm{Ord} \notin \mathcal{S}$ (otherwise $\mathrm{Ord} \in \mathrm{Ord}$). Applying $\mathsf{LoS}$ yields a class surjection $p : \mathrm{Ord} \twoheadrightarrow V$. Because $\mathrm{Ord}$ is well-ordered, setting $\mu(x) := \min\{\alpha \in \mathrm{Ord} \mid p(\alpha) = x\}$ defines a class injection $\mu : V \hookrightarrow \mathrm{Ord}$, which pulls the canonical well-ordering of $\mathrm{Ord}$ back to a global well-ordering of $V$, establishing $\mathsf{GC}$. To deduce $\mathrm{Rep}$ without invoking Hartogs ordinals, let $A \in \mathcal{S}$ and let $F : A \to V$ be a class function. If $F[A] \notin \mathcal{S}$, then $\mathsf{LoS}$ would furnish a class surjection $q : F[A] \twoheadrightarrow V$ and hence a composite class surjection $h := q \circ F : A \twoheadrightarrow V$. By G1--G8 and Subclass Separation on $A \in \mathcal{S}$, the diagonal subclass $D := \{a \in A \mid a \notin h(a)\} \subseteq A$ is a small set $D \in \mathcal{S} = V$, so $h(a_0) = D$ for some $a_0 \in A$, yielding the Russell--Cantor contradiction $a_0 \in D \iff a_0 \notin D$. Thus $F[A] \in \mathcal{S}$, proving $\mathrm{Rep}$ (whereupon Mostowski collapse of $\mu[V] \subseteq \mathrm{Ord}$ yields $\mathrm{Ord} \xrightarrow{\sim} V$).

Conversely, work in $\mathsf{GB}_0 + \mathrm{Rep} + \mathsf{GC} = \mathsf{NBG}$. If $C \in \mathcal{S}$, the Russell--Cantor diagonal argument just given shows no class surjection $C \twoheadrightarrow V$ can exist. If $C \notin \mathcal{S}$, Foundation and $\mathrm{Rep}$ give $V = \bigcup_{\alpha \in \mathrm{Ord}} V_\alpha$, and restricting $\mathsf{GC}$ to each set-sized power slice $\mathcal{P}(V_{\alpha+1} \setminus V_\alpha) \in \mathcal{S}$ assembles a global bijection $\Phi : V \xrightarrow{\sim} \mathrm{Ord}$. Then $\Phi[C] \subseteq \mathrm{Ord}$ is a proper class, hence unbounded in $\mathrm{Ord}$ (by Subclass Separation on bounded segments) and not order-isomorphic to any $\delta \in \mathrm{Ord}$ (by $\mathrm{Rep}$); sending each $\xi \in \Phi[C]$ to the order type of its set-sized initial segment $\Phi[C] \cap \xi \in \mathcal{S}$ (requiring no $\mathsf{ETR}_{\mathrm{Ord}}$) maps $C \xrightarrow{\sim} \mathrm{Ord} \xrightarrow{\Phi^{-1}} V$ bijectively onto $V$, verifying $\mathsf{LoS}$.

Finally, for $\mathsf{GB}_0 + \mathsf{GC} \not\vdash \mathsf{LoS}$, consider G\"odel's constructible rank $L_{\omega_\omega^L}$ with its definable classes $\mathrm{Def}(L_{\omega_\omega^L})$ (where $\omega_n^L$ denotes the $n$-th uncountable cardinal inside $L$). By G\"odel's condensation lemma applied inside $L$, every constructible subset of $x \in L_{\omega_n^L}$ lies in $L_{\omega_{n+1}^L} \subset L_{\omega_\omega^L}$, verifying Power Set alongside $\mathsf{Z}$, while the global $<_L$-order gives $\mathsf{GC}$ on a model of constructible ordinal height $\omega_\omega^L$. Yet the cofinal sequence $\{\omega_n^L : n \in \omega\}$ is a countable definable proper class admitting no surjection onto $L_{\omega_\omega^L}$, refuting both $\mathrm{Rep}$ and $\mathsf{LoS}$.
\end{proof}

\begin{table}[htbp]
\centering
\caption{Base-Fiber Map of Set-Level vs.\ Class-Level Spectra Along $\pi_{\mathrm{Set}}$}\label{tab:fibration-map}
\vspace{0.3em}
{\scriptsize\setlength{\tabcolsep}{2.5pt}\renewcommand{\arraystretch}{1.12}%
\begin{tabularx}{\textwidth}{@{}P{0.18\textwidth}P{0.12\textwidth}P{0.17\textwidth}YY@{}}
\toprule
\textbf{Class Theory} & \textbf{Set Base $\pi_{\mathrm{Set}}$} & \textbf{$\mathcal{L}_\in$-Conservative?} & \textbf{Separating Model} & \textbf{Categorical Residuum} \\
\midrule
$\mathsf{GBC}$ & $\mathsf{ZFC}$ & Yes ($\mathrm{Def}(M)$) & $M \models \forall p\,(V \neq \mathrm{HOD}(\{p\}))$ & Baseline left adjoint $\mathrm{Def}$ \\
$\mathsf{NBG}$ ($\mathsf{GBC}+\mathsf{GC}$) & $\mathsf{ZFC}$ & Yes (Felgner~\cite{Felgner71}) & $(M,\mathrm{Def}(M))$ vs.\ $M[G]$ & Unsplit $\mathbf{Card} \hookrightarrow \mathbf{Set}$ \\
$\mathsf{CT}^-[\mathsf{ZFC}]$ & $\mathsf{ZFC}$ & Yes (Fujimoto~\cite{Fujimoto12}) & Non-rec.-sat.\ $M \models \mathsf{ZFC}$ & Countable model non-expandability \\
$\mathsf{NBG}+\mathsf{ETR}_{\mathrm{Ord}}$ ($\subsetneq \mathsf{ETR}$) & $\mathsf{ZFC}$ / $T^*$ & \textbf{No} ($\mathrm{Con}(\mathsf{ZFC})$; yes/$T^*$) & Minimal $(M,\mathrm{Def}(M))$ & Class forcing $\Vdash_{\mathbb{P}}$~\cite[Thm.~1.2]{Gitman20} ($\mathrm{Ord}$-rec.; full $\mathsf{ETR} \vdash \mathrm{Con}$) \\
Kelley--Morse ($\mathsf{KM}$) & $\mathsf{ZFC}$ / $\mathrm{Th}_{\mathcal{L}_\in}(\mathsf{KM})$ & \textbf{No} ($\mathrm{Con}(T^*)$; yes/$\mathrm{Th}_{\mathcal{L}_\in}(\mathsf{KM})$) & $(V_\kappa, V_{\kappa+1})$ ($\kappa$ inacc.) & Impredicative $\Pi^1_\infty$ classes (Groth.\ univ.\ for Kan ext., Freyd) \\
$\mathsf{GBC}_0+\mathsf{GC}$ / $\mathsf{NBG}_{\neg\mathrm{Reg}}$ & $\mathsf{ZC}$ / $\mathsf{ZFC}_{\neg\mathrm{Reg}}$ & \textbf{No} (Glazer~\cite{Glazer23}; Felgner~\cite{Felgner76}) & Symm.\ $V_{\omega+\omega}$ ($\mathsf{ZC}$); FM / Rieger--Bernays ($\mathsf{ZFC}_{\neg\mathrm{Reg}}$) & Leaks $\mathcal{P}(\omega)\twoheadrightarrow W_2$ / $\mathsf{DC}_{\mathrm{Ord}}$ \\
\bottomrule
\end{tabularx}\par}
\end{table}

Table~\ref{tab:fibration-map} triggers an immediate plot twist for Volume~I itself.

\begin{proposition}[Countable Colimit Escape in $V_{\omega+\omega}$ Over $\mathsf{ZC}$ and $\mathsf{ETCS}$]\label{prop:etcs-colimit-escape}%
\uses{rem:zfc-minus-notation}%
Zermelo's cumulative rank universe $V_{\omega+\omega}$ models full Zermelo set theory with Choice $\mathsf{ZC}$ (including unbounded $\mathcal{L}_\in$-Separation, and \emph{a fortiori} its category of sets and functions models $\mathsf{Mac} \simeq \mathsf{ETCS}$), in which the countable tower $n \mapsto \mathcal{P}^n(\mathbb{N})$ is first-order definable over $\mathbb{N}$ yet its colimit $\bigcup_{n \in \mathbb{N}} \mathcal{P}^n(\mathbb{N}) \subsetneq V_{\omega+\omega} = \bigcup_{n \in \mathbb{N}} \mathcal{P}^n(V_\omega)$ fails to exist as a small object.
\end{proposition}

\begin{proof}
Because $\omega + \omega$ is a limit ordinal strictly above $\omega$, $V_{\omega+\omega}$ contains $\omega = \mathbb{N}$ and is closed under pairing, union, and power sets ($A \in V_{\omega+k} \implies \mathcal{P}(A) \in V_{\omega+k+1} \subset V_{\omega+\omega}$), inheriting $\mathrm{AC}$ from the metatheory. Moreover, for any $A \in V_{\omega+k}$ and \emph{any} subset $B \subseteq A$ carved out by an arbitrary unbounded $\mathcal{L}_\in$-formula over $(V_{\omega+\omega}, \in)$, we already have $B \in \mathcal{P}(A) \in V_{\omega+k+1} \subset V_{\omega+\omega}$, so $V_{\omega+\omega} \models \mathsf{ZC}$ (and hence models $\mathsf{Mac} \simeq \mathsf{ETCS}$). Inside $V_{\omega+\omega}$, each finite iterate $\mathcal{P}^n(\mathbb{N})$ has rank $\omega + n < \omega + \omega$ and $y = \mathcal{P}^n(\mathbb{N})$ is $\mathcal{L}_\in$-definable. However, every element $x \in V_{\omega+\omega}$ lies in some $V_{\omega+k}$ ($k < \omega$) of rank strictly below $\omega + \omega$, whereas $\bigcup_{n \in \mathbb{N}} \mathcal{P}^n(\mathbb{N})$ contains sets of rank $\omega + n$ for all $n \in \mathbb{N}$ and therefore escapes through the roof of the topos (strictly inside $V_{\omega+\omega} = \bigcup_{n \in \mathbb{N}} \mathcal{P}^n(V_\omega)$, since e.g.\ $V_\omega \in V_{\omega+1} \setminus \bigcup_{n \in \mathbb{N}} \mathcal{P}^n(\mathbb{N})$, while both $\{\mathcal{P}^n(\mathbb{N})\}_{n \in \mathbb{N}}$ and $\{V_{\omega+n}\}_{n \in \mathbb{N}}$ are rank-cofinal countable families escaping $\mathcal{S}$).
\end{proof}

One might suspect that losing rank $\omega + \omega$ and $\aleph_\omega$ affects only esoteric transfinite cardinal arithmetic. Harvey Friedman demolished that comfort in 1971~\cite{Friedman71} by proving that \textbf{Borel Determinacy}\index{Borel Determinacy}---whether every infinite two-player game on a Borel subset of Cantor space $2^{\mathbb{N}}$, an arena whose moves, strategies, and Borel payoff sets reside down in $V_{\omega+1}$ and $V_{\omega+2}$ inside $V_{\omega+\omega}$, admits a winning strategy---\textbf{cannot be proved in Zermelo set theory $\mathsf{ZC}$ (and hence cannot be proved in $\mathsf{ETCS}$)}. Synthesizing winning strategies for general Borel games ($\boldsymbol{\Sigma}^0_\alpha$ across $\alpha < \omega_1$) requires up to $\omega_1$ iterated power sets supplied by Replacement (already exceeding the bounded finite iterates $\{\mathcal{P}^n(\mathbb{N})\}_{n < \omega}$ available in $\mathsf{ZC}$).

Look at the resulting predicament on our supposedly settled ground floor: adjoining unbounded $\mathsf{GB}_0$ Subclass Separation already breaks $\mathcal{L}_\in$-conservativity over $\Delta_0$-bounded $\mathsf{ETCS} \simeq \mathsf{Mac}$~\cite{Mathias01}, and working over full $\mathsf{ZC}$ while casually invoking Global Choice overhead (``because $\mathsf{NBG}$ is conservative over $\mathsf{ZFC}$'') mutates small-set cardinal arithmetic underneath via Theorem~\ref{thm:glazer-nonconservativity}; yet remaining strictly inside $\mathsf{ETCS}$ or $\mathsf{ZC}$ leaves countable definable colimits escaping the universe and concrete Borel games in $V_{\omega+2}$ undecided. ``$\mathbf{Sets}$'' was never an indivisible atom---compelling us to place ``$\mathbf{Sets}$'' directly on the $(\mathcal{C}, \mathcal{S})$ dissecting table.

\section{\texorpdfstring{Dissecting ``$\mathbf{Sets}$'' in $(\mathcal{C}, \mathcal{S})$: Plumbing vs.\ Over-Compressed Existence}{Dissecting "Sets" in (C, S): Plumbing vs. Over-Compressed Existence}}\label{sec:ast-plumbing-compression}
\sectionepigraph{``The formal system may need to be changed. It would be one-sided to only `work in' a system, accepting it as a divine condemnation that we have to endure.''}{F.~William Lawvere, Categories Mailing List (October 20, 1998)}

Taking Lawvere's injunction seriously means asking why single-sorted $\mathsf{ZF}$ requires infinite first-order schemata for Separation and Replacement---because they tangle two distinct jobs together: \textbf{logical carving} (intersections, complements, pullbacks, and relational images) and \textbf{size compression} (certifying that a carved-out collection is small enough to belong to further collections). Passing to Joyal and Moerdijk's \textbf{Algebraic Set Theory (AST)}~\cite{JoyalMoerdijk95} separates those two jobs cleanly into an ambient class category $\mathcal{C}$ and a small-map ideal $\mathcal{S}$---provided we do not bake Replacement into $\mathcal{S}$ prematurely.

\begin{definition}[Predicative Class-Set Ambient $(\mathcal{C}, \mathcal{S})$ and Contravariant Powerclasses]\label{def:small-map-ambient}%
\uses{def:comparison-triple}%
A baseline \emph{predicative class-set ambient} $(\mathcal{C}, \mathcal{S})$ consists of:
\begin{enumerate}[label=\textup{(\roman*)},leftmargin=1.8em]
  \item an ambient Boolean regular (or, constructively, Heyting regular) category $\mathcal{C}$ of \emph{classes}, finitely generated in the two-sorted set/class case by G\"odel's eight predicative relational combinators \textbf{G1--G8}; and
  \item a stable subcategory $\mathcal{S} \subseteq \mathcal{C}$ of \emph{small maps} containing isomorphisms and terminal projections $0 \to 1$, $1 + 1 \to 1$, closed under composition, pullbacks along arbitrary morphisms of $\mathcal{C}$ (inducing Subclass Separation when $\mathrm{Mono} \subseteq \mathcal{S}$, i.e.\ subobject smallness $U \rightarrowtail B \in \mathcal{S}$ holds slice-internally across \emph{every} slice $\mathcal{C}/B$, not merely over global $1_{\mathcal{C}}$~\cite[\S I.1]{JoyalMoerdijk95}), diagonals $\Delta_C : C \rightarrowtail C \times C$ for small-element equality, and finite disjoint sums---\emph{without} assuming quotient or regular-image descent.
\end{enumerate}
An object $X \in \mathcal{C}$ is \emph{small} (a \emph{set}, written $X \in \mathcal{S}$) if $X \to 1$ belongs to $\mathcal{S}$. Under Joyal--Moerdijk's \textbf{Powerclass Axiom (P1)}~\cite[\S I.1]{JoyalMoerdijk95}---asserting that every class $C \in \mathcal{C}$ admits a powerclass object $\mathcal{P}_{\mathcal{S}}(C) \in \mathcal{C}$ with a small membership relation $\in_C \rightarrowtail C \times \mathcal{P}_{\mathcal{S}}(C)$ ($\pi_2 \in \mathcal{S}$) classifying $\mathcal{S}$-small subrelations $R \rightarrowtail D \times C$ ($\pi_1 \in \mathcal{S}$)---$\mathcal{P}_{\mathcal{S}} : \mathcal{C}^{\mathrm{op}} \to \mathcal{C}$, $f \mapsto f^*$, acts contravariantly and $\Omega := \mathcal{P}_{\mathcal{S}}(1)$ is the $\mathcal{S}$-small subobject classifier (coinciding with the full subobject classifier of $\mathcal{C}$ when all monomorphisms belong to $\mathcal{S}$).
\end{definition}

Because G\"odel's operations G1--G8 replace infinite comprehension schemata by finite relational algebra, $\mathcal{C}$ carries zero hidden ``overweight existence'' commitments. Running the existence axioms of set theory through $(\mathcal{C}, \mathcal{S})$ splits them into two sharply contrasting strata.

\subsection{\texorpdfstring{Stratum~I: Honest Lightweight Plumbing (\emph{Empty Set, Pairing, Finite Union})}{Stratum I: Honest Lightweight Plumbing (Empty Set, Pairing, Finite Union)}}\label{subsec:ast-stratum-i}

Categorically, Empty Set, Pairing, and Finite Union assert that $\mathcal{C}$ and $\mathcal{S}$ possess a strict initial object $0$, binary coproduct $1 + 1$, pushouts, and disjoint finite sums---the basic grammar of a \textbf{lextensive regular category}~\cite[Ch.~I]{JoyalMoerdijk95} (or, adjoining a Natural Numbers Object $\mathrm{NNO}$, Maietti and Joyal's \textbf{Arithmetic Universe}~\cite{Maietti05}, calibrated along Arabic Axis~1, ambient universe deformation, in Chapter~\ref{ch:base-fiber}, \S\ref{sec:axis1-deformations}, and along Roman Axes~II--III in Chapter~\ref{ch:panorama}, \S\ref{sec:predicativity-wtypes}). Weakening Stratum~I precipitates no cohomological or topological residua; it merely dismantles the ambient doctrine's ability to state comparison cells $\kappa : \Phi \Rightarrow \Psi$. These axioms pull their exact syntactic weight.

\subsection{\texorpdfstring{Stratum~II: Overweight Size-Compression Valves (\emph{Power Set, Replacement, Unbounded Separation})}{Stratum II: Overweight Size-Compression Valves (Power Set, Replacement, Unbounded Separation)}}\label{subsec:ast-stratum-ii}

Inside $(\mathcal{C}, \mathcal{S})$ under Powerclass Axiom (P1), every class $C \in \mathcal{C}$ \emph{already} possesses a powerclass $\mathcal{P}_{\mathcal{S}}(C) \in \mathcal{C}$, and every class morphism $f : I \to C$ \emph{already} possesses a regular image $\mathrm{im}(f) \rightarrowtail C$ in $\mathcal{C}$. The heavyweight existence axioms of $\mathsf{ZF}$ act strictly as \textbf{size-compression valves} forcing these class-sized constructions back down into $\mathcal{S}$.

\begin{proposition}[Pointillist Continuum in $H(\aleph_1)$ and Constructive $\Omega^X / 2^X$ Split]\label{prop:pointillist-continuum}%
\uses{def:small-map-ambient,rem:zfc-minus-notation}%
Consider the Power Set compression principle $X \in \mathcal{S} \implies \mathcal{P}_{\mathcal{S}}(X) \in \mathcal{S}$ inside $(\mathcal{C}, \mathcal{S})$:
\begin{enumerate}[label=\textup{(\roman*)},leftmargin=1.8em]
  \item Let $H(\aleph_1)$ be the universe of hereditarily countable sets equipped with its definable classes. Then $H(\aleph_1) \models \mathsf{ZFC}^-$~\cite{GitmanHamkinsJohnstone16} (satisfying Infinity, Collection/Replacement, and Choice while refuting Power Set): every individual real number $r \subseteq \mathbb{N}$ lies in $\mathcal{S}$, whereas Cantor space $2^{\mathbb{N}} = \mathcal{P}_{\mathcal{S}}(\mathbb{N})$ and the continuum $\mathbb{R}$ remain proper classes ($2^{\mathbb{N}} \notin \mathcal{S}$).
  \item Over constructive Aczel--Myhill set theory $\mathsf{CZF}$~\cite{Aczel78} (dropping Excluded Middle $\mathrm{LEM}$ along \textbf{Axis~II}, Constructive Logic), \emph{Exponentiation} ($2^X \in \mathcal{S}$, classifying decidable subobjects $\mathrm{Sub}_{\mathrm{dec}}(X)$) is strictly weaker than \emph{Power Set} ($\Omega^X = \mathcal{P}_{\mathcal{S}}(X) \in \mathcal{S}$, classifying arbitrary subobjects $\mathrm{Sub}(X)$), precipitating a nontrivial subobject-classifier size residuum $\Omega^X / 2^X$.
\end{enumerate}
\end{proposition}

\begin{proof}
For~(i), $\mathbb{N} = \omega \in H(\aleph_1)$, and for every $r \subseteq \mathbb{N}$ its transitive closure $r \cup \omega$ is countable, so $r \in H(\aleph_1)$. Likewise, any countable family of hereditarily countable sets has countable union and transitive closure, verifying Collection/Replacement in $H(\aleph_1)$. By Cantor's diagonal theorem, however, $\mathcal{P}_{\mathcal{S}}(\mathbb{N})$ is uncountable and cannot lie in $H(\aleph_1)$, remaining an honest proper class in $\mathcal{C} \setminus \mathcal{S}$.

For~(ii), inside a constructive Heyting $(\mathcal{C}, \mathcal{S})$ without $\mathrm{LEM}$, $2 = 1 + 1 \in \mathcal{S}$ embeds properly into $\Omega = \mathcal{P}_{\mathcal{S}}(1)$ along Axis~II ($\kappa : 1 + 1 \hookrightarrow \Omega$). Subset Collection in $\mathsf{CZF}$ implies Exponentiation ($Y^X \in \mathcal{S}$ for all $X, Y \in \mathcal{S}$, so $2^X \in \mathcal{S}$) while remaining $\Pi^0_2$-arithmetically conservative (equiconsistent, with identical proof-theoretic strength) over Kripke--Platek set theory with Infinity ($\mathsf{KP}\omega$) via Aczel's Martin-L\"of type-universe interpretation~\cite{Aczel78}. Forcing $\Omega \in \mathcal{S}$ ($\Omega^X \in \mathcal{S}$) impredicatively compresses all unbounded subobject predicates into $\mathcal{S}$, jumping far past $\mathsf{KP}\omega$ and full second-order arithmetic $\mathsf{Z}_2$.
\end{proof}

Next, Joyal and Moerdijk~\cite[Ch.~I, \S\S 1--3]{JoyalMoerdijk95} showed that \textbf{Replacement} corresponds precisely to closing $\mathcal{S}$ under quotients and regular images---the exact operation needed to make $\mathcal{P}_{\mathcal{S}}$ covariantly functorial!

\begin{definition}[Axis~VI Doctrinal Comparison Triple: Quotient-Image Descent and Cofinality]\label{def:axis-vi-replacement}%
\uses{def:comparison-triple,def:small-map-ambient}%
Over a baseline predicative class-set ambient $(\mathcal{C}, \mathcal{S})$ (Definition~\ref{def:small-map-ambient}), where an $I$-indexed family of small objects over $I \in \mathcal{S}$ is represented internally as a small map $p : E \to I$ in $\mathcal{S}$, \textbf{Axis~VI (Small-Indexed Cofinality / Replacement)}\index{Axis VI (Quotient-Image Descent / Replacement)} compares the class regular image $\Phi_{\mathrm{Im}}(I, f) := \mathrm{im}(f) \in \mathcal{C}$ of any morphism $f : I \to C$ with small domain $I \in \mathcal{S}$ against the smallness reflector $\mathcal{S} \hookrightarrow \mathcal{C}$, via the descent condition $\kappa_{\mathrm{Rep}} : \mathrm{im}(f) \in \mathcal{S}$ (inducing regular-quotient / coequalizer descent along regular epimorphisms $q : I \twoheadrightarrow Q$ in $\mathcal{C}$). Equivalently, Axis~VI asserts that covariant direct-image pushforward $\exists_g : \mathcal{P}_{\mathcal{S}}(C) \to \mathcal{P}_{\mathcal{S}}(D)$ lands in small subobjects for every $g : C \to D$ in $\mathcal{C}$ (promoting $\mathcal{P}_{\mathcal{S}} : \mathcal{C} \to \mathcal{C}$ to a covariant endofunctor) and, combined with Small Union ($I \in \mathcal{S} \wedge (p : E \to I) \in \mathcal{S} \implies E \in \mathcal{S}$) and coequalizer descent, compresses every small-diagram colimit $\varinjlim_{i \in I} F_i$ into $\mathcal{S}$.
\end{definition}

\begin{remark}[The Borel Determinacy Pincer and Martin's Proper-Class Unraveling Coalgebra]\label{rem:martin-borel-coalgebra}%
\uses{prop:pointillist-continuum,prop:etcs-colimit-escape,def:axis-vi-replacement}%
Comparing Proposition~\ref{prop:pointillist-continuum}(i) against Proposition~\ref{prop:etcs-colimit-escape} exposes a sharp reverse-mathematical pincer across Stratum~II. Inside an ambient universe $V \models \mathsf{ZFC}$, winning strategies on $\omega^\omega$ are hereditarily countable, so the true hereditarily countable universe $H(\aleph_1) \models \mathsf{ZFC}^-$ inherits full Borel Determinacy from $V$; what breaks already at finite Borel rank $\boldsymbol{\Sigma}^0_4$ is \emph{provability inside Power-Set-free $\mathsf{ZFC}^-$} itself, refuted in minimal or non-standard $\beta$-models of $\mathsf{ZFC}^-$ lacking $\mathcal{P}_{\mathcal{S}}(\mathbb{N})$~\cite{Friedman71,Martin75,MontalbanShore12,GitmanHamkinsJohnstone16}. Conversely, in Replacement-free $\mathsf{ZC}$ (and $\mathsf{ETCS} \simeq \mathsf{Mac}$), all reals, trees, and strategies $\sigma : \omega^{<\omega} \to \omega$ ($\in V_{\omega+1}$), as well as Borel subsets of Cantor space ($\in V_{\omega+2}$), already reside inside the true cumulative rank $V_{\omega+\omega} \subset V \models \mathsf{ZFC}$, which therefore inherits full Borel Determinacy from $V$. What breaks sharply at $\boldsymbol{\Sigma}^0_\omega$ is \emph{provability inside $\mathsf{ZC}$} (and $\mathsf{ETCS} \simeq \mathsf{Mac}$), refuted in non-standard or minimal models of $\mathsf{ZC}$~\cite{Friedman71}. Inside $(V_{\omega+\omega}, \mathrm{Def}(V_{\omega+\omega}))$, each finite iterate $\mathcal{P}^k(\mathbb{N})$ lies in $\mathcal{S}$---allowing $\mathsf{ZC}$ to prove $\boldsymbol{\Sigma}^0_k$-determinacy for every fixed meta-integer $k < \omega$~\cite{Martin75}---whereas the countable family $\{\mathcal{P}^n(\mathbb{N})\}_{n \in \mathbb{N}}$ escapes $\mathcal{S}$ and prevents Martin's inductive unraveling construction from internalizing inside $\mathcal{S}$. Inside $(\mathcal{C}, \mathcal{S})$, Tony Martin's inductive clopen unraveling of a $\boldsymbol{\Sigma}^0_\omega$ game still assembles into a well-defined \emph{proper-class covering tree in $\mathcal{C} \setminus \mathcal{S}$}, framing the class-to-set strategy-descent obstruction $\mathrm{Res}(\kappa_{\mathrm{Det}})$ as a flagship Volume~II problem. Meanwhile, \textbf{Separation / Reflection} stratifies into four distinct rungs: $\Delta_0$-pullbacks along $\top : 1 \to \Omega$ in $\mathsf{ETCS} \simeq \mathsf{Mac}$, full unbounded $\mathcal{L}_\in$-Separation in $\mathsf{ZC}$ (strictly stronger than $\mathsf{Mac}$ by Mathias~\cite{Mathias01}, yet satisfied in $V_{\omega+\omega}$ without Replacement), the full L\'evy--Montague Reflection Schema ($\Sigma_n$-Reflection to cumulative rank sets $V_\alpha$ for every meta-integer $n < \omega$, strictly beyond any bounded $\Sigma_k$-Reflection fragment) upgrading $\mathsf{ZC}$ to $\mathsf{ZFC}$, and impredicative class comprehension jumping from conservative $\mathsf{NBG}$ to $\mathsf{KM}$ ($\mathrm{Con}(\mathsf{NBG}+\mathsf{ETR}_{\mathrm{Ord}})$).
\end{remark}

\section{\texorpdfstring{Over-Quotienting Structure: The Anti-Existence Residua}{Over-Quotienting Structure: The Anti-Existence Residua}}\label{sec:anti-existence-residua}

Where Stratum~II manufactures residua by over-compressing proper-class constructions downward into $\mathcal{S}$, \textbf{Stratum~III} inverts the polarity completely. \textbf{Foundation (Regularity)} and \textbf{Extensionality} assert the existence of no sets whatsoever; they are \emph{structural prohibitions} that quotient away circular $\in$-graphs and internal symmetries. Relaxing either clamp uncovers universe-scale instances of the exact fixed-point ($\mu F \Rightarrow \nu F$) and truncation operators studied in Chapters~\ref{ch:delta-rings}, \ref{ch:higher-crm}, and~\ref{ch:realizability-sites}.

\begin{theorem}[Axis~VII Comparison: Foundation vs.\ Anti-Foundation as $\mathbf{Coalg}_{\mathrm{wf}}(\mathcal{P}_{\mathcal{S}}) \hookrightarrow \mathbf{Coalg}(\mathcal{P}_{\mathcal{S}})$]\label{thm:axis-vii-foundation}%
\lean{HomotopyCRM.SetsResiduum.pset_equiv_iff_bisimilar, HomotopyCRM.SetsResiduum.quine_atom_outside_wf}%
\leanwitness%
\uses{def:small-map-ambient,def:axis-vi-replacement}%
Let $(\mathcal{C}, \mathcal{S})$ be a Barr-exact class-set ambient with $\mathbb{N} \in \mathcal{S}$, Small Union, and Axis~VI image descent (Definition~\ref{def:axis-vi-replacement}), so that $\mathcal{P}_{\mathcal{S}} : \mathcal{C} \to \mathcal{C}$ acts covariantly by direct image $\exists_g$, and assume $\mathcal{P}_{\mathcal{S}}$ admits both an initial algebra $(\mu \mathcal{P}_{\mathcal{S}}, \mathrm{sup})$ and a terminal coalgebra $(\nu \mathcal{P}_{\mathcal{S}}, \mathrm{des})$. Let $\mathbf{Coalg}_{\mathrm{wf}}(\mathcal{P}_{\mathcal{S}}) \hookrightarrow \mathbf{Coalg}(\mathcal{P}_{\mathcal{S}})$ denote the full subdoctrine of well-founded $\mathcal{P}_{\mathcal{S}}$-coalgebras $(X, \gamma : X \to \mathcal{P}_{\mathcal{S}}(X))$ (satisfying internal $\in_\gamma$-induction). Then:
\begin{enumerate}[label=\textup{(\roman*)},leftmargin=1.8em]
  \item \textbf{Well-Founded Coalgebra Terminality (Osius~\cite{Osius74}):} The cumulative well-founded universe $V_{\mathsf{ZF}}$ satisfying Foundation\index{Axis VII (Foundation vs.\ Anti-Foundation)} is canonically isomorphic to the initial algebra $\mu \mathcal{P}_{\mathcal{S}}$, which via $\mathrm{sup}^{-1} : \mu \mathcal{P}_{\mathcal{S}} \xrightarrow{\sim} \mathcal{P}_{\mathcal{S}}(\mu \mathcal{P}_{\mathcal{S}})$ forms the terminal object $\nu^{\mathrm{wf}} \mathcal{P}_{\mathcal{S}} \simeq \mu \mathcal{P}_{\mathcal{S}}$ of $\mathbf{Coalg}_{\mathrm{wf}}(\mathcal{P}_{\mathcal{S}})$~\cite[Chs.~II--III]{JoyalMoerdijk95}.
  \item \textbf{Unrestricted Coalgebra Terminality (Aczel~\cite[Chs.~1--2, 6--7]{Aczel88}):} Peter Aczel's Anti-Foundation universe $V_{\mathsf{AFA}}$\index{Anti-Foundation Axiom ($\mathsf{AFA}$)}---in which every accessible pointed $\mathcal{S}$-small-branching graph $(G, \to, r)$ admits a unique decoration $d : G \to V$ satisfying $d(x) = \{d(y) \mid x \to y\}$---is canonically isomorphic to the terminal object $\nu \mathcal{P}_{\mathcal{S}}$ of $\mathbf{Coalg}(\mathcal{P}_{\mathcal{S}})$.
  \item \textbf{Strict Comparison Monomorphism and Factorization Criterion:} Passing along the doctrinal inclusion $\mathbf{Coalg}_{\mathrm{wf}}(\mathcal{P}_{\mathcal{S}}) \hookrightarrow \mathbf{Coalg}(\mathcal{P}_{\mathcal{S}})$ induces a canonical comparison monomorphism $\kappa_{\mathrm{Reg}} : \nu^{\mathrm{wf}} \mathcal{P}_{\mathcal{S}} \simeq \mu \mathcal{P}_{\mathcal{S}} \hookrightarrow \nu \mathcal{P}_{\mathcal{S}}$. Whenever $0 \not\cong 1$, $\kappa_{\mathrm{Reg}}$ is a \emph{strict} monomorphism whose non-well-founded residuum $\mathrm{Res}(\kappa_{\mathrm{Reg}}) = \nu \mathcal{P}_{\mathcal{S}} \setminus \mu \mathcal{P}_{\mathcal{S}}$ is anchored by the unique Quine atom\index{Quine atom ($q = \{q\}$)} $q = \{q\}$, and a coalgebra $(X, \gamma) \in \mathbf{Coalg}(\mathcal{P}_{\mathcal{S}})$ lies in $\mathbf{Coalg}_{\mathrm{wf}}(\mathcal{P}_{\mathcal{S}})$ (satisfies Foundation) if and only if its terminal decoration $d_\gamma : X \to \nu \mathcal{P}_{\mathcal{S}}$ factors through $\kappa_{\mathrm{Reg}} : \mu \mathcal{P}_{\mathcal{S}} \hookrightarrow \nu \mathcal{P}_{\mathcal{S}}$.
\end{enumerate}
\end{theorem}

\begin{proof}
By Lambek's Lemma (formalized as \texttt{lambek\_isIso} in Chapter~\ref{ch:delta-rings}, \S\ref{sec:delta-rings}), $\mathrm{sup} : \mathcal{P}_{\mathcal{S}}(\mu \mathcal{P}_{\mathcal{S}}) \xrightarrow{\sim} \mu \mathcal{P}_{\mathcal{S}}$ and $\mathrm{des} : \nu \mathcal{P}_{\mathcal{S}} \xrightarrow{\sim} \mathcal{P}_{\mathcal{S}}(\nu \mathcal{P}_{\mathcal{S}})$ are isomorphisms, modeling extensional membership via $x \in y \iff x \in \mathrm{sup}^{-1}(y)$ (respectively $x \in \mathrm{des}(y)$). Initiality of $\mu \mathcal{P}_{\mathcal{S}}$ yields $\in$-induction on $V_{\mathsf{ZF}} \cong \mu \mathcal{P}_{\mathcal{S}}$, and Osius's recursive coalgebra theorem~\cite{Osius74}~\cite[Chs.~II--III]{JoyalMoerdijk95} equips every $(X, \gamma) \in \mathbf{Coalg}_{\mathrm{wf}}(\mathcal{P}_{\mathcal{S}})$ with a unique coalgebra homomorphism $m_\gamma : (X, \gamma) \to (\mu \mathcal{P}_{\mathcal{S}}, \mathrm{sup}^{-1})$, proving $\nu^{\mathrm{wf}} \mathcal{P}_{\mathcal{S}} \simeq \mu \mathcal{P}_{\mathcal{S}}$ and~(i). Moreover, initiality of $(\mu \mathcal{P}_{\mathcal{S}}, \mathrm{sup})$ (equivalently, terminality of $(\nu \mathcal{P}_{\mathcal{S}}, \mathrm{des})$) induces the unique algebra-coalgebra homomorphism $\kappa_{\mathrm{Reg}} : \mu \mathcal{P}_{\mathcal{S}} \to \nu \mathcal{P}_{\mathcal{S}}$ with $\mathrm{des} \circ \kappa_{\mathrm{Reg}} = \exists_{\kappa_{\mathrm{Reg}}} \circ \mathrm{sup}^{-1}$. By $\in$-induction on $x \in \mu \mathcal{P}_{\mathcal{S}}$, if $\kappa_{\mathrm{Reg}}(u) = \kappa_{\mathrm{Reg}}(z) \implies u = z$ holds for all $u \in \mathrm{sup}^{-1}(x)$ and all $z \in \mu \mathcal{P}_{\mathcal{S}}$, then $\kappa_{\mathrm{Reg}}(x) = \kappa_{\mathrm{Reg}}(y)$ gives $\exists_{\kappa_{\mathrm{Reg}}}(\mathrm{sup}^{-1}(x)) = \exists_{\kappa_{\mathrm{Reg}}}(\mathrm{sup}^{-1}(y))$ and hence $\mathrm{sup}^{-1}(x) = \mathrm{sup}^{-1}(y)$ and $x = y$; thus $\kappa_{\mathrm{Reg}}$ is monic.

For~(ii) and~(iii), any accessible pointed small-branching graph $(G, \to, r)$ defines $(G, \gamma) \in \mathbf{Coalg}(\mathcal{P}_{\mathcal{S}})$ via $\gamma(x) := \{y \in G \mid x \to y\}$. Here $\mathbb{N} \in \mathcal{S}$, Small Union, and Axis~VI Replacement compress every reachable transitive subgraph $G_x = \bigcup_{n \in \mathbb{N}} \gamma^n(\{x\})$ into $\mathcal{S}$, while Barr-exactness of $\mathcal{C}$ (Chapter~\ref{ch:panorama}, Figure~\ref{fig:homological-lattice}) directly quotients the class of small APGs by maximum bisimulation to realize $\nu \mathcal{P}_{\mathcal{S}}$ in Foundation-free ambients without circularly invoking rank-based Scott's trick. Terminality of $(\nu \mathcal{P}_{\mathcal{S}}, \mathrm{des})$ yields a unique homomorphism $d_\gamma : (G, \gamma) \to (\nu \mathcal{P}_{\mathcal{S}}, \mathrm{des})$, i.e.\ $d_\gamma(x) = \{d_\gamma(y) \mid x \to y\}$ (Aczel's $\mathsf{AFA}$~\cite[Chs.~1--2]{Aczel88}). If $(X, \gamma) \in \mathbf{Coalg}_{\mathrm{wf}}(\mathcal{P}_{\mathcal{S}})$, uniqueness forces $d_\gamma = \kappa_{\mathrm{Reg}} \circ m_\gamma$. Conversely, if $d_\gamma = \kappa_{\mathrm{Reg}} \circ m$, monicity of $\kappa_{\mathrm{Reg}}$ makes $\exists_{\kappa_{\mathrm{Reg}}}$ monic, so $\exists_{\kappa_{\mathrm{Reg}}} \circ \exists_m \circ \gamma = \mathrm{des} \circ d_\gamma = \exists_{\kappa_{\mathrm{Reg}}} \circ \mathrm{sup}^{-1} \circ m$ implies $\mathrm{sup}^{-1} \circ m = \exists_m \circ \gamma$; for any $\in_\gamma$-inductive subobject $S \rightarrowtail X$, the Heyting/Boolean universal image $\forall_m(S) \rightarrowtail \mu \mathcal{P}_{\mathcal{S}}$ (Definition~\ref{def:small-map-ambient}(i)) is $\in$-inductive on $\mu \mathcal{P}_{\mathcal{S}}$, giving $\forall_m(S) = \mu \mathcal{P}_{\mathcal{S}}$ and $S = X$. Finally, whenever $0 \not\cong 1$, decorating $(1, \top : 1 \to \mathcal{P}_{\mathcal{S}}(1))$ yields $q = d_\top(\bullet) \in \nu \mathcal{P}_{\mathcal{S}}$ with $q = \{q\}$, refuting factorization through $\kappa_{\mathrm{Reg}}$ and witnessing $\nu \mathcal{P}_{\mathcal{S}} \setminus \mu \mathcal{P}_{\mathcal{S}} \neq \emptyset$.
\end{proof}

\begin{proposition}[Boffa Hypersets and Rieger--Bernays Slice vs.\ Global Torsors]\label{prop:boffa-torsors}%
\lean{HomotopyCRM.SetsResiduum.quineAPG_bisimilar_and_nonIso_twoCycle, HomotopyCRM.SetsResiduum.boffa_twin_torsor, HomotopyCRM.SetsResiduum.boffa_twin_addTorsor}%
\leanwitness%
\uses{thm:axis-vii-foundation}%
Weakening Aczel's $\mathsf{AFA}$ (which quotients coalgebras by maximum bisimulation) to Maurice Boffa's Anti-Foundation Axiom ($\mathsf{BAFA}$\index{BAFA@$\mathsf{BAFA}$ (Boffa Anti-Foundation Axiom)}~\cite{Boffa69}, requiring every injective decoration of a transitive subgraph of an extensional accessible pointed graph to extend to an injective decoration of the supergraph) lifts $\nu \mathcal{P}_{\mathcal{S}}$ to a universe $V_{\mathsf{BAFA}}$ admitting a proper class of distinct Quine atoms $a = \{a\} \neq b = \{b\}$. Then:
\begin{enumerate}[label=\textup{(\roman*)},leftmargin=1.8em]
  \item \textbf{Isomorphism vs.\ Bisimulation Fibers:} Mere bisimilarity $[A] = [B] \in \nu \mathcal{P}_{\mathcal{S}}$ does not imply $\in$-isomorphism (for instance, a Quine singleton $a = \{a\}$ with $|\mathrm{tc}(\{a\})| = 1$ and a mutual $2$-cycle $u = \{v\}, v = \{u\}$ with $|\mathrm{tc}(\{u\})| = 2$ share the same bisimulation image $q \in \nu \mathcal{P}_{\mathcal{S}}$ while $\mathrm{Iso}_\in(\mathrm{tc}(\{a\}), \mathrm{tc}(\{u\})) = \emptyset$). Restricted to hypersets $A, B \in V_{\mathsf{BAFA}}$ with $\in$-isomorphic transitive closures, $\mathrm{Iso}_\in(\mathrm{tc}(\{A\}), \mathrm{tc}(\{B\}))$ carries a free and transitive right $\mathrm{Aut}_\in(\mathrm{tc}(\{A\}))$-action.
  \item \textbf{Slice Twin Torsors vs.\ Global Parity Torsors Over $Q_0$:}\index{Boffa hyperset torsor} Fix any countably infinite small set $Q_0 \in \mathcal{S}$ of distinct Quine atoms in $V_{\mathsf{BAFA}}$ (obtained by decorating $\omega$ disjoint self-loops), and pass via Rieger--Bernays~\cite{Bernays54,Rieger57} $\in$-permutations fixing $V_{\mathsf{ZF}}$ pointwise on the cumulative subuniverse $V_{\mathsf{ZF}}(Q_0) \subseteq V_{\mathsf{BAFA}}$ to the equivariant permutation topos $\mathcal{E}_{Q_0} := \mathbf{Cont}(\mathrm{Sym}(Q_0)) \simeq \mathbf{Nom}$ (Chapter~\ref{ch:realizability-sites}, \S\ref{subsec:kihara-choice-nominal}). Over the unordered-pair orbit sheaf $[Q_0]^2 \in \mathcal{E}_{Q_0}$, the ordered-twin cover $Q_0^{\#(2)} \twoheadrightarrow [Q_0]^2$ is a pointless $\Delta(\mathbb{Z}/2\mathbb{Z})$-torsor in the \textbf{slice topos} $\mathcal{E}_{Q_0}/[Q_0]^2$,
  \[
    [Q_0^{\#(2)}] \;\neq\; 0 \;\in\; H^1\bigl(\mathcal{E}_{Q_0}/[Q_0]^2;\, \Delta(\mathbb{Z}/2\mathbb{Z})\bigr)
  \]
  (restricting on a single isolated twin pair $Q_2 = \{a, b\}$ to $\mathbf{Set}^{\mathbb{Z}/2\mathbb{Z}}$, Chapter~\ref{ch:choice-atlas}, Theorem~\ref{thm:coshep-svc}). Over the \textbf{global terminal object} $1_{\mathcal{E}_{Q_0}}$, Schreier--Ulam pro-discrete blindness forces $H^1(\mathcal{E}_{Q_0}; \Delta(\mathbb{Z}/2\mathbb{Z})) = 0$, and the non-vanishing global obstruction lifts to the odd-Quine-finset parity torsor $[\mathcal{T}_{\mathrm{odd}}(Q_0)] \neq 0 \in H^1(\mathcal{E}_{Q_0}; \mathrm{EvenFinset}(Q_0))$.
\end{enumerate}
\end{proposition}

\begin{proof}
Under $\mathsf{BAFA}$, extending the empty decoration to disjoint self-loops or to a directed $2$-cycle produces distinct Quine atoms $a = \{a\} \neq b = \{b\}$ and $u = \{v\} \neq v = \{u\}$; cardinality $|\mathrm{tc}(\{a\})| = 1 \neq 2 = |\mathrm{tc}(\{u\})|$ separates their $\in$-isomorphism classes inside $[q] \in \nu \mathcal{P}_{\mathcal{S}}$, whereas for $\in$-isomorphic closures composition gives a free transitive $\mathrm{Aut}_\in(\mathrm{tc}(\{A\}))$-action, proving~(i). For~(ii), passing to the cumulative hyperset subuniverse $V_{\mathsf{ZF}}(Q_0) \subseteq V_{\mathsf{BAFA}}$ generated over $V_{\mathsf{ZF}}$ from $Q_0$ as Quine-atom urelements (via Rieger--Bernays~\cite{Bernays54,Rieger57} and Boffa~\cite{Boffa69}), every permutation $\sigma \in \mathrm{Sym}(Q_0)$ extends uniquely by rank recursion ($\tilde{\sigma}(q) = \sigma(q)$ for $q \in Q_0$ and $\tilde{\sigma}(x) = \{\tilde{\sigma}(y) \mid y \in x\}$ otherwise) to an $\in$-automorphism of $V_{\mathsf{ZF}}(Q_0)$ fixing $V_{\mathsf{ZF}}$ pointwise. For any $\{a, b\} \in [Q_0]^2$, the transposition $(a\;b) \in \mathrm{Sym}(Q_0)$ stabilizes $\{a, b\}$ while swapping $(a, b) \leftrightarrow (b, a)$ in $Q_0^{\#(2)}$, blocking any $\mathrm{Sym}(Q_0)$-equivariant section $[Q_0]^2 \to Q_0^{\#(2)}$ and proving $[Q_0^{\#(2)}] \neq 0$ in $H^1(\mathcal{E}_{Q_0}/[Q_0]^2; \Delta(\mathbb{Z}/2\mathbb{Z}))$, while Chapter~\ref{ch:realizability-sites} (\S\ref{subsec:kihara-choice-nominal}, \texttt{NominalSimplicity.lean}) establishes both $H^1(\mathcal{E}_{Q_0}; \Delta(\mathbb{Z}/2\mathbb{Z})) = 0$ and $[\mathcal{T}_{\mathrm{odd}}(Q_0)] \neq 0 \in H^1(\mathcal{E}_{Q_0}; \mathrm{EvenFinset}(Q_0))$.
\end{proof}

Proposition~\ref{prop:boffa-torsors} acquires an illuminating homotopy-theoretic explanation as soon as we examine the companion Stratum~III clamp: \textbf{Extensionality}.

\begin{proposition}[{External Benchmark: Material Extensionality, Quotient Elimination, and Hereditary Truncation; Friedman~1973~\cite{Friedman73}, Gylterud~2018~\cite[Thm.~3.9]{Gylterud18}}]\label{prop:extensionality-truncation}%
\lean{HomotopyCRM.SetsResiduum.zfset_ext_quot_gate}%
\leanwitness%
\uses{thm:axis-vii-foundation,prop:boffa-torsors}%
Material Extensionality ($\forall z\,(z \in X \leftrightarrow z \in Y) \implies X = Y$) enforces both constructive quotient collapse and homotopical symmetry truncation:
\begin{enumerate}[label=\textup{(\roman*)},leftmargin=1.8em]
  \item \textbf{Constructive Equiconsistency vs.\ Direct Quotient Cost (Friedman~1973~\cite{Friedman73}):} Classical $\mathsf{ZF}$, intuitionistic $\mathsf{IZF}$, and Friedman's non-extensional intuitionistic set theory $\mathsf{IZF}_{\mathrm{non\text{-}ext}}$ are \emph{equiconsistent} ($\mathrm{Con}(\mathsf{IZF}_{\mathrm{non\text{-}ext}}) \implies \mathrm{Con}(\mathsf{ZF})$ via Friedman's recursive double-negation extensionality translation). Nevertheless, $\mathsf{IZF}$ is not directly set-interpretable in $\mathsf{IZF}_{\mathrm{non\text{-}ext}}$ without Collection or primitive quotients: Scott's classical trick of truncating each equivalence class to its small set of minimal-rank members $\{y \sim x \mid \forall z \sim x,\, \mathrm{rank}(y) \le \mathrm{rank}(z)\}$ (working in classical $\mathsf{ZF}$ without Choice) fails constructively because inhabited classes of ordinals need not have $\in$-least elements.
  \item \textbf{Univalent Iterative Multisets vs.\ Hereditary $(-1)$-Truncation (Gylterud~2018~\cite[Thm.~3.9]{Gylterud18}):}\index{Iterative multisets ($\mu \mathcal{P}_{\mathrm{Anima}}$)} Inside an $(\infty, 1)$-topos $\mathcal{X}$ with univalent universe $\mathcal{U}$~\cite[\S 6.1.6]{Lurie09}, let $M := \mu \mathcal{P}_{\mathrm{Anima}}$ be the initial algebra ($W$-type) of the polynomial endofunctor $\mathcal{P}_{\mathrm{Anima}}(X) := \sum_{A : \mathcal{U}} (A \to X)$ of \emph{iterative multisets}. Embedding any small $0$-type $A : \mathcal{U}_{\le 0}$ as the flat multiset $\underline{A} := \mathrm{sup}(A, \mathrm{const}_{\emptyset_M}) : M$ preserves internal symmetries via loop groups
  \[
    \pi_1(M, \underline{A}) \;\simeq\; \mathrm{Aut}(A).
  \]
  Applying naive outer $0$-truncation $\tau_{\le 0}(M)$ kills $\pi_1$ but \emph{retains element multiplicities} (since $X_1 := \mathrm{sup}(\mathbf{1}, \mathrm{const}_{\emptyset_M})$ and $X_2 := \mathrm{sup}(\mathbf{2}, \mathrm{const}_{\emptyset_M})$ satisfy $(X_1 =_M X_2) \simeq (\mathbf{1} \simeq \mathbf{2}) \simeq \emptyset$). Recovering the flat cumulative hierarchy of iterative sets $V_0 \hookrightarrow M$ requires \textbf{hereditary fiberwise $(-1)$-truncation}---restricting each branching map $f : A \to M$ to be a $(-1)$-truncated embedding $A \hookrightarrow M$ (or quotienting recursively by bisimulation)---which simultaneously eliminates multiplicities and forces $V_0$ to be $0$-truncated.
\end{enumerate}
\end{proposition}

\begin{remark}[Unifying Boffa Twin Torsors and Loop Groups via $\nu \mathcal{P}_{\mathrm{Anima}}$]\label{rem:nu-anima-fossil}%
\uses{prop:boffa-torsors,prop:extensionality-truncation}%
Propositions~\ref{prop:boffa-torsors} and~\ref{prop:extensionality-truncation} exhibit Boffa's twin Quine set $T = \{a, b\}$ as the \emph{$0$-truncated $1$-categorical fossil of higher homotopy loop groups}. In well-founded $M = \mu \mathcal{P}_{\mathrm{Anima}}$, even $\underline{\mathbf{2}}$ carries $\pi_1(M, \underline{\mathbf{2}}) \cong \mathbb{Z}/2\mathbb{Z}$, whereas in $V_{\mathsf{BAFA}}$ well-founded sets are rigid and symmetry survives only on non-well-founded twins $a = \{a\} \neq b = \{b\}$ whose Aczel quotient $\{q, q\} = \{q\} \in \nu \mathcal{P}_{\mathcal{S}}$ collapses multiplicity via $(-1)$-truncated image $\exists_d$. Concretely, inside an $(\infty, 1)$-topos $\mathcal{X}$ with nested univalent universes $\mathcal{U}_0 \in \mathcal{U}_1$~\cite[\S 6.1.6]{Lurie09}, set $\mathcal{P}_{\mathrm{Anima}}(X) := \sum_{A : \mathcal{U}_0} (A \to X) : \mathcal{U}_1 \to \mathcal{U}_1$ and $X_n := \mathcal{P}_{\mathrm{Anima}}^n(\mathbf{1})$. Because dependent sums over the constant base $\mathcal{U}_0$ ($|\omega^{\mathrm{op}}| \simeq *$) and exponentials $(A \to -)$ preserve sequential homotopy limits, Ahrens, Capriotti, and Spadotti~\cite{AhrensCapriottiSpadotti15} establish that the \textbf{terminal coalgebra ($M$-type) of non-well-founded univalent multisets}\index{Non-well-founded univalent multisets ($\nu \mathcal{P}_{\mathrm{Anima}}$)} exists at $\omega$ as $\nu \mathcal{P}_{\mathrm{Anima}} := \varprojlim_{n \in \omega^{\mathrm{op}}} X_n : \mathcal{U}_1$, with Lambek equivalence $\mathrm{des} : \nu \mathcal{P}_{\mathrm{Anima}} \xrightarrow{\sim} \mathcal{P}_{\mathrm{Anima}}(\nu \mathcal{P}_{\mathrm{Anima}})$. For the $1$-fold Quine multiset $q$ ($\mathrm{des}(q) = (\mathbf{1}, \mathrm{const}_q)$), finite unrollings $q_0 = * : \mathbf{1}$ and $q_{n+1} = (\mathbf{1}, \mathrm{const}_{q_n}) : X_{n+1}$ satisfy $\Omega(X_{n+1}, q_{n+1}) \simeq \mathrm{Aut}(\mathbf{1}) \times \Omega(X_n, q_n)^{\mathbf{1}} \simeq \mathbf{1}$ by induction on $n$, whence $\Omega(\nu \mathcal{P}_{\mathrm{Anima}}, q) \simeq \varprojlim_n \Omega(X_n, q_n) \simeq \mathbf{1}$. Consequently, for the $2$-fold Quine multiset $[q, q] := \mathrm{des}^{-1}(\mathbf{2}, \mathrm{const}_q)$---explicitly distinct from the $2$-branching self-loop $c = \mathrm{des}^{-1}(\mathbf{2}, \mathrm{const}_c)$, whose loop group is the infinite binary-tree automorphism group $\mathrm{Aut}(T_2)$---Lambek's equivalence and univalence ($\mathrm{const}_q \circ e^{-1} \equiv \mathrm{const}_q$ for $e : \mathbf{2} \simeq \mathbf{2}$) yield $\Omega(\nu \mathcal{P}_{\mathrm{Anima}}, [q, q]) \simeq \sum_{e : \mathbf{2} \simeq \mathbf{2}} \Omega(\nu \mathcal{P}_{\mathrm{Anima}}, q)^{\mathbf{2}} \simeq \mathrm{Aut}(\mathbf{2}) \cong \mathbb{Z}/2\mathbb{Z}$. Lifting Boffa's twin Quine graph $G_T = \{a, b, T\}$ ($a = \{a\}$, $b = \{b\}$, $T = \{a, b\}$) to an untruncated $\mathcal{P}_{\mathrm{Anima}}$-coalgebra without $(-1)$-truncation, terminal coalgebra uniqueness on $(\mathbf{1}, \mathrm{id}_{\mathbf{1}})$ sends $a, b \mapsto q$ and hence $d_{\mathrm{Anima}}(T) = \mathrm{des}^{-1}(\mathbf{2}, d_{\mathrm{Anima}} \circ \iota_{a,b}) = [q, q]$, directly uniting Axis~V with Axis~VII.
\end{remark}

\begin{table}[htbp]
\centering
\caption{Master Algebraic Set Theory Residuum Ledger Across $(\mathcal{C}, \mathcal{S})$}\label{tab:master-ast-ledger}
\vspace{0.3em}
{\scriptsize\setlength{\tabcolsep}{2.5pt}\renewcommand{\arraystretch}{1.12}%
\begin{tabularx}{\textwidth}{@{}P{0.15\textwidth}P{0.18\textwidth}P{0.15\textwidth}YY@{}}
\toprule
\textbf{Axiom Family} & \textbf{Role in $(\mathcal{C}, \mathcal{S})$} & \textbf{Polarity} & \textbf{Separating Model} & \textbf{Residuum $\mathrm{Res}(\kappa)$} \\
\midrule
Empty, Pair, Union & $0, 1+1$, finite sums in $\mathcal{C},\mathcal{S}$ & Plumbing & Non-extensive cats & None (ambient grammar) \\
Bernays G1--G8, $\Delta_0$-Sep & Finite limits + $\Delta_0$-pullbacks & Carving & Pre-regular categories & Missing subobject pullbacks \\
Unbounded Sep / $\mathsf{KM}$ & $\mathcal{L}_\in$ / class reflection to $\mathcal{S}$ & Compression & Mathias~\cite{Mathias01} $\mathsf{Mac} \subsetneq \mathsf{ZC}$; $\mathsf{NBG}$ vs.\ $\mathsf{KM}$ & Unreflected $\Pi_n^V$ / $\Pi^1_1$ invariants \\
Power Set & $\mathcal{P}_{\mathcal{S}}(X) \in \mathcal{S}$ ($\Omega^X$ small) & Compression & $H(\aleph_1)$ ($\mathsf{ZFC}^-$); $\mathsf{CZF}$ & Proper-class $2^{\mathbb{N}} \notin \mathcal{S}$; $\Omega^X / 2^X$ \\
Replacement (Axis~VI) & Image descent $\mathrm{im}(f) \in \mathcal{S}$ ($\exists_g$) & Compression & Zermelo $V_{\omega+\omega}$ ($\mathsf{ZC}$ / $\mathsf{ETCS}$) & Colimit escape; $\boldsymbol{\Sigma}^0_\omega$ Borel $\mathsf{ZC}$-unprovability \\
Limitation of Size ($\mathsf{LoS}$) & $C \notin \mathcal{S} \iff C \twoheadrightarrow V$ & Mega-Axiom & $(L_{\omega_\omega^L}, \mathrm{Def})$ over $\mathsf{GB}_0$; Easton & Bundles $\mathrm{Rep} + \mathsf{GC}$ ($\mathrm{Ord}\cong V$) \\
Foundation (Axis~VII) & $\mathbf{Coalg}_{\mathrm{wf}} \hookrightarrow \mathbf{Coalg}$ ($\nu^{\mathrm{wf}}\mathcal{P}_{\mathcal{S}} \simeq \mu\mathcal{P}_{\mathcal{S}}$) & Anti-Cycle & Aczel $V_{\mathsf{AFA}}$; Boffa $V_{\mathsf{BAFA}}$ & $\nu \mathcal{P}_{\mathcal{S}} \setminus \mu \mathcal{P}_{\mathcal{S}}$; slice/parity torsors \\
Extensionality & Hereditary $(-1)$-truncation & Anti-Symmetry & $\mathsf{IZF}_{\mathrm{non\text{-}ext}}$; Gylterud $M$ & Quotient cost; loops $\pi_1(M,\underline{A})$ \\
\bottomrule
\end{tabularx}\par}
\end{table}

\section{\texorpdfstring{Scale-Invariance, the Double-Clamped Universe, and Horizons Beyond ``$\mathbf{Sets}$''}{Scale-Invariance, the Double-Clamped Universe, and Horizons Beyond "Sets"}}\label{sec:scale-invariance-synthesis}

Up through Chapter~\ref{ch:realizability-sites}, a foundational skeptic could still raise a ``turtles all the way down'' objection: Categorical Reverse Mathematics might measure obstructions \emph{inside} an ambient topos or $(\infty, 1)$-category while leaving the meta-theoretic bedrock of ``$\mathbf{Sets}$'' to classical syntactic proof theory. Passing to Joyal--Moerdijk small-map pairs $(\mathcal{C}, \mathcal{S})$ and the baseline adjunction $\mathrm{Def} \dashv \pi_{\mathrm{Set}}$ closes that loop completely by establishing \textbf{reflexive scale-invariance} across the universe boundary.

\begin{theorem}[Reflexive Scale-Invariance of Double-Clamped $\mathbf{Set}_{\mathsf{ZFC}}$ and $\mathsf{NBG}$]\label{thm:scale-invariance}%
\uses{def:set-projection-fibration,prop:large-equiv-splitting,def:small-map-ambient,def:axis-vi-replacement,thm:axis-vii-foundation,prop:extensionality-truncation}%
Let $(\mathcal{C}, \mathcal{S})$ be a non-degenerate ($0 \not\cong 1$) Barr-exact predicative AST ambient (Definition~\ref{def:small-map-ambient}, with monomorphisms in $\mathcal{S}$ so $\Omega := \mathcal{P}_{\mathcal{S}}(1)$ classifies subobjects in $\mathcal{C}$) admitting $\mu \mathcal{P}_{\mathcal{S}}$ and $\nu \mathcal{P}_{\mathcal{S}}$ under Axis~VI image descent, with well-founded cumulative universe $V := \nu^{\mathrm{wf}} \mathcal{P}_{\mathcal{S}} \simeq \mu \mathcal{P}_{\mathcal{S}}$:
\[
  {\small\begin{aligned}
    \text{\textbf{Upper Clamps:}}\quad &\mathbb{N} \in \mathcal{S},\;\; 1+1 \xrightarrow{\sim} \Omega,\;\; \mathcal{P}_{\mathcal{S}}(\mathcal{S}) \subseteq \mathcal{S},\;\; \mathrm{im}(I \xrightarrow{f} C) \in \mathcal{S}, \\
    &\mathrm{IAC}_{\mathcal{S}}\;\; \bigl(\text{+\;}\mathrm{Sub}_{\mathcal{C}}(1)\cong 2,\; 1\text{ proj.; }\mathbf{Card}(V) \hookrightarrow \mathbf{Set}(V)\text{ for }\mathsf{NBG}\bigr) \\[0.25em]
    &\Downarrow \;\; \text{\emph{pinches $(V, \mathcal{C}, \mathcal{S})$ into}}\;\; \mathbf{Set}_{\mathsf{ZFC}} \;\;\Uparrow \\[0.25em]
    \text{\textbf{Lower Clamps:}}\quad &\mathbf{Coalg}_{\mathrm{wf}}(\mathcal{P}_{\mathcal{S}}) \hookrightarrow \mathbf{Coalg}(\mathcal{P}_{\mathcal{S}})\; (\nu^{\mathrm{wf}}\mathcal{P}_{\mathcal{S}} \simeq \mu\mathcal{P}_{\mathcal{S}} \subsetneq \nu\mathcal{P}_{\mathcal{S}}), \\
    &\text{Extensional }\mathcal{P}_{\mathcal{S}}\text{ (subobjects) + Osius pure-set fullness}
  \end{aligned}}
\]
\begin{enumerate}[label=\textup{(\roman*)},leftmargin=1.8em]
  \item \textbf{Internal Double-Clamped $\mathsf{ZF}/\mathsf{ZFC}$ and $\mathsf{GB}/\mathsf{GBC}/\mathsf{NBG}$ Soundness:} Imposing the four Choice-free upper clamps ($\mathbb{N} \in \mathcal{S}$, $1 + 1 \xrightarrow{\sim} \Omega := \mathcal{P}_{\mathcal{S}}(1)$, $\mathcal{P}_{\mathcal{S}}(\mathcal{S}) \subseteq \mathcal{S}$, and Axis~VI image descent $\mathrm{im}(I \xrightarrow{f} C) \in \mathcal{S}$) together with the lower structural clamps (extensional subobject membership via $\mathcal{P}_{\mathcal{S}}$ and Axis~VII Foundation restricting admissible set-coalgebras along $\mathbf{Coalg}_{\mathrm{wf}}(\mathcal{P}_{\mathcal{S}}) \hookrightarrow \mathbf{Coalg}(\mathcal{P}_{\mathcal{S}})$ so $V = \nu^{\mathrm{wf}} \mathcal{P}_{\mathcal{S}} \simeq \mu \mathcal{P}_{\mathcal{S}}$) makes small objects $\mathcal{E}_{\mathcal{S}} \subseteq \mathcal{C}$ a Boolean elementary topos with $\mathrm{NNO}$ and forces $\mathcal{C} \models ((V, \in) \models \mathsf{ZF})$ and internal $\mathcal{C} \models ((V, \mathrm{Sub}_{\mathcal{C}}(V)) \models \mathsf{GB})$. Every $X \in \mathcal{E}_{\mathcal{S}}$ is locally a pure set in $V$ (admitting a regular epimorphism $p : U \twoheadrightarrow 1$ in $\mathcal{E}_{\mathcal{S}}$ and a monomorphism $p^* X \rightarrowtail p^* V$ in $\mathcal{C}/U$) if and only if $(\mathcal{C}, \mathcal{S})$ satisfies \emph{local Osius pure-set fullness} (every $X \in \mathcal{S}$ embeds over some $U \twoheadrightarrow 1$ into a fiberwise-extensional well-founded coalgebra $(T, R) \in \mathbf{Coalg}_{\mathrm{wf}}(\mathcal{P}_{\mathcal{S}})$; failing on the atom object $\mathbb{A}$ of Fraenkel--Mostowski urelement models $\mathcal{V}_{\mathrm{FM}} \models \mathsf{ZFA} + \neg\mathrm{AC}$). Adjoining internal small-map choice $\mathrm{IAC}_{\mathcal{S}}$ upgrades the internal theory to $\mathcal{C} \models ((V, \in) \models \mathsf{ZFC})$ and $\mathcal{C} \models ((V, \mathrm{Sub}_{\mathcal{C}}(V)) \models \mathsf{GBC})$ (and $\mathcal{C} \models ((V, \mathrm{Sub}_{\mathcal{C}}(V)) \models \mathsf{NBG})$ iff $\mathbf{Card}(V) \hookrightarrow \mathbf{Set}(V)$ splits in $\mathbf{Cat}(\mathcal{C})$) and discharges local Osius pure-set fullness automatically.
  \item \textbf{External Well-Pointed Pinch to $\mathbf{Set}_{\mathsf{ZFC}}$ and Global $\mathcal{E}_{\mathcal{S}} \simeq \mathbf{El}(V)$:} Under the Choice-free upper and lower clamps together with local Osius pure-set fullness (automatically discharged under $\mathrm{IAC}_{\mathcal{S}}$), $\mathcal{E}_{\mathcal{S}}$ pinches externally to a well-pointed classical category of pure sets ($\mathcal{E}_{\mathcal{S}} \simeq \mathbf{El}(V)$, modeling $\mathsf{ETCS}$ + Osius--Joyal--Moerdijk Replacement: $\mathbf{Set}_{\mathsf{ZF}}$ without Choice and $\mathbf{Set}_{\mathsf{ZFC}}$ with $\mathrm{IAC}_{\mathcal{S}}$) if and only if $\mathcal{C}$ additionally satisfies external two-valuedness ($\mathrm{Sub}_{\mathcal{C}}(1) \cong \{0, 1\}$) and external projectivity of $1$ in $\mathcal{E}_{\mathcal{S}}$ (splitting epic covers $U \twoheadrightarrow 1$ globally and upgrading $\mathrm{IAC}_{\mathcal{S}}$ to external $\mathrm{AC}_{\mathcal{S}}$; failing on atomless Boolean-valued $(\mathbf{Sh}_{\mathrm{Class}}(B), \mathbf{Sh}_{\mathrm{Set}}(B))$ without two-valuedness and on $(\mathbf{Class}^{\mathbb{Z}/2\mathbb{Z}}, \mathbf{Set}^{\mathbb{Z}/2\mathbb{Z}})$ without external projectivity of $1$).
  \item \textbf{Residuum Scale-Invariance:} Each of the upper compression and extensionality clamps above---Infinity ($\mathbb{N} \in \mathcal{S}$, Axis~III), Excluded Middle ($1 + 1 \xrightarrow{\sim} \Omega$, Axis~II), Power Set ($\mathcal{P}_{\mathcal{S}}(\mathcal{S}) \subseteq \mathcal{S}$, Axis~II), Replacement / image descent (Axis~VI), Choice ($\mathrm{IAC}_{\mathcal{S}}$, external projectivity of $1$, and large splitting $\mathbf{Card}(V) \hookrightarrow \mathbf{Set}(V)$, Axis~I), and Extensionality (subobject $\mathcal{P}_{\mathcal{S}}$ and \texttt{Quot.sound} bisimulation, Axis~VII)---coincides with the vanishing $\mathrm{Res}(\kappa_A) \simeq 0$ of an explicit admissible comparison triple $\mathfrak{T}_A = (\mathcal{D}_0 \supseteq \mathcal{D}_1, \kappa_A) \in \mathrm{Dom}(\mathrm{Res})$, while Axis~VII Foundation corresponds strictly to restricting admissible set-coalgebras to $\mathbf{Coalg}_{\mathrm{wf}}(\mathcal{P}_{\mathcal{S}}) \hookrightarrow \mathbf{Coalg}(\mathcal{P}_{\mathcal{S}})$---that is, vanishing of the factorization obstruction across $\kappa_{\mathrm{Reg}} : \nu^{\mathrm{wf}}\mathcal{P}_{\mathcal{S}} \simeq \mu\mathcal{P}_{\mathcal{S}} \hookrightarrow \nu\mathcal{P}_{\mathcal{S}}$ for every admissible set-coalgebra $(X, \gamma)$ (Theorem~\ref{thm:axis-vii-foundation}(iii)), rather than forcing $\nu\mathcal{P}_{\mathcal{S}} \setminus \mu\mathcal{P}_{\mathcal{S}} \simeq 0$---all calibrated in Sections~\ref{sec:conservativity-kaleidoscope}--\ref{sec:anti-existence-residua} (Table~\ref{tab:master-ast-ledger}).
\end{enumerate}
\end{theorem}

\begin{proof}
For~(i), any morphism $f : U \to X$ in $\mathcal{C}$ between small objects $U, X \in \mathcal{S}$ factors as $U \xrightarrow{\Gamma_f} U \times X \xrightarrow{\mathrm{pr}_2} X$ where $\Gamma_f = (f \times \mathrm{id}_X)^*(\Delta_X) \in \mathcal{S}$ and $\mathrm{pr}_2 = X^*(U \to 1) \in \mathcal{S}$~\cite[Lemma~I.1.2]{JoyalMoerdijk95}, so $f \in \mathcal{S}$ and $\mathrm{Sub}_{\mathcal{E}_{\mathcal{S}}}(X) = \mathrm{Sub}_{\mathcal{S}}(X)$. Since $\mathcal{E}_{\mathcal{S}}$ has finite limits and power objects $\mathcal{P}_{\mathcal{S}}(X) \in \mathcal{S}$ for all $X \in \mathcal{S}$, Mikkelsen--Par\'e's theorem~\cite[Thm.~A2.2.3 \& Cor.~A2.2.4]{Johnstone02} makes $\mathcal{E}_{\mathcal{S}}$ an elementary topos with $\mathrm{NNO}$ $\mathbb{N} \in \mathcal{S}$, while $1 + 1 \xrightarrow{\sim} \Omega := \mathcal{P}_{\mathcal{S}}(1)$ classifying all subobjects of $\mathcal{C}$ makes both $\mathcal{C}$ and $\mathcal{E}_{\mathcal{S}}$ Boolean ($\forall_\pi = \neg\exists_\pi\neg$ along $Z \times V \to Z$). Under Axis~VI covariant direct image $\exists_g$, Lambek's isomorphism $\mathrm{sup} : \mathcal{P}_{\mathcal{S}}(V) \xrightarrow{\sim} V$ enforces Extensionality on $V = \nu^{\mathrm{wf}}\mathcal{P}_{\mathcal{S}} \simeq \mu\mathcal{P}_{\mathcal{S}}$ while Theorem~\ref{thm:axis-vii-foundation}(i) enforces Foundation ($\in$-induction). Composition closure of $\mathcal{S}$ together with Axis~VI image descent yields Pairing and Small Union (since $\pi_2 : \in_V \to V$ lies in $\mathcal{S}$, with fiber $\pi_2^{-1}(y) = \{x \mid x \in y\}$), $\mathcal{P}_{\mathcal{S}}(\mathcal{S}) \subseteq \mathcal{S}$ yields Power Set~\cite[Thm.~I.2.1]{JoyalMoerdijk95}, and $\mathbb{N} \in \mathcal{S}$ with strict order $\gamma_<(n) = \{k < n\}$ collapses via Osius terminality~\cite{Osius74} into $\omega \in V$ (Infinity). Boolean regular class operations (\textbf{G1--G8}) and pullback stability of $\mathcal{S}$ give Subclass Separation (hence unbounded $\mathcal{L}_\in$-Separation on $V$), while Axis~VI image descent yields Class Replacement, proving $\mathcal{C} \models ((V, \in) \models \mathsf{ZF})$ and $\mathcal{C} \models ((V, \mathrm{Sub}_{\mathcal{C}}(V)) \models \mathsf{GB})$. For Choice-free pure-set fullness, given any fiberwise-extensional well-founded coalgebra $(T \xrightarrow{\pi} U, R \rightarrowtail T \times_U T)$ in $\mathcal{C}/U$ ($T \in \mathcal{S}$), viewing $R \rightarrowtail T \times T$ as $(T, \tilde{\gamma}_T) \in \mathbf{Coalg}_{\mathrm{wf}}(\mathcal{P}_{\mathcal{S}})$ yields by Osius terminality a unique $m : T \to V$ with $\mathrm{sup}^{-1} \circ m = \exists_m \circ \tilde{\gamma}_T$; by $\in_R$-induction on $t \in T$, $\pi(t) = \pi(t') \wedge m(t) = m(t')$ forces $\tilde{\gamma}_T(t) = \tilde{\gamma}_T(t')$ and hence $t = t'$ by fiberwise extensionality, making $(\pi, m) : T \rightarrowtail U \times V = p^* V$ monic in $\mathcal{C}/U$. Conversely, any monomorphism $p^* X \rightarrowtail p^* V$ over $U$ generates a small fiberwise-transitive closure $T := \bigcup_{n \in \mathbb{N}} \bigcup^n_U \mathrm{im}(p^* X) \rightarrowtail U \times V$ via $\mathbb{N} \in \mathcal{S}$, Small Union, and Axis~VI Replacement, which restricts $\mathrm{sup}^{-1}$ to a fiberwise-extensional well-founded coalgebra containing $p^* X$. Adjoining $\mathrm{IAC}_{\mathcal{S}}$ yields $\mathcal{C} \models ((V, \in) \models \mathsf{ZFC})$ and $\mathcal{C} \models ((V, \mathrm{Sub}_{\mathcal{C}}(V)) \models \mathsf{GBC})$ (and $\mathcal{C} \models ((V, \mathrm{Sub}_{\mathcal{C}}(V)) \models \mathsf{NBG})$ iff $\mathbf{Card}(V) \hookrightarrow \mathbf{Set}(V)$ splits, Proposition~\ref{prop:large-equiv-splitting}), and by Zermelo's theorem in Boolean $\mathcal{E}_{\mathcal{S}} \models \mathrm{IAC}$~\cite[Prop.~D4.5.7]{Johnstone02} every $X \in \mathcal{S}$ carries a $U$-indexed well-ordering $<$ over its inhabited cover $U := \mathrm{WO}(X) \twoheadrightarrow 1$ in $\mathcal{E}_{\mathcal{S}}$, whose strict lower segments $\gamma_<(u, x) = \{(u, y) \mid y <_u x\}$ are fiberwise extensional and well-founded, automatically discharging local Osius pure-set fullness.

For~(ii), a non-degenerate Boolean elementary topos $\mathcal{E}_{\mathcal{S}}$ is well-pointed ($1$ generates $\mathcal{E}_{\mathcal{S}}$) if and only if $\mathrm{Sub}_{\mathcal{C}}(1) = \mathrm{Sub}_{\mathcal{E}_{\mathcal{S}}}(1) \cong \{0, 1\}$ and $1$ is externally projective in $\mathcal{E}_{\mathcal{S}}$~\cite[Prop.~VI.10.1]{MacLaneMoerdijk92} (for $f \neq g : X \rightrightarrows Y$, $X \setminus \mathrm{eq}(f, g) \not\cong 0$ has epic support onto $1$ by two-valuedness and hence a separating global point $1 \to X \setminus \mathrm{eq}(f, g)$ by projectivity of $1$). External projectivity of $1$ in $\mathcal{E}_{\mathcal{S}}$ splits every epic cover $p : U \twoheadrightarrow 1$ via a global point $s : 1 \to U$, pulling local pure-set embeddings $p^* X \rightarrowtail p^* V$ back to global monomorphisms $m_s : X \rightarrowtail V$ classified by $a_X := \mathrm{sup}(\ulcorner m_s \urcorner) : 1 \to V$ ($X \cong E_{a_X}$, establishing global $\mathcal{E}_{\mathcal{S}} \simeq \mathbf{El}(V)$) and upgrading $\mathrm{IAC}_{\mathcal{S}}$ to external $\mathrm{AC}_{\mathcal{S}}$ (Chapter~\ref{ch:choice-atlas}, Theorem~\ref{thm:iac-to-ac}), yielding Lawvere's $\mathsf{ETCS}$ plus Osius--Joyal--Moerdijk Replacement~\cite{Cole73,Mitchell72,Osius74}. Conversely, atomless $(\mathbf{Sh}_{\mathrm{Class}}(B), \mathbf{Sh}_{\mathrm{Set}}(B))$ models Boolean-valued $V^{(B)} \models \mathsf{ZFC}$ with $\mathrm{Sub}_{\mathcal{C}}(1) \cong B \not\cong \{0, 1\}$, while $(\mathbf{Class}^{\mathbb{Z}/2\mathbb{Z}}, \mathbf{Set}^{\mathbb{Z}/2\mathbb{Z}})$ satisfies two-valuedness and $\mathrm{IAC}_{\mathcal{S}}$ yet fails external projectivity of $1$ (Chapter~\ref{ch:choice-atlas}, Theorem~\ref{thm:coshep-svc}, leaving $V = \mu\mathcal{P}_{\mathcal{S}}$ with trivial $\mathbb{Z}/2\mathbb{Z}$-action so $X = (\{0, 1\}, \mathrm{swap}) \notin \mathbf{El}(V)$ globally). Finally, (iii) collects the explicit comparison triples calibrated across Sections~\ref{sec:conservativity-kaleidoscope}--\ref{sec:anti-existence-residua} (Table~\ref{tab:master-ast-ledger}).
\end{proof}

Classical $\mathbf{Set}_{\mathsf{ZFC}}$ is therefore neither minimal nor canonical; inside a Joyal--Moerdijk small-map ambient $(\mathcal{C}, \mathcal{S})$ it is a \textbf{double-clamped locus} carved out by a clean \textbf{duality of axiomatic excess}:
\begin{enumerate}[label=\textup{\arabic*.},leftmargin=1.8em]
  \item \textbf{Over-Compressing Existence and Logic (\emph{Infinity, Excluded Middle, Power Set, Replacement, Choice}):} Forcing infinite initial algebras $\mathbb{N}$, classical truth values $1 + 1 \xrightarrow{\sim} \Omega$, powerclasses $\mathcal{P}_{\mathcal{S}}(X)$, small-indexed images and countable colimits $\bigcup_{n \in \mathbb{N}} \mathcal{P}^n(\mathbb{N})$, and epimorphism sections back down into $\mathcal{S}$ (plus external two-valuedness and projectivity of $1$ in $\mathcal{E}_{\mathcal{S}}$ for well-pointed $\mathcal{E}_{\mathcal{S}} \simeq \mathbf{El}(V)$, and $\mathbf{Card}(V) \hookrightarrow \mathbf{Set}(V)$ for $\mathsf{NBG}$).
  \item \textbf{Over-Quotienting Structure (\emph{Foundation, Extensionality, Pure-Set Fullness}):} Restricting admissible set-coalgebras along $\mathbf{Coalg}_{\mathrm{wf}}(\mathcal{P}_{\mathcal{S}}) \hookrightarrow \mathbf{Coalg}(\mathcal{P}_{\mathcal{S}})$ ($\nu^{\mathrm{wf}} \mathcal{P}_{\mathcal{S}} \simeq \mu \mathcal{P}_{\mathcal{S}} \subsetneq \nu \mathcal{P}_{\mathcal{S}}$) and enforcing extensional subobjects plus Osius pure-set fullness. In the univalent setting of Proposition~\ref{prop:extensionality-truncation}(ii), inside an $(\infty, 1)$-topos $\mathcal{X}$ with nested univalent universes $\mathcal{U}_0 \in \mathcal{U}_1$~\cite[\S 6.1.6]{Lurie09}, Gylterud's hereditary fiberwise $(-1)$-truncation $V_0 \hookrightarrow M = \mu \mathcal{P}_{\mathrm{Anima}}$~\cite[Thm.~3.9]{Gylterud18} directly realizes $V_0 \cong \mu \mathcal{P}_{\mathcal{S}}$ inside $(\tau_{\le 0}\mathcal{X}_{\mathcal{U}_1}, \mathcal{S}_{\mathcal{U}_0})$ (via $\sum_{A : \mathcal{U}_0} (A \hookrightarrow X) \simeq \mathcal{P}_{\mathcal{S}}(X)$), while developing a universe-free elementary higher class-anima AST $(\mathcal{C}_\infty, \mathcal{S}_\infty)$ forms a central programmatic synthesis horizon for Volume~II.
\end{enumerate}

\subsection{\texorpdfstring{The Kernel Audit Gate: \texttt{HomotopyCRM.SetsResiduum}}{The Kernel Audit Gate: HomotopyCRM.SetsResiduum}}\label{subsec:sets-residuum-kernel-gate}

This double-clamped geometry\index{SetsResiduum@\texttt{SetsResiduum} Lean~4 module} maps directly onto the \texttt{\#print axioms} kernel audit gate enforced throughout our Lean~4 formalization suite \texttt{HomotopyCRM} (\texttt{leanprover/lean4:v4.28.0} over \texttt{mathlib4}). Inside Lean's Calculus of Inductive Constructions, universe-polymorphic \texttt{Classical.choice} realizes uniform global equivalence-splitting across cumulative universes \texttt{Type u}, while propositional extensionality (\texttt{propext}) and quotient soundness (\texttt{Quot.sound}, from which Lean derives function extensionality \texttt{funext}) implement $0$-truncated extensionality and recursive bisimulation quotienting---collapsing intensional tree presentations (\texttt{PSet}, modeling $\mu \mathcal{P}_{\mathrm{Anima}}$ at $0$-types) into extensional sets (\texttt{ZFSet}).

To separate all three kernel primitives cleanly under the monograph's \textbf{4-Layer Definition Check Gate} (matching the \texttt{\#guard\_msgs in \#print axioms} standard of the earlier formalization modules), our companion module \texttt{HomotopyCRM/\allowbreak SetsResiduum.lean} pairs its bundled two-sorted \texttt{ClassSetModel} adjunction (\texttt{defSetProjAdjunction}, \texttt{finToPowElemHom : SetElemHom}, \texttt{finToPowClassHom : ClassSetHom}) and $8$-part pullback counterexample with three deliberate design guardrails (Table~\ref{tab:sets-residuum-audit}). First, \texttt{card\_\allowbreak split\_\allowbreak iff\_\allowbreak global\_\allowbreak choice} avoids hidden \texttt{funext} (\texttt{Quot.sound}) dependencies by using pointwise homotopy-functor equalities (\texttt{forall x, ...}) over an abstract skeletal carrier, paired with \texttt{type\_\allowbreak equiv\_\allowbreak splitting\_\allowbreak global\_\allowbreak choice} on \texttt{CategoryTheory.Functor} over \texttt{Type u}. Second, because every term of Lean's inductive \texttt{PSet} is well-founded, \texttt{quine\_\allowbreak atom\_\allowbreak outside\_\allowbreak wf} embeds \texttt{PSet} into a pointed-graph coalgebra carrier \texttt{APG} and uses \texttt{PSet.rec} (\texttt{axiom}-free) to refute bisimilarity between the self-loop coalgebra \texttt{quineAPG} ($q = \{q\}$) and any \texttt{x : PSet}. Third, \texttt{zfset\_\allowbreak ext\_\allowbreak quot\_\allowbreak gate} invokes \texttt{Quotient.sound} directly over \texttt{PSet.Equiv} so its kernel footprint stays strictly \texttt{propext, Quot.sound} without pulling \texttt{Classical.choice}.

\begin{table}[htbp]
\centering
\caption{Kernel Axiom Audit Ledger for \texttt{HomotopyCRM.SetsResiduum}}\label{tab:sets-residuum-audit}
\vspace{0.3em}
{\scriptsize\setlength{\tabcolsep}{2.5pt}\renewcommand{\arraystretch}{1.12}%
\begin{tabularx}{\textwidth}{@{}P{0.27\textwidth}YP{0.21\textwidth}@{}}
\toprule
\textbf{Lean 4 Declaration} & \textbf{Calibrated Universe Clamp / Residuum} & \textbf{\texttt{\#print axioms} Gate} \\
\midrule
\texttt{defSetProjAdjunction}, \texttt{finToPowElemHom}, \texttt{finToPowClassHom} \& \texttt{naive\_\allowbreak pullback\_\allowbreak breaks\_\allowbreak subclass\_\allowbreak separation} & Definition~\ref{def:set-projection-fibration}: two-sorted \texttt{ClassSetModel} coreflection $\mathrm{Def} \dashv \pi_{\mathrm{Set}}$, elementary embedding \texttt{finToPowElemHom : SetElemHom} (\texttt{elem\_ext}), and $8$-part pullback Subclass-Separation counterexample (\texttt{finOmegaModel} $\hookrightarrow$ \texttt{powOmegaModel}, \texttt{evenPowClass}) & \texttt{propext}, \texttt{Classical.choice}, \texttt{Quot.sound} \\
\texttt{card\_\allowbreak split\_\allowbreak iff\_\allowbreak global\_\allowbreak choice} & Proposition~\ref{prop:large-equiv-splitting}: pointwise homotopy-functor splitting of skeletal $\mathbf{Card} \hookrightarrow \mathbf{Set}$ is equivalent to uniform small-set choice $G(X)=\eta_X^{-1}(\emptyset)$ & \texttt{(axiom-free)} \\
\texttt{type\_\allowbreak equiv\_\allowbreak splitting\_\allowbreak global\_\allowbreak choice} & Proposition~\ref{prop:large-equiv-splitting} + \texttt{CategoryTheory.Functor}: global functorial skeleton splitting of \texttt{Type u} via \texttt{Classical.choice} & \texttt{propext}, \texttt{Classical.choice}, \texttt{Quot.sound} \\
\texttt{pset\_\allowbreak equiv\_\allowbreak iff\_\allowbreak bisimilar} \& \texttt{quine\_\allowbreak atom\_\allowbreak outside\_\allowbreak wf} & Theorem~\ref{thm:axis-vii-foundation}: \texttt{PSet.Equiv} embeds faithfully into \texttt{APG} bisimulation while \texttt{quineAPG} ($q=\{q\}$) lies outside the image of \texttt{PSet} & \texttt{(axiom-free)} \\
\texttt{quineAPG\_\allowbreak bisimilar\_\allowbreak and\_\allowbreak nonIso\_\allowbreak twoCycle}, \texttt{boffa\_\allowbreak twin\_\allowbreak torsor} (\& \texttt{boffa\_\allowbreak twin\_\allowbreak addTorsor}) & Proposition~\ref{prop:boffa-torsors}: $2$-cycle non-isomorphic bisimulation witness and ordered vs.\ unordered twin Quine swap torsor (\texttt{Bool} core vs.\ Layer-3 \texttt{AddTorsor (ZMod 2)}) & \texttt{(axiom-free)} (Layer~3: \texttt{propext}, \texttt{Quot.sound}) \\
\texttt{zfset\_\allowbreak ext\_\allowbreak quot\_\allowbreak gate} & Proposition~\ref{prop:extensionality-truncation}: \texttt{Quotient.sound} over \texttt{PSet.Equiv} collapsing branch permutations and multiplicity $[\emptyset,\emptyset] \sim [\emptyset]$ & \texttt{propext}, \texttt{Quot.sound} \\
\bottomrule
\end{tabularx}\par}
\end{table}

\subsection{\texorpdfstring{Bridge to Chapters~\ref{ch:panorama}--\ref{ch:master-atlas} and Volume~II}{Bridge to Chapters 8--9 and Volume II}}\label{subsec:sets-bridge-ch8}

Above all, dissecting ``$\mathbf{Sets}$'' closes the constructive and topos-theoretic arc of Volume~I by naming what lay just beyond its elementary-topos horizon. Recording these universe clamps prepares the programmatic panorama and Master Reference Atlas of Chapters~\ref{ch:panorama}--\ref{ch:master-atlas} and defines the four-part research program of Volume~II (\emph{Universes, Classes, and Higher Fixed Points}):
\begin{itemize}[leftmargin=1.5em]
  \item \textbf{Axis~VI (Quotient-Image Descent and Borel Proof Unwinding):} Calibrates small-indexed image descent separating $\mathsf{ETCS} \simeq \mathsf{Mac}$ and $\mathsf{ZC}$ ($V_{\omega+\omega}$) from $\mathsf{ZFC}$, extracting the explicit strategy-descent residuum $\mathrm{Res}(\kappa_{\mathrm{Det}})$ on Martin's proper-class unraveling coalgebra in $\mathcal{C} \setminus \mathcal{S}$ at $\boldsymbol{\Sigma}^0_\omega$~\cite{Friedman71,Martin75} (Remark~\ref{rem:martin-borel-coalgebra}).
  \item \textbf{Axis~VII (Coalgebraic and Univalent Universe Fixed Points):} Develops the full spectrum between $V_{\mathsf{ZF}} = \mu \mathcal{P}_{\mathcal{S}}$, Aczel's $V_{\mathsf{AFA}} = \nu \mathcal{P}_{\mathcal{S}}$, Boffa's Rieger--Bernays torsors over $V_{\mathsf{BAFA}}$~\cite{Bernays54,Rieger57,Boffa69}, Gylterud's $M = \mu \mathcal{P}_{\mathrm{Anima}}$~\cite{Gylterud18}, and the terminal univalent multiset coalgebra $\nu \mathcal{P}_{\mathrm{Anima}}$ where $\Omega(\nu \mathcal{P}_{\mathrm{Anima}}, [q, q]) \cong \mathbb{Z}/2\mathbb{Z}$ unites Axes~V and~VII (Remark~\ref{rem:nu-anima-fossil}).
  \item \textbf{$2$-Categorical Calibration of the $(\infty, \infty)$-Ladder Without Universes:} Replaces Grothendieck universes ($V_\kappa$) across Chapter~\ref{ch:higher-crm}'s four-tier hierarchy $\mathrm{Cat}^L_{(\infty,\infty)} \subsetneq \mathrm{Cat}^\omega_{(\infty,\infty)} \subsetneq \mathrm{Cat}^{\mathrm{coind}}_{(\infty,\infty)} \subsetneq \mathrm{Cat}^R_{(\infty,\infty)}$~\cite{ORW26} with exact class calibrations across $\mathsf{GBC} \subsetneq \mathsf{NBG} \subsetneq \mathsf{NBG}+\mathsf{ETR}_{\mathrm{Ord}} \subsetneq \mathsf{NBG}+\mathsf{ETR} \subsetneq \mathsf{KM}$~\cite[Thm.~1.2]{Gitman20} (noting that Gitman et al.\ write ``$\mathsf{GBC}$'' for $\mathsf{NBG} = \mathsf{GBC} + \mathsf{GC}$, and distinguishing length-$\mathrm{Ord}$ $\mathsf{ETR}_{\mathrm{Ord}}$ from full $\mathsf{ETR}$ along arbitrary class well-orderings)---testing where $\alpha = 0$ stabilization on $\mathrm{Cat}_{\le \omega}$ (Theorem~\ref{thm:ordinal-rank-completeness}, including $\mathcal{Q}$) drops transfinite $\mathsf{ETR}_{\mathrm{Ord}}$ recursion down to $\mathsf{GBC}$, separating the $\mathrm{Con}(\mathsf{ZFC})$ threshold crossed already at $\mathsf{ETR}_\omega < \mathsf{ETR}_{\mathrm{Ord}}$ from full $\mathsf{ETR}$ and impredicative $\mathsf{KM}$ comprehension ($\mathrm{Con}(\mathsf{NBG}+\mathsf{ETR}_{\mathrm{Ord}})$).
  \item \textbf{Non-Conservative Base-Change Transgression:} Formulates a quantitative Leray-style transgression map along $\mathrm{Def} \dashv \pi_{\mathrm{Set}}$ (mirroring Chapter~\ref{ch:base-fiber}, Theorem~\ref{thm:quantitative-conservation}); because raw model-space pullbacks fail (\texttt{naive\_pullback\_breaks\_subclass\_separation} on \texttt{ClassSetModel}, mirrored on $M \prec V_\kappa$ in Definition~\ref{def:set-projection-fibration}(i)), a rigorous transgression carrier operates via Macro $\Sigma^1_0$-Morleyized classifying topoi $\gamma_I : \mathcal{E}[T_{\mathrm{Class}}] \to \mathcal{E}[T_{\mathrm{Set}}]$ coupled to Micro internal AST class-forcing sheaves $\gamma_\Sigma$, separating predicative $\mathcal{S}$-distributivity leaks (Glazer over $\mathsf{ZC}$~\cite{Glazer23}, Felgner over $\mathsf{ZFC}_{\neg\mathrm{Reg}}$~\cite{Felgner76}) from excision/reflection leaks ($H^0_c(U_{\mathrm{leak}}; \mathbb{Z})$ for $\mathsf{ETR}_{\mathrm{Ord}}$ and $\mathsf{KM}$), while vanishing under Replacement + Foundation over $\mathsf{ZFC}$.
\end{itemize}

The scare quotes around ``$\mathbf{Sets}$'' may now come off---not because a single monolithic foundation has prevailed, but because its hidden joints lie fully in view. What once appeared to be an unanalyzable meta-theoretic floor is simply the double-clamped locus inside Joyal--Moerdijk ambients $(\mathcal{C}, \mathcal{S})$---bridging to univalent $(\infty, 1)$-topos universes $(\mathcal{X}, \mathcal{U})$ and the Volume~II horizon of elementary class-anima ambients $(\mathcal{C}_\infty, \mathcal{S}_\infty)$---where existence-compression from above meets hereditary symmetry-truncation from below, and where every turn of the class kaleidoscope overhead obeys the exact same laws of necessity that govern the mathematics within.
}{\chapter{\texorpdfstring{Residuums and ``$\mathbf{Sets}$''}{Residuums and "Sets"}}
\chaptermark{Residuums and ``Sets''}
\label{ch:sets-residuum}
\chapterepigraph{``The inequalities in the world are like the combinations in a kaleidoscope; at every turn a fresh picture strikes the eye; and yet, in reality, you see only the same bits of glass as you saw before.''}{Arthur Schopenhauer, \emph{Counsels and Maxims}, \S 47}
\chapterepigraph{``Sometimes I have had the task to write reviews of articles, where set theoretic notions and theorems are used without any explanation of what kind of set theory is meant. It is a disagreeable job to write reviews in such cases. One does not know what the author really means.''}{Thoralf Skolem, \emph{A Critical Remark on Foundational Research} (1955)~\cite{Skolem55}}

Throughout Chapters~\ref{ch:foundations}--\ref{ch:realizability-sites}, Categorical Reverse Mathematics ($\mathrm{CRM}$) has calibrated mathematical necessity \emph{inside} variable universes of small maps~\cite{Lawvere69,Johnstone02}. Across all Five Foundational Coordinate Axes established in Chapter~\ref{ch:foundations} (Definitions~\ref{def:comparison-triple}--\ref{def:residuum} and Table~\ref{tab:five-axis-master})---Choice/Splitting (\textbf{Axis~I}), Constructive Logic (\textbf{Axis~II}), Arithmetic Infinity (\textbf{Axis~III}), Characteristic/$p$-Torsion (\textbf{Axis~IV}), and Directional Invertibility (\textbf{Axis~V})---doctrinal comparison triples $\mathfrak{T} = (\mathcal{D}_0 \supseteq \mathcal{D}_1,\; \kappa : \Phi \Rightarrow \Psi) \in \mathrm{Dom}(\mathrm{Res})$ were measured through concrete internal invariants: non-abelian torsor cohomology $H^1(\mathcal{E}/Y; G)$, derived limit phantoms $\varprojlim^1$, Heyting boundary lattices, $\delta$-ring lift torsors, and Postnikov obstructions. Throughout those calibrations, the background doctrine of ``$\mathbf{Sets}$'' remained quietly offstage---even in the face of Thoralf Skolem's warning~\cite{Skolem55} against unspecified background set theories---shielded by the folkloric consensus that Zermelo--Fraenkel set theory ($\mathsf{ZFC}$) and von Neumann--Bernays--G\"odel class theory ($\mathsf{NBG}$) are functionally interchangeable ways of naming the same mathematical bedrock.

Turning the residuum operator $\mathrm{Res}(\kappa)$ onto ``$\mathbf{Sets}$'' itself---via Andr\'e Joyal and Ieke Moerdijk's \textbf{Algebraic Set Theory}\index{Algebraic Set Theory} of class categories equipped with small-map ideals $(\mathcal{C}, \mathcal{S})$~\cite{JoyalMoerdijk95}---shatters that consensus along two orthogonal directions. First, holding an underlying universe of small sets completely rigid still leaves an entire vertical kaleidoscope of $2$-categorical class residua overhead. Second, and far more subversively, weakening the horizontal set axioms below full $\mathsf{ZFC}$ breaks class-to-set conservativity outright---exposing the structural plot twist that Lawvere's Elementary Theory of the Category of Sets ($\mathsf{ETCS}$)\index{ETCS@$\mathsf{ETCS}$ (Elementary Theory of the Category of Sets)}~\cite{Lawvere64}, the Boolean elementary-topos baseline inhabited across Chapters~\ref{ch:foundations}--\ref{ch:realizability-sites}, is bi-interpretable with Bounded Zermelo / Mac~Lane set theory ($\mathsf{Mac}$)~\cite{Cole73,Mitchell72,Osius74} rather than $\mathsf{ZFC}$ and has quietly operated without Replacement (and even without unbounded $\mathcal{L}_\in$-Separation) from page one.

\section{\texorpdfstring{The Conservativity Illusion and the Fixed-Set Kaleidoscope}{The Conservativity Illusion and the Fixed-Set Kaleidoscope}}\label{sec:conservativity-kaleidoscope}
\sectionepigraph{``Our present assumption of `one universe' is an adequate stopgap, not a forecast of the future.''}{Saunders Mac~Lane, \emph{Categories for the Working Mathematician}~\cite[Ch.~I, \S 6]{MacLane98}}

Despite Saunders Mac~Lane's caution against treating a single universe as final~\cite[Ch.~I, \S 6]{MacLane98}, foundational folklore routinely dismisses the gap between set theory and class theory on the strength of Ulrich Felgner's 1971 class-forcing theorem~\cite{Felgner71}: $\mathsf{NBG}$\index{NBG@$\mathsf{NBG}$ (von Neumann--Bernays--G\"odel class theory)}---formulated as Bernays--G\"odel class theory $\mathsf{GBC}$ augmented with Global Choice ($\mathsf{GC}$)---proves the exact same first-order $\mathcal{L}_\in$-sentences about \emph{sets} as $\mathsf{ZFC}$. Proof-theoretically, $\mathsf{NBG}$ stands to $\mathsf{ZFC}$ just as arithmetical comprehension $\mathrm{ACA}_0$ stands to Peano Arithmetic $\mathsf{PA}$ in classical Reverse Mathematics~\cite{Simpson09}.

The catch is that \textbf{first-order conservativity is strictly weaker than bi-interpretability} or categorical Morita equivalence~\cite{Caramello18}. To track class-to-set projections across both $\mathsf{ZFC}$ and weaker Replacement-free bases without circularity, we distinguish Replacement-free class algebra from full Bernays--G\"odel class theory at the outset.

\begin{definition}[Class-Set Model Spaces and the Baseline Adjunction $\mathrm{Def} \dashv \pi_{\mathrm{Set}}$]\label{def:set-projection-fibration}%
\lean{HomotopyCRM.SetsResiduum.defSetProjAdjunction, HomotopyCRM.SetsResiduum.finToPowElemHom, HomotopyCRM.SetsResiduum.finToPowClassHom, HomotopyCRM.SetsResiduum.naive_pullback_breaks_subclass_separation}%
\leanwitness%
\uses{def:comparison-triple,def:residuum}%
Let $\mathsf{GB}_0$\index{GB0@$\mathsf{GB}_0$ (Replacement-free class theory)} denote \emph{Replacement-free Bernays--G\"odel class theory}: the two-sorted $(\mathcal{L}_\in, \mathcal{C})$-theory over Zermelo set theory $\mathsf{Z}$ (Empty Set, Pairing, Union, Power Set, Separation, Infinity, Foundation, and Extensionality) whose class sort $\mathcal{C}$ is closed under G\"odel's eight finite predicative class operations \textbf{G1--G8} together with Subclass Separation ($C \subseteq A \in \mathcal{S} \implies C \in \mathcal{S}$, inducing full $\mathcal{L}_\in$-Separation on sets), \emph{omitting} Bernays--G\"odel Axiom~D (Class Replacement). Write:
\[
  \begin{aligned}
    \mathsf{GBC}_0 &:= \mathsf{GB}_0 + \mathrm{AC}, &\qquad
    \mathsf{GB} &:= \mathsf{GB}_0 + \mathrm{Rep}, \\
    \mathsf{GBC} &:= \mathsf{GB} + \mathrm{AC}, &\qquad
    \mathsf{NBG} &:= \mathsf{GBC} + \mathsf{GC}.
  \end{aligned}
\]
Following Definitions~\ref{def:comparison-triple}--\ref{def:residuum} of Chapter~\ref{ch:foundations}, let $\mathbf{Mod}(\mathsf{GB}_0)$ denote the category whose objects are two-sorted models $(M, \mathcal{C}) \models \mathsf{GB}_0$ and whose $1$-morphisms $(f, F) : (M, \mathcal{C}) \to (N, \mathcal{D})$ consist of an $\mathcal{L}_\in$-elementary embedding $f : M \hookrightarrow N$ paired with a $\Sigma^1_0$-elementary class map $F : \mathcal{C} \to \mathcal{D}$ preserving singleton parameters ($\{p\} \mapsto \{f(p)\}$), G\"odel operations \textbf{G1--G8}, and membership ($m \in C \iff f(m) \in F(C)$, restricting over $f = \mathrm{id}_M$ to fixed-set class inclusions $\mathcal{C} \subseteq \mathcal{C}'$), and let $\mathbf{Mod}(\mathsf{Z})$ be the category of models of $\mathsf{Z}$ and $\mathcal{L}_\in$-elementary embeddings.
\begin{enumerate}[label=\textup{(\roman*)},leftmargin=1.8em]
  \item The forgetful functor $\pi_{\mathrm{Set}} : \mathbf{Mod}(\mathsf{GB}_0) \twoheadrightarrow \mathbf{Mod}(\mathsf{Z})$ sending $(M, \mathcal{C}) \mapsto M$ is surjective on objects, with vertical fiber $\pi_{\mathrm{Set}}^{-1}(M)$ over $\mathrm{id}_M$ consisting of the poset of all $\mathsf{GB}_0$-realizations $(M, \mathcal{C})$ sharing the exact same underlying set universe $M$. (Note that naive class pullback $f^*\mathcal{D} := \{f^{-1}[D] \mid D \in \mathcal{D}\}$ along a proper $\mathcal{L}_\in$-elementary embedding $f : M \hookrightarrow V_\kappa$ does \emph{not} make $\pi_{\mathrm{Set}}$ a model-space Grothendieck fibration, as any small set $r \in \mathcal{P}(\omega)^{V_\kappa} \setminus M$ over countable $M \prec V_\kappa$ satisfies $f^{-1}[r] = r \notin M$, violating Subclass Separation while failing to commute with \textbf{G4} domain projection on $\{(n, r) \mid n \in r\}$; syntactically, however, deductive theory pairs $(T_{\mathrm{Set}}, T_{\mathrm{Class}}) \mapsto T_{\mathrm{Set}}$ do form a base-change opfibration via $T_{\mathrm{Class}} \mapsto T_{\mathrm{Class}} + \sigma$.)
  \item Sending each set model $M \models \mathsf{Z}$ to its first-order parameter-definable class expansion $\mathrm{Def}(M)$ defines a full coreflective baseline embedding $\mathrm{Def} : \mathbf{Mod}(\mathsf{Z}) \hookrightarrow \mathbf{Mod}(\mathsf{GB}_0)$ ($\pi_{\mathrm{Set}} \circ \mathrm{Def} = \mathrm{Id}_{\mathbf{Mod}(\mathsf{Z})}$) left adjoint to $\pi_{\mathrm{Set}}$:
  \[
    \mathrm{Def} \;\dashv\; \pi_{\mathrm{Set}} : \mathbf{Mod}(\mathsf{GB}_0) \twoheadrightarrow \mathbf{Mod}(\mathsf{Z}).
  \]
  Indeed, for any $(N, \mathcal{D}) \models \mathsf{GB}_0$ and $\mathcal{L}_\in$-elementary $f : M \hookrightarrow N$, $\mathcal{L}_\in$-elementarity and G\"odel's Class Existence Theorem send each definable class $\varphi^M(-, p) \in \mathrm{Def}(M)$ to a unique well-defined class $\varphi^N(-, f(p)) \in \mathrm{Def}(N) \subseteq \mathcal{D}$. Over $\mathbf{Mod}(\mathsf{ZFC}) \hookrightarrow \mathbf{Mod}(\mathsf{ZC})$, $\mathrm{Def} \dashv \pi_{\mathrm{Set}}$ restricts to $\mathbf{Mod}(\mathsf{GBC}) \twoheadrightarrow \mathbf{Mod}(\mathsf{ZFC})$.
\end{enumerate}
\end{definition}

Suppose we clamp an underlying countable model $M \models \mathsf{ZFC}$ completely rigid, freezing every set-level $\mathsf{ZFC}$ residuum to zero so that not a single element, ordinal, real number, or small morphism can vary. Even above this frozen set floor, the vertical fiber $\pi_{\mathrm{Set}}^{-1}(M)$ opens into an intricate second-order poset governing which proper classes $\mathcal{C}$ exist above $M$.

\begin{theorem}[{External Benchmark: Fixed-Set Class Kaleidoscope; Williams~2018~\cite[Ch.~3, \S 3.2]{Williams18}}]\label{thm:williams-kaleidoscope}%
\uses{def:set-projection-fibration}%
Fix any countable model $M \models \mathsf{ZFC}$. The partial order of $\mathsf{GBC}$-realizations $(M, \mathcal{C})$ having the exact same underlying set universe $M$, ordered by class inclusion $\mathcal{C} \subseteq \mathcal{C}'$, embeds every countable partial order.
\end{theorem}

Standing inside a fixed model $M \models \mathsf{ZFC}$ while varying $\mathcal{C}$ realizes Schopenhauer's kaleidoscope literally: the set-level bits of glass never move, yet turning through admissible class skies $(M, \mathcal{C})$ reconfigures which large categorical constructions succeed or obstruct.

Recall how Choice appears categorically at the small scale: complementing Chapter~\ref{ch:choice-atlas}'s epimorphism-splitting Choice-Geometry Atlas along Axis~I, Michael Makkai~\cite{Makkai96} observed that inside the $2$-category $\mathbf{Cat}(\mathcal{S})$ of small categories over $\mathbf{Set}$, $\mathrm{AC}$ is equivalent to requiring every \emph{weak categorical equivalence} (every functor that is fully faithful and essentially surjective on objects) to admit a strong quasi-inverse functor without passing to anafunctors. Crucially, the very first step up our vertical class fiber ($\mathsf{GBC} \subsetneq \mathsf{NBG}$) is this exact equivalence-splitting residuum shifted one universe upward from internal small categories $\mathbf{Cat}(\mathcal{S})$ to large categories $\mathbf{Cat}(\mathcal{C})$.

\begin{proposition}[Large Equivalence-Splitting Calibrates Global Choice]\label{prop:large-equiv-splitting}%
\lean{HomotopyCRM.SetsResiduum.card_split_iff_global_choice, HomotopyCRM.SetsResiduum.type_equiv_splitting_global_choice}%
\leanwitness%
\uses{def:set-projection-fibration}%
Work over $\mathsf{GBC}$\index{Global Choice ($\mathsf{GC}$)} in a class-set model $(M, \mathcal{C})$, and let $I : \mathbf{Card} \hookrightarrow \mathbf{Set}$ denote the full inclusion in $\mathbf{Cat}(\mathcal{C})$ of the skeletal subcategory of small von Neumann cardinals into small sets. Then:
\begin{enumerate}[label=\textup{(\roman*)},leftmargin=1.8em]
  \item $I : \mathbf{Card} \hookrightarrow \mathbf{Set}$ is a weak categorical equivalence (fully faithful and essentially surjective on objects).
  \item The existence of a quasi-inverse class functor $F : \mathbf{Set} \to \mathbf{Card}$ equipped with a class natural isomorphism $\eta : \mathrm{Id}_{\mathbf{Set}} \xrightarrow{\sim} I \circ F$ is strictly equivalent over $\mathsf{GBC}$ to Global Choice ($\mathsf{NBG} := \mathsf{GBC} + \mathsf{GC}$).
  \item Whenever $M \models \mathsf{ZFC} + \forall p\,(V \neq \mathrm{HOD}(\{p\}))$, as obtained by homogeneity of tail Easton class-forcing extensions~\cite{Easton70}, the minimal baseline fiber $(M, \mathrm{Def}(M)) \models \mathsf{GBC}$ violates $\mathsf{GC}$ and therefore leaves $I : \mathbf{Card} \hookrightarrow \mathbf{Set}$ unsplit. (Bare $V \neq \mathrm{HOD}$ does not suffice: e.g.\ adding a single Cohen real $c$ over $L$ gives $L[c] \models V \neq \mathrm{HOD}$, yet the parameter-definable well-ordering $<_{L[c]} \in \mathrm{Def}(L[c])$ still splits $I$.)
\end{enumerate}
\end{proposition}

\begin{proof}
Under set-level $\mathrm{AC}$ and Replacement in $\mathsf{GBC}$, every small set $X \in \mathbf{Set}$ admits a bijection with a unique von Neumann initial ordinal $|X| \in \mathbf{Card}$. Because $\mathbf{Card}$ is a full subcategory of $\mathbf{Set}$, $I$ is fully faithful and essentially surjective, proving~(i).

For~(ii), suppose first that $I$ splits via a class functor $F : \mathbf{Set} \to \mathbf{Card}$ and a class natural isomorphism $\eta : \mathrm{Id}_{\mathbf{Set}} \xrightarrow{\sim} I \circ F$. Since cardinals are skeletal, $F(X) = |X|$ for every set $X$. Whenever $X \neq \emptyset$, its cardinal $|X|$ is a nonzero von Neumann ordinal and hence satisfies $\emptyset \in |X|$. Because $\eta_X : X \xrightarrow{\sim} |X|$ is a bijection for each set $X$, the predicative class assignment
\[
  G(X) \;:=\; \eta_X^{-1}(\emptyset) \;\in\; X
\]
selects a uniform element of every nonempty small set $X \in \mathcal{S}$. Over $\mathsf{GBC}$, Foundation decomposes $V = \bigcup_{\alpha \in \mathrm{Ord}} V_\alpha$, so setting $\tilde{G}(C) := G(C \cap V_{\alpha_0})$ for $\alpha_0 := \min\{\alpha \in \mathrm{Ord} \mid C \cap V_\alpha \neq \emptyset\}$ immediately extends $G$ to a class-form Global Choice selector on all nonempty classes $C \in \mathcal{C}$, establishing $\mathsf{GC}$. Conversely, assume $\mathsf{GC}$. Applying global choice to the class family of nonempty bijection sets $\{\mathrm{Bij}(X, |X|)\}_{X \in \mathbf{Set}}$ simultaneously selects one canonical bijection $\eta_X : X \xrightarrow{\sim} |X|$ for every small set $X$. Defining $F(X) := |X|$ on objects and $F(f) := \eta_Y \circ f \circ \eta_X^{-1} : |X| \to |Y|$ on morphisms $f : X \to Y$ assembles $F : \mathbf{Set} \to \mathbf{Card}$ into a genuine class functor with natural isomorphism $\eta : \mathrm{Id}_{\mathbf{Set}} \xrightarrow{\sim} I \circ F$.

Finally, for~(iii), any Global Choice selector $G \in \mathrm{Def}(M)$ definable over $M$ with a set parameter $p \in M$ induces by transfinite recursion over each $V_\alpha$ a well-ordering of $V$ definable from $p$ and ordinals, forcing $M \models V = \mathrm{HOD}(\{p\})$. Notice that bare $V \neq \mathrm{HOD}$ allows parameter-definable well-orderings (for instance, $<_{L[c]} \in \mathrm{Def}(L[c])$ for a single Cohen real $c$ over $L$, where $L[c] \models V \neq \mathrm{HOD}$ yet $L[c] = \mathrm{HOD}(\{c\})$ still splits $I : \mathbf{Card} \hookrightarrow \mathbf{Set}$). By homogeneity of tail Easton class-forcing extensions~\cite{Easton70} (where every set parameter $p$ lies in an initial segment and is fixed by automorphisms of the tail forcing), models satisfying $M \models \mathsf{ZFC} + \forall p\,(V \neq \mathrm{HOD}(\{p\}))$ exist; in any such $M$, no parameter-definable global choice selector lies in $\mathrm{Def}(M)$, so $I : \mathbf{Card} \hookrightarrow \mathbf{Set}$ cannot split in $(M, \mathrm{Def}(M))$.
\end{proof}

Climbing above $\mathsf{NBG}$ along the vertical fiber calibrates the implicit class commitments invoked across modern $2$-category theory and homotopy theory:
\begin{itemize}[leftmargin=1.5em]
  \item \textbf{Class Forcing and Transfinite Ad\'amek Chains ($\mathsf{NBG} \subsetneq \mathsf{NBG} + \mathsf{ETR}_\omega \subsetneq \mathsf{NBG} + \mathsf{ETR}_{\mathrm{Ord}} \subsetneq \mathsf{NBG} + \mathsf{ETR}$):} Gitman, Hamkins, Holy, Schlicht, and Williams~\cite[Thm.~1.2]{Gitman20} (and Williams~\cite{Williams18}; note that Gitman et al.\ write ``$\mathsf{GBC}$'' for $\mathsf{NBG} = \mathsf{GBC} + \mathsf{GC}$) proved that the Class Forcing Theorem (definability of $\Vdash_{\mathbb{P}}$ and the Truth Lemma for proper-class forcing posets) is equivalent over $\mathsf{NBG} = \mathsf{GBC} + \mathsf{GC}$ to Elementary Transfinite Recursion of length $\mathrm{Ord}$ ($\mathsf{ETR}_{\mathrm{Ord}}$), failing in $\mathsf{NBG}$---and likewise encompassing $\mathrm{Ord}$-length Ad\'amek chains~\cite{Goldthorpe23} when encoded via first-order class step relations. By contrast, \emph{full $\mathsf{ETR}$} along arbitrary class well-orderings is strictly stronger, proving $\mathrm{Con}(\mathsf{NBG} + \mathsf{ETR}_{\mathrm{Ord}})$~\cite{Gitman20}, whereas constructing full Tarskian truth predicates relative to arbitrary class parameters holds already at $\mathsf{ETR}_\omega$, which already proves $\mathrm{Con}(\mathsf{ZFC})$ and triggers a syntactic $\mathcal{L}_\in$-reflection leak over $\mathsf{ZFC}$ while remaining $\mathcal{L}_\in$-conservative over its first-order trace $T^* := \mathrm{Th}_{\mathcal{L}_\in}(\mathsf{NBG} + \mathsf{ETR}_{\mathrm{Ord}})$ (strictly below full $\mathsf{ETR}$ and impredicative $\mathsf{KM}$).
  \item \textbf{Impredicative Class Comprehension vs.\ Grothendieck Universes ($\mathsf{NBG} + \mathsf{ETR}_{\mathrm{Ord}} \subsetneq \mathsf{NBG} + \mathsf{ETR} \subsetneq \mathsf{NBG} + \Pi^1_1\text{-}\mathsf{CA} \subsetneq \mathsf{KM}$):}\index{Kelley--Morse class theory ($\mathsf{KM}$)} By Freyd's theorem~\cite[Ch.~V, \S 2, Prop.~3]{MacLane98}, no small category admitting products indexed by its own morphism size can avoid collapsing to a preorder; hence $\mathbf{Set}$ cannot admit arbitrary proper-class-indexed limits, colimits, or unrestricted pointwise large Kan extensions even in Kelley--Morse ($\mathsf{KM}$), which instead require passing to a larger Grothendieck universe $V_\lambda \ni V_\kappa$ ($\kappa < \lambda$ inaccessible) as the ambient ``$\mathbf{Set}$''. Rather, $\mathsf{KM}$ calibrates full impredicative ($\Pi^1_\infty$) class comprehension over $V$---supporting complete class-lattice meets/joins and impredicative second-order categorical invariants over $V$---driving the syntactic reflection leak far past $T^*$ to $\mathrm{Th}_{\mathcal{L}_\in}(\mathsf{KM})$ (proving $\mathrm{Con}(\mathsf{NBG} + \mathsf{ETR})$ and hence $\mathrm{Con}(\mathsf{NBG} + \mathsf{ETR}_{\mathrm{Ord}})$ as modeled by two-sorted $(V_\kappa, V_{\kappa+1})$ for inaccessible $\kappa$), mirroring the classical climb from $\mathrm{ACA}_0$ through $\mathrm{ATR}_0$ and $\Pi^1_1\text{-}\mathrm{CA}_0$ to full second-order arithmetic $\mathsf{Z}_2$~\cite{Simpson09}.
\end{itemize}

\section{\texorpdfstring{Cracks in the Floor: Why ``$\mathbf{Sets}$'' Cannot Hold the Base}{Cracks in the Floor: Why "Sets" Cannot Hold the Base}}\label{sec:cracks-in-the-floor}

At this juncture, a pragmatic category theorist might still hope to quarantine the entire phenomenon upstairs: let specialists pay vertical class tolls in $\mathsf{NBG}$, $\mathsf{ETR}_{\mathrm{Ord}}$, or $\mathsf{KM}$ whenever they split large equivalences, iterate $\mathrm{Ord}$-chains, or evaluate impredicative class comprehensions (moving to higher Grothendieck universes whenever large limits or Kan extensions are formed), while ordinary small-map mathematics proceeds undisturbed inside a canonical ground floor of ``$\mathbf{Sets}$.''

Yet our \textbf{Base--Fiber Conservation Principle} (Chapter~\ref{ch:foundations}, Principle~\ref{prin:conservation}, and Chapter~\ref{ch:base-fiber}, Principle~\ref{thm:conservation-law}) prompts a pointed structural question about Felgner's 1971 conservativity theorem~\cite{Felgner71}: \emph{what prevented proper-class Global Choice from leaking down into set-level theorems in the first place?} Inspecting Felgner's forcing proof reveals two load-bearing set axioms acting as an insulating gasket underneath $\pi_{\mathrm{Set}}$: \textbf{Replacement}, which keeps the class forcing of set-sized partial well-orderings of $V$ $\kappa$-closed for every cardinal $\kappa$ so that no new small sets are added, and \textbf{Foundation}, which anchors cumulative rank induction across $V$.

Syntactic $\mathcal{L}_\in$-conservativity across $\pi_{\mathrm{Set}}$ is consequently not an intrinsic harmlessness of proper classes; it is a \emph{local flatness property} over $\mathbf{Mod}(\mathsf{ZFC})$ that shatters on wider small-set bases where Replacement fails:
\[
\begin{tikzcd}[row sep=2.2em, column sep=2.8em, font=\footnotesize]
\mathbf{Mod}(\mathsf{NBG})
  \arrow[r, hook, "\tilde{\iota}"]
  \arrow[d, "\substack{\pi_{\mathrm{Set}}|_{\mathsf{ZFC}}\\ \text{($\mathcal{L}_\in$-conserv.: Felgner 1971)}}"']
&
\mathbf{Mod}(\mathsf{GBC}_0 + \mathsf{GC})
  \arrow[d, "\substack{\pi_{\mathrm{Set}}|_{\mathsf{ZC}}\\ \text{(non-conserv.: Glazer 2023)}}"]
\\
\mathbf{Mod}(\mathsf{ZFC})
  \arrow[r, hook, "\iota"]
&
\mathbf{Mod}(\mathsf{ZC})
\end{tikzcd}
\]
Here horizontal inclusions run contravariantly to syntactic theory strength ($\mathsf{ZC} \subsetneq \mathsf{ZFC}$ induces $\mathbf{Mod}(\mathsf{ZFC}) \hookrightarrow \mathbf{Mod}(\mathsf{ZC})$). Moreover, unlike $\pi_{\mathrm{Set}} : \mathbf{Mod}(\mathsf{GBC}) \twoheadrightarrow \mathbf{Mod}(\mathsf{ZFC})$ (surjective via $\mathrm{Def}(M)$), neither vertical $\mathsf{GC}$-projection above is surjective on all models: $\pi_{\mathrm{Set}}|_{\mathsf{ZC}}$ omits $\mathsf{ZC}$-models of $\theta_{\mathrm{Glz}}$ by Theorem~\ref{thm:glazer-nonconservativity}, and $\pi_{\mathrm{Set}}|_{\mathsf{ZFC}}$ fails surjectivity on uncountable rather classless models of Easton's theory (cf.\ Shelah~\cite{Shelah78}). As soon as $\mathrm{CRM}$ relaxes the set base below $\mathsf{ZFC}$, class principles stop tracking the set base along three independent fault lines.

\begin{remark}[Set-Theoretic Notation Alignment: $\mathsf{Mac}$ vs.\ $\mathsf{ZC}$ and $\mathsf{ZFC}^-$ vs.\ $\mathsf{ZFC}_{\neg\mathrm{Reg}}$]\label{rem:zfc-minus-notation}%
\uses{def:set-projection-fibration}%
Two standard terminology collisions must be fixed cleanly before dissecting weaker set bases:
\begin{enumerate}[label=\textup{(\alph*)},leftmargin=1.8em]
  \item \textbf{Bounded Zermelo ($\mathsf{Mac}$) versus Full Zermelo ($\mathsf{ZC}$):}\index{Mac Lane set theory ($\mathsf{Mac}$)} Following Adrian Mathias~\cite{Mathias01}, $\mathsf{ZC}$ denotes \emph{full Zermelo set theory with Choice}, including the \emph{unbounded} first-order $\mathcal{L}_\in$-Separation schema, whereas \emph{Bounded Zermelo / Mac~Lane set theory} ($\mathsf{Mac}$, also denoted $\mathsf{ZBQC}$) restricts Separation to bounded ($\Delta_0$) formulas. Crucially, the Cole--Mitchell--Osius equivalence~\cite{Cole73,Mitchell72,Osius74} (\cite[Ch.~VI, \S 10]{MacLaneMoerdijk92}) proves that Lawvere's $\mathsf{ETCS}$~\cite{Lawvere64} is bi-interpretable with $\mathsf{Mac}$ ($\Delta_0$-Separation), \emph{not} full $\mathsf{ZC}$, and Mathias~\cite[pp.~107--234]{Mathias01} proved that $\mathsf{ZC}$ has strictly greater consistency strength than $\mathsf{Mac} \simeq \mathsf{ETCS}$.
  \item \textbf{Disambiguating $\mathsf{ZFC}^-$ versus $\mathsf{ZFC}_{\neg\mathrm{Reg}}$:} Set-theoretic literature overloads $\mathsf{ZFC}^-$ for both Power-Set-free $\mathsf{ZFC}$ and Foundation-free $\mathsf{ZFC}$~\cite{Felgner76}. Throughout this chapter and Chapter~\ref{ch:master-atlas}, we reserve $\mathsf{ZFC}^-$ strictly for \textbf{Power-Set-free} $\mathsf{ZFC}$ (formulated with Collection in place of functional Replacement~\cite{GitmanHamkinsJohnstone16}), and write $\mathsf{ZFC}_{\neg\mathrm{Reg}}$ and $\mathsf{NBG}_{\neg\mathrm{Reg}}$ when \textbf{Foundation (Regularity)} is dropped.
\end{enumerate}
\end{remark}

\begin{theorem}[{External Benchmark: Syntactic $\mathcal{L}_\in$-Non-Conservativity of Global Choice Over $\mathsf{ZC}$; Glazer~2023~\cite[Main Theorem]{Glazer23}}]\label{thm:glazer-nonconservativity}%
\uses{def:set-projection-fibration,rem:zfc-minus-notation}%
Global Choice is \emph{not} syntactically $\mathcal{L}_\in$-conservative over Zermelo set theory with Choice ($\mathsf{ZC}$, with or without Foundation, over $\mathsf{GBC}_0 + \mathsf{GC}$, i.e.\ $\mathsf{GB}_0 + \mathsf{AC} + \mathsf{GC}$): there exists an explicit first-order $\mathcal{L}_\in$-sentence $\theta_{\mathrm{Glz}} := (\mathsf{CH} \wedge \psi_{\mathrm{fnd}})$ consistent with $\mathsf{ZC}$ (realized in a symmetric extension model satisfying the Generalized Continuum Hypothesis $\mathsf{GCH}$) such that $\mathsf{GBC}_0 + \mathsf{GC} \vdash \neg\theta_{\mathrm{Glz}}$. Explicitly, $\psi_{\mathrm{fnd}}$ asserts that parameterless $\mathcal{L}_\in$-definable classes encode:
\begin{enumerate}[label=\textup{(\alph*)},leftmargin=1.8em]
  \item an $\omega$-indexed sequence $\langle X_n \rangle_{n < \omega}$ of nonempty countable sets (admitting no small uniform enumeration sequence in the absence of Replacement);
  \item a class $A$ separated via a binary relation $R$ by subsets of $\bigcup_{n < \omega} X_n$ (so $\{x \in \bigcup_n X_n \mid x \mathrel{R} a_0\} = \{x \in \bigcup_n X_n \mid x \mathrel{R} a_1\} \implies a_0 = a_1$ for $a_0, a_1 \in A$); and
  \item a class surjection $P_2 : A \twoheadrightarrow W_2$ onto the canonical small prewellordering $W_2 \in \mathcal{P}^5(\omega)$ of $\mathcal{P}^2(\omega)$ of second uncountable order type $\omega_2$ (``$\aleph_2$'').
\end{enumerate}
In $\mathsf{GBC}_0 + \mathsf{GC}$, Global Choice simultaneously selects surjections $f_n : \omega \twoheadrightarrow X_n$ for all $n < \omega$, allowing unbounded Subclass Separation (projecting out the unbounded existential quantifier $a \in A$ via \textbf{G4} and intersecting with $\mathcal{P}(\omega \times \omega) \times W_2 \in \mathcal{S}$) to carve out a small-set partial surjection $\mathcal{P}(\omega \times \omega) \rightharpoonup W_2$ sending $\{(m, n) \mid f_n(m) \mathrel{R} a\} \mapsto P_2(a)$ ($a \in A$), refuting $\mathsf{CH}$. (Separately, adjoining unbounded $\mathsf{GB}_0$ Subclass Separation already breaks $\mathcal{L}_\in$-conservativity over $\Delta_0$-bounded $\mathsf{Mac} \simeq \mathsf{ETCS}$~\cite{Mathias01}, whereas Theorem~\ref{thm:glazer-nonconservativity} proves $\mathsf{GC}$ itself breaks $\mathcal{L}_\in$-conservativity over full $\mathsf{ZC}$.)
\end{theorem}

Notice two structural subtleties highlighted by Theorem~\ref{thm:glazer-nonconservativity}. First, syntactic $\mathcal{L}_\in$-non-conservativity is strictly stronger than mere \emph{model non-expandability}: for example, compositional truth $\mathsf{CT}^-[\mathsf{ZFC}]$~\cite{Fujimoto12} is syntactically $\mathcal{L}_\in$-conservative over $\mathsf{ZFC}$ yet fails to expand countable non-recursively-saturated models $M \models \mathsf{ZFC}$ (and likewise $\mathsf{NBG} + \mathsf{ETR}_{\mathrm{Ord}}$ exhibits countable model non-expandability over its own first-order trace $T^*$, while already leaking $\mathrm{Con}(\mathsf{ZFC})$ over $\mathsf{ZFC}$~\cite[Thm.~1.2]{Gitman20}), whereas over $\mathsf{ZC}$, $\mathsf{GBC}_0 + \mathsf{GC}$ refutes the explicit first-order $\mathcal{L}_\in$-sentence $\theta_{\mathrm{Glz}}$ consistent with $\mathsf{ZC}$~\cite{Glazer23}. Second, inside any model $M$ of ordinal height $\mathrm{Ord}^M = \omega + \omega$, uncountable von Neumann initial ordinals do not exist at all (so Proposition~\ref{prop:large-equiv-splitting}(i)'s inclusion $I : \mathbf{Card} \hookrightarrow \mathbf{Set}$ also loses essential surjectivity once Replacement is dropped!), requiring uncountable cardinalities such as $\aleph_n$ to be represented via canonical small prewellorderings inside $\mathcal{P}^{n+3}(\omega)$.

Alongside Glazer's Replacement-free leak, dropping either Foundation or Choice exposes two further horizontal cracks:
\begin{itemize}[leftmargin=1.5em]
  \item \textbf{Dropping Foundation ($\mathsf{ZFC}_{\neg\mathrm{Reg}}$; Felgner~1976~\cite{Felgner76}):} Over Foundation-free set theory $\mathsf{ZFC}_{\neg\mathrm{Reg}} + \mathrm{Collection}$, Ulrich Felgner proved that $\mathsf{NBG}_{\neg\mathrm{Reg}}$ augmented with class Equipollence is non-conservative, forcing a strictly stronger set-level $\mathrm{Ord}$-indexed Dependent Choice schema $\mathsf{DC}_{\mathrm{Ord}}$.
  \item \textbf{Dropping Choice ($\mathsf{GB}_0$ vs.\ Limitation of Size $\mathsf{LoS}$):}\index{Limitation of Size ($\mathsf{LoS}$)} In von Neumann's original 1925 axiomatization of class theory, Replacement and Global Choice were fused into a single mega-axiom, \textbf{Limitation of Size} ($\mathsf{LoS}$): \emph{a class $C \in \mathcal{C}$ fails to belong to the smallness ideal $\mathcal{S}$ if and only if $C$ surjects onto the universe $V$}.
\end{itemize}

\begin{proposition}[Unbundling Limitation of Size Over $\mathsf{GB}_0$]\label{prop:los-smuggles-gc}%
\uses{def:set-projection-fibration}%
Over Replacement-free, Choice-free class theory $\mathsf{GB}_0$, von Neumann's Limitation of Size ($\mathsf{LoS}$) is strictly equivalent to the conjunction of Class Replacement ($\mathrm{Rep}$) and Global Choice ($\mathsf{GC}$):
\[
  \mathsf{GB}_0 \;\vdash\; \mathsf{LoS} \;\Longleftrightarrow\; (\mathrm{Rep} \wedge \mathsf{GC}).
\]
In particular, $\mathsf{GB}_0 + \mathsf{LoS}$ coincides with full $\mathsf{NBG}$ ($\mathrm{Ord} \xrightarrow{\sim} V$), smuggling Global Choice across Choice-free bases and satisfying $\mathsf{GB} + \mathsf{GC} \vdash \mathsf{LoS}$ whenever Replacement holds; conversely, over $\mathsf{GB}_0$ without Replacement, $\mathsf{GB}_0 + \mathsf{GC} \not\vdash \mathsf{LoS}$.
\end{proposition}

\begin{proof}
Work first in $\mathsf{GB}_0 + \mathsf{LoS}$. By predicative class comprehension, the class $\mathrm{Ord}$ of von Neumann ordinals is well-ordered by $\in$, and by the Burali-Forti paradox $\mathrm{Ord} \notin \mathcal{S}$ (otherwise $\mathrm{Ord} \in \mathrm{Ord}$). Applying $\mathsf{LoS}$ yields a class surjection $p : \mathrm{Ord} \twoheadrightarrow V$. Because $\mathrm{Ord}$ is well-ordered, setting $\mu(x) := \min\{\alpha \in \mathrm{Ord} \mid p(\alpha) = x\}$ defines a class injection $\mu : V \hookrightarrow \mathrm{Ord}$, which pulls the canonical well-ordering of $\mathrm{Ord}$ back to a global well-ordering of $V$, establishing $\mathsf{GC}$. To deduce $\mathrm{Rep}$ without invoking Hartogs ordinals, let $A \in \mathcal{S}$ and let $F : A \to V$ be a class function. If $F[A] \notin \mathcal{S}$, then $\mathsf{LoS}$ would furnish a class surjection $q : F[A] \twoheadrightarrow V$ and hence a composite class surjection $h := q \circ F : A \twoheadrightarrow V$. By G1--G8 and Subclass Separation on $A \in \mathcal{S}$, the diagonal subclass $D := \{a \in A \mid a \notin h(a)\} \subseteq A$ is a small set $D \in \mathcal{S} = V$, so $h(a_0) = D$ for some $a_0 \in A$, yielding the Russell--Cantor contradiction $a_0 \in D \iff a_0 \notin D$. Thus $F[A] \in \mathcal{S}$, proving $\mathrm{Rep}$ (whereupon Mostowski collapse of $\mu[V] \subseteq \mathrm{Ord}$ yields $\mathrm{Ord} \xrightarrow{\sim} V$).

Conversely, work in $\mathsf{GB}_0 + \mathrm{Rep} + \mathsf{GC} = \mathsf{NBG}$. If $C \in \mathcal{S}$, the Russell--Cantor diagonal argument just given shows no class surjection $C \twoheadrightarrow V$ can exist. If $C \notin \mathcal{S}$, Foundation and $\mathrm{Rep}$ give $V = \bigcup_{\alpha \in \mathrm{Ord}} V_\alpha$, and restricting $\mathsf{GC}$ to each set-sized power slice $\mathcal{P}(V_{\alpha+1} \setminus V_\alpha) \in \mathcal{S}$ assembles a global bijection $\Phi : V \xrightarrow{\sim} \mathrm{Ord}$. Then $\Phi[C] \subseteq \mathrm{Ord}$ is a proper class, hence unbounded in $\mathrm{Ord}$ (by Subclass Separation on bounded segments) and not order-isomorphic to any $\delta \in \mathrm{Ord}$ (by $\mathrm{Rep}$); sending each $\xi \in \Phi[C]$ to the order type of its set-sized initial segment $\Phi[C] \cap \xi \in \mathcal{S}$ (requiring no $\mathsf{ETR}_{\mathrm{Ord}}$) maps $C \xrightarrow{\sim} \mathrm{Ord} \xrightarrow{\Phi^{-1}} V$ bijectively onto $V$, verifying $\mathsf{LoS}$.

Finally, for $\mathsf{GB}_0 + \mathsf{GC} \not\vdash \mathsf{LoS}$, consider G\"odel's constructible rank $L_{\omega_\omega^L}$ with its definable classes $\mathrm{Def}(L_{\omega_\omega^L})$ (where $\omega_n^L$ denotes the $n$-th uncountable cardinal inside $L$). By G\"odel's condensation lemma applied inside $L$, every constructible subset of $x \in L_{\omega_n^L}$ lies in $L_{\omega_{n+1}^L} \subset L_{\omega_\omega^L}$, verifying Power Set alongside $\mathsf{Z}$, while the global $<_L$-order gives $\mathsf{GC}$ on a model of constructible ordinal height $\omega_\omega^L$. Yet the cofinal sequence $\{\omega_n^L : n \in \omega\}$ is a countable definable proper class admitting no surjection onto $L_{\omega_\omega^L}$, refuting both $\mathrm{Rep}$ and $\mathsf{LoS}$.
\end{proof}

\begin{table}[htbp]
\centering
\caption{Base-Fiber Map of Set-Level vs.\ Class-Level Spectra Along $\pi_{\mathrm{Set}}$}\label{tab:fibration-map}
\vspace{0.3em}
{\scriptsize\setlength{\tabcolsep}{2.5pt}\renewcommand{\arraystretch}{1.12}%
\begin{tabularx}{\textwidth}{@{}P{0.18\textwidth}P{0.12\textwidth}P{0.17\textwidth}YY@{}}
\toprule
\textbf{Class Theory} & \textbf{Set Base $\pi_{\mathrm{Set}}$} & \textbf{$\mathcal{L}_\in$-Conservative?} & \textbf{Separating Model} & \textbf{Categorical Residuum} \\
\midrule
$\mathsf{GBC}$ & $\mathsf{ZFC}$ & Yes ($\mathrm{Def}(M)$) & $M \models \forall p\,(V \neq \mathrm{HOD}(\{p\}))$ & Baseline left adjoint $\mathrm{Def}$ \\
$\mathsf{NBG}$ ($\mathsf{GBC}+\mathsf{GC}$) & $\mathsf{ZFC}$ & Yes (Felgner~\cite{Felgner71}) & $(M,\mathrm{Def}(M))$ vs.\ $M[G]$ & Unsplit $\mathbf{Card} \hookrightarrow \mathbf{Set}$ \\
$\mathsf{CT}^-[\mathsf{ZFC}]$ & $\mathsf{ZFC}$ & Yes (Fujimoto~\cite{Fujimoto12}) & Non-rec.-sat.\ $M \models \mathsf{ZFC}$ & Countable model non-expandability \\
$\mathsf{NBG}+\mathsf{ETR}_{\mathrm{Ord}}$ ($\subsetneq \mathsf{ETR}$) & $\mathsf{ZFC}$ / $T^*$ & \textbf{No} ($\mathrm{Con}(\mathsf{ZFC})$; yes/$T^*$) & Minimal $(M,\mathrm{Def}(M))$ & Class forcing $\Vdash_{\mathbb{P}}$~\cite[Thm.~1.2]{Gitman20} ($\mathrm{Ord}$-rec.; full $\mathsf{ETR} \vdash \mathrm{Con}$) \\
Kelley--Morse ($\mathsf{KM}$) & $\mathsf{ZFC}$ / $\mathrm{Th}_{\mathcal{L}_\in}(\mathsf{KM})$ & \textbf{No} ($\mathrm{Con}(T^*)$; yes/$\mathrm{Th}_{\mathcal{L}_\in}(\mathsf{KM})$) & $(V_\kappa, V_{\kappa+1})$ ($\kappa$ inacc.) & Impredicative $\Pi^1_\infty$ classes (Groth.\ univ.\ for Kan ext., Freyd) \\
$\mathsf{GBC}_0+\mathsf{GC}$ / $\mathsf{NBG}_{\neg\mathrm{Reg}}$ & $\mathsf{ZC}$ / $\mathsf{ZFC}_{\neg\mathrm{Reg}}$ & \textbf{No} (Glazer~\cite{Glazer23}; Felgner~\cite{Felgner76}) & Symm.\ $V_{\omega+\omega}$ ($\mathsf{ZC}$); FM / Rieger--Bernays ($\mathsf{ZFC}_{\neg\mathrm{Reg}}$) & Leaks $\mathcal{P}(\omega)\twoheadrightarrow W_2$ / $\mathsf{DC}_{\mathrm{Ord}}$ \\
\bottomrule
\end{tabularx}\par}
\end{table}

Table~\ref{tab:fibration-map} triggers an immediate plot twist for Volume~I itself.

\begin{proposition}[Countable Colimit Escape in $V_{\omega+\omega}$ Over $\mathsf{ZC}$ and $\mathsf{ETCS}$]\label{prop:etcs-colimit-escape}%
\uses{rem:zfc-minus-notation}%
Zermelo's cumulative rank universe $V_{\omega+\omega}$ models full Zermelo set theory with Choice $\mathsf{ZC}$ (including unbounded $\mathcal{L}_\in$-Separation, and \emph{a fortiori} its category of sets and functions models $\mathsf{Mac} \simeq \mathsf{ETCS}$), in which the countable tower $n \mapsto \mathcal{P}^n(\mathbb{N})$ is first-order definable over $\mathbb{N}$ yet its colimit $\bigcup_{n \in \mathbb{N}} \mathcal{P}^n(\mathbb{N}) \subsetneq V_{\omega+\omega} = \bigcup_{n \in \mathbb{N}} \mathcal{P}^n(V_\omega)$ fails to exist as a small object.
\end{proposition}

\begin{proof}
Because $\omega + \omega$ is a limit ordinal strictly above $\omega$, $V_{\omega+\omega}$ contains $\omega = \mathbb{N}$ and is closed under pairing, union, and power sets ($A \in V_{\omega+k} \implies \mathcal{P}(A) \in V_{\omega+k+1} \subset V_{\omega+\omega}$), inheriting $\mathrm{AC}$ from the metatheory. Moreover, for any $A \in V_{\omega+k}$ and \emph{any} subset $B \subseteq A$ carved out by an arbitrary unbounded $\mathcal{L}_\in$-formula over $(V_{\omega+\omega}, \in)$, we already have $B \in \mathcal{P}(A) \in V_{\omega+k+1} \subset V_{\omega+\omega}$, so $V_{\omega+\omega} \models \mathsf{ZC}$ (and hence models $\mathsf{Mac} \simeq \mathsf{ETCS}$). Inside $V_{\omega+\omega}$, each finite iterate $\mathcal{P}^n(\mathbb{N})$ has rank $\omega + n < \omega + \omega$ and $y = \mathcal{P}^n(\mathbb{N})$ is $\mathcal{L}_\in$-definable. However, every element $x \in V_{\omega+\omega}$ lies in some $V_{\omega+k}$ ($k < \omega$) of rank strictly below $\omega + \omega$, whereas $\bigcup_{n \in \mathbb{N}} \mathcal{P}^n(\mathbb{N})$ contains sets of rank $\omega + n$ for all $n \in \mathbb{N}$ and therefore escapes through the roof of the topos (strictly inside $V_{\omega+\omega} = \bigcup_{n \in \mathbb{N}} \mathcal{P}^n(V_\omega)$, since e.g.\ $V_\omega \in V_{\omega+1} \setminus \bigcup_{n \in \mathbb{N}} \mathcal{P}^n(\mathbb{N})$, while both $\{\mathcal{P}^n(\mathbb{N})\}_{n \in \mathbb{N}}$ and $\{V_{\omega+n}\}_{n \in \mathbb{N}}$ are rank-cofinal countable families escaping $\mathcal{S}$).
\end{proof}

One might suspect that losing rank $\omega + \omega$ and $\aleph_\omega$ affects only esoteric transfinite cardinal arithmetic. Harvey Friedman demolished that comfort in 1971~\cite{Friedman71} by proving that \textbf{Borel Determinacy}\index{Borel Determinacy}---whether every infinite two-player game on a Borel subset of Cantor space $2^{\mathbb{N}}$, an arena whose moves, strategies, and Borel payoff sets reside down in $V_{\omega+1}$ and $V_{\omega+2}$ inside $V_{\omega+\omega}$, admits a winning strategy---\textbf{cannot be proved in Zermelo set theory $\mathsf{ZC}$ (and hence cannot be proved in $\mathsf{ETCS}$)}. Synthesizing winning strategies for general Borel games ($\boldsymbol{\Sigma}^0_\alpha$ across $\alpha < \omega_1$) requires up to $\omega_1$ iterated power sets supplied by Replacement (already exceeding the bounded finite iterates $\{\mathcal{P}^n(\mathbb{N})\}_{n < \omega}$ available in $\mathsf{ZC}$).

Look at the resulting predicament on our supposedly settled ground floor: adjoining unbounded $\mathsf{GB}_0$ Subclass Separation already breaks $\mathcal{L}_\in$-conservativity over $\Delta_0$-bounded $\mathsf{ETCS} \simeq \mathsf{Mac}$~\cite{Mathias01}, and working over full $\mathsf{ZC}$ while casually invoking Global Choice overhead (``because $\mathsf{NBG}$ is conservative over $\mathsf{ZFC}$'') mutates small-set cardinal arithmetic underneath via Theorem~\ref{thm:glazer-nonconservativity}; yet remaining strictly inside $\mathsf{ETCS}$ or $\mathsf{ZC}$ leaves countable definable colimits escaping the universe and concrete Borel games in $V_{\omega+2}$ undecided. ``$\mathbf{Sets}$'' was never an indivisible atom---compelling us to place ``$\mathbf{Sets}$'' directly on the $(\mathcal{C}, \mathcal{S})$ dissecting table.

\section{\texorpdfstring{Dissecting ``$\mathbf{Sets}$'' in $(\mathcal{C}, \mathcal{S})$: Plumbing vs.\ Over-Compressed Existence}{Dissecting "Sets" in (C, S): Plumbing vs. Over-Compressed Existence}}\label{sec:ast-plumbing-compression}
\sectionepigraph{``The formal system may need to be changed. It would be one-sided to only `work in' a system, accepting it as a divine condemnation that we have to endure.''}{F.~William Lawvere, Categories Mailing List (October 20, 1998)}

Taking Lawvere's injunction seriously means asking why single-sorted $\mathsf{ZF}$ requires infinite first-order schemata for Separation and Replacement---because they tangle two distinct jobs together: \textbf{logical carving} (intersections, complements, pullbacks, and relational images) and \textbf{size compression} (certifying that a carved-out collection is small enough to belong to further collections). Passing to Joyal and Moerdijk's \textbf{Algebraic Set Theory (AST)}~\cite{JoyalMoerdijk95} separates those two jobs cleanly into an ambient class category $\mathcal{C}$ and a small-map ideal $\mathcal{S}$---provided we do not bake Replacement into $\mathcal{S}$ prematurely.

\begin{definition}[Predicative Class-Set Ambient $(\mathcal{C}, \mathcal{S})$ and Contravariant Powerclasses]\label{def:small-map-ambient}%
\uses{def:comparison-triple}%
A baseline \emph{predicative class-set ambient} $(\mathcal{C}, \mathcal{S})$ consists of:
\begin{enumerate}[label=\textup{(\roman*)},leftmargin=1.8em]
  \item an ambient Boolean regular (or, constructively, Heyting regular) category $\mathcal{C}$ of \emph{classes}, finitely generated in the two-sorted set/class case by G\"odel's eight predicative relational combinators \textbf{G1--G8}; and
  \item a stable subcategory $\mathcal{S} \subseteq \mathcal{C}$ of \emph{small maps} containing isomorphisms and terminal projections $0 \to 1$, $1 + 1 \to 1$, closed under composition, pullbacks along arbitrary morphisms of $\mathcal{C}$ (inducing Subclass Separation when $\mathrm{Mono} \subseteq \mathcal{S}$, i.e.\ subobject smallness $U \rightarrowtail B \in \mathcal{S}$ holds slice-internally across \emph{every} slice $\mathcal{C}/B$, not merely over global $1_{\mathcal{C}}$~\cite[\S I.1]{JoyalMoerdijk95}), diagonals $\Delta_C : C \rightarrowtail C \times C$ for small-element equality, and finite disjoint sums---\emph{without} assuming quotient or regular-image descent.
\end{enumerate}
An object $X \in \mathcal{C}$ is \emph{small} (a \emph{set}, written $X \in \mathcal{S}$) if $X \to 1$ belongs to $\mathcal{S}$. Under Joyal--Moerdijk's \textbf{Powerclass Axiom (P1)}~\cite[\S I.1]{JoyalMoerdijk95}---asserting that every class $C \in \mathcal{C}$ admits a powerclass object $\mathcal{P}_{\mathcal{S}}(C) \in \mathcal{C}$ with a small membership relation $\in_C \rightarrowtail C \times \mathcal{P}_{\mathcal{S}}(C)$ ($\pi_2 \in \mathcal{S}$) classifying $\mathcal{S}$-small subrelations $R \rightarrowtail D \times C$ ($\pi_1 \in \mathcal{S}$)---$\mathcal{P}_{\mathcal{S}} : \mathcal{C}^{\mathrm{op}} \to \mathcal{C}$, $f \mapsto f^*$, acts contravariantly and $\Omega := \mathcal{P}_{\mathcal{S}}(1)$ is the $\mathcal{S}$-small subobject classifier (coinciding with the full subobject classifier of $\mathcal{C}$ when all monomorphisms belong to $\mathcal{S}$).
\end{definition}

Because G\"odel's operations G1--G8 replace infinite comprehension schemata by finite relational algebra, $\mathcal{C}$ carries zero hidden ``overweight existence'' commitments. Running the existence axioms of set theory through $(\mathcal{C}, \mathcal{S})$ splits them into two sharply contrasting strata.

\subsection{\texorpdfstring{Stratum~I: Honest Lightweight Plumbing (\emph{Empty Set, Pairing, Finite Union})}{Stratum I: Honest Lightweight Plumbing (Empty Set, Pairing, Finite Union)}}\label{subsec:ast-stratum-i}

Categorically, Empty Set, Pairing, and Finite Union assert that $\mathcal{C}$ and $\mathcal{S}$ possess a strict initial object $0$, binary coproduct $1 + 1$, pushouts, and disjoint finite sums---the basic grammar of a \textbf{lextensive regular category}~\cite[Ch.~I]{JoyalMoerdijk95} (or, adjoining a Natural Numbers Object $\mathrm{NNO}$, Maietti and Joyal's \textbf{Arithmetic Universe}~\cite{Maietti05}, calibrated along Arabic Axis~1, ambient universe deformation, in Chapter~\ref{ch:base-fiber}, \S\ref{sec:axis1-deformations}, and along Roman Axes~II--III in Chapter~\ref{ch:panorama}, \S\ref{sec:predicativity-wtypes}). Weakening Stratum~I precipitates no cohomological or topological residua; it merely dismantles the ambient doctrine's ability to state comparison cells $\kappa : \Phi \Rightarrow \Psi$. These axioms pull their exact syntactic weight.

\subsection{\texorpdfstring{Stratum~II: Overweight Size-Compression Valves (\emph{Power Set, Replacement, Unbounded Separation})}{Stratum II: Overweight Size-Compression Valves (Power Set, Replacement, Unbounded Separation)}}\label{subsec:ast-stratum-ii}

Inside $(\mathcal{C}, \mathcal{S})$ under Powerclass Axiom (P1), every class $C \in \mathcal{C}$ \emph{already} possesses a powerclass $\mathcal{P}_{\mathcal{S}}(C) \in \mathcal{C}$, and every class morphism $f : I \to C$ \emph{already} possesses a regular image $\mathrm{im}(f) \rightarrowtail C$ in $\mathcal{C}$. The heavyweight existence axioms of $\mathsf{ZF}$ act strictly as \textbf{size-compression valves} forcing these class-sized constructions back down into $\mathcal{S}$.

\begin{proposition}[Pointillist Continuum in $H(\aleph_1)$ and Constructive $\Omega^X / 2^X$ Split]\label{prop:pointillist-continuum}%
\uses{def:small-map-ambient,rem:zfc-minus-notation}%
Consider the Power Set compression principle $X \in \mathcal{S} \implies \mathcal{P}_{\mathcal{S}}(X) \in \mathcal{S}$ inside $(\mathcal{C}, \mathcal{S})$:
\begin{enumerate}[label=\textup{(\roman*)},leftmargin=1.8em]
  \item Let $H(\aleph_1)$ be the universe of hereditarily countable sets equipped with its definable classes. Then $H(\aleph_1) \models \mathsf{ZFC}^-$~\cite{GitmanHamkinsJohnstone16} (satisfying Infinity, Collection/Replacement, and Choice while refuting Power Set): every individual real number $r \subseteq \mathbb{N}$ lies in $\mathcal{S}$, whereas Cantor space $2^{\mathbb{N}} = \mathcal{P}_{\mathcal{S}}(\mathbb{N})$ and the continuum $\mathbb{R}$ remain proper classes ($2^{\mathbb{N}} \notin \mathcal{S}$).
  \item Over constructive Aczel--Myhill set theory $\mathsf{CZF}$~\cite{Aczel78} (dropping Excluded Middle $\mathrm{LEM}$ along \textbf{Axis~II}, Constructive Logic), \emph{Exponentiation} ($2^X \in \mathcal{S}$, classifying decidable subobjects $\mathrm{Sub}_{\mathrm{dec}}(X)$) is strictly weaker than \emph{Power Set} ($\Omega^X = \mathcal{P}_{\mathcal{S}}(X) \in \mathcal{S}$, classifying arbitrary subobjects $\mathrm{Sub}(X)$), precipitating a nontrivial subobject-classifier size residuum $\Omega^X / 2^X$.
\end{enumerate}
\end{proposition}

\begin{proof}
For~(i), $\mathbb{N} = \omega \in H(\aleph_1)$, and for every $r \subseteq \mathbb{N}$ its transitive closure $r \cup \omega$ is countable, so $r \in H(\aleph_1)$. Likewise, any countable family of hereditarily countable sets has countable union and transitive closure, verifying Collection/Replacement in $H(\aleph_1)$. By Cantor's diagonal theorem, however, $\mathcal{P}_{\mathcal{S}}(\mathbb{N})$ is uncountable and cannot lie in $H(\aleph_1)$, remaining an honest proper class in $\mathcal{C} \setminus \mathcal{S}$.

For~(ii), inside a constructive Heyting $(\mathcal{C}, \mathcal{S})$ without $\mathrm{LEM}$, $2 = 1 + 1 \in \mathcal{S}$ embeds properly into $\Omega = \mathcal{P}_{\mathcal{S}}(1)$ along Axis~II ($\kappa : 1 + 1 \hookrightarrow \Omega$). Subset Collection in $\mathsf{CZF}$ implies Exponentiation ($Y^X \in \mathcal{S}$ for all $X, Y \in \mathcal{S}$, so $2^X \in \mathcal{S}$) while remaining $\Pi^0_2$-arithmetically conservative (equiconsistent, with identical proof-theoretic strength) over Kripke--Platek set theory with Infinity ($\mathsf{KP}\omega$) via Aczel's Martin-L\"of type-universe interpretation~\cite{Aczel78}. Forcing $\Omega \in \mathcal{S}$ ($\Omega^X \in \mathcal{S}$) impredicatively compresses all unbounded subobject predicates into $\mathcal{S}$, jumping far past $\mathsf{KP}\omega$ and full second-order arithmetic $\mathsf{Z}_2$.
\end{proof}

Next, Joyal and Moerdijk~\cite[Ch.~I, \S\S 1--3]{JoyalMoerdijk95} showed that \textbf{Replacement} corresponds precisely to closing $\mathcal{S}$ under quotients and regular images---the exact operation needed to make $\mathcal{P}_{\mathcal{S}}$ covariantly functorial!

\begin{definition}[Axis~VI Doctrinal Comparison Triple: Quotient-Image Descent and Cofinality]\label{def:axis-vi-replacement}%
\uses{def:comparison-triple,def:small-map-ambient}%
Over a baseline predicative class-set ambient $(\mathcal{C}, \mathcal{S})$ (Definition~\ref{def:small-map-ambient}), where an $I$-indexed family of small objects over $I \in \mathcal{S}$ is represented internally as a small map $p : E \to I$ in $\mathcal{S}$, \textbf{Axis~VI (Small-Indexed Cofinality / Replacement)}\index{Axis VI (Quotient-Image Descent / Replacement)} compares the class regular image $\Phi_{\mathrm{Im}}(I, f) := \mathrm{im}(f) \in \mathcal{C}$ of any morphism $f : I \to C$ with small domain $I \in \mathcal{S}$ against the smallness reflector $\mathcal{S} \hookrightarrow \mathcal{C}$, via the descent condition $\kappa_{\mathrm{Rep}} : \mathrm{im}(f) \in \mathcal{S}$ (inducing regular-quotient / coequalizer descent along regular epimorphisms $q : I \twoheadrightarrow Q$ in $\mathcal{C}$). Equivalently, Axis~VI asserts that covariant direct-image pushforward $\exists_g : \mathcal{P}_{\mathcal{S}}(C) \to \mathcal{P}_{\mathcal{S}}(D)$ lands in small subobjects for every $g : C \to D$ in $\mathcal{C}$ (promoting $\mathcal{P}_{\mathcal{S}} : \mathcal{C} \to \mathcal{C}$ to a covariant endofunctor) and, combined with Small Union ($I \in \mathcal{S} \wedge (p : E \to I) \in \mathcal{S} \implies E \in \mathcal{S}$) and coequalizer descent, compresses every small-diagram colimit $\varinjlim_{i \in I} F_i$ into $\mathcal{S}$.
\end{definition}

\begin{remark}[The Borel Determinacy Pincer and Martin's Proper-Class Unraveling Coalgebra]\label{rem:martin-borel-coalgebra}%
\uses{prop:pointillist-continuum,prop:etcs-colimit-escape,def:axis-vi-replacement}%
Comparing Proposition~\ref{prop:pointillist-continuum}(i) against Proposition~\ref{prop:etcs-colimit-escape} exposes a sharp reverse-mathematical pincer across Stratum~II. Inside an ambient universe $V \models \mathsf{ZFC}$, winning strategies on $\omega^\omega$ are hereditarily countable, so the true hereditarily countable universe $H(\aleph_1) \models \mathsf{ZFC}^-$ inherits full Borel Determinacy from $V$; what breaks already at finite Borel rank $\boldsymbol{\Sigma}^0_4$ is \emph{provability inside Power-Set-free $\mathsf{ZFC}^-$} itself, refuted in minimal or non-standard $\beta$-models of $\mathsf{ZFC}^-$ lacking $\mathcal{P}_{\mathcal{S}}(\mathbb{N})$~\cite{Friedman71,Martin75,MontalbanShore12,GitmanHamkinsJohnstone16}. Conversely, in Replacement-free $\mathsf{ZC}$ (and $\mathsf{ETCS} \simeq \mathsf{Mac}$), all reals, trees, and strategies $\sigma : \omega^{<\omega} \to \omega$ ($\in V_{\omega+1}$), as well as Borel subsets of Cantor space ($\in V_{\omega+2}$), already reside inside the true cumulative rank $V_{\omega+\omega} \subset V \models \mathsf{ZFC}$, which therefore inherits full Borel Determinacy from $V$. What breaks sharply at $\boldsymbol{\Sigma}^0_\omega$ is \emph{provability inside $\mathsf{ZC}$} (and $\mathsf{ETCS} \simeq \mathsf{Mac}$), refuted in non-standard or minimal models of $\mathsf{ZC}$~\cite{Friedman71}. Inside $(V_{\omega+\omega}, \mathrm{Def}(V_{\omega+\omega}))$, each finite iterate $\mathcal{P}^k(\mathbb{N})$ lies in $\mathcal{S}$---allowing $\mathsf{ZC}$ to prove $\boldsymbol{\Sigma}^0_k$-determinacy for every fixed meta-integer $k < \omega$~\cite{Martin75}---whereas the countable family $\{\mathcal{P}^n(\mathbb{N})\}_{n \in \mathbb{N}}$ escapes $\mathcal{S}$ and prevents Martin's inductive unraveling construction from internalizing inside $\mathcal{S}$. Inside $(\mathcal{C}, \mathcal{S})$, Tony Martin's inductive clopen unraveling of a $\boldsymbol{\Sigma}^0_\omega$ game still assembles into a well-defined \emph{proper-class covering tree in $\mathcal{C} \setminus \mathcal{S}$}, framing the class-to-set strategy-descent obstruction $\mathrm{Res}(\kappa_{\mathrm{Det}})$ as a flagship Volume~II problem. Meanwhile, \textbf{Separation / Reflection} stratifies into four distinct rungs: $\Delta_0$-pullbacks along $\top : 1 \to \Omega$ in $\mathsf{ETCS} \simeq \mathsf{Mac}$, full unbounded $\mathcal{L}_\in$-Separation in $\mathsf{ZC}$ (strictly stronger than $\mathsf{Mac}$ by Mathias~\cite{Mathias01}, yet satisfied in $V_{\omega+\omega}$ without Replacement), the full L\'evy--Montague Reflection Schema ($\Sigma_n$-Reflection to cumulative rank sets $V_\alpha$ for every meta-integer $n < \omega$, strictly beyond any bounded $\Sigma_k$-Reflection fragment) upgrading $\mathsf{ZC}$ to $\mathsf{ZFC}$, and impredicative class comprehension jumping from conservative $\mathsf{NBG}$ to $\mathsf{KM}$ ($\mathrm{Con}(\mathsf{NBG}+\mathsf{ETR}_{\mathrm{Ord}})$).
\end{remark}

\section{\texorpdfstring{Over-Quotienting Structure: The Anti-Existence Residua}{Over-Quotienting Structure: The Anti-Existence Residua}}\label{sec:anti-existence-residua}

Where Stratum~II manufactures residua by over-compressing proper-class constructions downward into $\mathcal{S}$, \textbf{Stratum~III} inverts the polarity completely. \textbf{Foundation (Regularity)} and \textbf{Extensionality} assert the existence of no sets whatsoever; they are \emph{structural prohibitions} that quotient away circular $\in$-graphs and internal symmetries. Relaxing either clamp uncovers universe-scale instances of the exact fixed-point ($\mu F \Rightarrow \nu F$) and truncation operators studied in Chapters~\ref{ch:delta-rings}, \ref{ch:higher-crm}, and~\ref{ch:realizability-sites}.

\begin{theorem}[Axis~VII Comparison: Foundation vs.\ Anti-Foundation as $\mathbf{Coalg}_{\mathrm{wf}}(\mathcal{P}_{\mathcal{S}}) \hookrightarrow \mathbf{Coalg}(\mathcal{P}_{\mathcal{S}})$]\label{thm:axis-vii-foundation}%
\lean{HomotopyCRM.SetsResiduum.pset_equiv_iff_bisimilar, HomotopyCRM.SetsResiduum.quine_atom_outside_wf}%
\leanwitness%
\uses{def:small-map-ambient,def:axis-vi-replacement}%
Let $(\mathcal{C}, \mathcal{S})$ be a Barr-exact class-set ambient with $\mathbb{N} \in \mathcal{S}$, Small Union, and Axis~VI image descent (Definition~\ref{def:axis-vi-replacement}), so that $\mathcal{P}_{\mathcal{S}} : \mathcal{C} \to \mathcal{C}$ acts covariantly by direct image $\exists_g$, and assume $\mathcal{P}_{\mathcal{S}}$ admits both an initial algebra $(\mu \mathcal{P}_{\mathcal{S}}, \mathrm{sup})$ and a terminal coalgebra $(\nu \mathcal{P}_{\mathcal{S}}, \mathrm{des})$. Let $\mathbf{Coalg}_{\mathrm{wf}}(\mathcal{P}_{\mathcal{S}}) \hookrightarrow \mathbf{Coalg}(\mathcal{P}_{\mathcal{S}})$ denote the full subdoctrine of well-founded $\mathcal{P}_{\mathcal{S}}$-coalgebras $(X, \gamma : X \to \mathcal{P}_{\mathcal{S}}(X))$ (satisfying internal $\in_\gamma$-induction). Then:
\begin{enumerate}[label=\textup{(\roman*)},leftmargin=1.8em]
  \item \textbf{Well-Founded Coalgebra Terminality (Osius~\cite{Osius74}):} The cumulative well-founded universe $V_{\mathsf{ZF}}$ satisfying Foundation\index{Axis VII (Foundation vs.\ Anti-Foundation)} is canonically isomorphic to the initial algebra $\mu \mathcal{P}_{\mathcal{S}}$, which via $\mathrm{sup}^{-1} : \mu \mathcal{P}_{\mathcal{S}} \xrightarrow{\sim} \mathcal{P}_{\mathcal{S}}(\mu \mathcal{P}_{\mathcal{S}})$ forms the terminal object $\nu^{\mathrm{wf}} \mathcal{P}_{\mathcal{S}} \simeq \mu \mathcal{P}_{\mathcal{S}}$ of $\mathbf{Coalg}_{\mathrm{wf}}(\mathcal{P}_{\mathcal{S}})$~\cite[Chs.~II--III]{JoyalMoerdijk95}.
  \item \textbf{Unrestricted Coalgebra Terminality (Aczel~\cite[Chs.~1--2, 6--7]{Aczel88}):} Peter Aczel's Anti-Foundation universe $V_{\mathsf{AFA}}$\index{Anti-Foundation Axiom ($\mathsf{AFA}$)}---in which every accessible pointed $\mathcal{S}$-small-branching graph $(G, \to, r)$ admits a unique decoration $d : G \to V$ satisfying $d(x) = \{d(y) \mid x \to y\}$---is canonically isomorphic to the terminal object $\nu \mathcal{P}_{\mathcal{S}}$ of $\mathbf{Coalg}(\mathcal{P}_{\mathcal{S}})$.
  \item \textbf{Strict Comparison Monomorphism and Factorization Criterion:} Passing along the doctrinal inclusion $\mathbf{Coalg}_{\mathrm{wf}}(\mathcal{P}_{\mathcal{S}}) \hookrightarrow \mathbf{Coalg}(\mathcal{P}_{\mathcal{S}})$ induces a canonical comparison monomorphism $\kappa_{\mathrm{Reg}} : \nu^{\mathrm{wf}} \mathcal{P}_{\mathcal{S}} \simeq \mu \mathcal{P}_{\mathcal{S}} \hookrightarrow \nu \mathcal{P}_{\mathcal{S}}$. Whenever $0 \not\cong 1$, $\kappa_{\mathrm{Reg}}$ is a \emph{strict} monomorphism whose non-well-founded residuum $\mathrm{Res}(\kappa_{\mathrm{Reg}}) = \nu \mathcal{P}_{\mathcal{S}} \setminus \mu \mathcal{P}_{\mathcal{S}}$ is anchored by the unique Quine atom\index{Quine atom ($q = \{q\}$)} $q = \{q\}$, and a coalgebra $(X, \gamma) \in \mathbf{Coalg}(\mathcal{P}_{\mathcal{S}})$ lies in $\mathbf{Coalg}_{\mathrm{wf}}(\mathcal{P}_{\mathcal{S}})$ (satisfies Foundation) if and only if its terminal decoration $d_\gamma : X \to \nu \mathcal{P}_{\mathcal{S}}$ factors through $\kappa_{\mathrm{Reg}} : \mu \mathcal{P}_{\mathcal{S}} \hookrightarrow \nu \mathcal{P}_{\mathcal{S}}$.
\end{enumerate}
\end{theorem}

\begin{proof}
By Lambek's Lemma (formalized as \texttt{lambek\_isIso} in Chapter~\ref{ch:delta-rings}, \S\ref{sec:delta-rings}), $\mathrm{sup} : \mathcal{P}_{\mathcal{S}}(\mu \mathcal{P}_{\mathcal{S}}) \xrightarrow{\sim} \mu \mathcal{P}_{\mathcal{S}}$ and $\mathrm{des} : \nu \mathcal{P}_{\mathcal{S}} \xrightarrow{\sim} \mathcal{P}_{\mathcal{S}}(\nu \mathcal{P}_{\mathcal{S}})$ are isomorphisms, modeling extensional membership via $x \in y \iff x \in \mathrm{sup}^{-1}(y)$ (respectively $x \in \mathrm{des}(y)$). Initiality of $\mu \mathcal{P}_{\mathcal{S}}$ yields $\in$-induction on $V_{\mathsf{ZF}} \cong \mu \mathcal{P}_{\mathcal{S}}$, and Osius's recursive coalgebra theorem~\cite{Osius74}~\cite[Chs.~II--III]{JoyalMoerdijk95} equips every $(X, \gamma) \in \mathbf{Coalg}_{\mathrm{wf}}(\mathcal{P}_{\mathcal{S}})$ with a unique coalgebra homomorphism $m_\gamma : (X, \gamma) \to (\mu \mathcal{P}_{\mathcal{S}}, \mathrm{sup}^{-1})$, proving $\nu^{\mathrm{wf}} \mathcal{P}_{\mathcal{S}} \simeq \mu \mathcal{P}_{\mathcal{S}}$ and~(i). Moreover, initiality of $(\mu \mathcal{P}_{\mathcal{S}}, \mathrm{sup})$ (equivalently, terminality of $(\nu \mathcal{P}_{\mathcal{S}}, \mathrm{des})$) induces the unique algebra-coalgebra homomorphism $\kappa_{\mathrm{Reg}} : \mu \mathcal{P}_{\mathcal{S}} \to \nu \mathcal{P}_{\mathcal{S}}$ with $\mathrm{des} \circ \kappa_{\mathrm{Reg}} = \exists_{\kappa_{\mathrm{Reg}}} \circ \mathrm{sup}^{-1}$. By $\in$-induction on $x \in \mu \mathcal{P}_{\mathcal{S}}$, if $\kappa_{\mathrm{Reg}}(u) = \kappa_{\mathrm{Reg}}(z) \implies u = z$ holds for all $u \in \mathrm{sup}^{-1}(x)$ and all $z \in \mu \mathcal{P}_{\mathcal{S}}$, then $\kappa_{\mathrm{Reg}}(x) = \kappa_{\mathrm{Reg}}(y)$ gives $\exists_{\kappa_{\mathrm{Reg}}}(\mathrm{sup}^{-1}(x)) = \exists_{\kappa_{\mathrm{Reg}}}(\mathrm{sup}^{-1}(y))$ and hence $\mathrm{sup}^{-1}(x) = \mathrm{sup}^{-1}(y)$ and $x = y$; thus $\kappa_{\mathrm{Reg}}$ is monic.

For~(ii) and~(iii), any accessible pointed small-branching graph $(G, \to, r)$ defines $(G, \gamma) \in \mathbf{Coalg}(\mathcal{P}_{\mathcal{S}})$ via $\gamma(x) := \{y \in G \mid x \to y\}$. Here $\mathbb{N} \in \mathcal{S}$, Small Union, and Axis~VI Replacement compress every reachable transitive subgraph $G_x = \bigcup_{n \in \mathbb{N}} \gamma^n(\{x\})$ into $\mathcal{S}$, while Barr-exactness of $\mathcal{C}$ (Chapter~\ref{ch:panorama}, Figure~\ref{fig:homological-lattice}) directly quotients the class of small APGs by maximum bisimulation to realize $\nu \mathcal{P}_{\mathcal{S}}$ in Foundation-free ambients without circularly invoking rank-based Scott's trick. Terminality of $(\nu \mathcal{P}_{\mathcal{S}}, \mathrm{des})$ yields a unique homomorphism $d_\gamma : (G, \gamma) \to (\nu \mathcal{P}_{\mathcal{S}}, \mathrm{des})$, i.e.\ $d_\gamma(x) = \{d_\gamma(y) \mid x \to y\}$ (Aczel's $\mathsf{AFA}$~\cite[Chs.~1--2]{Aczel88}). If $(X, \gamma) \in \mathbf{Coalg}_{\mathrm{wf}}(\mathcal{P}_{\mathcal{S}})$, uniqueness forces $d_\gamma = \kappa_{\mathrm{Reg}} \circ m_\gamma$. Conversely, if $d_\gamma = \kappa_{\mathrm{Reg}} \circ m$, monicity of $\kappa_{\mathrm{Reg}}$ makes $\exists_{\kappa_{\mathrm{Reg}}}$ monic, so $\exists_{\kappa_{\mathrm{Reg}}} \circ \exists_m \circ \gamma = \mathrm{des} \circ d_\gamma = \exists_{\kappa_{\mathrm{Reg}}} \circ \mathrm{sup}^{-1} \circ m$ implies $\mathrm{sup}^{-1} \circ m = \exists_m \circ \gamma$; for any $\in_\gamma$-inductive subobject $S \rightarrowtail X$, the Heyting/Boolean universal image $\forall_m(S) \rightarrowtail \mu \mathcal{P}_{\mathcal{S}}$ (Definition~\ref{def:small-map-ambient}(i)) is $\in$-inductive on $\mu \mathcal{P}_{\mathcal{S}}$, giving $\forall_m(S) = \mu \mathcal{P}_{\mathcal{S}}$ and $S = X$. Finally, whenever $0 \not\cong 1$, decorating $(1, \top : 1 \to \mathcal{P}_{\mathcal{S}}(1))$ yields $q = d_\top(\bullet) \in \nu \mathcal{P}_{\mathcal{S}}$ with $q = \{q\}$, refuting factorization through $\kappa_{\mathrm{Reg}}$ and witnessing $\nu \mathcal{P}_{\mathcal{S}} \setminus \mu \mathcal{P}_{\mathcal{S}} \neq \emptyset$.
\end{proof}

\begin{proposition}[Boffa Hypersets and Rieger--Bernays Slice vs.\ Global Torsors]\label{prop:boffa-torsors}%
\lean{HomotopyCRM.SetsResiduum.quineAPG_bisimilar_and_nonIso_twoCycle, HomotopyCRM.SetsResiduum.boffa_twin_torsor, HomotopyCRM.SetsResiduum.boffa_twin_addTorsor}%
\leanwitness%
\uses{thm:axis-vii-foundation}%
Weakening Aczel's $\mathsf{AFA}$ (which quotients coalgebras by maximum bisimulation) to Maurice Boffa's Anti-Foundation Axiom ($\mathsf{BAFA}$\index{BAFA@$\mathsf{BAFA}$ (Boffa Anti-Foundation Axiom)}~\cite{Boffa69}, requiring every injective decoration of a transitive subgraph of an extensional accessible pointed graph to extend to an injective decoration of the supergraph) lifts $\nu \mathcal{P}_{\mathcal{S}}$ to a universe $V_{\mathsf{BAFA}}$ admitting a proper class of distinct Quine atoms $a = \{a\} \neq b = \{b\}$. Then:
\begin{enumerate}[label=\textup{(\roman*)},leftmargin=1.8em]
  \item \textbf{Isomorphism vs.\ Bisimulation Fibers:} Mere bisimilarity $[A] = [B] \in \nu \mathcal{P}_{\mathcal{S}}$ does not imply $\in$-isomorphism (for instance, a Quine singleton $a = \{a\}$ with $|\mathrm{tc}(\{a\})| = 1$ and a mutual $2$-cycle $u = \{v\}, v = \{u\}$ with $|\mathrm{tc}(\{u\})| = 2$ share the same bisimulation image $q \in \nu \mathcal{P}_{\mathcal{S}}$ while $\mathrm{Iso}_\in(\mathrm{tc}(\{a\}), \mathrm{tc}(\{u\})) = \emptyset$). Restricted to hypersets $A, B \in V_{\mathsf{BAFA}}$ with $\in$-isomorphic transitive closures, $\mathrm{Iso}_\in(\mathrm{tc}(\{A\}), \mathrm{tc}(\{B\}))$ carries a free and transitive right $\mathrm{Aut}_\in(\mathrm{tc}(\{A\}))$-action.
  \item \textbf{Slice Twin Torsors vs.\ Global Parity Torsors Over $Q_0$:}\index{Boffa hyperset torsor} Fix any countably infinite small set $Q_0 \in \mathcal{S}$ of distinct Quine atoms in $V_{\mathsf{BAFA}}$ (obtained by decorating $\omega$ disjoint self-loops), and pass via Rieger--Bernays~\cite{Bernays54,Rieger57} $\in$-permutations fixing $V_{\mathsf{ZF}}$ pointwise on the cumulative subuniverse $V_{\mathsf{ZF}}(Q_0) \subseteq V_{\mathsf{BAFA}}$ to the equivariant permutation topos $\mathcal{E}_{Q_0} := \mathbf{Cont}(\mathrm{Sym}(Q_0)) \simeq \mathbf{Nom}$ (Chapter~\ref{ch:realizability-sites}, \S\ref{subsec:kihara-choice-nominal}). Over the unordered-pair orbit sheaf $[Q_0]^2 \in \mathcal{E}_{Q_0}$, the ordered-twin cover $Q_0^{\#(2)} \twoheadrightarrow [Q_0]^2$ is a pointless $\Delta(\mathbb{Z}/2\mathbb{Z})$-torsor in the \textbf{slice topos} $\mathcal{E}_{Q_0}/[Q_0]^2$,
  \[
    [Q_0^{\#(2)}] \;\neq\; 0 \;\in\; H^1\bigl(\mathcal{E}_{Q_0}/[Q_0]^2;\, \Delta(\mathbb{Z}/2\mathbb{Z})\bigr)
  \]
  (restricting on a single isolated twin pair $Q_2 = \{a, b\}$ to $\mathbf{Set}^{\mathbb{Z}/2\mathbb{Z}}$, Chapter~\ref{ch:choice-atlas}, Theorem~\ref{thm:coshep-svc}). Over the \textbf{global terminal object} $1_{\mathcal{E}_{Q_0}}$, Schreier--Ulam pro-discrete blindness forces $H^1(\mathcal{E}_{Q_0}; \Delta(\mathbb{Z}/2\mathbb{Z})) = 0$, and the non-vanishing global obstruction lifts to the odd-Quine-finset parity torsor $[\mathcal{T}_{\mathrm{odd}}(Q_0)] \neq 0 \in H^1(\mathcal{E}_{Q_0}; \mathrm{EvenFinset}(Q_0))$.
\end{enumerate}
\end{proposition}

\begin{proof}
Under $\mathsf{BAFA}$, extending the empty decoration to disjoint self-loops or to a directed $2$-cycle produces distinct Quine atoms $a = \{a\} \neq b = \{b\}$ and $u = \{v\} \neq v = \{u\}$; cardinality $|\mathrm{tc}(\{a\})| = 1 \neq 2 = |\mathrm{tc}(\{u\})|$ separates their $\in$-isomorphism classes inside $[q] \in \nu \mathcal{P}_{\mathcal{S}}$, whereas for $\in$-isomorphic closures composition gives a free transitive $\mathrm{Aut}_\in(\mathrm{tc}(\{A\}))$-action, proving~(i). For~(ii), passing to the cumulative hyperset subuniverse $V_{\mathsf{ZF}}(Q_0) \subseteq V_{\mathsf{BAFA}}$ generated over $V_{\mathsf{ZF}}$ from $Q_0$ as Quine-atom urelements (via Rieger--Bernays~\cite{Bernays54,Rieger57} and Boffa~\cite{Boffa69}), every permutation $\sigma \in \mathrm{Sym}(Q_0)$ extends uniquely by rank recursion ($\tilde{\sigma}(q) = \sigma(q)$ for $q \in Q_0$ and $\tilde{\sigma}(x) = \{\tilde{\sigma}(y) \mid y \in x\}$ otherwise) to an $\in$-automorphism of $V_{\mathsf{ZF}}(Q_0)$ fixing $V_{\mathsf{ZF}}$ pointwise. For any $\{a, b\} \in [Q_0]^2$, the transposition $(a\;b) \in \mathrm{Sym}(Q_0)$ stabilizes $\{a, b\}$ while swapping $(a, b) \leftrightarrow (b, a)$ in $Q_0^{\#(2)}$, blocking any $\mathrm{Sym}(Q_0)$-equivariant section $[Q_0]^2 \to Q_0^{\#(2)}$ and proving $[Q_0^{\#(2)}] \neq 0$ in $H^1(\mathcal{E}_{Q_0}/[Q_0]^2; \Delta(\mathbb{Z}/2\mathbb{Z}))$, while Chapter~\ref{ch:realizability-sites} (\S\ref{subsec:kihara-choice-nominal}, \texttt{NominalSimplicity.lean}) establishes both $H^1(\mathcal{E}_{Q_0}; \Delta(\mathbb{Z}/2\mathbb{Z})) = 0$ and $[\mathcal{T}_{\mathrm{odd}}(Q_0)] \neq 0 \in H^1(\mathcal{E}_{Q_0}; \mathrm{EvenFinset}(Q_0))$.
\end{proof}

Proposition~\ref{prop:boffa-torsors} acquires an illuminating homotopy-theoretic explanation as soon as we examine the companion Stratum~III clamp: \textbf{Extensionality}.

\begin{proposition}[{External Benchmark: Material Extensionality, Quotient Elimination, and Hereditary Truncation; Friedman~1973~\cite{Friedman73}, Gylterud~2018~\cite[Thm.~3.9]{Gylterud18}}]\label{prop:extensionality-truncation}%
\lean{HomotopyCRM.SetsResiduum.zfset_ext_quot_gate}%
\leanwitness%
\uses{thm:axis-vii-foundation,prop:boffa-torsors}%
Material Extensionality ($\forall z\,(z \in X \leftrightarrow z \in Y) \implies X = Y$) enforces both constructive quotient collapse and homotopical symmetry truncation:
\begin{enumerate}[label=\textup{(\roman*)},leftmargin=1.8em]
  \item \textbf{Constructive Equiconsistency vs.\ Direct Quotient Cost (Friedman~1973~\cite{Friedman73}):} Classical $\mathsf{ZF}$, intuitionistic $\mathsf{IZF}$, and Friedman's non-extensional intuitionistic set theory $\mathsf{IZF}_{\mathrm{non\text{-}ext}}$ are \emph{equiconsistent} ($\mathrm{Con}(\mathsf{IZF}_{\mathrm{non\text{-}ext}}) \implies \mathrm{Con}(\mathsf{ZF})$ via Friedman's recursive double-negation extensionality translation). Nevertheless, $\mathsf{IZF}$ is not directly set-interpretable in $\mathsf{IZF}_{\mathrm{non\text{-}ext}}$ without Collection or primitive quotients: Scott's classical trick of truncating each equivalence class to its small set of minimal-rank members $\{y \sim x \mid \forall z \sim x,\, \mathrm{rank}(y) \le \mathrm{rank}(z)\}$ (working in classical $\mathsf{ZF}$ without Choice) fails constructively because inhabited classes of ordinals need not have $\in$-least elements.
  \item \textbf{Univalent Iterative Multisets vs.\ Hereditary $(-1)$-Truncation (Gylterud~2018~\cite[Thm.~3.9]{Gylterud18}):}\index{Iterative multisets ($\mu \mathcal{P}_{\mathrm{Anima}}$)} Inside an $(\infty, 1)$-topos $\mathcal{X}$ with univalent universe $\mathcal{U}$~\cite[\S 6.1.6]{Lurie09}, let $M := \mu \mathcal{P}_{\mathrm{Anima}}$ be the initial algebra ($W$-type) of the polynomial endofunctor $\mathcal{P}_{\mathrm{Anima}}(X) := \sum_{A : \mathcal{U}} (A \to X)$ of \emph{iterative multisets}. Embedding any small $0$-type $A : \mathcal{U}_{\le 0}$ as the flat multiset $\underline{A} := \mathrm{sup}(A, \mathrm{const}_{\emptyset_M}) : M$ preserves internal symmetries via loop groups
  \[
    \pi_1(M, \underline{A}) \;\simeq\; \mathrm{Aut}(A).
  \]
  Applying naive outer $0$-truncation $\tau_{\le 0}(M)$ kills $\pi_1$ but \emph{retains element multiplicities} (since $X_1 := \mathrm{sup}(\mathbf{1}, \mathrm{const}_{\emptyset_M})$ and $X_2 := \mathrm{sup}(\mathbf{2}, \mathrm{const}_{\emptyset_M})$ satisfy $(X_1 =_M X_2) \simeq (\mathbf{1} \simeq \mathbf{2}) \simeq \emptyset$). Recovering the flat cumulative hierarchy of iterative sets $V_0 \hookrightarrow M$ requires \textbf{hereditary fiberwise $(-1)$-truncation}---restricting each branching map $f : A \to M$ to be a $(-1)$-truncated embedding $A \hookrightarrow M$ (or quotienting recursively by bisimulation)---which simultaneously eliminates multiplicities and forces $V_0$ to be $0$-truncated.
\end{enumerate}
\end{proposition}

\begin{remark}[Unifying Boffa Twin Torsors and Loop Groups via $\nu \mathcal{P}_{\mathrm{Anima}}$]\label{rem:nu-anima-fossil}%
\uses{prop:boffa-torsors,prop:extensionality-truncation}%
Propositions~\ref{prop:boffa-torsors} and~\ref{prop:extensionality-truncation} exhibit Boffa's twin Quine set $T = \{a, b\}$ as the \emph{$0$-truncated $1$-categorical fossil of higher homotopy loop groups}. In well-founded $M = \mu \mathcal{P}_{\mathrm{Anima}}$, even $\underline{\mathbf{2}}$ carries $\pi_1(M, \underline{\mathbf{2}}) \cong \mathbb{Z}/2\mathbb{Z}$, whereas in $V_{\mathsf{BAFA}}$ well-founded sets are rigid and symmetry survives only on non-well-founded twins $a = \{a\} \neq b = \{b\}$ whose Aczel quotient $\{q, q\} = \{q\} \in \nu \mathcal{P}_{\mathcal{S}}$ collapses multiplicity via $(-1)$-truncated image $\exists_d$. Concretely, inside an $(\infty, 1)$-topos $\mathcal{X}$ with nested univalent universes $\mathcal{U}_0 \in \mathcal{U}_1$~\cite[\S 6.1.6]{Lurie09}, set $\mathcal{P}_{\mathrm{Anima}}(X) := \sum_{A : \mathcal{U}_0} (A \to X) : \mathcal{U}_1 \to \mathcal{U}_1$ and $X_n := \mathcal{P}_{\mathrm{Anima}}^n(\mathbf{1})$. Because dependent sums over the constant base $\mathcal{U}_0$ ($|\omega^{\mathrm{op}}| \simeq *$) and exponentials $(A \to -)$ preserve sequential homotopy limits, Ahrens, Capriotti, and Spadotti~\cite{AhrensCapriottiSpadotti15} establish that the \textbf{terminal coalgebra ($M$-type) of non-well-founded univalent multisets}\index{Non-well-founded univalent multisets ($\nu \mathcal{P}_{\mathrm{Anima}}$)} exists at $\omega$ as $\nu \mathcal{P}_{\mathrm{Anima}} := \varprojlim_{n \in \omega^{\mathrm{op}}} X_n : \mathcal{U}_1$, with Lambek equivalence $\mathrm{des} : \nu \mathcal{P}_{\mathrm{Anima}} \xrightarrow{\sim} \mathcal{P}_{\mathrm{Anima}}(\nu \mathcal{P}_{\mathrm{Anima}})$. For the $1$-fold Quine multiset $q$ ($\mathrm{des}(q) = (\mathbf{1}, \mathrm{const}_q)$), finite unrollings $q_0 = * : \mathbf{1}$ and $q_{n+1} = (\mathbf{1}, \mathrm{const}_{q_n}) : X_{n+1}$ satisfy $\Omega(X_{n+1}, q_{n+1}) \simeq \mathrm{Aut}(\mathbf{1}) \times \Omega(X_n, q_n)^{\mathbf{1}} \simeq \mathbf{1}$ by induction on $n$, whence $\Omega(\nu \mathcal{P}_{\mathrm{Anima}}, q) \simeq \varprojlim_n \Omega(X_n, q_n) \simeq \mathbf{1}$. Consequently, for the $2$-fold Quine multiset $[q, q] := \mathrm{des}^{-1}(\mathbf{2}, \mathrm{const}_q)$---explicitly distinct from the $2$-branching self-loop $c = \mathrm{des}^{-1}(\mathbf{2}, \mathrm{const}_c)$, whose loop group is the infinite binary-tree automorphism group $\mathrm{Aut}(T_2)$---Lambek's equivalence and univalence ($\mathrm{const}_q \circ e^{-1} \equiv \mathrm{const}_q$ for $e : \mathbf{2} \simeq \mathbf{2}$) yield $\Omega(\nu \mathcal{P}_{\mathrm{Anima}}, [q, q]) \simeq \sum_{e : \mathbf{2} \simeq \mathbf{2}} \Omega(\nu \mathcal{P}_{\mathrm{Anima}}, q)^{\mathbf{2}} \simeq \mathrm{Aut}(\mathbf{2}) \cong \mathbb{Z}/2\mathbb{Z}$. Lifting Boffa's twin Quine graph $G_T = \{a, b, T\}$ ($a = \{a\}$, $b = \{b\}$, $T = \{a, b\}$) to an untruncated $\mathcal{P}_{\mathrm{Anima}}$-coalgebra without $(-1)$-truncation, terminal coalgebra uniqueness on $(\mathbf{1}, \mathrm{id}_{\mathbf{1}})$ sends $a, b \mapsto q$ and hence $d_{\mathrm{Anima}}(T) = \mathrm{des}^{-1}(\mathbf{2}, d_{\mathrm{Anima}} \circ \iota_{a,b}) = [q, q]$, directly uniting Axis~V with Axis~VII.
\end{remark}

\begin{table}[htbp]
\centering
\caption{Master Algebraic Set Theory Residuum Ledger Across $(\mathcal{C}, \mathcal{S})$}\label{tab:master-ast-ledger}
\vspace{0.3em}
{\scriptsize\setlength{\tabcolsep}{2.5pt}\renewcommand{\arraystretch}{1.12}%
\begin{tabularx}{\textwidth}{@{}P{0.15\textwidth}P{0.18\textwidth}P{0.15\textwidth}YY@{}}
\toprule
\textbf{Axiom Family} & \textbf{Role in $(\mathcal{C}, \mathcal{S})$} & \textbf{Polarity} & \textbf{Separating Model} & \textbf{Residuum $\mathrm{Res}(\kappa)$} \\
\midrule
Empty, Pair, Union & $0, 1+1$, finite sums in $\mathcal{C},\mathcal{S}$ & Plumbing & Non-extensive cats & None (ambient grammar) \\
Bernays G1--G8, $\Delta_0$-Sep & Finite limits + $\Delta_0$-pullbacks & Carving & Pre-regular categories & Missing subobject pullbacks \\
Unbounded Sep / $\mathsf{KM}$ & $\mathcal{L}_\in$ / class reflection to $\mathcal{S}$ & Compression & Mathias~\cite{Mathias01} $\mathsf{Mac} \subsetneq \mathsf{ZC}$; $\mathsf{NBG}$ vs.\ $\mathsf{KM}$ & Unreflected $\Pi_n^V$ / $\Pi^1_1$ invariants \\
Power Set & $\mathcal{P}_{\mathcal{S}}(X) \in \mathcal{S}$ ($\Omega^X$ small) & Compression & $H(\aleph_1)$ ($\mathsf{ZFC}^-$); $\mathsf{CZF}$ & Proper-class $2^{\mathbb{N}} \notin \mathcal{S}$; $\Omega^X / 2^X$ \\
Replacement (Axis~VI) & Image descent $\mathrm{im}(f) \in \mathcal{S}$ ($\exists_g$) & Compression & Zermelo $V_{\omega+\omega}$ ($\mathsf{ZC}$ / $\mathsf{ETCS}$) & Colimit escape; $\boldsymbol{\Sigma}^0_\omega$ Borel $\mathsf{ZC}$-unprovability \\
Limitation of Size ($\mathsf{LoS}$) & $C \notin \mathcal{S} \iff C \twoheadrightarrow V$ & Mega-Axiom & $(L_{\omega_\omega^L}, \mathrm{Def})$ over $\mathsf{GB}_0$; Easton & Bundles $\mathrm{Rep} + \mathsf{GC}$ ($\mathrm{Ord}\cong V$) \\
Foundation (Axis~VII) & $\mathbf{Coalg}_{\mathrm{wf}} \hookrightarrow \mathbf{Coalg}$ ($\nu^{\mathrm{wf}}\mathcal{P}_{\mathcal{S}} \simeq \mu\mathcal{P}_{\mathcal{S}}$) & Anti-Cycle & Aczel $V_{\mathsf{AFA}}$; Boffa $V_{\mathsf{BAFA}}$ & $\nu \mathcal{P}_{\mathcal{S}} \setminus \mu \mathcal{P}_{\mathcal{S}}$; slice/parity torsors \\
Extensionality & Hereditary $(-1)$-truncation & Anti-Symmetry & $\mathsf{IZF}_{\mathrm{non\text{-}ext}}$; Gylterud $M$ & Quotient cost; loops $\pi_1(M,\underline{A})$ \\
\bottomrule
\end{tabularx}\par}
\end{table}

\section{\texorpdfstring{Scale-Invariance, the Double-Clamped Universe, and Horizons Beyond ``$\mathbf{Sets}$''}{Scale-Invariance, the Double-Clamped Universe, and Horizons Beyond "Sets"}}\label{sec:scale-invariance-synthesis}

Up through Chapter~\ref{ch:realizability-sites}, a foundational skeptic could still raise a ``turtles all the way down'' objection: Categorical Reverse Mathematics might measure obstructions \emph{inside} an ambient topos or $(\infty, 1)$-category while leaving the meta-theoretic bedrock of ``$\mathbf{Sets}$'' to classical syntactic proof theory. Passing to Joyal--Moerdijk small-map pairs $(\mathcal{C}, \mathcal{S})$ and the baseline adjunction $\mathrm{Def} \dashv \pi_{\mathrm{Set}}$ closes that loop completely by establishing \textbf{reflexive scale-invariance} across the universe boundary.

\begin{theorem}[Reflexive Scale-Invariance of Double-Clamped $\mathbf{Set}_{\mathsf{ZFC}}$ and $\mathsf{NBG}$]\label{thm:scale-invariance}%
\uses{def:set-projection-fibration,prop:large-equiv-splitting,def:small-map-ambient,def:axis-vi-replacement,thm:axis-vii-foundation,prop:extensionality-truncation}%
Let $(\mathcal{C}, \mathcal{S})$ be a non-degenerate ($0 \not\cong 1$) Barr-exact predicative AST ambient (Definition~\ref{def:small-map-ambient}, with monomorphisms in $\mathcal{S}$ so $\Omega := \mathcal{P}_{\mathcal{S}}(1)$ classifies subobjects in $\mathcal{C}$) admitting $\mu \mathcal{P}_{\mathcal{S}}$ and $\nu \mathcal{P}_{\mathcal{S}}$ under Axis~VI image descent, with well-founded cumulative universe $V := \nu^{\mathrm{wf}} \mathcal{P}_{\mathcal{S}} \simeq \mu \mathcal{P}_{\mathcal{S}}$:
\[
  {\small\begin{aligned}
    \text{\textbf{Upper Clamps:}}\quad &\mathbb{N} \in \mathcal{S},\;\; 1+1 \xrightarrow{\sim} \Omega,\;\; \mathcal{P}_{\mathcal{S}}(\mathcal{S}) \subseteq \mathcal{S},\;\; \mathrm{im}(I \xrightarrow{f} C) \in \mathcal{S}, \\
    &\mathrm{IAC}_{\mathcal{S}}\;\; \bigl(\text{+\;}\mathrm{Sub}_{\mathcal{C}}(1)\cong 2,\; 1\text{ proj.; }\mathbf{Card}(V) \hookrightarrow \mathbf{Set}(V)\text{ for }\mathsf{NBG}\bigr) \\[0.25em]
    &\Downarrow \;\; \text{\emph{pinches $(V, \mathcal{C}, \mathcal{S})$ into}}\;\; \mathbf{Set}_{\mathsf{ZFC}} \;\;\Uparrow \\[0.25em]
    \text{\textbf{Lower Clamps:}}\quad &\mathbf{Coalg}_{\mathrm{wf}}(\mathcal{P}_{\mathcal{S}}) \hookrightarrow \mathbf{Coalg}(\mathcal{P}_{\mathcal{S}})\; (\nu^{\mathrm{wf}}\mathcal{P}_{\mathcal{S}} \simeq \mu\mathcal{P}_{\mathcal{S}} \subsetneq \nu\mathcal{P}_{\mathcal{S}}), \\
    &\text{Extensional }\mathcal{P}_{\mathcal{S}}\text{ (subobjects) + Osius pure-set fullness}
  \end{aligned}}
\]
\begin{enumerate}[label=\textup{(\roman*)},leftmargin=1.8em]
  \item \textbf{Internal Double-Clamped $\mathsf{ZF}/\mathsf{ZFC}$ and $\mathsf{GB}/\mathsf{GBC}/\mathsf{NBG}$ Soundness:} Imposing the four Choice-free upper clamps ($\mathbb{N} \in \mathcal{S}$, $1 + 1 \xrightarrow{\sim} \Omega := \mathcal{P}_{\mathcal{S}}(1)$, $\mathcal{P}_{\mathcal{S}}(\mathcal{S}) \subseteq \mathcal{S}$, and Axis~VI image descent $\mathrm{im}(I \xrightarrow{f} C) \in \mathcal{S}$) together with the lower structural clamps (extensional subobject membership via $\mathcal{P}_{\mathcal{S}}$ and Axis~VII Foundation restricting admissible set-coalgebras along $\mathbf{Coalg}_{\mathrm{wf}}(\mathcal{P}_{\mathcal{S}}) \hookrightarrow \mathbf{Coalg}(\mathcal{P}_{\mathcal{S}})$ so $V = \nu^{\mathrm{wf}} \mathcal{P}_{\mathcal{S}} \simeq \mu \mathcal{P}_{\mathcal{S}}$) makes small objects $\mathcal{E}_{\mathcal{S}} \subseteq \mathcal{C}$ a Boolean elementary topos with $\mathrm{NNO}$ and forces $\mathcal{C} \models ((V, \in) \models \mathsf{ZF})$ and internal $\mathcal{C} \models ((V, \mathrm{Sub}_{\mathcal{C}}(V)) \models \mathsf{GB})$. Every $X \in \mathcal{E}_{\mathcal{S}}$ is locally a pure set in $V$ (admitting a regular epimorphism $p : U \twoheadrightarrow 1$ in $\mathcal{E}_{\mathcal{S}}$ and a monomorphism $p^* X \rightarrowtail p^* V$ in $\mathcal{C}/U$) if and only if $(\mathcal{C}, \mathcal{S})$ satisfies \emph{local Osius pure-set fullness} (every $X \in \mathcal{S}$ embeds over some $U \twoheadrightarrow 1$ into a fiberwise-extensional well-founded coalgebra $(T, R) \in \mathbf{Coalg}_{\mathrm{wf}}(\mathcal{P}_{\mathcal{S}})$; failing on the atom object $\mathbb{A}$ of Fraenkel--Mostowski urelement models $\mathcal{V}_{\mathrm{FM}} \models \mathsf{ZFA} + \neg\mathrm{AC}$). Adjoining internal small-map choice $\mathrm{IAC}_{\mathcal{S}}$ upgrades the internal theory to $\mathcal{C} \models ((V, \in) \models \mathsf{ZFC})$ and $\mathcal{C} \models ((V, \mathrm{Sub}_{\mathcal{C}}(V)) \models \mathsf{GBC})$ (and $\mathcal{C} \models ((V, \mathrm{Sub}_{\mathcal{C}}(V)) \models \mathsf{NBG})$ iff $\mathbf{Card}(V) \hookrightarrow \mathbf{Set}(V)$ splits in $\mathbf{Cat}(\mathcal{C})$) and discharges local Osius pure-set fullness automatically.
  \item \textbf{External Well-Pointed Pinch to $\mathbf{Set}_{\mathsf{ZFC}}$ and Global $\mathcal{E}_{\mathcal{S}} \simeq \mathbf{El}(V)$:} Under the Choice-free upper and lower clamps together with local Osius pure-set fullness (automatically discharged under $\mathrm{IAC}_{\mathcal{S}}$), $\mathcal{E}_{\mathcal{S}}$ pinches externally to a well-pointed classical category of pure sets ($\mathcal{E}_{\mathcal{S}} \simeq \mathbf{El}(V)$, modeling $\mathsf{ETCS}$ + Osius--Joyal--Moerdijk Replacement: $\mathbf{Set}_{\mathsf{ZF}}$ without Choice and $\mathbf{Set}_{\mathsf{ZFC}}$ with $\mathrm{IAC}_{\mathcal{S}}$) if and only if $\mathcal{C}$ additionally satisfies external two-valuedness ($\mathrm{Sub}_{\mathcal{C}}(1) \cong \{0, 1\}$) and external projectivity of $1$ in $\mathcal{E}_{\mathcal{S}}$ (splitting epic covers $U \twoheadrightarrow 1$ globally and upgrading $\mathrm{IAC}_{\mathcal{S}}$ to external $\mathrm{AC}_{\mathcal{S}}$; failing on atomless Boolean-valued $(\mathbf{Sh}_{\mathrm{Class}}(B), \mathbf{Sh}_{\mathrm{Set}}(B))$ without two-valuedness and on $(\mathbf{Class}^{\mathbb{Z}/2\mathbb{Z}}, \mathbf{Set}^{\mathbb{Z}/2\mathbb{Z}})$ without external projectivity of $1$).
  \item \textbf{Residuum Scale-Invariance:} Each of the upper compression and extensionality clamps above---Infinity ($\mathbb{N} \in \mathcal{S}$, Axis~III), Excluded Middle ($1 + 1 \xrightarrow{\sim} \Omega$, Axis~II), Power Set ($\mathcal{P}_{\mathcal{S}}(\mathcal{S}) \subseteq \mathcal{S}$, Axis~II), Replacement / image descent (Axis~VI), Choice ($\mathrm{IAC}_{\mathcal{S}}$, external projectivity of $1$, and large splitting $\mathbf{Card}(V) \hookrightarrow \mathbf{Set}(V)$, Axis~I), and Extensionality (subobject $\mathcal{P}_{\mathcal{S}}$ and \texttt{Quot.sound} bisimulation, Axis~VII)---coincides with the vanishing $\mathrm{Res}(\kappa_A) \simeq 0$ of an explicit admissible comparison triple $\mathfrak{T}_A = (\mathcal{D}_0 \supseteq \mathcal{D}_1, \kappa_A) \in \mathrm{Dom}(\mathrm{Res})$, while Axis~VII Foundation corresponds strictly to restricting admissible set-coalgebras to $\mathbf{Coalg}_{\mathrm{wf}}(\mathcal{P}_{\mathcal{S}}) \hookrightarrow \mathbf{Coalg}(\mathcal{P}_{\mathcal{S}})$---that is, vanishing of the factorization obstruction across $\kappa_{\mathrm{Reg}} : \nu^{\mathrm{wf}}\mathcal{P}_{\mathcal{S}} \simeq \mu\mathcal{P}_{\mathcal{S}} \hookrightarrow \nu\mathcal{P}_{\mathcal{S}}$ for every admissible set-coalgebra $(X, \gamma)$ (Theorem~\ref{thm:axis-vii-foundation}(iii)), rather than forcing $\nu\mathcal{P}_{\mathcal{S}} \setminus \mu\mathcal{P}_{\mathcal{S}} \simeq 0$---all calibrated in Sections~\ref{sec:conservativity-kaleidoscope}--\ref{sec:anti-existence-residua} (Table~\ref{tab:master-ast-ledger}).
\end{enumerate}
\end{theorem}

\begin{proof}
For~(i), any morphism $f : U \to X$ in $\mathcal{C}$ between small objects $U, X \in \mathcal{S}$ factors as $U \xrightarrow{\Gamma_f} U \times X \xrightarrow{\mathrm{pr}_2} X$ where $\Gamma_f = (f \times \mathrm{id}_X)^*(\Delta_X) \in \mathcal{S}$ and $\mathrm{pr}_2 = X^*(U \to 1) \in \mathcal{S}$~\cite[Lemma~I.1.2]{JoyalMoerdijk95}, so $f \in \mathcal{S}$ and $\mathrm{Sub}_{\mathcal{E}_{\mathcal{S}}}(X) = \mathrm{Sub}_{\mathcal{S}}(X)$. Since $\mathcal{E}_{\mathcal{S}}$ has finite limits and power objects $\mathcal{P}_{\mathcal{S}}(X) \in \mathcal{S}$ for all $X \in \mathcal{S}$, Mikkelsen--Par\'e's theorem~\cite[Thm.~A2.2.3 \& Cor.~A2.2.4]{Johnstone02} makes $\mathcal{E}_{\mathcal{S}}$ an elementary topos with $\mathrm{NNO}$ $\mathbb{N} \in \mathcal{S}$, while $1 + 1 \xrightarrow{\sim} \Omega := \mathcal{P}_{\mathcal{S}}(1)$ classifying all subobjects of $\mathcal{C}$ makes both $\mathcal{C}$ and $\mathcal{E}_{\mathcal{S}}$ Boolean ($\forall_\pi = \neg\exists_\pi\neg$ along $Z \times V \to Z$). Under Axis~VI covariant direct image $\exists_g$, Lambek's isomorphism $\mathrm{sup} : \mathcal{P}_{\mathcal{S}}(V) \xrightarrow{\sim} V$ enforces Extensionality on $V = \nu^{\mathrm{wf}}\mathcal{P}_{\mathcal{S}} \simeq \mu\mathcal{P}_{\mathcal{S}}$ while Theorem~\ref{thm:axis-vii-foundation}(i) enforces Foundation ($\in$-induction). Composition closure of $\mathcal{S}$ together with Axis~VI image descent yields Pairing and Small Union (since $\pi_2 : \in_V \to V$ lies in $\mathcal{S}$, with fiber $\pi_2^{-1}(y) = \{x \mid x \in y\}$), $\mathcal{P}_{\mathcal{S}}(\mathcal{S}) \subseteq \mathcal{S}$ yields Power Set~\cite[Thm.~I.2.1]{JoyalMoerdijk95}, and $\mathbb{N} \in \mathcal{S}$ with strict order $\gamma_<(n) = \{k < n\}$ collapses via Osius terminality~\cite{Osius74} into $\omega \in V$ (Infinity). Boolean regular class operations (\textbf{G1--G8}) and pullback stability of $\mathcal{S}$ give Subclass Separation (hence unbounded $\mathcal{L}_\in$-Separation on $V$), while Axis~VI image descent yields Class Replacement, proving $\mathcal{C} \models ((V, \in) \models \mathsf{ZF})$ and $\mathcal{C} \models ((V, \mathrm{Sub}_{\mathcal{C}}(V)) \models \mathsf{GB})$. For Choice-free pure-set fullness, given any fiberwise-extensional well-founded coalgebra $(T \xrightarrow{\pi} U, R \rightarrowtail T \times_U T)$ in $\mathcal{C}/U$ ($T \in \mathcal{S}$), viewing $R \rightarrowtail T \times T$ as $(T, \tilde{\gamma}_T) \in \mathbf{Coalg}_{\mathrm{wf}}(\mathcal{P}_{\mathcal{S}})$ yields by Osius terminality a unique $m : T \to V$ with $\mathrm{sup}^{-1} \circ m = \exists_m \circ \tilde{\gamma}_T$; by $\in_R$-induction on $t \in T$, $\pi(t) = \pi(t') \wedge m(t) = m(t')$ forces $\tilde{\gamma}_T(t) = \tilde{\gamma}_T(t')$ and hence $t = t'$ by fiberwise extensionality, making $(\pi, m) : T \rightarrowtail U \times V = p^* V$ monic in $\mathcal{C}/U$. Conversely, any monomorphism $p^* X \rightarrowtail p^* V$ over $U$ generates a small fiberwise-transitive closure $T := \bigcup_{n \in \mathbb{N}} \bigcup^n_U \mathrm{im}(p^* X) \rightarrowtail U \times V$ via $\mathbb{N} \in \mathcal{S}$, Small Union, and Axis~VI Replacement, which restricts $\mathrm{sup}^{-1}$ to a fiberwise-extensional well-founded coalgebra containing $p^* X$. Adjoining $\mathrm{IAC}_{\mathcal{S}}$ yields $\mathcal{C} \models ((V, \in) \models \mathsf{ZFC})$ and $\mathcal{C} \models ((V, \mathrm{Sub}_{\mathcal{C}}(V)) \models \mathsf{GBC})$ (and $\mathcal{C} \models ((V, \mathrm{Sub}_{\mathcal{C}}(V)) \models \mathsf{NBG})$ iff $\mathbf{Card}(V) \hookrightarrow \mathbf{Set}(V)$ splits, Proposition~\ref{prop:large-equiv-splitting}), and by Zermelo's theorem in Boolean $\mathcal{E}_{\mathcal{S}} \models \mathrm{IAC}$~\cite[Prop.~D4.5.7]{Johnstone02} every $X \in \mathcal{S}$ carries a $U$-indexed well-ordering $<$ over its inhabited cover $U := \mathrm{WO}(X) \twoheadrightarrow 1$ in $\mathcal{E}_{\mathcal{S}}$, whose strict lower segments $\gamma_<(u, x) = \{(u, y) \mid y <_u x\}$ are fiberwise extensional and well-founded, automatically discharging local Osius pure-set fullness.

For~(ii), a non-degenerate Boolean elementary topos $\mathcal{E}_{\mathcal{S}}$ is well-pointed ($1$ generates $\mathcal{E}_{\mathcal{S}}$) if and only if $\mathrm{Sub}_{\mathcal{C}}(1) = \mathrm{Sub}_{\mathcal{E}_{\mathcal{S}}}(1) \cong \{0, 1\}$ and $1$ is externally projective in $\mathcal{E}_{\mathcal{S}}$~\cite[Prop.~VI.10.1]{MacLaneMoerdijk92} (for $f \neq g : X \rightrightarrows Y$, $X \setminus \mathrm{eq}(f, g) \not\cong 0$ has epic support onto $1$ by two-valuedness and hence a separating global point $1 \to X \setminus \mathrm{eq}(f, g)$ by projectivity of $1$). External projectivity of $1$ in $\mathcal{E}_{\mathcal{S}}$ splits every epic cover $p : U \twoheadrightarrow 1$ via a global point $s : 1 \to U$, pulling local pure-set embeddings $p^* X \rightarrowtail p^* V$ back to global monomorphisms $m_s : X \rightarrowtail V$ classified by $a_X := \mathrm{sup}(\ulcorner m_s \urcorner) : 1 \to V$ ($X \cong E_{a_X}$, establishing global $\mathcal{E}_{\mathcal{S}} \simeq \mathbf{El}(V)$) and upgrading $\mathrm{IAC}_{\mathcal{S}}$ to external $\mathrm{AC}_{\mathcal{S}}$ (Chapter~\ref{ch:choice-atlas}, Theorem~\ref{thm:iac-to-ac}), yielding Lawvere's $\mathsf{ETCS}$ plus Osius--Joyal--Moerdijk Replacement~\cite{Cole73,Mitchell72,Osius74}. Conversely, atomless $(\mathbf{Sh}_{\mathrm{Class}}(B), \mathbf{Sh}_{\mathrm{Set}}(B))$ models Boolean-valued $V^{(B)} \models \mathsf{ZFC}$ with $\mathrm{Sub}_{\mathcal{C}}(1) \cong B \not\cong \{0, 1\}$, while $(\mathbf{Class}^{\mathbb{Z}/2\mathbb{Z}}, \mathbf{Set}^{\mathbb{Z}/2\mathbb{Z}})$ satisfies two-valuedness and $\mathrm{IAC}_{\mathcal{S}}$ yet fails external projectivity of $1$ (Chapter~\ref{ch:choice-atlas}, Theorem~\ref{thm:coshep-svc}, leaving $V = \mu\mathcal{P}_{\mathcal{S}}$ with trivial $\mathbb{Z}/2\mathbb{Z}$-action so $X = (\{0, 1\}, \mathrm{swap}) \notin \mathbf{El}(V)$ globally). Finally, (iii) collects the explicit comparison triples calibrated across Sections~\ref{sec:conservativity-kaleidoscope}--\ref{sec:anti-existence-residua} (Table~\ref{tab:master-ast-ledger}).
\end{proof}

Classical $\mathbf{Set}_{\mathsf{ZFC}}$ is therefore neither minimal nor canonical; inside a Joyal--Moerdijk small-map ambient $(\mathcal{C}, \mathcal{S})$ it is a \textbf{double-clamped locus} carved out by a clean \textbf{duality of axiomatic excess}:
\begin{enumerate}[label=\textup{\arabic*.},leftmargin=1.8em]
  \item \textbf{Over-Compressing Existence and Logic (\emph{Infinity, Excluded Middle, Power Set, Replacement, Choice}):} Forcing infinite initial algebras $\mathbb{N}$, classical truth values $1 + 1 \xrightarrow{\sim} \Omega$, powerclasses $\mathcal{P}_{\mathcal{S}}(X)$, small-indexed images and countable colimits $\bigcup_{n \in \mathbb{N}} \mathcal{P}^n(\mathbb{N})$, and epimorphism sections back down into $\mathcal{S}$ (plus external two-valuedness and projectivity of $1$ in $\mathcal{E}_{\mathcal{S}}$ for well-pointed $\mathcal{E}_{\mathcal{S}} \simeq \mathbf{El}(V)$, and $\mathbf{Card}(V) \hookrightarrow \mathbf{Set}(V)$ for $\mathsf{NBG}$).
  \item \textbf{Over-Quotienting Structure (\emph{Foundation, Extensionality, Pure-Set Fullness}):} Restricting admissible set-coalgebras along $\mathbf{Coalg}_{\mathrm{wf}}(\mathcal{P}_{\mathcal{S}}) \hookrightarrow \mathbf{Coalg}(\mathcal{P}_{\mathcal{S}})$ ($\nu^{\mathrm{wf}} \mathcal{P}_{\mathcal{S}} \simeq \mu \mathcal{P}_{\mathcal{S}} \subsetneq \nu \mathcal{P}_{\mathcal{S}}$) and enforcing extensional subobjects plus Osius pure-set fullness. In the univalent setting of Proposition~\ref{prop:extensionality-truncation}(ii), inside an $(\infty, 1)$-topos $\mathcal{X}$ with nested univalent universes $\mathcal{U}_0 \in \mathcal{U}_1$~\cite[\S 6.1.6]{Lurie09}, Gylterud's hereditary fiberwise $(-1)$-truncation $V_0 \hookrightarrow M = \mu \mathcal{P}_{\mathrm{Anima}}$~\cite[Thm.~3.9]{Gylterud18} directly realizes $V_0 \cong \mu \mathcal{P}_{\mathcal{S}}$ inside $(\tau_{\le 0}\mathcal{X}_{\mathcal{U}_1}, \mathcal{S}_{\mathcal{U}_0})$ (via $\sum_{A : \mathcal{U}_0} (A \hookrightarrow X) \simeq \mathcal{P}_{\mathcal{S}}(X)$), while developing a universe-free elementary higher class-anima AST $(\mathcal{C}_\infty, \mathcal{S}_\infty)$ forms a central programmatic synthesis horizon for Volume~II.
\end{enumerate}

\subsection{\texorpdfstring{The Kernel Audit Gate: \texttt{HomotopyCRM.SetsResiduum}}{The Kernel Audit Gate: HomotopyCRM.SetsResiduum}}\label{subsec:sets-residuum-kernel-gate}

This double-clamped geometry\index{SetsResiduum@\texttt{SetsResiduum} Lean~4 module} maps directly onto the \texttt{\#print axioms} kernel audit gate enforced throughout our Lean~4 formalization suite \texttt{HomotopyCRM} (\texttt{leanprover/lean4:v4.28.0} over \texttt{mathlib4}). Inside Lean's Calculus of Inductive Constructions, universe-polymorphic \texttt{Classical.choice} realizes uniform global equivalence-splitting across cumulative universes \texttt{Type u}, while propositional extensionality (\texttt{propext}) and quotient soundness (\texttt{Quot.sound}, from which Lean derives function extensionality \texttt{funext}) implement $0$-truncated extensionality and recursive bisimulation quotienting---collapsing intensional tree presentations (\texttt{PSet}, modeling $\mu \mathcal{P}_{\mathrm{Anima}}$ at $0$-types) into extensional sets (\texttt{ZFSet}).

To separate all three kernel primitives cleanly under the monograph's \textbf{4-Layer Definition Check Gate} (matching the \texttt{\#guard\_msgs in \#print axioms} standard of the earlier formalization modules), our companion module \texttt{HomotopyCRM/\allowbreak SetsResiduum.lean} pairs its bundled two-sorted \texttt{ClassSetModel} adjunction (\texttt{defSetProjAdjunction}, \texttt{finToPowElemHom : SetElemHom}, \texttt{finToPowClassHom : ClassSetHom}) and $8$-part pullback counterexample with three deliberate design guardrails (Table~\ref{tab:sets-residuum-audit}). First, \texttt{card\_\allowbreak split\_\allowbreak iff\_\allowbreak global\_\allowbreak choice} avoids hidden \texttt{funext} (\texttt{Quot.sound}) dependencies by using pointwise homotopy-functor equalities (\texttt{forall x, ...}) over an abstract skeletal carrier, paired with \texttt{type\_\allowbreak equiv\_\allowbreak splitting\_\allowbreak global\_\allowbreak choice} on \texttt{CategoryTheory.Functor} over \texttt{Type u}. Second, because every term of Lean's inductive \texttt{PSet} is well-founded, \texttt{quine\_\allowbreak atom\_\allowbreak outside\_\allowbreak wf} embeds \texttt{PSet} into a pointed-graph coalgebra carrier \texttt{APG} and uses \texttt{PSet.rec} (\texttt{axiom}-free) to refute bisimilarity between the self-loop coalgebra \texttt{quineAPG} ($q = \{q\}$) and any \texttt{x : PSet}. Third, \texttt{zfset\_\allowbreak ext\_\allowbreak quot\_\allowbreak gate} invokes \texttt{Quotient.sound} directly over \texttt{PSet.Equiv} so its kernel footprint stays strictly \texttt{propext, Quot.sound} without pulling \texttt{Classical.choice}.

\begin{table}[htbp]
\centering
\caption{Kernel Axiom Audit Ledger for \texttt{HomotopyCRM.SetsResiduum}}\label{tab:sets-residuum-audit}
\vspace{0.3em}
{\scriptsize\setlength{\tabcolsep}{2.5pt}\renewcommand{\arraystretch}{1.12}%
\begin{tabularx}{\textwidth}{@{}P{0.27\textwidth}YP{0.21\textwidth}@{}}
\toprule
\textbf{Lean 4 Declaration} & \textbf{Calibrated Universe Clamp / Residuum} & \textbf{\texttt{\#print axioms} Gate} \\
\midrule
\texttt{defSetProjAdjunction}, \texttt{finToPowElemHom}, \texttt{finToPowClassHom} \& \texttt{naive\_\allowbreak pullback\_\allowbreak breaks\_\allowbreak subclass\_\allowbreak separation} & Definition~\ref{def:set-projection-fibration}: two-sorted \texttt{ClassSetModel} coreflection $\mathrm{Def} \dashv \pi_{\mathrm{Set}}$, elementary embedding \texttt{finToPowElemHom : SetElemHom} (\texttt{elem\_ext}), and $8$-part pullback Subclass-Separation counterexample (\texttt{finOmegaModel} $\hookrightarrow$ \texttt{powOmegaModel}, \texttt{evenPowClass}) & \texttt{propext}, \texttt{Classical.choice}, \texttt{Quot.sound} \\
\texttt{card\_\allowbreak split\_\allowbreak iff\_\allowbreak global\_\allowbreak choice} & Proposition~\ref{prop:large-equiv-splitting}: pointwise homotopy-functor splitting of skeletal $\mathbf{Card} \hookrightarrow \mathbf{Set}$ is equivalent to uniform small-set choice $G(X)=\eta_X^{-1}(\emptyset)$ & \texttt{(axiom-free)} \\
\texttt{type\_\allowbreak equiv\_\allowbreak splitting\_\allowbreak global\_\allowbreak choice} & Proposition~\ref{prop:large-equiv-splitting} + \texttt{CategoryTheory.Functor}: global functorial skeleton splitting of \texttt{Type u} via \texttt{Classical.choice} & \texttt{propext}, \texttt{Classical.choice}, \texttt{Quot.sound} \\
\texttt{pset\_\allowbreak equiv\_\allowbreak iff\_\allowbreak bisimilar} \& \texttt{quine\_\allowbreak atom\_\allowbreak outside\_\allowbreak wf} & Theorem~\ref{thm:axis-vii-foundation}: \texttt{PSet.Equiv} embeds faithfully into \texttt{APG} bisimulation while \texttt{quineAPG} ($q=\{q\}$) lies outside the image of \texttt{PSet} & \texttt{(axiom-free)} \\
\texttt{quineAPG\_\allowbreak bisimilar\_\allowbreak and\_\allowbreak nonIso\_\allowbreak twoCycle}, \texttt{boffa\_\allowbreak twin\_\allowbreak torsor} (\& \texttt{boffa\_\allowbreak twin\_\allowbreak addTorsor}) & Proposition~\ref{prop:boffa-torsors}: $2$-cycle non-isomorphic bisimulation witness and ordered vs.\ unordered twin Quine swap torsor (\texttt{Bool} core vs.\ Layer-3 \texttt{AddTorsor (ZMod 2)}) & \texttt{(axiom-free)} (Layer~3: \texttt{propext}, \texttt{Quot.sound}) \\
\texttt{zfset\_\allowbreak ext\_\allowbreak quot\_\allowbreak gate} & Proposition~\ref{prop:extensionality-truncation}: \texttt{Quotient.sound} over \texttt{PSet.Equiv} collapsing branch permutations and multiplicity $[\emptyset,\emptyset] \sim [\emptyset]$ & \texttt{propext}, \texttt{Quot.sound} \\
\bottomrule
\end{tabularx}\par}
\end{table}

\subsection{\texorpdfstring{Bridge to Chapters~\ref{ch:panorama}--\ref{ch:master-atlas} and Volume~II}{Bridge to Chapters 8--9 and Volume II}}\label{subsec:sets-bridge-ch8}

Above all, dissecting ``$\mathbf{Sets}$'' closes the constructive and topos-theoretic arc of Volume~I by naming what lay just beyond its elementary-topos horizon. Recording these universe clamps prepares the programmatic panorama and Master Reference Atlas of Chapters~\ref{ch:panorama}--\ref{ch:master-atlas} and defines the four-part research program of Volume~II (\emph{Universes, Classes, and Higher Fixed Points}):
\begin{itemize}[leftmargin=1.5em]
  \item \textbf{Axis~VI (Quotient-Image Descent and Borel Proof Unwinding):} Calibrates small-indexed image descent separating $\mathsf{ETCS} \simeq \mathsf{Mac}$ and $\mathsf{ZC}$ ($V_{\omega+\omega}$) from $\mathsf{ZFC}$, extracting the explicit strategy-descent residuum $\mathrm{Res}(\kappa_{\mathrm{Det}})$ on Martin's proper-class unraveling coalgebra in $\mathcal{C} \setminus \mathcal{S}$ at $\boldsymbol{\Sigma}^0_\omega$~\cite{Friedman71,Martin75} (Remark~\ref{rem:martin-borel-coalgebra}).
  \item \textbf{Axis~VII (Coalgebraic and Univalent Universe Fixed Points):} Develops the full spectrum between $V_{\mathsf{ZF}} = \mu \mathcal{P}_{\mathcal{S}}$, Aczel's $V_{\mathsf{AFA}} = \nu \mathcal{P}_{\mathcal{S}}$, Boffa's Rieger--Bernays torsors over $V_{\mathsf{BAFA}}$~\cite{Bernays54,Rieger57,Boffa69}, Gylterud's $M = \mu \mathcal{P}_{\mathrm{Anima}}$~\cite{Gylterud18}, and the terminal univalent multiset coalgebra $\nu \mathcal{P}_{\mathrm{Anima}}$ where $\Omega(\nu \mathcal{P}_{\mathrm{Anima}}, [q, q]) \cong \mathbb{Z}/2\mathbb{Z}$ unites Axes~V and~VII (Remark~\ref{rem:nu-anima-fossil}).
  \item \textbf{$2$-Categorical Calibration of the $(\infty, \infty)$-Ladder Without Universes:} Replaces Grothendieck universes ($V_\kappa$) across Chapter~\ref{ch:higher-crm}'s four-tier hierarchy $\mathrm{Cat}^L_{(\infty,\infty)} \subsetneq \mathrm{Cat}^\omega_{(\infty,\infty)} \subsetneq \mathrm{Cat}^{\mathrm{coind}}_{(\infty,\infty)} \subsetneq \mathrm{Cat}^R_{(\infty,\infty)}$~\cite{ORW26} with exact class calibrations across $\mathsf{GBC} \subsetneq \mathsf{NBG} \subsetneq \mathsf{NBG}+\mathsf{ETR}_{\mathrm{Ord}} \subsetneq \mathsf{NBG}+\mathsf{ETR} \subsetneq \mathsf{KM}$~\cite[Thm.~1.2]{Gitman20} (noting that Gitman et al.\ write ``$\mathsf{GBC}$'' for $\mathsf{NBG} = \mathsf{GBC} + \mathsf{GC}$, and distinguishing length-$\mathrm{Ord}$ $\mathsf{ETR}_{\mathrm{Ord}}$ from full $\mathsf{ETR}$ along arbitrary class well-orderings)---testing where $\alpha = 0$ stabilization on $\mathrm{Cat}_{\le \omega}$ (Theorem~\ref{thm:ordinal-rank-completeness}, including $\mathcal{Q}$) drops transfinite $\mathsf{ETR}_{\mathrm{Ord}}$ recursion down to $\mathsf{GBC}$, separating the $\mathrm{Con}(\mathsf{ZFC})$ threshold crossed already at $\mathsf{ETR}_\omega < \mathsf{ETR}_{\mathrm{Ord}}$ from full $\mathsf{ETR}$ and impredicative $\mathsf{KM}$ comprehension ($\mathrm{Con}(\mathsf{NBG}+\mathsf{ETR}_{\mathrm{Ord}})$).
  \item \textbf{Non-Conservative Base-Change Transgression:} Formulates a quantitative Leray-style transgression map along $\mathrm{Def} \dashv \pi_{\mathrm{Set}}$ (mirroring Chapter~\ref{ch:base-fiber}, Theorem~\ref{thm:quantitative-conservation}); because raw model-space pullbacks fail (\texttt{naive\_pullback\_breaks\_subclass\_separation} on \texttt{ClassSetModel}, mirrored on $M \prec V_\kappa$ in Definition~\ref{def:set-projection-fibration}(i)), a rigorous transgression carrier operates via Macro $\Sigma^1_0$-Morleyized classifying topoi $\gamma_I : \mathcal{E}[T_{\mathrm{Class}}] \to \mathcal{E}[T_{\mathrm{Set}}]$ coupled to Micro internal AST class-forcing sheaves $\gamma_\Sigma$, separating predicative $\mathcal{S}$-distributivity leaks (Glazer over $\mathsf{ZC}$~\cite{Glazer23}, Felgner over $\mathsf{ZFC}_{\neg\mathrm{Reg}}$~\cite{Felgner76}) from excision/reflection leaks ($H^0_c(U_{\mathrm{leak}}; \mathbb{Z})$ for $\mathsf{ETR}_{\mathrm{Ord}}$ and $\mathsf{KM}$), while vanishing under Replacement + Foundation over $\mathsf{ZFC}$.
\end{itemize}

The scare quotes around ``$\mathbf{Sets}$'' may now come off---not because a single monolithic foundation has prevailed, but because its hidden joints lie fully in view. What once appeared to be an unanalyzable meta-theoretic floor is simply the double-clamped locus inside Joyal--Moerdijk ambients $(\mathcal{C}, \mathcal{S})$---bridging to univalent $(\infty, 1)$-topos universes $(\mathcal{X}, \mathcal{U})$ and the Volume~II horizon of elementary class-anima ambients $(\mathcal{C}_\infty, \mathcal{S}_\infty)$---where existence-compression from above meets hereditary symmetry-truncation from below, and where every turn of the class kaleidoscope overhead obeys the exact same laws of necessity that govern the mathematics within.
}

\part{Programmatic Panorama, Open Horizons, and Formal Verification}

\chapter{Programmatic Panorama: External Benchmark Lattices Across Analysis, Algebra, Homotopy, and Cohesion}
\chaptermark{Programmatic Panorama: External Benchmark Lattices}
\label{ch:panorama}\label{ch:non-ac-lattices}
\chapterepigraph{``The introduction of numbers as coordinates by reference to the particular division scheme of the open one-dimensional continuum is an act of violence\dots''}{Hermann Weyl, \emph{Philosophy of Mathematics and Natural Science} (1949)}

\section{The Internal Continuum and Pointfree Spatiality Benchmarks}\label{sec:panorama-continuum-spatial}

\begin{remark}[Genre Signpost: External Benchmark Syntheses vs.\ the Proved Spine of Chapters~1--7]\label{rem:panorama-genre-signpost}%
\uses{def:comparison-triple,def:residuum}%
Unlike Chapters~\ref{ch:foundations}--\ref{ch:sets-residuum}---which establish self-contained proofs and Lean~4 kernel verifications for our novel categorical residuums across choice, $\delta$-rings, condensed smoothings, $(\infty,\infty)$-ranks, Fra\"iss\'e combinatorics, and algebraic set theory---this chapter serves explicitly as a \emph{programmatic synthesis panorama}. Every theorem and bifurcation lattice below translates established external theorems from constructive analysis, pointfree topology, non-abelian homological algebra, arithmetic pretopoi, synthetic homotopy theory, and differential cohesion into the unified comparison-triple coordinates of $\mathrm{Dom}(\mathrm{Res})$, providing calibrated external benchmarks for the open problems of Chapter~\ref{ch:master-atlas}.
\end{remark}

\subsection{The Internal Continuum and Real Analysis Lattice}

While the Axiom of Choice ($\mathrm{AC}$) provides the archetypal testing ground for Categorical Reverse Mathematics ($\mathrm{CRM}$), the residuum equivalence $\mathcal{D}_0 + \mathrm{Res}(\kappa) \iff \mathcal{D}_1$ applies across diverse mathematical structures whenever a classical equivalence bifurcates inside general categories or topoi. A primary example of this bifurcation occurs in the internal construction of the real line inside an elementary topos $\mathcal{E}$ equipped with a Natural Numbers Object ($\mathrm{NNO}$)~\cite{Johnstone82,Johnstone02}.

In classical mathematics over $\mathbf{Set}$ (with Countable Choice $\mathrm{CC}$ and the Law of the Excluded Middle $\mathrm{LEM}$), all standard completions of the rational numbers $\mathbb{Q}$ yield isomorphic fields. In an arbitrary elementary topos $\mathcal{E}$ with an $\mathrm{NNO}$, however, the continuum splits into a canonical chain of comparison morphisms:

\[
  \mathbb{Q}_{\mathcal{E}} \rightarrowtail \mathbb{R}_C \rightarrowtail \mathbb{R}_D \longrightarrow \mathbb{R}_M
\]

Here $\mathbb{R}_C$ denotes the \emph{Cauchy reals} (quotients of modulated Cauchy sequences of rationals), $\mathbb{R}_D$ denotes the \emph{Dedekind reals} (pairs of open, located, rounded lower and upper cuts $(L, U) \subseteq \mathbb{Q} \times \mathbb{Q}$), and $\mathbb{R}_M$ denotes the \emph{MacNeille reals} (the Dedekind-complete lattice order completion of $\mathbb{Q}$, which exists in any Boolean topos even when $\mathbb{R}_D$ is degenerate). Crucially, unlike $\mathbb{Q}_{\mathcal{E}} \rightarrowtail \mathbb{R}_C \rightarrowtail \mathbb{R}_D$, the canonical comparison map $\mathbb{R}_D \longrightarrow \mathbb{R}_M$ is a \emph{monomorphism} ($\mathbb{R}_D \rightarrowtail \mathbb{R}_M$, and indeed an isomorphism under bounded completeness) if and only if $\mathcal{E}$ satisfies the Weak Law of Excluded Middle ($\mathrm{WLEM}$: $\neg\phi \vee \neg\neg\phi = \top$, i.e., $\Omega_{\mathcal{E}}$ is a De Morgan / Stone lattice; see Theorem~\ref{thm:omega-dictionary})~\cite[\S D4.7, Thms.~D4.7.4--D4.7.12]{Johnstone02} (cf.~\cite[Ch.~III]{Johnstone82}).

Calibrating classical real analysis via the categorical residuum $(A \Rightarrow_{\mathbf{Topos}} B)$ isolates three fundamental geometric bridges:

\begin{enumerate}[label=\textbf{\arabic*.},leftmargin=1.8em]

  \item \textbf{The Cauchy-to-Dedekind Continuum Collapse ($\mathbb{R}_C \Rightarrow \mathbb{R}_D$).} In any topos $\mathcal{E}$ satisfying Countable Choice ($\mathrm{CC}$), one can extract a modulated rational Cauchy sequence from the locatedness of any Dedekind cut, forcing $\mathbb{R}_C \simeq \mathbb{R}_D$. However, $\mathrm{CC}$ is a choice principle rather than a pure geometric property of a site. If we forbid $\mathrm{CC}$ and ask for the purely topological residuum $\mathcal{R}(\mathcal{E})$ over a Grothendieck sheaf topos $\mathcal{E} \simeq \mathbf{Sh}(\mathbb{C}, J)$ with an $\mathrm{NNO}$ such that the canonical inclusion $\mathbb{R}_C \rightarrowtail \mathbb{R}_D$ is an isomorphism, $\mathrm{CRM}$ isolates $\omega$-completeness (open-cover refinement) of the site:

\[
  \top + \bigl(\text{Countable open-cover refinement of } (\mathbb{C}, J)\bigr) \iff \mathbb{R}_C \simeq \mathbb{R}_D
\]

Specifically, when every countable family of open coverings in the underlying locale admits a common refinement (as in sheaves over Baire or $\sigma$-compact zero-dimensional spaces), locally selected rational approximations glue into global Cauchy sequences without invoking $\mathrm{CC}$.

  \item \textbf{Approximate versus Exact Intermediate Value Theorem ($\varepsilon\text{-}\mathrm{IVT} \Rightarrow \mathrm{IVT}$).} Over Bishop-style constructive mathematics ($\mathrm{BISH}$)~\cite{Bishop67,Ishihara06} and in any elementary topos with an $\mathrm{NNO}$, the $\varepsilon$-Approximate Intermediate Value Theorem ($\varepsilon\text{-}\mathrm{IVT}$) holds unconditionally: for every uniformly continuous map $f: [0,1]_{\mathbb{R}_D} \to \mathbb{R}_D$ with $f(0) < 0 < f(1)$ and every $n \in \mathbb{N}$, bisection yields $x \in [0,1]$ with $|f(x)| < 2^{-n}$. Yet the \emph{Exact Intermediate Value Theorem} ($\mathrm{IVT}$: $\exists x \in [0,1],\, f(x) = 0$) fails in general sheaf topoi $\mathbf{Sh}(X)$, where continuously varying functions can cross zero on nowhere-dense boundaries without admitting a continuous root section~\cite[\S D4.7]{Johnstone02}. Computing the categorical residuum yields:

\[
  \varepsilon\text{-}\mathrm{IVT} + \mathcal{R}(\mathcal{E}) \iff \mathrm{IVT}
\]

where $\mathcal{R}(\mathcal{E})$ calibrates the exactness of Dedekind cuts against the \emph{zero-dimensionality or local connectedness} of the ambient site (or, internally, the Weak Law of Excluded Middle $\mathrm{WLEM}$ combined with countable cut-completeness).

  \item \textbf{Sequential Compactness to the Extreme Value Theorem.} In classical reverse mathematics over $\mathrm{RCA}_0$, the Extreme Value Theorem and Uniform Continuity of continuous functions on $[0,1]$ reverse to Weak K\"onig's Lemma ($\mathrm{WKL}_0$)~\cite[Ch.~IV]{Simpson09}. In $\mathrm{CRM}$, bridging sequential compactness ($A$) to the Extreme Value Theorem and Uniform Continuity on $[0,1]_{\mathcal{E}}$ ($B$) isolates the \emph{internal compactness of the Cantor locale $2^{\mathbb{N}}_{\mathcal{E}}$} inside $\mathcal{E}$, equivalent to the Fan Theorem over the internal binary tree~\cite[Chs.~III--IV]{Johnstone82} (cf.~\cite{Bishop67,Ishihara06}). (By contrast, the Baire locale $\mathbb{N}^{\mathbb{N}}_{\mathcal{E}}$ is \emph{never} compact in any non-trivial topos~\cite[Chs.~III--IV]{Johnstone82}.)

\end{enumerate}

\subsection{The Pointfree versus Spatial Duality Lattice}

A central result of topos theory is that many classical theorems of topology and functional analysis that require $\mathrm{AC}$ or the Boolean Prime Ideal Theorem ($\mathrm{BPI}$) for point-set topological spaces become \emph{completely constructive and choice-free} when formulated for pointfree locales (frames) inside an arbitrary topos $\mathcal{E}$~\cite{Johnstone82}.

A \emph{frame} $L$ in a topos $\mathcal{E}$ is a complete lattice satisfying the infinite distributive law $a \wedge \bigvee_{i \in I} b_i = \bigvee_{i \in I} (a \wedge b_i)$; the opposite category $\mathbf{Loc}(\mathcal{E}) = \mathbf{Frm}(\mathcal{E})^{\mathrm{op}}$ is the category of \emph{internal locales}. A point of a locale $L$ is a frame homomorphism $p^*: L \to \Omega_{\mathcal{E}}$, and $L$ is \emph{spatial} (has enough points) if distinct elements of $L$ are separated by points.

When we compare a pointfree localic theorem ($A$) to its classical spatial counterpart ($B$), $\mathrm{CRM}$ factors out set-theoretic choice and isolates \emph{internal spatiality} as the exact geometric residuum:

\begin{itemize}[leftmargin=1.5em]

  \item \textbf{Localic versus Spatial Tychonoff Theorem:} Peter Johnstone~\cite{Johnstone81} (\cite[Ch.~III, \S 1]{Johnstone82}) proved that the arbitrary product of compact locales $\prod_{i \in I} L_i$ is compact in \emph{any} elementary topos $\mathcal{E}$ without using $\mathrm{AC}$ or $\mathrm{BPI}$. Taking $A$ to be Johnstone's localic Tychonoff theorem and $B$ to be the classical point-set Tychonoff theorem for topological spaces in $\mathcal{E}$, the residuum $\mathcal{R}(\mathcal{E})$ is precisely the \emph{spatiality of compact product locales} (which, by Deligne's theorem~\cite[Exp.~VI, Appendice]{Deligne72}, is equivalent to internal $\mathrm{BPI}$ / logical compactness of $\Omega$).

  \item \textbf{Localic versus Spatial Hahn--Banach Theorem:} Christopher Mulvey and Joan Wick Pelletier~\cite{MulveyPelletier91} proved that for any normed space $E$ inside a Grothendieck topos $\mathcal{E}$, the \emph{localic Hahn--Banach theorem} ($A$) holds constructively: the locale of norm-preserving linear extensions of a functional on a subspace $M \subseteq E$ is always compact and non-trivial. Elevating $A$ to the classical point-set Hahn--Banach theorem ($B$: the existence of an actual global linear functional $f: E \to \mathbb{R}_D$) requires the residuum $\mathcal{R}(\mathcal{E})$ that compact regular locales in $\mathcal{E}$ possess global points.

\end{itemize}

Crucially, by the \emph{Base-Fiber Conservation Principle} (Chapter~\ref{ch:foundations}, Section~\ref{sec:base-fiber}, Principle~\ref{prin:conservation}), if we evaluate this spatiality residuum in a relative Grothendieck topos $\gamma: \mathcal{E} \to \mathcal{S}$ and \emph{remove} $\mathrm{BPI}$ from the metamathematical base topos $\mathcal{S}$, finitary coherence (every cover has a finite subcover) is no longer strong enough to guarantee points. The lost ultrafilter strength in $\mathcal{S}$ reappears in the fiber residuum $\mathcal{R}(\mathcal{E}/\mathcal{S})$, which must strengthen from finitary coherence to \emph{supercompactness}---the topological condition that every covering refines to a single generating element, producing points without base ultrafilters (proved in Chapter~\ref{ch:base-fiber}, Theorem~\ref{thm:deligne-deformed}).

\section{\texorpdfstring{$\Omega$-Geometry for Intermediate Logics and the Homological Exactness Lattice}{Omega-Geometry for Intermediate Logics and the Homological Exactness Lattice}}\label{sec:panorama-omega-homological}

\subsection{The Intermediate Logic and Subobject Classifier Benchmark}

In Constructive Reverse Mathematics over $\mathrm{BISH}$~\cite{Bishop67,Ishihara06}, numerous theorems of analysis and commutative algebra are calibrated by semi-classical logical principles lying strictly between Intuitionistic Propositional Calculus ($\mathrm{IPC}$) and the full Law of the Excluded Middle ($\mathrm{LEM}$). Categorical Reverse Mathematics translates this logical hierarchy directly into the order-algebraic geometry of the subobject classifier $\Omega$ and the topology of underlying sites, synthesizing the standard constructive-geometry correspondences due to Johnstone~\cite{Johnstone82,Johnstone02} and Mac~Lane--Moerdijk~\cite{MacLaneMoerdijk92}:

\begin{theorem}[The $\Omega$-Geometry Dictionary for Intermediate Logics]\label{thm:omega-dictionary}%
\uses{def:comparison-triple,def:residuum}%

Let $\mathcal{E}$ be an elementary topos with subobject classifier $\Omega$, and in the spatial case let $\mathcal{E} = \mathbf{Sh}(X)$ be sheaves on a topological space $X$. Then the residuums of semi-classical logical principles correspond to the following geometric invariants:

\begin{itemize}[leftmargin=1.5em]

  \item \textbf{De Morgan's Law ($\mathrm{WLEM}$:} $\neg\phi \vee \neg\neg\phi = \top$): Holds internally in $\mathcal{E}$ if and only if $\Omega_{\mathcal{E}}$ is an internal Stone (De Morgan) lattice (a Heyting algebra satisfying $\neg\phi \vee \neg\neg\phi = \top$, not to be confused with a topological Stone space); in $\mathbf{Sh}(X)$, this holds if and only if $X$ is \emph{extremally disconnected} (the closure of every open set is open)~\cite{Johnstone82}.

  \item \textbf{G\"odel--Dummett Linearity ($(\phi \to \psi) \vee (\psi \to \phi) = \top$):} Holds internally in $\mathcal{E}$ if and only if the Heyting algebra $\Omega$ is linearly ordered; in presheaf topoi $\mathbf{Set}^{\mathcal{P}^{\mathrm{op}}}$ over a poset $\mathcal{P}$, this holds if and only if $\mathcal{P}$ is a \emph{forest of trees} (every upward cone is linearly ordered).

  \item \textbf{Full Booleanness ($\mathrm{LEM}$:} $\phi \vee \neg\phi = \top$): Holds internally in $\mathcal{E}$ if and only if $\Omega \simeq 1+1$ is an internal Boolean algebra; in $\mathbf{Sh}(X)$ for a $T_0$ space $X$, this holds if and only if $X$ is discrete (or in localic topoi $\mathbf{Sh}(B)$, when the frame $B$ is a complete Boolean algebra).

\end{itemize}

\end{theorem}

This dictionary precisely calibrates the internal order-algebraic properties of the Dedekind reals $\mathbb{R}_D$ in a sheaf topos $\mathcal{E} = \mathbf{Sh}(X)$ (Johnstone, \emph{Sketches of an Elephant}, Section~D4.7~\cite{Johnstone02}). In any sheaf topos $\mathbf{Sh}(X)$, $\mathbb{R}_D$ constructively forms a commutative local ring ($x$ invertible or $1-x$ invertible) and an archimedean $f$-ring. However, two stronger classical properties of the continuum bifurcate along the $\Omega_{\mathcal{E}}$-geometry lattice:

\begin{itemize}[leftmargin=1.5em]

  \item \textbf{Weak Order Linearity and the MacNeille Collapse ($\mathbb{R}_D \simeq \mathbb{R}_M$):} Taking $A$ as the baseline theory of $\mathbb{R}_D$ and $B$ as \emph{weak order linearity} $\forall x \in \mathbb{R}_D\,(x \le 0 \vee x \ge 0)$ together with the order-completion isomorphism $\mathbb{R}_D \simeq \mathbb{R}_M$ with the MacNeille reals, the exact categorical residuum $(A \Rightarrow_{\mathbf{Sh}(X)} B)$ is the De Morgan (Stone) lattice structure of $\Omega_{\mathcal{E}}$ ($\mathrm{WLEM}: \neg\phi \vee \neg\neg\phi = \top$, i.e., $X$ is extremally disconnected)~\cite{Johnstone02}.

  \item \textbf{Geometric Field Dichotomy:} By contrast, the stronger \emph{geometric field dichotomy} $\forall x \in \mathbb{R}_D\,(x = 0 \vee x \text{ invertible})$ requires that zero-sets of continuous real-valued functions be open. This fails even on extremally disconnected spaces such as $\beta\mathbb{N}$, and requires open zero-sets ($X$ is a $P$-space, or Boolean $\Omega_{\mathcal{E}} \simeq 1+1$ on compact Hausdorff bases).

\end{itemize}

\subsection{The Homological and Categorical Exactness Lattice}

Because the formal definition of the categorical residuum $\mathrm{Res}(\kappa)$ on $\mathrm{Dom}(\mathrm{Res})$ (Definitions~\ref{def:comparison-triple}--\ref{def:residuum}) applies across arbitrary $2$-categories of structured doctrines---even when $\mathcal{D}_0$ consists of non-Cartesian or non-topos categories---$\mathrm{CRM}$ provides a rigorous foundation for calibrating the diagram lemmas of \emph{homological algebra}.

In classical homological algebra, diagram lemmas such as the Five Lemma, the $3 \times 3$ Lemma, and the Snake Lemma are proved inside \emph{abelian categories} ($\mathbf{Ab}$, $R\text{-}\mathbf{Mod}$), which enjoy biproducts, kernels, cokernels, and exactness of all monomorphisms and epimorphisms. However, non-abelian algebra ($\mathbf{Grp}$, $\mathbf{LieAlg}$, $\mathbf{Rng}$, crossed modules) operates in a genuine two-dimensional bifurcation lattice of exactness doctrines rather than a single linear chain (Figure~\ref{fig:homological-lattice}):

\begin{figure}[htbp]
\centering
\begin{tikzpicture}[
  >=Stealth,
  tier/.style={draw=CRMBlue!85,rounded corners=3pt,fill=CRMBlue!6,inner sep=4pt,font=\scriptsize\bfseries},
  wit/.style={draw=black!35,rounded corners=2pt,fill=black!3,inner sep=2.5pt,font=\scriptsize,align=center},
  dwn/.style={->,thick,draw=CRMBlue}
]
  \node[tier] (reg) at (0,5.2) {$\mathbf{RegCat}$};
  \node[tier] (barr) at (-2.35,3.8) {$\mathbf{BarrExact}$};
  \node[tier] (proto) at (2.35,3.8) {$\mathbf{PointedProtomodular}$};
  \node[tier] (semi) at (0,2.4) {$\mathbf{SemiAbelian}$};
  \node[tier] (ab) at (0,1.25) {$\mathbf{Abelian}$};
  \node[tier] (ab5) at (-2.1,0.0) {$\mathbf{AB5}$};
  \node[tier] (ab4s) at (2.1,0.0) {$\mathbf{AB4}^*$};

  \draw[dwn] (reg) -- node[above left,font=\scriptsize,pos=0.55]{eff.\ eq.\ rel.} (barr);
  \draw[dwn] (reg) -- node[above right,font=\scriptsize,pos=0.55]{Split 5} (proto);
  \draw[dwn] (barr) -- (semi);
  \draw[dwn] (proto) -- (semi);
  \draw[dwn] (semi) -- node[right,font=\scriptsize]{$+$\ additive} (ab);
  \draw[dwn] (ab) -- node[above left,font=\scriptsize,pos=0.55]{$\varinjlim$ exact} (ab5);
  \draw[dwn] (ab) -- node[above right,font=\scriptsize,pos=0.55]{$\prod$ exact} (ab4s);

  % Separating countermodel badges
  \node[wit,left=0.18cm of barr] (wbarr) {$\mathbf{Set},\;\mathbf{Set}_*$\\[-1pt]\tiny($\neg$Split 5)};
  \node[wit,right=0.14cm of proto] (wproto) {$\mathbf{TopGrp},\;\mathbf{HausGrp}$\\[-1pt]\tiny($\neg$eff.\ rel.)};
  \node[wit,right=0.3cm of semi] (wsemi) {$\mathbf{Grp},\;\mathbf{LieAlg}$\\[-1pt]\tiny(Short 5 $\not\Rightarrow$ Snake)};
  \node[wit,left=0.2cm of ab5] (wab5) {$\mathbf{Sh}(X,\mathbf{Ab})$\\[-1pt]\tiny($\neg\mathbf{AB4}^*,\;\varprojlim^1\neq 0$)};
  \node[wit,right=0.2cm of ab4s] (wab4s) {$(\mathbf{Ab})^{\mathrm{op}}$\\[-1pt]\tiny($\neg\mathbf{AB5}$)};
\end{tikzpicture}
\caption{Two-dimensional homological exactness bifurcation lattice with separating countermodels.}\label{fig:homological-lattice}
\end{figure}

Here downward arrows denote strict doctrine specialization ($\mathcal{D}_{\mathrm{top}} \supsetneq \mathcal{D}_{\mathrm{bot}}$). Applying the $\mathrm{CRM}$ residuum framework to diagram chasing and infinite limits reveals the exact structural cost separating each tier in Figure~\ref{fig:homological-lattice}:

\begin{itemize}[leftmargin=1.5em]

  \item \textbf{Orthogonality of Barr-Exactness and Bourn Protomodularity:} A category $\mathcal{A}$ is \emph{regular} if it has finite limits, coequalizers of kernel pairs, and pullback-stable regular epimorphisms. From $\mathbf{RegCat}$, exactness bifurcates along two genuinely incomparable axes:
  \begin{enumerate}[label=\textup{(\alph*)},leftmargin=1.6em]
    \item \emph{Barr-exactness}~\cite[\S\S A1.3--A1.4]{Johnstone02} demands that every internal equivalence relation is effective (a kernel pair). Every elementary topos ($\mathbf{Set}$, $\mathbf{Sh}(\mathbb{C}, J)$), pointed sets $\mathbf{Set}_*$, and commutative monoids $\mathbf{CMon}$ is Barr-exact, yet fails Bourn protomodularity (in $\mathbf{Set}_*$, the inclusion $\{*, a\} \hookrightarrow \{*, a, b\}$ pulls back to an isomorphism along the split epimorphism $\{*, a, b\} \twoheadrightarrow \{*, a\}$ sending $b \mapsto *$, violating the Split Short Five Lemma).
    \item \emph{Bourn protomodularity}~\cite{Bourn91} asserts for pointed regular categories that the Split Short Five Lemma holds for split extensions (equivalently, pullback functors along split epimorphisms reflect isomorphisms). Group-like categories ($\mathbf{Grp}$, $\mathbf{Loop}$ with inverses), torsion-free abelian groups $\mathbf{TFAB}$, and Hausdorff topological groups $\mathbf{HausGrp}$ are all pointed, regular, and Bourn protomodular, so the Short Five Lemma holds without abelianness. Yet neither $\mathbf{TFAB}$ nor $\mathbf{HausGrp}$ is Barr-exact: in $\mathbf{TFAB}$, the congruence $x \equiv y \pmod 2$ on $\mathbb{Z}$ is an internal equivalence relation whose coequalizer is $\mathbb{Z} \to 0$ (whose kernel pair is all of $\mathbb{Z} \times \mathbb{Z}$), while in $\mathbf{HausGrp}$ any non-closed normal subgroup defines a non-effective equivalence relation.
  \end{enumerate}

  \item \textbf{Semi-Abelian Intersection and the Snake Lemma Residuum:} By definition, a category is \emph{semi-abelian}~\cite{JMT02,Bourn91} (cf.~\cite[\S\S A1.3--A1.4]{Johnstone02}) if and only if it is pointed, Barr-exact, Bourn protomodular, and has finite coproducts:
  \[
    \begin{aligned}
      \mathbf{SemiAbelian} \;=\; \mathbf{BarrExact} &\cap \mathbf{Protomodular} \\
      &\cap \mathbf{Pointed} \cap \mathbf{FinCoproducts}.
    \end{aligned}
  \]
  While the Short Five Lemma ($A$) holds in any pointed protomodular category, the full \emph{Snake Lemma} ($B$---constructing the connecting homomorphism $\delta: \mathrm{Ker}(c) \to \mathrm{Coker}(a)$ and proving exactness of the six-term sequence) fails in general protomodular categories. Computing the homological residuum isolates \emph{Barr-exactness and normality/subobject modularity} ($\mathbf{SemiAbelian}$~\cite{JMT02,Bourn91}) for naturality of $\delta$, and full \emph{additivity} ($\mathbf{Ab}$-enrichment, where every monomorphism is normal) to ascend to $\mathbf{Abelian}$.

  \item \textbf{Divergence of Infinite Abelian Exactness ($\mathbf{AB5}$ vs.\ $\mathbf{AB4}^*$):} Within $\mathbf{Abelian}$, Grothendieck's two higher exactness axioms bifurcate along dual infinitary directions:
  \begin{enumerate}[label=\textup{(\alph*)},leftmargin=1.6em]
    \item \emph{Axiom $\mathbf{AB5}$} asserts that filtered colimits are exact (commute with finite limits). Together with a generator, $\mathbf{AB5}$ guarantees enough injectives for injective resolutions and right derived functors $R^q F$.
    \item \emph{Axiom $\mathbf{AB4}^*$} asserts that small (in particular countable) products of epimorphisms remain epimorphic (exactness of small products), which is the dual exactness condition governing inverse limits and the vanishing of derived inverse limits $\varprojlim^1$.
  \end{enumerate}
  In abelian sheaf categories $\mathbf{Sh}(X, \mathbf{Ab})$, $\mathbf{AB5}$ holds unconditionally while $\mathbf{AB4}^*$ fails on non-discrete spaces $X$, precipitating non-vanishing $\varprojlim^1$ obstructions (examined quantitatively in Chapter~\ref{ch:base-fiber}, Sections~4.4--\ref{sec:exotic-r4-e0}).

\end{itemize}

\section{Predicativity, Arithmetic Pretopoi, $W$-Types, and Higher Inductive Types}\label{sec:predicativity-wtypes}

\subsection{Constructive Pretopoi vs.\ Two-Sorted Hyperdoctrines (The Categorical ``Big Five'')}

Whereas the lower tiers of the classical Friedman--Simpson ``Big Five'' ($\mathrm{RCA}_0$ and $\mathrm{WKL}_0$) correspond to internal arithmetic and Cantor-locale compactness~\cite{Simpson09}, the upper tiers ($\mathrm{ACA}_0$, $\mathrm{ATR}_0$, and $\Pi^1_1\text{-}\mathrm{CA}_0$) calibrate arithmetical comprehension, transfinite bar recursion along well-orderings, and impredicative inductive definitions.

In Categorical Reverse Mathematics, calibrating this proof-theoretic climb requires separating two settings (Core Conjecture~5, \S\ref{sec:five-core-conjectures}). In single-sorted \emph{constructive} categories (without $\mathrm{LEM}$), where $1 + 1$ classifies only complemented decidable subobjects $\mathrm{Sub}_{\mathrm{dec}}(X) \subsetneq \mathrm{Sub}_{\mathcal{E}}(X)$, universes ascend strictly from predicative \emph{Arithmetic Universes} to impredicative \emph{Elementary Topoi}:

\[
  \mathbf{ArithUniv} \subsetneq \Pi\text{-}\mathbf{Pretopos} \subsetneq \Pi W\text{-}\mathbf{Pretopos} \subsetneq \mathbf{ElemTopos}.
\]

By contrast, imposing full classical logic ($\mathrm{LEM}$ on $\mathrm{Sub}_{\mathcal{E}}(X)$) on a single-sorted $\Pi$-pretopos forces every subobject $S \rightarrowtail X$ to be complemented ($S \amalg \neg S \xrightarrow{\sim} X$) and hence classified by $\Omega_{\mathcal{E}} := 1 + 1$, collapsing $\mathbf{Bool}\text{-}\Pi\text{-}\mathbf{Pretopos} = \mathbf{Bool}\text{-}\Pi W\text{-}\mathbf{Pretopos} = \mathbf{BoolElemTopos}$~\cite[Lemma~A1.4.3 \& Cor.~A2.1.6]{Johnstone02} (interpreting $\mathsf{Mac} \simeq \mathsf{ETCS}$, well above second-order arithmetic $\mathsf{Z}_2$). Retaining classical $\mathrm{LEM}$ on formulas without collapsing $1 + 1$ into an impredicative $\Omega$ therefore passes to two-sorted \emph{Lawvere hyperdoctrines} $p : \mathcal{P} \to \mathbb{B}$~\cite{Lawvere69}, calibrating each step by universal adjoints, initial algebras, or monotone fixed points:

\begin{enumerate}[label=\textbf{\arabic*.},leftmargin=1.8em]

  \item \textbf{Primitive Recursion to Arithmetical Comprehension ($\mathrm{RCA}_0 / \mathrm{WKL}_0 \to \mathrm{ACA}_0$).} Whereas $\mathrm{RCA}_0 \Rightarrow \mathrm{WKL}_0$ (both of proof-theoretic ordinal $\omega^\omega$) isolates Cantor-locale compactness on $2^{\mathbb{N}}$ (\S\ref{sec:panorama-continuum-spatial}; Theorem~\ref{thm:realizability-fan-wkl}), an \emph{Arithmetic Universe}\index{Arithmetic Universe} (Joyal--Maietti~\cite{Maietti05}) is a predicative list-arithmetic pretopos with parameterized $\mathrm{NNO}$, disjoint stable sums, and effective quotients, lacking $Y^X$ and $\Omega$. Ascending to ordinal $\varepsilon_0$ ($\mathrm{ACA}_0$)~\cite{Friedman75,Simpson09} corresponds constructively to \emph{local Cartesian closure} ($\Pi_{x:A} B(x)$ right adjoint to $f^*$, elevating $\mathbf{ArithUniv}$ to $\Pi\text{-}\mathbf{Pretopos}$, where $\mathsf{HA}^\omega \equiv_{\mathrm{Con}} \mathsf{PA}$) and hyperdoctrinally to full first-order quantifier comprehension $\exists_{\mathbb{N}} \dashv \pi^* \dashv \forall_{\mathbb{N}}$ reflected into $2^{\mathbb{N}}$.

  \item \textbf{Arithmetical Comprehension to Transfinite Recursion ($\mathrm{ACA}_0 \Rightarrow \mathrm{ATR}_0$ vs.\ $\Pi W$-Pretopoi).} Theorems at $\mathrm{ATR}_0$ (proof-theoretic ordinal $\Gamma_0$; cf.\ well-quasi-ordering benchmarks in~\cite{Friedman75,Simpson09}) iterate arithmetical comprehension along countable well-orderings. In a $\Pi$-pretopos, every container $f: B \to A$ defines a polynomial endofunctor $P_f(X) = \Sigma_{a \in A} X^{B_a}$ whose initial algebra is a \emph{Martin-L\"of $W$-type}\index{W-types (polynomial initial algebras)}~\cite{Maietti05}. Crucially, unrestricted subobject induction along $W$-types in a full $\Pi W$-pretopos interprets Aczel's $\mathsf{CZF} \equiv_{\mathrm{Con}} \mathsf{KP}\omega$~\cite{Aczel78} (Bachmann--Howard ordinal, strictly overshooting $\Gamma_0$); Core Conjecture~5 posits that restricting initial $P_f$-algebra recursion to \emph{predicative arithmetical fibers} isolates $\mathrm{ATR}_0$:

\[
  \begin{aligned}
    \mathrm{ACA}_0 + \bigl(\text{Arithmetical-Fiber Initial } &W\text{-Algebras for } P_f\bigr) \\
    &\iff \mathrm{ATR}_0 \quad (\text{Core Conj.~5}).
  \end{aligned}
\]

Adding unrestricted $W$-types promotes a constructive $\Pi$-pretopos to a predicative topos ($\Pi W$-pretopos), reaching $\mathsf{CZF}/\mathsf{KP}\omega$ without impredicative power sets.

  \item \textbf{Tree Induction to Impredicative Comprehension ($\mathrm{ATR}_0 \Rightarrow \Pi^1_1\text{-}\mathrm{CA}_0$ vs.\ $\mathbf{ElemTopos}$).} Theorems at $\Pi^1_1\text{-}\mathrm{CA}_0$ (such as Cantor--Bendixson or hyperarithmetical fixed points, ordinal $\psi_0(\Omega_\omega) \equiv_{\mathrm{Con}} \mathsf{ID}_{<\omega}$)~\cite{Friedman75,Simpson09} require impredicative extremal fixed points over $\mathcal{P}(\mathbb{N})$. Adding an unrestricted subobject classifier $\Omega$ and iterated power objects $\mathcal{P}^k(X) = \Omega^{\cdots\Omega^X}$ completes the constructive climb from $\Pi W\text{-}\mathbf{Pretopos}$ to $\mathbf{ElemTopos}$~\cite{Maietti05}, overshooting $\Pi^1_1\text{-}\mathrm{CA}_0$ into full second- and higher-order arithmetic ($\mathsf{Z}_2$ and $\mathsf{Mac} \simeq \mathsf{ETCS}$~\cite{Cole73}, Chapter~\ref{ch:sets-residuum}), whereas Core Conjecture~5 isolates \emph{monotone arithmetical Knaster--Tarski fixed points} $\mu\Phi, \nu\Phi$ on $\mathcal{P}(\mathbb{N})$ (or finitely iterated $W$-algebras) as the exact residuum for $\Pi^1_1\text{-}\mathrm{CA}_0$.

\end{enumerate}

\subsection{\texorpdfstring{From $0$-Truncated $W$-Types to Higher Inductive Types (HITs)}{From 0-Truncated W-Types to Higher Inductive Types (HITs)}}

In the preceding subsection, we formulated the principle that in a 1-topos or $\Pi$-pretopos, the categorical residuum bridging arithmetical comprehension ($\mathrm{ACA}_0$) to transfinite bar recursion ($\mathrm{ATR}_0$) is the existence of initial algebras for polynomial endofunctors (Martin-Löf $W$-types).

In an $(\infty, 1)$-topos $\mathcal{X}$, ordinary $W$-types generate $0$-truncated well-founded trees, which cannot construct higher spheres $S^n$, suspensions $\Sigma X$, Eilenberg--MacLane spaces $K(G, n)$, or homotopy pushouts $A \amalg^h_C B$. Taking $A$ as $0$-truncated $W$-type recursion and $B$ as Synthetic Homotopy Theory (such as the calculation $\pi_1(S^1) \simeq \mathbb{Z}$ or the Seifert--van Kampen theorem), the categorical residuum $(A \Rightarrow_{(\infty,1)\text{-}\mathbf{Cat}} B)$ is the existence of \emph{Higher Inductive Types (HITs)}---specifically, universal homotopy pushouts and initial algebras for polynomial endofunctors enriched with path constructors.

\section{\texorpdfstring{Univalence, Blakers--Massey Excision, and Directed Stabilization Benchmarks in $\mathrm{CatSp}$}{Univalence, Blakers-Massey Excision, and Directed Stabilization Benchmarks in CatSp}}\label{subsec:excision-catsp}

\subsection{Weak Classifiers to Voevodsky Univalence}

In an $(\infty, 1)$-category $\mathcal{X}$ with finite limits, a universe (or object classifier) for a class of $\kappa$-small morphisms is an object $\mathcal{U}$ equipped with a universal small fibration $\widetilde{\mathcal{U}} \to \mathcal{U}$ such that every $\kappa$-small morphism $X \to Y$ is a pullback of $\widetilde{\mathcal{U}} \to \mathcal{U}$~\cite[\S 6.1.6]{Lurie09}. In a weak higher pretopos, the classifying map $Y \to \mathcal{U}$ need only be weakly unique (surjective on equivalences).

Elevating weak classification ($A$, sufficient for basic type dependency) to Voevodsky's \emph{Univalence Axiom}~\cite[\S 2.10]{HoTT13} ($B$, asserting that the canonical map $\mathrm{Id}_{\mathcal{U}}(X, Y) \xrightarrow{\sim} \mathrm{Equiv}(X, Y)$ is an equivalence for all $X, Y: \mathcal{U}$, enabling structure identity principles and descent) isolates Colimit Universality and Descent (Giraud's $\infty$-Axioms)~\cite[\S\S 6.1.3, 6.1.6]{Lurie09} as the exact structural residuum on $\mathcal{X}$~\cite{GepnerKock17}.

\subsection{\texorpdfstring{Blakers--Massey Excision, Masuda's 3-Periodic Triad, and Gei\ss's Oriented Freudenthal Theorem}{Blakers-Massey Excision, Masuda's 3-Periodic Triad, and Geiss's Oriented Freudenthal Theorem}}

In classical algebraic topology, the Freudenthal Suspension Theorem ($\pi_{k+n}(S^n)$ stabilizes for $n > k+1$) and the Blakers--Massey Homotopy Excision Theorem (measuring how far a homotopy pushout square is from being a homotopy pullback square) are proved using singular chains or CW-approximation.

In Higher $\mathrm{CRM}$, these theorems hold constructively inside any $(\infty, 1)$-topos $\mathcal{X}$ without invoking $\mathrm{LEM}$ or $\mathrm{AC}$. When we weaken the ambient universe $\mathcal{X}$ to a general lex $(\infty, 1)$-category with pushouts and ask for the categorical residuum $\mathcal{R}(\mathcal{X})$ that bridges basic homotopy pushouts ($A$) to the \emph{Blakers--Massey Excision and Freudenthal Suspension Theorems} ($B$), Anel--Biedermann--Finster--Joyal~\cite[\S\S 3--4, Thm.~4.1.1]{ABFJ20} proved that the exact residuum is \emph{universality of colimits and modality descent} (the $(\infty, 1)$-pretopos exactness axioms).

When passing to stable and directed stable homotopy, Giraud strict initiality~\cite[Thm.~6.1.0.6, \S 6.1.1]{Lurie09} forbids any non-trivial ``stable $(\infty, 1)$-topos'' and stabilization $\Sigma^\infty_+(*) \simeq \mathbb{S}$ in $\mathrm{Pr}^L$~\cite[\S\S 1.4.4, 4.8.2]{Lurie17} (never on $U$ or $\mathcal{U}$) forces directed stabilization to pass to \emph{categorical spectra}\index{Categorical spectra@Categorical spectra $\mathrm{CatSp}$} $\mathrm{CatSp} := \mathrm{Sp}(\mathrm{Cat}^R_*)$ along $\Sigma_{\mathrm{red}} \dashv \Omega_{\mathrm{dir}}$ across $\mathrm{CatSp}^L \subsetneq \mathrm{CatSp}^\omega \subsetneq \mathrm{CatSp}^{\mathrm{coind}} \subsetneq \mathrm{CatSp}^R$ (Chapter~\ref{ch:higher-crm}, \S\ref{subsec:crm-discovery-engine-inf-inf};~\cite{Masuda26,Heine26,ORW26}):
\begin{itemize}[leftmargin=1.5em]
  \item \textbf{Masuda v2 Steiner Partial-Duality Correction and 3-Periodic Lax/Bifiber Triad:} In~\cite[Section~2.4, Footnote~5]{Masuda26} (crediting Marcus Nicolas), Masuda withdraws v1's Proposition~2.4.17 regarding odd/even partial duals $(-)^{\mathrm{op}}, (-)^{\mathrm{co}}$ on strong Steiner complexes---refuted by the pinched $2$-cube $\square^2 \cup_{\{10, 01\}} \square^1$---retaining monoidal preservation of $\infty\mathbf{Gaunt}_{\mathrm{Ste}}$ under the total dual $(-)^{\circ} = (-)^{\mathrm{coop}}$ and establishing a \emph{3-periodic lax/bifiber triad} in $\mathrm{CatSp}$ (\cite[Theorem~5.3.1 and Remark~5.3.4]{Masuda26}): any morphism $f : \Omega_{\mathrm{dir}} X \to Y$ induces $\overrightarrow{\mathrm{cof}}(f) \xrightarrow{\sim} Z \xrightarrow{\sim} \overrightarrow{\mathrm{fib}}(\Sigma_{\mathrm{dir}} f)$ with $Y \to Z \to \Sigma_{\mathrm{dir}} X$ a non-lax \emph{bifiber sequence}, cycling with strict period~$3$ through lax cofiber sequences ($\overrightarrow{\mathrm{cof}}$), bifiber sequences, and lax fiber sequences ($\overrightarrow{\mathrm{fib}}$).
  \item \textbf{Gei\ss's Oriented Freudenthal Suspension Theorem:} Extending Freudenthal stability to directed $\omega$-categories without group-completion collapse, Gei\ss~\cite[Theorem~5.1 and Corollary~5.14]{Geiss26} proves that for every pointed $n$-connected $\omega$-category $\mathcal{C}$, the directed unit $\mathcal{C} \to \Omega_{\mathrm{dir}}\Sigma_{\mathrm{red}}\mathcal{C}$ is $2n$-connected\index{Oriented Freudenthal Suspension Theorem}---established via a $2$-shifted $T^2$ comparison and the $(n + m + 2)$-connected Cartesian pushout-product bound (\cite[Theorem~4.13]{Geiss26}) on the directed James construction $\Omega_{\mathrm{dir}}\Sigma_{\mathrm{red}}\mathcal{C} \simeq \mathcal{J}(\mathcal{C}) = \varinjlim_k J_k(\mathcal{C})$---stabilizing $\pi_{r+k}(\Sigma_{\mathrm{red}}^k \mathcal{C})$ for all $k \ge \max(0, r - 2n)$ without inverting directed $1$-cells ($\tau_0(\Sigma^\infty_{\mathrm{dir}} S^0) \cong \mathbb{N}$ vs.\ $\pi_0(\mathbb{S}) \cong \mathbb{Z}$).
\end{itemize}

\section{Synthetic Differential Geometry, Cohesive Quadruples, and Higher Gauge Theory}\label{sec:cohesive-lattice}

\subsection{Kock--Lawvere Infinitesimals, Microlinearity, and $1$-Cohesion vs.\ Tinyness}

An important feature of Categorical Reverse Mathematics over classical Friedman--Simpson reverse mathematics is that $\mathrm{CRM}$ can calibrate mathematical theories that are formally inconsistent with classical logic ($\mathrm{LEM}$) and the Axiom of Choice ($\mathrm{AC}$). The premier example is \emph{Synthetic Differential Geometry (SDG)}\index{Synthetic Differential Geometry}, developed by F.~William Lawvere, Anders Kock~\cite{Kock06}, and Ieke Moerdijk--Gonzalo Reyes~\cite{MoerdijkReyes91} inside smooth Grothendieck topoi (such as the Cahiers topos or Dubuc's topos of $C^\infty$-rings).

Inside a smooth topos $\mathcal{E}$ equipped with a commutative line object $R$, one defines the object of first-order nilpotent infinitesimals $D = \{x \in R \mid x^2 = 0\}$. The foundational \emph{Kock--Lawvere Axiom} ($A$) asserts that the canonical evaluation map~\cite{Kock06,MoerdijkReyes91}:

\[
  \alpha: R \times R \xrightarrow{\sim} R^D, \quad (a, b) \mapsto \bigl(d \mapsto a + b \cdot d\bigr)
\]

is an isomorphism in $\mathcal{E}$, meaning every curve on the infinitesimal tangent disk $D$ is uniquely affine. Because $D$ contains non-zero indistinguishables ($\neg\forall d \in D,\, d = 0$ yet $\forall d \in D,\, \neg\neg(d = 0)$), the Kock--Lawvere axiom immediately refutes $\mathrm{LEM}$ and hence refutes $\mathrm{AC}$ by Diaconescu's theorem~\cite{Diaconescu75}.

Within this non-Boolean smooth universe $\mathcal{E}$, $\mathrm{CRM}$ calibrates the deductive hierarchy of differential geometry and cohomology via two exact geometric residuums:

\begin{enumerate}[label=\textbf{\arabic*.},leftmargin=1.8em]

  \item \textbf{First-Order Calculus to Jet Bundles and Curvature (The Microlinearity Residuum).} While the Kock--Lawvere axiom ($A$) suffices for elementary single-variable derivatives on $R$, it does not guarantee that general smooth manifolds $M \in \mathcal{E}$ possess well-behaved tangent bundles $M^D$, Lie brackets of vector fields, affine connections, or Riemann curvature tensors ($B$). In a general smooth topos, higher infinitesimal objects---such as $D(2) = \{(d_1, d_2) \in D \times D \mid d_1 d_2 = 0\}$---form finite colimits of Weil algebras in the opposite category of infinitesimal spaces. Computing the categorical residuum $(A \Rightarrow_{\mathbf{SmoothTopos}} B)$ isolates \emph{microlinearity}~\cite{Kock06,MoerdijkReyes91}: the exactness condition that the functor $M^{(-)}$ sends canonical finite colimits of infinitesimal varieties (such as $D \amalg_{\{0\}} D \to D(2)$) to finite limits (pullbacks $M^{D(2)} \xrightarrow{\sim} M^D \times_M M^D$):

\[
  \begin{aligned}
    &\text{Kock--Lawvere Axiom } (A) + \text{Microlinearity} \\
    &\quad\Longleftrightarrow \text{Connections, Lie Brackets, and Curvature } (B)
  \end{aligned}
\]

  \item \textbf{Local Infinitesimal Calculus to Global Cohomology (Cohesion vs.\ Tinyness Across Smooth Regimes).} Even when a smooth topos $\mathcal{E}$ satisfies the Kock--Lawvere axiom and microlinearity (supporting local differential forms and infinitesimal simplices $D^n$), global theorems such as Stokes' Theorem and the de Rham isomorphism between synthetic singular cohomology and differential-form cohomology ($B$) can fail. Following Moerdijk and Reyes~\cite{MoerdijkReyes91}, solving $A + \mathcal{R}(\mathcal{E}) \iff B$ bifurcates into two distinct geometric residuum regimes depending on the presentation site of $\mathcal{E}$:

  \begin{itemize}[leftmargin=1.5em]

    \item \textbf{Manifold and Analytic Smooth Topoi (Integration, Archimedean Glueing, and Cohesion):} In well-adapted Grothendieck smooth topoi equipped with open-cover topologies---such as Dubuc's topos $\mathcal{G}$ and the Cahiers topos---the infinitesimal disk $D$ is \emph{not} tiny, since a right adjoint $(-)^D \dashv (-)_D$ would force $(-)^D$ to preserve open covers. In this regime, global Stokes' and de Rham theorems require the \emph{Integration Axiom}, \emph{Archimedean open-cover glueing}, and \emph{Lawvere's cohesive adjoint quadruple}~\cite{Lawvere07,Schreiber13}, where over the base topos $\mathcal{S}$, the canonical geometric morphism $\gamma: \mathcal{E} \to \mathcal{S}$ (with direct image $\Gamma = \gamma_*$) extends to a \emph{cohesive quadruple of adjoint functors}:

\[
  \Pi \dashv \mathrm{Disc} \dashv \Gamma \dashv \mathrm{Codisc}
\]

where the left adjoint $\Pi$ (the fundamental $\infty$-groupoid / path-component functor) preserves finite products, bridging local infinitesimal contractibility to global homotopy and de Rham cohomology.

    \item \textbf{Algebraic and Simplicial Smooth Presheaf Topoi (Tinyness of $D$):} In algebraic or finitely presented $C^\infty$-ring presheaf topoi $\mathbf{Set}^{\mathbb{L}^{\mathrm{op}}}$, synthetic cochain duality is bridged instead by the \emph{tinyness (internal atomicity) of the infinitesimal disk $D$}~\cite{Kock06,MoerdijkReyes91}: the tangent-bundle exponentiation functor $(-)^D: \mathcal{E} \to \mathcal{E}$ possesses a right adjoint $(-)_D$ ($(-)^D \dashv (-)_D$), making $D$ an internally projective atom so that infinitesimal cochains dualize directly into global differential forms.

  \end{itemize}

\end{enumerate}

\subsection{\texorpdfstring{Differential Cohesive $(\infty, 1)$-Topoi and the Differential Cohomology Hexagon}{Differential Cohesive (infinity, 1)-Topoi and the Differential Cohomology Hexagon}}

In the preceding subsection, we saw that Lawvere's~\cite{Lawvere07} cohesive adjoint quadruple $\CohesiveQuad$ over a 1-topos $\mathcal{E} \to \mathcal{S}$ bridges local infinitesimal calculus to global de Rham cohomology. Passing to an $(\infty, 1)$-topos $\mathcal{X} \to \infty\text{-}\mathbf{Grpd}$ equipped with Urs Schreiber's \emph{differential cohesion}\index{Differential cohesion}~\cite[\S 4.1]{Schreiber13} ($\CohesiveQuad$ combined with an infinitesimal shape modality $\mathrm{Im} \dashv \mathrm{Inf}$) unlocks the categorical reverse mathematics of \emph{higher gauge theory and Chern--Weil theory}:

\begin{itemize}[leftmargin=1.5em]

  \item \textbf{Local $n$-Connections to Global Principal $\infty$-Bundles:} Taking $A$ as local Lie $\infty$-algebra valued differential forms on coordinate charts and $B$ as global principal $\infty$-bundles with connection (such as $B^2 U(1)$ bundle gerbes in string and M-theory), the residuum $\mathcal{R}(\mathcal{X})$ is cohesive descent and $\Pi$-product preservation~\cite[\S 4]{Schreiber13}.

  \item \textbf{Flat Connections to Differential Cohomology Hexagon:} Bridging local degree-$k$ curvature forms ($A$) to the global Simons--Sullivan differential cohomology hexagon~\cite{SimonsSullivan08} ($B$) isolates the exactness of the cohesive fracture square coupling the flat modality $\flat = \mathrm{Disc} \circ \Gamma$ to the de~Rham and infinitesimal shape modalities $\flat_{\mathrm{dR}}$ and $\mathrm{Im}$~\cite[\S 4]{Schreiber13}.

\end{itemize}

\chapter{Master Atlas of Open Problems, Methodological Blueprints, and Research Horizons}
\chaptermark{Master Atlas of Open Problems and Research Horizons}
\label{ch:master-atlas}
\chapterepigraph{``Who of us would not be glad to lift the veil behind which the future lies hidden; to cast a glance at the next advances of our science and at the secrets of its development during future centuries?''}{David Hilbert, \emph{Mathematical Problems}, ICM Paris (1900)}

\section{The Unified Architecture of Categorical Reverse Mathematics}

Across Chapters~\ref{ch:foundations}--\ref{ch:sets-residuum} and the external benchmark panorama of Chapter~\ref{ch:panorama}, Categorical Reverse Mathematics ($\mathrm{CRM}$) has been developed as a structural framework that measures the deductive gap between mathematical theorems in terms of the geometry of abstract universes and doctrinal comparisons. Where classical Friedman--Simpson reverse mathematics fixes the single formal subsystem $\mathrm{RCA}_0$ and projects theorems onto a one-dimensional ladder of set-existence axioms, and Ishihara's constructive reverse mathematics decomposes theorems over $\mathrm{BISH}$ into constructive cores plus logical or choice principles, $\mathrm{CRM}$ evaluates the residuum operator $\mathrm{Res}$ across arbitrary doctrinal comparison triples
\[
  \mathfrak{T} = \bigl(\mathcal{D}_0 \supseteq \mathcal{D}_1,\; \kappa : \Phi \Longrightarrow \Psi\bigr) \;\in\; \mathrm{Dom}(\mathrm{Res})
\]
in the sense of Chapter~\ref{ch:foundations} (Definitions~\ref{def:comparison-triple}--\ref{def:residuum} and Table~\ref{tab:five-axis-master}), calibrating the universal obstruction $\mathrm{Res}(\kappa)$ measuring how $\kappa$ fails to be an equivalence outside $\mathcal{D}_1$. In the topos-theoretic specialization $\mathcal{D}_0 = \mathbf{Topos}$, a deductive asymmetry $\mathbf{Mod}(B) \subsetneq \mathbf{Mod}(A)$ defines the \emph{categorical residuum} $\mathcal{R} = (A \Rightarrow_{\mathbf{Topos}} B)$ as the relative pseudocomplement satisfying:

\[
  \mathbf{Mod}(A) \cap \mathbf{Mod}(\mathcal{R}) = \mathbf{Mod}(B)
\]

Equivalently, $\mathcal{R}(\mathcal{E})$ is the weakest orthogonal structural invariant on $\mathcal{E}$ validating the Fundamental Equivalence $A + \mathcal{R}(\mathcal{E}) \iff B$. Because $\mathcal{R}(\mathcal{E})$ is formulated via universal categorical properties rather than second-order arithmetic syntax, any proven reversal transfers functorially across the major domains of topos and higher category theory: Grothendieck sheaf and localic topoi (Chapters~1--2, 4, and~\ref{ch:panorama}), Joyal $\delta$-rings (Chapter~\ref{ch:delta-rings}), higher $(\infty, 1)$-topoi, univalent type universes, and directed $(\infty, \infty)$-categories (Chapter 5), computational realizability and Fraïssé permutation topoi (Chapter 6), and Joyal--Moerdijk algebraic set-theoretic small-map ambients $(\mathcal{C}, \mathcal{S})$ alongside univalent $(\infty, 1)$-topos universes $(\mathcal{X}, \mathcal{U})$ (Chapter~\ref{ch:sets-residuum}).

To prevent the Triviality Trap---where the material implication $A \to B$ is offered as a tautological solution---every valid residuum $\mathcal{R}(\mathcal{E})$ must obey the Admissibility Criterion by belonging to one of four orthogonal structural families:

\begin{itemize}[leftmargin=1.5em]

  \item \textbf{Classifier Geometry:} Order-theoretic and modal properties of $\Omega$, including Booleanness ($\Omega \simeq 1+1$), De Morgan/Stone lattices, Gödel--Dummett linearity, compactness, and Lawvere--Tierney topologies $j: \Omega \to \Omega$.

  \item \textbf{Generator and Dimensional Exactness:} Projectivity, covering, and dimensional-localization properties, such as external projectivity of $1_{\mathcal{E}}$, localicity (subobjects of $1_{\mathcal{E}}$ generate $\mathcal{E}$), $0$-truncated generation in $(\infty, 1)$-topoi, $\mathrm{COSHEP}$, and reflective localization exactness across the four-tier tower $\mathrm{Cat}^L_{(\infty, \infty)} \subsetneq \mathrm{Cat}^\omega_{(\infty, \infty)} \subsetneq \mathrm{Cat}^{\mathrm{coind}}_{(\infty, \infty)} \subsetneq \mathrm{Cat}^R_{(\infty, \infty)} := \varprojlim_\kappa \mathrm{Cat}_{(\infty, n)}$ in $(\infty, \infty)$-categories.

  \item \textbf{Site Topology:} Combinatorial and localic invariants of presentation sites $(\mathbb{C}, J)$, such as the Right Ore condition on cospans, groupoid sites, finitary coherence, supercompactness, zero-dimensionality, and $\omega$-presentability.

  \item \textbf{Cohomological Vanishing:} Homological and higher-gerbe invariants, including Barr-exactness, subobject modularity, torsor vanishing $H^1(\mathcal{E}/Y,\allowbreak G) = 0$, Giraud non-abelian gerbe vanishing $H^2_{\mathrm{gerbe}}(\mathcal{X}/Y,\allowbreak \Pi_1) = 0$ ($n = 0$), Postnikov $k$-invariants $H^{n+2}(\mathcal{X}/Y,\allowbreak \Pi_{n+1}) = 0$ ($n \ge 1$), and continuous group cohomology $H^1_{\mathrm{cont}}(U,\allowbreak S_k) = 0$.

\end{itemize}

\section{Methodological Blueprint for Computing a Categorical Residuum}

Solving the residuum equivalence $A + \mathcal{R}(\mathcal{E}) \iff B$ for a new pair of mathematical theorems follows a five-step analytical protocol distilled from our case studies across Chapters~\ref{ch:foundations}--\ref{ch:sets-residuum}:

\begin{enumerate}[label=\textbf{Step \arabic*:},leftmargin=4.2em]

  \item \textbf{Internalization:} Cast $A$ and $B$ in the internal language of $\mathcal{E}$ to fix universe classes $\mathbf{Mod}(B) \subsetneq \mathbf{Mod}(A)$.

  \item \textbf{Countermodel Separation:} Exhibit a universe $\mathcal{E} \in \mathbf{Mod}(A) \setminus \mathbf{Mod}(B)$, such as $\mathbf{Sh}_{\neg\neg}(\mathbb{A})$, $\Eff$, or $\mathbf{Cont}(\mathrm{Aut}(\mathcal{M}))$.

  \item \textbf{Obstruction Extraction:} Diagnose the failure of $B$ via cohomology $H^1(\mathcal{E}/Y, G)$, Postnikov $k$-invariants, or Lawvere--Tierney operators $j: \Omega \to \Omega$.

  \item \textbf{Orthogonal Characterization:} Isolate the weakest admissible invariant $\mathcal{R}(\mathcal{E})$ that kills the obstruction without restating $A \to B$.

  \item \textbf{Base-Fiber Calibration:} Deform $\mathcal{E}$ (Axis 1) and the base topos $\mathcal{S}$ (Axis 2) via the Order-Theoretic Base-Fiber Conservation Principle (Principle~\ref{thm:conservation-law}).

\end{enumerate}

In Step 5, testing the reversal along a bounded geometric morphism $\gamma: \mathcal{E} \to \mathcal{S}$ determines whether $\mathcal{R}(\mathcal{E})$ is \emph{absolute} (base-invariant, like Diaconescu's $\mathrm{IAC} \implies \mathrm{LEM}$) or \emph{relative} (base-sensitive, like the Barr--Diaconescu--Freyd or Deligne--$\mathrm{BPI}$ theorems). Whenever metamathematical axioms $\mathcal{M}(\mathcal{S})$ are weakened in the base $\mathcal{S}$ from $\mathcal{M}_1(\mathcal{S})$ to $\mathcal{M}_2(\mathcal{S})$ ($\mathbf{Mod}(\mathcal{M}_1) \subsetneq \mathbf{Mod}(\mathcal{M}_2)$), the Order-Theoretic Base-Fiber Conservation Principle (Principle~\ref{thm:conservation-law}) dictates a contravariant strengthening of the fiber residuum:
\[
  \mathbf{Mod}(\mathcal{R}_2(\mathcal{E}/\mathcal{S})) \subsetneq \mathbf{Mod}(\mathcal{R}_1(\mathcal{E}/\mathcal{S}))
\]
transferring the missing metamathematical strength into heightened topological, atomic, or cohomological rigidity on $\mathcal{E}/\mathcal{S}$.

\section{The Five Core Conjectures of the CRM Program}\label{sec:five-core-conjectures}

At the core of the Categorical Reverse Mathematics research agenda stand five core conjectures that connect topos geometry, sheaf cohomology, locale theory, and higher homotopy theory:

\begin{enumerate}[label=\textbf{\arabic*.},leftmargin=1.8em]

  \item \textbf{Cohomological Vanishing:} Extend the unconditionally proved reductions for $A \in \{\top, \mathrm{IBPI}, \mathrm{IDC}, \mathrm{SVC}\}$ to a uniform assignment $\mathcal{G}_A(Y) \subseteq \mathbf{Grp}(\mathcal{E}/Y)$ across all intermediate choice axioms (such as $\mathrm{WISC}$, $\mathrm{OP}$, and $\mathrm{AC}_\lambda$) where $(A \Rightarrow_{\mathbf{Topos}} \mathrm{AC})$ reduces to Booleanness, external projectivity of $1_{\mathcal{E}}$, and $H^1(\mathcal{E}/Y, G) = 0$.

  \item \textbf{Sequential Presentability:} Prove that \'etale $\omega$-presentability of localic groupoid isotropy (extending the proved subcountable-step epimorphism bridge and normal pro-discrete $\mathbf{Cont}(\mathcal{G})$ collapse) is the minimal residuum bridging $\mathrm{IDC}$ to $\mathrm{IAC}$ in Boolean Grothendieck topoi with an $\mathrm{NNO}$ over $\mathbf{Set}_{\mathrm{AC}}$ (contrasting external base-$\mathrm{DC}$ $\mathrm{CCC}$ locales in Theorem~\ref{thm:bdf-deformed}).

  \item \textbf{Zero-Dimensionality Residuum:} Show Booleanness plus \'etale clopen-partition refinement of coherent epimorphic covers (\'etale internal zero-dimensionality) is the minimal residuum elevating $\mathrm{IBPI}$ to $\mathrm{IAC}$ in coherent Grothendieck topoi over $\mathbf{Set}_{\mathrm{AC}}$ (and to external $\mathrm{AC}$ when $1_{\mathcal{E}}$ is externally projective, killing $H^1(\mathcal{E}/Y, K_X)$).

  \item \textbf{Postnikov Stepwise Vanishing:} Complement the unconditionally proved abelian rungs $(\mathrm{AC}_n \Rightarrow \mathrm{AC}_{n+1})$ for $n \ge 1$ ($0$-truncated effacement of $H^{n+2}(\mathcal{X}/Y; A)$) and Giraud \emph{lien} equivalence at $n = 0$ by determining whether $0$-truncated effacement of $H^2_{\mathrm{gerbe}}(\mathcal{X}/Y, L(G))$ restricted strictly to global $0$-truncated group objects $G \in \mathbf{Grp}(\tau_{\le 0}(\mathcal{X}/Y))$ suffices for $(\mathrm{AC}_0 \Rightarrow \mathrm{AC}_1)$ across arbitrary $(\infty, 1)$-topoi.

  \item \textbf{Adjoint-Closure and Fixed-Point Hierarchy:} Formulate two-sorted Lawvere hyperdoctrine~\cite{Lawvere69} and constructive pretopos comparison triples calibrating all four adjacent ``Big Five'' transitions~\cite{Friedman75,Simpson09}---Cantor-locale compactness ($\mathrm{RCA}_0 \Rightarrow \mathrm{WKL}_0$), quantifier comprehension adjoints $\exists_{\mathbb{N}} \dashv \pi^* \dashv \forall_{\mathbb{N}}$ ($\mathrm{WKL}_0 \Rightarrow \mathrm{ACA}_0$), arithmetical-fiber $W$-recursion ($\mathrm{ACA}_0 \Rightarrow \mathrm{ATR}_0$ at $\Gamma_0$, avoiding the $\mathsf{CZF}/\mathsf{KP}\omega$ overshoot of unrestricted $\Pi W$-induction~\cite{Aczel78}), and monotone Knaster--Tarski fixed points ($\mathrm{ATR}_0 \Rightarrow \Pi^1_1\text{-}\mathrm{CA}_0$, below full $\Omega$ / $\mathsf{Mac} \simeq \mathsf{ETCS}$ power-object iteration~\cite{Cole73,Maietti05})---without triggering Johnstone's single-sorted Boolean collapse $\mathbf{Bool}\text{-}\Pi\text{-}\mathbf{Pretopos} = \mathbf{BoolElemTopos}$~\cite[\S A2.1]{Johnstone02}.

\end{enumerate}

Resolving these five conjectures will complete the structural dictionary between classical set-theoretic axioms and categorical geometry. In particular, Conjectures~\ref{conj:cohom-vanishing} and~\ref{conj:postnikov-stepwise}, bookended by the proved abelian rungs of~\S\ref{sec:postnikov-residuum} and Theorem~\ref{thm:limit-residuum}, assemble 1-categorical and $(\infty, 1)$-categorical choice into a single graded delooping-obstruction tower over $\mathcal{E} = \tau_{\le 0}\mathcal{X}$: clamped at degree~$0$ by $\Omega_{\mathcal{E}} \simeq 1+1$ and external projectivity of $1_{\mathcal{E}}$, governed at degree~$1$ by $BG$ torsors $H^1(\mathcal{E}/Y, G)$ (proved for $A \in \{\top, \mathrm{IBPI}, \mathrm{IDC}, \mathrm{SVC}\}$, open for intermediate $A$), at degree~$2$ ($n = 0$) by Giraud $1$-gerbes $H^2_{\mathrm{gerbe}}(\mathcal{X}/Y, \mathfrak{L})$~\cite{Giraud71} and representation $1$-gerbes $R(E_1)$ over $L(G)$ (proved for \emph{liens}, bounded-nilpotent bands, and \'etale $\mathbf{Sh}_\infty(B)_{/W}$), at degrees $n + 2 \ge 3$ ($n \ge 1$) unconditionally by $K(A, n+2)$ classifying $H^{n+2}(\mathcal{X}/Y; A)$~\cite[\S7.2.2]{Lurie09}, and at $n = \omega$ ($\forall n < \infty,\, \mathrm{AC}_n \Rightarrow \mathrm{AC}_\infty$) by hypercompleteness and $\mathrm{DC}$ killing $\infty$-connected and $\varprojlim^1$ phantoms.

\section{The Seventeen Open Problems of Categorical Reverse Mathematics}\label{sec:fourteen-open-problems}

Beyond the five core conjectures, the structural vulnerabilities, kernel-verified witnesses, and multi-axis bridges uncovered across Chapters~\ref{ch:foundations}--\ref{ch:sets-residuum} organize into seventeen open problems\index{Seventeen Open Problems (Horizons I--VIII)} spanning eight research horizons:

\subsection*{Horizon I: Topos Choice Geometry and Multi-Axis Residuum Trade-Offs}

\begin{problem}[Pareto Residuum Frontiers in Non-Linear Topos Lattices]\label{prob:open-1}%
\uses{def:comparison-triple,def:residuum,def:doctrinal-philosophy,prop:anti-casuistry-schism}%
Building on Doctrinal Philosophies $\mathfrak{P} = (\mathbb{D}, \mathcal{P})$ (Definition~\ref{def:doctrinal-philosophy}) and the constructive finite Pareto Schism $\{\mathrm{Proj}_{\mathrm{sub}}, \mathrm{Localic}\}$ (and full-topos $(\top \Rightarrow \mathrm{AC})$ vs.\ $(\mathrm{IAC} \Rightarrow \mathrm{AC})$ splits) of Proposition~\ref{prop:anti-casuistry-schism}, determine: \textup{(a)} intrinsic functorial criteria on mixed-variance Chu conversions $(F, G)$ in $\mathbf{Phil}$~\cite{Shulman20,HHLN23} preserving minimal Pareto frontiers $\mathrm{Res}_{\mathcal{P}}^{\min}$ against Dogma, Blindness, and non-topos axes; \textup{(b)} general small-site characterizations of hereditary subterminal projectivity $\mathrm{SS} = \mathrm{Proj}_{\mathrm{sub}}$~\cite{FreydScedrov90} beyond monoid right-zero ideals~\cite{HemelaerRogers21} and its $3^k 2^m$ slice subtopos strata over $\mathbf{Set}^{\mathcal{M}_1}/X$; and \textup{(c)} how the $\{\mathrm{Proj}_{\mathrm{sub}}, \mathrm{Localic}\}$ Pareto frontier deforms over choiceless base topoi $\mathcal{S}$.
\end{problem}

\begin{problem}[Iterated Tower Conservation Across Composite Geometric Morphisms]\label{prob:open-2}
Extend the Base--Fiber Conservation Principle (Principle~\ref{thm:conservation-law}, Chapter~\ref{ch:base-fiber}) and its quantitative $5$-term Leray sequence to arbitrary $n$-stage towers of bounded geometric morphisms $\mathcal{E}_n \xrightarrow{\gamma_n} \mathcal{E}_{n-1} \to \cdots \xrightarrow{\gamma_1} \mathcal{S}$, classifying the exact transfer of metamathematical strength from $\mathcal{S}$ into iterated fiber invariants across all intermediate stages $\mathcal{E}_k$.
\end{problem}

\subsection*{Horizon II: Computability, Realizability, and Lawvere--Tierney Modalities}

\begin{problem}[Weihrauch-Degree and Lawvere--Tierney Topology Duality]\label{prob:open-3}
Given Kihara's~\cite{Kihara23} lattice isomorphism $\mathbf{LTT}(\Eff) \cong \mathcal{W}_{\mathrm{gen,ext}}$ and Maschio--Trotta's~\cite{MaschioTrotta26} extended Weihrauch tripos $\mathsf{eW}$ recovering relative realizability topoi as $\mathbf{Sh}_j$ (Theorem~\ref{thm:ltt-residuum}, \cite{Hyland82,LeeVanOosten13,vanOosten08}; cf.~\cite{LietzStreicher12}), determine which single-call classical Weihrauch cones in $\mathcal{W}$ correspond to $j$-hereditary subtopos conditions and characterize the explicit PCA-oracle, $\cap$-preserving, and $(\cup, \cap)$-preserving sublattices of $\mathbf{LTT}(\Eff)$.
\end{problem}

\begin{problem}[Minimal Oracle and Non-Standard Modalities Beyond $j_{\mathrm{Lif}}$]\label{prob:open-4}
Explicitly characterize the minimal Lawvere--Tierney topologies $j_{(A \Rightarrow B)} : \Omega \to \Omega$ in $\mathbf{LTT}(\Eff)$ realizing intermediate semi-classical and combinatorial principles lying strictly between recursive arithmetic ($\mathrm{CT} + \mathrm{MP}$) and $\mathrm{ACA}_0$---including Weak Weak K\"onig's Lemma ($\mathrm{WWKL}$), Ramsey's theorem for pairs ($\mathrm{RT}^2_2$), and $\Pi^1_1$-avoidance principles---over relative realizability topoi $\mathbf{RT}(\mathcal{K}_1^O)$.
\end{problem}

\subsection*{Horizon III: Atomic Choice Spectra, Nominal Torsors, and Site Combinatorics}

\begin{problem}[Higher-Cardinal Stability and Kihara's Five-Rung Choice Spectrum]\label{prob:open-5}
Calibrate choice thresholds across Boolean atomic Grothendieck topoi $\mathbf{Sh}(\mathcal{K}_\lambda^{\mathrm{op}}, J_{\mathrm{at}}) \simeq \mathbf{Cont}(\mathrm{Aut}(\mathcal{M}))$ along Kihara's~\cite{Kihara26} five-rung spectrum $\mathrm{LEM} \subsetneq \mathrm{DC}_\omega + \mathrm{AC}_{\le \lambda} \subsetneq \mathrm{AC}_{\mathrm{Set}} \subsetneq \mathrm{IAC} + \mathrm{COSHEP} \subsetneq \mathrm{EAC}$ (Theorem~\ref{thm:permutation-topos} and Section~\ref{subsec:kihara-choice-nominal}) for singular and higher regular cardinals $\lambda$, separating discrete $\mathbf{Set}^G$ ($G \neq \{1\}$ at Rung~3) from non-discrete oligomorphic $\mathbf{Cont}(\mathrm{Aut}(\mathcal{M}))$ (where external projectivity of $1$, localicity, $\mathrm{COSHEP}$, $\mathrm{IAC}$, and $\mathrm{AC}$ all collapse to $\mathcal{U}(\mathcal{G}) = \{\mathcal{G}\}$).
\end{problem}

\begin{problem}[Nominal Non-Constant Sheaf Cohomology and Site-Level Ore Bridges]\label{prob:open-6}
Compute higher sheaf cohomology $H^k(\mathbf{Nom}; \mathcal{F})$ and profinite index-$k$ obstructions $H^1_{\mathrm{cont}}(U, S_k)$ for non-constant sheaves in $\mathbf{Cont}(\mathrm{Aut}(\mathcal{M}))$ beyond the parity torsor $[\mathcal{T}_{\mathrm{odd}}] = \delta(1) \neq 0 \in H^1(\mathbf{Nom}; \mathrm{EvenFinset})$ (\texttt{HomotopyCRM/\allowbreak NominalSimplicity.lean}), and classify small sites $(\mathbb{C}, J)$ bridging Right Ore $\mathrm{WLEM}$ to atomic $\mathrm{LEM}$ without groupoidifying directed embeddings (Theorems~\ref{thm:site-dictionary} and~\ref{thm:fraisse-atomic-escape}).
\end{problem}

\subsection*{Horizon IV: Directed Combinatorial Geometry and Higher Cubical Amalgamation}

\begin{problem}[The Directed $n \ge 3$ Cubical Phase Transition Beyond $F_n$]\label{prob:open-7}
Determine the minimal structural relaxation of flow-acyclicity (\texttt{IsFlowAcyclic}) or two-clause flow-convexity (\texttt{IsFlowConvex}: \texttt{cell\_\allowbreak mem} + \texttt{inter\_\allowbreak mem}) on $\mathbf{Poset}_{\mathrm{og,fin}}^{\mathrm{closed}}$ that restores all-$n$ ($n \ge 3$) cubical $n$-amalgamation beyond binary $n = 2$ \texttt{sap\_\allowbreak flowConvex} (\texttt{HomotopyCRM/\allowbreak DirectedFraisseSAP.lean}), circumventing the facet-supported directed $n$-cycle obstruction $F_n$ of Theorem~\ref{thm:directed-cubical-amalgamation} (\texttt{HomotopyCRM/\allowbreak CubicalAmalgamation.lean}) while preserving unbiased pinched-$2$-cube acyclicity under \texttt{FlowStep.dim\_\allowbreak eq}~\cite{Masuda26}.
\end{problem}

\begin{problem}[Galois-Stable vs.\ Unstable Higher $n$-Amalgamation Obstruction Calculus]\label{prob:open-8}
Develop a unified internal obstruction invariant governing the epimorphicity of the disjoint matching projection $\mathrm{Match}_m^{\mathrm{disj}}(D) \twoheadrightarrow \mathbf{y}_{D(\emptyset)}$ (Proposition~\ref{prop:matching-amalgamation}) across the $2 \times 2$ square separating $\mathrm{ACF}_0$, Goodrick--Kim--Kolesnikov~\cite{GKK14} abelian $(m-2)$-cocycle covers $T_m$ on $\partial\Delta^{m-1} \simeq S^{m-2}$, the Rado graph $\mathrm{RG}$, and Henson's $K_m$-free homogeneous graphs $\mathcal{H}_m$~\cite{Kihara26}.
\end{problem}

\subsection*{Horizon V: Differential Cohesion, Synthetic Homotopy, and Postnikov Choice Towers}

\begin{problem}[Postnikov Stepwise Gerbe and $k$-Invariant Vanishing Thresholds]\label{prob:open-9}
Resolve the remaining non-\'etale, non-nilpotent, and site-geometric open strata of Theorem~\ref{thm:postnikov-stepwise-proved} and Conjecture~\ref{conj:postnikov-stepwise} across arbitrary Grothendieck $(\infty, 1)$-topoi $\mathcal{X}$~\cite{Lurie09}: determine whether the representation $1$-gerbe $R(E_1) = 1_{\mathcal{X}/Y} \times_{\mathbf{Lien}^\kappa} \mathcal{M}_* \twoheadrightarrow 1_{\mathcal{X}/Y}$ splits beyond locally bounded-nilpotent bands and \'etale $\mathbf{Sh}_\infty(B)_{/W}$ topoi so that $0$-truncated effacement of $H^2_{\mathrm{gerbe}}(\mathcal{X}/Y, L(G))$~\cite{Giraud71} over global group objects $G \in \mathbf{Grp}(\tau_{\le 0}(\mathcal{X}/Y))$ suffices unconditionally for $(\mathrm{AC}_0 \Rightarrow \mathrm{AC}_1)$, characterize $0$-truncated effacement at all $n \ge 0$ via $(n+1)$-localic / hyperconnected site modalities~\cite[\S6.4.5]{Lurie09}, and classify non-hypercomplete $(\infty, 1)$-topoi testing the independence of Postnikov convergence from sequential $\mathrm{DC}$ in Theorem~\ref{thm:limit-residuum}.
\end{problem}

\begin{problem}[Cohesive Gauge Bridges and Differential Cohomological Descent]\label{prob:open-10}
Apply Schreiber's differential cohesive modality structure~\cite{Lawvere07,MoerdijkReyes91,Schreiber13} ($\CohesiveQuad$ together with $\mathrm{Im} \dashv \mathrm{Inf}$) inside an $(\infty, 1)$-topos $\mathcal{X}$ to extract exact CRM residuums bridging local Lie $\infty$-connections and flat curvature forms to global Chern--Weil classes, $B^n U(1)$ bundle gerbes, and the Simons--Sullivan differential cohomology fracture square~\cite{SimonsSullivan08,Schreiber13}.
\end{problem}

\subsection*{Horizon VI: Directed $(\infty, \infty)$-Localizations, Phantoms, and $\mathrm{CatSp}$}

\begin{problem}[Reflective Localization Boundaries Across the Four-Tier $(\infty, \infty)$-Ladder and \texttt{ICaTT} Adequacy]\label{prob:open-11}
Complete the residuum atlas across Ozornova--Rovelli--Walde and Gepner--Heine's strict four-tier chain inside $\mathrm{Cat}^R_{(\infty, \infty)} := \varprojlim_\kappa \mathrm{Cat}_{(\infty, n)}$~\cite{ORW26,GepnerHeine26},
\[
  \mathrm{Cat}^L_{(\infty, \infty)} \subsetneq \mathrm{Cat}^\omega_{(\infty, \infty)} \subsetneq \mathrm{Cat}^{\mathrm{coind}}_{(\infty, \infty)} \subsetneq \mathrm{Cat}^R_{(\infty, \infty)}
\]
(enforcing Simon Henry's erratum withdrawing Goldthorpe's earlier~\cite{Goldthorpe23} terminal-coalgebra claim on $\omega$-categories with weakly coinductive equivalences while retaining inductive initial algebras on $\mathrm{SymMon}^{\mathrm{Pres}}_\infty$, relative perfection $A_{\mathrm{perf}}$~\cite{Joyal85,BhattScholze22} without touching general prisms $(A, I)$, and $\mathbf{ICaTT}$ Grothendieck--Maltsiniotis clan localness at Tier~3 with $\mathrm{IDC}_\omega$ stream adequacy in Theorem~\ref{thm:coinductive-stream-choice}~\cite{BCM26,FinsterMimram17}). Following the negative resolution of Conjecture~C1 via the gaunt Noetherian rank-$\omega$ step-product $\mathcal{Q}$ (Theorems~\ref{thm:ordinal-rank-completeness} and~\ref{thm:classifier-scc-trapping}--\ref{thm:step-product-counterexample}, \texttt{OrdinalRankCompleteness.lean}, \texttt{StepProductPhantom.lean}, and Proposition~\ref{prop:lali-ceiling}), identify natural non-abelian internal conditions weakening $1$-acyclicity that force $\mathcal{C} \in \mathrm{Cat}^L_{(\infty, \infty)}$ inside $\mathrm{Cat}_{\le \omega}$.
\end{problem}

\begin{problem}[Directed Stabilization in $\mathrm{CatSp}$ Without Group-Completion Collapse]\label{prob:open-12}
Classify all $\pi_{\le 0}$-induced group-completion and $\Sigma_{\mathrm{red}} \dashv \Omega_{\mathrm{dir}}$ stabilization residuums in $\mathrm{CatSp} := \mathrm{Sp}(\mathrm{Cat}^R_*)$ across $\mathrm{CatSp}^L \subsetneq \mathrm{CatSp}^\omega \subsetneq \mathrm{CatSp}^{\mathrm{coind}} \subsetneq \mathrm{CatSp}^R$ (Sections~\ref{subsec:crm-discovery-engine-inf-inf} and~\ref{subsec:excision-catsp};~\cite{Masuda26,Heine26,Geiss26}), extending the non-collapsing deloopings $\mathbf{B}^\infty(\mathbf{FinSet}^{\simeq})$ and $\mathbf{B}^\infty(M_4^\#)$ to $4$-manifold cobordism categorical spectra without Wall stable diffeomorphism collapse or Giraud strict-initiality collapse.
\end{problem}

\subsection*{Horizon VII: Borel Complexity, Smooth $4$-Manifolds, and Condensed Tail Residua}

\begin{problem}[Light Condensed $\mathrm{Ext}^1_{\mathrm{CondAb}}$ vs.\ Definable $\varprojlim^1_{\mathrm{def}}$ on Exotic $\mathbb{R}^4$'s]\label{prob:open-13}
Compare the non-splitting light condensed torsor class $[q] \neq 0 \in \mathrm{Ext}^1_{\mathrm{CondAb}}(\mathcal{Q}_{E_0}, \bigoplus_{\mathbb{N}} \mathbb{F}_2)$ in $\mathbf{Sh}(\mathrm{Prof}_{\aleph_0})$ from the Double-Solvent Theorem\index{Double-Solvent Theorem} (Theorem~\ref{thm:double-solvent-r4}, Chapter~\ref{ch:base-fiber}, Section~\ref{sec:exotic-r4-e0}) with Gompf--Panagiotopoulos's definable derived limit $\varprojlim^1_{\mathrm{def}}$ along compact exhaustions $K_0 \Subset K_1 \Subset \cdots \subset \mathbb{R}^4$~\cite{GompfPanagiotopoulos26,BergfalkSmythe25}, testing whether condensed $\mathrm{Ext}^1_{\mathrm{CondAb}}$ yields a complete invariant distinguishing radial from non-radial exotic $\mathbb{R}^4$'s between Freedman's topological collapse and $\mathrm{ZFC}$ Vitali splitting---and further formalizing the parallel with the $(\infty,\infty)$-step-product phantom $X_\bullet \in RL(\mathcal{Q}) \setminus \mathcal{Q}$ (Remark~\ref{rem:step-product-e0-twin}).
\end{problem}

\begin{problem}[Pro-Oligomorphic Blindness and Turbulent Classifiers Beyond $E_{S_\infty}$]\label{prob:open-14}
Calibrate the Borel complexity transition separating Rung~4 countably classifiable moduli ($E_\infty \le_B E \le_B E_{S_\infty}$, reduced via $S_\infty$ logic actions~\cite{BeckerKechris96,BergfalkSmythe25} to nominal / Fra\"iss\'e quotient stacks over $\mathbf{Nom} \simeq \mathbf{Cont}(\mathrm{Sym}(\mathbb{N}))$) from Rung~5 turbulent moduli ($\not\le_B E_{S_\infty}$, such as complex $n$-manifolds for $n \ge 2$ up to biholomorphism, Table~\ref{tab:crm-borel-dictionary};~\cite{HjorthKechris00,Hjorth00,GompfPanagiotopoulos26,Kihara26}): determine what condensed, cohesive, or $(\infty, 1)$-topos invariants replace continuous $\mathbf{Cont}(\mathrm{Aut}(\mathcal{M}))$ classifying functors once pro-oligomorphic blindness collapses every definable nominal invariant.
\end{problem}

\subsection*{Horizon VIII: Algebraic Set Theory, Proper-Class Unraveling, and Universe Fixed Points}

\begin{problem}[Explicit Class-Coalgebra Residuum of Borel Determinacy]\label{prob:open-15}
Extract the explicit class-level obstruction invariant $\mathrm{Res}(\kappa_{\mathrm{Det}})$ blocking small strategy descent on Tony Martin's~\cite{Martin75} proper-class inductive unraveling coalgebras inside Joyal--Moerdijk~\cite{JoyalMoerdijk95} $(\mathcal{C}, \mathcal{S})$ (Chapter~\ref{ch:sets-residuum}, Remark~\ref{rem:martin-borel-coalgebra}), calibrating the sharp pincer between Power-Set-free $\mathsf{ZFC}^-$ (whose provability of Borel Determinacy breaks at $\boldsymbol{\Sigma}^0_4$ in minimal or non-standard $\beta$-models lacking $\mathcal{P}(\mathbb{N})$, in contrast to full ambient $H(\aleph_1) \subset V \models \mathsf{ZFC}$~\cite{Friedman71,MontalbanShore12,GitmanHamkinsJohnstone16}, Proposition~\ref{prop:pointillist-continuum}) and $(V_{\omega+\omega}, \mathrm{Def}(V_{\omega+\omega})) \models \mathsf{ZC}$ (failing at $\boldsymbol{\Sigma}^0_{\omega+1}$ via $\{\mathcal{P}^n(\mathbb{N})\}_{n \in \mathbb{N}} \notin \mathcal{S}$~\cite{Friedman71}, Proposition~\ref{prop:etcs-colimit-escape}).
\end{problem}

\begin{problem}[Class-Theoretic Calibration of the Four-Tier $(\infty, \infty)$-\allowbreak Localization Ladder]\label{prob:open-16}
Calibrate Ozornova--Rovelli--Walde and Gepner--Heine's four-tier reflective localization chain $\mathrm{Cat}^L_{(\infty,\infty)} \subsetneq\allowbreak \mathrm{Cat}^\omega_{(\infty,\infty)} \subsetneq\allowbreak \mathrm{Cat}^{\mathrm{coind}}_{(\infty,\infty)} \subsetneq\allowbreak \mathrm{Cat}^R_{(\infty,\infty)}$ \cite{ORW26,GepnerHeine26} and Goldthorpe's~\cite{Goldthorpe23} inductive $\Lambda$-iterated enrichment initial algebras along $\mathsf{GBC} \subsetneq\allowbreak \mathsf{NBG} \subsetneq\allowbreak \mathsf{NBG}+\mathsf{ETR}_{\mathrm{Ord}} \subsetneq\allowbreak \mathsf{KM}$ (and $\mathsf{GBC}+\mathsf{ETR}_{\mathrm{Ord}}$)~\cite{Gitman20,Williams18} without Grothendieck universes, testing whether $\alpha = 0$ stabilization on locally finite-dimensional $\mathrm{Cat}_{\le \omega}$ (Theorem~\ref{thm:ordinal-rank-completeness}) drops the metatheoretic cost from $\mathsf{ETR}_{\mathrm{Ord}}$ down to $\mathsf{GBC}$ below the $\mathrm{Con}(\mathsf{ZFC})$ reflection threshold already crossed at $\mathsf{ETR}_\omega$ (distinct from impredicative $\mathrm{Con}(\mathsf{GBC}+\mathsf{ETR}_{\mathrm{Ord}})$ at $\mathsf{KM}$).
\end{problem}

\begin{problem}[Leray Transgression for $\pi_{\mathrm{Set}}$ and Univalent Coalgebras $\nu\mathcal{P}_{\mathrm{Anima}}$]\label{prob:open-17}
Construct a quantitative transgression differential on $\pi_{\mathrm{Set}}$ (Definition~\ref{def:set-projection-fibration}) governing when vertical class torsors ($\mathbf{Card} \hookrightarrow \mathbf{Set}$ splitting, Proposition~\ref{prop:large-equiv-splitting}) remain $\mathcal{L}_\in$-conservative~\cite{Felgner71} versus leaking small-set theorems over $\mathsf{ZC}$~\cite{Glazer23} (Theorem~\ref{thm:glazer-nonconservativity}) and $\mathsf{ZFC}_{\neg\mathrm{Reg}}$~\cite{Felgner76}, and construct the d\'ecalage functor from the terminal coalgebra $\nu\mathcal{P}_{\mathrm{Anima}}$ of non-well-founded univalent multisets ($\Omega(\nu\mathcal{P}_{\mathrm{Anima}}, [q,q]) \cong \mathbb{Z}/2\mathbb{Z}$, Remark~\ref{rem:nu-anima-fossil}) to Rieger--Bernays permutation topoi $\mathbf{Cont}(\mathrm{Sym}(Q_0))$~\cite{Bernays54,Rieger57,Boffa69} (Proposition~\ref{prop:boffa-torsors}). Because naive model-space pullbacks fail to preserve subclass separation (\texttt{naive\_\allowbreak pullback\_\allowbreak breaks\_\allowbreak subclass\_\allowbreak separation}), lift the transgression onto a coupled two-level architecture---pairing Macro $\Sigma^1_0$-Morleyized classifying topoi $\gamma_I \colon \mathcal{E}[\mathbb{T}_{\mathrm{Class}}] \to \mathcal{E}[\mathbb{T}_{\mathrm{Set}}]$ whose compactly supported cohomology $H^0_c(U_{\mathrm{leak}}; \mathbb{Z})$ records syntactic reflection leaks ($\mathsf{ETR}_{\mathrm{Ord}}$, $\mathsf{KM}$) with Micro internal $\mathsf{AST}$ class-forcing sheaves $\gamma_\Sigma$ isolating predicative $\mathcal{S}$-distributivity leaks (Glazer over $\mathsf{ZC}$, Felgner over $\mathsf{ZFC}_{\neg\mathrm{Reg}}$).
\end{problem}

\section{\texorpdfstring{Machine-Checked Lean~4 Formalization Ledger (\texttt{HomotopyCRM})}{Machine-Checked Lean 4 Formalization Ledger (HomotopyCRM)}}\label{subsec:lean4-verification-ledger}

To anchor these categorical, algebraic, and combinatorial separations against vacuous definitions or hidden metamathematical assumptions, core components of Chapters~\ref{ch:foundations}, \ref{ch:delta-rings}, \ref{ch:base-fiber}, \ref{ch:higher-crm}, \ref{ch:realizability-sites}, and~\ref{ch:sets-residuum} are formalized in Lean~4 over \texttt{mathlib4} inside the \texttt{HomotopyCRM}\index{HomotopyCRM@\texttt{HomotopyCRM} Lean~4 library} library under a strict four-layer kernel verification gate (concrete \texttt{Nontrivial} carrier inhabitation, bidirectional positive/negative predicate separation, stock \texttt{mathlib4} typeclass projection with derived boundary lemmas, and \texttt{\#guard\_msgs in \#print axioms} permitting only \texttt{propext}, \texttt{Classical.choice}, and \texttt{Quot.sound} with zero \texttt{sorry}, custom axioms, or \texttt{opaque} terms; 19,961 verified lines across eleven production modules, audited via 216 \texttt{\#guard\_msgs in \#print axioms} checks):
\begin{itemize}[leftmargin=1.5em]
  \item \texttt{HomotopyCRM/\allowbreak InfinityResiduum.lean} (949 lines, 20 checks): Formalizes Chapter~1 (\S\ref{sec:finset-infinity}) finite dynamical lassos \texttt{LassoCarrier p q} / \texttt{lassoDynSys}, non-initiality in $\mathbf{FinSet}$ (\texttt{isEmpty\_\allowbreak hom\_\allowbreak of\_\allowbreak fintype\_\allowbreak to\_\allowbreak succ\_\allowbreak lasso}), coherent transition maps \texttt{lassoTransition}, profinite surjection \texttt{profiniteLassoToZHat\_\allowbreak surjective} ($\widehat{\mathbb{N}}_{\mathrm{lasso}} \twoheadrightarrow \widehat{\mathbb{Z}}$) with non-natural $-1$ separation (\texttt{zHatToProfiniteLasso\_\allowbreak neg\_\allowbreak one\_\allowbreak not\_\allowbreak mem\_\allowbreak range\_\allowbreak nat}), Pigeonhole triviality vs.\ nominal atom crossing (\texttt{dedekindResiduum\_\allowbreak eq\_\allowbreak empty\_\allowbreak of\_\allowbreak fintype}, \texttt{nominalDedekindResiduum\_\allowbreak eq\_\allowbreak empty}, \texttt{not\_\allowbreak hasFiniteSupportMap\_\allowbreak nat\_\allowbreak succ}), and biproduct-tail cokernels (\texttt{TailCokernel}, \texttt{subsingleton\_\allowbreak tailCokernel\_\allowbreak of\_\allowbreak fintype}, \texttt{nontrivial\_\allowbreak tailCokernel\_\allowbreak nat}).
  \item \texttt{HomotopyCRM/\allowbreak DeltaRingPerfection.lean} (2,113 lines, 26 checks): Formalizes Goldthorpe's slice endofunctor \texttt{sliceEndofunctor F c}, Lambek's isomorphism \texttt{lambek\_\allowbreak isIso}, fixed-point coalgebra equivalence (\texttt{FixedPoint\allowbreak CoalgebraUnder}, \texttt{theoremA\_\allowbreak homEquiv}), sequential colimit relative perfection $A_{\mathrm{perf}} = \varinjlim_\phi A$ with shift chirality $\phi_{A_{\mathrm{perf}}} = s^{-1}$ (\texttt{theoremB\_\allowbreak natEndo\_\allowbreak app\_\allowbreak eq\_\allowbreak shift\_\allowbreak inv}), \texttt{DeltaRing2Cat} with Joyal's polynomial cofactor $C_2(x, y) = -xy$, concrete $2$-torsion witness $R_2 = \mathbb{Z}[\eta]/(2\eta, \eta^2)$ with rigid automorphism group (\texttt{r2\_\allowbreak aut\_\allowbreak unique}), $\mathrm{Der}_\phi$-pseudo-torsor (\texttt{deltaRing2\_\allowbreak sub\_\allowbreak frobDerivation2}), non-isomorphic lifts (\texttt{r2Delta0\_\allowbreak not\_\allowbreak iso\_\allowbreak r2Delta1}), and Nygaard filtration separation on \texttt{DualInt} ($\mathbb{Z}[\varepsilon]/(\varepsilon^2)$).
  \item \texttt{HomotopyCRM/\allowbreak BorelTailResiduum.lean} (682 lines, 19 checks): Formalizes Chapter~4 (\S\ref{sec:exotic-r4-e0}) $E_0$ Cantor tail exact sequence $0 \to (\mathbb{N} \to_0 \mathrm{ZMod}\;2) \to (\mathbb{N} \to \mathrm{ZMod}\;2) \xrightarrow{q} \mathcal{Q}_{E_0} \to 0$ (\texttt{exact\_\allowbreak finFlipIncl\_\allowbreak e0QuotientMk}), classical $\mathsf{ZFC}$ right-inverse Hamel splitting via \texttt{Submodule.exists\_\allowbreak isCompl} (\texttt{exists\_\allowbreak zfc\_\allowbreak linear\_\allowbreak section\_\allowbreak e0Quotient}), basis-free Cantor product-topology coset density (\texttt{dense\_\allowbreak evenE0Coset}, \texttt{dense\_\allowbreak e0Coset}) forcing $T_1$ invariant constancy (\texttt{continuous\_\allowbreak e0Invariant\_\allowbreak const}, \texttt{not\_\allowbreak exists\_\allowbreak continuous\_\allowbreak t1\_\allowbreak e0Separator}), and the $E_0^{\mathrm{even}}$ parity torsor dichotomy separating continuity impossibility (\texttt{not\_\allowbreak exists\_\allowbreak continuous\_\allowbreak flipAlternating}) from discontinuous \texttt{Classical.choice} linear extension (\texttt{exists\_\allowbreak zfc\_\allowbreak linear\_\allowbreak flipAlternating}, \texttt{exists\_\allowbreak zfc\_\allowbreak linear\_\allowbreak section\_\allowbreak evenToE0Quotient}).
  \item \texttt{HomotopyCRM/\allowbreak OrdinalRankCompleteness.lean} (1,774 lines, 18 checks): Formalizes universe-safe globular state coalgebras \texttt{GlobularInfSystem S}, transfinite ordinal invertibility (\texttt{IsEquiv}, \texttt{IsCoindEquiv}, \texttt{IsOrdinalInv}, \texttt{IsOmegaInv}), Theorem~\ref{thm:ordinal-rank-completeness} (\texttt{theorem9A\_\allowbreak ordinal\_\allowbreak rank\_\allowbreak completeness}, \texttt{theorem9A\_\allowbreak omega\_\allowbreak completeness}, \texttt{orw\_\allowbreak question558\_\allowbreak stabilizes\_\allowbreak at\_\allowbreak zero}), unified rank separation $0 < 1 < 2 < \omega < \omega + 1$ in \texttt{OmegaTowerState} alongside \texttt{CoindWalkState} (\texttt{Checks.\allowbreak strict\_\allowbreak four\_\allowbreak tier\_\allowbreak equivalence\_\allowbreak separation}), and Proposition~\ref{prop:lali-ceiling} with sharp $+2$ join-cone boundary \texttt{JoinSuspState.joinCone K} (\texttt{prop11A\_\allowbreak lali\_\allowbreak stabilization\_\allowbreak above\_\allowbreak ceiling}, \texttt{prop11A\_\allowbreak ceiling\_\allowbreak sharp\_\allowbreak at\_\allowbreak max\_\allowbreak add\_\allowbreak two}).
  \item \texttt{HomotopyCRM/\allowbreak StepProductPhantom.lean} (3,505 lines, 22 checks): Formalizes $[1]$-classifier reachability reflection and single-SCC tower trapping (Theorem~\ref{thm:classifier-scc-trapping}: \texttt{theorem12B\_\allowbreak classifier\_\allowbreak scc\_\allowbreak trapping}, \texttt{theorem12B\_\allowbreak tower\_\allowbreak trapped\_\allowbreak in\_\allowbreak single\_\allowbreak scc}), $1$-acyclic Postnikov completeness (Theorem~\ref{thm:ff-acyclic-completeness}: \texttt{theorem12C\_\allowbreak antisymmetric\_\allowbreak postnikov\_\allowbreak completeness}), inductive gaunt commutative absorbing monoids (\texttt{Absorbing\allowbreak Monoid\allowbreak Carrier}, \texttt{isHereditarily\allowbreak Gaunt\_\allowbreak monoid}, \texttt{absorbing\allowbreak Monoid\_\allowbreak collapse\_\allowbreak le}), $\alpha = 0$ $\omega$-completeness on $\mathcal{Q}$ (\texttt{StepState.\allowbreak qFull}) (\texttt{stepProductQ\_\allowbreak omega\_\allowbreak completeness}, \texttt{stepProductQ\_\allowbreak rank\_\allowbreak hom}, \texttt{homQ\_\allowbreak card\_\allowbreak eq}), and the pure Postnikov phantom $X_\bullet \in RL(\mathcal{Q}) \setminus \mathcal{Q}$ negatively resolving Conjecture~C1 (Theorem~\ref{thm:step-product-counterexample}: \texttt{theorem12D\_\allowbreak stepProductQ\_\allowbreak phantom\_\allowbreak not\_\allowbreak realized}, \texttt{stepProduct\allowbreak Postnikov\allowbreak System}, \texttt{conjecture\_\allowbreak c1\_\allowbreak disproof\_\allowbreak summary}).
  \item \texttt{HomotopyCRM/\allowbreak DirectedFraisseSAP.lean} (2,125 lines, 13 checks): Formalizes \texttt{Oriented\allowbreak GradedPoset}, \texttt{SetPushout\allowbreak Square}, Steiner cycle obstructions (\texttt{steiner\_\allowbreak cycle\_\allowbreak obstruction}), two-clause \texttt{IsFlowConvex} (\texttt{cell\_\allowbreak mem} + \texttt{inter\_\allowbreak mem}) embeddings, 1-sided cobase-change (\texttt{Amalgamated\allowbreak Poset\allowbreak Square.\allowbreak isFlowReach\allowbreak Reflecting\_\allowbreak j₁}, \texttt{Amalgamated\allowbreak Poset\allowbreak Square.\allowbreak isFlowConvex\_\allowbreak j₁}), unbiased equal-dimension \texttt{FlowStep} acyclicity on Marcus Nicolas's pinched $2$-cube, and binary Strong Amalgamation (\texttt{sap\_\allowbreak flowConvex}).
  \item \texttt{HomotopyCRM/\allowbreak Nominal\allowbreak Simplicity.lean} (1,275 lines, 20 checks): Formalizes proper-open coset global-section emptiness (\texttt{isEmpty\_\allowbreak global\allowbreak Sections\_\allowbreak quotient\_\allowbreak of\_\allowbreak lt\_\allowbreak top}, \texttt{Checks.\allowbreak layer2\_\allowbreak sep\_\allowbreak guardrail5\_\allowbreak coset\_\allowbreak global\allowbreak Sections}), open-normal simplicity of $\mathrm{Sym}(\mathbb{N})$ (\texttt{open\_\allowbreak normal\_\allowbreak eq\_\allowbreak top}, \texttt{open\_\allowbreak finite\allowbreak Index\_\allowbreak eq\_\allowbreak top}), and non-vanishing odd-finset parity torsor cohomology $[\mathcal{T}_{\mathrm{odd}}] = \delta(1) \neq 0 \in \allowbreak H^1(\mathbf{Nom}; \allowbreak \mathrm{Even\allowbreak Finset})$ (\texttt{isEmpty\_\allowbreak t\_\allowbreak odd\_\allowbreak sections\_\allowbreak empty}, \v{C}ech $1$-cocycle \texttt{Checks.\allowbreak derived\_\allowbreak cech\_\allowbreak cocycle\_\allowbreak atom\allowbreak Section}).
  \item \texttt{HomotopyCRM/\allowbreak Cubical\allowbreak Amalgamation.lean} (1,636 lines, 15 checks): Formalizes \texttt{Punctured\allowbreak Cube n}, the quotient-free \texttt{Nat.find} minimal-support carrier \texttt{Cubical\allowbreak Amalgam D} (\texttt{Cubical\allowbreak Poset\allowbreak Diagram.\allowbreak IsMinimal\allowbreak Cell}, \texttt{Cubical\allowbreak Amalgam.\allowbreak exists\_\allowbreak minimal\allowbreak Ancestor}), all-$n$ disjoint $n$-amalgamation on $\mathbf{Poset}_{\mathrm{og,fin}}^{\mathrm{closed}}$ (\texttt{cubical\_\allowbreak ndap\_\allowbreak closed\_\allowbreak oriented\_\allowbreak poset}, \texttt{hasCubical\allowbreak NDAP\_\allowbreak iff\_\allowbreak is\allowbreak Intersective}), and the all-$n \ge 3$ flow-acyclic convex obstruction $F_n$ (Theorem~\ref{thm:directed-cubical-amalgamation}: \texttt{facet\allowbreak Cycle\allowbreak Diagram}, \texttt{no\_\allowbreak flow\allowbreak Acyclic\_\allowbreak cubical\_\allowbreak amalgam\_\allowbreak of\_\allowbreak three\_\allowbreak le}, \texttt{theorem10B\_\allowbreak cubical\_\allowbreak amalgamation\_\allowbreak dichotomy}, \texttt{Checks.\allowbreak cubical\allowbreak Amalgam\allowbreak Facet\allowbreak Cycle\allowbreak Order\allowbreak Iso}).
  \item \texttt{HomotopyCRM/\allowbreak NominalICaTT.lean} (1,524 lines, 21 checks): Formalizes Lemma~\ref{lem:polygraphic-support} and Theorem~\ref{thm:coinductive-stream-choice} calibrating $\mathbf{ICaTT}$ coinductive invertibility streams: classical Dependent Choice surjectivity onto step-sound relations (\texttt{surjective\_\allowbreak toRel\_\allowbreak of\_\allowbreak stepSound}, \texttt{globularInvTree\_\allowbreak isCoindEquiv}), two-sided finite polygraphic support equality on \texttt{CaTTTerm} syntax trees (\texttt{cattTerm\_\allowbreak exact\_\allowbreak support}, \texttt{cattTerm\_\allowbreak freshCell\_\allowbreak exact\_\allowbreak support}), and Schanuel's nominal permutation model on $\mathbb{A} = \mathbb{N}$ where every finite stage \texttt{FreshChain n} splits locally (\texttt{exists\_\allowbreak freshChain\_\allowbreak supported\_\allowbreak of\_\allowbreak card\_\allowbreak le}, \texttt{hasFiniteSupport\_\allowbreak freshChain}) while no infinite fresh stream or root witness stream admits finite support (\texttt{not\_\allowbreak hasFiniteSupport\_\allowbreak freshStream}, \texttt{isEmpty\_\allowbreak nominal\_\allowbreak freshRootStream}, \texttt{theorem9B\_\allowbreak choice\_\allowbreak calibration\_\allowbreak of\_\allowbreak coinductive\_\allowbreak streams}).
  \item \texttt{HomotopyCRM/\allowbreak SetsResiduum.lean} (994 lines, 16 checks): Formalizes Chapter~\ref{ch:sets-residuum} universe clamps and class-set residua: the baseline definable-class coreflection $\mathrm{Def} \dashv \pi_{\mathrm{Set}}$ and naive pullback separation (\texttt{defSetProjAdjunction}, \texttt{naive\_\allowbreak pullback\_\allowbreak breaks\_\allowbreak subclass\_\allowbreak separation}), axiom-free equivalence between pointwise homotopy-functorial $\mathbf{Card} \hookrightarrow \mathbf{Set}$ splitting and uniform small-set choice alongside \texttt{CategoryTheory.Skeleton} functorial realization (\texttt{card\_\allowbreak split\_\allowbreak iff\_\allowbreak global\_\allowbreak choice}, \texttt{type\_\allowbreak equiv\_\allowbreak splitting\_\allowbreak global\_\allowbreak choice}), faithful embedding of \texttt{PSet} into pointed coalgebras \texttt{APG} with non-bisimilarity of \texttt{quineAPG} ($q = \{q\}$) (\texttt{pset\_\allowbreak equiv\_\allowbreak iff\_\allowbreak bisimilar}, \texttt{quine\_\allowbreak atom\_\allowbreak outside\_\allowbreak wf}), Boffa's $2$-cycle non-isomorphic bisimulation witness and ordered twin Quine $\mathbb{Z}/2\mathbb{Z}$-torsor (\texttt{quineAPG\_\allowbreak bisimilar\_\allowbreak and\_\allowbreak nonIso\_\allowbreak twoCycle}, \texttt{boffa\_\allowbreak twin\_\allowbreak torsor}, \texttt{boffa\_\allowbreak twin\_\allowbreak addTorsor}), and the \texttt{ZFSet} extensional quotient gate collapsing multiplicity and branch permutations (\texttt{zfset\_\allowbreak ext\_\allowbreak quot\_\allowbreak gate}).
  \item \texttt{HomotopyCRM/\allowbreak DoctrinalFibration.lean}\index{DoctrinalFibration@\texttt{DoctrinalFibration} Lean~4 module} (3,384 lines, 26 checks): Formalizes Chapter~\ref{ch:foundations} cloven doctrinal fibrations and involutive Beck--Chevalley comparison squares (\texttt{DoctrinalFibration}, \texttt{MateDoubleSquare}, \texttt{mateDoubleSquare\_\allowbreak mate\_\allowbreak involutive}); derived tri-modal defect witnesses separating $(-1)$-connected inhabitation from $(-2)$-contractibility on $\mathbb{Z} \twoheadrightarrow \mathbb{Z}/2\mathbb{Z}$ and pro-tower phantoms (\texttt{intToZMod2TypeMode\_\allowbreak sectionSpace\_\allowbreak nontrivial}, \texttt{not\_\allowbreak intToZMod2TypeMode\_\allowbreak contractible}, \texttt{isEmpty\_\allowbreak intToZMod2AbMode\_\allowbreak sectionSpace}, \texttt{profiniteLassoTowerMode\_\allowbreak hasPhantom}, \texttt{stepProductTowerMode\_\allowbreak hasPhantom}); Doctrinal Philosophies, Inquisitor's Impossibility, and the finite constructive Pareto Schism $\{\mathrm{Proj}_{\mathrm{sub}}, \mathrm{Localic}\}$ via $\mathcal{M}_1$, the $3$-element V-poset, $D_4$, Freyd cover $(\mathbf{B}(\mathbb{Z}/2\mathbb{Z}))_\bot$, and $3^k 2^m$ slice Lawvere--Tierney enumeration (\texttt{DoctrinalPhilosophy}, \texttt{inquisitor\_\allowbreak impossibility\_\allowbreak casuistry}, \texttt{monoid\_\allowbreak proj1\_\allowbreak iff\_\allowbreak rightZero}, \texttt{m1\_\allowbreak projSub\_\allowbreak and\_\allowbreak not\_\allowbreak localic}, \texttt{vPoset\_\allowbreak localic\_\allowbreak not\_\allowbreak proj1}, \texttt{diamond\_\allowbreak proj1\_\allowbreak not\_\allowbreak projSub}, \texttt{freydCover\_\allowbreak dogma\_\allowbreak banishment}, \texttt{pareto\_\allowbreak incomparable\_\allowbreak projSub\_\allowbreak localic}, \texttt{M1SliceLTT.\allowbreak m1\_\allowbreak slice\_\allowbreak lt\_\allowbreak count}); Chapter~\ref{ch:delta-rings}'s three-stage Mode~(iii) $\delta$-ring lift tower culminating in $16$ lifts and $6$ $\mathrm{GL}_2(\mathbb{F}_2)$-orbits on $R_{2,2}$ (\texttt{delta\_\allowbreak mul\_\allowbreak p\_\allowbreak eq\_\allowbreak frob\_\allowbreak sub}, \texttt{frob\_\allowbreak annihilates\_\allowbreak p\_\allowbreak torsion}, \texttt{empty\_\allowbreak deltaLift\_\allowbreak R3}, \texttt{isEmpty\_\allowbreak deltaRing2\_\allowbreak zmod4}, \texttt{empty\_\allowbreak deltaLift\_\allowbreak R220}, \texttt{empty\_\allowbreak deltaLift\_\allowbreak R2Idemp}, \texttt{r2\_\allowbreak lift\_\allowbreak card\_\allowbreak eq\_\allowbreak two}, \texttt{r22\_\allowbreak lift\_\allowbreak card\_\allowbreak eq\_\allowbreak sixteen}, \texttt{r22\_\allowbreak iso\_\allowbreak orbits\_\allowbreak eq\_\allowbreak six}); and Chapter~\ref{ch:base-fiber}'s $\mathrm{SL}_2(\mathbb{F}_5) \twoheadrightarrow A_5$ Schur transgression injection (\texttt{commutator\_\allowbreak sl2\_\allowbreak f5\_\allowbreak eq\_\allowbreak top}, \texttt{sl2\_\allowbreak f5\_\allowbreak leray\_\allowbreak transgression\_\allowbreak injective}).
\end{itemize}

\section{Master Reference Atlas of Categorical Reverse Mathematics}

\noindent{\scriptsize\setlength{\tabcolsep}{2.6pt}\renewcommand{\arraystretch}{1.12}%
\begin{xltabular}{\textwidth}{@{}P{0.11\textwidth}P{0.17\textwidth}P{0.18\textwidth}P{0.18\textwidth}Y@{}}
\toprule
\textbf{Chapter} & \textbf{Universe Class} & \textbf{Weak Principle $A$} & \textbf{Strong Theorem $B$} & \textbf{Categorical Residuum $\mathcal{R}$} \\
\midrule
\endfirsthead
\toprule
\textbf{Chapter} & \textbf{Universe Class} & \textbf{Weak Principle $A$} & \textbf{Strong Theorem $B$} & \textbf{Categorical Residuum $\mathcal{R}$} \\
\midrule
\endhead
\bottomrule
\endfoot
Ch.~1 (\S\ref{sec:finset-infinity}) & Finitary Topos ($\neg\mathsf{Inf}$) & $\mathbf{FinSet} \models \mathrm{LEM} + \mathrm{AC}$ & $\mathrm{NNO}$ ($\mu(\mathrm{Id}+1)$) & Profinite lasso $\widehat{\mathbb{N}}_{\mathrm{lasso}} \notin \mathbf{FinSet}$, Dedekind gap, $\prod A / \bigoplus A$ \\
Ch.~2 & Grothendieck Topos & Local Choice $\top$ & External Choice $\mathrm{AC}$ & Booleanness and Localicity \\
Ch.~2 & Elementary Topos & Internal Choice $\mathrm{IAC}$ & External Choice $\mathrm{AC}$ & External projectivity of $1_{\mathcal{E}}$ \\
Ch.~2 & Elementary Topos & Projective Covers $\mathrm{COSHEP}$ & External Choice $\mathrm{AC}$ & $\mathrm{SVC}_{\mathrm{ret}}$, Booleanness, and external projectivity of $1_{\mathcal{E}}$ \\
Ch.~2 & Coherent Topos (over $\mathbf{Set}_{\mathrm{AC}}$) & Internal Prime Ideal ($\mathrm{IBPI}$) & Internal Choice ($\mathrm{IAC}$) (and $\mathrm{AC}$ when $1_{\mathcal{E}}$ is projective) & Booleanness and \'etale zero-dimensionality (Conj.~\ref{conj:bpi-ac}) \\
Ch.~2 & Boolean Topos ($\mathrm{NNO}$, over $\mathbf{Set}_{\mathrm{AC}}$) & Dependent Choice $\mathrm{IDC}$ & Internal Choice $\mathrm{IAC}$ & \'Etale $\omega$-presentability of groupoid isotropy (Conj.~\ref{conj:dc-ac}) \\
Ch.~2 (\S\ref{sec:symmetric-bv-zf}) & Symmetric $\mathrm{HS}_{\mathcal{F}} \subset V^{\mathbb{B}}$ / $\mathbf{Sh}_{\mathcal{G},\mathcal{F}}(\mathbb{B})$ & Pure $\mathsf{ZF}$ ($\Delta_0$-absoluteness + Symmetry Lemma) & $\neg\mathrm{AC}_\omega$, $\mathrm{DC}+\neg\mathrm{AC}$, $\mathrm{BPI}+\neg\mathrm{AC}$, $\mathrm{SVC}$ & Normal subgroup filter $\mathcal{F}$ (Cohen, Feferman--L\'evy, Jech $\sigma$-closed, Halpern--L\'evy, Blass~\cite{Jech73,HalpernLevy71,Blass79}) \\
Ch.~\ref{ch:panorama} & Sheaf Topos ($\mathrm{NNO}$) & Baseline $\top$ & Isomorphism $\mathbb{R}_C \simeq \mathbb{R}_D$ & Countable open-cover refinement \\
Ch.~\ref{ch:panorama} & Elementary Topos & Localic Tychonoff & Spatial Tychonoff & Internal spatiality of locales \\
Ch.~\ref{ch:panorama} & Regular Category & Short Five Lemma & Full Snake Lemma & Barr-exactness and modularity \\
Ch.~\ref{ch:panorama} & Pretopos / Hyperdoctrine & Comprehension $\mathrm{ACA}_0$ & Recursion $\mathrm{ATR}_0$ / $\Pi^1_1\text{-}\mathrm{CA}_0$ & Arithmetical $W$-algebras \& monotone fixed points (Core Conj.~5, \S\ref{sec:five-core-conjectures}; full $\Pi W, \Omega$ reach $\mathsf{CZF}, \mathsf{Mac}$) \\
Ch.~\ref{ch:panorama} & Smooth Topos & Local infinitesimal calculus & Global Stokes and de Rham & Cohesive $\CohesiveQuad$ (analytic) or tiny $D$ (algebraic) \\
Ch.~\ref{ch:delta-rings} (\S\ref{sec:delta-rings}) & $\mathbf{FrobLiftRing}_p$ / $\delta$-$\mathbf{Ring}$ & Torsion-free or $2$-torsion $R_2, R_{2,2}$ & Relative perfection $A_{\mathrm{perf}} = \varinjlim_\phi A$ & Joyal $C_p(x,y)$, $3$-stage lift tower ($2$ on $R_2$; $16$/$6$ orbits on $R_{2,2}$), Nygaard separation (not on $(R)_{\Prism}$) \\
Ch.~4 & Relative Topos $\mathcal{E}/\mathcal{S}$ & Base without $\mathrm{BPI}$ & Spatial Compactness & Supercompact 1-element covers \\
Ch.~4 (\S\ref{sec:exotic-r4-e0}) & Condensed / Smooth $4$-Manifold & Compact $=_{\mathbb{N}}$ vs.\ $\mathbf{Top}(\mathbb{R}^4)$ & Radial exotic $\mathbb{R}^4$'s ($E_0$) & $\mathcal{Q}_{E_0} = 2^{\mathbb{N}}\sslash E_0$: $\pi_0^{\mathrm{Pol}} \simeq 1$, $[q] \neq 0 \in \mathrm{Ext}^1_{\mathrm{CondAb}}$ \& $\varprojlim^1_{\mathrm{def}}$, $C^*(\mathcal{G}_{E_0}) \cong \mathrm{UHF}(2^\infty)$ ($K_0 \cong \mathbb{Z}[\tfrac{1}{2}]$) \\
Ch.~5 & $(\infty, 1)$-Topos $\mathcal{X}$ & Set Choice $\mathrm{AC}_0$ & Infinity-Choice $\mathrm{AC}_\infty$ & $0$-truncated sets cover $\mathcal{X}$ \\
Ch.~5 & $(\infty, 1)$-Topos $\mathcal{X}$ & $n$-Choice $\mathrm{AC}_n$ & $(n+1)$-Choice $\mathrm{AC}_{n+1}$ & $0$-truncated effacement of $H^{n+2}(\mathcal{X}/Y; A)$ ($n \ge 1$, proved) and Giraud $\mathfrak{L}$-gerbes ($n = 0$; group bands $L(G)$ Thm.~\ref{thm:postnikov-stepwise-proved}, Conj.~\ref{conj:postnikov-stepwise}) \\
Ch.~5 & $(\infty, 1)$-Topos $\mathcal{X}$ & Finite $\forall n,\, \mathrm{AC}_n$ & Infinity-Choice $\mathrm{AC}_\infty$ & Postnikov convergence and $\mathrm{DC}$ \\
Ch.~5 & $(\infty, \infty)$-Category $\mathcal{C} \in \mathrm{Cat}^R_{(\infty, \infty)}$ & Coinductive bi-invertibility ($\mathrm{Cat}^{\mathrm{coind}}_{(\infty, \infty)}$) & Inductive / Postnikov equivalence ($\mathrm{Cat}^L_{(\infty, \infty)}$) & Coinductive hypercompleteness across $\mathrm{Cat}^L_{(\infty, \infty)} \subsetneq \mathrm{Cat}^\omega_{(\infty, \infty)} \subsetneq \mathrm{Cat}^{\mathrm{coind}}_{(\infty, \infty)} \subsetneq \mathrm{Cat}^R_{(\infty, \infty)}$ (\texttt{ICaTT} Grothendieck--Maltsiniotis clan $\mathcal{E}_1$-localness cuts out Tier~3; Conjecture~C1 resolved negatively by Theorems~\ref{thm:classifier-scc-trapping}--\ref{thm:step-product-counterexample} via step-product $\mathcal{Q}$) \\
Ch.~5 / Ch.~\ref{ch:panorama} & Pointed $\mathrm{Cat}^R_*$ \& Directed Delooping & $\Sigma_{\mathrm{red}} \dashv \Omega_{\mathrm{dir}}$, $\mathbf{B}^\infty(\mathbf{FinSet}^{\simeq})$, $\mathbf{B}^\infty(M_4^\#)$ & $2n$-connected Oriented Freudenthal, 3-periodic triad, $\mathbb{S}_{\mathrm{top}}$ / Wall collapse & Cartesian pushout-product $(n+m+2)$-connectivity~\cite{Geiss26}, $\mathrm{CatSp} = \mathrm{Sp}(\mathrm{Cat}^R_*)$~\cite{Masuda26}, and $\pi_{\le 0}$ group-completion triggering $\mathbb{N} \to \mathbb{Z}$ (Section~\ref{subsec:crm-discovery-engine-inf-inf}) \\
Ch.~5 / Ch.~6 & Directed Combinatorial Sites \& Clan Topoi & Relational $E_{\mathrm{rel}} = \nu\Phi$ \& binary \texttt{sap\_\allowbreak flowConvex} & \texttt{ICaTT} stream $E_{\mathrm{rel}} = \mathrm{im}(\mathrm{InvStr})$ \& $n \ge 3$ flow-acyclic $n$-AP & $\mathrm{IDC}_\omega$ (fails at Rung~0 on $\mathrm{Free}_{\mathbf{CaTT}}(E[\mathbb{A}])$ in $\mathbf{Nom}$, Theorem~\ref{thm:coinductive-stream-choice}); two-clause \texttt{IsFlowConvex} on $n=2$ (keeping pinched $2$-cube $P, P^{\mathrm{op}}, P^{\mathrm{co}}$ acyclic) vs.\ $n \ge 3$ directed $n$-cycle obstruction $F_n$ (Theorem~\ref{thm:directed-cubical-amalgamation}) \\
Ch.~6 & Subtopoi of $\Eff$ & Computable Principle $A$ & Stronger Principle $B$ & Minimal operator $j_{(A \Rightarrow B)}$ \\
Ch.~6 & Atomic Topos $\mathbf{Cont}(\mathcal{G})$ & Kihara 5-Rung Choice Spectrum~\cite{Kihara26} & $\mathrm{IAC} + \mathrm{COSHEP}$ (Rung~3) vs.\ $\mathrm{EAC}$ (Rung~4) & Minimal open $\{1\} \le G$ in discrete $\mathbf{Set}^G$ vs.\ collapse of $\mathrm{EAC}$, localicity, $\mathrm{COSHEP}$, $\mathrm{IAC}$, $\mathrm{AC}$ to $\mathcal{U}(\mathcal{G}) = \{\mathcal{G}\}$ in non-discrete $\mathbf{Cont}(\mathrm{Aut}(\mathcal{M}))$ \\
Ch.~6 & Fraïssé / Nominal Topos & Constant $\mathrm{FAC}_k$ ($H^1(\mathbf{Nom};\allowbreak \Delta(\mathbb{Z}/2\mathbb{Z})) = 0$) & Non-constant torsor splitting / stream adequacy & Vanishing of $H^1(\mathbf{Nom}; \mathrm{EvenFinset})$ ($[\mathcal{T}_{\mathrm{odd}}] = \delta(1) \neq 0$) and $\mathrm{IDC}_\omega$ (Theorem~\ref{thm:coinductive-stream-choice}) \\
Ch.~\ref{ch:sets-residuum} & Class-Set Fiber $\mathbf{Mod}(\mathsf{GBC})$ & Set-Level $\mathsf{ZFC}$ ($\mathrm{Def} \dashv \pi_{\mathrm{Set}}$) & Global Choice $\mathsf{NBG}$ & Large equivalence splitting $\mathbf{Card} \hookrightarrow \mathbf{Set}$ in $\mathbf{Cat}(\mathcal{C})$ (Prop.~\ref{prop:large-equiv-splitting}; $\mathsf{LoS} \iff \mathrm{Rep} \wedge \mathsf{GC}$ over $\mathsf{GB}_0$, Prop.~\ref{prop:los-smuggles-gc}; leaks over $\mathsf{ZC}$ via~\cite{Glazer23}) \\
Ch.~\ref{ch:sets-residuum} & Predicative AST $(\mathcal{C}, \mathcal{S})$ & $\mathsf{ETCS} \simeq \mathsf{Mac}$ or $\mathsf{ZC}$ ($V_{\omega+\omega}$) & Full $\mathsf{ZFC}$ & Axis~VI quotient/image descent $\mathrm{im}(f) \in \mathcal{S}$ ($\exists_g$ covariant; colimit $\bigcup_n \mathcal{P}^n(\mathbb{N}) \notin V_{\omega+\omega}$ and $\boldsymbol{\Sigma}^0_{\omega+1}$ Borel Determinacy gap~\cite{Friedman71}) \\
Ch.~\ref{ch:sets-residuum} & Barr-Exact AST $(\mathcal{C}, \mathcal{S})$ & Foundation-Free $\mathsf{ZFC}_{\neg\mathrm{Reg}}$ & Well-Founded $V_{\mathsf{ZF}}$ ($\mathsf{ZFC}$) & Axis~VII $\mathbf{Coalg}_{\mathrm{wf}}(\mathcal{P}_{\mathcal{S}}) \hookrightarrow \mathbf{Coalg}(\mathcal{P}_{\mathcal{S}})$ inducing strict $\kappa_{\mathrm{Reg}} : \nu^{\mathrm{wf}}\mathcal{P}_{\mathcal{S}} \simeq \mu\mathcal{P}_{\mathcal{S}} \hookrightarrow \nu\mathcal{P}_{\mathcal{S}}$ (Quine atom $q = \{q\}$; Boffa slice twin $[Q_0^{\#(2)}] \neq 0$ \& parity $[\mathcal{T}_{\mathrm{odd}}(Q_0)] \neq 0$ torsors in $\mathbf{Cont}(\mathrm{Sym}(Q_0))$, Prop.~\ref{prop:boffa-torsors}) \\
Ch.~\ref{ch:sets-residuum} & Univalent $(\mathcal{X}, \mathcal{U})$ / AST $(\mathcal{C}, \mathcal{S})$ & $\mathsf{IZF}_{\mathrm{non\text{-}ext}}$ / $M = \mu\mathcal{P}_{\mathrm{Anima}}$ & Iterative Sets $V_0$ ($\mathbf{Set}_{\mathsf{ZFC}}$) & Fiberwise $(-1)$-truncation $V_0 \hookrightarrow M$ killing $\pi_1(M, \underline{A}) \simeq \mathrm{Aut}(A)$~\cite{Gylterud18}, Osius pure-set fullness \& well-pointedness, and \texttt{Quot.sound} bisimulation gate (Prop.~\ref{prop:extensionality-truncation}, Thm.~\ref{thm:scale-invariance}) \\
\end{xltabular}\par}

\section{Epilogue: Mathematical Truth as Obstruction Invariant}\label{sec:epilogue-truth}

To read mathematics through the lens of Categorical Reverse Mathematics is to abandon a one-dimensional picture of axiomatic truth in favor of structural covariance. What separates intuitionistic arithmetic from classical analysis, constructive locales from Stone spatiality, finite-rank $(\infty,\infty)$-categories from Postnikov completeness, or compact smooth $n$-manifolds from radial exotic $\mathbb{R}^4$'s is never a single undifferentiated ``extra axiom.'' Rather, every theorem $B$ relative to a weaker core $A$ over a doctrinal comparison triple $\mathfrak{T} = (\mathcal{D}_0 \supseteq \mathcal{D}_1, \kappa) \in \mathrm{Dom}(\mathrm{Res})$ projects a sharp geometric, cohomological, or modal shadow---the residuum $\mathrm{Res}(\kappa)$---measuring precisely where ambient symmetry or missing completeness blocks global gluing.

Across nine chapters, five canonical anchor witnesses have materialized this obstruction calculus on the page and in the kernel of Lean~4: the odd-finset parity torsor $[\mathcal{T}_{\mathrm{odd}}] = \delta(1) \neq 0 \in H^1(\mathbf{Nom}; \mathrm{EvenFinset})$ blocking external projectivity of $1_{\mathbf{Nom}}$ even in a Boolean atomic topos; the $2$-torsion $\delta$-ring $R_2 = \mathbb{Z}[\eta]/(2\eta, \eta^2)$ ($\mathbb{Z}\oplus\mathbb{F}_2$ additively) and torsion-free \texttt{DualInt} pinning relative slice perfection $A_{\mathrm{perf}}$ precisely against the Nygaard filtration; the gaunt Noetherian rank-$\omega$ step-product $\mathcal{Q} \subset \prod_{k \ge 1}\mathcal{A}_k$ whose absorbing off-diagonals birth the pure Postnikov phantom $X_\bullet \in RL(\mathcal{Q}) \setminus \mathcal{Q}$ refuting Conjecture~C1; the directed facet cycles $F_n$ erecting an uncrossable $n \ge 3$ cubical amalgamation wall past binary two-clause flow-convexity; and the light condensed smooth-$4$-manifold modulus $\mathcal{Q}_{E_0} = 2^{\mathbb{N}} \sslash E_0$ with CAR convolution algebra $\mathrm{UHF}(2^\infty)$ ($K_0 \cong \mathbb{Z}[\tfrac{1}{2}]$), whose $\mathrm{Ext}^1_{\mathrm{CondAb}}$ tail class mirrors $X_\bullet$ between Freedman topological collapse and $\mathrm{ZFC}$ Vitali splitting.

Backing these theorems stands the 19,961-line, eleven-module \texttt{HomotopyCRM} library in Lean~4 (\texttt{leanprover/lean4:v4.28.0}) over \texttt{mathlib4}, audited under a strict four-layer kernel verification gate with zero \texttt{sorry}, zero custom axioms, and zero \texttt{opaque} declarations across 216 \texttt{\#guard\_msgs in \#print axioms} checks. Whether the seventeen open problems across Horizons~I--VIII yield new Pareto frontiers, Weihrauch--Lawvere--Tierney dualities, higher cubical invariants, or condensed exotic-$\mathbb{R}^4$ obstructions, the guiding principle remains unchanged: mathematical necessity is never an accident of syntax, but the geometry of the obstructions that vanish when universes align. Dissecting ``$\mathbf{Sets}$'' itself as a double-clamped locus inside $(\mathcal{C}, \mathcal{S})$---bridging to univalent universes $(\mathcal{X}, \mathcal{U})$ and a Volume~II elementary class-anima horizon $(\mathcal{C}_\infty, \mathcal{S}_\infty)$---thus completes the foundational arc of Volume~I and launches Volume~II (\emph{Universes, Classes, and Higher Fixed Points}).

\backmatter
\setcounter{secnumdepth}{-1}
\bookmarksetup{startatroot}
\begin{thebibliography}{999}
\addcontentsline{toc}{chapter}{Bibliography}

\bibitem{ABFJ20}
M.~Anel, G.~Biedermann, E.~Finster, and A.~Joyal,
\emph{A generalized Blakers--Massey theorem},
Journal of Topology \textbf{13} (2020), no.~4, 1521--1553.

\bibitem{Aczel78}
P.~Aczel,
\emph{The type theoretic interpretation of constructive set theory},
in \emph{Logic Colloquium '77}, Studies in Logic and the Foundations of Mathematics, vol.~96, North-Holland, Amsterdam, 1978, pp.~55--66.

\bibitem{Aczel88}
P.~Aczel,
\emph{Non-Well-Founded Sets},
CSLI Lecture Notes, vol.~14, Stanford University, Stanford, CA, 1988.

\bibitem{AhrensCapriottiSpadotti15}
B.~Ahrens, P.~Capriotti, and R.~Spadotti,
\emph{Non-wellfounded trees in Homotopy Type Theory},
13th International Conference on Typed Lambda Calculi and Applications (TLCA 2015), LIPIcs, vol.~38, Schloss Dagstuhl--Leibniz-Zentrum f\"ur Informatik, 2015, pp.~17--30 (\href{https://arxiv.org/abs/1504.02949}{\texttt{arXiv:1504.02949}}).

\bibitem{Barr74}
M.~Barr,
\emph{Toposes without points},
Journal of Pure and Applied Algebra \textbf{5} (1974), 265--280.

\bibitem{BarrattPriddy72}
M.~G.~Barratt and S.~B.~Priddy,
\emph{On the homology of non-connected monoids and their associated groups},
Commentarii Mathematici Helvetici \textbf{47} (1972), 1--14.

\bibitem{BCM26}
T.~Benjamin, C.~Champin, and I.~Markakis,
\emph{A type theory for invertibility in weak $\omega$-categories},
\href{https://arxiv.org/abs/2602.16602}{\texttt{arXiv:2602.16602}}, 2026.

\bibitem{BeckerKechris96}
H.~Becker and A.~S.~Kechris,
\emph{The Descriptive Set Theory of Polish Group Actions},
London Mathematical Society Lecture Note Series, vol.~232, Cambridge University Press, Cambridge, 1996.

\bibitem{BergfalkSmythe25}
J.~Bergfalk and I.~B.~Smythe,
\emph{Manifold classification from the descriptive viewpoint},
preprint (2025), \href{https://arxiv.org/abs/2512.24996}{\texttt{arXiv:2512.24996}}.

\bibitem{Bernays54}
P.~Bernays,
\emph{A system of axiomatic set theory---Part~VII},
Journal of Symbolic Logic \textbf{19} (1954), no.~2, 81--96.

\bibitem{BhattScholze22}
B.~Bhatt and P.~Scholze,
\emph{Prisms and prismatic cohomology},
Annals of Mathematics \textbf{196} (2022), no.~3, 1135--1275 (\href{https://arxiv.org/abs/1905.08229v4}{\texttt{arXiv:1905.08229v4}}).

\bibitem{Bishop67}
E.~Bishop,
\emph{Foundations of Constructive Analysis},
McGraw-Hill, New York, 1967.

\bibitem{Blass79}
A.~Blass,
\emph{Injectivity, projectivity, and the axiom of choice},
Transactions of the American Mathematical Society \textbf{255} (1979), 31--59.

\bibitem{Blass83}
A.~Blass,
\emph{Cohomology detects failures of the axiom of choice},
Transactions of the American Mathematical Society \textbf{279} (1983), no.~1, 257--269.

\bibitem{Boffa69}
M.~Boffa,
\emph{Sur la th\'eorie des ensembles sans axiome de fondement},
Bulletin de la Soci\'et\'e Math\'ematique de Belgique \textbf{21} (1969), 16--56.

\bibitem{Bourn91}
D.~Bourn,
\emph{Normalization equivalence, kernel equivalence and affine categories},
Category Theory (Como, 1990), Lecture Notes in Mathematics, vol.~1488, Springer, Berlin, 1991, pp.~43--62.

\bibitem{Caramello18}
O.~Caramello,
\emph{Theories, Sites, Toposes: Relating and Studying Mathematical Theories through Topos-Theoretic `Bridges'},
Oxford University Press, Oxford, 2018.

\bibitem{Cole73}
J.~C.~Cole,
\emph{Categories of sets and models of set theory},
in \emph{The Proceedings of the Bertrand Russell Memorial Logic Conference} (Uldum, 1971), Leeds, 1973, pp.~351--399.

\bibitem{Connes94}
A.~Connes,
\emph{Noncommutative Geometry},
Academic Press, San Diego, 1994.

\bibitem{Deligne72}
P.~Deligne,
\emph{Appendice~: Crit\`ere d'existence de points},
in A.~Grothendieck and J.-L.~Verdier, \emph{Conditions de finitude. Topos et sites fibr\'es. Applications aux questions de passage \`a la limite}, Exp.~VI, Th\'eorie des Topos et Cohomologie \'Etale des Sch\'emas (SGA 4), Tome 2, Lecture Notes in Mathematics, vol.~270, Springer, Berlin, 1972, pp.~341--365.

\bibitem{Diaconescu75}
R.~Diaconescu,
\emph{Axiom of choice and complementation},
Proceedings of the American Mathematical Society \textbf{51} (1975), 176--178.

\bibitem{Easton70}
W.~B.~Easton,
\emph{Powers of regular cardinals},
Annals of Mathematical Logic \textbf{1} (1970), no.~2, 139--178 (Ph.D.\ thesis, Princeton University, 1964).

\bibitem{Felgner71}
U.~Felgner,
\emph{Comparison of the axioms of local and universal choice},
Fundamenta Mathematicae \textbf{71} (1971), 43--62.

\bibitem{Felgner76}
U.~Felgner,
\emph{Choice functions on sets and classes},
in \emph{Sets and Classes: On the Work by Paul Bernays}, Studies in Logic and the Foundations of Mathematics, vol.~84, North-Holland, Amsterdam, 1976, pp.~217--255.

\bibitem{FinsterMimram17}
E.~Finster and S.~Mimram,
\emph{A type-theoretical definition of weak $\omega$-categories},
Proceedings of the 32nd Annual ACM/IEEE Symposium on Logic in Computer Science (LICS 2017), IEEE, 2017, pp.~1--12 (\href{https://arxiv.org/abs/1706.02866}{\texttt{arXiv:1706.02866}}).

\bibitem{FourmanHyland79}
M.~P.~Fourman and J.~M.~E.~Hyland,
\emph{Sheaf models for analysis},
in \emph{Applications of Sheaves} (Durham, 1977), Lecture Notes in Mathematics, vol.~753, Springer, Berlin, 1979, pp.~280--301.

\bibitem{Fraisse54}
R.~Fra\"iss\'e,
\emph{Sur l'extension aux relations de quelques propri\'et\'es des ordres},
Annales Scientifiques de l'\'Ecole Normale Sup\'erieure (3) \textbf{71} (1954), no.~4, 363--388.

\bibitem{Freyd72}
P.~J.~Freyd,
\emph{Aspects of topoi},
Bulletin of the Australian Mathematical Society \textbf{7} (1972), 1--76.

\bibitem{FreydScedrov90}
P.~J.~Freyd and A.~\v{S}\v{c}edrov,
\emph{Categories, Allegories},
North-Holland Mathematical Library, vol.~39, North-Holland, Amsterdam, 1990.

\bibitem{Friedman71}
H.~Friedman,
\emph{Higher set theory and mathematical practice},
Annals of Mathematical Logic \textbf{2} (1971), no.~3, 325--357.

\bibitem{Friedman73}
H.~Friedman,
\emph{The consistency of classical set theory relative to a set theory with intuitionistic logic},
Journal of Symbolic Logic \textbf{38} (1973), no.~2, 315--319.

\bibitem{Friedman75}
H.~Friedman,
\emph{Some systems of second order arithmetic and their use},
Proceedings of the International Congress of Mathematicians (Vancouver, 1974), vol.~1, 1975, pp.~235--242.

\bibitem{Fujimoto12}
K.~Fujimoto,
\emph{Classes and truths in set theory},
Annals of Pure and Applied Logic 163 (2012), no.~11, 1484--1523.

\bibitem{Geiss26}
T.~Gei\ss,
\emph{The oriented Freudenthal suspension theorem},
\href{https://arxiv.org/abs/2610.38949}{\texttt{arXiv:2610.38949}}, 2026.

\bibitem{GepnerHaugseng15}
D.~Gepner and R.~Haugseng,
\emph{Enriched $\infty$-categories via non-symmetric $\infty$-operads},
Advances in Mathematics \textbf{279} (2015), 575--716.

\bibitem{GepnerHeine26}
D.~Gepner and H.~Heine,
\emph{Homotopy posets, Postnikov towers, and hypercompletions of $\infty$-categories},
\href{https://arxiv.org/abs/2603.09903}{\texttt{arXiv:2603.09903}}, 2026.

\bibitem{GepnerKock17}
D.~Gepner and J.~Kock,
\emph{Univalence in locally cartesian closed $\infty$-categories},
Forum Mathematicum \textbf{29} (2017), no.~3, 617--652 (\href{https://arxiv.org/abs/1208.1749}{\texttt{arXiv:1208.1749}}).

\bibitem{Giraud71}
J.~Giraud,
\emph{Cohomologie non ab\'elienne},
Grundlehren der mathematischen Wissenschaften, vol.~179, Springer, Berlin, 1971.

\bibitem{GitmanHamkinsJohnstone16}
V.~Gitman, J.~D.~Hamkins, and T.~A.~Johnstone,
\emph{What is the theory $\mathrm{ZFC}$ without power set?},
Mathematical Logic Quarterly \textbf{62} (2016), no.~4--5, 391--406 (\href{https://arxiv.org/abs/1110.2430}{\texttt{arXiv:1110.2430}}).

\bibitem{Gitman20}
V.~Gitman, J.~D.~Hamkins, P.~Holy, P.~Schlicht, and K.~Williams,
\emph{The exact strength of the class forcing theorem},
Journal of Symbolic Logic \textbf{85} (2020), no.~3, 869--905.

\bibitem{Glazer23}
E.~Glazer,
\emph{Global choice is not conservative over local choice for Zermelo set theory},
preprint (2023; rev.\ v3 2024), \href{https://arxiv.org/abs/2312.11902}{\texttt{arXiv:2312.11902}}.

\bibitem{Glimm60}
J.~G.~Glimm,
\emph{On a certain class of operator algebras},
Transactions of the American Mathematical Society \textbf{95} (1960), 318--340.

\bibitem{Glimm61}
J.~Glimm,
\emph{Type~I $C^*$-algebras},
Annals of Mathematics (2) \textbf{73} (1961), 572--612.

\bibitem{Goldthorpe23}
Z.~Goldthorpe,
\emph{Homotopy theories of $(\infty, \infty)$-categories as universal fixed points with respect to enrichment},
Internat.\ Math.\ Res.\ Notices (2023), no.~22, 19592--19640 (\href{https://arxiv.org/abs/2307.00442}{\texttt{arXiv:2307.00442}}, v5 2024; inductive initial-algebra fixed point on $\mathrm{SymMon}^{\mathrm{Pres}}_\infty$ and relative slice fixed points in Proposition~2.3.13; earlier terminal-coalgebra claims on $\omega$-categories with weakly coinductive equivalences withdrawn per S.~Henry's erratum); see also \emph{Sheaves of $(\infty, \infty)$-categories}, \href{https://arxiv.org/abs/2403.06926v4}{\texttt{arXiv:2403.06926v4}}, 2024.

\bibitem{GKK14}
J.~Goodrick, B.~Kim, and A.~Kolesnikov,
\emph{Type-amalgamation properties and polygroupoids in stable theories},
J.\ Math.\ Logic \textbf{15} (2015), no.~1, 1550004 (\href{https://arxiv.org/abs/1404.1525}{\texttt{arXiv:1404.1525}}).

\bibitem{GompfPanagiotopoulos26}
R.~E.~Gompf and A.~Panagiotopoulos,
\emph{On classifying smoothings of $\mathbb{R}^4$---from the beginning to the present} and \emph{On the difficulty of classifying open manifolds},
ETH Z\"urich / Harvard CMSA Michael Freedman Topology Seminar announcements (2025--2026).

\bibitem{Gylterud18}
H.~R.~Gylterud,
\emph{From multisets to sets in homotopy type theory},
Journal of Symbolic Logic \textbf{83} (2018), no.~3, 1132--1146 (\href{https://arxiv.org/abs/1612.05468}{\texttt{arXiv:1612.05468}}).

\bibitem{HalpernLevy71}
J.~D.~Halpern and A.~L\'evy,
\emph{The Boolean prime ideal theorem does not imply the axiom of choice},
in \emph{Axiomatic Set Theory}, Proceedings of Symposia in Pure Mathematics, vol.~XIII, Part~I, American Mathematical Society, Providence, RI, 1971, pp.~83--134.

\bibitem{HKL90}
L.~A.~Harrington, A.~S.~Kechris, and A.~Louveau,
\emph{A Glimm--Effros dichotomy for Borel equivalence relations},
Journal of the American Mathematical Society \textbf{3} (1990), 903--928.

\bibitem{HHLN23}
R.~Haugseng, F.~Hebestreit, S.~Linskens, and J.~Nuiten,
\emph{Lax monoidal adjunctions, two-variable fibrations and the calculus of mates},
Proceedings of the London Mathematical Society (3) \textbf{127} (2023), no.~4, 889--957.

\bibitem{Heine26}
H.~Heine,
\emph{Stable homotopy theory of higher categories},
\href{https://arxiv.org/abs/2605.05195v2}{\texttt{arXiv:2605.05195v2}}, 2026.

\bibitem{HemelaerRogers21}
J.~Hemelaer and M.~Rogers,
\emph{Monoid properties as invariants of toposes of monoid actions},
Applied Categorical Structures \textbf{29} (2021), 379--413.

\bibitem{HenryLoubaton25}
S.~Henry and F.~Loubaton,
\emph{An inductive model structure for strict $\infty$-categories},
Journal of Pure and Applied Algebra \textbf{229} (2025), no.~1, 107859 (\href{https://arxiv.org/abs/2301.11424}{\texttt{arXiv:2301.11424}}).

\bibitem{HJP80}
J.~M.~E.~Hyland, P.~T.~Johnstone, and A.~M.~Pitts,
\emph{Tripos theory},
Mathematical Proceedings of the Cambridge Philosophical Society \textbf{88} (1980), 205--232.

\bibitem{Hjorth00}
G.~Hjorth,
\emph{Classification and Orbit Equivalence Relations},
Mathematical Surveys and Monographs, vol.~75, American Mathematical Society, Providence, RI, 2000.

\bibitem{HjorthKechris00}
G.~Hjorth and A.~S.~Kechris,
\emph{The complexity of the classification of Riemann surfaces and complex manifolds},
Illinois Journal of Mathematics \textbf{44} (2000), no.~1, 104--137.

\bibitem{Hodges93}
W.~Hodges,
\emph{Model Theory},
Encyclopedia of Mathematics and its Applications, vol.~42, Cambridge University Press, Cambridge, 1993.

\bibitem{HoTT13}
The Univalent Foundations Program,
\emph{Homotopy Type Theory: Univalent Foundations of Mathematics},
Institute for Advanced Study, Princeton, NJ, 2013.

\bibitem{Hyland82}
J.~M.~E.~Hyland,
\emph{The effective topos},
The L.~E.~J.~Brouwer Centenary Symposium (Noordwijkerhout, 1981), Studies in Logic and the Foundations of Mathematics, vol.~110, North-Holland, Amsterdam, 1982, pp.~165--216.

\bibitem{Ishihara06}
H.~Ishihara,
\emph{Reverse mathematics in Bishop's constructive mathematics},
Philosophia Scientiae (2006), Cahier Sp\'ecial 6, 43--59.

\bibitem{JMT02}
G.~Janelidze, L.~M\'arki, and W.~Tholen,
\emph{Semi-abelian categories},
Journal of Pure and Applied Algebra \textbf{168} (2002), no.~2--3, 367--386.

\bibitem{Jech73}
T.~J.~Jech,
\emph{The Axiom of Choice},
Studies in Logic and the Foundations of Mathematics, vol.~75, North-Holland, Amsterdam, 1973.

\bibitem{Johnstone81}
P.~T.~Johnstone,
\emph{Tychonoff's theorem without the axiom of choice},
Fundamenta Mathematicae \textbf{113} (1981), no.~1, 21--35.

\bibitem{Johnstone82}
P.~T.~Johnstone,
\emph{Stone Spaces},
Cambridge Studies in Advanced Mathematics, vol.~3, Cambridge University Press, Cambridge, 1982.

\bibitem{Johnstone02}
P.~T.~Johnstone,
\emph{Sketches of an Elephant: A Topos Theory Compendium},
Vols.~1--2, Oxford Logic Guides, vols.~43--44, Oxford University Press, Oxford, 2002.

\bibitem{Joyal85}
A.~Joyal,
\emph{$\delta$-anneaux et vecteurs de Witt},
Comptes Rendus Math\'ematiques de l'Acad\'emie des Sciences, La Soci\'et\'e Royale du Canada \textbf{7} (1985), no.~3, 177--182.

\bibitem{JoyalMoerdijk95}
A.~Joyal and I.~Moerdijk,
\emph{Algebraic Set Theory},
London Mathematical Society Lecture Note Series, vol.~220, Cambridge University Press, Cambridge, 1995.

\bibitem{Kechris95}
A.~S.~Kechris,
\emph{Classical Descriptive Set Theory},
Graduate Texts in Mathematics, vol.~156, Springer-Verlag, New York, 1995.

\bibitem{Kihara23}
T.~Kihara,
\emph{Lawvere--Tierney topologies for computability theorists},
Transactions of the American Mathematical Society, Series B 10 (2023), no.~2, 48--85.

\bibitem{Kihara26}
T.~Kihara,
\emph{Atomic toposes, choice spectra and model-theoretic stability},
\href{https://arxiv.org/abs/2609.21111}{\texttt{arXiv:2609.21111}}, 2026.

\bibitem{Kleene52}
S.~C.~Kleene,
\emph{Recursive functions and intuitionistic mathematics},
in \emph{Proceedings of the International Congress of Mathematicians} (Cambridge, Mass., 1950), vol.~1, American Mathematical Society, Providence, RI, 1952, pp.~679--685.

\bibitem{Kock06}
A.~Kock,
\emph{Synthetic Differential Geometry},
2nd ed., London Mathematical Society Lecture Note Series, vol.~333, Cambridge University Press, Cambridge, 2006.

\bibitem{Lawvere64}
F.~W.~Lawvere,
\emph{An elementary theory of the category of sets},
Proceedings of the National Academy of Sciences \textbf{52} (1964), no.~6, 1506--1511.

\bibitem{Lawvere69}
F.~W.~Lawvere,
\emph{Adjointness in foundations},
Dialectica \textbf{23} (1969), no.~3--4, 281--296; reprinted in Reprints in Theory and Applications of Categories, no.~16 (2006), 1--16.

\bibitem{Lawvere07}
F.~W.~Lawvere,
\emph{Axiomatic cohesion},
Theory and Applications of Categories \textbf{19} (2007), no.~3, 41--49.

\bibitem{LeeVanOosten13}
S.~Lee and J.~van~Oosten,
\emph{Basic subtoposes of the effective topos},
Annals of Pure and Applied Logic \textbf{164} (2013), no.~9, 866--883.

\bibitem{LietzStreicher12}
P.~Lietz and T.~Streicher,
\emph{Realizability models refuting Ishihara's boundedness principle},
Annals of Pure and Applied Logic \textbf{163} (2012), no.~12, 1803--1807.

\bibitem{Lurie09}
J.~Lurie,
\emph{Higher Topos Theory},
Annals of Mathematics Studies, vol.~170, Princeton University Press, Princeton, NJ, 2009.

\bibitem{Lurie17}
J.~Lurie,
\emph{Higher Algebra},
preprint (September 2017), \href{https://www.math.ias.edu/~lurie/papers/HigherAlgebra.pdf}{Institute for Advanced Study}.

\bibitem{MacLane98}
S.~Mac~Lane,
\emph{Categories for the Working Mathematician},
2nd ed., Graduate Texts in Mathematics, vol.~5, Springer, New York, 1998 (1st ed., 1971).

\bibitem{MacLaneMoerdijk92}
S.~Mac~Lane and I.~Moerdijk,
\emph{Sheaves in Geometry and Logic: A First Introduction to Topos Theory},
Universitext, Springer, New York, 1992.

\bibitem{Maietti05}
M.~E.~Maietti,
\emph{Modular correspondence between dependent type theories and categories including pretopoi and topoi},
Mathematical Structures in Computer Science \textbf{15} (2005), no.~6, 1089--1149.

\bibitem{Makkai96}
M.~Makkai,
\emph{Avoiding the axiom of choice in general category theory},
Journal of Pure and Applied Algebra \textbf{108} (1996), no.~2, 109--173.

\bibitem{Martin75}
D.~A.~Martin,
\emph{Borel determinacy},
Annals of Mathematics (2) \textbf{102} (1975), no.~2, 363--371.

\bibitem{MaschioTrotta26}
S.~Maschio and D.~Trotta,
\emph{A topos for extended Weihrauch degrees},
Annals of Pure and Applied Logic (2026), \href{https://arxiv.org/abs/2505.08697}{\texttt{arXiv:2505.08697}}.

\bibitem{Masuda26}
N.~Masuda,
\emph{The algebra of categorical spectra},
\href{https://arxiv.org/abs/2605.03114v2}{\texttt{arXiv:2605.03114v2}}, 2026.

\bibitem{Mathias01}
A.~R.~D.~Mathias,
\emph{The strength of Mac~Lane set theory},
Annals of Pure and Applied Logic \textbf{110} (2001), no.~1--3, 107--234.

\bibitem{Mitchell72}
W.~J.~Mitchell,
\emph{Boolean topoi and the theory of sets},
Journal of Pure and Applied Algebra \textbf{2} (1972), no.~4, 261--274.

\bibitem{MoerdijkReyes91}
I.~Moerdijk and G.~E.~Reyes,
\emph{Models for Smooth Infinitesimal Analysis},
Springer, New York, 1991.

\bibitem{MontalbanShore12}
A.~Montalb\'an and R.~A.~Shore,
\emph{The limits of determinacy in second-order arithmetic},
Proceedings of the London Mathematical Society (3) \textbf{104} (2012), no.~2, 223--252.

\bibitem{MulveyPelletier91}
C.~J.~Mulvey and J.~W.~Pelletier,
\emph{A globalization of the Hahn--Banach theorem},
Advances in Mathematics \textbf{89} (1991), no.~1, 1--59.

\bibitem{Osius74}
G.~Osius,
\emph{Categorical set theory: a characterization of the category of sets},
Journal of Pure and Applied Algebra \textbf{4} (1974), no.~1, 79--119.

\bibitem{ORW26}
V.~Ozornova, M.~Rovelli, and T.~Walde,
\emph{Cores and localizations of $(\infty, \infty)$-categories} (\href{https://arxiv.org/abs/2603.11005v2}{\texttt{arXiv:2603.11005v2}}),
\emph{About the contractibility of the walking coinductive equivalence} (\href{https://arxiv.org/abs/2608.28355}{\texttt{arXiv:2608.28355}}), and
\emph{Models of $n$-dimensional categories as $n \to \infty$} (\href{https://arxiv.org/abs/2610.38698}{\texttt{arXiv:2610.38698}}), 2026.

\bibitem{Quillen71}
D.~Quillen,
\emph{Cohomology of groups},
in \emph{Actes du Congr\`es International des Math\'ematiciens} (Nice, 1970), Tome~2, Gauthier-Villars, Paris, 1971, pp.~47--51.

\bibitem{Rieger57}
L.~Rieger,
\emph{A contribution to G\"odel's axiomatic set theory, I},
Czechoslovak Mathematical Journal \textbf{7 (82)} (1957), no.~3, 323--357.

\bibitem{Schreiber13}
U.~Schreiber,
\emph{Differential cohomology in a cohesive $\infty$-topos},
\href{https://arxiv.org/abs/1310.7930}{\texttt{arXiv:1310.7930}}, 2013.

\bibitem{Serafin23}
L.~Serafin,
\emph{Ultrafilters, transversals, and the hat game},
preprint (2023), \href{https://arxiv.org/abs/2311.07886}{\texttt{arXiv:2311.07886}}.

\bibitem{Shulman08}
M.~Shulman,
\emph{Framed bicategories and monoidal fibrations},
Theory and Applications of Categories \textbf{20} (2008), no.~18, 650--738.

\bibitem{Shulman20}
M.~Shulman,
\emph{The 2-Chu-Dialectica construction and the polycategory of multivariable adjunctions},
Theory and Applications of Categories \textbf{35} (2020), no.~4, 89--136.

\bibitem{SimonsSullivan08}
J.~Simons and D.~Sullivan,
\emph{Axiomatic characterization of ordinary differential cohomology},
Journal of Topology \textbf{1} (2008), no.~1, 45--56 (\href{https://arxiv.org/abs/math/0701077}{\texttt{arXiv:math/0701077}}).

\bibitem{Simpson09}
S.~G.~Simpson,
\emph{Subsystems of Second Order Arithmetic},
2nd ed., Perspectives in Logic, Cambridge University Press, Cambridge, 2009.

\bibitem{Skolem55}
T.~Skolem,
\emph{A critical remark on foundational research} (1955),
in \emph{Selected Works in Logic}, Universitetsforlaget, Oslo, 1970, pp.~581--586.

\bibitem{vanOosten08}
J.~van~Oosten,
\emph{Realizability: An Introduction to its Categorical Side},
Studies in Logic and the Foundations of Mathematics, vol.~152, Elsevier, Amsterdam, 2008.

\bibitem{Williams18}
K.~J.~Williams,
\emph{The Structure of Models of Second-Order Set Theories},
Ph.D.\ Dissertation, CUNY Graduate Center, New York, 2018.

\end{thebibliography}

\printindex

\end{document}
